\documentclass{report}


% GEMINI

% ==========================================
% FONTS & ENCODING
% ==========================================
\usepackage[utf8]{inputenc}
\usepackage[T1]{fontenc}
% Body font. Comment the next line out to fall back to plain Computer Modern.
\usepackage{lmodern}          % Latin Modern -- has a full sans family
\usepackage{anyfontsize}
\renewcommand{\sfdefault}{lmss}
% --- Body font: switch to Palatino by uncommenting the next line (single-line change) ---
% \usepackage{mathpazo}

\usepackage[english]{babel}
\usepackage{csquotes}

% --- HYPHENATION / LINE-BREAKING TUNING ---
\hyphenpenalty=100
\exhyphenpenalty=100
\doublehyphendemerits=10000
\finalhyphendemerits=5000
\tolerance=1500
\emergencystretch=3em
\hfuzz=0.25ex

% ==========================================
% CORE PACKAGES
% ==========================================
\usepackage{amsmath}
\usepackage{amssymb}
\usepackage{amsthm}
\usepackage{aliascnt}   % lets lemma/corollary/... share the theorem counter
                        % while keeping distinct names for cleveref
\usepackage{mathtools}
\usepackage{etoolbox}
\usepackage[dvipsnames, x11names]{xcolor}
\usepackage[top=15mm, bottom=15mm, includehead, includefoot,
            headheight=9mm, headsep=3.5mm, footskip=9mm,
            nomarginpar]{geometry} % try 25 mm

% change enumerations:
\usepackage{enumitem, tabstackengine}
\usepackage{multirow}
\usepackage{booktabs}

% use tikz:
\usepackage{tikz}
\usepackage{tikz-cd}
\usepackage{tikz-3dplot}
\usepackage{subcaption}
\usetikzlibrary{matrix}
\usetikzlibrary{positioning}
\usetikzlibrary{arrows.meta}
\usetikzlibrary{shapes}
\usetikzlibrary{shadows}
\usetikzlibrary{babel, calc, patterns}
\usepackage{tcolorbox}
\tcbuselibrary{skins, breakable}
\usepackage{fontawesome5}

% change page-header/footer:
\usepackage{fancyhdr}
\usepackage{extramarks}
\usepackage{lastpage}

% ==========================================
% GLOBAL COLOR PALETTE
% ==========================================
\definecolor{ThemePurple}{HTML}{6A3199}
\colorlet{ThemePrimary}{ThemePurple}
\colorlet{ThemeSecondary}{MidnightBlue}
\colorlet{ThemeWeekNumber}{OliveGreen}
\colorlet{ThemeAccent1}{BurntOrange}

% --- Theorem environments ---
\colorlet{EnvTheorem}{ThemePurple}
\definecolor{GoldOrange}{HTML}{C26B00}
\colorlet{EnvAINote}{GoldOrange}

% --- Proofs ---
\definecolor{ThemeGreen}{HTML}{008837}
% \colorlet{ProofTint}{OliveGreen}
\colorlet{ProofTint}{ThemeGreen}
\colorlet{ProofQEDText}{ProofTint}
\colorlet{ProofQEDBg}{white}

% --- Running text ---
% \colorlet{TextEmphColor}{MidnightBlue}
\definecolor{EmphBlue}{HTML}{008CC8}
\colorlet{TextEmphColor}{EmphBlue}
% \colorlet{TextBoldColor}{Brown!80!black}   % Espresso brown
\definecolor{TextBoldColor}{HTML}{A8651E}    % Warm Bronze Brown
\colorlet{TextMetaNote}{Gray}
\colorlet{LinkColor}{MidnightBlue}
\colorlet{FileColor}{Magenta}

% --- Lists ---
% \colorlet{ItemizeBullet}{BurntOrange}
% \colorlet{EnumerateBullet}{MidnightBlue}
\definecolor{ListAmber}{HTML}{BD2400}
\colorlet{ItemizeBullet}{ListAmber}
\colorlet{EnumerateBullet}{ListAmber}

% --- Headings ---
\colorlet{PartTitleText}{ThemePurple!95!black}
\colorlet{PartNumberText}{OliveGreen}
\colorlet{ChapTitleText}{MidnightBlue!90!black}
\colorlet{ChapNumberText}{OliveGreen}
\colorlet{ChapLine}{MidnightBlue!80!black}
\colorlet{SecTitleColor}{MidnightBlue}
\colorlet{SecNumberColor}{OliveGreen}
\colorlet{SubSecTitleColor}{MidnightBlue}
\colorlet{SubSubSecTitleColor}{TextBoldColor}
\colorlet{HeaderFooterFaint}{OliveGreen!45!Gray}
\colorlet{HeaderFooterLine}{HeaderFooterFaint!50}

% --- TOC entry colours ---
% These must be *named* colours (not `A!95!B' expressions), because they are
% handed to hyperref's linkcolor as well -- otherwise the link colour would
% override the entry colour and the whole TOC would come out midnight blue.
\colorlet{TocPartColor}{ThemePurple!85!black}
\colorlet{TocChapColor}{ThemePurple!95!black}
\colorlet{TocSecColor}{TextEmphColor}
\colorlet{TocSubsecColor}{TextBoldColor}
\colorlet{TocSubsubColor}{gray}
\colorlet{TocLabelColor}{ThemePurple!85!black}

% --- EMPHASIS STYLE OVERRIDE ---
\DeclareTextFontCommand{\emph}{\sffamily\bfseries\color{TextEmphColor}}
\RenewDocumentCommand{\textbf}{m}{{\sffamily\bfseries\color{TextBoldColor}#1}}
% NB: no \setlist[enumerate] here. There used to be one setting `font=...', but the fuller
% \setlist[enumerate] in the LIST FORMATTING block below sets `font=' again and, enumitem
% keys being last-wins, silently overrode it. It was dead configuration that looked live.
% The surviving declaration is the one place to change enumerate styling.

% ==========================================
% HYPERLINKS & CLEVEREF  (hyperref must come before cleveref)
% ==========================================
\usepackage{hyperref}
\hypersetup{
  colorlinks=true,
  linkcolor=LinkColor,
  filecolor=FileColor,
  urlcolor=LinkColor,
  citecolor=LinkColor,
  bookmarksopen=true
}
\usepackage{bookmark}   % tolerates skipped heading levels (e.g. \chapter -> \subsection)
\usepackage[nameinlink, capitalize]{cleveref}

% ==========================================
% TYPOGRAPHY & SPACING UNITS
% ==========================================
\linespread{1.05}

\newlength{\customparskip}
\setlength{\customparskip}{3pt plus 1.5pt minus 0.5pt}

\setlength{\parindent}{0pt}
\setlength{\parskip}{\customparskip}

\newlength{\customenvspace}
\setlength{\customenvspace}{2.0ex plus 0.5ex minus 0.2ex}

% Inline math breathing room
\setlength{\mathsurround}{1pt}

% ...but NOT inside TikZ. A node's boundary is derived from the width of the box
% it contains, and \mathsurround pads every math box on both sides. TikZ then
% anchors arrows on the padded boundary while the glyphs sit elsewhere, so
% tikz-cd arrows overshoot their targets and nodes collide. Must be 0 in there.
\AtBeginEnvironment{tikzpicture}{\setlength{\mathsurround}{0pt}}
\AtBeginEnvironment{tikzcd}{\setlength{\mathsurround}{0pt}}

% Matrix and table rows
\renewcommand{\arraystretch}{1.15}

% Proportional math display spacing
\AtBeginDocument{
  \setlength{\abovedisplayskip}{1.2\customparskip plus 0.5ex minus 0.2ex}
  \setlength{\belowdisplayskip}{1.2\customparskip plus 0.5ex minus 0.2ex}
  \setlength{\abovedisplayshortskip}{0.5\customparskip plus 0.2ex}
  \setlength{\belowdisplayshortskip}{0.8\customparskip plus 0.2ex minus 0.1ex}
}

% RETIRED 2026-08-09: \omitted{...}, a red-text macro from a dead error-flagging convention.
% It was marked "kept for future use if needed" and had zero uses in content/. Errors are
% flagged in an `ainote`, which is the whole system; see style.md.

% Custom commands for italicized terminology and "quotes".
% \newterm    -> English quotes ``...''   (the default; this is an English text)
% \germanterm -> German quotes  \glqq...\grqq  for a German word MENTIONED in a sentence,
%                where the word itself is the subject: \germanterm{Abschluss} is the closure.
% \germanfor  -> the MIRRORING form, for giving the German of an English term just introduced:
%                a \newterm{compact} set \germanfor{kompakte Menge}
%                prints  a "compact" set (German: \glqq kompakte Menge\grqq).
%                It carries its own parentheses, so do NOT wrap it in a literal pair --- that
%                was the old \germanterm convention and it is what this macro replaces.
\definecolor{QuoteRed}{HTML}{BD2400}
\newcommand{\newterm}[1]{\textcolor{QuoteRed}{``}\textcolor{TextBoldColor}{\textit{#1}}\textcolor{QuoteRed}{''}}
\newcommand{\germanterm}[1]{\textcolor{QuoteRed}{\glqq}\textcolor{TextBoldColor}{\textit{#1}}\textcolor{QuoteRed}{\grqq}}
\newcommand{\germanfor}[1]{(German:\ \germanterm{#1})}
\newcommand{\qt}[1]{\textit{\textcolor{QuoteRed}{``}#1\textcolor{QuoteRed}{''}}}
\newcommand{\inlinemetanote}[1]{\textit{[\textcolor{TextMetaNote}{#1}]}}
\newcommand{\hint}[1]{\textcolor{OliveGreen}{\textbf{Hint:}} #1}

% Redefine \bigbreak to ensure the following paragraph is not indented
\let\oldbigbreak\bigbreak
\renewcommand{\bigbreak}{\oldbigbreak\noindent}

% ==========================================
% HEADING FORMATTING & SPACING
% ==========================================
\usepackage[explicit]{titlesec}
\usepackage{titletoc}

% Headings hang into the left margin, staggered by level.
\newlength{\chapouthang}\setlength{\chapouthang}{-3.2ex}
\newlength{\secouthang}\setlength{\secouthang}{-2.4ex}
\newlength{\subsecouthang}\setlength{\subsecouthang}{-1.6ex}
\newlength{\subsubsecouthang}\setlength{\subsubsecouthang}{-0.8ex}


% 0. Custom Part Format: a title page carrying "PART III" over the part name.
% Parts group the ~25 topic chapters into the seven thematic blocks of the course;
% they do NOT reset \thechapter, so chapter (and therefore theorem) numbering runs
% continuously across the whole document -- which is what \thetheorem below assumes.
\titleformat{name=\part}[display]
{\normalfont\sffamily\bfseries\filcenter}
{\fontsize{15pt}{18pt}\selectfont\mdseries\color{PartNumberText}\MakeUppercase{\partname\ \thepart}}
{2.4ex}
{\fontsize{26pt}{32pt}\selectfont\color{PartTitleText}#1}
[{\vspace{1.5ex}{\color{ChapLine}\rule{0.5\textwidth}{1pt}}}]

\titlespacing*{\part}{0pt}{*4}{*3}

% Since \chapter exists in this class, titlesec's compat layer classifies \part as
% its "page" class (\titleclass{\part}{page}), whose implementation \ttl@page@ii
% issues a \newpage right after the title block -- that ejected a fresh page before
% any following text could appear, swallowing the part-intro prose paragraphs below
% onto their own near-empty page. Drop that forced break so the prose shares the
% title page; \chapter's own \clearpage still starts the next chapter fresh.
% (Document is oneside/non-openright, so the twoside/openright blank-page branch
% a few lines later in \ttl@page@ii never fires and is unaffected.)
\makeatletter
\patchcmd{\ttl@page@ii}{\newpage}{}{}{\ClassWarning{report}{Could not patch \string\ttl@page@ii}}
\makeatother

% 1. Custom Chapter Format: Title (chapter X)
\titleformat{name=\chapter}[block]  % [block] forces it onto a single line
{\normalfont\sffamily\Large\bfseries\color{ChapTitleText}}  % Main style
{}                           % Empty standard label
{0ex}                        % Separation
% #1 is the title. Suffix is \large (smaller) and un-bolded
{#1 {\large\mdseries\color{ChapNumberText}(\chaptername\ \thechapter)}}
[{\color{ChapLine}\rule{\textwidth}{0.8pt}}]

\titleformat{name=\chapter,numberless}[block]
{\normalfont\sffamily\Large\bfseries\color{ChapTitleText}}
{}
{0ex}
{#1}
[{\color{ChapLine}\rule{\textwidth}{0.8pt}}]

% 2. Space-efficient Chapter Headings
\titlespacing*{\chapter}{\chapouthang}{-7ex}{2.2ex}

\setcounter{secnumdepth}{3}
\setcounter{tocdepth}{3}

\renewcommand{\thesection}{\thechapter.\alph{section}}
\renewcommand{\thesubsection}{\thesection.\arabic{subsection}}
\renewcommand{\thesubsubsection}{\thesubsection.\arabic{subsubsection}}

% 3. Custom Section Format: Title (Section X.a)
\titleformat{name=\section}
{\normalfont\sffamily\Large\bfseries\color{SecTitleColor}}
{}
{0ex}
{#1 {\large\mdseries\color{SecNumberColor}(Section \thesection)}}

\titleformat{name=\section,numberless}
{\normalfont\sffamily\Large\bfseries\color{SecTitleColor}}
{}
{0ex}
{#1}

% 4. Custom Subsection Format: Title (Subsection X.a.Y)
\titleformat{name=\subsection}
{\normalfont\sffamily\large\bfseries\color{SubSecTitleColor}}
{}
{0ex}
{#1 {\normalsize\mdseries\color{SecNumberColor}(Subsection \thesubsection)}}

\titleformat{name=\subsection,numberless}
{\normalfont\sffamily\large\bfseries\color{SubSecTitleColor}}
{}
{0ex}
{#1}

% 5. Custom Subsubsection Format: Title (Subsubsection X.a.Y.Z)
\titleformat{name=\subsubsection}
{\normalfont\sffamily\normalsize\bfseries\color{SubSubSecTitleColor}}
{}
{0ex}
{#1 {\small\mdseries\color{SecNumberColor}(Subsubsection \thesubsubsection)}}

\titleformat{name=\subsubsection,numberless}
{\normalfont\sffamily\normalsize\bfseries\color{SubSubSecTitleColor}}
{}
{0ex}
{#1}

% 4. Space-efficient Heading Spacing
\titlespacing*{\section}{\secouthang}{3.2ex plus 0.8ex minus .2ex}{1.3ex plus .2ex}
\titlespacing*{\subsection}{\subsecouthang}{2.7ex plus 0.8ex minus .2ex}{1.1ex plus .2ex}
\titlespacing*{\subsubsection}{\subsubsecouthang}{2.2ex plus 0.8ex minus .2ex}{0.9ex plus .2ex}

% ==========================================
% TABLE OF CONTENTS TYPOGRAPHY & COLORS
% ==========================================
% Width reserved for the page number at the right edge.
\contentsmargin{4.7ex}

% Staggered indentation per level; dotted leaders in the faint header colour.
% Parts carry no page number and no dotted leader: they head a block of chapters
% rather than pointing at a page worth turning to.
\titlecontents{part}[0ex]
  {\addvspace{5.2ex}\hypersetup{linkcolor=TocPartColor}\sffamily\bfseries\Large\color{TocPartColor}}
  {\color{PartNumberText}\MakeUppercase{\partname}\ \thecontentslabel\hspace{1.9ex}\color{TocPartColor}}
  {}
  {}

\titlecontents{chapter}[0ex]
  {\addvspace{3.7ex}\hypersetup{linkcolor=TocChapColor}\sffamily\bfseries\large\color{TocChapColor}}
  {\color{ChapNumberText}\MakeUppercase{\chaptername}\ \thecontentslabel\hspace{1.9ex}\color{TocChapColor}}
  {}
  {\color{HeaderFooterLine}\titlerule*[0.6pc]{.}\color{ChapNumberText}\contentspage}

\titlecontents{section}[4.2ex]
  {\addvspace{1.2ex}\hypersetup{linkcolor=TocSecColor}\sffamily\bfseries\color{TocSecColor}}
  {\color{TocLabelColor}\thecontentslabel\hspace{1.9ex}\color{TocSecColor}}
  {}
  {\color{HeaderFooterLine}\titlerule*[0.6pc]{.}\color{TocSecColor}\contentspage}

\titlecontents{subsection}[9.3ex]
  {\addvspace{0.6ex}\hypersetup{linkcolor=TocSubsecColor}\sffamily\color{TocSubsecColor}}
  {\color{TocSubsecColor}\thecontentslabel\hspace{1.9ex}\color{TocSubsecColor}}
  {}
  {\color{HeaderFooterLine}\titlerule*[0.6pc]{.}\color{TocSubsecColor}\contentspage}

\titlecontents{subsubsection}[15.1ex]
  {\addvspace{0.35ex}\hypersetup{linkcolor=TocSubsubColor}\sffamily\small\bfseries\color{TocSubsubColor}}
  {\color{TocSubsubColor}\thecontentslabel\hspace{1.9ex}\color{TocSubsubColor}}
  {}
  {\normalfont\color{HeaderFooterLine}\titlerule*[0.6pc]{.}\color{TocSubsubColor}\contentspage}

% ==========================================
% LIST FORMATTING
% ==========================================
\setlist{
  topsep=0pt, 
  partopsep=0pt, 
  parsep=0.5\customparskip, 
  itemsep=0pt, 
}

\setlist[itemize]{
  label=\textcolor{ItemizeBullet}{\textbullet}
}

\setlist[enumerate]{
  font=\colorlet{TextBoldColor}{EnumerateBullet}, 
  label=\textcolor{EnumerateBullet}{\bfseries(\alph*)}, 
  ref=(\alph*)
}

% No \setlist[description]: the `description' environment is forbidden by this project's
% house style (use itemize with the name bolded inline instead), and nothing in content/
% uses it. Configuring an environment that must never appear only invited someone to
% reach for it.

% Abstände bei geschachtelten Listen 
\setlist[itemize,2]{topsep=0pt, itemsep=0pt, parsep=0pt}
\setlist[enumerate,2]{topsep=0pt, itemsep=0pt, parsep=0pt}

% --- FIX: LISTEN DIREKT NACH THEOREM-HEADING ---
\makeatletter
\providecommand{\fixlistaftertheorem}{%
  \if@inlabel 
    \leavevmode \par 
  \fi
}
\BeforeBeginEnvironment{itemize}{\fixlistaftertheorem}
\BeforeBeginEnvironment{enumerate}{\fixlistaftertheorem}
\BeforeBeginEnvironment{description}{\fixlistaftertheorem}
\makeatother

% ==========================================
% THEOREM CONFIGURATION
% ==========================================
% THEOREM CONFIGURATION
% ==========================================
% 1. Theorems, Lemmata, Propositions, Corollaries, Claims (\faGem)
\newtheoremstyle{purplethm}%
{\customenvspace}{\customenvspace}{\normalfont}{}{\color{EnvTheorem}\sffamily\bfseries}{:}{0.5em}%
{\textcolor{EnvTheorem}{\faGem}\hspace{0.4em}%
 \thmname{\textcolor{EnvTheorem}{\sffamily\bfseries#1}}\thmnumber{\sffamily\bfseries\ #2}%
 \thmnote{\ \textcolor{EnvTheorem}{\bfseries(}{\normalfont\sffamily\itshape#3}\textcolor{EnvTheorem}{\bfseries)}}}

% 2. Definitions (\faBookOpen)
\newtheoremstyle{defthm}%
{\customenvspace}{\customenvspace}{\normalfont}{}{\color{EnvTheorem}\sffamily\bfseries}{:}{0.5em}%
{\textcolor{EnvTheorem}{\faBookOpen}\hspace{0.4em}%
 \thmname{\textcolor{EnvTheorem}{\sffamily\bfseries#1}}\thmnumber{\sffamily\bfseries\ #2}%
 \thmnote{\ \textcolor{EnvTheorem}{\bfseries(}{\normalfont\sffamily\itshape#3}\textcolor{EnvTheorem}{\bfseries)}}}

% 3. Questions & Answers (\faQuestionCircle)
\newtheoremstyle{questionthm}%
{\customenvspace}{\customenvspace}{\normalfont}{}{\color{MidnightBlue}\sffamily\bfseries}{:}{0.5em}%
{\textcolor{MidnightBlue}{\faQuestionCircle}\hspace{0.4em}%
 \thmname{\textcolor{MidnightBlue}{\sffamily\bfseries#1}}\thmnumber{\sffamily\bfseries\ #2}%
 \thmnote{\ \textcolor{MidnightBlue}{\bfseries(}{\normalfont\sffamily\itshape#3}\textcolor{MidnightBlue}{\bfseries)}}}

% 4. Remarks & Notes (\faInfoCircle)
\newtheoremstyle{remarkthm}%
{\customenvspace}{\customenvspace}{\normalfont}{}{\color{ThemeGreen}\sffamily\bfseries}{:}{0.5em}%
{\textcolor{ThemeGreen}{\faInfoCircle}\hspace{0.4em}%
 \thmname{\textcolor{ThemeGreen}{\sffamily\bfseries#1}}\thmnumber{\sffamily\bfseries\ #2}%
 \thmnote{\ \textcolor{ThemeGreen}{\bfseries(}{\normalfont\sffamily\itshape#3}\textcolor{ThemeGreen}{\bfseries)}}}

% 5. Examples & Exercises (\faLightbulb)
\newtheoremstyle{examplethm}%
{\customenvspace}{\customenvspace}{\normalfont}{}{\color{EnvTheorem}\sffamily\bfseries}{:}{0.5em}%
{\textcolor{EnvTheorem}{\faLightbulb}\hspace{0.4em}%
 \thmname{\textcolor{EnvTheorem}{\sffamily\bfseries#1}}\thmnumber{\sffamily\bfseries\ #2}%
 \thmnote{\ \textcolor{EnvTheorem}{\bfseries(}{\normalfont\sffamily\itshape#3}\textcolor{EnvTheorem}{\bfseries)}}}

\theoremstyle{purplethm}

\newtheorem{theorem}{Theorem}

\renewcommand{\thetheorem}{%
  \thechapter.\ifnum\value{section}>0 \alph{section}.\fi\arabic{theorem}}

\makeatletter
\AtBeginDocument{%
  \renewcommand{\theHtheorem}{\thechapter.\alph{section}.\arabic{theorem}}%
  \renewcommand{\theHequation}{\thechapter.\alph{section}.\arabic{equation}}%
  \@namedef{theHlemma}{\theHtheorem}%
  \@namedef{theHcorollary}{\theHtheorem}%
  \@namedef{theHdefinition}{\theHtheorem}%
  \@namedef{theHproposition}{\theHtheorem}%
  \@namedef{theHexercise}{\theHtheorem}%
  \@namedef{theHclaim}{\theHtheorem}%
  \@namedef{theHnotation}{\theHtheorem}%
  \@namedef{theHremark}{\theHtheorem}%
  \@namedef{theHexample}{\theHtheorem}%
  \@namedef{theHsummary}{\theHtheorem}%
  \@namedef{theHwarmup}{\theHtheorem}%
  \@namedef{theHquestion}{\theHtheorem}%
  \@namedef{theHanswer}{\theHtheorem}%
  \@namedef{theHimportantremark}{\theHtheorem}%
  \@namedef{theHgoals}{\theHtheorem}%
  \@namedef{theHconclusion}{\theHtheorem}%
  \@namedef{theHaiexample}{\theHtheorem}%
  \@namedef{theHaiexercise}{\theHtheorem}%
}
\makeatother

\makeatletter
\@addtoreset{theorem}{chapter}
\makeatother

\setcounter{MaxMatrixCols}{20}

% --- NUMBERED ENVIRONMENTS ---
\theoremstyle{purplethm}
\newaliascnt{lemma}{theorem}
\newtheorem{lemma}[lemma]{Lemma}
\aliascntresetthe{lemma}

\newaliascnt{corollary}{theorem}
\newtheorem{corollary}[corollary]{Corollary}
\aliascntresetthe{corollary}

\newaliascnt{proposition}{theorem}
\newtheorem{proposition}[proposition]{Proposition}
\aliascntresetthe{proposition}

\newaliascnt{claim}{theorem}
\newtheorem{claim}[claim]{Claim}
\aliascntresetthe{claim}

\theoremstyle{defthm}
\newaliascnt{definition}{theorem}
\newtheorem{definition}[definition]{Definition}
\aliascntresetthe{definition}

\theoremstyle{examplethm}
\newaliascnt{exercise}{theorem}
\newtheorem{exercise}[exercise]{Exercise}
\aliascntresetthe{exercise}

\newaliascnt{example}{theorem}
\newtheorem{example}[example]{Example}
\aliascntresetthe{example}

\newaliascnt{warmup}{theorem}
\newtheorem{warmup}[warmup]{Warm up}
\aliascntresetthe{warmup}

\newaliascnt{aiexample}{theorem}
\newtheorem{aiexample}[aiexample]{AI-Example}
\aliascntresetthe{aiexample}

\newaliascnt{aiexercise}{theorem}
\newtheorem{aiexercise}[aiexercise]{AI-Exercise}
\aliascntresetthe{aiexercise}

\theoremstyle{questionthm}
\newaliascnt{question}{theorem}
\newtheorem{question}[question]{Question}
\aliascntresetthe{question}

\newaliascnt{answer}{theorem}
\newtheorem{answer}[answer]{Answer}
\aliascntresetthe{answer}

\theoremstyle{remarkthm}
\newaliascnt{notation}{theorem}
\newtheorem{notation}[notation]{Notation}
\aliascntresetthe{notation}

\newaliascnt{remark}{theorem}
\newtheorem{remark}[remark]{Remark}
\aliascntresetthe{remark}

\newaliascnt{summary}{theorem}
\newtheorem{summary}[summary]{Summary}
\aliascntresetthe{summary}

\newaliascnt{importantremark}{theorem}
\newtheorem{importantremark}[importantremark]{Important remark}
\aliascntresetthe{importantremark}

\newaliascnt{goals}{theorem}
\newtheorem{goals}[goals]{Goals}
\aliascntresetthe{goals}

\newaliascnt{conclusion}{theorem}
\newtheorem{conclusion}[conclusion]{Conclusion}
\aliascntresetthe{conclusion}

\newtheoremstyle{ainotethm}%
{\customenvspace}{\customenvspace}{\normalfont}{}{\color{EnvAINote}\sffamily\bfseries}{:}{0.5em}%
{\textcolor{EnvAINote}{\faRobot}\hspace{0.4em}%
 \thmname{\textcolor{EnvAINote}{\sffamily\bfseries#1}}\thmnumber{\sffamily\bfseries\ #2}%
 \thmnote{\ \textcolor{EnvAINote}{\bfseries(}{\normalfont\sffamily\itshape#3}\textcolor{EnvAINote}{\bfseries)}}}

\theoremstyle{ainotethm}
\newtheorem{ainote}{AI-Note}[chapter]
\theoremstyle{purplethm}

% ==========================================
% CLEVEREF NAMES
% ==========================================
\crefname{theorem}{Theorem}{Theorems}
\Crefname{theorem}{Theorem}{Theorems}
\crefname{lemma}{Lemma}{Lemmata}
\Crefname{lemma}{Lemma}{Lemmata}
\crefname{corollary}{Corollary}{Corollaries}
\Crefname{corollary}{Corollary}{Corollaries}
\crefname{definition}{Definition}{Definitions}
\Crefname{definition}{Definition}{Definitions}
\crefname{proposition}{Proposition}{Propositions}
\Crefname{proposition}{Proposition}{Propositions}
\crefname{exercise}{Exercise}{Exercises}
\Crefname{exercise}{Exercise}{Exercises}
\crefname{part}{Part}{Parts}
\Crefname{part}{Part}{Parts}
\crefname{chapter}{Chapter}{Chapters}
\Crefname{chapter}{Chapter}{Chapters}
\crefname{section}{Section}{Sections}
\Crefname{section}{Section}{Sections}
\crefname{equation}{Equation}{Equations}
\Crefname{equation}{Equation}{Equations}
\crefname{figure}{Figure}{Figures}
\Crefname{figure}{Figure}{Figures}
\crefname{table}{Table}{Tables}
\Crefname{table}{Table}{Tables}
\crefname{item}{Part}{Parts}
\Crefname{item}{Part}{Parts}

% Sonnet 5 (Medium): cleveref names for the newly-numbered environments above.
\crefname{claim}{Claim}{Claims}
\Crefname{claim}{Claim}{Claims}
\crefname{notation}{Notation}{Notations}
\Crefname{notation}{Notation}{Notations}
\crefname{remark}{Remark}{Remarks}
\Crefname{remark}{Remark}{Remarks}
\crefname{example}{Example}{Examples}
\Crefname{example}{Example}{Examples}
\crefname{summary}{Summary}{Summaries}
\Crefname{summary}{Summary}{Summaries}
\crefname{warmup}{Warm up}{Warm ups}
\Crefname{warmup}{Warm up}{Warm ups}
\crefname{question}{Question}{Questions}
\Crefname{question}{Question}{Questions}
\crefname{answer}{Answer}{Answers}
\Crefname{answer}{Answer}{Answers}
\crefname{importantremark}{Important remark}{Important remarks}
\Crefname{importantremark}{Important remark}{Important remarks}
\crefname{goals}{Goals}{Goals}
\Crefname{goals}{Goals}{Goals}
\crefname{conclusion}{Conclusion}{Conclusions}
\Crefname{conclusion}{Conclusion}{Conclusions}
\crefname{ainote}{AI-Note}{AI-Notes}
\Crefname{ainote}{AI-Note}{AI-Notes}
\crefname{aiexample}{AI-Example}{AI-Examples}
\Crefname{aiexample}{AI-Example}{AI-Examples}
\crefname{aiexercise}{AI-Exercise}{AI-Exercises}
\Crefname{aiexercise}{AI-Exercise}{AI-Exercises}

% ==========================================
% PROOF ENVIRONMENT (with pen nib icon & shadowed Q.E.D. box)
% ==========================================
\newcommand{\qedbox}{%
  \tikz[baseline=(Q.base)]\node[draw=ProofTint, fill=ProofQEDBg, line width=1pt, rounded corners=2pt, inner xsep=4pt, inner ysep=2pt, font=\sffamily\bfseries\small, text=ProofQEDText, drop shadow={shadow xshift=0.5ex, shadow yshift=-0.5ex, fill=ProofTint, opacity=0.4}] (Q) {Q.E.D.};%
}
\renewcommand{\qedsymbol}{\qedbox}

\makeatletter
\renewenvironment{proof}[1][\proofname]{%
  \par\if@nobreak\vspace{-1ex}\else\addvspace{2.2ex plus 0.5ex minus 0.2ex}\fi
  \noindent\textcolor{ProofTint}{\faPenNib}\hspace{0.5em}%
  {\normalfont\sffamily\bfseries\color{ProofTint}#1:}%
  \par\nobreak\vspace{0.8ex plus .1ex}\vspace{-\parskip}
  \@afterindentfalse\@afterheading
}{%
  \unskip\nobreak\hfill\penalty50\hskip2em\null\nobreak\hfill
  \qedbox
  {\parfillskip=0pt\par}\vspace{2ex}
}
\makeatother

\newenvironment{exercisesolution}[1][Solution]{%
  \begin{proof}[#1]%
}{%
  \end{proof}%
}

% ==========================================
% RETIRED: the week-skeleton macros
% ==========================================
% This document used to be organised one chapter per teaching week, and carried three
% macros to support that: \session{Monday/Friday} (a styled day marker), \exercisesheet{N}
% (an unnumbered heading opening the week's problem sheet), and \continuedfrom{label}
% (a back-pointer for a topic split across two weeks).
%
% All three are gone, together with the structure that needed them. Chapters are topics
% now, so there is no day to mark; the problem sheets are dissolved, each problem sitting
% with the material it tests (see content/appendix-d-problem-sheets.tex, which restores the
% sheet view); and a topic no longer gets cut in half by a week boundary, so nothing needs
% to point backwards. The \continuedfrom on compactness is the case that made the argument:
% \section{Compactness --- continued} and its pointer both vanished the moment weeks 2 and 3
% were merged into \cref{ch:compactness}.
%
% Do not reinstate these without reinstating the week structure -- see gemini.md.

% RETIRED 2026-08-09: \begin{figurestub}{FIG-W02-01} ... \end{figurestub}, a grey placeholder
% box that carried a drawing spec until the figure behind it was drawn. Every figure is now
% drawn and render-checked, the environment had zero uses left in content/, and leaving a
% placeholder mechanism defined invites new placeholders. Do not reinstate it: a figure that
% cannot be drawn yet does not go into the document at all.

% --- MATH OPERATORS ---
\newcommand{\M}{\mathcal{M}}
\DeclareMathOperator{\ColS}{Cols}
\DeclareMathOperator{\RowS}{Rows}
\DeclareMathOperator{\Eig}{Eig}
\DeclareMathOperator{\End}{End}
\DeclareMathOperator{\Tr}{Tr}
\DeclareMathOperator{\Sp}{Sp}
\DeclareMathOperator{\rank}{rank}
\DeclareMathOperator{\sgn}{sgn}
\DeclareMathOperator{\Hom}{Hom}
\DeclareMathOperator{\id}{id}
\DeclareMathOperator{\GL}{GL}
\DeclareMathOperator{\Fib}{Fib}
\DeclareMathOperator{\Img}{Im}
\DeclareMathOperator{\Sol}{Sol}
\DeclareMathOperator{\Spec}{Spec}
\DeclareMathOperator{\blockdiag}{blockdiag}
\DeclareMathOperator{\adj}{adj}
\DeclareMathOperator{\diag}{diag}
\DeclareMathOperator{\dist}{dist}
\DeclareMathOperator{\Rea}{Re}   % real part; \Re is taken by the fraktur symbol
\DeclareMathOperator{\diam}{diam}
\DeclareMathOperator{\supp}{supp}
\DeclareMathOperator{\vol}{vol}
\DeclareMathOperator{\divg}{div}
\DeclareMathOperator{\curl}{curl}
\DeclareMathOperator{\Jac}{J}
\DeclareMathOperator{\Hess}{Hess}
\DeclareMathOperator{\grad}{grad}

% Matrix groups. Declared as operators so that "O", "U", "SO", "SU" stay upright
% and are never mistaken for the italic variables O, U (which this text also uses
% -- e.g. "let $O \in \Orth(n)$").
\DeclareMathOperator{\Orth}{O}
\DeclareMathOperator{\Unit}{U}
\DeclareMathOperator{\SOrth}{SO}
\DeclareMathOperator{\SUnit}{SU}

% ==========================================
% TRANSPOSE  --  one switch for the whole document
% ==========================================
% Usage: A\transp, (PAQ)\transp, \overline{v}\transp.
%
% This text uses `T' constantly for linear maps ($T$, $T_A$, $[T]$), so an
% italic $A^T$ reads as an index or a map name. An upright sans-serif T cannot
% be confused with it and survives being shrunk to superscript size, which is
% why it is preferred here over a bold T (bold at 70% mostly reads as smudged,
% and bold+colour is already reserved for \textbf / \emph in this document).
%
% To switch the whole document to a bold transpose, change \mathsf to \mathbf
% on the next line -- nothing else needs touching.
\newcommand{\transposesymbol}{\mathsf{T}}
\newcommand{\transp}{^{\transposesymbol}}


% ==========================================
% HEADER & FOOTER
% ==========================================
% Header: chapter (left) / page x of y (right).
% Footer: current section (left) / current subsection (right).
\pagestyle{fancy}
\fancyhf{}

% \parbox[b] lets a multi-line header grow upwards, so it always sits flush
% with the page number on the rule.
\fancyhead[L]{\parbox[b]{0.8\textwidth}{\sloppy\raggedright\small\sffamily\color{HeaderFooterFaint}\leftmark}}
\fancyhead[R]{\parbox[b]{0.18\textwidth}{\sloppy\raggedleft\small\sffamily\color{HeaderFooterFaint}\textbf{\thepage} / \pageref*{LastPage}}}

\fancyfoot[L]{\parbox[t]{0.47\textwidth}{\sloppy\raggedright\small\sffamily\color{HeaderFooterFaint}\rightmark}}
\fancyfoot[R]{\parbox[t]{0.47\textwidth}{\sloppy\raggedleft\small\sffamily\color{HeaderFooterFaint}\lastxmark}}

\renewcommand{\headrule}{\hbox to\headwidth{\color{HeaderFooterLine}\leaders\hrule height 0.5pt\hfill}}
\renewcommand{\footrule}{\hbox to\headwidth{\color{HeaderFooterLine}\leaders\hrule height 0.5pt\hfill}}

% Chapter openings switch to 'plain' automatically: no rules, centred page number.
\fancypagestyle{plain}{%
  \fancyhf{}%
  \fancyfoot[C]{\small\sffamily\color{HeaderFooterFaint}\textbf{\thepage} / \pageref*{LastPage}}%
  \renewcommand{\headrule}{}%
  \renewcommand{\footrule}{}%
}

% What goes into the marks.
% LaTeX's \appendix switches \thechapter to letters and sets \@chapapp, but leaves
% \chaptername alone -- so the heading, the TOC entry and the running header all announced
% "Chapter A --- Ordinary Differential Equations". Renaming \chaptername inside \appendix
% fixes the heading and the header, since both now ask for it by name rather than spelling
% "Chapter" out.
%
% The TOC needs the second line as well. \tableofcontents typesets the .toc file at the very
% front of the document, long before \appendix is ever executed, so at that moment
% \chaptername is still "Chapter". Writing the redefinition *into* the .toc makes the switch
% happen at the right position in the list instead.
\appto\appendix{%
  \addtocontents{toc}{\protect\renewcommand{\protect\chaptername}{\protect\appendixname}}%
  \renewcommand{\chaptername}{\appendixname}}

\renewcommand{\chaptermark}[1]{\markboth{\chaptername\ \thechapter\ \ ---\ \ #1}{}}
\renewcommand{\sectionmark}[1]{%
  \markright{\thesection\ \ ---\ \ #1}%
  \extramarks{}{}%
}
\renewcommand{\subsectionmark}[1]{%
  \extramarks{#1}{#1}%
}

\begin{document}

\begin{titlepage}
\thispagestyle{empty}
\begin{tikzpicture}[remember picture, overlay]
  % Dezent geometrische Akzentformen oben rechts
  \fill[ThemeSecondary!12, rounded corners=18pt] ($(current page.north east) + (-6.5cm,-2.5cm)$) rectangle ($(current page.north east) + (-1.2cm,-7.8cm)$);
  \fill[ThemeWeekNumber!20, rounded corners=12pt] ($(current page.north east) + (-4.8cm,-4.2cm)$) rectangle ($(current page.north east) + (-0.6cm,-9.0cm)$);
\end{tikzpicture}

\vspace*{4.5cm}
\noindent\begin{minipage}{\textwidth}
    {\sffamily\fontsize{40pt}{48pt}\selectfont\bfseries\color{ThemePrimary!95!black} Analysis II \par}
    \vspace{1.2cm}
    {\color{ThemeWeekNumber}\rule{\linewidth}{2.5pt}\par}
    \vspace{0.2cm}
    {\color{ThemeAccent1}\rule{0.6\linewidth}{1.5pt}\par}
    \vspace{2.0cm}
    {\sffamily\fontsize{18pt}{24pt}\selectfont\bfseries\color{TextEmphColor} Several Variables --- Exercise Class Notes \& Transcript \par}
    \vspace{0.6cm}
    {\sffamily\fontsize{15pt}{22pt}\selectfont\color{TextBoldColor} Lecturer: \textbf{Prof. Joaquim Serra} \par}
    \vspace{0.1cm}
    {\sffamily\fontsize{14pt}{20pt}\selectfont\color{TextBoldColor} Teaching Assistant: \textbf{Corsin Nick} \par}
    \vspace{0.3cm}
    {\sffamily\fontsize{14pt}{20pt}\selectfont\color{OliveGreen} Department of Mathematics -- ETH Zürich \par}
    \vspace{2.5cm}
    {\sffamily\fontsize{14pt}{20pt}\selectfont\bfseries\color{ThemeWeekNumber}\faGraduationCap\quad Spring Semester 2026 (401-1262-07L) \par}
\end{minipage}
\vfill
\noindent\begin{minipage}{\textwidth}
    \begin{tcolorbox}[
        arc=6pt,
        boxrule=0.7pt,
        colback=MidnightBlue!3!white,
        colframe=MidnightBlue!30!white,
        left=10pt, right=10pt, top=8pt, bottom=8pt
    ]
    {\sffamily\bfseries\small\color{MidnightBlue!90!black}\faInfoCircle\quad Notice \& Transcription Disclaimer \par}
    \vspace{0.25cm}
    \sffamily\small\color{black!75}
    \begin{itemize}[leftmargin=1.5em, itemsep=3pt, topsep=0pt, parsep=0pt]
        \item[\textcolor{MidnightBlue!80!black}{\small\faRobot}] \textbf{AI-Generated Notes:} This document was automatically generated using AI/LLM systems from the hand-written notes.
        \item[\textcolor{MidnightBlue!80!black}{\small\faUserShield}] \textbf{Unofficial Document:} This is \emph{not} official course material of ETH Zürich or the lecturers, and has \emph{not} been manually corrected or proofread by the staff (unreviewed draft).
        \item[\textcolor{MidnightBlue!80!black}{\small\faExclamationTriangle}] \textbf{Disclaimer:} Artificial intelligence can make errors (e.g., in notation, indices, or logical derivations). Use at your own risk --- without warranty of any kind regarding accuracy, completeness, or relevance to exams.
    \end{itemize}
    \end{tcolorbox}
\end{minipage}
\vspace*{0.5cm}
\end{titlepage}
% TOC entry spacing is now handled by the \titlecontents blocks in the preamble.
\tableofcontents % This generates the TOC
\hypersetup{linkcolor=LinkColor} % restore after the per-level TOC link colours
\cleardoublepage

% ==========================================
% CONTENT
% ==========================================
% The document is organised by topic, not by teaching week. Each chapter lives in its own
% directory under content/ as a thin 00-chapter.tex stub (\chapter + intro + \input lines)
% plus one file per \section. Every moved file records its provenance in an
% "% Originally:" comment at the top; the "% Source:" comments are unchanged.
% Chapters still carrying a week-NN.tex name have not been migrated yet.

\part{Foundations}

% Generator: Claude Opus 5 (High) --- part-level framing prose, new content, not tied to any
% tutor source. Rewritten 2026-08-09 in the fuller register style.md now prescribes for \part.
\emph{Analysis II} begins where \emph{Analysis I} stopped, and the first thing to settle is how
much of it carries over. The definitions of convergence, of continuity, and of a set being open
survive the move to several variables essentially verbatim. What does not survive is every
statement that quietly used the ordering of $\mathbb{R}$, and each of those has to be rebuilt on
a new hypothesis. This part draws that line, fixes the notation for $\mathbb{R}^n$, and rehearses
two habits used constantly later on: sketching a subset of $\mathbb{R}^n$ from the inequalities
defining it, and reading a function of two variables off its level sets. It also proves Young's
inequality, which looks like an isolated exercise here and turns out to be the engine behind the
Hölder and Minkowski inequalities in the next part.

% ============================================================
% Chapter 1 --- Prerequisites and a first look at R^n
% Originally: content/week-01.tex (no single tutor; assembled from several
%             tutors' recap slides, since Corsin's notes start a week later)
% ============================================================
\chapter{Prerequisites and a First Look at \texorpdfstring{$\mathbb{R}^n$}{Rn}}
\label{ch:prerequisites}

% Transition: Claude Opus 5 (Medium)
This chapter collects the \emph{Analysis I} material the course takes for granted, together with
the $\mathbb{R}^n$ warm-up that precedes any of its theory. It has no single tutor behind it:
Corsin's own class notes open with the second week of teaching, and no tutor's notes cover the
opening week on its own terms, since that week consists mostly of a recap and an exercise in
reading and drawing subsets of $\mathbb{R}^n$. The material below is therefore assembled from the
recap slides several tutors independently found worth revisiting. It also proves Young's
inequality (\cref{prop:youngs_inequality}), a fact used repeatedly once norms and Hölder-type
estimates appear.

% Originally: content/week-01.tex, \section{From Analysis I to several variables}
% Source: Maarten Cnoops/Woche 1.pdf, p. 1
\section{From Analysis I to several variables}
\label{sec:analysis1_recap}

Several of the definitions you will meet in the coming chapters are exactly the definitions you
already know from \emph{Analysis I}, only stripped of their reliance on the real line. A
\newterm{sequence} (\germanterm{Folge}) $(x_k)_{k \in \mathbb{N}}$ in a set $X$ converging to
$x \in X$, a function being continuous at a point, a set being open or closed --- all of these
notions carry over to $\mathbb{R}^n$ and, as \cref{sec:structured_spaces} will show, to much more
general spaces than $\mathbb{R}^n$. What does \emph{not} carry over automatically is anything
that relied on the \emph{order} of $\mathbb{R}$: there is no \qt{$\leq$} on $\mathbb{R}^n$ for
$n \geq 2$, so statements like the intermediate value theorem or \qt{every bounded sequence has a
convergent subsequence} need new proofs, or new hypotheses, once we leave $\mathbb{R}$.

\begin{remark}
That last fact --- bounded does \emph{not} imply convergent, but bounded sequences in
$\mathbb{R}^n$ do always have a convergent \emph{subsequence} --- is exactly the statement of the
Bolzano--Weierstrass theorem, revisited via \cref{ex:1.1} below. It
reappears as the Heine--Borel theorem once compactness is introduced in \cref{sec:compactness}.
\end{remark}

% Supplement: Linus Lüchinger/Slides/Slides-02-23.pdf
\begin{ainote}
Linus Lüchinger's first exercise-session slides collect the most common mistakes on the first
problem sheet, whose problems are distributed through this chapter. Two are worth flagging before
attempting \cref{ex:1.1}: on part 2, most students first guessed $f(x) = \log|x| + c$ for a single
constant $c$ --- wrong, since the two constants on $(-\infty,0)$ and $(0,\infty)$ may genuinely
differ (see \cref{sec:prerequisites_solutions}). On part 4, most students assumed a function with
everywhere-positive second derivative cannot have a global maximum; it can, at the boundary of the
domain.
\end{ainote}

% Source: exercises/Ex1_Analysis2_eng.pdf, p. 1
\begin{exercise}[1.1 --- Recap Questions]
\label{ex:1.1}
\begin{enumerate}[label=\textbf{\arabic*.}]
  \item Let $\{x_k\}_{k \in \mathbb{N}} \subset [1,2] \subset \mathbb{R}$ be a sequence. Is it
    necessarily convergent?
  \item Describe all differentiable functions $f : \mathbb{R} \setminus \{0\} \to \mathbb{R}$ such
    that $f'(x) = 1/x$ for all $x \neq 0$.
  \item Let $f : \mathbb{R} \to \mathbb{R}$ be continuous and bounded. Is it true that
    \[
      \int_0^{\pi/2} f(\cos(x))\,dx = \int_0^1 \frac{f(t)}{\sqrt{1-t^2}}\,dt \, ?
    \]
  \item Let $f : [0,\infty) \to \mathbb{R}$ have $f''(x) > 0$ for all $x \geq 0$. Can $f$ have a
    global maximum point?
  \item Let $\phi : \mathbb{R}^2 \to \mathbb{R}^2$ be a linear map and let
    $Q := [0,1] \times [0,1]$ be the unit square. How does $\phi(Q)$ look like? What are the
    possibilities?
  \item Consider the vector space $V$ of polynomials of degree less than $4$, and the function
    $D : V \to V$ which associates to a polynomial $p$ its derivative $p'$. Is $D$ a linear map?
    If so, what is the dimension of its null space?
\end{enumerate}
\end{exercise}

% Originally: content/week-01.tex, \section{Young's inequality}
% Authored for these notes around the statement of problem 1.4; no tutor's notes cover it.
\section{Young's inequality}
\label{sec:youngs_inequality}

The following inequality is not part of the lecture's main line of development, but it is the
cleanest route to comparing different norms on $\mathbb{R}^n$ (see \cref{sec:cauchy_schwarz}) and
to the Hölder-type estimates that appear whenever two $L^p$-type quantities are played off against
each other. It is worth proving once and citing from then on.

\begin{proposition}[Young's inequality]
\label{prop:youngs_inequality}
% Source: exercises/Ex1_Analysis2_eng.pdf, Problem 1.4
Let $a, b > 0$ and let $p, q > 1$ satisfy $\tfrac{1}{p} + \tfrac{1}{q} = 1$. Then
\[
  ab \leq \frac{a^p}{p} + \frac{b^q}{q}.
\]
\end{proposition}

\begin{remark}
The case $p = q = 2$ recovers the familiar inequality $ab \leq \tfrac{a^2}{2} + \tfrac{b^2}{2}$,
which is just a rewriting of $(a-b)^2 \geq 0$.
\end{remark}

\begin{proof}
Fix conjugate exponents $p, q > 1$, and define $g : (0,\infty) \to \mathbb{R}$ by
$g(t) := \tfrac{t}{p} + \tfrac{1}{q} - t^{1/p}$. We first show that $g \geq 0$, with equality only
at $t = 1$. Differentiating,
\[
  g'(t) = \frac{1}{p}\left(1 - t^{-(p-1)/p}\right).
\]
Since $p > 1$, the exponent $-(p-1)/p$ is negative, so $t^{-(p-1)/p} > 1$ for $t \in (0,1)$ and
$t^{-(p-1)/p} < 1$ for $t > 1$. Hence $g' < 0$ on $(0,1)$ and $g' > 0$ on $(1,\infty)$, so $g$
attains a global minimum at $t = 1$. As $g(1) = \tfrac{1}{p} + \tfrac{1}{q} - 1 = 0$, we obtain
\begin{equation}
\label{eq:young_scalar_key}
  t^{1/p} \leq \frac{t}{p} + \frac{1}{q} \qquad \text{for all } t > 0.
\end{equation}
Now fix $a, b > 0$ and apply \eqref{eq:young_scalar_key} with $t := a^p / b^q$. Multiplying both
sides by $b^q > 0$ gives
\[
  \left(\frac{a^p}{b^q}\right)^{1/p} b^q \ \leq\ \frac{a^p}{p} + \frac{b^q}{q}.
\]
It remains to identify the left-hand side, and this is where the conjugacy of $p$ and $q$ enters.
Using $1 - \tfrac{1}{p} = \tfrac{1}{q}$,
\[
  \left(\frac{a^p}{b^q}\right)^{1/p} b^q
    \ =\ a\,b^{-q/p}\,b^{q}
    \ =\ a\,b^{\,q\left(1 - \frac{1}{p}\right)}
    \ =\ a\,b^{\,q \cdot \frac{1}{q}}
    \ =\ ab,
\]
which is exactly the asserted inequality.
\end{proof}

% Generator: Claude Opus 5 (High)
\begin{remark}[From Young to Hölder]
\label{rem:young_to_hoelder}
The exercise sheet also asks for the vector version
\begin{equation}
\label{eq:young_vector}
  \langle a,b\rangle \leq \frac{\|a\|_p^p}{p} + \frac{\|b\|_q^q}{q}, \qquad a,b \in \mathbb{R}^n,
\end{equation}
where $\|x\|_p := \left(\sum_i |x_i|^p\right)^{1/p}$. This is \emph{not} yet Hölder's inequality:
it is the scalar case applied coordinatewise and summed (see the solution of \cref{ex:1.4}
below), and the right-hand side is a sum of $p$-th and $q$-th powers rather than the product
$\|a\|_p\|b\|_q$. The step from one to the other is a normalisation, and it is worth seeing
once, because the same trick recurs whenever an inequality is homogeneous.

Assume $a,b \neq 0$ and apply \eqref{eq:young_vector} to the rescaled vectors
$a/\|a\|_p$ and $b/\|b\|_q$, each of which has $p$- resp.\ $q$-norm equal to $1$:
\[
  \frac{\langle a,b\rangle}{\|a\|_p\,\|b\|_q} \ \leq\ \frac{1}{p} + \frac{1}{q} \ =\ 1,
\]
which is exactly \newterm{Hölder's inequality} \germanfor{Höldersche Ungleichung}
$\langle a,b\rangle \leq \|a\|_p\|b\|_q$. The conjugacy condition
$\tfrac1p+\tfrac1q=1$, which so far looked like an arbitrary bookkeeping device, is what makes
the two constants add up to exactly $1$ here.

This is also a first hint of why $p$-norms other than $p=2$ are worth naming:
\Cref{sec:cauchy_schwarz} only treats $p = q = 2$, where Hölder's inequality is
Cauchy--Schwarz, but the same machinery bounds $\langle a, b\rangle$ for every conjugate pair
$(p,q)$.
\end{remark}

% Source: exercises/Ex1_Analysis2_eng.pdf, p. 2
\begin{exercise}[1.4 --- Young's inequality]
\label{ex:1.4}
Let $a,b > 0$, and let $p,q>1$ with $\tfrac{1}{p}+\tfrac{1}{q}=1$. Prove that
$ab \leq \tfrac{a^p}{p} + \tfrac{b^q}{q}$.
\begin{enumerate}[label=\textbf{\arabic*.}]
  \item Prove that for all $t>0$, $t^{1/p} \leq \tfrac{t}{p} + \tfrac{1}{q}$.
  \item Multiply the inequality by $b^q$ and choose an appropriate value of $t$ to conclude.
\end{enumerate}
\textbf{A bonus:} show that $\langle a,b\rangle \leq \tfrac{\|a\|_p^p}{p} + \tfrac{\|b\|_q^q}{q}$
for all $a,b \in \mathbb{R}^n$.
\end{exercise}

% Originally: content/week-01.tex, \section{Visualizing sets in R^2 and R^3}
% Authored for these notes; no single source. The remark at the end of the section draws on
% several tutors' opening-week slides (the live Desmos 3D / GeoGebra demonstrations).
\section{Visualizing sets in \texorpdfstring{$\mathbb{R}^2$}{R2} and \texorpdfstring{$\mathbb{R}^3$}{R3}}
\label{sec:visualizing_sets}

Much of the geometric intuition for this course --- open sets, level sets, submanifolds --- rests
on a skill that has nothing to do with any theorem: reading a set-builder description and
producing an accurate picture, or going the other way. Two techniques are worth naming
explicitly, both used repeatedly in the solutions below.

\begin{itemize}
  \item \textbf{Slicing.} To sketch a set $S \subseteq \mathbb{R}^3$ defined by an equation
    involving $z$, fix $z = t$ and sketch the resulting planar slice $S \cap \{z = t\}$ as $t$
    varies. Stacking the slices gives the full picture.
  \item \textbf{Level sets.} If a set is defined by $f(x,y) = t$ for a fixed $t$, first work out
    the level sets $\{f = t\}$ of $f$ for varying $t$; many natural functions (e.g.
    $f(x,y) = \max\{|x|,|y|\}$, whose level sets are squares, or $f(x,y) = \sqrt{x^2+y^2}$, whose
    level sets are circles) have level sets that are easy to name, which turns a 3D sketch into a
    stack of named 2D shapes.
\end{itemize}

Both techniques are applied to \cref{ex:1.2,ex:1.3} below, and worked out in
\cref{sec:prerequisites_solutions}.

\begin{remark}
Several tutors' slides (e.g.\ those using Desmos 3D or GeoGebra to plot these sets live) note that
students commonly guess the wrong growth rate for sets like
$\{(x,y,z) : z = \sqrt{1+x^2+y^2}\}$, sketching a cone instead of a hyperboloid-like bowl that
steepens towards slope $1$; the slicing technique above makes the mistake visible immediately, since the slice
at height $t$ only exists for $t \geq 1$ and has radius $\sqrt{t^2-1}$, not $t$.
\end{remark}

% Source: exercises/Ex1_Analysis2_eng.pdf, p. 1
\begin{exercise}[1.2 --- From equations to drawings]
\label{ex:1.2}
Sketch the following subsets of $\mathbb{R}^2$, $\mathbb{R}^3$.
\begin{enumerate}[label=\textbf{\arabic*.}]
  \item $A := \{(x,y) : x^2+y^2 < 1,\ y \geq x^2\}$
  \item $B := \{(x,y) : y^2 = x^3 - x\}$
  \item $C := \{(x,y) : \tfrac{\pi}{6} < \arctan(y/x) < \tfrac{\pi}{3},\ 0 < x < y\}$
  \item $D := \{(x,y,z) : \sqrt{x^2+y^2} = 1-z,\ 0 < z < 1\}$
  \item $E := \{(x,y,z) : z = \sqrt{1+x^2+y^2}\}$
  \item $F := \{(x,y,z) : z = \max\{|x|,|y|\}\}$
\end{enumerate}
\end{exercise}

% Source: exercises/Ex1_Analysis2_eng.pdf, p. 1
\begin{exercise}[1.3 --- From drawings to equations]
\label{ex:1.3}
Express in terms of functions and equations the following sets.
\begin{enumerate}[label=\textbf{\arabic*.}]
  \item The coordinates of a point in the plane that goes endlessly back and forth from the
    origin to the point $(-1,1)$.
  \item A planar domain with two holes in it.
  \item A bi-dimensional disk sitting inside the plane $\{x+y+z=0\}$ in $\mathbb{R}^3$.
  \item A Pac-Man shaped domain in the plane.
  \item A pyramid, whose base is a square, sitting in $\mathbb{R}^3$.
\end{enumerate}
\end{exercise}

% Originally: content/week-01.tex, \section{Solutions}
\section{Solutions}
\label{sec:prerequisites_solutions}

% No tutor sampled for sheet 1 left a full worked solution to transcribe; most of their
% opening-week material is session logistics or "gotcha" summaries. The solutions below were
% written for this document and cross-checked against exercises/Sol1_Analysis2_eng.pdf.

\begin{exercisesolution}[Solution of \cref{ex:analysis1_recap}]
\label{sol:analysis1_recap}
% Generator: Claude Sonnet 5 (Medium)
\begin{enumerate}[label=\textbf{(\alph*)}]
  \item No. The sequence $x_k := \tfrac{3}{2} + \tfrac{1}{2}(-1)^k$ lies in $[1,2]$ but
    oscillates between $1$ and $2$ without settling, so it does not converge (it has two distinct
    accumulation points, $1$ and $2$). What \emph{is} true, by the Bolzano--Weierstrass theorem, is
    that every bounded sequence in $\mathbb{R}$ has a convergent subsequence.
  \item On each of the intervals $(-\infty,0)$ and $(0,\infty)$ separately, $f'(x) = 1/x$ forces
    $f(x) = \log|x| + c$ for some constant $c$, since two antiderivatives on a common interval
    differ by a constant. The subtlety is that $\mathbb{R}\setminus\{0\}$ is \emph{disconnected}
    into these two intervals, so the two constants need not agree:
    \[
      f(x) =
      \begin{cases}
        \alpha + \log(x) & x > 0, \\
        \beta + \log(-x) & x < 0,
      \end{cases}
      \qquad \alpha, \beta \in \mathbb{R} \text{ arbitrary}.
    \]
    The answer $f(x) = \log|x| + c$ for a \emph{single} constant $c$ is wrong; connectedness of
    the domain is exactly what would be needed to force $\alpha = \beta$, and it is precisely
    what \cref{sec:connectedness} formalizes.
  \item Yes. Substituting $t = \cos(x)$, $dt = -\sin(x)\,dx$, and $\sin(x) = \sqrt{1-t^2}$ for
    $x \in [0,\pi/2]$ turns the left integral into the right one; this is the one-variable change
    of variables formula, which the change-of-variables theorem
    (\cref{sec:change_of_variables}) generalizes to $\mathbb{R}^n$.
  \item Yes: take $f(x) := e^{-x}$. Then $f''(x) = e^{-x} > 0$ for all $x$, yet $f$ attains its
    global maximum on $[0,\infty)$ at the boundary point $x=0$. Everywhere-positive second
    derivative (strict convexity) rules out an \emph{interior} local maximum, but says nothing
    whatever about the boundary of the domain.
  \item $\phi(Q)$ is the parallelogram spanned by $\phi(1,0)$ and $\phi(0,1)$, with vertices at
    $0$, $\phi(1,0)$, $\phi(0,1)$, and $\phi(1,0)+\phi(0,1)$. Two degenerate cases occur: if
    $\phi(1,0)$ and $\phi(0,1)$ are linearly dependent (parallel or one is zero) but not both
    zero, the parallelogram collapses to a line segment; if $\phi \equiv 0$, it collapses to the
    single point $\{0\}$.
  \item Yes: differentiation is linear on functions in general, hence in particular on the
    subspace of polynomials of degree $< 4$. With respect to the basis
    $\mathcal{B} = \{1, X, X^2, X^3\}$, the map $D$ sends $1 \mapsto 0$, $X \mapsto 1$,
    $X^2 \mapsto 2X$, $X^3 \mapsto 3X^2$, giving matrix
    \[
      [D]_{\mathcal{B},\mathcal{B}} =
      \begin{pmatrix} 0&1&0&0 \\ 0&0&2&0 \\ 0&0&0&3 \\ 0&0&0&0 \end{pmatrix}.
    \]
    This matrix has rank $3$ (its three nonzero rows are independent), so by the rank--nullity
    theorem the null space has dimension $4 - 3 = 1$, and it is spanned by the constant
    polynomials.
\end{enumerate}
\end{exercisesolution}

% Supplement: Linus Lüchinger/Slides/Slides-02-23.pdf
% Moved here from 01-from-analysis-i.tex, where it stood above the exercise and gave away
% parts (b) and (d) before the reader had attempted them.
\begin{remark}[The two parts most often answered wrongly]
Linus Lüchinger's first exercise-session slides record where this sheet went wrong most often, and
two of the parts above account for most of it. Part \textbf{(b)} was one: the single-constant
answer $f(x) = \log\lvert x\rvert + c$ was the usual first guess, and it fails for the reason just
given, namely that the domain falls apart into two intervals. Part \textbf{(d)} was the other: a
function with
everywhere positive second derivative was widely assumed to admit no global maximum at all, when
in fact it may attain one at a boundary point. Both mistakes have the same shape. A hypothesis
that controls the \emph{interior} of a domain was read as though it controlled the domain.
\end{remark}

\begin{exercisesolution}[Solution of \cref{ex:equations_to_drawings}]
\label{sol:equations_to_drawings}
% Generator: Claude Sonnet 5 (Medium)
\begin{enumerate}[label=\textbf{(\alph*)}]
  \item $A$ is the intersection of the open unit disk with the region on or above the parabola
    $y = x^2$, a lens-shaped region:
    \begin{center}
    \begin{tikzpicture}[scale=1.8]
      \begin{scope}
        \clip (0,0) circle (1);
        \clip plot[domain=-1.1:1.1,samples=60] (\x,{\x*\x}) -- (1.1,1.3) -- (-1.1,1.3) -- cycle;
        \fill[blue!25] (-1.3,-1.3) rectangle (1.3,1.3);
      \end{scope}
      \draw[blue!60!black] (0,0) circle (1);
      \draw[blue!60!black,domain=-1:1,samples=60] plot (\x,{\x*\x});
      \node at (0,-0.55) {\footnotesize $x^2+y^2<1$};
      \node at (0,1.05) {\footnotesize $y \geq x^2$};
    \end{tikzpicture}
    \end{center}
  \item Since $x^3-x = x(x+1)(x-1) \geq 0$ exactly for $x \in [-1,0]\cup[1,\infty)$, the curve
    $y = \pm\sqrt{x^3-x}$ is only real on these two intervals; $B$ is the union of the graphs of
    $f(x) := \sqrt{x^3-x}$ and $-f(x)$, giving a closed loop over $[-1,0]$ and an unbounded branch
    over $[1,\infty)$ growing like $x^{3/2}$. Note that the two branches meet exactly where
    $f = 0$, i.e.\ at the three roots $x = -1, 0, 1$, which is why the loop closes up and the
    right branch starts with a vertical tangent.
    \begin{center}
    % Generator: Claude Opus 5 (High) --- FIG-W01-01.
    % Sign check: x^3-x = x(x-1)(x+1) >= 0 on [-1,0] and on [1,infty).
    % Max of x^3-x on [-1,0] is at x = -1/sqrt3 = -0.577, value 0.385, so |y| <= 0.62 there.
    \begin{tikzpicture}[>=Latex, scale=1.15]
      \draw[->] (-1.6,0) -- (2.15,0) node[right, font=\small] {$x$};
      \draw[->] (0,-2.25) -- (0,2.45) node[above, font=\small] {$y$};
      \draw (-1,0.06) -- (-1,-0.06) node[below left, font=\footnotesize] {$-1$};
      \draw (1,0.06) -- (1,-0.06) node[below right, font=\footnotesize] {$1$};

      \draw[blue!60!black, thick, domain=-1:0, samples=120]
        plot (\x, {sqrt(abs(\x*\x*\x-\x))});
      \draw[blue!60!black, thick, domain=-1:0, samples=120]
        plot (\x, {-sqrt(abs(\x*\x*\x-\x))});
      \draw[blue!60!black, thick, domain=1:1.8, samples=120]
        plot (\x, {sqrt(abs(\x*\x*\x-\x))});
      \draw[blue!60!black, thick, domain=1:1.8, samples=120]
        plot (\x, {-sqrt(abs(\x*\x*\x-\x))});

      \fill (-1,0) circle (1.3pt);
      \fill (0,0) circle (1.3pt);
      \fill (1,0) circle (1.3pt);
      % Labels are kept clear of the curve: the loop never exceeds |y| = 0.62, the strip
      % 0 < x < 1 is empty, and the right branch reaches y = 1.75 only at x = 1.63.
      \node[font=\footnotesize, blue!60!black, align=center] at (-0.72,1.3)
        {closed loop\\ over $[-1,0]$};
      % Both kept off the curve: at height y the right branch sits at x with x^3-x = y^2,
      % so y = 1.30 puts it at x = 1.44 and the label starts only at x = 1.95.
      \node[font=\footnotesize, blue!60!black, align=left, anchor=west] at (1.95,1.30)
        {unbounded branch,\\ $\sim x^{3/2}$};
      \node[font=\footnotesize, gray, align=left, anchor=west] at (-1.62,-1.55)
        {no real points\\ for $0<x<1$};
    \end{tikzpicture}
    \end{center}
    % Generator: Claude Sonnet 5 (Medium) --- the solution continues below; only the figure
    % above is an addition.
  \item Since $0 < x < y$, points of $C$ lie in the first octant with $y > x$; the condition
    $\pi/6 < \arctan(y/x) < \pi/3$ says the angle to the positive $x$-axis is between $30^\circ$
    and $60^\circ$. So $C$ is the open angular sector between the lines $y=x$ and $y=\sqrt{3}\,x$:
    \begin{center}
    \begin{tikzpicture}[scale=1.4]
      \fill[blue!20] (0,0) -- (1.6,1.6) -- (0.92,1.6) -- cycle;
      \draw[blue!60!black] (0,0) -- (1.6,1.6) node[right] {\footnotesize $y=x$};
      \draw[blue!60!black] (0,0) -- (0.92,1.6) node[above] {\footnotesize $y=\sqrt{3}x$};
      \draw[->] (-0.2,0) -- (1.8,0) node[right] {\footnotesize $x$};
      \draw[->] (0,-0.2) -- (0,1.8) node[above] {\footnotesize $y$};
      \draw (0.5,0) arc (0:45:0.5);
      \node at (0.65,0.22) {\footnotesize $30^\circ$--$60^\circ$};
    \end{tikzpicture}
    \end{center}
  \item Slicing at height $z=t \in (0,1)$: the equation $\sqrt{x^2+y^2} = 1-t$ describes a circle
    of radius $1-t$ centered on the $z$-axis. As $t$ ranges over $(0,1)$, the radius shrinks
    linearly from $1$ to $0$, so $D$ is the lateral surface of a right circular cone of height $1$
    over the unit circle, apex at $(0,0,1)$, base circle excluded (the inequalities are strict).
  \item Slicing at height $z=t$: the equation $\sqrt{1+x^2+y^2}=t$ has solutions only for
    $t \geq 1$, and the slice is then a circle of radius $\sqrt{t^2-1}$. Note that the radius is
    $\sqrt{t^2-1}$ and not $t$; this is the single most common slip on this part, and it turns the
    answer into a cone. $E$ is instead a bowl-shaped surface of revolution with vertex at
    $(0,0,1)$, obtained by rotating the curve $y = \sqrt{1+x^2}$ about the $z$-axis. It approaches
    slope $1$ for large $x,y$ without ever attaining it, which is exactly how it differs from the
    cone $z = \sqrt{x^2+y^2}$ that it is so often mistaken for.
  \item The level sets $\{\max\{|x|,|y|\}=t\}$ are the boundaries of the squares $[-t,t]^2$.
    Stacking these square outlines at height $z=t$ as $t$ grows from $0$ upward sweeps out four
    flat triangular faces meeting at the apex $(0,0,0)$. So $F$ is an upward-opening
    \emph{square cone}: the exact analogue of the circular cone of part \textbf{(d)}, with the
    round level circles of $\sqrt{x^2+y^2}$ replaced by the square level curves of
    $\max\{|x|,|y|\}$.
\end{enumerate}

Parts \textbf{(d)}, \textbf{(e)} and \textbf{(f)} are the same construction carried out three
times, stacking level sets, and are best compared side by side. In each panel the dotted curve is
one intermediate level set:

\begin{center}
% Generator: Claude Opus 5 (High) --- FIG-W01-02.
% Oblique projection: a horizontal circle of radius R is drawn as an ellipse (R and 0.34*R).
% (4) apex (0,0,1) at the top, radius 1-z, so radius 1 at z=0: a downward-opening cone.
% (5) z = sqrt(1+x^2+y^2): vertex (0,0,1), radius sqrt(z^2-1); at z=2 the radius is sqrt3=1.73.
% (6) z = max(|x|,|y|): apex at the origin, level sets are SQUARES of half-side z.
\begin{tikzpicture}[>=Latex, scale=0.95]

  % ---------- Panel (4): circular cone, apex up ----------
  \begin{scope}[shift={(0,0)}]
    \draw[->] (0,-0.35) -- (0,2.5) node[above, font=\footnotesize] {$z$};
    \draw[blue!70!black, thick] (0,2.0) -- (-1.3,0);
    \draw[blue!70!black, thick] (0,2.0) -- (1.3,0);
    \draw[blue!70!black, thick, dashed] (0,0) ellipse (1.3 and 0.44);
    \draw[blue!70!black, dotted] (0,1.0) ellipse (0.65 and 0.22);
    \fill (0,2.0) circle (1.5pt) node[above right, font=\footnotesize] {$(0,0,1)$};
    \node[font=\footnotesize, align=center] at (0,-1.0)
      {\textbf{(d)} $\sqrt{x^2+y^2}=1-z$\\ radius shrinks $1 \to 0$};
  \end{scope}

  % ---------- Panel (5): bowl, opening upward ----------
  \begin{scope}[shift={(4.4,0)}]
    \draw[->] (0,-0.35) -- (0,2.5) node[above, font=\footnotesize] {$z$};
    % asymptotic cone z = |x| (slope 1), drawn from the origin
    \draw[gray!65, dashed] (-1.9,1.9) -- (0,0) -- (1.9,1.9);
    \draw[blue!70!black, thick, domain=-1.75:1.75, samples=80]
      plot (\x, {sqrt(1+\x*\x)});
    \draw[blue!70!black, thick] (0,2.02) ellipse (1.75 and 0.30);
    \draw[blue!70!black, dotted] (0,1.41) ellipse (1.0 and 0.20);
    \fill (0,1.0) circle (1.5pt) node[below right, font=\footnotesize] {$(0,0,1)$};
    \node[gray!70!black, font=\footnotesize, anchor=west] at (1.2,0.45) {$z=|x|$};
    \node[font=\footnotesize, align=center] at (0,-1.0)
      {\textbf{(e)} $z=\sqrt{1+x^2+y^2}$\\ radius $\sqrt{z^2-1}$, \emph{not} $z$};
  \end{scope}

  % ---------- Panel (6): square cone, apex at the origin ----------
  \begin{scope}[shift={(8.8,0)}]
    \draw[->] (0,-0.35) -- (0,2.5) node[above, font=\footnotesize] {$z$};
    % the four corners of the square [-t,t]^2 at height t=2, in the same oblique projection
    \coordinate (c1) at (1.3,2.44);
    \coordinate (c2) at (1.3,1.56);
    \coordinate (c3) at (-1.3,1.56);
    \coordinate (c4) at (-1.3,2.44);
    \draw[blue!70!black, thick] (c1) -- (c2) -- (c3) -- (c4) -- cycle;
    \foreach \c in {(c1),(c2),(c3),(c4)}{\draw[blue!70!black, thick] (0,0) -- \c;}
    \draw[blue!70!black, dotted] (0.65,1.22) -- (0.65,0.78) -- (-0.65,0.78) -- (-0.65,1.22) -- cycle;
    \fill (0,0) circle (1.5pt) node[below right, font=\footnotesize] {$(0,0,0)$};
    \node[font=\footnotesize, align=center] at (0,-1.0)
      {\textbf{(f)} $z=\max\{|x|,|y|\}$\\ squares instead of circles};
  \end{scope}
\end{tikzpicture}
\end{center}
\end{exercisesolution}

\begin{exercisesolution}[Solution of \cref{ex:drawings_to_equations}]
\label{sol:drawings_to_equations}
% Generator: Claude Sonnet 5 (Medium)
\begin{enumerate}[label=\textbf{(\alph*)}]
  \item A point moving back and forth along the segment from $(0,0)$ to $(-1,1)$ can be
    parametrized as $t \mapsto (-f(t), f(t))$ for $t \geq 0$, where $f : [0,\infty) \to [0,1]$
    oscillates endlessly; $f(t) := \tfrac{1}{2}(1+\sin t)$ is one valid choice.
  \item Remove two small disjoint squares from a large square, e.g.
    $\{10 > \max\{|x|,|y|\} > 1\} \setminus \{\max\{|x-5|,|y-6|\} \leq 1\}$.
  \item Intersect the closed unit ball with the plane: $\{x^2+y^2+z^2 \leq 1\} \cap \{x+y+z=0\}$.
  \item Remove a $60^\circ$ wedge from a disk:
    $\{x^2+y^2 \leq 1\} \setminus \{|\arctan(y/x)| \leq \pi/6,\ x>0\}$.
  \item Combining the square level sets from \cref{ex:equations_to_drawings} part \textbf{(f)} with the conical
    slicing from part \textbf{(d)}: $\{1-z = \max\{|x|,|y|\},\ 0 \leq z \leq 1\}$ is a pyramid of
    height $1$ over the square $[-1,1]^2$, apex at $(0,0,1)$.
\end{enumerate}

The two planar answers are worth drawing, since in both cases the description \qt{remove a piece}
translates directly into a set difference:

\begin{center}
% Generator: Claude Opus 5 (High) --- FIG-W01-03.
\begin{tikzpicture}[>=Latex, scale=1.0]

  % ---------- Panel 2: square annulus with a second square removed ----------
  % Matches the formula above exactly. The outer square is max{|x|,|y|} < 10 (drawn with
  % half-side 1.5, i.e. one figure unit = 6.67 units of the formula); the FIRST hole is the
  % central max{|x|,|y|} <= 1, the SECOND is the square of half-side 1 centred at (5,6),
  % i.e. at (0.75, 0.90) in figure coordinates. Both holes therefore have half-side 0.15.
  \begin{scope}[shift={(0,0)}]
    \fill[blue!12] (-1.5,-1.5) rectangle (1.5,1.5);
    \draw[blue!60!black, thick] (-1.5,-1.5) rectangle (1.5,1.5);
    \fill[white] (-0.15,-0.15) rectangle (0.15,0.15);
    \draw[blue!60!black, thick] (-0.15,-0.15) rectangle (0.15,0.15);
    \fill[white] (0.60,0.75) rectangle (0.90,1.05);
    \draw[blue!60!black, thick] (0.60,0.75) rectangle (0.90,1.05);
    \node[font=\scriptsize, anchor=east] at (-0.26,0) {hole at $0$};
    \node[font=\scriptsize, anchor=west] at (0.99,0.90) {hole at $(5,6)$};
    \node[font=\scriptsize, align=center, anchor=north] at (0,-1.72)
      {\textbf{2.} a domain with two holes:\\
       $\{10>\max\{|x|,|y|\}>1\}$\\ minus a unit square at $(5,6)$};
  \end{scope}

  % ---------- Panel 4: Pac-Man ----------
  \begin{scope}[shift={(5.5,0)}]
    % disk minus the 60-degree wedge |arctan(y/x)| <= pi/6, x > 0
    \fill[blue!12] (0,0) -- (30:1.5) arc (30:330:1.5) -- cycle;
    \draw[blue!60!black, thick] (0,0) -- (30:1.5) arc (30:330:1.5) -- cycle;
    % the removed wedge, dashed
    \draw[red!70!black, thick, dashed] (0,0) -- (30:1.5);
    \draw[red!70!black, thick, dashed] (0,0) -- (-30:1.5);
    \draw[red!70!black, dotted] (0.55,0) arc (0:30:0.55);
    \node[red!70!black, font=\footnotesize, anchor=west] at (1.0,0.0) {$60^\circ$ wedge removed};
    \fill (0,0) circle (1.2pt);
    \node[font=\scriptsize, align=center, anchor=north] at (0,-1.72)
      {\textbf{4.} Pac-Man: the closed unit disk\\ minus the wedge
       $\{|\arctan(y/x)|\leq\pi/6,\ x>0\}$};
  \end{scope}
\end{tikzpicture}
\end{center}
\end{exercisesolution}

\begin{exercisesolution}[Solution of \cref{ex:youngs_inequality}]
\label{sol:youngs_inequality}
% Generator: Claude Sonnet 5 (Medium); display re-broken into align* by Claude Opus 5 (High).
Parts \textbf{(a)} and \textbf{(b)} are exactly the proof of \cref{prop:youngs_inequality} above.
For the bonus vector version, apply the scalar inequality coordinatewise, with $|a_i|$ in place of
$a$ and $|b_i|$ in place of $b$, and sum over $i$:
\begin{align*}
  \langle a,b\rangle
    &\ \leq\ \sum_{i=1}^n |a_i|\,|b_i|
     \ \leq\ \sum_{i=1}^n \left(\frac{|a_i|^p}{p} + \frac{|b_i|^q}{q}\right) \\[0.4ex]
    &\ =\ \frac{1}{p}\sum_{i=1}^n |a_i|^p + \frac{1}{q}\sum_{i=1}^n |b_i|^q
     \ =\ \frac{\|a\|_p^p}{p} + \frac{\|b\|_q^q}{q}.
\end{align*}
Note that the two sums separate only because the right-hand side of Young's inequality is a
\emph{sum} of powers rather than a product; this is precisely the feature that
\cref{rem:young_to_hoelder} then removes by rescaling.
\end{exercisesolution}



% Chapter order follows the dependency chain, which is also Corsin's own teaching order:
% metrics -> open/closed sets -> sequences -> continuity -> completeness -> compactness ->
% connectedness. Completeness sits after continuity because a contraction is defined as a
% Lipschitz map with constant < 1.
\part{Metric Spaces and Topology}

% Generator: Claude Opus 5 (High) --- part-level framing prose, new content.
Very little of one-variable analysis actually used the arithmetic of $\mathbb{R}$. Convergence,
continuity and compactness need only a notion of distance, and this part takes that observation
seriously by developing all three for an arbitrary metric space. The payoff is not generality for
its own sake. Function spaces such as $C^0([0,1],\mathbb{R})$ are metric spaces too, so once the
theory is stated at this level, the same arguments apply to a sequence of functions and to a
sequence of points alike. Two properties do most of the work: completeness, which guarantees that
a limit one has constructed actually exists, and compactness, which allows an infinite problem to
be replaced by a finite one. Almost every existence proof in the remaining five parts leans on
one of them.

% ============================================================
% Chapter 2 --- Metric, normed and inner-product spaces
% Originally: content/week-02.tex (Corsin Week 2, Monday) --- the hierarchy of structured
%             spaces and the metric axioms
%           + content/week-03.tex (Corsin Week 3, Friday) --- Cauchy--Schwarz and the
%             geometry of the inner product
% ============================================================
\chapter{Metric, Normed and Inner-Product Spaces}
\label{ch:metric_spaces}

% Transition: Claude Opus 5 (High)
% Corsin taught the two halves of this chapter a week apart, opening the later one with a recap
% of the definitions given here. sec:inner_products_recap keeps that recap, now as a pointer back
% to sec:structured_spaces rather than a repetition.
How much structure does a set need before one can do analysis on it? This chapter answers that by
taking $\mathbb{R}^n$ apart into layers. Distance comes first, and distance alone already carries
convergence and continuity. Adding a linear structure gives a norm, which measures the length of
one vector rather than the gap between two. Adding an inner product on top of that gives angles,
and with them orthogonal projection. The layers nest, each one strictly inside the last, and the
inequality that makes the innermost layer work at all is Cauchy--Schwarz. A last section settles
the question all of this raises: on $\mathbb{R}^n$ the choice of norm turns out not to matter for
anything topological, and it matters a great deal that this is a theorem rather than an
assumption.

% Originally: content/week-02.tex, \section{Structured spaces} (Corsin Week 2, Monday)

\section{Structured spaces}
\label{sec:structured_spaces}

This semester, you will be introduced to certain structures that can be given to a set, which in
this context is also called a \newterm{space} \germanfor{Raum} and which we will denote by $X$:

\begin{itemize}
  \item \textbf{0. Topological spaces:} fourth semester.
  % German mirror deliberately omitted here: this bullet is a preview, and the first *formal*
  % introduction of the term is \cref{def:metric_space}, which carries the \germanterm pair.
  \item \textbf{1. Metric spaces} $(X,\ d : X \times X \to \mathbb{R}_{\geq 0})$
    \begin{itemize}
      \item Defines a \textbf{distance} between points.
      \item $X$ is \textbf{not} necessarily a vector space.
      \item Covered in \textit{Analysis}.
    \end{itemize}
  \item \textbf{2. Normed metric spaces} $(X,\ \|\cdot\| : X \to \mathbb{R}_{\geq 0})$
    \begin{itemize}
      \item Defines the \textbf{length} of a vector.
      \item Presupposes a linear structure: $X$ must be a \textbf{vector space}.
      \item Covered in \textit{Linalg / Analysis}.
    \end{itemize}
  \item \textbf{3. Inner product spaces} $(X,\ \langle\cdot,\cdot\rangle : X \times X \to \mathbb{R})$
    \begin{itemize}
      \item Defines \textbf{both} lengths \textbf{and} angles!
      \item Super important!
      \item Covered in \textit{Linalg} (and sometimes used in \textit{Analysis}). $X$ is a \textbf{vector space}.
    \end{itemize}
\end{itemize}

\begin{center}
\begin{tikzpicture}
\fill[red!25] (0,0) ellipse (3.2 and 2.4);
\fill[yellow!35] (0,-0.1) ellipse (2.4 and 1.8);
\fill[green!25] (0,-0.2) ellipse (1.6 and 1.2);
\fill[blue!25] (0,-0.3) ellipse (0.8 and 0.6);
\draw[red!70!black] (0,0) ellipse (3.2 and 2.4);
\draw[yellow!70!black] (0,-0.1) ellipse (2.4 and 1.8);
\draw[green!60!black] (0,-0.2) ellipse (1.6 and 1.2);
\draw[blue!70!black] (0,-0.3) ellipse (0.8 and 0.6);
\node[red!70!black] at (0,2.0) {\footnotesize Topological spaces};
\node[olive!60!black] at (0,1.35) {\footnotesize Metric spaces};
\node[green!60!black] at (0,0.75) {\footnotesize Normed vector spaces};
\node[blue!70!black] at (0,-0.3) {\footnotesize Inner product spaces};
\end{tikzpicture}
\end{center}

You will see that an \textbf{inner product} $\langle\cdot,\cdot\rangle$ induces a \textbf{norm}
by the formula
\[\|x\| := \sqrt{\langle x, x\rangle}, \qquad x \in X.\]
Furthermore, a \textbf{norm} induces a \textbf{metric} by
\[d(x,y) := \|x - y\|, \qquad x, y \in X.\]

% Moved here from Week 3 (\S Cauchy--Schwarz): Corsin defines inner products and norms on the
% Friday of Week 3, but the two supplements below -- and the two displayed formulas above --
% already use both, so the definitions are stated here, at the point of first use. Corsin's own
% development of the geometry (angle, projection, the second proof of Cauchy--Schwarz) stays
% where he taught it, in \cref{sec:cauchy_schwarz}.
\begin{definition}[Inner product / scalar product]
\label{def:inner_product}
Let $V$ be a \textbf{vector space} over $\mathbb{R}$. An \newterm{inner product}
\germanfor{Skalarprodukt} on $V$ is a map
\[
\langle\cdot,\cdot\rangle : V \times V \to \mathbb{R}
\]
satisfying, for all $u, v, w \in V$ and $\alpha, \beta \in \mathbb{R}$:
\begin{enumerate}[label=\textbf{(\alph*)}]
  \item \textbf{Bilinearity:}
    $\langle \alpha u + \beta v, w\rangle = \alpha\langle u,w\rangle + \beta\langle v,w\rangle$,
    and equivalently in the second variable.
  \item \textbf{Symmetry:} $\langle u,v\rangle = \langle v,u\rangle$.
  \item \textbf{Definiteness:} $\langle u,u\rangle \geq 0$, and $\langle u,u\rangle = 0 \iff u = 0$.
\end{enumerate}
\end{definition}

\begin{definition}[Norm]
\label{def:norm}
Let $V$ be a \textbf{real} vector space. A \newterm{norm} \germanfor{Norm} on $V$ is a map
$\|\cdot\| : V \to [0,\infty)$ satisfying, for all $u, v \in V$ and $\alpha \in \mathbb{R}$:
\begin{enumerate}[label=\textbf{(\alph*)}]
  \item \textbf{Definiteness:} $\|v\| \geq 0$, and $\|v\| = 0 \iff v = 0$.
  \item \textbf{Homogeneity:} $\|\alpha v\| = |\alpha|\|v\|$.
  \item \textbf{$\Delta$-inequality:} $\|u+v\| \leq \|u\| + \|v\|$.
\end{enumerate}
\end{definition}

% Supplement: Analysis_II_Script_v1.pdf, p. 31
\begin{ainote}
The official script (Remark 9.94, p.\ 31 onward) adopts the opposite notational convention from
this document: from that point on \emph{it} writes the Euclidean norm as $|x|$ and reserves
$\|\cdot\|$ for non-standard norms. This document instead uses $\|\cdot\|$
(\texttt{\textbackslash lVert}\ldots\texttt{\textbackslash rVert}) uniformly for every norm,
Euclidean or not, per the house style, and keeps that convention throughout. Neither choice is
more standard than the other; this note exists only to warn anyone cross-referencing the script
directly that a bare $|\cdot|$ there and a bare $\|\cdot\|$ here can both denote \qt{the}
Euclidean norm. The warning matters most in the script's chapters 13--14, where the $|\cdot|$
convention is used densely.
\end{ainote}

% Source: Jérôme Paschoud/Notizen/Woche 4.1 Normen, das Differential.pdf, p. 1
% Generator: Gemini 3.1 Pro (High)
\begin{example}[Hilbert--Schmidt norm on matrix spaces]
\label{ex:hilbert_schmidt_norm}
A very important example of a norm on the space of real $n \times m$ matrices
$\mathbb{R}^{n \times m}$ is the \newterm{Hilbert--Schmidt norm} (also known as the Frobenius
norm). It is defined by
\[
\|M\|_2 := \sqrt{\Tr(M\transp M)},
\]
where $\Tr$ denotes the trace of a matrix. Only the diagonal of $M\transp M$ is needed, so the
off-diagonal entries are left unevaluated below. For a $2 \times 3$ matrix
$M = \left(\begin{smallmatrix} a & b & c \\ d & e & f \end{smallmatrix}\right)$, we have
\begin{align*}
M\transp M
  &= \begin{pmatrix} a & d \\ b & e \\ c & f \end{pmatrix}
     \begin{pmatrix} a & b & c \\ d & e & f \end{pmatrix} \\[0.6ex]
  &= \begin{pmatrix} a^2+d^2 & * & * \\ * & b^2+e^2 & * \\ * & * & c^2+f^2 \end{pmatrix}.
\end{align*}
The trace is the sum of the diagonal entries, so $\|M\|_2 = \sqrt{a^2+b^2+c^2+d^2+e^2+f^2}$. This
norm is therefore nothing but the standard Euclidean norm of the matrix entries read as a single
vector in $\mathbb{R}^{6}$, or in general in $\mathbb{R}^{nm}$. The inner product inducing it is
$\langle A, B \rangle := \Tr(A\transp B)$.
\end{example}

% Transition: Claude Opus 5 (High)
Comparing the three lists of axioms is already informative. A norm has exactly the three
properties a metric has, restated for a single vector instead of a pair of points, plus
homogeneity. That extra axiom has no metric analogue at all, since stating it requires scalar
multiplication. An inner product goes further and replaces the triangle inequality by
bilinearity, which is far stronger, and then recovers the triangle inequality as a \emph{theorem}
(\cref{prop:induced_norm_and_metric}) rather than assuming it. That is the precise sense in which
each layer carries more structure than the one outside it.

% The two blocks that follow supplement Corsin's outline with material from other tutors. They
% belong here logically, being about the innermost layer of the hierarchy just drawn, even though
% Corsin himself develops inner-product geometry only later (sec:cauchy_schwarz), where his own
% proof is kept. That prop:cauchy_schwarz_geometric_proof and thm:cauchy_schwarz are the same
% computation twice over is a mathematical point, and is made where it belongs, in
% rem:two_proofs_cauchy_schwarz.

% Supplement: Sascha Brack/Class Notes/Week_02_Notes_Monday_Updated.pdf, p. 4
\begin{remark}[Geometric angle and projection in inner product spaces]
\label{rem:inner_product_angle_geometry}
In an inner product space $(X, \langle\cdot,\cdot\rangle)$ over $\mathbb{R}$, the Cauchy--Schwarz inequality $|\langle x, y\rangle| \leq \|x\|\,\|y\|$ guarantees that for any non-zero vectors $x, y \in X$,
\[
-1 \leq \frac{\langle x, y\rangle}{\|x\|\,\|y\|} \leq 1.
\]
This allows us to formally define the \newterm{angle} $\Theta(x,y) \in [0,\pi]$ between $x$ and $y$ via
\[
\Theta(x,y) := \arccos\left(\frac{\langle x, y\rangle}{\|x\|\,\|y\|}\right).
\]
Geometrically, when $\|x\| = \|y\| = 1$, the inner product $\langle x, y\rangle$ is the signed
length of the orthogonal projection of $y$ onto the direction of $x$. In that reading
Cauchy--Schwarz becomes visually obvious: a projection of a unit vector cannot be longer than the
unit vector itself, so $\langle x,y\rangle \leq 1$, with equality exactly when $y$ already points
along $x$.
\begin{center}
% Generator: Claude Opus 5 (High) --- redrawn. The previous version parked a 6cm text block
% inside the picture, which pushed the figure off-centre and stated in words what the drawing
% should state by itself. Both vectors are unit, so they sit on the dashed unit circle; the
% angle arc now shows Theta, and the projection foot P = (cos 50, 0) is marked with a right
% angle. Angle 50 degrees: cos(50) = 0.643, sin(50) = 0.766.
\begin{tikzpicture}[>=Latex, scale=2.7]
  \coordinate (O) at (0,0);
  \coordinate (X) at (1,0);
  \coordinate (Y) at ({cos(50)},{sin(50)});
  \coordinate (P) at ({cos(50)},0);

  \draw[->, gray!55] (-0.18,0) -- (1.32,0);
  \draw[->, gray!55] (0,-0.18) -- (0,1.12);

  % both vectors are unit, so they end on this circle
  \draw[gray!40, dashed] (O) circle (1);

  % angle between x and y
  \draw[green!45!black, thick] (0.32,0) arc (0:50:0.32);
  \node[green!45!black, font=\footnotesize] at (0.50,0.17) {$\Theta$};

  % dropped perpendicular, with a right-angle mark at the foot
  \draw[blue!55, dashed] (Y) -- (P);
  \draw[blue!55] ({cos(50)-0.07},0) -- ({cos(50)-0.07},0.07) -- ({cos(50)},0.07);

  % The projection has to be measured BELOW the axis: drawn on the axis it would be collinear
  % with the arrow for x and simply disappear underneath it.
  \draw[orange!90!black, dotted] (0,0) -- (0,-0.17);
  \draw[orange!90!black, dotted] ({cos(50)},0) -- ({cos(50)},-0.17);
  \draw[orange!90!black, line width=0.9pt, <->] (0,-0.17) -- ({cos(50)},-0.17);
  \node[orange!90!black, font=\footnotesize, below] at ({cos(50)/2},-0.18)
    {$\langle x,y\rangle$};

  \draw[blue!60!black, thick, ->] (O) -- (X) node[above right, font=\small] {$x$};
  \draw[blue!60!black, thick, ->] (O) -- (Y) node[above left, font=\small] {$y$};

  \node[gray!55!black, font=\footnotesize, anchor=west] at (1.0,0.92) {$\|x\|=\|y\|=1$};
\end{tikzpicture}
\end{center}
\end{remark}

% Quelle: Diego Torres Tejeda/Notes - 09.03 - Diego Torres Tejeda.pdf, pp. 1--2
\begin{proposition}[Geometric derivation of the Cauchy--Schwarz inequality]
\label{prop:cauchy_schwarz_geometric_proof}
Let $(V, \langle\cdot,\cdot\rangle)$ be a real inner product space. Then for all $v, w \in V$,
\[
|\langle v, w\rangle| \leq \|v\|\,\|w\|.
\]
\end{proposition}

\begin{proof}
If $w = 0$, both sides vanish and the inequality holds trivially. For $w \neq 0$, consider the squared distance from $v$ to the points of the line $L := \operatorname{span}_{\mathbb{R}}\{w\}$, parametrized by $\lambda \mapsto \lambda w$:
\[
f(\lambda) := \|v - \lambda w\|^2 = \langle v - \lambda w, v - \lambda w\rangle = \|v\|^2 - 2\lambda\langle v, w\rangle + \lambda^2\|w\|^2.
\]
Since $f(\lambda) \geq 0$ for all $\lambda \in \mathbb{R}$, $f$ is a convex quadratic polynomial in $\lambda$ with positive leading coefficient $\|w\|^2 > 0$. Differentiating or completing the square shows that $f$ attains its global minimum at $\lambda_0 := \frac{\langle v, w\rangle}{\|w\|^2}$. Evaluating $f$ at $\lambda_0$ yields
\[
0 \leq f(\lambda_0) = \|v\|^2 - 2\frac{\langle v, w\rangle^2}{\|w\|^2} + \frac{\langle v, w\rangle^2}{\|w\|^2} = \|v\|^2 - \frac{\langle v, w\rangle^2}{\|w\|^2}.
\]
Multiplying by $\|w\|^2$ and taking square roots gives $|\langle v, w\rangle| \leq \|v\|\,\|w\|$. Geometrically, $\pi_L(v) := \lambda_0 w = \frac{\langle v, w\rangle}{\|w\|^2}w$ is the unique orthogonal projection of $v$ onto the line $L$, achieving the minimal distance $d(v, L) = \sqrt{f(\lambda_0)}$.
\end{proof}

% Generator: Claude Opus 5 (High) --- moved here from Week 3 together with \cref{def:norm}:
% Cauchy--Schwarz is now available (immediately above), which is the only input the triangle
% inequality needs, so the two arrows of the hierarchy can be justified where they are drawn.
\begin{proposition}[Inner product $\Rightarrow$ norm $\Rightarrow$ metric]
\label{prop:induced_norm_and_metric}
Let $V$ be a real vector space.
\begin{enumerate}[label=\textbf{(\alph*)}]
  \item If $\langle\cdot,\cdot\rangle$ is an inner product on $V$ (\cref{def:inner_product}),
    then $\|v\| := \sqrt{\langle v,v\rangle}$ is a norm on $V$ (\cref{def:norm}).
  \item If $\|\cdot\|$ is a norm on $V$, then $d(x,y) := \|x-y\|$ is a metric on $V$
    (\cref{def:metric_space}).
\end{enumerate}
\end{proposition}

\begin{proof}
\textbf{(a)} Definiteness of the norm is definiteness of the inner product, and homogeneity
follows from bilinearity: $\|\alpha v\|^2 = \langle\alpha v,\alpha v\rangle =
\alpha^2\|v\|^2$. For the triangle inequality, expand and apply
\cref{prop:cauchy_schwarz_geometric_proof}:
\begin{align*}
  \|u+v\|^2 &= \|u\|^2 + 2\langle u,v\rangle + \|v\|^2 \\
            &\leq \|u\|^2 + 2\lvert\langle u,v\rangle\rvert + \|v\|^2 \\
            &\leq \|u\|^2 + 2\|u\|\,\|v\| + \|v\|^2 \ =\ \big(\|u\|+\|v\|\big)^2,
\end{align*}
and taking square roots gives $\|u+v\| \leq \|u\|+\|v\|$.

\textbf{(b)} Definiteness: $d(x,y)=0$ means $\|x-y\|=0$, hence $x=y$. Symmetry: homogeneity
with $\alpha=-1$ gives $\|x-y\| = \|-(y-x)\| = \|y-x\|$. Triangle inequality: writing
$x-z = (x-y)+(y-z)$,
\[d(x,z) = \|(x-y)+(y-z)\| \leq \|x-y\| + \|y-z\| = d(x,y)+d(y,z). \qedhere\]
\end{proof}

\begin{remark}[Neither arrow reverses]
\label{rem:norm_not_from_inner_product}
The nesting drawn above is strict at both steps.
\begin{itemize}
  \item \textbf{Not every metric comes from a norm.} The discrete metric
    (\cref{ex:metric_examples}\textbf{(a)}) on a vector space cannot: a norm-induced metric
    satisfies $d(\alpha x, 0) = |\alpha|\,d(x,0)$, which fails as soon as
    $\alpha \neq 0,\pm1$.
  \item \textbf{Not every norm comes from an inner product.} A norm induced by an inner product
    obeys the \newterm{parallelogram law}
    $\|u+v\|^2 + \|u-v\|^2 = 2\|u\|^2 + 2\|v\|^2$ (expand both sides by bilinearity). The
    supremum norm on $\mathbb{R}^2$ violates it: for $u := (1,0)$ and $v := (0,1)$ the left side
    is $1+1 = 2$ and the right side is $2+2 = 4$.
\end{itemize}
Consequently, each layer of the hierarchy really does carry strictly more structure than the one
outside it. The picture at the start of this section is a genuine nesting of four distinct
classes, not a chain of equivalent descriptions of the same thing.
\end{remark}

% Source: Corsin Nick/Class Notes/Week 2.pdf, pp. 1--13

% Originally: content/week-02.tex, \section{Metric spaces} (Corsin Week 2, Monday)

\section{Metric spaces}
\label{sec:metric_spaces}

\begin{definition}[Metric space]
\label{def:metric_space}
A \newterm{metric space} (\germanterm{metrischer Raum}) $(X,d)$ is any set $X$ with a
\newterm{metric} (\germanterm{Metrik}) $d : X \times X \to \mathbb{R}_{\geq 0}$ such that for all
$x, y, z \in X$:
\begin{enumerate}[label=\textbf{(\alph*)}]
  \item \textbf{Definiteness:} $d(x,y) \geq 0$, with equality if and only if $x = y$. In words: the
    distance is non-negative, and a point has zero distance to itself.
  \item \textbf{Symmetry:} $d(x,y) = d(y,x)$, i.e., $x$ has the same distance to $y$ as $y$ to $x$.
  \item \textbf{Triangle inequality ($\Delta$-inequality):} $d(x,z) \leq d(x,y) + d(y,z)$ ---
    \qt{detours make the way longer}.
\end{enumerate}
\end{definition}

\begin{ainote}
Corsin's gloss on definiteness reads \qt{the distance positive}; the formula $d \geq 0$ is
correct, but the gloss omits the ``non-''. Corrected in the prose above.
\end{ainote}

\begin{center}
\begin{tikzpicture}[>=Latex]
\draw[->] (-0.3,0) -- (4.2,0);
\draw[->] (0,-0.3) -- (0,3);
\coordinate (x) at (0.2,0.2);
\coordinate (y) at (1.6,2.0);
\coordinate (z) at (3.6,2.3);
\draw[blue,thick,->] (x) -- (y);
\draw[blue,thick,->] (y) -- (z);
\draw[red,thick,->] (x) -- (z);
\fill (x) circle (1.5pt) node[below]{$x$};
\fill (y) circle (1.5pt) node[above]{$y$};
\fill (z) circle (1.5pt) node[above]{$z$};
\node[red] at (2.3,0.7) {\footnotesize $d(x,z)$};
\node[blue] at (0.6,1.3) {\footnotesize $d(x,y)$};
\node[blue] at (2.7,2.4) {\footnotesize $d(y,z)$};
\node[red,below] at (2,-0.6) {\small $d(x,z) \le d(x,y)+d(y,z)$};
\end{tikzpicture}
\end{center}

% Generator: Claude Opus 5 (Medium)
\begin{aiexample}[Each axiom rules out something, and none is redundant]
\label{ex:metric_axioms_are_needed}
Three functions on $\mathbb{R}$, each failing exactly one axiom of
\Cref{def:metric_space}. It is worth checking these once: a definition only becomes memorable
once you know what it is keeping out.
\begin{enumerate}[label=\textbf{(\alph*)}]
  \item \textbf{Definiteness fails:} $d(x,y) := |x^2-y^2|$. This is symmetric and satisfies the
    triangle inequality (it is $|\cdot|$ composed with $x \mapsto x^2$), but
    $d(-1,1) = 0$ while $-1 \neq 1$. Distinct points at distance zero: the space cannot tell
    $-1$ from $1$, and no sequence could ever have a unique limit.
  \item \textbf{Symmetry fails:} $d(x,y) := \max\{\,y-x,\ 2(x-y)\,\}$, i.e.\ moving to the right
    costs $1$ per unit and moving to the left costs $2$. Here $d(0,1)=1$ but $d(1,0)=2$, so the
    return trip is more expensive than the outward one. Definiteness still holds ---
    $d(x,y) = 0$ forces $y-x \leq 0 \leq y-x$, hence $x=y$ --- and so does the triangle
    inequality, since $d(x,y) = \varphi(y-x)$ for $\varphi(t) := \max\{t,-2t\}$, and $\varphi$
    is a maximum of two linear functions vanishing at $0$, hence convex and positively
    homogeneous, hence subadditive. (Such a $d$ is called a \newterm{quasi-metric} and is
    genuinely useful elsewhere --- but $B_r(x)$ would then depend on which argument the centre
    sits in.)
  \item \textbf{Triangle inequality fails:} $d(x,y) := (x-y)^2$. Definiteness and symmetry are
    clear, but
    \[d(0,2) = 4 \quad > \quad d(0,1)+d(1,2) = 1+1 = 2,\]
    so the detour through $1$ is \emph{shorter} than going direct --- exactly the
    \qt{detours make the way longer} slogan reversed.
\end{enumerate}
\end{aiexample}

\subsection{Examples}

% Transition: Claude Opus 5 (High)
Three axioms are few enough that a great many things qualify as metrics, and the point of the
list below is exactly that breadth. Two of the examples are not what one would draw on a
blackboard as \qt{distance}: the discrete metric measures nothing but \qt{same or different},
and the supremum metric measures the distance between two \emph{functions}. Both are legitimate,
and it is worth letting them stretch the intuition now --- the second, in particular, is the
setting in which uniform convergence turns out to be an ordinary case of
\cref{def:convergent_sequence} (\cref{ex:uniform_convergence}).

% Source: Corsin Nick/Class Notes/Week 2.pdf, pp. 1--13
\begin{example}[Standard metrics on $\mathbb{R}^n$]
\label{ex:metric_examples}
\begin{enumerate}[label=\textbf{(\alph*)}]
  \item Let $X$ be a set. The \newterm{discrete metric} is given by
    \[d(x,y) := \begin{cases} 1 & \text{if } x \neq y \\ 0 & \text{if } x = y \end{cases}\]
    Important for counterexamples!
  \item Let $X := \mathbb{R}^n$. Then we have many choices for a metric; some common ones are:
  \begin{enumerate}[label=\textbf{(\roman*)}]
    \item \textbf{Manhattan metric:}
      \[d_1(x,y) := |x_1 - y_1| + \dots + |x_n - y_n| = \sum_{k=1}^{n} |x_k - y_k|\]
    \item \textbf{Standard metric:}
      \[d_2(x,y) := \sqrt{(x_1-y_1)^2 + \dots + (x_n-y_n)^2} = \sqrt{\sum_{k=1}^{n}(x_k - y_k)^2}\]
    \item \textbf{Supremum metric:}
      \[d_3(x,y) := \sup_{k = 1,\dots,n} |x_k - y_k|\]
  \end{enumerate}
\end{enumerate}
\end{example}

\begin{center}
\begin{tikzpicture}[>=Latex, scale=1.1]
  \draw[->] (-0.3,0) -- (4.5,0) node[right] {$x_1$};
  \draw[->] (0,-0.3) -- (0,3.5) node[above] {$x_2$};
  \coordinate (X) at (0.8, 0.8);
  \coordinate (Y) at (3.8, 2.8);
  \coordinate (Corner) at (3.8, 0.8);

  \draw[dotted, gray] (0, 0.8) node[left, font=\footnotesize] {$x_2$} -- (X) -- (0.8, 0) node[below, font=\footnotesize] {$x_1$};
  \draw[dotted, gray] (0, 2.8) node[left, font=\footnotesize] {$y_2$} -- (Y) -- (3.8, 0) node[below, font=\footnotesize] {$y_1$};

  % The two legs are just coordinate differences -- neither one IS a metric.
  \draw[orange!90!black, ultra thick] (X) -- (Corner) -- (Y);
  \node[below, font=\small, orange!90!black] at (2.3,0.72) {$a := |x_1-y_1|$};
  \node[right, font=\small, orange!90!black] at (3.9,1.8) {$b := |x_2-y_2|$};

  \draw[blue!80!black, very thick] (X) -- (Y);
  \node[above left, font=\small, blue!80!black] at (2.9,2.1) {$d_2 = \sqrt{a^2+b^2}$};

  \fill (X) circle (2pt) node[above left] {$x$};
  \fill (Y) circle (2pt) node[above right] {$y$};

  % d_3 picks out the longer leg; here a > b.
  \draw[red!80!black, thick, <->] (0.8,-0.95) -- node[below, font=\small] {$d_3 = \max\{a,b\} = a$} (3.8,-0.95);
  \draw[red!80!black, dotted] (0.8,0) -- (0.8,-0.95);
  \draw[red!80!black, dotted] (3.8,0) -- (3.8,-0.95);

  \node[orange!90!black, font=\small, anchor=west] at (0.1,3.25) {$d_1 = a+b$ \ (the full orange path)};
\end{tikzpicture}
\end{center}

Note that $d_1$ is the \textbf{sum} of the two legs and $d_3$ the \textbf{larger} of them ---
neither is a single leg on its own. Only $d_2$ is the straight-line distance you would draw by
reflex.

In general,
\[
d(x,y) := \sqrt[m]{\sum_{k=1}^{n} (x_k - y_k)^m}
\]
is a metric on $\mathbb{R}^n$ for any $m$.

\begin{ainote}
Two caveats Corsin leaves out. The powers must be taken of \textbf{absolute values},
$\sqrt[m]{\sum_k|x_k-y_k|^m}$ --- otherwise for odd $m$ the radicand can be negative. And it is
a metric only for $m \geq 1$: for $m \in (0,1)$ the triangle inequality fails, as
\Cref{ex:4.3}\textbf{.6} shows with $x=e_1$, $y=e_2$, where
$|x+y|_m = 2^{1/m} > 2 = |x|_m+|y|_m$. Written as stated above.
\end{ainote}

\begin{example}[continued]
\begin{enumerate}[label=\textbf{(\alph*)}, start=3]
  \item Let $X := C^0([0,1], \mathbb{R})$ be the set of continuous functions $[0,1] \to \mathbb{R}$.
    Then we can define the \newterm{supremum metric} for $f, g \in X$:
    \[
    d(f,g) := \sup_{x \in [0,1]} |f(x) - g(x)|
    \]
    This is well defined --- i.e.\ the supremum is a real number and not $+\infty$ --- because
    $|f-g|$ is continuous on the compact interval $[0,1]$ and therefore bounded, by the extreme
    value theorem from \emph{Analysis I} (revisited in this generality as
    \cref{lem:sequential_compactness_extreme_values} below). Compactness of the domain is doing
    real work here: on $C^0(\mathbb{R},\mathbb{R})$ the same formula would produce $\infty$ for
    $f(x) := x$, $g := 0$.
  \item Let $X := S^2 := \{(x_1,x_2,x_3) \in \mathbb{R}^3 : x_1^2 + x_2^2 + x_3^2 = 1\}$. $S^n$ is
    called the \newterm{$n$-dimensional sphere}, or \newterm{$n$-sphere}. The closest path between
    two points on the sphere is an \textbf{arc of a great circle} (a \textit{geodesic}). This
    defines a metric, similar to how straight lines define a metric on $\mathbb{R}^n$.
\end{enumerate}
\end{example}

\begin{center}
\begin{tikzpicture}[scale=1.1]
  \begin{scope}[shift={(0,0)}]
    \draw[thick] (0,0) circle (1.4);
    \draw[dashed, gray] (1.4,0) arc (0:180:1.4 and 0.45);
    \draw[gray] (-1.4,0) arc (180:360:1.4 and 0.45);
    \coordinate (x1) at (-0.4,-0.6);
    \coordinate (y1) at (0.3,0.9);
    \draw[blue!80!black, thick] (x1) to[out=55, in=230] (y1);
    \fill[red!80!black] (x1) circle (2pt) node[left] {$x$};
    \fill[red!80!black] (y1) circle (2pt) node[right] {$y$};
    \node[blue!80!black, font=\footnotesize] at (-0.3,0.3) {$d(x,y)$};
    \node[below, font=\small] at (0,-1.5) {Distance on sphere};
  \end{scope}

  \begin{scope}[shift={(4.5,0)}]
    \draw[thick] (0,0) circle (1.4);
    \draw[dashed, gray] (1.4,0) arc (0:180:1.4 and 0.45);
    \draw[gray] (-1.4,0) arc (180:360:1.4 and 0.45);
    \coordinate (x2) at (-0.4,-0.6);
    \coordinate (y2) at (0.3,0.9);
    \draw[blue!80!black, thick] (x2) to[out=55, in=230] (y2);
    \draw[orange!90!black, thick, dashed] (x2) -- (y2);
    \fill[red!80!black] (x2) circle (2pt) node[left] {$x$};
    \fill[red!80!black] (y2) circle (2pt) node[right] {$y$};
    \node[blue!80!black, font=\tiny] at (-0.6,0.3) {sphere arc};
    \node[orange!90!black, font=\tiny] at (0.2,0.0) {Euclidean chord};
    \node[below, font=\small] at (0,-1.5) {Sphere arc vs.\ Euclidean chord};
  \end{scope}
\end{tikzpicture}
\end{center}

\begin{exercise}[Supremum metric]
\label{ex:sup_metric_is_metric}
Show that the supremum metric from \cref{ex:metric_examples} \textbf{(c)} is a metric.
\end{exercise}

% Kept inline: solved directly in Corsin Nick/Class Notes/Week 2.pdf, p. 6.
\begin{exercisesolution}[Proof of \cref{ex:sup_metric_is_metric}]
Clearly, symmetry and definiteness hold. For the triangle inequality, let $f, g, h \in C^0([0,1])$ and
set $p := f - h$ and $q := h - g$. Then for all $x \in [0,1]$, we have:
\[|p(x)| \leq \sup|p| \qquad \text{and} \qquad |q(x)| \leq \sup|q|.\]
Adding the inequalities, we obtain for all $x \in [0,1]$:
\begin{align*}
|f(x) - g(x)| &= |p(x) + q(x)| \\
&\leq |p(x)| + |q(x)| \\
&\leq \sup|p| + \sup|q| \\
&= \sup|f - h| + \sup|h - g|.
\end{align*}
\end{exercisesolution}

% Quelle: Jérôme Paschoud/Notizen/Woche 2.1 Metrische Räume.pdf, p. 2
\begin{exercise}[$L^1$-metric]
\label{ex:integral_metric_is_metric}
Let $X := C^0([0,1], \mathbb{R})$ be the set of continuous functions $[0,1] \to \mathbb{R}$.
For $f,g \in X$, define
\[d(f,g) := \int_0^1 \lvert f(t)-g(t)\rvert\,dt.\]
Is $(X,d)$ a metric space? Provide a proof.
\end{exercise}

% Generator: Gemini 3.6 Pro
\begin{exercisesolution}[Proof of \cref{ex:integral_metric_is_metric}]
Yes, it is a metric space. We verify the three axioms of a metric:
\begin{enumerate}[label=\textbf{(\alph*)}]
  \item \textbf{Definiteness:} Clearly $d(f,g) \geq 0$ since the integrand $\lvert f(t)-g(t)\rvert$ is non-negative. If $f = g$, then $d(f,f) = \int_0^1 0\,dt = 0$. Conversely, suppose $d(f,g) = 0$. Since $\lvert f-g\rvert$ is a continuous, non-negative function, its integral over $[0,1]$ can only be zero if the integrand is identically zero. Thus $f(t) = g(t)$ for all $t \in [0,1]$, which means $f = g$.
  \item \textbf{Symmetry:} $d(f,g) = \int_0^1 \lvert f(t)-g(t)\rvert\,dt = \int_0^1 \lvert g(t)-f(t)\rvert\,dt = d(g,f)$.
  \item \textbf{Triangle inequality:} Let $f, g, h \in X$. The triangle inequality for real numbers yields $\lvert f(t)-h(t)\rvert \leq \lvert f(t)-g(t)\rvert + \lvert g(t)-h(t)\rvert$ for all $t \in [0,1]$. Integrating both sides, we obtain:
  \begin{align*}
  d(f,h) &= \int_0^1 \lvert f(t)-h(t)\rvert\,dt \\
  &\leq \int_0^1 \left(\lvert f(t)-g(t)\rvert + \lvert g(t)-h(t)\rvert\right)\,dt \\
  &= \int_0^1 \lvert f(t)-g(t)\rvert\,dt + \int_0^1 \lvert g(t)-h(t)\rvert\,dt \\
  &= d(f,g) + d(g,h).
  \end{align*}
\end{enumerate}
\end{exercisesolution}

% The old chapter closed here with a transition into "Open and closed sets"; that topic is now
% \cref{ch:open_closed}, whose own opening paragraph says the same thing, so the bridge is gone.

% --- exercises relocated here from the dissolved sheets ---

% Transition: Claude Opus 5 (Medium)
\begin{ainote}[Where the exercise-sheet problems live]
The official problem sheets are not reproduced as blocks anywhere in this document. Each problem
is quoted verbatim (from \texttt{exercises/ExN\_Analysis2\_eng.pdf}; attribution: Prof.\ Joaquim
Serra, D-MATH, ETH Z\"urich) and placed with the material it tests, so a problem from sheet 3 may
well appear several chapters away from a problem from sheet 2. Its solution is then in the
\emph{Solutions} section of that same chapter. Corsin's priority tags (\textbf{important},
\textbf{semi-important}, \textbf{optional}) and the official \textbf{(*)} difficulty marker are
carried along in each problem's title.
\end{ainote}

\begin{notation}
% Originally: content/exercise-sheets/week-02.tex, the sheet's own preamble.
Some problems in this document have a closed-answer format, similar to what you might find on the
final exam. \qt{Multiple Choice} means that zero, one or more answers can be correct. Questions
marked with (*) are a bit more complex, you might want to skip them at the first read.
\end{notation}

% Originally: content/week-02.tex, end of \section{Compactness}

\begin{aiexercise}[Equivalence of Manhattan and Supremum Metrics]
\label{ex:ai_metric_equivalence}
% Generator: Gemini 3.6 Flash (Medium)
Let $x := (x_1, x_2) \in \mathbb{R}^2$. Show that the Manhattan metric $d_1(x,0) := |x_1| + |x_2|$ and the supremum metric $d_3(x,0) := \max\{|x_1|, |x_2|\}$ satisfy the sharp inequalities:
\[
d_3(x,0) \leq d_1(x,0) \leq 2 \, d_3(x,0).
\]
\end{aiexercise}

% Originally: content/exercise-sheets/week-02.tex, problem 2.1
% Source: exercises/Ex2_Analysis2_eng.pdf, p. 1

\begin{exercise}[2.1 --- Examples and Non-examples of Metric spaces]
\label{ex:2.1}
Which of the following pairs are metric spaces? Prove it or provide a counterexample.
\begin{enumerate}[label=\textbf{\arabic*.}]
  \item $(B(X), d)$, where $B(X)$ denotes the set of all bounded functions from a non-empty set
    $X$ to $\mathbb{R}$ and
    $$d(f,g) := \sup_{x \in X} |f(x) - g(x)|.$$
  \item $(\mathbb{Q}_+, d)$, where $\mathbb{Q}_+$ are the positive rational numbers and
    $d(x,y) := \left|\tfrac{1}{x} - \tfrac{1}{y}\right|$.
  \item $(\mathbb{R}^2, d)$, where $d(x,y) := (x_1 - y_1)^2 + |x_2 - y_2|$.
  \item $(\mathbb{R}^2, d)$, where $d(x,y) := |x_1 - y_1|^{1/2} + |x_2 - y_2|$.
  \item $(\mathbb{R}^{n \times n}, d)$ with
    $d(X,Y) := \left(\Tr\{(X-Y)\transp(X-Y)\}\right)^{1/2}$.
  \item \textbf{(*)} $(\mathbb{R}^{n \times n}, d)$ with
    $$d(X,Y) := \sup\{|v\transp(X-Y)v| : v \in \mathbb{R}^n, \|v\| = 1\},$$
    and $\mathbb{R}^{n\times n}$ denoting the set of square matrices.
  \item \textbf{(*)} $(\mathbb{R}^2/\mathbb{Z}^2, d)$ where the flat 2-dimensional torus
    $\mathbb{R}^2/\mathbb{Z}^2$ is the set of equivalence classes of pairs of real numbers under
    the equivalence relation
    $$x, y \in \mathbb{R}^2, \quad x \sim y \iff x_1 - y_1 \in \mathbb{Z},\ x_2 - y_2 \in \mathbb{Z},$$
    and $d([x],[y]) := \inf_{k,h \in \mathbb{Z}} \|x - y + (k,h)\|$, with $\|\cdot\|$ denoting the
    Euclidean distance and $[x] \in \mathbb{R}^2/\mathbb{Z}^2$ denoting the equivalence class of
    $x \in \mathbb{R}^2$.
\end{enumerate}

\textit{Official hints:} \textbf{2.1.3} and \textbf{2.1.4} --- ignore what happens in the second
variable, it is there only to distract you. \textbf{2.1.5} --- rewrite, for a general matrix
$A := X - Y$, what $\Tr\{A\transp A\}^{1/2}$ actually means; it should look familiar.
\textbf{2.1.6} --- see what happens if $(X-Y)$ is anti-symmetric. \textbf{2.1.7} --- convince
yourself first that the ``inf'' is in fact a ``min''.
\end{exercise}

% Quelle: Diego Torres Tejeda/Notes - 02.03 - Diego Torres Tejeda.pdf, p. 2
\begin{ainote}[Diego's geometric intuition for the flat torus $\mathbb{R}^2/\mathbb{Z}^2$]
Diego's notes highlight the geometric picture behind part \textbf{2.1.7}: every equivalence class in $\mathbb{R}^2/\mathbb{Z}^2$ has a representative in $[0,1] \times [0,1]$. Identifying opposite boundaries $(0,x) \sim (1,x)$ and $(x,0) \sim (x,1)$ folds the unit square into a cylinder and subsequently into a 2D torus $\mathbb{T}^2 := \mathbb{R}^2/\mathbb{Z}^2$. The distance $d([x],[y])$ is simply the shortest Euclidean distance between points on this folded torus.
\end{ainote}

% Originally: content/exercise-sheets/week-02.tex, problem 2.5
% Source: exercises/Ex2_Analysis2_eng.pdf, p. 1

\begin{exercise}[2.5 --- Product of metric spaces]
\label{ex:2.5}
Let $(X, d_X)$ and $(Y, d_Y)$ be a pair of metric spaces. Recall that the set of ordered pairs
$(x,y)$ with $x \in X$ and $y \in Y$ is denoted by $X \times Y$. Consider the following functions
$X \times Y \to [0,\infty)$:
$$
\begin{aligned}
d_1((x,y),(x',y')) &:= \max\{d_X(x,x'), d_Y(y,y')\} \\
d_2((x,y),(x',y')) &:= d_X(x,x') + d_Y(y,y') \\
d_3((x,y),(x',y')) &:= \sqrt{d_X(x,x')^2 + d_Y(y,y')^2}.
\end{aligned}
$$
\begin{enumerate}[label=\textbf{\arabic*.}]
  \item Show that they are all valid distance functions on $X \times Y$.
  \item Show that they are all equivalent, i.e.\ there is a number $C > 0$ such that
    $$d_1((x,y),(x',y')) \leq C\,d_2((x,y),(x',y')) \leq C^2 d_3((x,y),(x',y')) \leq C^3 d_1((x,y),(x',y'))$$
    for all $x, x' \in X$, $y, y' \in Y$.
  \item Show that a sequence $(x_n, y_n) \to (x,y)$ with respect to $(X \times Y, d_3)$ if and
    only if $x_n \to x$ with respect to $d_X$ and $y_n \to y$ with respect to $d_Y$.
\end{enumerate}
\textit{Official hint:} don't get distracted by the abstract set-up --- you already know all
these things for $\mathbb{R} \times \mathbb{R}$; start from there, then rewrite those arguments in
this general framework.
\end{exercise}

% Originally: content/week-03.tex, \section{Cauchy--Schwarz} (Corsin Week 3, Friday)

\section{Cauchy--Schwarz}
\label{sec:cauchy_schwarz}

\subsection{Recap: inner products and norms}
\label{sec:inner_products_recap}

% Transition: Claude Opus 5 (High)
Corsin opens his Friday session by recalling the two definitions, which this chapter has already
stated in \cref{sec:structured_spaces}, where the hierarchy of structured spaces was drawn and
where the first supplements needed them. Rather than repeat them, here is what to have in mind:
\begin{itemize}
  \item An \emph{inner product} (\cref{def:inner_product}) is \textbf{bilinear},
    \textbf{symmetric} and \textbf{definite}. It measures length \emph{and} angle.
  \item A \emph{norm} (\cref{def:norm}) is \textbf{definite}, \textbf{homogeneous} and
    satisfies the \textbf{triangle inequality}. It measures length only.
  \item Each induces the next, via $\|v\| := \sqrt{\langle v,v\rangle}$ and
    $d(x,y) := \|x-y\|$, and neither implication reverses
    (\cref{prop:induced_norm_and_metric}, \cref{rem:norm_not_from_inner_product}).
\end{itemize}
What this section adds is the \emph{geometry}: why $\langle x,y\rangle$ deserves to be read as an
angle at all, and the resulting second proof of Cauchy--Schwarz.

\subsection{Geometric meaning of the inner product}

In $\mathbb{R}^2$, the following identity holds:
\[
\langle x,y\rangle := x_1y_1 + x_2y_2 = \|x\|\|y\|\cos\theta
\]
where $\theta$ is the angle between $x, y \in \mathbb{R}^2$.

To see this, conveniently choose our basis vectors such that $x = (x_1, 0)$, $y = (y_1, y_2)$.

\begin{center}
\begin{tikzpicture}[>=Latex, scale=1.1]
  \begin{scope}[shift={(0,0)}]
    \draw[->] (-0.3,0) -- (2.2,0);
    \draw[->] (0,-0.3) -- (0,2.2);
    \draw[blue!80!black, thick, ->] (0,0) -- (1.5,0.6) node[right] {$x$};
    \draw[orange!90!black, thick, ->] (0,0) -- (0.8,1.8) node[above] {$y$};
    \draw[purple, thick] (0.5,0.2) arc (21.8:66.0:0.5);
    \node[purple, font=\footnotesize] at (0.7,0.45) {$\theta$};
  \end{scope}

  \draw[->, thick] (2.4,1.0).. controls (2.8,1.3).. (3.3,1.0) node[midway, above, red, font=\tiny] {\shortstack{change of basis\\(rotation)}};

  % Same lengths and same angle as the left panel -- a rotation preserves both.
  % Left:  |x| = 1.616 at 21.80 deg, |y| = 1.970 at 66.04 deg, so theta = 44.24 deg.
  % Right: x rotated onto the axis at (1.616, 0), y at 44.24 deg with |y| unchanged.
  \begin{scope}[shift={(3.8,0)}]
    \draw[->] (-0.3,0) -- (2.5,0);
    \draw[->] (0,-0.3) -- (0,2.2);
    \draw[blue!80!black, thick, ->] (0,0) -- (1.616,0) node[below] {$x$};
    \draw[orange!90!black, thick, ->] (0,0) -- (1.411,1.375) node[above right] {$y$};
    \draw[dotted, gray] (1.411,1.375) -- (1.411,0);
    \draw[purple, thick] (0.5,0) arc (0:44.24:0.5);
    \node[purple, font=\footnotesize] at (0.72,0.22) {$\theta$};
  \end{scope}
\end{tikzpicture}
\end{center}

Then
\[
\langle x,y\rangle = x_1y_1 = x_1\|y\|\cos\theta = \|x\|\|y\|\cos\theta.
\]

For more rigour, you will show in Linalg that the rotations in $\mathbb{R}^2$ are given by
matrices $R$ satisfying $\langle Rx, Ry\rangle = \langle x,y\rangle$.

In a general inner product space $V$, this becomes the \textbf{definition} of the angle $\theta$.
It then makes sense to define the \newterm{projection onto $v$}, $\pi_v$:
\[
\pi_v(u) := \frac{\langle u,v\rangle}{\|v\|^2}\,v = \|u\|\cos\theta\,\frac{v}{\|v\|}.
\]

\begin{center}
\begin{tikzpicture}[>=Latex, scale=1.2]
  \draw[->] (-0.3,0) -- (3.5,0);
  \draw[->] (0,-0.3) -- (0,2.5);

  \draw[orange!90!black, thick, ->] (0,0) -- (3.2,0.8) node[right] {$v$};
  \draw[blue!80!black, thick, ->] (0,0) -- (1.2,2.0) node[above left] {$u$};

  \coordinate (P) at (1.60, 0.40);
  \coordinate (U) at (1.2, 2.0);

  \draw[purple!80!black, thick, dashed] (U) -- (P);
  \fill[purple!80!black] (P) circle (2pt) node[below right, font=\small] {$\pi_v(u)$};

  \draw[green!60!black, thick] (0.5,0.125) arc (14:59:0.5);
  \node[green!60!black, font=\small] at (0.6,0.3) {$\theta$};
\end{tikzpicture}
\end{center}

In words, $\pi_v(u)$ is the part of $u$ which points along $v$. In particular,
\[
\|\pi_v(u)\| = \|u\|\,|\cos\theta| = \frac{|\langle u,v\rangle|}{\|v\|}.
\]
In this sense, Cauchy--Schwarz says something visually obvious. Reading the picture above, all
three of the following are the same statement:
\[
\frac{|\langle u,v\rangle|}{\|v\|} \leq \|u\| \iff \|\pi_v(u)\| \leq \|u\| \iff \|u\|\,|\cos\theta| \leq \|u\|.
\]

\begin{remark}[Hint for the proof of Cauchy--Schwarz]
\label{rem:proof_hint_cauchy_schwarz}
Use definiteness and expand the term
\[
\langle u - \pi_v(u),\ u - \pi_v(u)\rangle = \dots
\]
where $u - \pi_v(u)$ is the \textbf{perpendicular part} of $u$ to $v$.
\end{remark}

% Sonnet 5 (Medium)
\subsection{The Cauchy--Schwarz inequality}

% Corsin's own notes stop at the hint above (rem:proof_hint_cauchy_schwarz) and never state the
% inequality itself as a formal result, which is notable given that it names the section. The
% theorem below fills that in, carrying out exactly the computation he hints at.

\begin{theorem}[Cauchy--Schwarz inequality]
\label{thm:cauchy_schwarz}
Let $V$ be an inner product space (\cref{def:inner_product}) with induced norm
$\|\cdot\| := \sqrt{\langle\cdot,\cdot\rangle}$ (\cref{def:norm}). Then for all $u,v \in V$,
\[|\langle u,v\rangle| \leq \|u\|\,\|v\|,\]
with equality if and only if $u,v$ are linearly dependent.
\end{theorem}

\begin{proof}
If $v = 0$, both sides vanish and the claim holds trivially, so assume $v \neq 0$. Following
\Cref{rem:proof_hint_cauchy_schwarz}, definiteness (\cref{def:inner_product}) gives
\[
  0 \ \leq\ \langle u-\pi_v(u),\ u-\pi_v(u)\rangle
    \ =\ \langle u,u\rangle - 2\langle u,\pi_v(u)\rangle + \langle \pi_v(u),\pi_v(u)\rangle.
\]
Using $\pi_v(u) = \tfrac{\langle u,v\rangle}{\|v\|^2}v$ and bilinearity, the two right-hand terms
are
\begin{align*}
  \langle u,\pi_v(u)\rangle &= \frac{\langle u,v\rangle^2}{\|v\|^2}, \\[0.3ex]
  \langle \pi_v(u),\pi_v(u)\rangle
    &= \frac{\langle u,v\rangle^2}{\|v\|^4}\langle v,v\rangle
     = \frac{\langle u,v\rangle^2}{\|v\|^2}.
\end{align*}
Note that the two are equal, so substituting them leaves a single copy behind:
\[
  0 \ \leq\ \|u\|^2 - \frac{2\langle u,v\rangle^2}{\|v\|^2} + \frac{\langle u,v\rangle^2}{\|v\|^2}
    \ =\ \|u\|^2 - \frac{\langle u,v\rangle^2}{\|v\|^2}.
\]
This rearranges to $\langle u,v\rangle^2 \leq \|u\|^2\|v\|^2$, that is,
$|\langle u,v\rangle| \leq \|u\|\,\|v\|$. By definiteness, equality holds if and only if
$u - \pi_v(u) = 0$, which says precisely that $u$ is a scalar multiple of $v$.
\end{proof}

% Generator: Claude Opus 5 (High)
\begin{remark}[The same proof twice]
\label{rem:two_proofs_cauchy_schwarz}
This is the second proof of Cauchy--Schwarz in these notes;
\cref{prop:cauchy_schwarz_geometric_proof} gave another, following Diego Torres
Tejeda. The two look different at first. One minimises the quadratic
$f(\lambda) := \|v-\lambda w\|^2$ over $\lambda$, while the other expands
$\|u-\pi_v(u)\|^2 \geq 0$. But the minimiser found in the first is
$\lambda_0 = \tfrac{\langle v,w\rangle}{\|w\|^2}$, which is exactly the coefficient defining
$\pi_w(v)$. Consequently, the first proof \emph{derives} the projection by optimisation, and the
second \emph{postulates} it and checks that it works. Keeping both is worthwhile precisely because
the pair shows where the projection comes from. It is not a lucky guess, but the closest point of
the line to $v$.
\end{remark}


% --- exercises relocated here from the dissolved sheets ---

% Originally: content/exercise-sheets/week-03.tex, problem 3.7
% Source: exercises/Ex3_Analysis2_eng.pdf

\begin{exercise}[3.7 --- Cauchy--Schwarz]
\label{ex:3.7}
Let $V$ be a vector space over $\mathbb{R}$, let $\langle\cdot,\cdot\rangle$ be an inner product
on $V$, and let $\|\cdot\| : V \to \mathbb{R}$ be given by $\|v\| = \sqrt{\langle v,v\rangle}$.
Then the inequality
\begin{equation*}
|\langle v,w\rangle| \leq \|v\|\|w\| \tag{1}
\end{equation*}
holds for all $v, w \in V$. Furthermore, equality in (1) holds if and only if $v$ and $w$ are
linearly dependent.
\end{exercise}

\begin{remark}
Corsin devotes the whole second half of Friday (\cref{sec:cauchy_schwarz}) to this problem,
including the geometric reading and an explicit proof hint.
\end{remark}

% Originally: content/exercise-sheets/week-03.tex, problem 3.8
% Source: exercises/Ex3_Analysis2_eng.pdf

\begin{exercise}[3.8 --- Inner product and norm]
\label{ex:3.8}
Let $V$ be a vector space over $\mathbb{R}$, let $\langle -,-\rangle$ be an inner product on $V$.
The map defined by
$$\|\cdot\| : V \to \mathbb{R}, \qquad \|v\| = \sqrt{\langle v,v\rangle}$$
satisfies the triangular inequality and is a norm.
\end{exercise}

% Originally: content/exercise-sheets/week-03.tex, problem 3.9
% Source: exercises/Ex3_Analysis2_eng.pdf

\begin{exercise}[3.9 --- All norms are equivalent in $\mathbb{R}^n$]
\label{ex:3.9}
Let $|\cdot|$ denote the standard Euclidean norm in $\mathbb{R}^n$ and let
$f : \mathbb{R}^n \to [0,\infty)$ be another norm (that is, a function satisfying the properties
of \textit{Definition 9.89}).
\begin{enumerate}[label=\textbf{\arabic*.}]
  \item Expressing $x$ in a basis and using the ``abstract'' properties that $f$ must have, show
    that there is a constant $C_1 > 0$ such that $f(x) \leq C_1|x|$ for all $x \in \mathbb{R}^n$.
  \item Show that $f$ is continuous in $\mathbb{R}^n$ (with respect to the standard distance of
    $\mathbb{R}^n$!).
  \item Show that there is a number $c_2 > 0$ such that $f(x) \geq c_2$ for all $|x| = 1$.
  \item Conclude that $f(x) \geq c_2|x|$ for all $x \in \mathbb{R}^n$.
  \item Show that if $\tilde f$ is yet another norm, then there is $C > 0$ such that
    $C^{-1}f(x) \leq \tilde f(x) \leq Cf(x)$ for all $x \in \mathbb{R}^n$.
\end{enumerate}
\exhint[Official hints]{\textbf{3.9.2:} recall that from the definition it follows that
$|f(x) - f(y)| \leq f(x-y)$, and that Lipschitz functions are always continuous.
\textbf{3.9.3:} apply the Weierstrass theorem to $f$ on $S := \{x \in \mathbb{R}^n : |x| = 1\}$
and use that $f$, being a norm, is nondegenerate. \textbf{3.9.5:} if $f$ is equivalent to $|\cdot|$ and
$\tilde f$ is equivalent to $|\cdot|$, it follows by transitivity that $f$ is equivalent to
$\tilde f$.}
\end{exercise}

\begin{remark}
Hint 3.9.3 is exactly Corsin's reason \#2 for caring about compactness
(\cref{sec:why_compactness_matters}), namely the Weierstrass extreme value theorem applied on the
compact unit sphere.
\end{remark}

% Supplement: Analysis_II_Script_v1.pdf, p. 35
\begin{exercise}[The polynomial ring $\mathbb{R}{[}x{]}$ is infinite-dimensional]
\label{ex:polynomials_infinite_dimensional}
Show that the vector space $\mathbb{R}[x]$ of polynomials with real coefficients is infinite
dimensional, and exhibit a basis.
\end{exercise}

\begin{remark}
This is exactly the hypothesis \cref{ex:3.9} needs and cannot do without.
\Cref{rem:norm_equivalence_needs_finite_dimension} already gave $C^0([0,1],\mathbb{R})$ as an
infinite-dimensional space where norm equivalence fails. Note that $\mathbb{R}[x]$ is the
algebraically simplest such example, with an explicit basis and no analysis required at all.
\end{remark}

% Originally: content/exercise-sheets/week-03.tex, problem 3.10
% Source: exercises/Ex3_Analysis2_eng.pdf

\begin{exercise}[3.10 --- Hilbert--Schmidt norm of the composition]
\label{ex:3.10}
Take two linear functions $\varphi : \mathbb{R}^d \to \mathbb{R}^n$ and
$\psi : \mathbb{R}^n \to \mathbb{R}^m$, and denote with $\Phi$, $\Psi$ the matrices that represent
them in the canonical bases. Recall that the linear map
$\psi \circ \varphi : \mathbb{R}^d \to \mathbb{R}^m$ is represented in these bases by the matrix
$\Psi \cdot \Phi$. Show that
$$\|\Psi \cdot \Phi\| \leq \|\Psi\|\|\Phi\|,$$
where $\|\cdot\|$ is the Hilbert--Schmidt norm of a matrix (see \textit{Definition 9.96} in the
notes).

\exhint[Official hint]{Recall that $\|M\|^2$ is the sum of the squares of the entries of $M$.
Notice that $(\Psi\cdot\Phi)^i_j$ ($i$-th row, $j$-th column) is the scalar product of $\Psi^i$
and $\Phi_j$, which are vectors of $\mathbb{R}^n$. Apply Cauchy--Schwarz to each of them.}
\end{exercise}

% Originally: content/week-04.tex, \exercisesheet{4}
% Source: exercises/Ex4_Analysis2_eng.pdf

\begin{exercise}[4.3 --- $p$-norms \textnormal{(semi-important; parts 4 and 6 important)}]
\label{ex:4.3}
For $p \geq 1$ and $x \in \mathbb{R}^n$ define the \textbf{$p$-norm} of $x$ as
\[|x|_p := \Big(\sum_{i=1}^{n} |x_i|^p\Big)^{1/p}.\]
\begin{enumerate}[label=\textbf{\arabic*.}]
  \item For $n = 2$ and $p = 1, 2, 10$ sketch the sets $\{x \in \mathbb{R}^2 : |x|_p \leq 1\}$.
  \item For a given $x \in \mathbb{R}^n$, compute the limit
    $|x|_\infty := \lim_{p\to\infty}|x|_p$.
  \item Using an appropriate inequality that you have seen in class, prove that
    \[\Big(\sum_{i=1}^{n} a_i^{p-1}b_i\Big)^2 \leq \Big(\sum_{i=1}^{n}a_i^p\Big)\Big(\sum_{i=1}^{n}a_i^{p-2}b_i^2\Big),\]
    whenever $a_i, b_i$ are $n$-tuples of positive numbers.
  \item Fix $x, y \in \mathbb{R}^n$ and consider the function $f : [0,1] \to \mathbb{R}$ defined
    as
    \begin{equation*}
    f(t) := |tx + (1-t)y|_p = \Big(\sum_{i=1}^{n}|tx_i + (1-t)y_i|^p\Big)^{1/p}. \tag{1}
    \end{equation*}
    Show that $f$ is convex. You may assume that the coordinates of $x$ and $y$ are all strictly
    positive and use the inequality of the previous point.
  \item Deduce from the previous point that the triangular inequality holds, i.e.\
    $|x+y|_p \leq |x|_p + |y|_p$ for all $x,y \in \mathbb{R}^n$.
  \item What happens for $p \in (0,1)$?
\end{enumerate}
\exhint[Official hints]{\textbf{4.3.3:} use Cauchy--Schwarz and the fact that
$a_i^{p-1}b_i = a_i^{p/2}\cdot a_i^{(p-2)/2}b_i$. \textbf{4.3.4:} show that $f''(t) \geq 0$; it
is not the simplest derivative, but if you get it right the inequality in 4.3.3 will be just what
you need. \textbf{4.3.5:} write the convexity inequality between $x$, $y$ and the middle point
$(x+y)/2$. To get the general case from the one with positive coordinates, observe that for
$x \in \mathbb{R}^n$, denoting $\hat x := (|x_1|,\dots,|x_n|)$, we have
$|x+y|_p \leq |\hat x + \hat y|_p$ and $|\hat x|_p = |x|_p$ for all $x,y \in \mathbb{R}^n$.}
\end{exercise}

% Originally: content/exercise-sheets/week-04.tex, problem 4.4 --- never previously built
% Source: exercises/Ex4_Analysis2_eng.pdf

\begin{exercise}[4.4 --- $p$-means \textnormal{(optional)}]
\label{ex:4.4}
For $x \in \mathbb{R}^n$ with positive coordinates and $p \neq 0$ define the \textbf{$p$-mean} as
$$\mu_p(x) := \left(\frac{x_1^p + \dots + x_n^p}{n}\right)^{1/p}.$$
\begin{enumerate}[label=\textbf{\arabic*.}]
  \item Compute the limits $p \to \pm\infty$, $p \to 0$ and define accordingly
    $\mu_{-\infty}(x)$, $\mu_0(x)$, $\mu_{+\infty}(x)$.
  \item For any $n$-tuple of numbers $a_i > 0$ show that
    $$\sum_{i=1}^{n} \frac{a_i\log(a_i)}{a_1 + \dots + a_n} \geq \log\left(\frac{a_1 \dots + a_n}{n}\right).$$
  \item For a fixed $x$, show that the function $f : \mathbb{R}\to\mathbb{R}$, given by
    $f(t) := \mu_t(x)$, is continuous and increasing.
  \item Prove the Arithmetic--Geometric inequality and Arithmetic--Quadratic inequality:
    $$n(x_1x_2\cdots x_n)^{1/n} \leq x_1 + \dots + x_n, \qquad (x_1+\dots+x_n)^2 \leq n(x_1^2+\dots+x_n^2).$$
  \item \textbf{(*)} Is $f$ continuously differentiable in the whole $\mathbb{R}$?
\end{enumerate}
\exhint[Official hints]{\textbf{4.4.2:} combine the concavity of $\log(\cdot)$ and the
Cauchy--Schwarz inequality; use the generalisation of the concavity inequality: for any concave
$f : \mathbb{R}\to\mathbb{R}$,
$f(\lambda_1x_1 + \dots + \lambda_nx_n) \geq \lambda_1f(x_1)+\dots+\lambda_nf(x_n)$ for all
$x_i \in \mathbb{R}$, $0 \leq \lambda_i \leq 1$ with $\lambda_1+\dots+\lambda_n = 1$.
\textbf{4.4.3:} it might be more convenient to work with $\log f(t)$.}
\end{exercise}

% Originally: new section, 2026-08-09. Not a move --- this material had no home in the document
% before. Exercise 3.9 asked the reader to prove that all norms on R^n are equivalent, and the
% solution in 99-solutions.tex proved it, but the statement itself was nowhere in the body, so
% nothing could \cref it and the reader met the result only as homework.
% Supplement: lec_notes.pdf, pp. 31--33 (Definition 9.105, Theorem 9.107, Corollary 9.109)
% Generator: Claude Opus 5 (High)

\section{Equivalence of norms}
\label{sec:equivalence_of_norms}

The previous sections put several norms on the same vector space: the Euclidean norm, the
$p$-norms, the supremum norm. A natural worry follows. Convergence, openness, continuity and
compactness were all defined from a metric, and on a normed space that metric came from the norm.
So does changing the norm change which sequences converge and which sets are open?

On $\mathbb{R}^n$ the answer is no, and the reason is worth isolating, because it is one of the
few places in this course where finite-dimensionality is doing essential work.

% Supplement: lec_notes.pdf, p. 31 (Definition 9.105)
\begin{definition}[Equivalent norms]
\label{def:equivalent_norms}
Let $V$ be a vector space over $\mathbb{R}$ and let $\|\cdot\|_1$ and $\|\cdot\|_2$ be two norms
on $V$. They are called \newterm{equivalent} \germanfor{äquivalent}, or \emph{comparable}, if
there are constants $A, B > 0$ with
\[\|v\|_1 \leq A\|v\|_2 \qquad\text{and}\qquad \|v\|_2 \leq B\|v\|_1 \qquad\text{for all } v \in V.\]
\end{definition}

\begin{remark}[Why two inequalities and not one]
A single inequality $\|v\|_1 \leq A\|v\|_2$ says only that $\|\cdot\|_2$-convergence implies
$\|\cdot\|_1$-convergence: it makes the first norm the coarser of the two. Equivalence is the
two-sided statement, and it is exactly what is needed for the two norms to have the same null
sequences, and hence the same convergent sequences, open sets and compact sets. That equivalence
of norms really is an equivalence relation is part 5 of \cref{ex:3.9}.
\end{remark}

% Supplement: lec_notes.pdf, p. 32 (Example 9.106)
\begin{example}[The $1$-norm against the maximum norm]
\label{ex:one_norm_vs_max_norm}
On $\mathbb{R}^n$,
\[\|v\|_\infty \leq \|v\|_1 \qquad\text{and}\qquad \|v\|_1 \leq n\|v\|_\infty,\]
so the two are equivalent with $A = 1$ and $B = n$. The left inequality holds because the largest
$|v_i|$ is one of the terms of the sum $\sum_i|v_i|$; the right because each of the $n$ terms of
that sum is at most $\max_i|v_i|$.

The constant $n$ on the right is sharp, as $v = (1,1,\dots,1)$ shows, and it is the first sign of
what goes wrong later: the constants in \cref{thm:all_norms_equivalent} are allowed to depend on
the dimension, and there is no dimension to depend on in an infinite-dimensional space.
\end{example}

% Supplement: lec_notes.pdf, p. 32 (Theorem 9.107)
\begin{theorem}[All norms on $\mathbb{R}^n$ are equivalent]
\label{thm:all_norms_equivalent}
Any two norms on $\mathbb{R}^n$ are equivalent in the sense of \cref{def:equivalent_norms}.
\end{theorem}

\begin{proof}
This is \cref{ex:3.9}, whose five parts are exactly the five steps of the argument, and it is
proved in full in the solutions to that exercise. In outline, with $|\cdot|$ the Euclidean norm
and $f$ an arbitrary norm: expanding $x$ in the standard basis and applying the triangle
inequality, homogeneity and Cauchy--Schwarz gives $f(x) \leq C_1|x|$; that bound makes $f$
Lipschitz, hence continuous, for the Euclidean metric; Weierstrass then gives $f$ a positive
minimum $c_2$ on the Euclidean unit sphere, which is compact by Heine--Borel
(\cref{thm:heine_borel}); and homogeneity spreads that bound from the sphere to all of
$\mathbb{R}^n$, giving $f(x) \geq c_2|x|$.
\end{proof}

\begin{remark}[Where each hypothesis is used]
\label{rem:where_finite_dimension_enters}
Both halves of the proof use finite-dimensionality, and they use it differently. The first
inequality expands $x$ in a \emph{finite} basis, so the sum of $n$ terms it produces is a sum and
not a series. The second applies Weierstrass on the unit sphere, which needs that sphere to be
\emph{compact}, and by Heine--Borel it is, in $\mathbb{R}^n$ and nowhere else.
\Cref{rem:norm_equivalence_needs_finite_dimension} exhibits the failure in
$C^0([0,1],\mathbb{R})$, and \cref{ex:polynomials_infinite_dimensional} gives the algebraically
simplest infinite-dimensional space to have in mind.
\end{remark}

% Supplement: lec_notes.pdf, p. 33 (9.108)
\begin{corollary}[Any finite-dimensional real vector space]
\label{cor:all_norms_equivalent_finite_dim}
Let $V$ be a real vector space of finite dimension $n$. Then any two norms on $V$ are equivalent.
\end{corollary}

\begin{proof}
Fix a basis $\mathcal{B} = \{e_1,\dots,e_n\}$ of $V$. The coordinate map
$\iota_{\mathcal{B}} : \mathbb{R}^n \to V$, $(x_1,\dots,x_n) \mapsto x_1e_1+\dots+x_ne_n$, is a
linear isomorphism, so a norm $\|\cdot\|$ on $V$ pulls back to the norm
$x \mapsto \|\iota_{\mathcal{B}}(x)\|$ on $\mathbb{R}^n$, and the two-sided inequalities of
\cref{def:equivalent_norms} are preserved in both directions under this identification. Applying
\cref{thm:all_norms_equivalent} on $\mathbb{R}^n$ and transporting back gives the claim.
\end{proof}

% Supplement: lec_notes.pdf, p. 33 (Corollary 9.109)
\begin{corollary}[All norms give the same topology]
\label{cor:norms_same_topology}
Topological properties of a subset $E \subseteq \mathbb{R}^n$ (openness, closedness, convergence
of sequences, compactness, connectedness) do not depend on which norm was used to define the
metric.
\end{corollary}

\begin{proof}
Let $d_1, d_2$ be the metrics induced by two norms. By \cref{thm:all_norms_equivalent} there is a
constant $C > 0$ with $C^{-1}d_1(x,y) \leq d_2(x,y) \leq Cd_1(x,y)$ for all $x,y$, since
$d_i(x,y) = \|x-y\|_i$ and the norms are equivalent. Hence every $d_1$-ball contains a
$d_2$-ball about the same centre and conversely, which by \cref{def:open_closed} makes a set open
for $d_1$ exactly when it is open for $d_2$. The two metrics therefore generate the same
topology, and every property defined purely in terms of the open sets agrees.
\end{proof}

\begin{ainote}[What this licenses, and what it does not]
\label{ainote:norm_equivalence_scope}
\Cref{cor:norms_same_topology} is the reason this document is free to say \qt{open}, \qt{compact}
or \qt{convergent} in $\mathbb{R}^n$ without ever naming a norm, and it is worth knowing that the
licence has been earned rather than assumed.

Three things it does \emph{not} give you.
\begin{itemize}
  \item It is about \emph{topological} properties only. Metric quantities are not preserved:
    the unit ball genuinely changes shape with the norm (see the three drawn in
    \cref{sec:unit_balls}), distances change, and so do Lipschitz constants. A $1$-Lipschitz map
    for one norm need not be $1$-Lipschitz for another, only Lipschitz with some other constant.
  \item It says nothing in infinite dimensions, where the conclusion is false
    (\cref{rem:where_finite_dimension_enters}).
  \item It does not make the choice of norm irrelevant in practice. Estimates are still written
    in whichever norm makes the constants come out cleanly, which is why the supremum norm is the
    one used on $C(K,\mathbb{R}^m)$ in \cref{sec:ascoli_arzela} even though that space is
    infinite-dimensional and the theorem above does not reach it.
\end{itemize}
\end{ainote}

% Originally: the \section{Solutions} of content/week-02.tex and content/week-03.tex

\section{Solutions}
\label{sec:metric_spaces_solutions}

\begin{ainote}[How the solutions in this document were produced]
Where a tutor solves a problem in their own notes, that solution is kept \emph{inline}, directly
after the statement, mirroring how they present it --- each such case carries a
\texttt{\% Kept inline} comment naming the source page. Everything else is collected in a
\emph{Solutions} section like this one at the end of the chapter that hosts the problem. Those
were worked independently and only then checked against the official
\texttt{exercises/SolN\_Analysis2\_eng.pdf}; wherever the comparison turned up anything worth
flagging, it is flagged in an AI-Note next to the solution.
\end{ainote}

\begin{exercisesolution}[Proof of \cref{ex:ai_metric_equivalence}]
% Generator: Gemini 3.6 Flash (Medium)
For any $x = (x_1, x_2) \in \mathbb{R}^2$, we have $\max\{|x_1|, |x_2|\} \leq |x_1| + |x_2|$, giving $d_3(x,0) \leq d_1(x,0)$. Conversely, since $|x_1| \leq \max\{|x_1|, |x_2|\}$ and $|x_2| \leq \max\{|x_1|, |x_2|\}$, summing these yields $|x_1| + |x_2| \leq 2 \max\{|x_1|, |x_2|\}$, i.e., $d_1(x,0) \leq 2 \, d_3(x,0)$. Both bounds are sharp: equality on the left holds for $x=(1,0)$, and equality on the right holds for $x=(1,1)$.
\end{exercisesolution}

\begin{exercisesolution}[Solution to \cref{ex:2.1}]
% Generator: Claude Opus 5 (High)
Throughout, non-negativity and symmetry are immediate in every part and are not belaboured; the
interest is in \textbf{definiteness} and the \textbf{triangle inequality}.
\begin{enumerate}[label=\textbf{\arabic*.}]
  \item \textbf{Yes.} Boundedness of $f$ and $g$ makes the supremum finite, so $d$ really maps
    into $\mathbb{R}$. Definiteness: $d(f,g)=0$ forces $|f(x)-g(x)| = 0$ for every $x$, i.e.\
    $f=g$. For the triangle inequality, estimate pointwise and only then take the supremum:
    \[|f(x)-h(x)| \leq |f(x)-g(x)| + |g(x)-h(x)| \leq d(f,g) + d(g,h) \quad \text{for all } x,\]
    so the right-hand side is an upper bound for $|f-h|$ and hence dominates $d(f,h)$. This is
    \cref{ex:sup_metric_is_metric} verbatim, with $[0,1]$ replaced by an arbitrary set $X$ and
    continuity replaced by boundedness --- note that it was boundedness, not continuity, that
    the argument there actually used.
  \item \textbf{Yes.} Let $\varphi(x) := 1/x$, an injective map $\mathbb{Q}_+ \to \mathbb{Q}_+$,
    and observe that $d(x,y) = |\varphi(x)-\varphi(y)|$. Symmetry and the triangle inequality
    are inherited from $|\cdot|$ without any work, and definiteness is exactly injectivity of
    $\varphi$: from $d(x,y)=0$ we get $1/x = 1/y$, hence $x=y$. This is the same pullback
    construction as \cref{ex:completeness_not_topological}, and the same warning applies ---
    pulling a metric back along an injection \emph{always} produces a metric, which is why the
    interesting question about such a $d$ is never whether it is one, but whether it is
    complete.
  \item \textbf{No} --- the triangle inequality fails. Ignore the second coordinate, as the
    official hint suggests, and take three collinear points $x := (0,0)$, $y := (1,0)$ and
    $z := (2,0)$:
    \[d(x,z) = 2^2 = 4 \quad > \quad d(x,y) + d(y,z) = 1 + 1 = 2.\]
    Squaring produces the same failure as in
    \cref{ex:metric_axioms_are_needed}\textbf{(c)}: a function growing faster than linearly
    makes one long step more expensive than two short ones.
  \item \textbf{Yes.} A sum of two metrics acting on separate coordinates is a metric, so
    everything reduces to the first coordinate, where the claim is that
    $(s,t) \mapsto |s-t|^{1/2}$ satisfies the triangle inequality on $\mathbb{R}$. That follows
    from \textbf{subadditivity of the square root}, $\sqrt{a+b} \leq \sqrt a + \sqrt b$ for
    $a,b \geq 0$, which holds because squaring the right-hand side gives
    $a+b+2\sqrt{ab} \geq a+b$. Hence
    \[|x_1-z_1|^{1/2} \leq \big(|x_1-y_1| + |y_1-z_1|\big)^{1/2}
      \leq |x_1-y_1|^{1/2} + |y_1-z_1|^{1/2}.\]
    Contrast this with part 3: $\sqrt{\cdot}$ is concave and vanishes at $0$, hence subadditive,
    while $(\cdot)^2$ is convex, hence super-additive. That single distinction settles both
    parts.
  \item \textbf{Yes.} Write $A := X-Y$. Then
    \[\Tr(A\transp A) = \sum_{i=1}^n (A\transp A)_{ii} = \sum_{i,j=1}^n A_{ji}^2,\]
    so $d(X,Y)$ is nothing but the Euclidean distance between $X$ and $Y$ regarded as vectors of
    $\mathbb{R}^{n^2}$, and it is a metric by
    \cref{prop:induced_norm_and_metric}\textbf{(b)}. The associated norm is the
    \newterm{Frobenius norm} (\germanterm{Frobeniusnorm}), also called the Hilbert--Schmidt norm.
  \item \textbf{No} --- and here it is \textbf{definiteness} that fails, not the triangle
    inequality. Following the hint, suppose $A := X-Y$ is \emph{anti-symmetric},
    $A\transp = -A$. For any $v$ the number $v\transp Av$ is a $1\times1$ matrix, hence equal to
    its own transpose:
    \[v\transp A v = (v\transp A v)\transp = v\transp A\transp v = -\,v\transp A v,\]
    so $v\transp Av = 0$ for \emph{every} $v$. Taking $n=2$ with
    \[X := \begin{pmatrix} 0 & 1 \\ -1 & 0\end{pmatrix}, \qquad
      Y := \begin{pmatrix} 0 & 0 \\ 0 & 0\end{pmatrix}\]
    gives $X \neq Y$ but $d(X,Y) = 0$. Every other axiom does hold, so $d$ is a
    \newterm{pseudo-metric} (\germanterm{Pseudometrik}): it cannot separate two matrices
    differing by an anti-symmetric one. Equivalently, it restricts to a genuine metric on the
    symmetric matrices.
  \item \textbf{Yes.} Three things need checking, and the first is the one that is easy to skip.

    \textit{Well defined.} Replacing $x$ by $x+(m,n)$ and $y$ by $y+(m',n')$ shifts $x-y$ by the
    integer vector $(m-m',\,n-n')$, which merely re-indexes the infimum. So $d([x],[y])$ does
    not depend on the representatives chosen.

    \textit{The infimum is a minimum.} Only finitely many integer vectors $(k,h)$ satisfy
    $\lVert x-y+(k,h)\rVert \leq \lVert x-y\rVert$, because $\mathbb{Z}^2$ is discrete; the
    infimum is therefore taken over a finite non-empty set of candidates and is attained. This
    is what the official hint asks you to settle first, and it is what makes the next step
    legitimate.

    \textit{Definiteness.} If $d([x],[y]) = 0$ then, the minimum being attained,
    $x-y+(k,h) = 0$ for some integers $k,h$, that is $x \sim y$ and $[x] = [y]$. Note exactly
    where \qt{the infimum is a minimum} was needed: an infimum equal to $0$ but not attained
    would say nothing at all.

    \textit{Triangle inequality.} Choose minimisers $(k_1,h_1)$ for the pair $(x,y)$ and
    $(k_2,h_2)$ for $(y,z)$. Their sum is a legitimate candidate for the pair $(x,z)$, so
    \begin{align*}
      d([x],[z]) &\leq \big\lVert x-z+(k_1{+}k_2,\ h_1{+}h_2)\big\rVert \\
      &= \big\lVert \big(x-y+(k_1,h_1)\big) + \big(y-z+(k_2,h_2)\big)\big\rVert \\
      &\leq d([x],[y]) + d([y],[z]).
    \end{align*}
\end{enumerate}
\end{exercisesolution}

\begin{ainote}
Parts 3, 4 and 6 are the ones worth revisiting, because each breaks (or rescues) a different
axiom and together they cover every way the definition can fail. Part 3 kills the triangle
inequality by growing too fast; part 4 rescues it by growing slowly enough; part 6 satisfies the
triangle inequality perfectly and fails on definiteness instead. Read alongside
\cref{ex:metric_axioms_are_needed}, they make the case that the three axioms are genuinely
independent.
\end{ainote}

\begin{exercisesolution}[Solution to \cref{ex:2.5}]
% Generator: Claude Opus 5 (High)
Write $s := d_X(x,x')$ and $t := d_Y(y,y')$ throughout, so that the three candidates are
$d_1 = \max\{s,t\}$, $d_2 = s+t$ and $d_3 = \sqrt{s^2+t^2}$ --- that is, the supremum, Manhattan
and standard metrics of \cref{ex:metric_examples} applied to the pair $(s,t)$. Almost everything
below is that one observation.
\begin{enumerate}[label=\textbf{\arabic*.}]
  \item \textbf{All three are metrics.} Each vanishes exactly when $s = t = 0$, i.e.\ when
    $x=x'$ and $y=y'$, which is definiteness; symmetry is inherited from $d_X$ and $d_Y$. For
    the triangle inequality, set
    \[a := \big(d_X(x,x'),\,d_Y(y,y')\big), \quad
      b := \big(d_X(x',x''),\,d_Y(y',y'')\big), \quad
      c := \big(d_X(x,x''),\,d_Y(y,y'')\big),\]
    all lying in $[0,\infty)^2$. The triangle inequalities in $X$ and in $Y$ say precisely that
    $c \leq a+b$ \emph{componentwise}. Each of $\max$, the sum and $\lVert\cdot\rVert_2$ is
    monotone on $[0,\infty)^2$ and subadditive, so
    \[N(c) \ \leq\ N(a+b) \ \leq\ N(a) + N(b)\]
    for $N \in \{\max,\ \textstyle\sum,\ \lVert\cdot\rVert_2\}$ --- which is the required
    inequality in all three cases at once.
  \item \textbf{Equivalence.} For all $s,t \geq 0$,
    \[\max\{s,t\} \ \leq\ \sqrt{s^2+t^2} \ \leq\ s+t \ \leq\ 2\max\{s,t\},\]
    i.e.\ $d_1 \leq d_3 \leq d_2 \leq 2d_1$; the middle inequality holds because
    $(s+t)^2 = s^2+t^2+2st \geq s^2+t^2$. Hence $C := 2$ satisfies the chain in the statement:
    $d_1 \leq 2d_2$, then $2d_2 \leq 4d_3$ (since $d_2 \leq \sqrt2\,d_3$ by Cauchy--Schwarz,
    and $\sqrt2 \leq 2$), and finally $4d_3 \leq 8d_1$ (since $d_3 \leq \sqrt2\,d_1$).
  \item \textbf{Convergence is coordinatewise.} From part 2,
    \[\max\big\{d_X(x_n,x),\ d_Y(y_n,y)\big\}
      \ \leq\ d_3\big((x_n,y_n),(x,y)\big)
      \ \leq\ d_X(x_n,x) + d_Y(y_n,y).\]
    If the middle term tends to $0$, the left inequality forces both coordinates to converge; if
    both coordinates converge, the right inequality forces the middle term to $0$. This is
    \cref{prop:Rn_coordinatewise_complete}\textbf{(a)} with $n = 2$ and the two factors allowed
    to be arbitrary metric spaces. By part 2 the same conclusion holds for $d_1$ and $d_2$,
    since equivalent metrics have the same convergent sequences.
\end{enumerate}
\end{exercisesolution}

\begin{exercisesolution}[Proof of \cref{ex:3.7}]
% Generator: Claude Opus 5 (High)
This is \cref{thm:cauchy_schwarz}, proved in \cref{sec:cauchy_schwarz} above by expanding
$\lVert u - \pi_v(u)\rVert^2 \geq 0$, including the equality case. A second, independent proof
--- minimising $\lambda \mapsto \lVert v-\lambda w\rVert^2$ --- is
\cref{prop:cauchy_schwarz_geometric_proof} in Week 2. \Cref{rem:two_proofs_cauchy_schwarz}
explains why the two are the same computation approached from opposite ends.

\begin{remark}[A third route: the discriminant]
The official solution gives a third proof, worth recording because it is the one that
generalises most mechanically. By definiteness, the quadratic polynomial
\[\lambda \ \longmapsto\ \langle \lambda v + w,\ \lambda v + w\rangle
  = \lambda^2\langle v,v\rangle + 2\lambda\langle v,w\rangle + \langle w,w\rangle\]
is non-negative for every real $\lambda$. A real quadratic with positive leading coefficient is
non-negative exactly when its discriminant is $\leq 0$, i.e.
\[\big(2\langle v,w\rangle\big)^2 - 4\langle v,v\rangle\langle w,w\rangle \leq 0,\]
which rearranges to the inequality. Equality forces the discriminant to vanish, so the
polynomial has a real root $\lambda_0$, and definiteness turns
$\langle \lambda_0 v + w, \lambda_0 v + w\rangle = 0$ into $\lambda_0 v + w = 0$ --- the linear
dependence. Note that this is the previous proof with the geometry stripped out: $\lambda_0$ is
again the projection coefficient, now arriving as a root rather than as a minimiser.
\end{remark}
\end{exercisesolution}

\begin{exercisesolution}[Proof of \cref{ex:3.8}]
% Generator: Claude Opus 5 (High)
This is \cref{prop:induced_norm_and_metric}\textbf{(a)}, proved in Week 2 immediately after
\cref{prop:cauchy_schwarz_geometric_proof} --- which is exactly the input the triangle inequality
needs, and the reason the two problems appear adjacently on this sheet. In one line: expanding
$\lVert u+v\rVert^2 = \lVert u\rVert^2 + 2\langle u,v\rangle + \lVert v\rVert^2$ and bounding the
middle term by $2\lVert u\rVert\lVert v\rVert$ turns the right-hand side into
$(\lVert u\rVert + \lVert v\rVert)^2$.
\end{exercisesolution}

\begin{exercisesolution}[Solution to \cref{ex:3.9}]
% Generator: Claude Opus 5 (High)
Write $e_1,\dots,e_n$ for the standard basis and $|\cdot|$ for the Euclidean norm.
\begin{enumerate}[label=\textbf{\arabic*.}]
  \item Expand $x = \sum_{i=1}^n x_ie_i$ and use the triangle inequality and homogeneity of $f$
    (\cref{def:norm}):
    \[f(x) \ \leq\ \sum_{i=1}^n f(x_ie_i) \ =\ \sum_{i=1}^n |x_i|\,f(e_i)
      \ \leq\ \Big(\max_i f(e_i)\Big)\sum_{i=1}^n|x_i|
      \ \leq\ \Big(\max_i f(e_i)\Big)\sqrt n\,|x|,\]
    the last step by Cauchy--Schwarz applied to $(|x_1|,\dots,|x_n|)$ and $(1,\dots,1)$. So
    $C_1 := \sqrt n \max_i f(e_i)$ works, and $C_1 > 0$ because $e_i \neq 0$ forces
    $f(e_i) > 0$ by definiteness.
  \item The reverse triangle inequality for $f$ --- namely $f(x) \leq f(y) + f(x-y)$ and its
    mirror image, giving $|f(x)-f(y)| \leq f(x-y)$ --- combines with part 1 to yield
    \[|f(x)-f(y)| \ \leq\ f(x-y) \ \leq\ C_1|x-y|.\]
    So $f$ is $C_1$-Lipschitz with respect to the \emph{Euclidean} distance
    (\cref{def:lipschitz_continuous}), hence continuous (\cref{rem:continuity_hierarchy}). Note
    which distance is which here: the claim is continuity of $f$ as a function on the metric
    space $(\mathbb{R}^n, |\cdot|)$, and part 1 is what makes that meaningful.
  \item The unit sphere $S := \{x : |x| = 1\}$ is closed (preimage of $\{1\}$ under the
    continuous $|\cdot|$) and bounded, hence \textbf{compact} by \cref{thm:heine_borel}. By part
    2 the function $f$ is continuous on $S$, so it attains a minimum there
    (\cref{item:extreme_value_theorem}), say $c_2 := f(x_0)$ with $|x_0| = 1$. Since $x_0 \neq 0$,
    definiteness gives $c_2 = f(x_0) > 0$.
  \item For $x \neq 0$ the vector $x/|x|$ lies in $S$, so $f(x/|x|) \geq c_2$, and homogeneity
    gives
    \[f(x) = |x|\,f\!\left(\frac{x}{|x|}\right) \ \geq\ c_2\,|x|.\]
    For $x = 0$ both sides vanish. Together with part 1, $c_2|x| \leq f(x) \leq C_1|x|$.
  \item Apply parts 1 and 4 to both norms: $c_2|x| \leq f(x) \leq C_1|x|$ and
    $\tilde c_2|x| \leq \tilde f(x) \leq \tilde C_1|x|$, all four constants positive. Then
    \[\tilde f(x) \ \leq\ \tilde C_1|x| \ \leq\ \frac{\tilde C_1}{c_2}\,f(x),
      \qquad\qquad
      \tilde f(x) \ \geq\ \tilde c_2|x| \ \geq\ \frac{\tilde c_2}{C_1}\,f(x),\]
    so $C := \max\big\{\tilde C_1/c_2,\ C_1/\tilde c_2\big\}$ satisfies
    $C^{-1}f \leq \tilde f \leq Cf$. Equivalence of norms is thus an equivalence relation, and
    the Euclidean norm serves as the common reference point.
\end{enumerate}

\begin{remark}[Where finite dimension is indispensable]
\label{rem:norm_equivalence_needs_finite_dimension}
The result is false in infinite dimensions, and the proof shows exactly which step breaks: part 1
expands $x$ in a \emph{finite} basis and part 3 uses compactness of the unit sphere. In an
infinite-di\-men\-sional normed space the closed unit ball is never compact, and one can write
down inequivalent norms --- for instance $\lVert\cdot\rVert_\infty$ and
$\int_0^1|\cdot|$ on $C^0([0,1])$, where a tall thin spike has supremum norm $1$ and integral
norm arbitrarily small. Every appearance of \qt{all norms on $\mathbb{R}^n$ are equivalent} later
in the course silently uses $n < \infty$.
\end{remark}
\end{exercisesolution}

\begin{exercisesolution}[Proof of \cref{ex:3.10}]
% Generator: Claude Opus 5 (High)
Write $\Psi^i \in \mathbb{R}^n$ for the $i$-th \emph{row} of $\Psi$ and
$\Phi_j \in \mathbb{R}^n$ for the $j$-th \emph{column} of $\Phi$, so that by the definition of
matrix multiplication
\[(\Psi\Phi)_{ij} = \langle \Psi^i,\ \Phi_j\rangle.\]
Cauchy--Schwarz (\cref{thm:cauchy_schwarz}) applied to this pair of vectors of $\mathbb{R}^n$
gives $\big|(\Psi\Phi)_{ij}\big|^2 \leq \lVert\Psi^i\rVert^2\,\lVert\Phi_j\rVert^2$. Summing over
all entries and factorising the resulting double sum,
\[\lVert\Psi\Phi\rVert^2
  = \sum_{i=1}^m\sum_{j=1}^d \big|(\Psi\Phi)_{ij}\big|^2
  \ \leq\ \sum_{i=1}^m\sum_{j=1}^d \lVert\Psi^i\rVert^2\lVert\Phi_j\rVert^2
  = \left(\sum_{i=1}^m\lVert\Psi^i\rVert^2\right)\left(\sum_{j=1}^d\lVert\Phi_j\rVert^2\right)
  = \lVert\Psi\rVert^2\,\lVert\Phi\rVert^2,\]
using that the Hilbert--Schmidt norm squared is the sum of the squares of all entries --- read
row by row for $\Psi$ and column by column for $\Phi$. Taking square roots gives
$\lVert\Psi\Phi\rVert \leq \lVert\Psi\rVert\lVert\Phi\rVert$.

\begin{ainote}
The Hilbert--Schmidt norm is the same object as the Frobenius norm met in
\cref{ex:2.1} part 5 on the previous sheet, where it was shown to be the Euclidean norm of the
matrix read as a vector of $\mathbb{R}^{n^2}$. The inequality just proved says that this norm is
\newterm{submultiplicative} (\germanterm{submultiplikativ}), which is exactly the property that
makes it useful for estimating compositions of linear maps --- and it is what will let
$\lVert Df\rVert$ be handled like an absolute value in the mean value inequality of Week 4
(\cref{rem:mvt_fails_vector_valued}).
\end{ainote}
\end{exercisesolution}


% ============================================================
% Chapter 3 --- Open and closed sets
% Originally: content/week-02.tex, \section{Open and closed sets} plus the Friday
%             unit-ball warm-up (Corsin Week 2)
% ============================================================
\chapter{Open and Closed Sets}
\label{ch:open_closed}

% Transition: Claude Opus 5 (Medium)
With a distance in hand, the next question is which subsets a metric singles out. This chapter
introduces open balls, open and closed sets, and the interior, closure and boundary of a set, and
ends by looking at what the unit ball of each of the three standard metrics actually looks like
--- three visibly different shapes that nonetheless induce exactly the same open sets.

% Originally: content/week-02.tex, \section{Open and closed sets} (Corsin Week 2, Monday)

% Retitled from "Open and closed sets" now that the chapter carries that name; the label is
% unchanged, so every existing \cref still resolves.
\section{Open balls, open and closed sets}
\label{sec:open_closed_sets}

\begin{definition}[Open ball]
\label{def:open_ball}
Let $(X,d)$ be a metric space, let $x \in X$, and let $r \in \mathbb{R}$. Then we define the
\newterm{open ball} of radius $r$ as
\[B_r(x) := \{\,\boxed{y \in X}\, : d(y,x) < r\,\}\]
\end{definition}

\begin{importantremark}
Corsin marks the boxed part in red as an \qt{important subtlety}: the ambient set $X$ in
$y \in X$ is what makes openness relative to the space, and it is exactly what the common
misconception on \cref{rem:common_misconception} turns on.
\end{importantremark}

\begin{ainote}
\Cref{def:open_ball} is stated for $r \in \mathbb{R}$ rather than $r > 0$. Harmless here, but
loose.
\end{ainote}

% Moved here from before \cref{def:open_ball} (review item D1): it uses $B_r(x)$,
% "open" and "clopen", none of which were defined at its former position.
\begin{aiexample}[Open Balls in the Discrete Metric]
% Generator: Gemini 3.6 Flash (Medium)
Let $(X, d_{\text{disc}})$ be a discrete metric space. The open ball $B_r(x)$ centered at
$x \in X$ with radius $r > 0$ is
\[
B_r(x) = \begin{cases} \{x\} & \text{if } 0 < r \leq 1, \\ X & \text{if } r > 1. \end{cases}
\]
Hence every singleton $\{x\}$ is open in the discrete metric, and consequently every subset of
$X$ is both open and closed (\newterm{clopen}). This is exactly what \Cref{ex:X_empty_clopen}
below asks you to verify for $X$ and $\emptyset$ in a general metric space.
\end{aiexample}

% Generator: Claude Opus 5 (Medium)
\begin{remark}[Two degenerate balls worth knowing]
Taking \cref{def:open_ball} literally at the edges of its range is a cheap way to test whether
you have read it correctly.
\begin{itemize}
  \item If $r \leq 0$, then $B_r(x) = \emptyset$, and \emph{not} $\{x\}$. The condition is
    $d(y,x) < r$, strictly, and $d(x,x) = 0 \not< 0$. So even the centre is missing from
    $B_0(x)$.
  \item In the discrete metric $B_1(x) = \{x\}$, while $B_{1.01}(x) = X$: a ball can jump from a
    single point to the entire space without passing through anything in between. Radii do not
    have to behave like radii.
\end{itemize}
These are the two cases that make \qt{shrink the ball a little} arguments go wrong if applied
without thinking.
\end{remark}

% Generator: Claude Opus 5 (Medium)
\begin{aiexercise}[Open balls are open]
\label{ex:ai_open_balls_are_open}
The name promises it, but \cref{def:open_ball} does not say it. Show that in any metric space
$(X,d)$ the ball $B_r(x)$ is an open set in the sense of \cref{def:open_closed}.

\textit{Hint:} given $y \in B_r(x)$, the number $\varepsilon := r - d(y,x)$ is strictly positive.
Draw the picture, then let the triangle inequality do the work.
\end{aiexercise}

\begin{definition}[Open and closed sets]
\label{def:open_closed}
Then, we call a subset $U \subseteq X$ \newterm{open} (\germanterm{offen}) if and only if
\[\forall y \in U \ \exists \varepsilon > 0 \text{ such that } B_\varepsilon(y) \subseteq U.\]
And we call $A \subseteq X$ \newterm{closed} (\germanterm{abgeschlossen}) if and only if
$X \setminus A$ is open.
\end{definition}

\begin{center}
\begin{tikzpicture}[scale=1.2]
  % Open Ball
  \draw[dashed, thick] (0,0) circle (1.2);
  \fill (0,0) circle (1.5pt) node[below] {$x$};
  \node[anchor=west] at (0.8, 1.2) {\footnotesize ``boundary'' not included};
  \node[below] at (0, -1.5) {\textbf{Open Ball}};

  % Open set
  \begin{scope}[xshift=4cm]
    \draw[dashed, thick, fill=blue!5] 
      plot[smooth cycle, tension=0.7] coordinates {(-1,-0.5) (-0.5,1) (1,0.5) (0.5,-1)};
    \fill (-0.2,0.2) circle (1pt);
    \draw[dashed, blue!80] (-0.2,0.2) circle (0.3);
    \fill (0.5,-0.3) circle (1pt);
    \draw[dashed, blue!80] (0.5,-0.3) circle (0.2);
    \fill (-0.6,-0.1) circle (1pt);
    \draw[dashed, blue!80] (-0.6,-0.1) circle (0.15);
    \node[right, font=\bfseries] at (1, 0.5) {open};
    \node[below] at (0, -1.5) {\textbf{Open Set}};
  \end{scope}

  % Not Open set (Closed / boundary included)
  \begin{scope}[xshift=7cm]
    \draw[very thick, fill=gray!10] 
      plot[smooth cycle, tension=0.7] coordinates {(-1,-0.5) (-0.5,1) (0.5,0.8) (1,-0.2) (0,-1)};
    \fill (1,-0.2) circle (1.5pt);
    \draw[dashed, purple!80] (1,-0.2) circle (0.4);
    \node[purple, right] at (1.4, -0.2) {\footnotesize doesn't work $\Rightarrow$ not open};
    \node[below] at (0, -1.5) {\textbf{Not Open (e.g., closed)}};
  \end{scope}
\end{tikzpicture}
\end{center}

% Supplement: Analysis_II_Script_v1.pdf, p. 14
\begin{example}[Clopen sets other than $X$ and $\emptyset$]
\label{ex:clopen_nontrivial}
\Cref{ex:X_empty_clopen} shows $X$ and $\emptyset$ are always clopen, but nothing in
\cref{def:open_closed} forces these to be the \emph{only} clopen subsets of $X$. Take
$X := (0,1)\cup(2,3) \subset \mathbb{R}$ with the standard metric: both $(0,1)$ and $(2,3)$ are
clopen \emph{in $X$}. Each is open in $\mathbb{R}$, hence open in $X$; and each is the complement
in $X$ of the other, so each is also closed in $X$. Whether a space has clopen subsets besides
$X$ and $\emptyset$ turns out to be exactly the question \cref{ch:connectedness} is organised
around. A space with none is called \newterm{connected}.
\end{example}

% Supplement: Sascha Brack/Class Notes/Week_02_Notes_Friday_Updated.pdf, p. 1
\begin{proposition}[How open and closed sets combine]
\label{prop:open_closed_operations}
Let $(X,d)$ be a metric space.
\begin{enumerate}[label=\textbf{(\alph*)}]
  \item \label{item:arbitrary_union_open}
    \textbf{Arbitrary} unions of open sets are open: if $\{U_i\}_{i \in I}$ are open, so is
    $\bigcup_{i\in I}U_i$.
  \item \label{item:finite_intersection_open}
    \textbf{Finite} intersections of open sets are open: if $U_1,\dots,U_k$ are open, so is
    $\bigcap_{i=1}^{k}U_i$.
  \item \textbf{Arbitrary} intersections of closed sets are closed, and \textbf{finite} unions
    of closed sets are closed.
\end{enumerate}
\end{proposition}

\begin{proof}
\textbf{(a)} Let $x \in \bigcup_i U_i$. Then $x \in U_{i_0}$ for some $i_0$, and since $U_{i_0}$
is open there is $\varepsilon>0$ with $B_\varepsilon(x) \subseteq U_{i_0} \subseteq \bigcup_iU_i$.

\textbf{(b)} Let $x \in \bigcap_{i=1}^kU_i$. For each $i$ pick $\varepsilon_i>0$ with
$B_{\varepsilon_i}(x) \subseteq U_i$, and set $\varepsilon := \min\{\varepsilon_1,\dots,\varepsilon_k\}$.
Then $\varepsilon > 0$, and \emph{this is the only place finiteness is used}: a minimum of
finitely many positive numbers is positive, whereas an infimum of infinitely many can be $0$.
Since $B_\varepsilon(x) \subseteq U_i$ for every $i$, we conclude that
$B_\varepsilon(x) \subseteq \bigcap_iU_i$.

\textbf{(c)} Take complements and apply \textbf{(a)}, \textbf{(b)} via De Morgan:
$X\setminus\bigcap_iA_i = \bigcup_i(X\setminus A_i)$ and
$X\setminus\bigcup_{i=1}^kA_i = \bigcap_{i=1}^k(X\setminus A_i)$.
\end{proof}

% Generator: Claude Opus 5 (Medium)
\begin{importantremark}[The asymmetry is real, not an artefact of the proof]
\label{rem:finiteness_is_necessary}
It is tempting to read \qt{finite} in \textbf{(b)} as laziness. It is not: an infinite
intersection of open sets genuinely need not be open. In $\mathbb{R}$,
\[\bigcap_{n=1}^{\infty}\left(-\tfrac1n,\ \tfrac1n\right) = \{0\},\]
an intersection of open intervals which is not open. The proof shows exactly where it breaks:
$\varepsilon := \min_i\varepsilon_i$ is the whole argument, and an infimum of infinitely many
positive numbers can be $0$. Dually, $\bigcup_{n\geq1}\big[\tfrac1n,1\big] = (0,1]$ is an
infinite union of closed sets that is not closed.

\begin{ainote}
These four closure properties are exactly the axioms of a \newterm{topology}, which is how the
subject is set up when one drops the metric altogether (\cref{sec:structured_spaces},
\qt{fourth semester}). That is why they are worth isolating: they are the only facts about open
sets that survive into the general theory.
\end{ainote}
\end{importantremark}

% Supplement: Sascha Brack/Class Notes/Week_02_Notes_Friday_Updated.pdf, p. 1
\begin{example}[Open and closed sets in $\mathbb{R}$]
\label{ex:open_closed_R}
Consider the following subsets of $\mathbb{R}$ with the standard metric:
\begin{enumerate}[label=\textbf{(\alph*)}]
  \item \textbf{Is $\mathbb{N}$ open, closed, or neither?} It is \textbf{closed}, since its
    complement $\mathbb{R} \setminus \mathbb{N} = (-\infty, 1) \cup \bigcup_{n=1}^\infty (n, n+1)$
    is a union of open intervals, which is open by
    \Cref{prop:open_closed_operations}\textbf{\ref{item:arbitrary_union_open}}. It is not open,
    because no open ball around any $n \in \mathbb{N}$ is contained in $\mathbb{N}$.
  \item \textbf{Let $A := \{\frac{1}{n} \mid n = 1, 2, \dots\}$.} The set $A$ is \textbf{neither}
    open nor closed. It is not open for the same reason as $\mathbb{N}$. It is not closed because
    $0$ is an accumulation point of $A$, but $0 \notin A$, meaning its complement
    $\mathbb{R} \setminus A$ is not open (any open ball around $0$ intersects $A$).
  \item \textbf{Let $F := \{\frac{1}{n} \mid n = 1, 2, \dots\} \cup \{0\}$.} The set $F$ is
    \textbf{closed}. By adding the limit point $0$, the complement $\mathbb{R} \setminus F$
    becomes a union of open intervals, which is open.
\end{enumerate}
\end{example}

% Supplement: Linus Lüchinger/Slides/Slides-02-23.pdf, slide 20
% Linus's self-test list. The three entries Sascha already works above (N, {1/n},
% {1/n} u {0}) are dropped; what remains is the part that is not already answered.
\begin{exercise}[A longer self-test]
\label{ex:open_closed_self_test}
Linus sets the following list as a self-test at the end of his first session. Each subset of
$\mathbb{R}$ carries the standard metric; decide for each whether it is \textbf{open},
\textbf{closed}, \textbf{both}, or \textbf{neither}.
\begin{enumerate}[label=\textbf{(\alph*)}]
  \item $[0,\infty)$
  \item $\mathbb{Q}$
  \item $\{x - \lfloor x \rfloor : x \in \mathbb{R}\}$
  \item $S := \{x \in \mathbb{R} : x\cdot 2^p \in [1,2] \text{ for some integer } p \geq 0\}$
  \item $T := \{x \in \mathbb{R} : x\cdot 2^p \in (1,2) \text{ for some integer } p \geq 0\}$
  \item $\{x_n : n \in \mathbb{N}\}$, where $x_0 := 1$ and
    $x_n := \tfrac12\big(x_{n-1} + x_{n-1}^2\big)$
  \item $\{\Rea(c) : c \in \mathbb{C},\ \zeta(c) = 0\}$, where $\zeta$ is the Riemann zeta
    function
\end{enumerate}
Compare \textbf{(d)} with \textbf{(e)} before reading the solution: the two differ only in one
bracket.
\end{exercise}

% Quelle: Jérôme Paschoud/Notizen/Woche 2.1 Metrische Räume.pdf, p. 5
\begin{exercise}[Relative openness]
\label{ex:relative_openness_subsets}
Let $X := (\{0\} \times \mathbb{R}) \cup (0,1)^2$ be equipped with the subspace metric from $\mathbb{R}^2$. Consider the following subsets of $X$:
\begin{align*}
  A &:= \{(x,y) \in \mathbb{R}^2 \mid x = 0, y \in (2,3)\} \\
  B &:= \left\{(x,y) \in \mathbb{R}^2 \ \middle|\ \left(x - \tfrac12\right)^2 + \left(y - \tfrac12\right)^2 < \tfrac{1}{16}\right\}
\end{align*}
Are $A$ and $B$ open with respect to $X$? Are they open with respect to $\mathbb{R}^2$?
\end{exercise}

% Generator: Gemini 3.6 Pro
\begin{exercisesolution}[Proof of \cref{ex:relative_openness_subsets}]
\begin{itemize}
  \item $A$ is \textbf{open with respect to $X$}, but \textbf{not open with respect to $\mathbb{R}^2$}. It is open in $X$ because $A = X \cap U$ where $U = (-\varepsilon, \varepsilon) \times (2,3)$ is an open set in $\mathbb{R}^2$. It is not open in $\mathbb{R}^2$ because any open ball around a point in $A$ will contain points with $x \neq 0$, which are not in $A$.
  \item $B$ is \textbf{open with respect to both $X$ and $\mathbb{R}^2$}. $B$ is exactly the open ball $B_{1/4}((\tfrac12, \tfrac12))$ in $\mathbb{R}^2$, which is open in $\mathbb{R}^2$. Since $B \subseteq (0,1)^2 \subset X$, we have $B = X \cap B$, so $B$ is also open relative to $X$.
\end{itemize}
\end{exercisesolution}

% Originally: content/week-02.tex, the second half of \section{Open and closed sets}
% (Corsin Week 2, Monday) --- promoted to a section of its own, since the chapter
% otherwise consisted of one section carrying the chapter's own title.
\section{Interior, closure and boundary}
\label{sec:interior_closure_boundary}

% Sonnet 5 (Medium)
\begin{definition}[Interior, closure, and boundary]
\label{def:interior_closure_boundary}
Let $(X,d)$ be a metric space and $A \subseteq X$. The \newterm{interior}
(\germanterm{Inneres}) of $A$ is
\[\operatorname{int}(A) := \{x \in A : \exists\, r>0 \text{ such that } B_r(x) \subseteq A\},\]
the largest open subset of $X$ contained in $A$. The \newterm{closure} (\germanterm{Abschluss}) of
$A$ is
\[\overline{A} := \{x \in X : B_r(x) \cap A \neq \emptyset \text{ for all } r > 0\},\]
the smallest closed subset of $X$ containing $A$. The \newterm{boundary} (\germanterm{Rand}) of
$A$ is $\partial A := \overline{A} \setminus \operatorname{int}(A)$.
\end{definition}

\begin{aiexample}[Interior, Closure, and Boundary of an Interval]
% Generator: Gemini 3.6 Flash (Medium)
Applying \cref{def:interior_closure_boundary} to $A := (0,1] \subset \mathbb{R}$ under the
standard Euclidean metric:
\begin{itemize}
  \item The \textbf{interior} is $\operatorname{int}(A) = (0,1)$, the largest open set contained in $A$.
  \item The \textbf{closure} is $\overline{A} = [0,1]$, the smallest closed set containing $A$.
  \item The \textbf{boundary} is $\partial A = \overline{A} \setminus \operatorname{int}(A) = \{0, 1\}$.
\end{itemize}
Note that $1 \in A$ while $0 \notin A$: the boundary is not required to lie inside the set, nor
outside it. That is the point of defining it as a difference of two other sets rather than as
\qt{the edge of $A$}.
\end{aiexample}

% Generator: Claude Opus 5 (High) --- FIG-W02-01: the three notions in one picture.
\begin{center}
\begin{tikzpicture}[scale=1.25]
  % A blob A. int(A) is the shaded open part; the solid curve is dA; closure = both together.
  \draw[blue!70!black, thick, fill=blue!14, smooth cycle, tension=0.75]
    plot coordinates {(-1.7,-0.9) (-1.9,0.7) (-0.3,1.4) (1.4,1.0) (1.8,-0.6) (0.2,-1.3)};

  \node[blue!70!black, font=\small] at (-0.15,0.05) {$\operatorname{int}(A)$};

  % a point of the interior, with a ball that fits
  \fill (-1.0,0.35) circle (1.2pt);
  \draw[green!55!black, dashed] (-1.0,0.35) circle (0.42);
  \node[green!55!black, font=\scriptsize, below] at (-1.0,-0.15) {ball fits};

  % a boundary point: every ball meets A and its complement
  \fill[red!80!black] (1.61,0.32) circle (1.4pt);
  \draw[red!80!black, dashed] (1.61,0.32) circle (0.42);
  \node[red!80!black, font=\scriptsize, anchor=west] at (2.15,0.32) {no ball fits};

  \node[blue!70!black, font=\small, anchor=east] at (-2.05,1.05) {$\partial A$};
  \draw[blue!70!black, ->] (-2.0,0.95) -- (-1.55,0.62);

  \node[font=\footnotesize, align=center] at (0,-2.05)
    {$\operatorname{int}(A)$: points with a ball inside $A$ \quad$\cdot$\quad
     $\partial A$: points every ball of which meets\\
     both $A$ and $X\setminus A$ \quad$\cdot$\quad
     $\overline{A} = \operatorname{int}(A) \cup \partial A$};
\end{tikzpicture}
\end{center}

\begin{remark}[Interior and closure are dual]
\label{rem:interior_closure_dual}
The two constructions are not independent: taking complements swaps them,
\[X \setminus \operatorname{int}(A) = \overline{X\setminus A},
  \qquad\qquad
  X \setminus \overline{A} = \operatorname{int}(X\setminus A).\]
Both read off \cref{def:interior_closure_boundary} directly, since $x$ fails to have a ball inside
$A$ exactly when every ball around $x$ meets $X\setminus A$. In particular $\partial A =
\partial(X\setminus A)$: a set and its complement share their boundary, which is why the
boundary of the open disk and of its closed complement are the same circle.
\end{remark}

% Source: Corsin Nick/Class Notes/Week 2.pdf, p. 7
\begin{exercise}[$X$ and $\emptyset$ are clopen]
\label{ex:X_empty_clopen}
Let $(X,d)$ be a metric space. Show that $X$ and $\emptyset$ are both open and closed.
\end{exercise}

% Kept inline: solved directly in Corsin Nick/Class Notes/Week 2.pdf, p. 7.
\begin{exercisesolution}[Proof of \cref{ex:X_empty_clopen}]
$X$ is open, since $B_r(x) \subseteq X$ for all $r > 0$ and $x \in X$
(\cref{def:open_ball}). The empty set contains no
points and is therefore vacuously open. Since $X$ and $\emptyset$ are each other's complement in
$X$, they are also closed.
\end{exercisesolution}

% Supplement: Sascha Brack/Class Notes/Week_02_Notes_Friday_Updated.pdf, p. 2
\begin{example}[Interior, closure, and boundary in $\mathbb{R}$ and $\mathbb{R}^2$]
\label{ex:interior_closure_boundary_more}
Applying \cref{def:interior_closure_boundary} to various sets:
\begin{enumerate}[label=\textbf{(\alph*)}]
  \item Let $\Omega := \mathbb{Z} \subset \mathbb{R}$. Then $\operatorname{int}(\Omega) = \emptyset$ (no interval fits inside $\mathbb{Z}$), $\overline{\Omega} = \mathbb{Z}$ (it is already closed), and $\partial\Omega = \mathbb{Z} \setminus \emptyset = \mathbb{Z}$.
  \item Let $\Omega := B_1(0) \subset \mathbb{R}^2$. Then $\operatorname{int}(\Omega) = B_1(0)$ (it is already open), the closure is the closed disk $\overline{\Omega} = \{x \in \mathbb{R}^2 \mid d(x,0) \leq 1\}$, and the boundary is the unit circle $\partial\Omega = \{x \in \mathbb{R}^2 \mid d(x,0) = 1\}$.
  \item Let $\Omega := \mathbb{Q} \subset \mathbb{R}$. Then $\operatorname{int}(\Omega) = \emptyset$ (every open interval contains irrationals), $\overline{\Omega} = \mathbb{R}$ (every open interval contains rationals, so every real number is a closure point), and $\partial\Omega = \mathbb{R} \setminus \emptyset = \mathbb{R}$.
\end{enumerate}
\end{example}

% Supplement: Analysis_II_Script_v1.pdf, p. 15
\begin{exercise}[Closure and boundary of a ball in $\mathbb{R}^n$]
\label{ex:closure_boundary_of_ball}
For balls $B_r(x)$ in $\mathbb{R}^n$ with the Euclidean metric, prove that
\[\overline{B_r(x)} = \{y \in \mathbb{R}^n : d(x,y) \leq r\}\]
and deduce $\partial B_r(x) = \{y \in \mathbb{R}^n : d(x,y) = r\}$.
\end{exercise}

\begin{remark}
Item \textbf{(b)} of \cref{ex:interior_closure_boundary_more} above is exactly the case $r=1$,
$x=0$, $n=2$ of this exercise, worked here for general $x$, $r$ and $n$.
\end{remark}

% Quelle: Jérôme Paschoud/Notizen/Woche 2.1 Metrische Räume.pdf, p. 6
\begin{exercise}[Interior, closure, and boundary of a mixed set]
\label{ex:interior_closure_boundary_mixed}
Let $E := (\{0\} \times \mathbb{R}) \cup (0,1)^2 \subset \mathbb{R}^2$. Find the interior $\operatorname{int}(E)$, the closure $\overline{E}$, and the boundary $\partial E$ with respect to the standard metric on $\mathbb{R}^2$.
\end{exercise}

% Generator: Gemini 3.6 Pro
\begin{exercisesolution}[Proof of \cref{ex:interior_closure_boundary_mixed}]
\begin{itemize}
  \item \textbf{Interior:} $\operatorname{int}(E) = (0,1)^2$. The line $\{0\} \times \mathbb{R}$ contains no 2D open balls, so it contributes nothing to the interior. The open square $(0,1)^2$ is already open in $\mathbb{R}^2$.
  \item \textbf{Closure:} $\overline{E} = (\{0\} \times \mathbb{R}) \cup [0,1]^2$. The line $\{0\} \times \mathbb{R}$ is closed. The closure of the open square is the closed square $[0,1]^2$. The closure of a finite union is the union of the closures.
  \item \textbf{Boundary:} $\partial E = \overline{E} \setminus \operatorname{int}(E) = (\{0\} \times \mathbb{R}) \cup \big([0,1]^2 \setminus (0,1)^2\big)$. The second term is the perimeter of the square.
\end{itemize}
\end{exercisesolution}

% Originally: content/week-02.tex, the Friday warm-up following \subsection{Sequences and convergence}
% Corsin used this to reopen the topic on the Friday; it needs \cref{def:open_ball}, so it is
% placed here, after the open/closed material, rather than where the week happened to put it.
\section{Unit balls of the three standard metrics}
\label{sec:unit_balls}

\begin{warmup}
What is the shape of $B_1(0) \subseteq \mathbb{R}^2$ with the supremum metric? With the Manhattan
metric?
\end{warmup}

\begin{answer}
Supremum: the axis-aligned \textbf{open square} $(-1,1)^2$. Manhattan: the \textbf{open diamond}
$\{|x_1|+|x_2| < 1\}$. In both cases the strict inequality in \cref{def:open_ball} matters: the
points $(\pm 1,0)$ and $(0,\pm 1)$ sit on the boundary and are \emph{not} in the ball. They are
the corners of the diamond and the edge midpoints of the square, and are drawn hollow below.
\end{answer}


\begin{center}
\begin{tikzpicture}[>=Latex, scale=1.2]
  \begin{scope}[shift={(0,0)}]
    \draw[->] (-1.5,0) -- (1.5,0) node[right] {$x_1$};
    \draw[->] (0,-1.5) -- (0,1.5) node[above] {$x_2$};
    \draw[blue!80!black, thick, dashed, fill=blue!10] (-1,-1) rectangle (1,1);
    % Hollow: (1,0) and (0,1) have d_3 = 1, so they are NOT in the open ball B_1(0).
    \draw[orange!90!black, thick, fill=white] (1,0) circle (2pt);
    \node[orange!90!black, above right, font=\footnotesize] at (1,0) {$(1,0)$};
    \draw[orange!90!black, thick, fill=white] (0,1) circle (2pt);
    \node[orange!90!black, above right, font=\footnotesize] at (0,1) {$(0,1)$};
    \node[below, font=\small, align=center] at (0,-1.6)
      {Supremum metric $d_3$:\\ open square $(-1,1)^2$};
  \end{scope}

  \begin{scope}[shift={(4.4,0)}]
    \draw[->] (-1.5,0) -- (1.5,0) node[right] {$x_1$};
    \draw[->] (0,-1.5) -- (0,1.5) node[above] {$x_2$};
    \draw[blue!80!black, thick, dashed, fill=blue!10] (1,0) -- (0,1) -- (-1,0) -- (0,-1) -- cycle;
    % Same here: d_1((1,0),0) = 1, so the corners are boundary points, not elements.
    \draw[orange!90!black, thick, fill=white] (1,0) circle (2pt);
    \node[orange!90!black, above right, font=\footnotesize] at (1,0) {$(1,0)$};
    \draw[orange!90!black, thick, fill=white] (0,1) circle (2pt);
    \node[orange!90!black, above right, font=\footnotesize] at (0,1) {$(0,1)$};
    \node[below, font=\small, align=center] at (0,-1.6)
      {Manhattan metric $d_1$:\\ open diamond};
  \end{scope}

  \begin{scope}[shift={(8.8,0)}]
    \draw[->] (-1.5,0) -- (1.5,0) node[right] {$x_1$};
    \draw[->] (0,-1.5) -- (0,1.5) node[above] {$x_2$};
    \draw[blue!80!black, thick, dashed, fill=blue!10] (0,0) circle (1);
    \draw[orange!90!black, thick, fill=white] (1,0) circle (2pt);
    \node[orange!90!black, above right, font=\footnotesize] at (1,0) {$(1,0)$};
    \draw[orange!90!black, thick, fill=white] (0,1) circle (2pt);
    \node[orange!90!black, above right, font=\footnotesize] at (0,1) {$(0,1)$};
    \node[below, font=\small, align=center] at (0,-1.6)
      {Standard metric $d_2$:\\ open disk};
  \end{scope}
\end{tikzpicture}
\end{center}

% Generator: Claude Opus 5 (High)
\begin{remark}[The unit ball remembers which metric it came from]
\label{rem:unit_balls_distinguish_metrics}
Three metrics, three unmistakably different unit balls: the square for $d_3$, the disk for
$d_2$, the diamond for $d_1$. This is the quickest way to tell the three apart, and it also makes
\cref{ex:ai_metric_equivalence} visible, since the inclusions
\[B^{d_1}_1(0) \ \subseteq\ B^{d_2}_1(0) \ \subseteq\ B^{d_3}_1(0)\]
are precisely the inequalities $d_3 \leq d_2 \leq d_1$ read backwards. The larger the metric, the
harder it is to be within distance $1$, and hence the smaller the ball. The same family, continued
to all $p \geq 1$, is sketched in \cref{ex:4.3}\textbf{.1}, where the square reappears as the
limiting case $p \to \infty$.

None of this changes which sets are open. All three metrics have the same open sets, as
\cref{ex:ai_metric_equivalence} shows. The balls look different, and the topology does not notice.
\end{remark}



% Originally: content/week-02.tex, end of \section{Compactness}

\begin{aiexercise}[Boundary of the Open Unit Disk]
\label{ex:ai_disk_boundary}
% Generator: Gemini 3.6 Flash (Medium)
Let $D := \{(x,y) \in \mathbb{R}^2 \mid x^2 + y^2 < 1\}$ be the open unit disk in $\mathbb{R}^2$. Show that the boundary of $D$ is the unit circle $\partial D = \{(x,y) \in \mathbb{R}^2 \mid x^2 + y^2 = 1\}$.
\end{aiexercise}



% --- exercises relocated here from the dissolved sheet 2 ---

% Originally: content/exercise-sheets/week-02.tex, problem 2.2
% Source: exercises/Ex2_Analysis2_eng.pdf, p. 1

\begin{exercise}[2.2 --- Multiple choice (closure)]
\label{ex:2.2}
Let $(X,d)$ be a metric space, and $Y_1, Y_2 \subset X$ subsets. Select all the statements below
that are necessarily true.
\begin{enumerate}[label=\textbf{(\alph*)}]
  \item $\overline{Y_1 \cup Y_2} = \overline{Y_1} \cup \overline{Y_2}$
  \item $\overline{Y_1} \cap \overline{Y_2} \subset \overline{Y_1 \cap Y_2}$
  \item $\overline{Y_1 \cap Y_2} \subset \overline{Y_1} \cap \overline{Y_2}$
  \item $\overline{Y_1 \cap Y_2} = \overline{Y_1} \cap \overline{Y_2}$
\end{enumerate}
\exhint[Official hints]{\textbf{(b):} compare with \cref{ex:2.1}.3. \textbf{(d):} play with a set
of three points (a triangle).}
\end{exercise}

% Originally: content/exercise-sheets/week-02.tex, problem 2.3
% Source: exercises/Ex2_Analysis2_eng.pdf, p. 1

\begin{exercise}[2.3 --- Multiple choice (distance from a set)]
\label{ex:2.3}
Let $(X,d)$ be a metric space, and $A \subset X$ a non-empty subset. We define the function
``distance from $A$'' as
\[
  d(\cdot, A) : X \to [0,\infty), \qquad d(x,A) := \inf_{a \in A} d(x,a).
\]
Select all the statements below that are necessarily true.
\begin{enumerate}[label=\textbf{(\alph*)}]
  \item If $A$ is closed and $x \in A^c$, then $d(x,A) > 0$.
  \item The set $M := \{x \in X : d(x,A) \geq 1\}$ is closed in $X$.
  \item For $x, y \in X$, $d(x,A) \leq d(x,y) + d(y,A)$ holds.
  \item If $A^\circ$ is non-empty and $x \in X$, then $d(x,A) = d(x, A^\circ)$.
\end{enumerate}
\exhint[Official hint]{First convince yourself with a drawing that the name of this function is
appropriate; then use the characterization of open/closed sets with sequences.}
\end{exercise}

% Originally: content/exercise-sheets/week-02.tex, problem 2.4
% Source: exercises/Ex2_Analysis2_eng.pdf, p. 1

\begin{exercise}[2.4 --- Boundary, Interior etc.]
\label{ex:2.4}
Determine the interior, closure, and boundary of the following subsets $Y$ of $\mathbb{R}$, for
the standard topology on $\mathbb{R}$. No need to justify the answer.
\begin{center}
\begin{tabular}{ll ll}
(1) & $Y = [0,1]$ & (2) & $Y = \mathbb{Q}$ \\
(3) & $Y = \emptyset$ & (4) & $Y = (0,1)$ \\
(5) & $Y = [-1,1) \setminus \{0\}$ & (6) & $Y = [0,\infty)$ \\
(7) & $Y = \{0\}$ & (8) & $Y = \left\{\tfrac{1}{n} \mid n \in \mathbb{N}\setminus\{0\}\right\}$
\end{tabular}
\end{center}
\end{exercise}

% Originally: the \section{Solutions} of content/week-02.tex

\section{Solutions}
\label{sec:open_closed_solutions}

\begin{exercisesolution}[Proof of \cref{ex:ai_open_balls_are_open}]
% Generator: Claude Opus 5 (Medium)
Let $y \in B_r(x)$, so $d(y,x) < r$, and set $\varepsilon := r - d(y,x) > 0$. We claim that
$B_\varepsilon(y) \subseteq B_r(x)$, which is exactly what \cref{def:open_closed} requires.

Let $z \in B_\varepsilon(y)$, so $d(z,y) < \varepsilon$. By the triangle inequality
(\cref{def:metric_space}),
\[d(z,x) \leq d(z,y) + d(y,x) < \varepsilon + d(y,x) = \big(r-d(y,x)\big) + d(y,x) = r,\]
so $z \in B_r(x)$. Since $y \in B_r(x)$ was arbitrary, $B_r(x)$ is open.

\begin{remark}
The entire content is the choice $\varepsilon := r - d(y,x)$: the radius left over after
travelling from the centre out to $y$. Note that it shrinks to $0$ as $y$ approaches the edge of
the ball, so no \emph{single} $\varepsilon$ works for every $y$ at once. That is the same
\qt{the constant depends on the point} phenomenon which separates continuity from uniform
continuity in \cref{rem:continuity_vs_uniform_continuity}, met here for the first time.
\end{remark}
\end{exercisesolution}

\begin{exercisesolution}[Proof of \cref{ex:ai_disk_boundary}]
% Generator: Gemini 3.6 Flash (Medium)
Since $D$ is open, its interior is $\operatorname{int}(D) = D$. The closure of $D$ is $\overline{D} = \{(x,y) \in \mathbb{R}^2 \mid x^2 + y^2 \leq 1\}$. By \cref{def:interior_closure_boundary}, the boundary is:
\begin{align*}
\partial D &= \overline{D} \setminus \operatorname{int}(D) \\
&= \{(x,y) \in \mathbb{R}^2 \mid x^2 + y^2 \leq 1\} \setminus \{(x,y) \in \mathbb{R}^2 \mid x^2 + y^2 < 1\} \\
&= \{(x,y) \in \mathbb{R}^2 \mid x^2 + y^2 = 1\}.
\end{align*}
\end{exercisesolution}

% ===========================================================================
% Official sheet 2 (2.1--2.7). Derived independently, then checked against
% exercises/Sol2_Analysis2_eng.pdf.
% ===========================================================================

\begin{exercisesolution}[Solution to \cref{ex:2.2}]
% Generator: Claude Opus 5 (High)
\textbf{(a) and (c) are necessarily true; (b) and (d) are not.}

\textbf{(a) True.} Both inclusions come from general properties of the closure. Since
$Y_i \subseteq Y_1 \cup Y_2$ and the closure is monotone,
$\overline{Y_i} \subseteq \overline{Y_1\cup Y_2}$ for $i=1,2$, which gives \qt{$\supseteq$}.
Conversely $\overline{Y_1}\cup\overline{Y_2}$ is a \emph{finite} union of closed sets, hence
closed (\cref{prop:open_closed_operations}\textbf{(c)}), and it contains $Y_1\cup Y_2$; since
$\overline{Y_1\cup Y_2}$ is the smallest closed set containing $Y_1 \cup Y_2$
(\cref{def:interior_closure_boundary}), \qt{$\subseteq$} follows.

\textbf{(c) True,} and it is the easy half: $Y_1\cap Y_2 \subseteq Y_i$ gives
$\overline{Y_1\cap Y_2} \subseteq \overline{Y_i}$ for each $i$ by monotonicity, hence
$\overline{Y_1\cap Y_2} \subseteq \overline{Y_1}\cap\overline{Y_2}$.

\textbf{(b) False,} and therefore \textbf{(d) False} as well, (d) being (b) and (c) combined. In
$\mathbb{R}$ take
\[Y_1 := (0,1), \qquad Y_2 := (1,2).\]
Then $\overline{Y_1}\cap\overline{Y_2} = [0,1]\cap[1,2] = \{1\}$, whereas
$Y_1\cap Y_2 = \emptyset$ and so $\overline{Y_1\cap Y_2} = \emptyset$. The two sets meet only in
the limit, and the closure of the intersection cannot see a point belonging to neither set.

\begin{remark}[Where finiteness enters, again]
Part \textbf{(a)} is false for infinite unions, for exactly the reason given in
\cref{rem:finiteness_is_necessary}. Taking $Y_q := \{q\}$ for each $q \in \mathbb{Q}$,
\[\bigcup_{q\in\mathbb{Q}}\overline{Y_q} = \mathbb{Q}
  \qquad\text{but}\qquad
  \overline{\bigcup_{q\in\mathbb{Q}}Y_q} = \overline{\mathbb{Q}} = \mathbb{R}.\]
So the honest statement is: the closure commutes with \emph{finite} unions, and only
one-directionally with intersections. That is the same asymmetry that runs through the whole of
\cref{prop:open_closed_operations}.
\end{remark}
\end{exercisesolution}

\begin{exercisesolution}[Solution to \cref{ex:2.3}]
% Generator: Claude Opus 5 (High)
\textbf{(a), (b) and (c) are necessarily true; (d) is not.} It is efficient to prove
\textbf{(c)} first, since \textbf{(b)} then follows from it.

\textbf{(c) True.} Fix $x,y \in X$. For every $a \in A$ the triangle inequality gives
$d(x,a) \leq d(x,y) + d(y,a)$, and the left-hand side is bounded below by $d(x,A)$, so
\[d(x,A) \leq d(x,y) + d(y,a) \qquad\text{for every } a \in A.\]
Taking the infimum over $a$ on the right yields $d(x,A) \leq d(x,y) + d(y,A)$. Exchanging $x$
and $y$ and combining,
\[\big|d(x,A) - d(y,A)\big| \leq d(x,y),\]
so $d(\cdot,A)$ is \textbf{$1$-Lipschitz} (\cref{def:lipschitz_continuous}), in particular
continuous --- which is what the remaining parts use.

\textbf{(a) True.} If $A$ is closed then $A^c$ is open, so for $x \in A^c$ there is $r>0$ with
$B_r(x) \subseteq A^c$. No point of $A$ lies in that ball, i.e.\ $d(x,a) \geq r$ for every
$a \in A$, and taking the infimum gives $d(x,A) \geq r > 0$.

\textbf{(b) True.} By \textbf{(c)} the function $d(\cdot,A)$ is continuous, and
$M = d(\cdot,A)^{-1}\big([1,\infty)\big)$ is the preimage of a closed set under a continuous
map, hence closed (\cref{def:continuous}\textbf{(c)}, applied to complements).

\textbf{(d) False.} In $\mathbb{R}$ take $A := [0,1]\cup\{2\}$, whose interior
$A^\circ = (0,1)$ is non-empty as required. At the point $x := 2$,
\[d(x,A) = 0 \quad\text{(because } 2 \in A\text{)}, \qquad\qquad
  d(x,A^\circ) = \inf_{a\in(0,1)}|2-a| = 1.\]
The isolated point $2$ of $A$ contributes to $d(\cdot,A)$ but is invisible to $A^\circ$. The
hypothesis \qt{$A^\circ$ is non-empty} is far too weak, because it constrains $A$ nowhere near
$x$; the statement would become true under a hypothesis such as $A = \overline{A^\circ}$.
\end{exercisesolution}

\begin{exercisesolution}[Solution to \cref{ex:2.4}]
% Generator: Claude Opus 5 (High)
All with respect to the standard metric on $\mathbb{R}$, using
\cref{def:interior_closure_boundary}.
\begin{center}
\begin{tabular}{@{}c l l l l@{}}
& \textbf{$Y$} & \textbf{$\operatorname{int}(Y)$} & \textbf{$\overline{Y}$} & \textbf{$\partial Y$} \\
\hline
(1) & $[0,1]$ & $(0,1)$ & $[0,1]$ & $\{0,1\}$ \\
(2) & $\mathbb{Q}$ & $\emptyset$ & $\mathbb{R}$ & $\mathbb{R}$ \\
(3) & $\emptyset$ & $\emptyset$ & $\emptyset$ & $\emptyset$ \\
(4) & $(0,1)$ & $(0,1)$ & $[0,1]$ & $\{0,1\}$ \\
(5) & $[-1,1)\setminus\{0\}$ & $(-1,0)\cup(0,1)$ & $[-1,1]$ & $\{-1,0,1\}$ \\
(6) & $[0,\infty)$ & $(0,\infty)$ & $[0,\infty)$ & $\{0\}$ \\
(7) & $\{0\}$ & $\emptyset$ & $\{0\}$ & $\{0\}$ \\
(8) & $\left\{\tfrac1n : n \geq 1\right\}$ & $\emptyset$ &
  $\left\{\tfrac1n : n \geq 1\right\}\cup\{0\}$ &
  $\left\{\tfrac1n : n \geq 1\right\}\cup\{0\}$ \\
\end{tabular}
\end{center}

Three rows repay a second look.
\begin{itemize}
  \item \textbf{(1) versus (4).} The closed interval and the open interval have the same
    interior, the same closure and the same boundary. These three operations simply cannot
    distinguish a set from any other set squeezed between its interior and its closure.
  \item \textbf{(5).} Removing the single point $0$ from $[-1,1)$ adds $0$ to the boundary
    without changing the closure: a removed interior point becomes a boundary point. Note also
    that $1 \in \partial Y$ although $1 \notin Y$, while $-1 \in \partial Y$ and $-1 \in Y$ ---
    the boundary is indifferent to membership.
  \item \textbf{(2) and (8).} Both have empty interior, so $\partial Y = \overline{Y}$: once a
    set has no interior, the boundary swallows the entire closure. That is the general reason
    for $\partial\mathbb{Q} = \mathbb{R}$ in
    \cref{ex:interior_closure_boundary_more}\textbf{(c)}, and not an accident of $\mathbb{Q}$.
\end{itemize}
\end{exercisesolution}


% ============================================================
% Chapter 4 --- Sequences and convergence
% Originally: content/week-02.tex, \subsection{Sequences and convergence} (Corsin Week 2,
%             Monday) --- promoted from a subsection to a chapter of its own.
% ============================================================
\chapter{Sequences and Convergence}
\label{ch:sequences}

% Transition: Claude Opus 5 (High)
Open and closed sets describe a space statically. They say which points sit comfortably inside a
set and which sit on its edge, but they say it all at once, with no notion of approach. Sequences
make the same information dynamic, and in a metric space the two descriptions turn out to be
equivalent: a set is closed exactly when it catches the limits of its own convergent sequences.

That equivalence is what makes this short chapter worth its own place. Almost every later
argument that needs to produce a point produces it as a limit, and the sequential criterion is
what converts that limit back into a statement about open or closed sets. The chapter also
isolates a fact easy to pass over, namely that convergence depends only on the topology and not
on the metric generating it, so two visibly different metrics can have exactly the same
convergent sequences.

% Originally: content/week-02.tex, \subsection{Sequences and convergence} (Corsin Week 2, Monday)

\section{Sequences and convergence}
\label{sec:sequences_and_convergence}

% Transition: Claude Opus 5 (High)
Open and closed sets describe a space statically: which points sit comfortably inside a set and
which sit on its edge. Sequences describe the same thing dynamically, by asking where a moving
point ends up. The two descriptions turn out to be equivalent
(\cref{prop:open_closed_via_sequences}), and having both is what makes metric spaces pleasant to
work in. One description is usually easier to verify, the other easier to apply. Note that the
definition below is word for word the one from \emph{Analysis I}, with $|x_n - x|$ replaced by
$d(x_n,x)$. That substitution is the pattern for the whole chapter.

% Source: Corsin Nick/Class Notes/Week 2.pdf, pp. 1--13
\begin{definition}[Convergent sequence]
\label{def:convergent_sequence}
A \newterm{sequence} $(x_n)_{n=0}^{\infty}$ in a metric space $(X,d)$ is a function
$x : \mathbb{N} \to X$. We say that $(x_n)_{n=0}^{\infty}$ \newterm{converges} to $x \in X$ if
\[\forall \varepsilon > 0 \ \exists N \in \mathbb{N} \text{ such that } d(x_n, x) < \varepsilon \quad \forall n \geq N.\]
\end{definition}

\begin{center}
\begin{tikzpicture}[scale=1.2]
  % Axes
  \draw[->, thick] (-0.2, 0) -- (3, 0);
  \draw[->, thick] (0, -0.2) -- (0, 2);
  \node[anchor=south west] at (0, 2) {$\mathbb{R}^2$};

  % Sequence points
  \fill (0.5, 0.4) circle (1.5pt) node[below] {$x_1$};
  \fill (1.0, 1.5) circle (1.5pt) node[left] {$x_2$};
  \fill (1.8, 1.7) circle (1.5pt) node[above] {$x_3$};
  \fill (2.2, 0.5) circle (1.5pt) node[below] {$x_4$};
  \fill (2.6, 1.0) circle (1.5pt) node[right] {$x_5$};
  \node at (2.4, 1.6) {etc.};
  
  % Path
  \draw[->, purple, dashed, thick] (0.6, 0.5) to[bend left=20] (0.9, 1.4);
  \draw[->, purple, dashed, thick] (1.1, 1.6) to[bend left=20] (1.7, 1.7);
  \draw[->, purple, dashed, thick] (1.8, 1.5) to[bend right=20] (2.1, 0.6);
  \draw[->, purple, dashed, thick] (2.3, 0.6) to[bend right=20] (2.5, 0.9);
  \draw[->, purple, dotted, thick] (2.6, 1.1) to[bend left=20] (2.4, 1.5);
\end{tikzpicture}
\end{center}

% Supplement: Analysis_II_Script_v1.pdf, p. 7
\begin{notation}[\qt{Eventually}]
\label{not:eventually}
Let $(x_n)_{n\geq0}$ be a sequence in a set $X$ and let $\mathcal{P}$ be a property that a point
of $X$ may or may not have. One says $x_n$ satisfies $\mathcal{P}$ \newterm{eventually}
(\germanterm{schlie\ss lich}) if there exists $N \in \mathbb{N}$ such that $\mathcal{P}(x_n)$
holds for all $n \geq N$, or equivalently if $\mathcal{P}(x_n)$ fails for at most finitely many
$n$. In this vocabulary, \cref{def:convergent_sequence} reads: $(x_n)_n$ converges to $x$ if, for
every $\varepsilon>0$, $x_n$ eventually lies in $B_\varepsilon(x)$. The same compression applies
throughout the rest of this chapter and the next three.
\end{notation}

% Supplement: Sascha Brack/Class Notes/Week_02_Notes_Monday_Updated.pdf, p. 7
\begin{example}[A convergent sequence in $\mathbb{R}^3$]
\label{ex:convergent_sequence_R3}
Define a sequence $(x_n)_{n=1}^{\infty}$ in $\mathbb{R}^3$ by $x_n := \left(\frac{1}{n}, \frac{2}{n}, 1\right)$ for all $n \in \mathbb{N}$. Does a limit exist? If yes, what is it?
Yes, the sequence converges. Its limit can be computed coordinate-wise:
\[
\lim_{n\to\infty} x_n = \left(\lim_{n\to\infty} \frac{1}{n}, \lim_{n\to\infty} \frac{2}{n}, \lim_{n\to\infty} 1\right) = (0,0,1).
\]
\end{example}

% Sonnet 5 (Medium)
\begin{proposition}[Uniqueness of limits]
\label{prop:uniqueness_of_limits}
Let $(X,d)$ be a metric space. If a sequence $(x_n)_{n=0}^{\infty}$ in $X$ converges to both
$x \in X$ and $x' \in X$, then $x = x'$.
\end{proposition}

\begin{proof}
\begin{center}
\begin{tikzpicture}
  \fill (0,0) circle (1.5pt) node[left] {$x$};
  \draw[green!60!black, dashed, thick] (0,0) circle (1);
  \node[green!60!black, above left] at (0.8, 0.8) {$\varepsilon$};
  
  \fill (2,0) circle (1.5pt) node[right] {$y$};
  
  \fill[purple] (0.5,0.2) circle (1pt) node[below] {\tiny $x_n$};
  \fill[purple] (0.3,-0.4) circle (1pt);
  \fill[purple] (-0.2,0.5) circle (1pt);
  \fill[purple] (-0.4,-0.2) circle (1pt);
  \fill[purple] (0.1,0.1) circle (1pt);
  
  \node[anchor=west, text width=6cm] at (3, 0) {
    \small Proof idea: Let $\varepsilon > 0$. Since $x_n \to x$, for large $n$, $d(x_n, x) < \varepsilon$. For $\varepsilon$ small enough, $y$ will be outside the green circle $\bigcirc$ so the $x_n$ will never get closer to $y$.
  };
\end{tikzpicture}
\end{center}
Fix $\varepsilon > 0$. By \cref{def:convergent_sequence}, there exist $N, N' \in \mathbb{N}$ such
that $d(x_n,x) < \varepsilon/2$ for all $n \geq N$ and $d(x_n,x') < \varepsilon/2$ for all
$n \geq N'$. For any $n \geq \max\{N,N'\}$, the triangle inequality (\cref{def:metric_space})
gives
\[d(x,x') \leq d(x,x_n) + d(x_n,x') < \varepsilon.\]
Since $\varepsilon > 0$ was arbitrary, $d(x,x') = 0$, so $x = x'$ by definiteness
(\cref{def:metric_space}).
\end{proof}

This justifies speaking of \qt{the} limit of a convergent sequence, as we do implicitly
throughout.

\begin{definition}[Accumulation point]
\label{def:accumulation_point}
Let $(x_n)_{n=0}^{\infty}$ be a sequence in a metric space $(X,d)$. A point $x \in X$ is an
\newterm{accumulation point} (\germanterm{H\"aufungspunkt}) of $(x_n)_{n=0}^{\infty}$ if some
subsequence $(x_{n_k})_{k=0}^{\infty}$ converges to $x$.
\end{definition}

% Supplement: Sascha Brack/Class Notes/Week_02_Notes_Monday_Updated.pdf, p. 3
\begin{remark}[A claim in the lecture script that is not fully correct]
The official lecture script states that a sequence converges to $x$ if and only if $x$ is its
\emph{only} accumulation point (\cref{def:accumulation_point}). Sascha Brack's class notes flag
this as false in general, with two counterexamples worth recording. First, a sequence can have a
unique accumulation point without converging: on $[0,1)$, let $x_{2n+1} := \tfrac1n$ and
$x_{2n} := 1-\tfrac1n$; the only accumulation point is $0$, yet $(x_n)$ does not converge, since
the terms $x_{2n} \to 1 \notin [0,1)$ leave every neighbourhood of $0$ infinitely often. Second,
even in $\mathbb{R}$ itself, the sequence with $x_n := \tfrac1n$ for odd $n$ and $x_n := n$ for
even $n$ has $0$ as its only accumulation point but is unbounded, hence divergent. What \emph{is}
true is the easier direction: if $(x_n)$ converges to $x$, then $x$ is its unique accumulation
point (this follows directly from \cref{prop:uniqueness_of_limits} together with the fact that
every subsequence of a convergent sequence converges to the same limit). The converse needs an
extra hypothesis, such as compactness of the ambient space (\cref{sec:compactness}) together with
some control ruling out the first counterexample's boundary-escaping behaviour.
\end{remark}

\begin{proposition}[Open/closed sets via sequences]
\label{prop:open_closed_via_sequences}
\begin{enumerate}[label=\textbf{(\alph*)}]
  \item $U \subseteq X$ is \textbf{open} if and only if, for every convergent sequence
    (\cref{def:convergent_sequence}) in $X$ with limit in $U$, the sequence eventually lies in
    $U$.
  \item $A \subseteq X$ is \textbf{closed} if and only if every convergent sequence in $A$ has its
    limit in $A$.
\end{enumerate}
\end{proposition}

% Generator: Claude Opus 5 (High) --- Corsin states this without proof, but it is the
% workhorse behind the uniform-convergence exercise below, so the argument is filled in here.
\begin{proof}
\begin{enumerate}[label=\textbf{(\alph*)}]
  \item \textbf{($\Rightarrow$)} Let $U$ be open and let $x_n \to x \in U$. Choose
    $\varepsilon>0$ with $B_\varepsilon(x) \subseteq U$ (\cref{def:open_closed}). By
    \cref{def:convergent_sequence} there is $N$ with $d(x_n,x)<\varepsilon$ for all $n \geq N$,
    so $x_n \in B_\varepsilon(x) \subseteq U$ from $N$ onwards.

    \textbf{($\Leftarrow$)} By contraposition. If $U$ is not open, some $x \in U$ has
    $B_{1/n}(x) \not\subseteq U$ for every $n \geq 1$, so we may pick
    $x_n \in B_{1/n}(x)\setminus U$. Then $d(x_n,x) < \tfrac1n \to 0$, so $x_n \to x \in U$, and
    yet \emph{no} term of the sequence lies in $U$.
  \item \textbf{($\Rightarrow$)} Let $A$ be closed, $(x_n)$ a sequence in $A$ with $x_n \to x$.
    If $x \notin A$, then $x$ lies in the open set $X\setminus A$, so by \textbf{(a)} the
    sequence eventually lies in $X\setminus A$, contradicting $x_n \in A$ for all $n$. Hence
    $x \in A$.

    \textbf{($\Leftarrow$)} We show $X\setminus A$ is open. If it were not, the same
    construction as in \textbf{(a)} yields $x \in X\setminus A$ and points
    $x_n \in B_{1/n}(x) \cap A$. Then $(x_n)$ is a sequence in $A$ converging to $x$, so the
    hypothesis forces $x \in A$, which is the required contradiction.
\end{enumerate}
\end{proof}

\begin{remark}[The shrinking-ball construction]
\label{rem:shrinking_ball_trick}
Both \qt{$\Leftarrow$} directions used the same device: negating \qt{some ball fits inside $U$}
gives a bad point in $B_{1/n}(x)$ for \emph{every} $n$, and collecting those points produces a
sequence converging to $x$. It is worth recognising, because it is the standard bridge from a
statement about balls to a statement about sequences, and it is exactly what makes
\cref{def:open_closed} and \cref{prop:open_closed_via_sequences} interchangeable in practice. In
the usual division of labour one proves closedness by the sequence criterion and openness by the
ball criterion.
Note that it needs the radii $\tfrac1n$ to shrink to $0$, which is where the metric enters;
in a general topological space the equivalence can fail.
\end{remark}

% Supplement: Analysis_II_Script_v1.pdf, p. 17
\begin{lemma}[Convergence is a topological notion]
\label{lem:convergence_is_topological}
Let $(X,d)$ be a metric space. A sequence $(x_n)_n$ in $X$ converges to $x$ if and only if, for
every open set $U \ni x$, $(x_n)_n$ eventually (\cref{not:eventually}) lies in $U$.
\end{lemma}

\begin{proof}
If $x_n \to x$ and $U \ni x$ is open, some $\varepsilon>0$ has $B_\varepsilon(x)\subseteq U$
(\cref{def:open_closed}), and $x_n$ eventually lies in $B_\varepsilon(x)$
(\cref{def:convergent_sequence}), hence in $U$. Conversely, taking $U := B_\varepsilon(x)$ itself
for each $\varepsilon>0$ recovers exactly \cref{def:convergent_sequence}.
\end{proof}

% Supplement: Analysis_II_Script_v1.pdf, p. 17
\begin{corollary}[Two metrics with the same topology have the same convergent sequences]
\label{cor:same_topology_same_convergence}
Let $X$ be a set endowed with two metrics $d_1$ and $d_2$. Then $(X,d_1)$ and $(X,d_2)$ have the
same convergent sequences if and only if the topologies they generate coincide.
\end{corollary}

\begin{proof}
By \cref{lem:convergence_is_topological}, whether $x_n \to x$ depends only on the collection of
open sets, that is, only on the topology $d_1$ or $d_2$ generates. Consequently, equal topologies
give equal convergent sequences.

For the converse, suppose $U \subseteq X$ is open with respect to $d_1$ but not with respect to
$d_2$; we exhibit a sequence distinguishing the two. Since $U$ is not $d_2$-open, some $x \in U$
has, for every $k \geq 0$, a point $x_k \in B^{d_2}_{2^{-k}}(x) \setminus U$. Then $x_k \to x$
with respect to $d_2$ (\cref{def:convergent_sequence}, since $2^{-k}\to0$), so by hypothesis also
$x_k \to x$ with respect to $d_1$. But $U$ is $d_1$-open and contains $x$, so
\cref{lem:convergence_is_topological} forces $(x_k)_k$ to eventually lie in $U$, contradicting
$x_k \notin U$ for every $k$. So every $d_1$-open set is $d_2$-open, and by symmetry the two
topologies coincide.
\end{proof}

\begin{remark}
This is the precise sense in which \qt{only the topology matters, not the metric itself} for
convergence. Two visibly different metrics, even one bounded and one unbounded, generate the same
convergent sequences exactly when they generate the same open sets, and one never has to compare
the metrics directly. \Cref{ex:2.5}\textbf{.2}'s three product metrics are one instance of this; so is
\cref{prop:concave_reshaping_metric} together with \cref{ex:reshaping_same_convergence}.
\end{remark}

% Source: Corsin Nick/Class Notes/Week 2.pdf, pp. 1--13
% (Re-cited: the proof and remark above this point are additions, not transcription.)
\begin{example}[Open, closed, and neither in $\mathbb{R}^2$]
\label{ex:open_closed_neither_R2}
In the euclidean metric space $\mathbb{R}^2$, and with
\cref{prop:open_closed_via_sequences} in hand, each answer is a one-line sequence argument:
\begin{enumerate}[label=\textbf{(\alph*)}]
  \item $[0,1]^2$ is \textbf{closed}: a sequence in $[0,1]^2$ converges coordinatewise
    (\cref{prop:Rn_coordinatewise_complete}), and each coordinate limit stays in $[0,1]$ because
    weak inequalities survive limits.
  \item $(0,1)^2$ is \textbf{open}: given $x \in (0,1)^2$, the distance from $x$ to the nearest
    edge is some $\varepsilon>0$, and $B_\varepsilon(x) \subseteq (0,1)^2$.
  \item $[0,1] \times (0,1)$ is \textbf{neither open nor closed}: it is not open because no ball
    around $(0,\tfrac12)$ fits inside it, and it is not closed because
    $(\tfrac12, \tfrac1n) \to (\tfrac12, 0)$ is a convergent sequence in the set whose limit is
    not.
\end{enumerate}
The picture below shows exactly these three witnesses. In \textbf{(a)} the escaping sequence stays
trapped, whereas in \textbf{(c)} one sequence escapes through the open edge while another does
not.
\end{example}

\begin{center}
\begin{tikzpicture}[scale=1.4]
  % Panel 1: Closed
  \begin{scope}[shift={(0,0)}]
    \draw[thick, fill=blue!5] (0,0) rectangle (1.8,1.8);
    \draw[purple, thick, dotted, ->] (0.5,0.6).. controls (0.8,1.2) and (1.2,1.3).. (1.75,1.75);
    \node[purple, font=\footnotesize, left] at (0.6,0.6) {$x_n$};
    \node[red, font=\small] at (1.8,1.8) {$\times$};
    \node[below, font=\small] at (0.9,-0.1) {(a) $[0,1]^2$ (\textbf{closed})};
  \end{scope}

  % Panel 2: Open
  \begin{scope}[shift={(2.8,0)}]
    \draw[thick, dashed, fill=blue!5] (0,0) rectangle (1.8,1.8);
    \draw[purple, thick, dotted, ->] (0.4,1.4).. controls (0.8,1.6) and (1.2,1.7).. (1.75,1.8);
    \node[purple, font=\footnotesize, left] at (0.4,1.4) {$x_n$};
    \node[below, font=\small] at (0.9,-0.1) {(b) $(0,1)^2$ (\textbf{open})};
  \end{scope}

  % Panel 3: Neither
  \begin{scope}[shift={(5.6,0)}]
    \fill[blue!5] (0,0) rectangle (1.8,1.8);
    \draw[thick] (0,0) -- (0,1.8);
    \draw[thick] (1.8,0) -- (1.8,1.8);
    \draw[thick, dashed] (0,1.8) -- (1.8,1.8);
    \draw[thick, dashed] (0,0) -- (1.8,0);

    \draw[orange!90!black, thick, dotted, ->] (0.6,0.4).. controls (0.8,1.1) and (1.0,1.5).. (1.2,1.8);
    \node[orange!90!black, font=\footnotesize, left] at (0.6,0.4) {$y_n$};

    \draw[purple, thick, dotted, ->] (0.4,1.2).. controls (0.9,0.8) and (1.3,0.3).. (1.75,0.05);
    \node[purple, font=\footnotesize, left] at (0.4,1.2) {$x_n$};
    \node[red, font=\small] at (1.8,0) {$\times$};

    \node[below, font=\small] at (0.9,-0.1) {(c) $[0,1]\times(0,1)$ (\textbf{neither})};
  \end{scope}
\end{tikzpicture}
\end{center}

\begin{importantremark}[Common misconception!]
\label{rem:common_misconception}
In the euclidean metric space $X = [0,1]^2$, the set $[0,1]^2$ \textbf{is open}!
\begin{center}
\begin{tabular}{p{0.42\linewidth} p{0.42\linewidth}}
\textbf{$[0,1]^2 \subseteq \mathbb{R}^2$} & \textbf{$[0,1]^2 \subseteq [0,1]^2$} \\
Not open, since $B_\varepsilon(x)$ is not contained in $[0,1]^2$ if $x$ is on the boundary. &
Open, since the full space is always open, as shown before.
\end{tabular}
\end{center}
\end{importantremark}

\begin{center}
\begin{tikzpicture}[scale=1.5]
  % Panel 1: Ambient R^2
  \begin{scope}[shift={(0,0)}]
    \draw[thick, fill=gray!10] (0,0) rectangle (1.6,1.6);
    \fill[red] (1.6, 0.8) circle (1.5pt) node[right, font=\footnotesize] {$x$};
    \draw[orange!90!black, thick, dashed] (1.6,0.8) circle (0.55);
    \node[below, font=\small] at (0.8,-0.2) {$[0,1]^2 \subseteq \mathbb{R}^2$ (\textbf{not open})};
  \end{scope}

  % Panel 2: Ambient [0,1]^2
  \begin{scope}[shift={(3.2,0)}]
    \draw[thick, fill=gray!10] (0,0) rectangle (1.6,1.6);
    \begin{scope}
      \clip (0,0) rectangle (1.6,1.6);
      \draw[orange!90!black, thick, dashed, fill=orange!15] (1.6,0.8) circle (0.55);
    \end{scope}
    \fill[red] (1.6, 0.8) circle (1.5pt) node[right, font=\footnotesize] {$x$};
    \node[below, font=\small] at (0.8,-0.2) {$[0,1]^2 \subseteq [0,1]^2$ (\textbf{open})};
  \end{scope}
\end{tikzpicture}
\end{center}

\begin{exercise}[Uniform convergence]
\label{ex:uniform_convergence}
\begin{enumerate}[label=\textbf{(\alph*)}]
  \item Write down the definition of convergence of a sequence $(f_n)_{n=0}^{\infty}$ to $f$ in
    the metric space $X$ of bounded functions $[0,1] \to \mathbb{R}$ with the supremum metric.
  \item You know the resulting expression from Analysis 1. What is it called?
  \item \textit{(Optional)} Use the fact that $C^0([0,1]) \subseteq X$ is closed and
    \Cref{prop:open_closed_via_sequences} on \qt{sequentially closed} to conclude a big result
    from Analysis 1.
\end{enumerate}
\end{exercise}

% Kept inline: solved directly in Corsin Nick/Class Notes/Week 2.pdf, pp. 9--10.
\begin{exercisesolution}[Solution to \cref{ex:uniform_convergence}]
\begin{enumerate}[label=\textbf{(\alph*)}]
  \item $\displaystyle \forall \varepsilon > 0 \ \exists N \in \mathbb{N} : \sup_{x \in [0,1]} |f_n(x) - f(x)| < \varepsilon \quad \forall n \geq N$.
  \item This is \newterm{uniform convergence} (\germanterm{gleichm\"assige Konvergenz})!
  \item $C^0([0,1])$ is closed in the space of bounded functions with the supremum metric.

    \textit{Proof.} Let $f \in X \setminus C^0([0,1])$, i.e., $f : [0,1] \to \mathbb{R}$ is
    bounded and has at least one discontinuity at a point $x_0 \in [0,1]$. That is, we find
    $\varepsilon > 0$ such that for all $\delta > 0$ there exists $y \in (x_0 - \delta, x_0 +
    \delta)$ with $|f(x_0) - f(y)| \geq \varepsilon$. Then the $\tfrac{\varepsilon}{3}$-ball in
    $X$,
    \[
    B_{\varepsilon/3}(f) := \Big\{ g \in X : \sup_{x \in [0,1]} |f(x) - g(x)| < \tfrac{\varepsilon}{3} \Big\},
    \]
    is contained in $X \setminus C^0([0,1])$.

    \textit{Why:} let $g \in B_{\varepsilon/3}(f)$ and let $\delta > 0$ be arbitrary. Choose
    $y \in (x_0-\delta,\,x_0+\delta)$ with $|f(x_0)-f(y)| \geq \varepsilon$. Then
    \[\varepsilon \leq |f(x_0)-f(y)|
      \leq \underbrace{|f(x_0)-g(x_0)|}_{<\ \varepsilon/3}
         + |g(x_0)-g(y)|
         + \underbrace{|g(y)-f(y)|}_{<\ \varepsilon/3},\]
    so $|g(x_0)-g(y)| > \varepsilon/3$. As $\delta$ was arbitrary, $g$ is discontinuous at
    $x_0$ too, i.e.\ $g \notin C^0([0,1])$. The whole point of the $\tfrac{\varepsilon}{3}$ is
    that two of the three terms are spent on moving from $f$ to $g$ and back, leaving a third
    of the jump behind.

    Given this fact, \cref{prop:open_closed_via_sequences} ensures that the uniform limit $f$ of
    $(f_n)_{n=0}^{\infty}$ in $C^0([0,1])$ is also in $C^0([0,1])$, i.e., continuous, given that
    $f : [0,1] \to \mathbb{R}$ is bounded. This is the case, since $f_n$ is bounded for all
    $n \in \mathbb{N}$, and with $\varepsilon = 1$ we find $N \in \mathbb{N}$ such that
    \[
    \sup_{x \in [0,1]} |f_N(x) - f(x)| < 1,
    \]
    which implies that $f$ is bounded.
\end{enumerate}
\end{exercisesolution}

\begin{center}
\begin{tikzpicture}[>=Latex, scale=1.4]
  \draw[->] (-0.2,0) -- (3.2,0) node[right] {$x$};
  \draw[->] (0,-0.2) -- (0,2.2) node[above] {$y$};
  \draw (0,0.05) -- (0,-0.05) node[below, font=\footnotesize] {$0$};
  \draw (2.6,0.05) -- (2.6,-0.05) node[below, font=\footnotesize] {$1$};

  \draw[blue!80!black, thick] plot[domain=0:2.6, samples=50] (\x, {1.0 + 0.3*sin(\x r * 2.5)});
  \node[blue!80!black, right] at (2.6, 1.06) {$f(x)$};

  \draw[red!80!black, thick, dotted] plot[domain=0:2.6, samples=50] (\x, {1.0 + 0.35 + 0.3*sin(\x r * 2.5)});
  \draw[red!80!black, thick, dotted] plot[domain=0:2.6, samples=50] (\x, {1.0 - 0.35 + 0.3*sin(\x r * 2.5)});

  \draw[orange!90!black, <->, thick] (1.3, 0.5) -- node[right, font=\footnotesize] {$\frac{2\varepsilon}{3}$} (1.3, 1.2);
  \node[red!80!black, font=\small, above] at (1.3, 1.8) {$B_{\varepsilon/3}(f)$ ($\varepsilon$-tube)};
\end{tikzpicture}
\end{center}


% ============================================================
% Chapter 5 --- Continuity
% Originally: content/week-02.tex, \section{Continuity} (Corsin Week 2, Friday)
%           + content/week-03.tex, the Lipschitz block that sat inside
%             \subsection{Why do we care about compactness?} (Corsin Week 3, Monday)
% ============================================================
\chapter{Continuity}
\label{ch:continuity}

% Transition: Claude Opus 5 (High)
% Lipschitz continuity is drawn from a second tutor's slides. It reached Corsin's own notes only
% obliquely, as a constant used for contraction mappings; sec:lipschitz_continuity names it
% properly and places it in the hierarchy.
With distance, shape and convergence in place, we can finally ask the question the whole
apparatus was built for. What does it mean for a map \emph{between} two metric spaces to respect
that structure?

Three answers are given, and they are equivalent: the $\varepsilon$--$\delta$ condition inherited
from \emph{Analysis I}, a sequential condition, and a topological one asking that preimages of
open sets be open. Having all three is not redundance. The first is what one verifies for a
concrete function, the second is what combines with the sequential arguments of the previous
chapter, and the third is the only one that survives when the metric is taken away. The chapter
then adds the two strengthenings needed later, uniform continuity and Lipschitz continuity, and
locates them in a hierarchy of increasingly demanding conditions.

% Originally: content/week-02.tex, \section{Continuity} (Corsin Week 2, Friday)

% Source: Corsin Nick/Class Notes/Week 2.pdf, pp. 1--13
\section{Continuity}
\label{sec:continuity}

\begin{definition}[Continuous map]
\label{def:continuous}
For $(X, d_X)$ and $(Y, d_Y)$ metric spaces, a function $f : X \to Y$ is called \newterm{continuous}
\germanfor{stetig} if one of the following equivalent conditions holds:
\begin{enumerate}[label=\textbf{(\alph*)}]
  \item \textbf{($\varepsilon$--$\delta$):} $\forall x_0 \in X, \varepsilon > 0 \ \exists \delta > 0$ such that for all $x \in X$:
    \[d(x_0, x) < \delta \implies d(f(x_0), f(x)) < \varepsilon\]
  \item \textbf{(Sequentially continuous):} for any sequence $(x_n)_{n=0}^{\infty}$ in $X$ with
    limit $x \in X$, $\lim_{n\to\infty} f(x_n) = f(x)$ (and the limit exists).
  \item \textbf{(Topological):} for all $U \subseteq Y$ open, $f^{-1}(U) \subseteq X$ is open
    (\cref{def:open_closed}).
\end{enumerate}
\end{definition}

\begin{ainote}
Corsin writes $d(x_0,x) < \varepsilon$ in the ($\varepsilon$--$\delta$) hypothesis above, which
must read $\delta$, since otherwise the quantified $\delta$ is never used at all. This is a slip
of the pen, and it is corrected in \cref{def:continuous}.
\end{ainote}

% Generator: Claude Opus 5 (High) --- \cref{def:continuous} asserts the equivalence of three
% conditions without proving it; the cycle is short and uses only results already available.
\begin{proposition}[The three definitions of continuity agree]
\label{prop:continuity_equivalences}
The conditions \textbf{(a)}, \textbf{(b)} and \textbf{(c)} of \cref{def:continuous} are
equivalent, so \cref{def:continuous} is well posed.
\end{proposition}

\begin{proof}
We prove the cycle
$\text{\textbf{(a)}} \Rightarrow \text{\textbf{(b)}} \Rightarrow \text{\textbf{(c)}}
\Rightarrow \text{\textbf{(a)}}$.

\textbf{(a) $\Rightarrow$ (b).} Let $x_n \to x$ and let $\varepsilon>0$. Pick $\delta>0$ as in
\textbf{(a)} for the point $x$, then $N$ with $d_X(x_n,x)<\delta$ for $n \geq N$. For those $n$,
$d_Y(f(x_n),f(x)) < \varepsilon$, so $f(x_n) \to f(x)$.

\textbf{(b) $\Rightarrow$ (c).} Let $U \subseteq Y$ be open and suppose $f^{-1}(U)$ were not.
By the shrinking-ball construction (\cref{rem:shrinking_ball_trick}) there are
$x \in f^{-1}(U)$ and points $x_n \in B_{1/n}(x)$ with $x_n \notin f^{-1}(U)$. Then $x_n \to x$,
so $f(x_n) \to f(x) \in U$ by \textbf{(b)}, and since $U$ is open,
\cref{prop:open_closed_via_sequences}\textbf{(a)} forces $f(x_n) \in U$ for large $n$. That is,
$x_n \in f^{-1}(U)$, which is a contradiction.

\textbf{(c) $\Rightarrow$ (a).} Fix $x_0 \in X$ and $\varepsilon>0$. The ball
$B_\varepsilon(f(x_0)) \subseteq Y$ is open (\cref{ex:ai_open_balls_are_open}), so
$f^{-1}\big(B_\varepsilon(f(x_0))\big)$ is open by \textbf{(c)} and contains $x_0$. Hence there
is $\delta>0$ with $B_\delta(x_0) \subseteq f^{-1}\big(B_\varepsilon(f(x_0))\big)$, which says
precisely that $d_X(x_0,x)<\delta$ implies $d_Y(f(x_0),f(x))<\varepsilon$.
\end{proof}

% Generator: Claude Opus 5 (High) --- FIG-05-01. The chapter defines continuity three ways and
% had no picture at all; its only figure was the Lipschitz cone, two sections later.
%
% GEOMETRY. Nothing here is a computation, so there is no arithmetic to check; what matters is
% that the drawing asserts only true things. Two deliberate choices:
%   * The image of B_delta(x_0) is drawn as an irregular blob strictly INSIDE B_epsilon(f(x_0)),
%     not as a disc and not filling it. Continuity says the image is contained in the target
%     ball; it says nothing about the image being round, being centred, or exhausting the ball.
%     A figure drawing f(B_delta) as a concentric disc would assert all three.
%   * The blob touches neither the boundary of B_epsilon nor f(x_0)'s dot, since the inclusion
%     is into the OPEN ball and f(x_0) is interior to the image, not on its edge.
% The same picture reads condition (c) backwards: B_delta sits inside the preimage of the open
% ball B_epsilon, which is the containment the proof of (c) => (a) produces.
\begin{center}
\begin{tikzpicture}[>=Latex, scale=1.0]
  % ---- the source space ----
  \draw[gray!55] plot[smooth cycle, tension=0.75] coordinates
    {(-2.5,1.15) (-0.85,1.5) (0.15,0.55) (-0.35,-1.2) (-1.9,-1.45) (-2.95,-0.25)};
  \node[gray!55!black, font=\small] at (-2.72,1.28) {$X$};
  \fill[MidnightBlue!8] (-1.45,0.05) circle (0.62);
  \draw[MidnightBlue!70, thick] (-1.45,0.05) circle (0.62);
  \fill[MidnightBlue!70] (-1.45,0.05) circle (1.4pt);
  \node[MidnightBlue!70, font=\scriptsize, anchor=north east] at (-1.5,-0.02) {$x_0$};
  \node[MidnightBlue!70, font=\scriptsize, anchor=south] at (-1.45,0.70) {$B_\delta(x_0)$};

  % ---- the map ----
  \draw[->, thick] (0.55,0.05) -- (1.95,0.05) node[midway, above, font=\small] {$f$};

  % ---- the target space ----
  \draw[gray!55] plot[smooth cycle, tension=0.75] coordinates
    {(2.6,1.3) (4.4,1.5) (5.5,0.4) (4.9,-1.3) (3.3,-1.5) (2.35,-0.3)};
  \node[gray!55!black, font=\small] at (5.35,1.15) {$Y$};
  \fill[BrickRed!7] (3.95,0.05) circle (0.78);
  \draw[BrickRed!80, thick] (3.95,0.05) circle (0.78);
  \node[BrickRed!80, font=\scriptsize, anchor=south] at (3.95,0.86)
    {$B_\varepsilon(f(x_0))$};

  % The image: irregular, strictly inside, not centred. All three are deliberate.
  \fill[OliveGreen!14] plot[smooth cycle, tension=0.8] coordinates
    {(3.72,0.42) (4.25,0.3) (4.38,-0.12) (3.95,-0.42) (3.55,-0.18)};
  \draw[OliveGreen, thick] plot[smooth cycle, tension=0.8] coordinates
    {(3.72,0.42) (4.25,0.3) (4.38,-0.12) (3.95,-0.42) (3.55,-0.18)};
  \fill[BrickRed!80] (3.95,0.02) circle (1.4pt);
  % Anchored below the dot and inside the blob. Anchored east it extended leftwards across the
  % blob's own outline at x = 3.55, so the label sat half in and half out of the region it names.
  \node[BrickRed!80, font=\scriptsize, anchor=north] at (3.95,-0.04) {$f(x_0)$};
  \node[OliveGreen, font=\scriptsize, anchor=north west] at (4.35,-0.35)
    {$f\big(B_\delta(x_0)\big)$};
\end{tikzpicture}
\end{center}

Condition \textbf{(a)} is the statement that the green region fits inside the red ball, and that
a $\delta$ making it fit can be found however small the red ball is made. Condition \textbf{(c)}
is the same picture read from right to left: $B_\delta(x_0)$ lies inside
$f^{-1}\big(B_\varepsilon(f(x_0))\big)$, which is exactly the containment produced in the
\textbf{(c)} $\Rightarrow$ \textbf{(a)} step above. Note what the drawing does not claim: the
image need not be round, need not be centred at $f(x_0)$, and need not fill the ball.

\begin{remark}[Which of the three to reach for]
Condition \textbf{(c)} is the only one that never mentions the metric. That is why it, and not the
$\varepsilon$--$\delta$ formulation, is what survives into general topology
(\cref{sec:structured_spaces}). Condition \textbf{(b)} is usually the fastest way to prove a
function is \emph{dis}continuous, since a single badly chosen sequence suffices; this is exactly
how \cref{ex:partial_derivatives_not_continuous} works. Condition \textbf{(a)} is the
one to use when a quantitative modulus is wanted, as in
\cref{def:uniformly_continuous} immediately below.
\end{remark}

% Sonnet 5 (Medium)
\begin{definition}[Uniform continuity]
\label{def:uniformly_continuous}
For $(X,d_X)$ and $(Y,d_Y)$ metric spaces, a function $f : X \to Y$ is \newterm{uniformly
continuous} \germanfor{gleichm\"assig stetig} if
\[\forall \varepsilon > 0 \ \exists \delta > 0 \ \forall x, x' \in X : \quad d_X(x,x') < \delta \ \text{ implies } \ d_Y(f(x),f(x')) < \varepsilon.\]
\end{definition}

\begin{importantremark}[Where the quantifiers sit]
\label{rem:continuity_vs_uniform_continuity}
Compare \cref{def:continuous} \textbf{(a)} with \cref{def:uniformly_continuous}: the
\emph{formulas} differ only in the order of two quantifiers, and that order is the whole
content.
\begin{itemize}
  \item \textbf{Continuity:} $\forall x_0\ \forall\varepsilon\ \exists\delta$. The point $x_0$
    is chosen \emph{before} $\delta$, so $\delta$ may depend on it, and a different $\delta$ is
    allowed at every point.
  \item \textbf{Uniform continuity:} $\forall\varepsilon\ \exists\delta\ \forall x$. Now
    $\delta$ is chosen \emph{before} seeing the point, so one and the same $\delta$ must work
    everywhere at once.
\end{itemize}
Uniform continuity is therefore strictly stronger. The standard witness is
$f(x) := 1/x$ on $(0,1)$: it is continuous, but near $0$ the graph steepens without bound, so
the $\delta$ that works at $x_0 = \tfrac12$ is hopeless at $x_0 = \tfrac{1}{1000}$, and no
single $\delta$ survives. \Cref{prop:compact_implies_uniformly_continuous} below says that on a
\emph{compact} space this can never happen. The gap between the two notions closes exactly when
there is no room left to run off towards a bad point.
\end{importantremark}

\begin{example}[$f(x) = x^2$]
Let $f : \mathbb{R} \to \mathbb{R}, x \mapsto x^2$. Let $a < b \in \mathbb{R}$. Then
\begin{itemize}
  \item for $a < 0 < b$: $\quad f^{-1}((a,b)) = (-\sqrt{b}, \sqrt{b})$, which is open.
  \item for $a < b < 0$: $\quad f^{-1}((a,b)) = \emptyset$.
  \item for $0 < a < b$: $\quad f^{-1}((a,b)) = (-\sqrt{b}, -\sqrt{a}) \cup (\sqrt{a}, \sqrt{b})$.
\end{itemize}
\textit{Why is it enough to check only these cases?} Hint:
$f^{-1}(U_1 \cup U_2) = f^{-1}(U_1) \cup f^{-1}(U_2)$.
\end{example}

% Source: Jérôme Paschoud/Notizen/Woche 4.1 Normen, das Differential.pdf, p. 2
% Generator: Gemini 3.1 Pro (High)
\begin{example}[Limits in $\mathbb{R}^2$ depend on the path]
\label{ex:limits_in_r2}
When computing limits in $\mathbb{R}^n$, one must ensure that the limit exists and is identical along \emph{all} possible paths approaching the point. For example, consider the function $f: \mathbb{R}^2 \setminus \{(0,0)\} \to \mathbb{R}$ given by $f(x,y) = \frac{x^2}{x^2+y^2}$.
If we approach the origin along the $x$-axis (setting $y=0$ first), we get
\[
\lim_{x \to 0} f(x, 0) = \lim_{x \to 0} \frac{x^2}{x^2} = 1.
\]
However, if we approach along the $y$-axis (setting $x=0$ first), we get
\[
\lim_{y \to 0} f(0, y) = \lim_{y \to 0} \frac{0}{y^2} = 0.
\]
Since these limits disagree, $\lim_{(x,y) \to (0,0)} f(x,y)$ does not exist. A robust method to check limits at the origin uniformly across all directions is to switch to \newterm{polar coordinates} ($x = r\cos\varphi, y = r\sin\varphi$). If the limit as $r \to 0$ is independent of $\varphi$, then the function converges. (This is explored further in \cref{ex:3.2}).
\end{example}

\begin{question}
Let $(X, d_X)$ and $(Y, d_Y)$ be metric spaces, where $d_Y$ is the discrete metric. Is
\emph{every} function $f : (X,d_X) \to (Y,d_Y)$ continuous?
\end{question}

\begin{answer}
\textbf{No.} For example, take the spaces $X = Y = \mathbb{R}$, $d_X$ the standard metric, and
$f = \id : X \to Y, z \mapsto z$. Then $\{0\} \subseteq Y$ is open, but
$\{0\} = \id^{-1}(\{0\})$ is not open in $(X, d_X)$.
\end{answer}

\begin{question}
And what if the roles are reversed, with $d_X$ the discrete metric and $d_Y$ arbitrary?
\end{question}

\begin{answer}
\textbf{Any} subset of $X$ is open; therefore, $f^{-1}(U)$ is open for any open $U \subseteq Y$. So
every such $f$ is continuous.
\end{answer}

% The old chapter continued from here into "Completeness"; that topic is now
% \cref{ch:completeness}, whose opening paragraph carries this transition.


% --- exercises relocated here from the dissolved sheets ---

% Originally: content/exercise-sheets/week-02.tex, problem 2.6
% Source: exercises/Ex2_Analysis2_eng.pdf, p. 1

\begin{exercise}[Continuity of the distance]
\label{ex:2.6}
Let $(X,d)$ be a metric space, and endow $X \times X$ with the product distance
$d_2(x,y) := d(x_1,y_1) + d(x_2,y_2)$. Show that the distance function
$d : X \times X \to \mathbb{R}$ is continuous with respect to $d_2$ and, more precisely, that it
is 1-Lipschitz. Is $d$ also continuous with respect to the distance
$d_1(x,y) := \max\{d(x_1,y_1), d(x_2,y_2)\}$? Is $d$ also Lipschitz continuous with respect to the
distance $d_3(x,y) := \sqrt{d(x_1,y_1)^2 + d(x_2,y_2)^2}$? (Notation consistent with
\cref{ex:2.5}.)
\exinfo{This exercise is Problem 2.6 of Problem Sheet 2, Analysis II, Spring Semester 2026.}
\end{exercise}

% Originally: content/exercise-sheets/week-02.tex, problem 2.7
% Source: exercises/Ex2_Analysis2_eng.pdf, p. 1

\begin{exercise}[Continuity of the composition]
\label{ex:2.7}
Let $X, Y, Z$ be metric spaces and let $f : X \to Y$, $g : Y \to Z$ be continuous functions. Show
that $g \circ f : X \to Z$ is continuous using at least two of the three equivalent definitions of
continuity seen in class.
\exinfo{This exercise is Problem 2.7 of Problem Sheet 2, Analysis II, Spring Semester 2026.}
\end{exercise}

\begin{remark}
The three equivalent definitions of continuity referenced in \cref{ex:2.7} are exactly the ones
stated in \cref{sec:continuity} below (\(\varepsilon\)--\(\delta\), sequential, topological).
\end{remark}

% Originally: content/exercise-sheets/week-03.tex, problem 3.2
% Source: exercises/Ex3_Analysis2_eng.pdf

\begin{exercise}[Problems at the origin]
\label{ex:3.2}
Show that there is no continuous function $f : \mathbb{R}^2 \to \mathbb{R}$ such that
$f(x,y) = xy/(x^2+y^2)$ for all $(x,y) \neq (0,0)$. On the other hand, show that there is exactly
one continuous function $g : \mathbb{R}^2 \to \mathbb{R}$ such that
$g(x,y) = xy/\sqrt{x^2+y^2}$ for all $(x,y) \neq (0,0)$.

\exhint[Official hint]{If a function is continuous at $0$ and $x_k \to 0$ and $y_k \to 0$, then
$f(x_k)$ and $f(y_k)$ must converge to the same value.}
\exinfo{This exercise is Problem 3.2 of Problem Sheet 3, Analysis II, Spring Semester 2026.}
\end{exercise}

% Quelle: Jérôme Paschoud/Notizen/Woche 2.2 Topologie und Stetigkeit.pdf, p. 2
\begin{exercise}[Distance to a point is continuous]
\label{ex:distance_to_point_continuous}
Let $(X, d)$ be a metric space and fix $a \in X$. Show that the function $f : X \to [0,\infty)$ defined by $f(x) := d(x, a)$ is continuous.
\end{exercise}

% Generator: Gemini 3.6 Pro
\begin{exercisesolution}[Proof of \cref{ex:distance_to_point_continuous}]
We use $\varepsilon$--$\delta$ continuity. Let $x \in X$ and $\varepsilon > 0$. We need to find $\delta > 0$ such that for all $x' \in X$ with $d(x, x') < \delta$, we have $|f(x) - f(x')| < \varepsilon$.
Using the reverse triangle inequality, we have
\[|f(x) - f(x')| = |d(x, a) - d(x', a)| \leq d(x, x').\]
Therefore, we can simply choose $\delta := \varepsilon$. Then $d(x, x') < \delta$ directly implies $|f(x) - f(x')| < \varepsilon$, meaning $f$ is continuous at $x$. Since $x$ was arbitrary, $f$ is continuous everywhere.
\end{exercisesolution}

% Supplement: Analysis_II_Script_v1.pdf, p. 19
\begin{corollary}[Closed sets are exactly the zero sets of a distance function]
\label{cor:closed_iff_zero_of_distance}
Let $(X,d)$ be a metric space and $\emptyset \neq E \subseteq X$. The function
$f_E(x) := d(x,E) = \inf_{z\in E}d(x,z)$ is $1$-Lipschitz, which is \cref{ex:2.3}\textbf{(c)},
proved there by the same reverse-triangle-inequality argument as
\cref{ex:distance_to_point_continuous} above. Moreover,
\[E \text{ is closed} \quad\Longleftrightarrow\quad E = \{x \in X : f_E(x) = 0\}.\]
\end{corollary}

\begin{proof}
The inclusion $E \subseteq \{f_E = 0\}$ always holds, since $z \in E$ gives
$f_E(z) \leq d(z,z) = 0$. So only the reverse inclusion carries content, and it is exactly where
closedness is needed.

\textbf{($\Rightarrow$)} Let $E$ be closed and $f_E(x) = 0$. By definition of the infimum, choose
$z_n \in E$ with $d(x,z_n) \to 0$; then $z_n \to x$, and closedness of $E$ gives $x \in E$
(\cref{prop:open_closed_via_sequences}\textbf{(b)}).

\textbf{($\Leftarrow$)} Suppose $E = \{f_E = 0\}$ and let $(z_n)$ be a sequence in $E$ with
$z_n \to x$. Then $f_E(x) \leq d(x,z_n) \to 0$, so $f_E(x) = 0$, i.e.\ $x \in E$ by hypothesis.
\Cref{prop:open_closed_via_sequences}\textbf{(b)} then gives that $E$ is closed.
\end{proof}

\begin{remark}
This sharpens \cref{ex:2.3}\textbf{(a)}, which showed only one direction ($E$ closed
$\Rightarrow$ $f_E > 0$ off $E$) for a general subset. It is exactly closedness, and nothing
weaker, that lets $f_E$ detect membership in $E$ by vanishing.
\end{remark}

% Quelle: Jérôme Paschoud/Notizen/Woche 2.2 Topologie und Stetigkeit.pdf, p. 3
\begin{exercise}[Evaluation map is continuous]
\label{ex:evaluation_map_continuous}
Let $X := C^0([0,1])$ equipped with the supremum metric $d(f,g) := \sup_{t \in [0,1]} |f(t) - g(t)|$. Define the function $F : X \to \mathbb{R}$ by $F(f) := f(0)$. Show that $F$ is continuous.
\end{exercise}

% Generator: Gemini 3.6 Pro
\begin{exercisesolution}[Proof of \cref{ex:evaluation_map_continuous}]
Let $f \in X$ and $\varepsilon > 0$. Choose $\delta := \varepsilon$. If $g \in X$ satisfies $d(f, g) < \delta$, then by the definition of the supremum metric, we have $|f(t) - g(t)| < \delta$ for all $t \in [0,1]$.
In particular, for $t = 0$, we get
\[|F(f) - F(g)| = |f(0) - g(0)| \leq \sup_{t \in [0,1]} |f(t) - g(t)| = d(f, g) < \delta = \varepsilon.\]
Thus $F$ is continuous at $f$. Since $f$ was arbitrary, $F$ is continuous on all of $X$.
\end{exercisesolution}

% Originally: content/week-03.tex, the Lipschitz block inside \subsection{Why do we care about compactness?}

\section{Lipschitz continuity}
\label{sec:lipschitz_continuity}

% Quelle: Linus Lüchinger/Slides/Slides-03-02.pdf, pp. 10--13 -- Sonnet 5 (Medium)
% Lüchinger's exercise slides motivate and name the condition and give a quick self-test; both
% are reproduced here, reformulated, since they fill a genuine gap in Corsin's own notes.
\begin{ainote}
The course never formally introduces Lipschitz continuity, even though a Lipschitz constant is
already used informally for contraction mappings (\cref{def:contraction_mapping}, below). The
definition and the hierarchy below are therefore supplied here, and will not be found under this
name in the lecture material.
\end{ainote}

\begin{definition}[Lipschitz continuity]
\label{def:lipschitz_continuous}
Let $(X,d_X)$, $(Y,d_Y)$ be metric spaces. A function $f : X \to Y$ is \newterm{Lipschitz
continuous} (\germanterm{Lipschitz-stetig}) if there exists $L \geq 0$, a \newterm{Lipschitz
constant} for $f$, such that
\[d_Y(f(x),f(x')) \leq L\cdot d_X(x,x') \quad \text{for all } x,x' \in X.\]
\end{definition}

\begin{remark}[Hierarchy of continuity notions]
\label{rem:continuity_hierarchy}
Lipschitz continuity is stronger than uniform continuity
(\cref{def:uniformly_continuous}), which in turn is stronger than continuity
(\cref{def:continuous}). For the first implication, assume $L > 0$ (for $L = 0$ the function is
constant and there is nothing to prove); then $\delta := \varepsilon/L$ satisfies the definition
of uniform continuity, since
\[d_X(x,x') < \frac{\varepsilon}{L} \quad \text{gives} \quad d_Y(f(x),f(x')) \leq L\,d_X(x,x') < \varepsilon\]
for \emph{every} pair $x,x'$. This is one and the same $\delta$ at every point, which is exactly
what uniform continuity asks for. The second implication is immediate: a $\delta$ that works at
all points in particular works at each one.

Both converses fail. Every continuously differentiable function on a compact \emph{interval} is
Lipschitz continuous, with $L := \sup|f'|$ by the mean value theorem. The word \emph{interval}
matters here, since the mean value theorem needs the segment between the two points to lie in the
domain. But $f(x) := \sqrt{x}$ on $[0,\infty)$ is uniformly continuous and \emph{not} Lipschitz.
A Lipschitz constant would force $\delta$ to shrink only linearly in $\varepsilon$, whereas here
$|\sqrt{x}-\sqrt{x'}| < \varepsilon$ near $0$ already requires $\delta$ of order
$\varepsilon^2$. The reason is visible in the graph, since the tangent at $0$ is vertical and no
finite slope $L$ can dominate it.
\end{remark}

% Sonnet 5 (Medium)
\begin{center}
\begin{tikzpicture}[scale=1.0]
  \draw[blue!70!black, thick] (0,0) circle (2.2);
  \node[blue!70!black, font=\small] at (0,-1.9) {continuous};
  \draw[orange!90!black, thick] (0,0) circle (1.5);
  \node[orange!90!black, font=\small] at (0,-1.2) {uniformly cont.};
  \draw[purple, thick] (0,0) circle (0.8);
  \node[purple, font=\small] at (0,0) {Lipschitz};
  \fill[orange!90!black] (1.1,0.9) circle (1.5pt) node[right, font=\scriptsize] {$\sqrt{x}$};
\end{tikzpicture}
\end{center}

% Supplement: Linus Lüchinger/Slides/Slides-03-02.pdf, slide 13
% Was an `aiexercise` labelled ex:ai_continuity_classification. It is Linus's exercise, not
% one invented here, so by the ai*-prefix rule it is a plain `exercise` with a provenance
% comment; only the solution below is authored. His domain for the first row was [0.5,1];
% widened to [0,infty) so that the row separates uniform continuity from Lipschitz, which is
% the distinction \cref{rem:continuity_hierarchy} has just drawn.
\begin{exercise}[Classifying continuity]
\label{ex:continuity_classification}
For each of the following functions, determine whether it is continuous, uniformly continuous,
and/or Lipschitz continuous on the given domain:
\begin{center}
\begin{tabular}{ll}
\textbf{Function} & \textbf{Domain} \\
\hline
$\sqrt{x}$ & $[0,\infty)$ \\
$1/x$ & $(0,1)$ \\
$\sgn(x)$ & $(-1,1)$ \\
$|x|^{3/2}$ & $(-1,1)$
\end{tabular}
\end{center}
\end{exercise}


% --- exercises relocated here from the dissolved sheet 3 ---

% Originally: content/exercise-sheets/week-03.tex, problem 3.1
% Source: exercises/Ex3_Analysis2_eng.pdf

\begin{exercise}[3.1 --- Lipschitz or not]
\label{ex:3.1}
Give an example of:
\begin{enumerate}[label=\textbf{(\alph*)}]
  \item A $\tfrac{1}{2}$-Lipschitz function $f : [0,1) \to [0,1)$ that has no fixed point.
  \item A map $f : \mathbb{R}^2 \to \mathbb{R}^2$ that has no fixed point, but is an isometry,
    i.e.\ $\|f(x) - f(x')\| = \|x - x'\|$ for all $x, x' \in \mathbb{R}^2$.
\end{enumerate}
\end{exercise}

% Originally: content/exercise-sheets/week-03.tex, problem 3.6
% Source: exercises/Ex3_Analysis2_eng.pdf

\begin{exercise}[3.6 --- Multiple choice (uniform continuity)]
\label{ex:3.6}
Select all the statements below that are true.
\begin{enumerate}[label=\textbf{(\alph*)}]
  \item The function $(x,y) \mapsto x+y$ is uniformly continuous in $\mathbb{R}^2$.
  \item The function $(x,y) \mapsto xy$ is uniformly continuous in $\mathbb{R}^2$.
  \item The function $(x,y) \mapsto x+y$ is uniformly continuous in $[0,1]^2$.
  \item The function $(x,y) \mapsto xy$ is uniformly continuous in $[0,1]^2$.
\end{enumerate}
\end{exercise}

\begin{remark}
Corsin's third reason to care about compactness (\cref{sec:why_compactness_matters}) is precisely
that a continuous function on a compact space is uniformly continuous, which settles
\cref{ex:3.6}(c) and (d) at once.
\end{remark}

% Originally: content/exercise-sheets/week-03.tex, problem 3.11
% Source: exercises/Ex3_Analysis2_eng.pdf

\begin{exercise}[3.11 --- (*) Uniform continuity and moduli of continuity]
\label{ex:3.11}
Let $(X,d_X)$ and $(Y,d_Y)$ be metric spaces and let $f : X \to Y$.

A \textbf{modulus of continuity} is a function $\omega : [0,\infty) \to [0,\infty]$ such that
$$\omega(0) = 0, \qquad \omega \text{ is nondecreasing}, \qquad \lim_{t \downarrow 0}\omega(t) = 0.$$
(Here $[0,\infty]$ means we also allow the value $+\infty$.) We say that $f$ \textbf{has modulus
of continuity} $\omega$ if
$$d_Y\big(f(x), f(y)\big) \leq \omega\big(d_X(x,y)\big) \qquad \forall x,y \in X.$$
\begin{enumerate}[label=\textbf{\arabic*.}]
  \item Prove that $f$ is uniformly continuous on $X$ if and only if $f$ has a modulus of
    continuity.
  \item \textit{(Bonus 1)} Assume $X \subset \mathbb{R}$ is an interval with the usual distance
    $d_X(x,y) = |x-y|$. Show that the modulus constructed in (1) can be chosen
    \textbf{subadditive}, i.e.\ $\omega(t_1+t_2) \leq \omega(t_1) + \omega(t_2)$ for all
    $t_1, t_2 \geq 0$.
  \item \textit{(Bonus 2)} Assume $X \subset \mathbb{R}^n$ is \textbf{convex}, meaning: for all
    $x_1, x_2 \in X$ and all $\lambda \in [0,1]$, $(1-\lambda)x_1 + \lambda x_2 \in X$. With
    $d_X(x,y) = \|x-y\|$ (Euclidean distance), show again that the modulus from (1) can be chosen
    subadditive.
\end{enumerate}

\begin{remark}[official]
In general metric spaces a modulus of continuity may have infinite value for finite $t > 0$.
Example: $X = \mathbb{Z}$ with the usual distance, $Y = \mathbb{R}$, $f(n) = n^2$. This $f$ is
uniformly continuous on $\mathbb{Z}$ (because points closer than 1 must be equal), but no finite
$\omega(1)$ can bound $|f(n+1) - f(n)| = 2n+1$ for all $n$. Allowing the value $+\infty$ avoids
this issue, and in many familiar settings (e.g.\ intervals/convex sets in $\mathbb{R}^n$) the
resulting $\omega_f(t)$ is finite for all $t$.
\end{remark}

\textit{Official hints:} \textbf{3.11.1:} define the ``worst oscillation at scale $t$''
$\omega_f(t) := \sup\{d_Y(f(x),f(y)) : d_X(x,y) \leq t\}$. \textbf{3.11.2 / 3.11.3:} if
$d_X(x,y) \leq t_1 + t_2$, choose an intermediate point $z$ on the segment from $x$ to $y$ so
that $d_X(x,z) \leq t_1$ and $d_X(z,y) \leq t_2$, then use the triangle inequality in $Y$.
\end{exercise}

% Originally: the \section{Solutions} of content/week-02.tex and content/week-03.tex

\section{Solutions}
\label{sec:continuity_solutions}

% Transition: Claude Opus 5 (High)
Not every solution in this chapter is collected here. \Cref{ex:distance_to_point_continuous}
(\cpageref{ex:distance_to_point_continuous}) and \cref{ex:evaluation_map_continuous}
(\cpageref{ex:evaluation_map_continuous}) are solved where they appear.

\begin{exercisesolution}[Solution to \cref{ex:2.6}]
% Generator: Claude Opus 5 (High)
\textbf{$d$ is $1$-Lipschitz with respect to $d_2$.} Let $x = (x_1,x_2)$ and $y = (y_1,y_2)$ be
points of $X\times X$. Applying the triangle inequality in $X$ twice,
\[d(x_1,x_2) \ \leq\ d(x_1,y_1) + d(y_1,y_2) + d(y_2,x_2),\]
so $d(x_1,x_2) - d(y_1,y_2) \leq d(x_1,y_1) + d(x_2,y_2) = d_2(x,y)$. Exchanging the roles of
$x$ and $y$ bounds the opposite difference in the same way, hence
\[\big|d(x_1,x_2) - d(y_1,y_2)\big| \ \leq\ d_2(x,y),\]
which is $1$-Lipschitz continuity (\cref{def:lipschitz_continuous}) and therefore continuity
(\cref{rem:continuity_hierarchy}).

\textbf{For $d_1$ and $d_3$: yes, still Lipschitz, but the constant changes.} By
\cref{ex:2.5} part 2 we have $d_2 \leq 2d_1$ and $d_2 \leq \sqrt2\,d_3$, so
\[\big|d(x_1,x_2)-d(y_1,y_2)\big| \leq 2\,d_1(x,y)
  \qquad\text{and}\qquad
  \big|d(x_1,x_2)-d(y_1,y_2)\big| \leq \sqrt2\,d_3(x,y).\]
This is a general principle rather than a computation: \textbf{Lipschitz continuity survives
passage to an equivalent metric}, with the constant multiplied by the comparison factor.
Continuity survives too, being topological, and by \cref{ex:2.5} the three product metrics
induce the same topology.

% Reclassified from ainote to remark: sharpness of the constant is mathematics.
\begin{remark}[The constant $2$ cannot be improved]
The constant $2$ for $d_1$ is sharp, so $d$ is genuinely \emph{not} $1$-Lipschitz with respect
to $d_1$. Take $X = \mathbb{R}$ with $x := (0,0)$ and $y := (-1,1)$. Then $d(x_1,x_2) = 0$ and
$d(y_1,y_2) = 2$, while $d_1(x,y) = \max\{1,1\} = 1$, so the difference equals exactly
$2 = 2\,d_1(x,y)$. This is worth checking, because the problem asks only \emph{whether} $d$ is
Lipschitz for $d_1$ and $d_3$, and it is easy to answer \qt{yes, with the same constant} without
noticing that the constant has in fact changed.
\end{remark}
\end{exercisesolution}

\begin{exercisesolution}[Solution to \cref{ex:2.7}]
% Generator: Claude Opus 5 (High)
All three arguments are short; the point of giving more than one is to see how differently the
same fact reads through the three definitions of \cref{def:continuous}.

\textbf{Topological (\cref{def:continuous}\textbf{(c)}).} This is the shortest, because
preimages compose. Let $W \subseteq Z$ be open. Then
\[(g\circ f)^{-1}(W) = f^{-1}\big(g^{-1}(W)\big),\]
and $g^{-1}(W)$ is open in $Y$ by continuity of $g$, so its preimage under $f$ is open in $X$ by
continuity of $f$. Hence $g\circ f$ is continuous.

\textbf{Sequential (\cref{def:continuous}\textbf{(b)}).} Let $x_n \to x$ in $X$. Continuity of
$f$ gives $f(x_n) \to f(x)$ in $Y$, and continuity of $g$ applied to \emph{that} sequence gives
$g(f(x_n)) \to g(f(x))$ in $Z$.

\textbf{$\varepsilon$--$\delta$ (\cref{def:continuous}\textbf{(a)}).} Fix $x_0 \in X$ and
$\varepsilon > 0$. Continuity of $g$ at $f(x_0)$ supplies $\eta > 0$ such that
$d_Y(f(x_0),y) < \eta$ implies $d_Z\big(g(f(x_0)),g(y)\big) < \varepsilon$. Continuity of $f$ at
$x_0$, applied to \emph{that} $\eta$, supplies $\delta > 0$ such that $d_X(x_0,x) < \delta$
implies $d_Y(f(x_0),f(x)) < \eta$. Chaining the two, $d_X(x_0,x) < \delta$ implies
$d_Z\big((g\circ f)(x_0),(g\circ f)(x)\big) < \varepsilon$.

\begin{remark}[Why the topological proof is the short one]
All three arguments have the same shape, applying the hypothesis for $g$ first and then feeding
the result to $f$. Only the topological version needs no quantifiers at all, because the set
identity $(g\circ f)^{-1} = f^{-1}\circ g^{-1}$ does the bookkeeping the other two proofs carry
out by hand. This is a fair advertisement for \cref{def:continuous}\textbf{(c)}, which can look
like the least intuitive of the three on first meeting: it is the formulation in which
composition is trivial, and that is exactly why general topology takes it as \emph{the}
definition.
\end{remark}
\end{exercisesolution}

\begin{exercisesolution}[Solution to \cref{ex:continuity_classification}]
% Generator: Claude Sonnet 5 (Medium)
\begin{itemize}
  \item $\sqrt{x}$ on $[0,\infty)$: \textbf{uniformly continuous, not Lipschitz.} Uniform
    continuity follows from continuity on $[0,1]$ (compact, \cref{item:uniform_continuity_compact})
    together with $\tfrac12$-Lipschitz behaviour on $[1,\infty)$ (since $|\sqrt x - \sqrt y| \leq
    \tfrac12|x-y|$ there by the mean value theorem, as $|(\sqrt{\cdot})'| = \tfrac{1}{2\sqrt x}
    \leq \tfrac12$ for $x \geq 1$). Not Lipschitz near $0$: for $x_n :=
    1/n^2 \to 0$, $\dfrac{|\sqrt{x_n}-\sqrt0|}{|x_n-0|} = n \to \infty$, so no global $L$ works.
  \item $1/x$ on $(0,1)$: \textbf{continuous, neither uniformly nor Lipschitz continuous.} For
    $x_n := 1/n$, $y_n := 1/(2n)$, $|x_n-y_n| \to 0$ but $|1/x_n - 1/y_n| = n \to \infty$, so no
    uniform $\delta(\varepsilon)$ (and hence no Lipschitz constant) can work near $0$.
  \item $\sgn(x)$ on $(-1,1)$: \textbf{none of the three}, being discontinuous at $0$.
  \item $|x|^{3/2}$ on $(-1,1)$: \textbf{Lipschitz, hence also uniformly continuous and
    continuous.} $|x|^{3/2}$ is $C^1$ on $(-1,1)\setminus\{0\}$ with $\big|\tfrac{d}{dx}|x|^{3/2}\big|
    = \tfrac32|x|^{1/2} \leq \tfrac32$, and the bound extends across $0$ by direct estimation, so
    $L := \tfrac32$ works globally.
\end{itemize}
\end{exercisesolution}

\begin{exercisesolution}[Solution to \cref{ex:3.1}]
% Generator: Claude Opus 5 (High)
Both parts ask for a map with no fixed point, and each defeats one hypothesis of
\cref{thm:banach_fixed_point} while keeping the other, which is exactly what
\cref{rem:banach_hypotheses} predicts must be possible.
\begin{enumerate}[label=\textbf{(\alph*)}]
  \item Take $f(x) := \tfrac{x+1}{2}$. It maps $[0,1)$ into $[\tfrac12,1) \subseteq [0,1)$, and
    $|f(x)-f(y)| = \tfrac12|x-y|$, so it is $\tfrac12$-Lipschitz. A fixed point would satisfy
    $x = \tfrac{x+1}{2}$, i.e.\ $x = 1$, which is not in $[0,1)$. \textbf{What fails is
    completeness:} $[0,1)$ is not a complete metric space, the iterates
    $x_{n+1} = f(x_n)$ march towards the missing endpoint $1$, and \cref{thm:banach_fixed_point}
    does not apply.
  \item Take the translation $f(x) := x + (1,0)$. It is an isometry, since
    $|f(x)-f(x')| = |x-x'|$ exactly, and $f(x) = x$ would force $(1,0) = 0$. \textbf{What
    fails is the strict contraction:} an isometry has Lipschitz constant $L = 1$, not $L < 1$,
    and $\mathbb{R}^2$ is perfectly complete. Together with part \textbf{(a)} this shows neither
    hypothesis of the fixed point theorem can be dropped.
\end{enumerate}
\end{exercisesolution}

\begin{exercisesolution}[Solution to \cref{ex:3.2}]
% Generator: Claude Opus 5 (High)
\textbf{No continuous extension exists for $xy/(x^2+y^2)$.} Suppose $f$ were continuous on
$\mathbb{R}^2$ and agreed with $xy/(x^2+y^2)$ off the origin. Approach the origin along two
different lines:
\[f\!\left(\tfrac1k,\tfrac1k\right) = \frac{1/k^2}{2/k^2} = \frac12,
  \qquad\qquad
  f\!\left(\tfrac1k,0\right) = 0 \quad\text{for every } k.\]
Both sequences tend to $(0,0)$, so continuity would force
$f(0,0) = \tfrac12$ and $f(0,0) = 0$ simultaneously. Hence no such $f$ exists. (This is the same
function, and the same argument, as \cref{ex:partial_derivatives_not_continuous},
where it reappears to show that the existence of both partial derivatives implies nothing.)

\textbf{Exactly one continuous extension exists for $xy/\sqrt{x^2+y^2}$.} Write
$(x,y) = (r\cos\theta, r\sin\theta)$ with $r > 0$. Then
\[g(x,y) = \frac{r^2\cos\theta\sin\theta}{r} = r\cos\theta\sin\theta,
  \qquad\text{so}\qquad |g(x,y)| \leq \frac r2 = \frac{|(x,y)|}{2}.\]
The bound is uniform in $\theta$, so $g(x,y) \to 0$ as $(x,y) \to (0,0)$, and setting
$g(0,0) := 0$ produces a continuous function on all of $\mathbb{R}^2$.

\textit{Uniqueness} is the part worth spelling out, and it needs no computation at all: the value
at the origin is \emph{forced}. Since $(0,0)$ lies in the closure of
$\mathbb{R}^2\setminus\{(0,0)\}$, any two continuous extensions agree on a dense subset of
$\mathbb{R}^2$, and two continuous functions agreeing on a dense set agree everywhere. To see the
latter, take a sequence from the dense set converging to the point in question and apply
\cref{def:continuous}\textbf{(b)} to both. So the extension is unique whenever it exists, which
is also why the first half of the problem had only to exhibit \emph{two} conflicting limits.
\end{exercisesolution}

\begin{exercisesolution}[Solution to \cref{ex:3.6}]
% Generator: Claude Opus 5 (High)
\textbf{(a), (c) and (d) are true; only (b) is false.}

\textbf{(a) True.} Addition is Lipschitz, hence uniformly continuous
(\cref{rem:continuity_hierarchy}):
\[\big|(x+y)-(x'+y')\big| \leq |x-x'| + |y-y'| \leq \sqrt2\,\big|(x,y)-(x',y')\big|,\]
the last step by Cauchy--Schwarz. So $\delta := \varepsilon/\sqrt2$ works everywhere at once.

\textbf{(b) False.} Multiplication is continuous but not uniformly so on the unbounded domain
$\mathbb{R}^2$. Take
\[p_k := \left(k,\ \tfrac1k\right), \qquad q_k := \left(k,\ 0\right).\]
Then $|p_k - q_k| = \tfrac1k \to 0$, while the values are $k \cdot \tfrac1k = 1$ and
$0$, a constant gap of $1$. So no $\delta$ can work for $\varepsilon < 1$. The failure is exactly
the one described in \cref{rem:continuity_vs_uniform_continuity}: the function steepens without
bound as one runs off to infinity, so the local $\delta$ shrinks without a positive floor.

\textbf{(c) and (d) True,} and both for free: $[0,1]^2$ is closed and bounded, hence
\textbf{compact} by \cref{thm:heine_borel}, and a continuous function on a compact metric space
is automatically uniformly continuous by
\cref{prop:compact_implies_uniformly_continuous} (Corsin's third reason to care about
compactness, \cref{item:uniform_continuity_compact}). No estimate is needed for either.

\begin{remark}
The pair \textbf{(b)} versus \textbf{(d)} is the whole lesson: the \emph{same} function $xy$ is
uniformly continuous on $[0,1]^2$ and not on $\mathbb{R}^2$. Uniform continuity is therefore not
a property of a formula but of a formula \emph{together with its domain}, and compactness of the
domain is what removes the room to run off to a bad point.
\end{remark}
\end{exercisesolution}

\begin{exercisesolution}[Solution to \cref{ex:3.11}]
% Generator: Claude Opus 5 (High)
\begin{enumerate}[label=\textbf{\arabic*.}]
  \item \textbf{($\Leftarrow$) A modulus of continuity forces uniform continuity.} Suppose
    $d_Y(f(x),f(y)) \leq \omega(d_X(x,y))$ for all $x,y$, with $\omega$ as in the statement. Let
    $\varepsilon > 0$. Since $\omega(t) \to 0$ as $t \downarrow 0$, there is $\delta > 0$ with
    $\omega(\delta) < \varepsilon$. If $d_X(x,y) < \delta$ then, $\omega$ being nondecreasing,
    \[d_Y\big(f(x),f(y)\big) \leq \omega\big(d_X(x,y)\big) \leq \omega(\delta) < \varepsilon.\]
    The $\delta$ did not depend on $x$ or $y$, which is precisely
    \cref{def:uniformly_continuous}.

    \textbf{($\Rightarrow$) Uniform continuity produces one.} Following the official hint, define
    the \newterm{modulus of continuity of $f$} as its worst oscillation at scale $t$:
    \[\omega_f(t) := \sup\big\{\,d_Y(f(x),f(y))\ :\ x,y \in X,\ d_X(x,y) \leq t\,\big\}
      \ \in\ [0,\infty].\]
    It has all three required properties. $\omega_f(0) = 0$, because $d_X(x,y) \leq 0$ forces
    $x=y$ by definiteness and then $d_Y(f(x),f(y)) = 0$. It is nondecreasing, because enlarging
    $t$ enlarges the set over which the supremum is taken. And $\omega_f(t) \to 0$ as
    $t \downarrow 0$: given $\varepsilon > 0$, uniform continuity supplies $\delta > 0$ with
    $d_X(x,y) < \delta \Rightarrow d_Y(f(x),f(y)) < \varepsilon$, so every $t < \delta$ has
    $\omega_f(t) \leq \varepsilon$. Finally $f$ has modulus $\omega_f$ by construction, since
    $d_Y(f(x),f(y))$ is one of the numbers the supremum defining $\omega_f(d_X(x,y))$ ranges
    over.
  \item \textbf{(Bonus 1) On an interval, $\omega_f$ is already subadditive.} Let
    $t_1,t_2 \geq 0$ and take any $x,y \in X$ with $|x-y| \leq t_1+t_2$. Choose an intermediate
    point $z$ on the segment from $x$ to $y$ as follows: if $|x-y| \leq t_1$ put $z := y$;
    otherwise let $z$ be the point of that segment at distance exactly $t_1$ from $x$, so that
    $|z-y| = |x-y| - t_1 \leq t_2$. Either way $z \in X$, because $X$ is an interval and $z$ lies
    between $x$ and $y$, and $|x-z| \leq t_1$, $|z-y| \leq t_2$. The triangle inequality in $Y$
    now gives
    \[d_Y\big(f(x),f(y)\big) \ \leq\ d_Y\big(f(x),f(z)\big) + d_Y\big(f(z),f(y)\big)
      \ \leq\ \omega_f(t_1) + \omega_f(t_2).\]
    Taking the supremum over all admissible $x,y$ yields
    $\omega_f(t_1+t_2) \leq \omega_f(t_1) + \omega_f(t_2)$.
  \item \textbf{(Bonus 2) Convexity is exactly what the argument needed.} Repeat part 2 verbatim,
    now reading $|x-y|$ as the Euclidean norm on $\mathbb{R}^n$ rather than the absolute value on
    $\mathbb{R}$: the only property of an interval that was used is that the segment from $x$ to
    $y$ lies inside $X$, and that is the definition of convexity. The intermediate point is
    $z := x + \min\{t_1,|x-y|\}\,\frac{y-x}{|y-x|}$ for $x \neq y$, which lies on the segment and
    satisfies the same two estimates.
\end{enumerate}

\begin{remark}[Why the value $+\infty$ has to be allowed]
The official remark accompanying this problem gives the reason, and it is worth restating: on
$X = \mathbb{Z} \subseteq \mathbb{R}$ the map $f(n) := n^2$ is uniformly continuous for the
trivial reason that any two distinct points are at distance at least $1$. Take
$\delta := \tfrac12$, and the implication holds vacuously. Yet
$|f(n+1)-f(n)| = 2n+1$ is unbounded, so $\omega_f(1) = +\infty$. Parts 2 and 3 rule this out by
hypothesis rather than by luck: on a convex set the intermediate-point construction lets a large
displacement be broken into many small ones, so $\omega_f(t) \leq \lceil t/s\rceil\,\omega_f(s)$
is finite as soon as some $\omega_f(s)$ is. A space like $\mathbb{Z}$ has no intermediate points
to subdivide with.
\end{remark}
\end{exercisesolution}


% ============================================================
% Chapter 6 --- Completeness and the Banach fixed point theorem
% Originally: content/week-02.tex, \section{Completeness} (Corsin Week 2, Friday)
%           + content/week-03.tex, \subsection{Contraction mappings and the Banach fixed
%             point theorem} (Corsin Week 3, Friday)
%
% Placed after \cref{ch:continuity} because a contraction is by definition a Lipschitz map
% with constant < 1, so \cref{def:lipschitz_continuous} must already exist. Corsin's own
% order was the same; only the week boundary between the two halves is gone.
% ============================================================
\chapter{Completeness and the Banach Fixed Point Theorem}
\label{ch:completeness}

% Transition: Claude Opus 5 (High)
Having examined how maps behave between metric spaces, we return to a property of a single space.
Does every sequence that \qt{looks like} it ought to converge (\cref{def:convergent_sequence})
actually do so? A space where the answer is yes is called complete.

Completeness is the first genuinely \emph{metric} rather than topological property met in this
part, and the distinction is sharp: two metrics can generate exactly the same open sets while
only one of them is complete. What completeness is usually good for is producing objects one
cannot write down. The Banach fixed point theorem at the end of the chapter is the standard
instrument for that, and it reappears as the engine of both the inverse function theorem and the
Picard--Lindelöf existence theorem later on.

% Originally: content/week-02.tex, \section{Completeness} (Corsin Week 2, Friday)

\section{Completeness}
\label{sec:completeness}

\begin{definition}[Cauchy sequence and complete metric space]
\label{def:complete_metric_space}
Let $(X,d)$ be a metric space and $(x_n)_{n=0}^{\infty}$ a sequence in $X$. We say that
$(x_n)_{n=0}^{\infty}$ is \newterm{Cauchy} if
\[\forall \varepsilon > 0 \ \exists N \in \mathbb{N} \text{ such that } d(x_n, x_m) < \varepsilon \quad \forall n, m \geq N.\]
$(X,d)$ is a \newterm{complete metric space} (\germanterm{vollst\"andiger metrischer Raum}) if
every Cauchy sequence in $X$ converges with respect to $d$.
\end{definition}

% Supplement: Sascha Brack/Class Notes/Week_02_Notes_Monday_Updated.pdf, pp. 7--9
\begin{proposition}[Convergence in $\mathbb{R}^n$ is coordinatewise; $\mathbb{R}^n$ is complete]
\label{prop:Rn_coordinatewise_complete}
Equip $\mathbb{R}^n$ with the standard metric $d_2$, and write
$x_m = (x_{m,1},\dots,x_{m,n})$.
\begin{enumerate}[label=\textbf{(\alph*)}]
  \item $x_m \to x$ in $\mathbb{R}^n$ if and only if $x_{m,i} \to x_i$ in $\mathbb{R}$ for every
    $i = 1,\dots,n$.
  \item $(\mathbb{R}^n, d_2)$ is complete.
\end{enumerate}
\end{proposition}

\begin{proof}
Everything follows from the two-sided estimate, valid for every $j$:
\[|x_{m,j}-x_j| \ \leq\ \lVert x_m-x\rVert \ =\ \sqrt{\sum_{i=1}^n|x_{m,i}-x_i|^2}
  \ \leq\ \sqrt{n}\,\max_i|x_{m,i}-x_i|.\]

\textbf{(a)} If $\lVert x_m-x\rVert \to 0$, the left inequality forces each coordinate to
converge. Conversely, if every coordinate converges then $\max_i|x_{m,i}-x_i| \to 0$ (a maximum
of \emph{finitely} many null sequences), so the right inequality forces
$\lVert x_m - x\rVert \to 0$.

\textbf{(b)} Let $(x_m)$ be Cauchy in $\mathbb{R}^n$. The same left inequality gives
$|x_{m,j}-x_{k,j}| \leq \lVert x_m-x_k\rVert$, so each coordinate sequence $(x_{m,j})_m$ is
Cauchy in $\mathbb{R}$. Since $\mathbb{R}$ is complete --- this is the completeness axiom from
\emph{Analysis I}, and the only input the proof has --- each converges, say
$x_{m,j} \to x_j$. Setting $x := (x_1,\dots,x_n)$, part \textbf{(a)} gives $x_m \to x$.
\end{proof}

\begin{ainote}
Worth noting how little is proved here: completeness of $\mathbb{R}^n$ is completeness of
$\mathbb{R}$, applied $n$ times. That is also where the argument would fail in infinite
dimensions --- $\max_i$ over infinitely many coordinates need not tend to $0$ even when each
does, which is exactly the phenomenon behind the $\varepsilon$-net argument in
\Cref{lem:sequentially_compact_totally_bounded}.
\end{ainote}

% Source: Corsin Nick/Class Notes/Week 2.pdf, pp. 1--13
\begin{example}[Completeness examples]
\begin{itemize}
  \item The Euclidean spaces $(\mathbb{R}^n, \langle\cdot,\cdot\rangle_{\text{eucl}})$ are
    \textbf{complete}.
  \item Any \textbf{non-closed} subset $U \subseteq \mathbb{R}^n$ (standard metric) and the
    rationals $\mathbb{Q}$ are \textbf{not complete}.
  \item $\big(C^0([0,1]), \sup_{x \in [0,1]}|\cdot|\big)$ is \textbf{complete} (Zorich, p.\ 22).
\end{itemize}
\end{example}

% Generator: Claude Opus 5 (Medium)
\begin{remark}[Completeness is not the same as closedness]
\label{rem:complete_vs_closed}
The second bullet is easy to over-read as \qt{complete means closed}. It does not, and the
difference is worth pinning down now (because it recurs).

\textbf{Closedness is relative, completeness is intrinsic.} \qt{Closed} only means anything once
you say \emph{closed in what} (i.e., it is a statement about how a set sits inside an ambient space).
Completeness asks nothing about an ambient space --- only whether the Cauchy sequences of the
space converge \emph{within it}. So:
\begin{itemize}
  \item $\mathbb{Q}$ is \textbf{closed in $\mathbb{Q}$} --- every space is closed in itself
    (\cref{ex:X_empty_clopen}) --- and yet \textbf{not complete}: the sequence
    $1,\,1.4,\,1.41,\,1.414,\dots$ is Cauchy in $\mathbb{Q}$ and converges to nothing there.
  \item $(0,1)$ is \textbf{not closed in $\mathbb{R}$} and also not complete, which is the
    second bullet's case --- but the two failures have different causes, and the first does not
    imply the second in general.
\end{itemize}

What \emph{is} true is a one-directional statement, and only inside a complete ambient space:
if $X$ is complete and $A \subseteq X$, then
\[A \text{ closed in } X \iff A \text{ complete}.\]
This is why the second bullet is correct \emph{as stated for $\mathbb{R}^n$} --- it silently uses
completeness of the ambient $\mathbb{R}^n$. Drop that hypothesis and the equivalence dies, as
$\mathbb{Q}$ inside $\mathbb{Q}$ shows.

\begin{ainote}
The same relative-versus-intrinsic distinction is the one behind
\Cref{rem:common_misconception} ($[0,1]^2$ is open in itself, not in $\mathbb{R}^2$), and it
returns as \cref{rem:boundedness_not_topological}. Compactness, by contrast, is
intrinsic like completeness: a subset $K$ is compact exactly when the metric space
$(K, d|_{K\times K})$ is, with no reference to the ambient space at all --- which is precisely
the reduction used in \cref{def:sequential_topological_compactness}.
\end{ainote}
\end{remark}

% Quelle: Diego Torres Tejeda/Notes - 06.03 - Diego Torres Tejeda.pdf, p. 2
\begin{exercise}[Completeness is not a topological property]
\label{ex:completeness_not_topological}
Corsin does not address whether completeness survives a change of metric; Diego Torres Tejeda
sets this exercise, which settles it. Let $X := (0,1]$, let $d_{\text{eucl}}(x,y) := |x-y|$ be
the standard metric on $X$, and define
\[d(x,y) := \left|\frac1x - \frac1y\right|, \qquad x,y \in X.\]
\begin{enumerate}[label=\textbf{(\alph*)}]
  \item Show that $d$ is a metric on $X$.
  \item Show that $(X,d_{\text{eucl}})$ is \textbf{not} complete.
  \item Show that $(X,d)$ \textbf{is} complete.
  \item Show that $d_{\text{eucl}}$ and $d$ induce the \textbf{same topology} on $X$.
\end{enumerate}
\textit{Diego's hint for} \textbf{(c)}: let $(x_n)$ be $d$-Cauchy. Then $(1/x_n)$ is Cauchy for
the standard metric, so it has a limit $\ell \in \mathbb{R}$; show $\ell \geq 1$ and that
$x_n \to 1/\ell$.
\end{exercise}

\begin{ainote}
Diego states the exercise on $X = (0,\infty)$, but there part \textbf{(c)} is false: on
$(0,\infty)$ the map $x\mapsto1/x$ is a bijection \emph{onto} $(0,\infty)$, so $(X,d)$ is
isometric to $\big((0,\infty),|\cdot|\big)$, which is not complete --- concretely $x_n := n$
satisfies $d(x_n,x_m) = |\tfrac1n-\tfrac1m| \to 0$ yet has no limit in $X$. His own hint asks to
show $\ell > 0$, and that is precisely the step that fails. Restricting to a bounded interval
repairs it: on $(0,1]$ the image is $[1,\infty)$, which \emph{is} closed in $\mathbb{R}$ and
hence complete. Typeset with $X := (0,1]$ above.
\end{ainote}

% Sonnet 5 (Medium)
Compactness will turn out to be a strictly stronger finiteness-type property than completeness:
every compact metric space is complete (\cref{ex:compact_finite_complete}), but the converse
fails --- $\mathbb{R}^n$ itself is complete, yet (being unbounded) not compact, as we will see
next (\cref{thm:sequential_compactness}, \cref{sec:heine_borel}).

% Source: Corsin Nick/Class Notes/Week 2.pdf, pp. 1--13

% Supplement: Analysis_II_Script_v1.pdf, pp. 11--13
\section{The completion of a metric space}
\label{sec:completion_of_a_metric_space}

% This section has no counterpart in Corsin's notes. It is transcribed directly from the official
% script (S9.1.3, printed pp. 11--13), rather than enriching an existing tutor passage the way the
% rest of this document's supplements do. Included because script Exercises 9.33--9.35 walking
% through the construction are short, self-contained, and exactly the kind of material this
% chapter collects. Not included: the script's own use of the construction to build R from Q
% (S9.1.4), which it marks "extra; cf. Grundstrukturen" and which belongs to the sibling project.
% Transition: Claude Opus 5 (High)
\Cref{rem:complete_vs_closed} left one question open. We know $\mathbb{Q}$ is not complete and
$\mathbb{R}$ is, and that $\mathbb{R}$ is built to repair exactly that defect. Is the repair a
one-off trick special to the reals, or can any incomplete space be fixed the same way?

It can, and the construction below is the general answer. Every metric space embeds isometrically,
and densely, into a complete one, obtained by taking the space of its own Cauchy sequences and
declaring two of them equal when they ought to have the same limit. The reals from the rationals
is then just the first instance.

\begin{definition}[Completion of a metric space]
\label{def:completion_of_metric_space}
Let $(X,d)$ be a metric space. Write $\mathcal{C}_X$ for the set of all Cauchy sequences in $X$,
and define an equivalence relation on $\mathcal{C}_X$ by
\[(x_n)_{n=0}^\infty \sim (y_n)_{n=0}^\infty \quad :\Longleftrightarrow\quad \lim_{n\to\infty}d(x_n,y_n) = 0.\]
The quotient set $\overline X := \mathcal{C}_X/{\sim}$, equipped with the map
\[\overline d\big([(x_n)_{n=0}^\infty],\,[(y_n)_{n=0}^\infty]\big) := \lim_{n\to\infty}d(x_n,y_n),\]
is called the \newterm{completion} (\germanterm{Vervollst\"andigung}) of $(X,d)$. The injection
$\iota : X \to \overline X$ sending $x$ to the class of the constant sequence $(x,x,x,\dots)$ is
called the \newterm{canonical embedding}. For all $x,y \in X$,
\[d(x,y) = \overline d\big(\iota(x),\iota(y)\big),\]
which in particular shows $\iota$ is injective.
\end{definition}

% Supplement: Analysis_II_Script_v1.pdf, p. 11
\begin{exercise}[Well-definedness of the completion]
\label{ex:completion_well_defined}
Show that the objects introduced in \cref{def:completion_of_metric_space} are well defined. In
particular, verify that $\overline d$ is indeed a metric on $\overline X$.
\end{exercise}

% Supplement: Analysis_II_Script_v1.pdf, p. 11
\begin{exercise}[The completion is complete]
\label{ex:completion_is_complete}
As the name suggests, the completion $(\overline X,\overline d)$ of a metric space is complete,
meaning every Cauchy sequence in $\overline X$ converges. A sequence in $\overline X$ is
essentially a sequence of sequences, i.e.\
$\big[(x_{m,n})_{n=0}^\infty\big]_{m=0}^\infty$. Show that $\big[(x_{n,n})_{n=0}^\infty\big]$ is a
limit of this sequence.
\end{exercise}

% Supplement: Analysis_II_Script_v1.pdf, p. 11
\begin{exercise}[Universal property of the completion]
\label{ex:completion_universal_property}
Let $(X,d)$ be a metric space with completion $(\overline X,\overline d)$. Let $(Y,d_Y)$ be a
complete metric space, and let $f : X \to Y$ be a function such that
\[d(x,y) = d_Y\big(f(x),f(y)\big) \qquad \text{for all } x,y \in X.\]
Show that there exists a unique function $\overline f : \overline X \to Y$ such that
$f = \overline f \circ \iota_X$ and
\[\overline d(\overline x,\overline y) = d_Y\big(\overline f(\overline x),\overline f(\overline y)\big)
  \qquad \text{for all } \overline x,\overline y \in \overline X.\]
This can be interpreted as: \qt{$\overline X$ is the \emph{smallest} complete metric space
containing $X$.}
\end{exercise}

\begin{remark}
\Cref{ex:completion_universal_property} is what makes \qt{the} completion well-posed rather than
just \qt{a} completion: any other complete space that $X$ embeds into isometrically receives a
unique isometric copy of $\overline X$ as well, via $\overline f$. This is the same
uniqueness-up-to-isomorphism pattern behind \qt{the} real numbers as \qt{the} complete ordered
field, and it is exactly how the script uses this construction in $\S9.1.4$, taking
$X = \mathbb{Q}$.
\end{remark}

% Originally: content/week-03.tex, \subsection{Contraction mappings and the Banach fixed point theorem}

% Sonnet 5 (Medium)
\section{Contraction mappings and the Banach fixed point theorem}
\label{sec:banach_fixed_point}

% Not in Corsin's own notes either, but needed to justify the appeal to the "Banach Fixed Point
% Theorem" in the solution to ex:ai_contraction_cos below, and a standard companion to
% completeness (def:complete_metric_space), which is why it sits in this chapter.
% Transition: Claude Opus 5 (High)
Completeness guarantees that certain limits exist, but on its own it does not say what they are.
The following theorem is the standard way of turning that guarantee into a construction: it
identifies a single point, proves it exists, proves it is unique, and supplies an algorithm that
converges to it. Almost every existence proof in the second half of this course goes through it.

\begin{definition}[Contraction mapping]
\label{def:contraction_mapping}
Let $(X,d)$ be a metric space. A map $f : X \to X$ is a \newterm{contraction}
\germanfor{Kontraktion} if it is Lipschitz continuous (\cref{def:lipschitz_continuous}) with
some Lipschitz constant $L \in [0,1)$, i.e.,
\[d(f(x),f(y)) \leq L \cdot d(x,y) \quad \text{for all } x,y \in X.\]
\end{definition}

\begin{theorem}[Banach fixed point theorem]
\label{thm:banach_fixed_point}
Let $(X,d)$ be a nonempty complete metric space (\cref{def:complete_metric_space}) and let
$f : X \to X$ be a contraction (\cref{def:contraction_mapping}). Then $f$ has a unique
\newterm{fixed point} \germanfor{Fixpunkt} $x^* \in X$, i.e., $f(x^*) = x^*$; moreover, for any
$x_0 \in X$, the iterates $x_{n+1} := f(x_n)$ converge to $x^*$.
\end{theorem}

% Supplement: Sascha Brack/Class Notes/Week_02_Notes_Monday_Updated.pdf, p. 13
\begin{proof}[Proof idea]
The core idea is to create a Cauchy sequence by repeatedly applying the contraction mapping $f$.
Fix an arbitrary starting point $x_0 \in X$, and define the sequence by $x_{n+1} := f(x_n)$ for all $n \geq 0$.
Comparing consecutive terms gives:
\[
d(x_{n+1}, x_n) = d(f(x_n), f(x_{n-1})) \leq \lambda \cdot d(x_n, x_{n-1}).
\]
Iterating this inequality yields $d(x_{n+1}, x_n) \leq \lambda^n d(x_1, x_0)$. Since $\lambda < 1$, one can use the triangle inequality and the geometric series to show that the distance between any two terms $x_m$ and $x_n$ (for $m > n$) becomes arbitrarily small as $n \to \infty$. Thus, the sequence is Cauchy.
Because $X$ is a complete metric space, this sequence converges to some limit $x^* \in X$. Since contraction mappings are continuous, one can pass the limit inside the function to conclude $f(x^*) = x^*$.
\end{proof}

% Generator: Claude Opus 5 (High) --- Sascha's proof idea establishes existence only; uniqueness
% is two lines and is half of what the theorem claims.
\begin{proof}[Uniqueness of the fixed point]
Suppose $f(x) = x$ and $f(y) = y$. Then
\[d(x,y) = d\big(f(x),f(y)\big) \leq L\,d(x,y),\]
so $(1-L)\,d(x,y) \leq 0$. Since $L<1$ we have $1-L>0$, forcing $d(x,y) \leq 0$ and hence
$d(x,y)=0$, i.e.\ $x=y$.
\end{proof}

% Generator: Claude Opus 5 (High) --- FIG-06-01, cobweb diagrams for the fixed point iteration.
% This chapter had no figure at all, and the theorem's proof is a construction (iterate f from
% an arbitrary starting point) whose whole content is visual in one variable.
% Arithmetic check. LEFT: f(x) = 0.45x + 0.35, a contraction with L = 0.45 < 1. Fixed point
% x* = 0.35/(1-0.45) = 0.35/0.55 = 0.6364. From x0 = 0.05 the iterates are
%   0.05 -> 0.3725 -> 0.51763 -> 0.58293 -> 0.61232 -> 0.62554, climbing to x*.
% RIGHT: f(x) = x + 0.12, slope exactly 1, so L = 1 and NOT a contraction. The graph is
% parallel to the diagonal and never meets it, so there is no fixed point at all. From
% x0 = 0.10 the iterates are 0.22, 0.34, 0.46, 0.58, 0.70, 0.82, marching off to the right.
% This is the second bullet of rem:banach_hypotheses (the T(x) = x+1 example) drawn on [0,1].
\begin{center}
\begin{tikzpicture}[x={(3.4cm,0)}, y={(0,3.4cm)}, >=Latex,
    web/.style={BrickRed, thick},
    fgraph/.style={blue!65!black, very thick}]

  % ---------------- left: a genuine contraction ----------------
  \begin{scope}
    \draw[->, gray!70] (0,0) -- (1.1,0) node[right, font=\footnotesize, text=black] {$x$};
    \draw[->, gray!70] (0,0) -- (0,1.1);
    \draw[gray!55, dashed] (0,0) -- (1,1) node[above left, font=\scriptsize, text=black] {$y=x$};
    \draw[fgraph] (0,0.35) -- (1,0.80) node[right, font=\scriptsize] {$f$};

    \draw[web] (0.05,0.05) -- (0.05,0.3725) -- (0.3725,0.3725) -- (0.3725,0.51763)
      -- (0.51763,0.51763) -- (0.51763,0.58293) -- (0.58293,0.58293) -- (0.58293,0.61232)
      -- (0.61232,0.61232) -- (0.61232,0.62554);

    \fill[black] (0.6364,0.6364) circle (1.7pt);
    \node[font=\scriptsize, anchor=west] at (0.68,0.60) {$x^{*}$};
    \draw[gray!70] (0.05,0.02) -- (0.05,-0.02) node[below, font=\scriptsize, text=black] {$x_0$};
    \node[font=\scriptsize, align=center, anchor=north] at (0.5,-0.15)
      {contraction, $L=0.45$\\ iterates climb to the unique $x^{*}$};
  \end{scope}

  % ---------------- right: slope exactly 1, no fixed point ----------------
  \begin{scope}[shift={(1.42,0)}]
    \draw[->, gray!70] (0,0) -- (1.1,0) node[right, font=\footnotesize, text=black] {$x$};
    \draw[->, gray!70] (0,0) -- (0,1.1);
    \draw[gray!55, dashed] (0,0) -- (1,1) node[below right, font=\scriptsize, text=black] {$y=x$};
    \draw[fgraph] (0,0.12) -- (0.88,1.0) node[above left, font=\scriptsize] {$f$};

    \draw[web] (0.10,0.10) -- (0.10,0.22) -- (0.22,0.22) -- (0.22,0.34) -- (0.34,0.34)
      -- (0.34,0.46) -- (0.46,0.46) -- (0.46,0.58) -- (0.58,0.58) -- (0.58,0.70)
      -- (0.70,0.70) -- (0.70,0.82) -- (0.82,0.82) -- (0.82,0.94);
    \draw[web, ->] (0.82,0.94) -- (0.94,0.94);

    \draw[gray!70] (0.10,0.02) -- (0.10,-0.02) node[below, font=\scriptsize, text=black] {$x_0$};
    \node[font=\scriptsize, align=center, anchor=north] at (0.5,-0.15)
      {$L=1$: graph parallel to $y=x$\\ no fixed point, iterates escape};
  \end{scope}
\end{tikzpicture}
\end{center}

The left panel is the proof in miniature. Each step moves vertically to the graph of $f$ and then
horizontally back to the diagonal, and because the graph is shallower than the diagonal, the steps
shrink by a factor of at least $L$ every time. That is precisely the geometric-series estimate the
proof runs on. The right panel shows what a single missing hypothesis costs: with slope exactly
$1$ the graph never meets the diagonal, so no fixed point exists and the same iteration marches
off instead of converging.

\begin{remark}[Where each hypothesis is used]
\label{rem:banach_hypotheses}
Both hypotheses of \cref{thm:banach_fixed_point} are indispensable, and each fails in a
different way.
\begin{itemize}
  \item \textbf{Completeness} is what supplies the limit. On the incomplete space
    $X := (0,1]$, the map $f(x) := x/2$ is a contraction with $L=\tfrac12$ and has no fixed
    point in $X$, since the iterates march towards $0$, which is exactly the point missing.
  \item \textbf{$L<1$ strictly} is what forces the iterates together. The map
    $f(x) := x + \tfrac{1}{1+e^{x}}$ on the complete space $\mathbb{R}$ satisfies the strict
    inequality $|f(x)-f(y)| < |x-y|$ for $x \neq y$, yet has no fixed point, since
    $f(x)>x$ everywhere. A supremum of $1$ that is never attained is still not good enough.
\end{itemize}
The uniqueness proof above shows the second point cleanly: the whole argument is the division by
$1-L$, which is exactly what $L=1$ forbids.
\end{remark}

% Supplement: Analysis_II_Script_v1.pdf, p. 20
\begin{exercise}[Both hypotheses of Banach's theorem are sharp]
\label{ex:banach_hypotheses_sharp}
Find examples for:
\begin{enumerate}[label=\textbf{(\alph*)}]
  \item A Lipschitz contraction $T : X \to X$ on a non-complete metric space $X$ that has no
    fixed point.
  \item A complete metric space $(X,d)$ and an \newterm{isometry} (a mapping $T : X \to X$ with
    $d(T(x_1),T(x_2)) = d(x_1,x_2)$ for all $x_1,x_2$) that has no fixed point.
\end{enumerate}
\end{exercise}

\begin{remark}
Part \textbf{(b)} sharpens \cref{rem:banach_hypotheses}'s second bullet: that example already
showed a Lipschitz constant of exactly $1$ can defeat the theorem, but its map $T$ was not an
isometry ($|T(x)-T(y)| < |x-y|$ strictly, just with no uniform gap). An isometry is the most
rigid failure imaginable, not shrinking distances even a little, and it \emph{still} need not have
a fixed point, even on a complete space. Note that unboundedness is doing real work in the
counterexample: on a \emph{compact} space an isometry into itself is automatically surjective, and
even that is not enough to force a fixed point. An irrational rotation of the circle $S^1$, which
is compact and an isometry, has none either. Fixed points of isometries are a
genuinely separate question from \cref{thm:banach_fixed_point}, not one it settles as a special
case.
\end{remark}

\begin{aiexercise}[{Contraction Mapping on $[0,1]$}]
\label{ex:ai_contraction_cos}
% Generator: Gemini 3.6 Flash (Medium)
Let $f \colon [0,1] \to [0,1]$ be defined by $f(x) := \frac{1}{2}\cos(x)$.
\begin{enumerate}[label=\textbf{(\alph*)}]
  \item Show that $f$ is a contraction mapping (\cref{def:contraction_mapping}) with Lipschitz
    constant $L = \frac{1}{2}\sin(1) < 1$.
  \item Deduce, using \cref{thm:banach_fixed_point}, that $f$ has a unique fixed point
    $x^* \in [0,1]$.
\end{enumerate}
\end{aiexercise}

% Supplement: Linus Lüchinger/Slides/Slides-03-02.pdf, slide 16;
%             Linus Lüchinger/Slides/Slides-03-05.pdf, slide 3
\begin{exercise}[Does the fixed-point theorem apply?]
\label{ex:bfpt_applicability_table}
\cref{rem:banach_hypotheses} shows that each hypothesis is needed by exhibiting one failure
apiece. Linus drills the same point the other way round, as a checklist: for each row below,
decide whether \cref{thm:banach_fixed_point} \emph{applies} to $T : X \to X$, and separately
determine the set of fixed points of $T$ in $X$. The two questions have different answers more
often than one expects.
\begin{center}
\begin{tabular}{lll}
\textbf{$T(x)$} & \textbf{$X$} & \textbf{Applies? / fixed points} \\
\hline
$\sin(x)$ & $\mathbb{R}$ & \\
$\tfrac{9}{10}\cos(x)$ & $\mathbb{R}$ & \\
$x^2$ & $\mathbb{R}$ & \\
$\tfrac12x^2$ & $(-1,1)$ & \\
$\tfrac12\big(x + \tfrac2x\big)$ & $\mathbb{Q}\cap(1,2)$ &
\end{tabular}
\end{center}
\end{exercise}


% --- exercise relocated here from the dissolved sheet 3 ---

% Originally: content/exercise-sheets/week-03.tex, problem 3.5
% Source: exercises/Ex3_Analysis2_eng.pdf

\begin{exercise}[Multiple choice (fixed points)]
\label{ex:3.5}
Select all the statements below that are true.
\begin{enumerate}[label=\textbf{(\alph*)}]
  \item If you spread a map of Zurich on your desk, then one point of the desk will coincide with
    its representation on the map (in an ideal world).
  \item If $f : [0,1] \to [0,2]$ is $\tfrac{1}{2}$-Lipschitz, then $f$ has a fixed point.
  \item If $f : [0,2] \to [0,1]$ is $\tfrac{1}{2}$-Lipschitz, then $f$ has a fixed point.
  \item If $f : [0,1] \cup [2,3] \to [0,1] \cup [2,3]$ is continuously differentiable with
    $|f'(x)| < 1$ for all $x \in [0,1]\cup[2,3]$, then it has a fixed point.
  \item \textbf{(*)} If $f : [0,1] \to [0,1]$ is differentiable with $|f'(x)| < 1$ for all
    $x \in (0,1)$, then it has a fixed point.
\end{enumerate}
\exinfo{This exercise is Problem 3.5 of Problem Sheet 3, Analysis II, Spring Semester 2026.}
\end{exercise}

% Originally: the \section{Solutions} of content/week-02.tex and content/week-03.tex

\section{Solutions}
\label{sec:completeness_solutions}

\begin{exercisesolution}[Proof of \cref{ex:cauchy_elementary_facts}]
% Generator: Claude Sonnet 5 (High)
\textbf{(a) Bounded.} Take $\varepsilon=1$: there is $N$ with $d(x_n,x_m)<1$ for all
$n,m\geq N$; in particular $d(x_n,x_N)<1$ for $n\geq N$. Setting
$M := \max\{d(x_0,x_N),\dots,d(x_{N-1},x_N),\,1\}$, the triangle inequality gives
$d(x_n,x_0) \leq d(x_n,x_N)+d(x_N,x_0) \leq 1+M$ for every $n$, so $(x_n)_n \subseteq B_{1+M}(x_0)$.

\textbf{(b) Convergent $\Rightarrow$ Cauchy.} If $x_n \to x$, fix $\varepsilon>0$ and choose $N$
with $d(x_n,x)<\varepsilon/2$ for $n\geq N$. Then for $n,m\geq N$,
$d(x_n,x_m) \leq d(x_n,x)+d(x,x_m) < \varepsilon$.

\textbf{(c) Cauchy + convergent subsequence $\Rightarrow$ convergent.} Let $(x_n)$ be Cauchy and
$x_{n_k} \to x$. The \qt{only if} direction is part \textbf{(b)} applied to the whole sequence.
For \qt{if}: fix $\varepsilon>0$. Cauchyness gives $N$ with $d(x_n,x_m)<\varepsilon/2$ for
$n,m\geq N$; convergence of the subsequence gives $K$ with $d(x_{n_k},x)<\varepsilon/2$ for
$k\geq K$. Since $n_k \to \infty$, pick $k\geq K$ with $n_k \geq N$. Then for every $n\geq N$,
\[d(x_n,x) \leq d(x_n,x_{n_k}) + d(x_{n_k},x) < \tfrac{\varepsilon}{2}+\tfrac{\varepsilon}{2} = \varepsilon,\]
so $x_n \to x$.
\end{exercisesolution}

\begin{exercisesolution}[Solution to \cref{ex:completeness_not_topological}]
% Generator: Claude Opus 5 (Medium)
Write $\varphi : (0,1] \to [1,\infty)$, $\varphi(x) := 1/x$. It is a bijection, and by
construction $d(x,y) = |\varphi(x)-\varphi(y)|$. Consequently, $d$ is nothing but the Euclidean
metric of $[1,\infty)$ transported back to $(0,1]$ along $\varphi$. Every part follows from that.

\begin{enumerate}[label=\textbf{(\alph*)}]
  \item Symmetry and the triangle inequality are inherited from $|\cdot|$ directly. For
    definiteness, $d(x,y)=0$ forces $\varphi(x)=\varphi(y)$, and $\varphi$ is injective, so
    $x=y$.
  \item The sequence $x_n := \tfrac1n$ lies in $(0,1]$ and is Cauchy for
    $d_{\text{eucl}}$, being Cauchy in $\mathbb{R}$. Its limit in $\mathbb{R}$ is $0 \notin X$,
    and limits in $\mathbb{R}$ are unique (\cref{prop:uniqueness_of_limits}), so it has no limit
    in $X$.
  \item Let $(x_n)$ be $d$-Cauchy. Then $\big(\varphi(x_n)\big)$ is Cauchy in $\mathbb{R}$,
    hence converges to some $\ell \in \mathbb{R}$. Since $\varphi(x_n) \geq 1$ for every $n$,
    also $\ell \geq 1$; in particular $\ell \neq 0$, so $x := 1/\ell$ lies in $(0,1]$. Then
    \[d(x_n,x) = \big|\varphi(x_n) - \varphi(x)\big| = \big|\varphi(x_n)-\ell\big|
      \longrightarrow 0,\]
    so $x_n \to x$ in $(X,d)$. Note where the boundedness of $X$ entered: it is what forces
    $\ell \geq 1$ rather than merely $\ell \geq 0$.
  \item $\varphi$ and $\varphi^{-1}(t) = 1/t$ are continuous, so $\varphi$ is a
    homeomorphism. A set is $d$-open exactly when its $\varphi$-image is open in $[1,\infty)$,
    exactly when it is $d_{\text{eucl}}$-open. Concretely, the same two-sided-ball argument as in
    \Cref{ex:heine_borel_fails}\textbf{(b)} works: on $(0,1]$ one has
    $|x-y| \leq d(x,y) \leq \tfrac{1}{x y}\,|x-y|$ with $\tfrac{1}{xy}$ locally bounded.
\end{enumerate}

\begin{remark}
So $d_{\text{eucl}}$ and $d$ have the \emph{same open sets}, the same convergent sequences, the
same continuous functions and the same compact sets, and yet one is complete and the other is
not. \textbf{Completeness is a property of the metric, not of the topology it generates.}

This completes a trio worth holding together, all three proved by the same device of comparing
two metrics that induce one topology:
\begin{itemize}
  \item \textbf{Compactness is topological.} It is defined by open covers alone
    (\cref{def:compact_metric_space}), so it cannot distinguish two metrics with the same open
    sets.
  \item \textbf{Boundedness is not:} see \cref{rem:boundedness_not_topological}.
  \item \textbf{Completeness is not:} this exercise.
\end{itemize}
That is exactly why \qt{compact if and only if closed and bounded} needs the standard metric, and
why the correct general replacement is \qt{complete and totally bounded}
(\cref{rem:compact_iff_complete_totally_bounded}). Both ingredients of the replacement are
non-topological, and only their combination is.
\end{remark}
\end{exercisesolution}

\begin{exercisesolution}[Proof of \cref{ex:completion_well_defined}]
% Generator: Claude Sonnet 5 (High)
\textbf{$\sim$ is an equivalence relation on $\mathcal{C}_X$.} Reflexivity is
$d(x_n,x_n) = 0 \to 0$. Symmetry is symmetry of $d$. For transitivity, if
$(x_n)\sim(y_n)$ and $(y_n)\sim(z_n)$, the triangle inequality gives
$0 \leq d(x_n,z_n) \leq d(x_n,y_n)+d(y_n,z_n) \to 0+0 = 0$, so $(x_n)\sim(z_n)$.

\textbf{$\overline d$ is well-defined.} First, the limit defining $\overline d$ exists: for
$(x_n),(y_n) \in \mathcal{C}_X$, the reverse triangle inequality gives
\[\big|d(x_n,y_n) - d(x_m,y_m)\big| \leq d(x_n,x_m) + d(y_n,y_m),\]
and the right-hand side $\to 0$ as $n,m\to\infty$ since $(x_n)$ and $(y_n)$ are each Cauchy, so
$\big(d(x_n,y_n)\big)_n$ is a Cauchy sequence of \emph{real numbers}, hence convergent. Note that
this is the completeness axiom for $\mathbb{R}$ itself from \emph{Analysis I}, and not an instance
of \cref{def:complete_metric_space} being proved circularly. Second, the limit
does not depend on the choice of representative: if $(x_n)\sim(x_n')$ and $(y_n)\sim(y_n')$, the
same reverse triangle inequality gives
$\big|d(x_n,y_n)-d(x_n',y_n')\big| \leq d(x_n,x_n')+d(y_n,y_n') \to 0$, so the two limits agree.

\textbf{$\overline d$ is a metric.} Fix representatives $(x_n)\in\overline x$, $(y_n)\in\overline
y$, $(z_n)\in\overline z$.
\begin{itemize}
  \item \textbf{Definiteness.} $\overline d(\overline x,\overline y) = 0$ means
    $\lim_n d(x_n,y_n) = 0$, i.e.\ exactly $(x_n)\sim(y_n)$, i.e.\ $\overline x = \overline y$.
  \item \textbf{Symmetry.} Immediate from $d(x_n,y_n)=d(y_n,x_n)$ for every $n$.
  \item \textbf{Triangle inequality.} Take $n\to\infty$ in
    $d(x_n,z_n) \leq d(x_n,y_n) + d(y_n,z_n)$, using that limits preserve weak inequalities.
\end{itemize}
Finally $d(x,y) = \overline d(\iota(x),\iota(y))$ holds because the constant sequences
$(x,x,\dots)$ and $(y,y,\dots)$ give $d(x_n,y_n) = d(x,y)$ for every $n$, so the limit is
$d(x,y)$ itself; and this identity makes $\iota$ injective, since $d(x,y)=0 \iff x=y$.
\end{exercisesolution}

\begin{exercisesolution}[Proof of \cref{ex:completion_is_complete}]
% Generator: Claude Sonnet 5 (High)
Let $(\xi_m)_{m=0}^\infty$ be Cauchy in $\overline X$, with $\xi_m = [(x_{m,n})_{n=0}^\infty]$ for
some Cauchy sequence $(x_{m,n})_n$ in $X$ representing $\xi_m$. Replacing $(x_{m,n})_n$ by a tail
of itself if necessary, which is harmless because a Cauchy sequence and any of its tails represent
the same class, differing in only finitely many terms, we may assume each representative is already
\newterm{regularized}:
\[d(x_{m,n},x_{m,n'}) < \tfrac1m \qquad \text{for all } n,n' \geq 0.\]
Set $y_m := x_{m,m}$, the diagonal sequence the exercise asks about. Letting $n'\to\infty$ in the
regularization bound (with $n=m$) gives
\begin{equation}
\overline d\big(\iota(y_m),\,\xi_m\big) = \lim_{n\to\infty} d(x_{m,m},x_{m,n}) \ \leq\ \tfrac1m.
\label{eq:completion_diagonal_close}
\end{equation}

\textbf{$(y_m)_m$ is Cauchy in $X$.} By the triangle inequality for $\overline d$ (already
established in \cref{ex:completion_well_defined}) and \eqref{eq:completion_diagonal_close},
\[d(y_m,y_{m'}) = \overline d\big(\iota(y_m),\iota(y_{m'})\big)
  \ \leq\ \overline d(\iota(y_m),\xi_m) + \overline d(\xi_m,\xi_{m'}) + \overline d(\xi_{m'},\iota(y_{m'}))
  \ \leq\ \tfrac1m + \overline d(\xi_m,\xi_{m'}) + \tfrac1{m'}.\]
Given $\varepsilon>0$, Cauchyness of $(\xi_m)_m$ gives $M_1$ with
$\overline d(\xi_m,\xi_{m'}) < \varepsilon/3$ for $m,m'\geq M_1$; taking also
$M \geq \max\{M_1,\,3/\varepsilon\}$ makes all three terms $< \varepsilon/3$ for $m,m'\geq M$, so
$(y_m)_m$ is Cauchy in $X$ and defines a class $\overline x := [(y_m)_m] \in \overline X$.

\textbf{$\xi_m \to \overline x$.} Using \eqref{eq:completion_diagonal_close} again and the same
triangle inequality,
\[\overline d(\xi_m,\overline x) \ \leq\ \overline d(\xi_m,\iota(y_m)) + \overline d(\iota(y_m),\overline x)
  \ \leq\ \tfrac1m + \lim_{k\to\infty}d(y_m,y_k).\]
The first term $\to 0$ as $m\to\infty$ by construction, and the second does too, being (for each
fixed $m$) bounded by $\sup_{k\geq m}d(y_m,y_k)$, which $\to0$ as $m\to\infty$ since $(y_m)_m$ is
Cauchy. Hence $\overline d(\xi_m,\overline x) \to 0$, i.e.\ $\xi_m \to \overline x$ in
$\overline X$, as required.
\end{exercisesolution}

\begin{exercisesolution}[Proof of \cref{ex:completion_universal_property}]
% Generator: Claude Sonnet 5 (High)
\textbf{Construction.} For $\overline x = [(x_n)_n] \in \overline X$, the sequence $(f(x_n))_n$
is Cauchy in $Y$, since $d_Y(f(x_n),f(x_m)) = d(x_n,x_m) \to 0$ as $(x_n)_n$ is Cauchy in $X$.
As $Y$ is complete, $f(x_n) \to y$ for a unique $y \in Y$
(\cref{prop:uniqueness_of_limits}); define $\overline f(\overline x) := y$.

\textbf{Well-defined.} If $(x_n)\sim(x_n')$, then $d_Y(f(x_n),f(x_n')) = d(x_n,x_n') \to 0$, so
$(f(x_n))_n$ and $(f(x_n'))_n$ have the same limit in $Y$ by uniqueness of limits; hence
$\overline f(\overline x)$ does not depend on the representative chosen.

\textbf{$f = \overline f \circ \iota_X$.} For $x\in X$, $\iota_X(x)$ is represented by the
constant sequence $(x,x,\dots)$, so $\overline f(\iota_X(x)) = \lim_n f(x) = f(x)$.

\textbf{$\overline f$ satisfies the displayed identity.} For $\overline x = [(x_n)_n]$,
$\overline y = [(y_n)_n]$, continuity of $d_Y$ in both arguments (itself $1$-Lipschitz, by the
same reverse-triangle-inequality argument as in \cref{ex:completion_well_defined}) gives
\[d_Y\big(\overline f(\overline x),\overline f(\overline y)\big)
  = d_Y\Big(\lim_n f(x_n),\ \lim_n f(y_n)\Big)
  = \lim_n d_Y\big(f(x_n),f(y_n)\big)
  = \lim_n d(x_n,y_n) = \overline d(\overline x,\overline y).\]

\textbf{Uniqueness.} Let $g : \overline X \to Y$ satisfy the same two conditions. The displayed
identity makes $g$ (like $\overline f$) $1$-Lipschitz, hence continuous. For
$\overline x = [(x_n)_n]$, the sequence $\iota_X(x_n)$ converges to $\overline x$ in
$\overline X$: indeed $\overline d(\iota_X(x_n),\overline x) = \lim_k d(x_n,x_k) \to 0$ as
$n\to\infty$, since $(x_n)_n$ is Cauchy. By continuity of $g$ and $f = g\circ\iota_X$,
\[g(\overline x) = \lim_n g(\iota_X(x_n)) = \lim_n f(x_n) = \overline f(\overline x),\]
using uniqueness of limits in $Y$ once more for the last equality. As $\overline x$ was
arbitrary, $g = \overline f$.
\end{exercisesolution}

\begin{exercisesolution}[Proof of \cref{ex:banach_hypotheses_sharp}]
% Generator: Claude Sonnet 5 (High)
\begin{enumerate}[label=\textbf{(\alph*)}]
  \item This is \cref{rem:banach_hypotheses}'s first bullet: $X := (0,1]$ is not complete,
    and $T(x) := x/2$ is a $\tfrac12$-contraction on $X$ with iterates marching towards the
    missing point $0$, hence no fixed point in $X$.
  \item Take $X := \mathbb{R}$ (complete, \cref{prop:Rn_coordinatewise_complete}) with the
    translation $T(x) := x+1$. Then $d(T(x),T(y)) = |(x+1)-(y+1)| = |x-y| = d(x,y)$, so $T$ is an
    isometry, and in particular Lipschitz with constant $L=1$ rather than $L<1$. Consequently,
    this does not contradict \cref{thm:banach_fixed_point}. A fixed point would need $x+1=x$,
    which is impossible, so $T$ has none.
\end{enumerate}
\end{exercisesolution}

\begin{exercisesolution}[Proof of \cref{ex:kobel_local_banach_fixed_point}]
% Generator: Claude Opus 5 (High)
Write $K := \overline{B}_r(x_0)$. The whole content of the exercise is that $f$ restricts to a
self-map of $K$, after which \cref{thm:banach_fixed_point} finishes the job.

\textbf{$K$ is complete.} A closed ball is a closed subset of $X$, and a closed subset of a
complete metric space is complete: a Cauchy sequence in $K$ is Cauchy in $X$, hence converges in
$X$, and its limit lies in $K$ because $K$ is closed. It is also non-empty, since $x_0 \in K$.

\textbf{$f(K) \subseteq K$.} Let $x \in K$, so $d(x,x_0) \le r$. Since $K \subseteq E$, both $x$
and $x_0$ lie in the domain of $f$, and the triangle inequality together with the contraction
estimate gives
\[d\big(f(x),x_0\big) \ \le\ d\big(f(x),f(x_0)\big) + d\big(f(x_0),x_0\big)
  \ \le\ L\,d(x,x_0) + (1-L)r \ \le\ Lr + (1-L)r \ =\ r,\]
so $f(x) \in K$. The hypothesis $d(f(x_0),x_0) \le (1-L)r$ is used exactly once, here, and it is
what makes the two terms add up to $r$ rather than to something larger.

\textbf{Conclusion.} The restriction $f|_K : K \to K$ is a contraction with the same constant $L$
on the non-empty complete metric space $K$, so \cref{thm:banach_fixed_point} supplies a unique
$\bar x \in K$ with $f(\bar x) = \bar x$. Since $K \subseteq E$, this $\bar x$ is a fixed point in
$E$, as required.

\textbf{Uniqueness comes for free, and in all of $E$.} \Cref{thm:banach_fixed_point} gives
uniqueness only within $K$, but the two-line argument used for the theorem itself needs nothing
beyond the contraction estimate, which holds on all of $E$: if $f(x)=x$ and $f(y)=y$ with
$x,y \in E$, then $d(x,y) = d(f(x),f(y)) \le L\,d(x,y)$, and $L<1$ forces $d(x,y)=0$. So $\bar x$
is the only fixed point of $f$ in $E$, even though the existence argument saw only the ball.
\end{exercisesolution}

\begin{exercisesolution}[Proof of \cref{ex:ai_contraction_cos}]
% Generator: Gemini 3.6 Flash (Medium)
\begin{enumerate}[label=\textbf{(\alph*)}]
  \item By the Mean Value Theorem, for any $x, y \in [0,1]$, there exists $\xi \in (0,1)$ such that $|f(x) - f(y)| = |f'(\xi)| |x-y| = \frac{1}{2}\sin(\xi) |x-y|$. Since $\sin(x)$ is strictly increasing on $[0,1]$, $\sup_{\xi \in [0,1]} \frac{1}{2}\sin(\xi) = \frac{1}{2}\sin(1) \approx 0.4207 < 1$. Thus $f$ is a contraction (\cref{def:contraction_mapping}) with $L = \frac{1}{2}\sin(1) < 1$.
  \item Since $[0,1]$ equipped with the Euclidean metric is a complete metric space (\cref{def:complete_metric_space}) and $f([0,1]) \subseteq [0,1]$, \cref{thm:banach_fixed_point} guarantees the existence and uniqueness of a fixed point $x^* \in [0,1]$ satisfying $f(x^*) = x^*$.
\end{enumerate}
\end{exercisesolution}

\begin{exercisesolution}[Solution to \cref{ex:3.5}]
% Generator: Claude Opus 5 (High)
\textbf{(a), (c) and (e) are true; (b) and (d) are false.}

\begin{enumerate}[label=\textbf{(\alph*)}]
  \item \textbf{True,} and this is the classic illustration of \cref{thm:banach_fixed_point}. The
    map of Zurich lying on the desk defines a function from the (compact, hence complete) region
    of the city depicted onto a smaller region of the desk, and that function is a similarity: a
    rotation and translation composed with a scaling by the map's ratio $L < 1$. It is therefore
    a contraction of a complete metric space into itself, and so has exactly one fixed point: one
    spot on the map lying directly above the point of the city it represents.
  \item \textbf{False.} The hypothesis says nothing about $f$ mapping its domain into itself, and
    the codomain $[0,2]$ is strictly larger than the domain $[0,1]$. Take
    \[f(x) := \tfrac{x}{2} + \tfrac32, \qquad f : [0,1] \to [\tfrac32,2] \subseteq [0,2],\]
    which is $\tfrac12$-Lipschitz. A fixed point would need $x = \tfrac x2 + \tfrac32$, i.e.\
    $x=3$, which is not in $[0,1]$.
  \item \textbf{True.} Now the inclusion runs the right way: $f$ maps $[0,2]$ into
    $[0,1] \subseteq [0,2]$, so $f$ is a self-map of the complete space $[0,2]$
    (closed and bounded in $\mathbb{R}$, hence compact, hence complete by
    \cref{ex:compact_finite_complete}), and it is a contraction with $L = \tfrac12 < 1$.
    \Cref{thm:banach_fixed_point} applies directly.
  \item \textbf{False,} and this is the instructive one: the domain is \emph{disconnected}, and
    the condition $|f'| < 1$ is purely local, so it cannot see across the gap. Define
    \[f(x) := \begin{cases} \tfrac52 & x \in [0,1], \\ \tfrac12 & x \in [2,3].\end{cases}\]
    Then $f$ is continuously differentiable with $f' \equiv 0$, it maps
    $[0,1]\cup[2,3]$ into itself, and it has no fixed point, because each of the two pieces is
    sent into the \emph{other} one. Compare the local-to-global reading of
    \cref{rem:local_to_global}: a pointwise derivative bound upgrades to a global contraction
    only along paths that stay inside the domain, and here there are none.
  \item \textbf{True,} but \emph{not} by the fixed point theorem, and that is the point of the
    star. The hypothesis $|f'(x)| < 1$ gives no uniform $L < 1$, since the supremum of $|f'|$ may
    well equal $1$, so \cref{thm:banach_fixed_point} is unavailable. Existence follows from the
    intermediate value theorem instead. Put $g(x) := f(x) - x$, which is continuous on $[0,1]$.
    Since $f$ takes values in $[0,1]$,
    \[g(0) = f(0) \geq 0 \qquad\text{and}\qquad g(1) = f(1) - 1 \leq 0,\]
    so $g$ vanishes somewhere in $[0,1]$ by \cref{cor:intermediate_value_theorem}, and that zero
    is a fixed point of $f$. Note that the derivative hypothesis was never used: \emph{every}
    continuous self-map of $[0,1]$ has a fixed point. What $|f'| < 1$ buys is
    \textbf{uniqueness}, by the same two-line argument as in the proof of uniqueness for
    \cref{thm:banach_fixed_point} combined with the mean value theorem.
\end{enumerate}
\end{exercisesolution}

\begin{exercisesolution}[Solution to \cref{ex:bfpt_applicability_table}]
% Generator: Claude Opus 5 (High)
Recall that \cref{thm:banach_fixed_point} needs three things at once: $X$ complete, $T$ mapping
$X$ into $X$, and a Lipschitz constant $\lambda < 1$ that is \emph{actually attained as a bound},
not merely approached.
\begin{itemize}
  \item \textbf{$T(x) = \sin(x)$ on $\mathbb{R}$. Does not apply; fixed points $\{0\}$.}
    $\mathbb{R}$ is complete and $T$ maps it into itself, but
    $\lambda = \sup_{x}\lvert\cos x\rvert = 1$, attained at $x=0$. So the theorem is unavailable.
    A fixed point exists all the same: $\sin x = x$ holds at $x=0$ and nowhere else, since
    $\lvert\sin x\rvert < \lvert x\rvert$ for $x \neq 0$. This is the first lesson of the table:
    the theorem is \emph{sufficient}, never necessary.
  \item \textbf{$T(x) = \tfrac{9}{10}\cos(x)$ on $\mathbb{R}$. Applies; one fixed point.}
    Here $\lvert T'(x)\rvert = \tfrac{9}{10}\lvert\sin x\rvert \leq \tfrac{9}{10} < 1$, so by the
    mean value theorem $T$ is a $\tfrac{9}{10}$-contraction; $\mathbb{R}$ is complete and
    $T(\mathbb{R}) \subseteq [-\tfrac9{10},\tfrac9{10}] \subseteq \mathbb{R}$. All three
    hypotheses hold, so there is exactly one fixed point. It has no closed form; iterating from
    any starting value gives $x^* \approx 0.693$.
  \item \textbf{$T(x) = x^2$ on $\mathbb{R}$. Does not apply; fixed points $\{0,1\}$.}
    $T$ is not Lipschitz at all on $\mathbb{R}$, since $T'(x) = 2x$ is unbounded. That it cannot
    apply is visible without any computation: $x^2 = x$ has the two solutions $0$ and $1$, and
    the theorem asserts \emph{uniqueness}, so any set of hypotheses implying it must fail here.
  \item \textbf{$T(x) = \tfrac12x^2$ on $(-1,1)$. Does not apply; fixed points $\{0\}$.}
    Two hypotheses fail together. The domain $(-1,1)$ is not complete, and
    $\lambda = \sup_{\lvert x\rvert<1}\lvert T'(x)\rvert = \sup\lvert x\rvert = 1$. This supremum
    is not attained, but a supremum of $1$ is still not $<1$, which is exactly the trap of
    \cref{rem:banach_hypotheses}. The map does send $(-1,1)$ into $[0,\tfrac12) \subseteq
    (-1,1)$. Solving $\tfrac12x^2 = x$ gives $x \in \{0,2\}$, of which only $0$ lies in the
    domain, so a unique fixed point exists regardless.
  \item \textbf{$T(x) = \tfrac12\big(x+\tfrac2x\big)$ on $\mathbb{Q}\cap(1,2)$. Does not apply;
    \emph{no} fixed points.} This row is the interesting one, and it is the Babylonian iteration
    for $\sqrt2$. Everything works except completeness:
    \begin{itemize}
      \item $T'(x) = \tfrac12 - \tfrac1{x^2}$, so on $(1,2)$ we have
        $-\tfrac12 \leq T'(x) \leq \tfrac14$ and hence $\lambda = \tfrac12 < 1$: a genuine
        contraction.
      \item $T$ maps $X$ into $X$: it sends rationals to rationals, and on $(1,2)$ its minimum
        is $T(\sqrt2) = \sqrt2 \approx 1.414$ while $T(1) = T(2) = \tfrac32$, so
        $T\big((1,2)\big) \subseteq [\sqrt2,\tfrac32] \subseteq (1,2)$.
      \item But $\mathbb{Q}\cap(1,2)$ is \textbf{not complete}.
    \end{itemize}
    And there is indeed no fixed point: $x = \tfrac12(x+\tfrac2x)$ forces $x^2 = 2$, so
    $x = \sqrt2 \notin \mathbb{Q}$. The iterates are the familiar rational approximations
    $\tfrac32, \tfrac{17}{12}, \tfrac{577}{408}, \dots$, which form a Cauchy sequence in $X$
    converging to a point that $X$ does not contain.

    This is worth contrasting with \cref{rem:banach_hypotheses}, where completeness failed on
    $(0,1]$ with the missing point sitting at an endpoint one is already suspicious of. Here the
    domain is an interval of rationals, the hole is invisible, and the contraction is a
    thoroughly practical algorithm. Completeness is not a technicality.
\end{itemize}
\end{exercisesolution}

\begin{exercisesolution}[Solution of \cref{ex:incomplete_metric_space_example}]
% Generator: Claude Opus 5 (High)
Take $X := (0,1]$ with the Euclidean distance $d(x,y) := \lvert x-y\rvert$ inherited from
$\mathbb{R}$. The sequence $x_k := 1/k$ lies in $X$ and is Cauchy, since $\lvert x_k - x_l\rvert \le
\max\{1/k, 1/l\} \to 0$. But it does not converge in $X$: its only candidate limit in $\mathbb{R}$
is $0$, and $0 \notin X$, so no point of $X$ can be its limit. Hence $(X,d)$ is not complete.

The same mechanism gives the other standard answers. Taking $X := \mathbb{Q}$ with
$d(x,y) := \lvert x-y\rvert$, any rational sequence converging to $\sqrt2$ is Cauchy without a limit
in $X$; taking $X := \mathbb{R} \setminus \{0\}$ works for the same reason as $(0,1]$.
\begin{ainote}
What is being exploited each time is that completeness is \emph{not} a property of the distance
alone: $(0,1]$ and $[0,1]$ carry the same formula for $d$, and one is complete while the other is
not. A Cauchy sequence is a sequence that is trying to converge, and the question is only whether
the space has kept the point it is aiming at. This is exactly the gap that the completion of a
metric space (\cref{sec:completion_of_a_metric_space}) is built to fill.
\end{ainote}
\end{exercisesolution}

\begin{exercisesolution}[Solution of \cref{ex:kobel_cauchy_single_choice}]
% Generator: Claude Opus 5 (High)
Only \textbf{(d)} is correct.
\begin{enumerate}[label=\textbf{(\alph*)}]
  \item \textbf{False.} A Cauchy sequence converges precisely when the ambient space has kept the
    point it is aiming at, and that is the definition of completeness
    (\cref{def:complete_metric_space}), not a theorem about Cauchy sequences.
    \Cref{ex:incomplete_metric_space_example} is a counterexample in one line.
  \item \textbf{False}, and this is the option the question is built around. Closedness is not a
    property a space has, only a property a subset has \emph{relative to} something larger: every
    metric space is closed in itself, so reading (b) as a hypothesis on $X$ alone makes it say
    nothing at all, and $X := (0,1]$ is again a counterexample. What is true is the relative
    statement: a closed subset of a \emph{complete} space is complete, and the ambient
    completeness is the part (b) omits.
  \item \textbf{False.} An interval need not be complete either. On $X := (0,1)$ the sequence
    $a_i := 1/i$ is Cauchy with no limit in $X$. Only the closed bounded intervals, and the closed
    unbounded ones, are complete.
  \item \textbf{True.} A compact metric space is sequentially compact
    (\cref{thm:sequential_compactness}), so $(a_i)$ has a convergent subsequence, and a Cauchy
    sequence with a convergent subsequence converges
    (\cref{ex:cauchy_elementary_facts}\textbf{(c)}). Hence every compact metric space is complete,
    which is \cref{ex:compact_finite_complete}.
\end{enumerate}
\end{exercisesolution}

\begin{exercisesolution}[Solution of \cref{ex:banach_hypotheses_true_false}]
% Generator: Claude Opus 5 (High)
\begin{enumerate}[label=\textbf{(\alph*)}]
  \item \textbf{True.} For $x,y \in [0,1]$,
    $\lvert x^2 - y^2\rvert = \lvert x+y\rvert\,\lvert x-y\rvert \le 2\lvert x-y\rvert$,
    since $x + y \le 2$ on $[0,1]$. So $2$ is a Lipschitz constant. It is even the optimal one, as
    $\lvert f'(x)\rvert = 2x$ approaches $2$ at $x = 1$.
  \item \textbf{False.} The space $X = [0,1]$ is complete, and $f$ does map $X$ into $X$; what fails
    is the contraction hypothesis. A contraction needs a Lipschitz constant \emph{strictly less
    than $1$}, and by \textbf{(a)} the best available constant here is $2$. Concretely, taking
    $x = 1$ and $y = 1 - \varepsilon$ gives
    $\lvert f(x) - f(y)\rvert = (2-\varepsilon)\varepsilon > \lvert x - y\rvert$ for small
    $\varepsilon > 0$: near $x = 1$ the map \emph{expands} distances.
  \item \textbf{True} --- and note the plural. Solving $x^2 = x$ on $[0,1]$ gives $x = 0$ and
    $x = 1$, so $f$ has exactly two fixed points.
\end{enumerate}
\begin{remark}
Parts \textbf{(b)} and \textbf{(c)} together are the point of the exercise. The Banach fixed point
theorem (\cref{thm:banach_fixed_point}) is a sufficient condition, not a necessary one: its
hypotheses fail here, and fixed points exist anyway. What the failure does cost is the
\emph{uniqueness} half of the conclusion, and sure enough there are two of them. Reading \qt{the
hypotheses fail} as \qt{there is no fixed point} is the standard mistake.
\end{remark}
\end{exercisesolution}

\begin{exercisesolution}[Proof of \cref{ex:banach_fixed_point_linear}]
% Generator: Gemini 3.6 Flash (High)
% Correction: Claude Opus 5 (High) --- rewritten for R^n. Flash's version worked in an abstract
% metric space and wrote d(f(x)/lambda, f(y)/lambda) = d(f(x),f(y))/lambda, which is meaningless
% there and false for a general metric even when scaling does make sense. It is the homogeneity of
% the Euclidean norm that is doing the work, so the norm has to be visible in the computation.
% Confirmed against WS22_Lösung.pdf p. 16, which runs the identical norm computation and likewise
% names the threshold lambda > L. The official solution agrees line for line.
\begin{enumerate}[label=\textbf{(\alph*)}]
  \item Fix $x,y \in \mathbb{R}^n$. Since the Euclidean norm is homogeneous and $f$ is
    $L$-Lipschitz,
\[
\lvert g(x) - g(y)\rvert = \left\lvert \frac{f(x) - f(y)}{\lambda} \right\rvert
  = \frac{1}{\lambda}\,\lvert f(x) - f(y)\rvert \leq \frac{L}{\lambda}\,\lvert x - y\rvert .
\]
Set $k := L/\lambda$. Since $\lambda > L > 0$ we have $k < 1$, and therefore $g$ is a contraction
with contraction constant $k$.

  \item The space $\mathbb{R}^n$ is complete, so the Banach Fixed Point Theorem
    (\cref{thm:banach_fixed_point}) applies to $g$ and yields a unique $x \in \mathbb{R}^n$ with
    $g(x) = x$. Multiplying $\lambda^{-1} f(x) = x$ by $\lambda$ turns this into $f(x) = \lambda x$,
    and the two equations have exactly the same solutions. Hence $f(x) = \lambda x$ has a unique
    solution.
\end{enumerate}
\begin{ainote}
The hypothesis is what makes this work, and it is worth naming: a Lipschitz $f$ need not be a
contraction, and need have no fixed point at all (translation $x \mapsto x + v$ is $1$-Lipschitz).
Dividing by $\lambda > L$ buys the missing factor. This is also why the original exam question says
\qt{for all sufficiently large $\lambda$} rather than naming a threshold: any $\lambda > L$ works,
but $L$ is not part of the data the question hands you.
\end{ainote}
\end{exercisesolution}

\begin{exercisesolution}[Proof of \cref{ex:examHS24_TF_banach}]
% Generator: Gemini 3.1 Pro (High)
\begin{enumerate}[label=\textbf{(\alph*)}]
  \item \textbf{True:} The derivative is $f'(x) = -\sin(x)$. Since $\lvert f'(x)\rvert = \lvert\sin x\rvert \le 1$ for all $x \in \mathbb{R}$, the mean value theorem gives $\lvert f(x) - f(y)\rvert \le 1 \cdot \lvert x-y\rvert$, so $f$ has Lipschitz constant $1$.
  \item \textbf{False:} The Banach fixed point theorem requires a \emph{strict} contraction, that is, a Lipschitz constant $\lambda < 1$. Here, the tightest possible constant is $\sup_{x \in [-\pi,\pi]} \lvert f'(x)\rvert = 1$ (attained at $x = \pm \pi/2$), so the map is non-expansive but not strictly contractive.
  \item \textbf{True:} The function $f(x) = \cos x$ is continuous and maps $[-\pi,\pi]$ into $[-1,1] \subseteq [-\pi,\pi]$. Since $g(x) := \cos x - x$ has $g(-\pi) = \pi - 1 > 0$ and $g(\pi) = -1 - \pi < 0$, the intermediate value theorem supplies a root, which is a fixed point of $f$.
\end{enumerate}
\end{exercisesolution}




% ============================================================
% Chapter 6 --- Compactness
% Originally: content/week-02.tex, \section{Compactness} (Corsin Week 2, Friday)
%           + content/week-03.tex, \section{Compactness --- continued} and
%             \subsection{Why do we care about compactness?} (Corsin Week 3, Monday)
%
% These were one topic split across two teaching weeks, joined in the old files by a
% \section{Compactness --- continued} heading and a \continuedfrom back-pointer. Both are
% gone: the material is continuous here, so nothing needs to point backwards.
% ============================================================
\chapter{Compactness}
\label{ch:compactness}

% Transition: Claude Opus 5 (Medium)
Compactness is the least self-explanatory definition in this part, so it is worth saying in
advance what it is \emph{for}. Finite sets are easy to handle: one can take a maximum, choose a
smallest radius, check finitely many cases. Infinite sets are not. Compactness is the property
that lets an infinite space be treated as if it were finite --- not because it is small, but
because every attempt to spread it out over infinitely many open pieces collapses back to
finitely many.

The chapter defines compactness by open covers, proves it equivalent to sequential compactness,
records the Heine--Borel theorem for $\mathbb{R}^n$ and how badly it fails elsewhere, and closes
with the three payoffs that make the definition worth having.

% Originally: content/week-02.tex, \section{Compactness} (Corsin Week 2, Friday)

\section{Compactness}
\label{sec:compactness}

% Transition: Claude Opus 5 (High)
Compactness is the least self-explanatory definition in this chapter, so it is worth saying in
advance what it is \emph{for}. Finite sets are easy to handle: one can take a maximum, choose a
smallest radius, check finitely many cases. Infinite sets are not. Compactness is the property
that lets an infinite space be treated as if it were finite --- not because it is small, but
because every attempt to spread it out over infinitely many open pieces collapses back to
finitely many. \Cref{sec:why_compactness_matters} collects the three payoffs
(continuous images, attained extrema, uniform continuity); the definition below is what buys
them.

\begin{definition}[Compact metric space]
\label{def:compact_metric_space}
A metric space $(X,d)$ is \newterm{compact} (\germanterm{kompakt}) if every open cover
(by sets open in the sense of \cref{def:open_closed}) has a finite subcover.
\end{definition}

\begin{notation}[Open covers]
\label{not:open_cover}
An \newterm{open cover} (\germanterm{offene \"Uberdeckung}) of $X$ is a family
$\{U_i\}_{i \in I}$ of open sets with $X = \bigcup_{i\in I}U_i$; the index set $I$ is allowed to
be arbitrary, in particular uncountable. A \newterm{subcover} is a subfamily
$\{U_i\}_{i \in J}$, $J \subseteq I$, that still covers $X$, and it is \emph{finite} when $J$
is. The order of quantifiers is what carries the content: \emph{every} cover must admit
\emph{some} finite subcover. Exhibiting one cover with a finite subcover proves nothing --- any
space has the cover $\{X\}$ --- whereas exhibiting one cover \emph{without} a finite subcover
disproves compactness outright, which is how \cref{ex:heine_borel_fails} and the
open-interval example below both work.
\end{notation}

% Source: Corsin Nick/Class Notes/Week 2.pdf, pp. 1--13

\begin{theorem}[Sequential compactness]
\label{thm:sequential_compactness}
A metric space $(X,d)$ is compact if and only if every sequence in $X$ has a convergent
(\cref{def:convergent_sequence}) subsequence in $X$.
\end{theorem}

\begin{ainote}
Neither Corsin nor the official lecture notes prove this here; the proof is given below as
\Cref{thm:sequential_topological_compactness_equivalence}, once the tools it needs (total
boundedness) are available.
\end{ainote}

\begin{center}
\begin{tikzpicture}[scale=1.1]
  \draw[blue!80!black, thick, fill=blue!5] plot[smooth cycle, tension=0.7] coordinates {(-2, -1.2) (-1.5, 1.5) (1.2, 1.8) (2.2, -0.5) (1.0, -1.5)};
  \node[blue!80!black, font=\large] at (-2.55, 1.55) {$X$};

  % The three sets must genuinely cover X -- ellipses chosen so every point of the
  % blob lies inside at least one of them (that is the whole point of a cover).
  \draw[red!80!black, thick, dotted, fill=red!15, fill opacity=0.35] (-1.15,0.05) ellipse (1.6 and 1.75);
  \node[red!80!black, font=\small] at (-1.75, -1.35) {$U_1$};

  \draw[orange!90!black, thick, dotted, fill=orange!20, fill opacity=0.35] (0.5,0.95) ellipse (1.5 and 1.25);
  \node[orange!90!black, font=\small] at (1.35, 1.75) {$U_2$};

  \draw[purple!80!black, thick, dotted, fill=purple!15, fill opacity=0.35] (1.15,-0.55) ellipse (1.5 and 1.3);
  \node[purple!80!black, font=\small] at (2.1, -1.35) {$U_3$};

  \node[below, font=\small] at (0, -2.1) {$X = U_1 \cup U_2 \cup U_3$};
\end{tikzpicture}
\end{center}

% Sonnet 5 (Medium)

% Originally: content/week-02.tex, \subsection{Compact subsets} (Corsin Week 2, Friday)

\subsection{Compact subsets}

A subset $K$ of a metric space $(X,d)$ is compact if the metric space $(K, d|_{K \times K})$ is
compact (\cref{def:compact_metric_space}). This is equivalent to the definition in the script:

\begin{definition}[9.63, quoted from the official lecture notes]
\label{def:sequential_topological_compactness}
Let $(X,d)$ be a metric space. A subset $K \subset X$ is called \textbf{compact} if one of the
following equivalent conditions hold:
\begin{enumerate}[label=\textbf{(\alph*)}]
  \item \label{item:seq_compact_def963}
    $K$ is \textbf{sequentially compact}: every sequence $(x_n)_{n\in\mathbb{N}}$ in $K$ has
    a subsequence that is convergent (\cref{def:convergent_sequence}) in $K$.
  \item \label{item:top_compact_def963}
    $K$ is \textbf{topologically compact}: every family of open sets
    (\cref{def:open_closed}) $\{U_i\}_{i \in I}$ that cover $K$, has a finite subcover.
\end{enumerate}
\end{definition}

\begin{center}
\begin{tikzpicture}[scale=1.1]
  \draw[blue!80!black, thick, fill=blue!5] plot[smooth cycle, tension=0.7] coordinates {(-2.2, -1.5) (-1.8, 1.8) (1.8, 1.8) (2.4, -0.8) (0.8, -1.8)};
  \node[blue!80!black, font=\large] at (-1.9, 1.5) {$X$};

  \draw[violet, thick, fill=violet!25] plot[smooth cycle, tension=0.7] coordinates {(-0.8, -0.5) (-0.5, 0.8) (0.8, 0.5) (0.6, -0.6)};
  \node[violet, font=\large] at (0.2, 0.1) {$K$};

  % As above: the three sets must cover all of K (they need not, and here do not,
  % cover all of X -- that is exactly the difference between the two pictures).
  \draw[green!60!black, thick, dotted] (-0.45,-0.1) ellipse (0.85 and 0.95);
  \node[green!60!black, font=\small] at (-1.15, -0.85) {$U_1$};

  \draw[orange!90!black, thick, dotted] (0.15,0.55) ellipse (0.85 and 0.75);
  \node[orange!90!black, font=\small] at (0.55, 1.45) {$U_2$};

  \draw[magenta!90!black, thick, dashed] (0.55,-0.25) ellipse (0.8 and 0.8);
  \node[magenta!90!black, font=\small] at (1.35, -0.85) {$U_3$};

  \node[below, font=\small] at (0, -2.2) {$K \subseteq U_1 \cup U_2 \cup U_3 \subseteq X$};
\end{tikzpicture}
\end{center}

% Quelle: Diego Torres Tejeda/Addendum - Sequential and Topological Compactness.pdf -- Sonnet 5 (Medium)

% Originally: content/week-02.tex, \subsection{Sequential and topological compactness agree}

\subsection{Sequential and topological compactness agree}

% Diego Torres Tejeda's supplementary note "Addendum - Sequential and Topological Compactness.pdf"
% is the one typeset document among his materials, the rest being handwritten scans. Its proof is
% adapted below into this project's notation and house style.
\begin{ainote}
Neither this equivalence nor \cref{thm:sequential_compactness} above is proved in the
transcribed notes. A complete, self-contained proof is supplied here
from a second tutor's addendum. By the reduction already noted above, compactness of
$K \subseteq X$ being compactness of the metric space $(K,d|_{K\times K})$, it suffices to treat a
whole metric space $X$ rather than a subset $K$.
\end{ainote}

\begin{lemma}[Continuous functions attain extrema on sequentially compact spaces]
\label{lem:sequential_compactness_extreme_values}
Let $(X,d)$ be sequentially compact (\cref{item:seq_compact_def963}) and $f : X \to \mathbb{R}$
continuous. Then $f$ attains a maximum and a minimum on $X$.
\end{lemma}

\begin{proof}
It suffices to find a minimum (apply the result to $-f$ for the maximum). Let
$m := \inf_{x\in X} f(x) \in [-\infty,\infty)$ and choose a sequence $(x_n)_{n}$ with
$f(x_n) \to m$. By sequential compactness, some subsequence $x_{n_k} \to x \in X$; by continuity,
$f(x) = \lim_k f(x_{n_k}) = m$. In particular $m \neq -\infty$, so $f$ attains its minimum $m$ at
$x$.
\end{proof}

\begin{remark}
This is exactly the extreme value theorem (\cref{item:extreme_value_theorem}) proved from first
principles, without presupposing \cref{thm:sequential_compactness}. Avoiding that appeal is
essential here, since that theorem is precisely what we are about to prove.
\end{remark}

\begin{definition}[Totally bounded]
\label{def:totally_bounded}
A metric space $(X,d)$ is \newterm{totally bounded} \germanfor{total beschr\"ankt} if for
every $\varepsilon > 0$ there exist finitely many points $x_1,\dots,x_n \in X$ such that
\[X = \bigcup_{k=1}^{n} B_\varepsilon(x_k).\]
\end{definition}

% Supplement: Analysis_II_Script_v1.pdf, p. 22
\begin{exercise}[Totally bounded is strictly stronger than bounded]
\label{ex:totally_bounded_implies_bounded}
Let $X$ be a metric space. Show that if $X$ is totally bounded, then it is bounded, i.e.\
$\sup_{x,x'\in X}d(x,x') < \infty$. Find an example of a bounded metric space that is
\emph{not} totally bounded.
\end{exercise}

\begin{exercisesolution}[Proof of \cref{ex:totally_bounded_implies_bounded}]
% Generator: Claude Sonnet 5 (High)
\textbf{Totally bounded $\Rightarrow$ bounded.} Apply \cref{def:totally_bounded} with
$\varepsilon := 1$: there are $x_1,\dots,x_n$ with $X = \bigcup_{k=1}^n B_1(x_k)$. Let
$M := \max_{i,j}d(x_i,x_j)$ (finite, a maximum over finitely many pairs). For any
$x,x' \in X$, choose $i,j$ with $x \in B_1(x_i)$, $x' \in B_1(x_j)$; the triangle inequality
gives $d(x,x') \leq d(x,x_i)+d(x_i,x_j)+d(x_j,x') < 1+M+1$, so $X$ is bounded.

\textbf{Not conversely.} Let $X$ be any infinite set with the discrete metric. Then
$\sup_{x,x'}d(x,x') = 1 < \infty$, so $X$ is bounded. But $X$ is not totally bounded: for
$\varepsilon := \tfrac12$, every ball $B_{1/2}(x) = \{x\}$ (\cref{def:open_ball}), so no
\emph{finite} union of such balls can cover the infinite set $X$.
\end{exercisesolution}

\begin{remark}
This is the same asymmetry \cref{lem:sequentially_compact_totally_bounded} proves in one
direction only: sequential compactness forces \emph{both} boundedness and total boundedness, and
this exercise shows the two are genuinely different hypotheses. Total boundedness is the strictly
stronger of the two, and the gap opens up in infinite dimensions.
\end{remark}

\begin{lemma}[Sequential compactness gives (total) boundedness]
\label{lem:sequentially_compact_totally_bounded}
Let $(X,d)$ be sequentially compact. Then $X$ is bounded and totally bounded
(\cref{def:totally_bounded}).
\end{lemma}

\begin{proof}
\textbf{Bounded.} Fix $x_0 \in X$; the function $f(y) := d(x_0,y)$ is continuous, so by
\Cref{lem:sequential_compactness_extreme_values} it attains a maximum $M \geq 0$ on $X$. Then
$X = B_{M+1}(x_0)$.

\textbf{Totally bounded.} Fix $\varepsilon > 0$. Construct a sequence inductively: let
$x_1 \in X$ be arbitrary, and, having chosen $x_1,\dots,x_n$, let
$U_n := \bigcup_{k=1}^n B_\varepsilon(x_k)$. If $U_n = X$, stop, since we then have the desired
finite cover. Otherwise pick $x_{n+1} \in X\setminus U_n$, so that $d(x_{n+1},x_k) \geq \varepsilon$ for
all $k \leq n$. If this process never terminates, sequential compactness gives a convergent, hence
Cauchy, subsequence $(x_{n_k})_k$; but $d(x_{n_k},x_{n_{k+1}}) \geq \varepsilon$ for all $k$
contradicts the Cauchy property. So, the process must terminate after finitely many steps, giving
the desired finite $\varepsilon$-net.

% Sonnet 5 (Medium)
\begin{center}
\begin{tikzpicture}[scale=0.9]
  \draw[blue!70!black, thick] plot[smooth cycle, tension=0.7] coordinates
    {(-2,-1) (-1.5,1.3) (1.2,1.5) (2.1,-0.4) (0.8,-1.6)};
  % Total boundedness means the finitely many balls really do cover X, so the
  % radius must be large enough to reach the corners of the blob.
  \foreach \p in {(-1.2,0.1),(0.0,1.0),(1.2,0.0),(-0.1,-0.9)} {
    \fill \p circle (1.5pt);
    \draw[orange!90!black, dashed] \p circle (1.45);
  }
  \node[font=\scriptsize, below left] at (-1.2,0.1) {$x_1$};
  \node[font=\scriptsize, above right] at (0.0,1.0) {$x_2$};
  \node[font=\scriptsize, right] at (1.2,0.0) {$x_3$};
  \node[font=\scriptsize, below] at (-0.1,-0.9) {$x_4$};
  \node[font=\scriptsize, orange!90!black, below] at (0,-2.6) {$X = B_\varepsilon(x_1)\cup\dots\cup B_\varepsilon(x_4)$};
\end{tikzpicture}
\end{center}
\end{proof}

\begin{theorem}[Sequential and topological compactness agree]
\label{thm:sequential_topological_compactness_equivalence}
A metric space $(X,d)$ is sequentially compact (\cref{item:seq_compact_def963}) if and only if it
is topologically compact (\cref{def:compact_metric_space}). In particular,
\Cref{thm:sequential_compactness} holds.
\end{theorem}

\begin{proof}
\textbf{($\Leftarrow$)} Suppose $X$ is not sequentially compact, so some sequence
$(x_n)_{n\in\mathbb{N}}$ has no convergent subsequence. Then no $x \in X$ is the limit of any
subsequence, so there exist $\varepsilon_x > 0$ and $N_x \in \mathbb{N}$ such that
$B_{\varepsilon_x}(x) \cap \{x_n : n \geq N_x\} = \emptyset$. The balls
$\{B_{\varepsilon_x}(x)\}_{x \in X}$ cover $X$, so topological compactness gives finitely many
$x_1,\dots,x_m$ with $X = \bigcup_{k=1}^m B_{\varepsilon_{x_k}}(x_k)$. With
$N := \max\{N_{x_1},\dots,N_{x_m}\}$, every $x_n$ with $n \geq N$ then avoids
$B_{\varepsilon_{x_1}}(x_1) \cup \dots \cup B_{\varepsilon_{x_m}}(x_m) = X$, which is absurd.
Therefore, $X$ is sequentially compact.

\begin{center}
\begin{tikzpicture}[scale=1.5]
  % Space K
  \draw[very thick, fill=gray!10] plot[smooth cycle, tension=0.7] coordinates {(-1.5, -0.5) (-1.8, 0.5) (-0.5, 1) (1.2, 0.8) (1.8, 0.2) (1.2, -0.8) (0, -1)};
  
  % Sequence points and path
  \fill[purple] (-1.2, 0.2) circle (1pt) node[left] {\tiny $x_1$};
  \fill[purple] (-0.8, -0.6) circle (1pt) node[below] {\tiny $x_2$};
  \fill[purple] (0.2, 0.6) circle (1pt) node[above] {\tiny $x_3$};
  \fill[purple] (1.0, -0.5) circle (1pt) node[right] {\tiny $x_4$};
  \draw[->, purple, dashed] (-1.2, 0.2) to[bend right=15] (-0.8, -0.6);
  \draw[->, purple, dashed] (-0.8, -0.6) to[bend left=15] (0.2, 0.6);
  \draw[->, purple, dashed] (0.2, 0.6) to[bend right=15] (1.0, -0.5);
  
  % Accumulation point
  \fill[green!60!black] (0.8, 0.2) circle (1.5pt) node[above right] {\small $x$};
  \draw[green!60!black, dashed, thick] (0.8, 0.2) circle (0.35);
  
  % Subsequence converging
  \fill[purple] (0.6, -0.2) circle (1pt);
  \fill[purple] (0.9, 0) circle (1pt);
  \fill[purple] (0.7, 0.3) circle (1pt);
  \draw[->, purple, dotted, thick] (1.0, -0.5) to[bend left=10] (0.6, -0.2);
  \draw[->, purple, dotted, thick] (0.6, -0.2) to[bend right=10] (0.9, 0);
  \draw[->, purple, dotted, thick] (0.9, 0) to[bend left=10] (0.7, 0.3);
  \draw[->, purple, dotted, thick] (0.7, 0.3) -- (0.78, 0.22);
  
  \node[green!60!black, anchor=west] at (1.3, 0.4) {\footnotesize accumulation point};
  \node[purple, anchor=west] at (1.3, 0.1) {\footnotesize subsequence $x_{n_k} \to x$};
\end{tikzpicture}
\end{center}

\textbf{($\Rightarrow$)} Suppose $X$ is sequentially compact and let $\{U_i\}_{i\in I}$ be an open
cover; we may assume no single $U_i$ equals $X$ (else $\{U_i\}$ is itself a finite subcover).
Define
\[r(x) := \tfrac12\sup\{r \geq 0 : B_r(x) \subseteq U_i \text{ for some } i \in I\}, \qquad x \in X.\]
\textit{$r$ is finite and positive.} Since $\{U_i\}$ covers $X$ and each $U_i$ is open, some
$U_{i_0} \ni x$ contains a ball around $x$, so $r(x) > 0$; and since $X$ is bounded
(\cref{lem:sequentially_compact_totally_bounded}) but no $U_i$ equals $X$, $r(x) < \infty$.

\textit{$r$ is Lipschitz, hence continuous.} (That is, $|r(x)-r(y)| \leq L\,d(x,y)$ for a
constant $L$; the notion is named and placed in context as \cref{def:lipschitz_continuous},
but only the displayed estimate is used here.) Write
$R(x) := \sup\{\rho \geq 0 : B_\rho(x) \subseteq U_i \text{ for some } i\}$, so that
$r = \tfrac12 R$; it is $R$ that the argument bounds, and the factor $\tfrac12$ must be carried
through.

Let $x,y \in X$ with $d(x,y) < r(x)$. Since $r(x) < R(x)$, we have
$B_{r(x)}(x) \subseteq U_i$ for some $i$, so $y \in U_i$, and the triangle inequality gives
\[B_{r(x)-d(x,y)}(y) \subseteq B_{r(x)}(x) \subseteq U_i
\qquad\Longrightarrow\qquad R(y) \geq r(x)-d(x,y).\]
Halving, $r(y) = \tfrac12R(y) \geq \tfrac12\big(r(x)-d(x,y)\big)$, which also holds trivially
when $d(x,y) \geq r(x)$ since $r(y) \geq 0$. By symmetry,
\[|r(x)-r(y)| \leq \tfrac12\,d(x,y),\]
so $r$ is $\tfrac12$-Lipschitz, and in particular continuous, which is all that is needed below.

By \cref{lem:sequential_compactness_extreme_values}, $r$ attains a minimum
$\varepsilon := \min_{x\in X} r(x) > 0$ on $X$. Then $B_\varepsilon(x) \subseteq U_i$ for some $i$,
for every $x \in X$. By \cref{lem:sequentially_compact_totally_bounded}, $X$ is totally bounded,
so there are $x_1,\dots,x_n \in X$ with $X = \bigcup_{k=1}^n B_\varepsilon(x_k)$; choosing, for
each $k$, an index $i_k$ with $B_\varepsilon(x_k) \subseteq U_{i_k}$, the sets
$U_{i_1},\dots,U_{i_n}$ form a finite subcover of $X$.
\end{proof}

\begin{remark}[Total boundedness, more generally]
\label{rem:compact_iff_complete_totally_bounded}
Diego's note records a sharper characterization as a further exercise: $(X,d)$ is compact if and
only if it is complete and totally bounded. \Cref{lem:sequentially_compact_totally_bounded} gives
one direction (compact $\implies$ complete, by \cref{ex:compact_finite_complete} below, and
totally bounded); for the converse, a Cauchy subsequence can be extracted from any sequence using
total boundedness (cover by finitely many $\tfrac1k$-balls for $k=1,2,\dots$ and pass to nested
subsequences), and it converges by completeness.
\end{remark}

% Quelle: Diego Torres Tejeda/Addendum - Sequential and Topological Compactness.pdf, \S4(b)
\begin{definition}[Lebesgue number]
\label{def:lebesgue_number}
Let $(X,d)$ be a metric space and $\{U_i\}_{i\in I}$ an open cover of $X$. A number
$\varepsilon > 0$ is a \newterm{Lebesgue number} \germanfor{Lebesgue-Zahl} for the cover if
\[\forall x \in X\ \exists i \in I: B_\varepsilon(x) \subseteq U_i.\]
\end{definition}

\begin{remark}[Lebesgue's covering lemma]
The proof of \cref{thm:sequential_topological_compactness_equivalence} ($\Rightarrow$) shows
more than stated: on a compact metric space, \emph{every} open cover has a Lebesgue number (take
$\varepsilon := \min_{x\in X} r(x)$ as constructed there). This is known as \newterm{Lebesgue's
covering lemma} \germanfor{Lebesguesches \"Uberdeckungslemma}.
\end{remark}

\begin{proposition}[Compactness implies uniform continuity]
\label{prop:compact_implies_uniformly_continuous}
Let $(X,d_X)$, $(Y,d_Y)$ be metric spaces with $X$ compact, and let $f : X \to Y$ be continuous.
Then $f$ is uniformly continuous (\cref{def:uniformly_continuous}).
\end{proposition}

\begin{proof}
% Generator: Claude Sonnet 5 (Medium), after Diego Torres Tejeda/Addendum, §4(b)
Let $\varepsilon > 0$. The sets $f^{-1}(B_{\varepsilon/2}(y))$, $y \in Y$, are open
(\cref{def:continuous}) and cover $X$, so they admit a Lebesgue number $\delta > 0$
(\cref{def:lebesgue_number}). For any $x,x' \in X$ with $d_X(x,x') < \delta$, the ball
$B_\delta(x)$ lies inside some $f^{-1}(B_{\varepsilon/2}(y))$, so $x' \in B_\delta(x)$ gives
$f(x),f(x') \in B_{\varepsilon/2}(y)$, hence $d_Y(f(x),f(x')) < \varepsilon$ by the triangle
inequality. Since $\delta$ did not depend on $x$, $f$ is uniformly continuous.
\end{proof}

% Supplement: Sascha Brack/Class Notes/Week_02_Notes_Friday_Updated.pdf, p. 6
\begin{proposition}[A closed subset of a compact set is compact]
\label{prop:closed_subset_of_compact}
Let $(X,d)$ be a metric space, $K \subseteq X$ compact, and $A \subseteq K$ closed in $X$. Then
$A$ is compact.
\end{proposition}

\begin{proof}
Let $\mathcal{U}$ be an open cover of $A$. Since $A$ is closed, $X\setminus A$ is open, so
\[\mathcal{U}' := \mathcal{U} \cup \{X\setminus A\}\]
is an open cover of $K$, covering $A$ by assumption and everything outside $A$ by the added
set. Compactness of $K$ gives a finite subcover
$U_{1},\dots,U_{m} \in \mathcal{U}'$ of $K$, hence of $A$.

\begin{center}
\begin{tikzpicture}[scale=1.2]
  % Compact set K
  \draw[very thick] plot[smooth cycle, tension=0.7] coordinates {(-1.5, -1) (-2, 0.5) (0, 1.2) (2, 0.5) (1.5, -1.2)};
  \node[above left, font=\large] at (-1.8, 0.5) {$K$};
  
  % Closed set A
  \draw[thick, fill=gray!20] (0, 0.5) -- (-0.6, 0) -- (0, -0.5) -- (0.6, 0) -- cycle;
  \node[font=\large] at (0, 0) {$A$};
  
  % Open cover V
  \draw[dashed, blue] (-1, 0.8) rectangle (1.2, -0.8);
  \node[blue, right] at (1.2, 0.5) {$\mathcal{V}$ covers $K$, so also $A$};
\end{tikzpicture}
\end{center}

Now simply discard $X\setminus A$ from that list if it occurs: it contains no point of $A$, so
the remaining sets still cover $A$. They all belong to $\mathcal{U}$, and there are finitely
many. As $\mathcal{U}$ was arbitrary, $A$ is compact.
\end{proof}

% Reclassified from ainote to remark: a proof technique worth naming is mathematics.
\begin{remark}[Enlarge the cover, then discard the extra set]
The trick just used is the standard way to transfer compactness downward, and it is worth
remembering in the form \emph{enlarge the cover by the complement of the set you care about, then
throw that complement away again}. It reappears whenever one needs a compact subset of a compact
set. Note that the hypothesis \qt{closed} is essential, since $(0,1) \subseteq [0,1]$ sits inside
a compact set without being compact.
\end{remark}

% Source: Corsin Nick/Class Notes/Week 2.pdf, p. 15
\begin{exercise}[Compact sets are complete]
\label{ex:compact_finite_complete}
Let $(X,d)$ be a metric space.
\begin{enumerate}[label=\textbf{(\alph*)}]
  \item Show that every finite subset $\{x_1, \dots, x_n\} \subseteq X$ is compact.
  \item If $X$ is compact, is it complete (\cref{def:complete_metric_space})?
\end{enumerate}
\end{exercise}

% Kept inline: solved directly in Corsin Nick/Class Notes/Week 2.pdf, p. 15. (The chapter-level
% "pp. 1--13" Source comment above predates this exercise; the PDF actually runs to p. 16+.)
\begin{exercisesolution}[Proof of \cref{ex:compact_finite_complete}]
\begin{enumerate}[label=\textbf{(\alph*)}]
  \item Let $\{U_\alpha\}_{\alpha \in I}$ be an open cover of $\{x_1, \dots, x_n\}$. For each
    $i = 1, \dots, n$, pick $U_{\alpha_i}$ such that $x_i \in U_{\alpha_i}$. Then
    $\{U_{\alpha_i}\}_{i=1,\dots,n}$ is a finite subcover.
  \item \textbf{Yes.} Let $(x_n)_{n=0}^{\infty}$ be a Cauchy sequence in $X$. Then, because $X$ is
    compact, we can extract a convergent subsequence such that $\lim_{i \to \infty} x_{n_i} = x \in X$.

    Fix $\varepsilon > 0$. We find $N_1 \in \mathbb{N}$ such that
    \[d(x_n, x_m) < \varepsilon \quad \forall n, m > N_1\]
    and $N_2 \in \mathbb{N}$ such that
    \[d(x, x_{n_i}) < \varepsilon \quad \forall i > N_2.\]
    Let $N := \max\{N_1, N_2\}$. For any $k > N$, we have:
    \[d(x, x_k) \leq d(x, x_{n_k}) + d(x_{n_k}, x_k) < 2\varepsilon\]
    where we used the fact that $n_k \geq k > N$. Thus, any Cauchy sequence converges.
\end{enumerate}
\end{exercisesolution}

% Moved here from after the two aiexercises below (review item D3): its "above" pointed
% past them to this solution, which is what it actually comments on.
\begin{ainote}
Corsin writes $n_k > k$ above; for a subsequence the standing fact is $n_k \geq k$. Immaterial to
the argument.
\end{ainote}

% Source: Jérôme Paschoud/Notizen/Woche 3.1 Banach und Kompaktheit.pdf, p. 3
% Generator: Gemini 3.6 Flash (High)
\begin{exercise}[Compactness of a convergent sequence]
\label{ex:convergent_sequence_compact}
Let $(X,d)$ be a metric space. Let $(x_n)_{n=0}^\infty$ be a sequence such that $x_n \to \bar{x}$. Prove that the set
\[E := \{x_n \mid n \in \mathbb{N}\} \cup \{\bar{x}\}\]
is compact.
\end{exercise}

\begin{exercisesolution}[Proof of \cref{ex:convergent_sequence_compact}]
Using topological compactness: Let $\{U_i\}_{i \in I}$ be an open cover of $E$. Since $\bar{x} \in E$, there exists some $i_* \in I$ such that $\bar{x} \in U_{i_*}$. Because $U_{i_*}$ is open and $x_n \to \bar{x}$, there exists $N \in \mathbb{N}$ such that $x_n \in U_{i_*}$ for all $n \geq N$. For the remaining finitely many points $x_0, \dots, x_{N-1}$, we can pick sets $U_{i_0}, \dots, U_{i_{N-1}}$ from the cover containing each of them. Then $\{U_{i_*}, U_{i_0}, \dots, U_{i_{N-1}}\}$ is a finite subcover of $E$.

Using sequential compactness: Let $(y_k)_{k=0}^\infty$ be a sequence in $E$. If $(y_k)$ takes values in a finite subset of $E$, it has a constant (hence convergent) subsequence. Otherwise, it takes infinitely many distinct values from $E$, which means it must contain a subsequence of $(x_n)$ progressing towards $\bar{x}$. In either case, it contains a convergent subsequence with limit in $E$. Thus $E$ is sequentially compact, hence compact.
\end{exercisesolution}

% Source: Jérôme Paschoud/Notizen/Woche 3.1 Banach und Kompaktheit.pdf, pp. 3, 6
% Generator: Gemini 3.6 Flash (High)
\begin{exercise}[Compactness of a finite set grid]
\label{ex:finite_set_compact}
Let 
\[X := \big\{ (x,y) \in \mathbb{R}^2 \mid x,y \in \mathbb{Z} \wedge |x|,|y| \leq 3 \big\}.\]
Show that $X$ is compact.
\end{exercise}

\begin{exercisesolution}[Proof of \cref{ex:finite_set_compact}]
The set $X$ contains exactly $7 \times 7 = 49$ points, so it is finite. By \cref{ex:compact_finite_complete}, every finite subset of a metric space is compact. Alternatively, using sequential compactness (as hinted in the notes), any sequence in $X$ must take at least one value infinitely often (by the pigeonhole principle), thus it contains a constant, hence convergent, subsequence.
\end{exercisesolution}

% Originally: content/week-03.tex, \section{Compactness --- continued} (Corsin Week 3, Monday)

% Corsin taught this a week after the definition, and opened it as "Compactness --- continued"
% with a back-pointer. In a topic chapter the two sit together, so the heading names its content
% and the \continuedfrom pointer is gone.
\section{Heine--Borel}
\label{sec:heine_borel}

\begin{theorem}[Heine--Borel]
\label{thm:heine_borel}
A subset $K \subseteq \mathbb{R}^n$ equipped with the \textbf{standard metric} is
\textbf{compact} if and only if it is \newterm{closed} and \newterm{bounded}
\germanfor{beschr\"ankt}.
\end{theorem}

\begin{ainote}
Neither Corsin nor this document proves \cref{thm:heine_borel}; it is proved in the lecture, and
the full argument is longer than anything else in these two chapters. It is worth knowing what
goes into it, though, since every ingredient has already appeared:
\begin{itemize}
  \item \textbf{($\Rightarrow$)} is short and general. A compact metric space is bounded and
    complete (\cref{lem:sequentially_compact_totally_bounded},
    \cref{ex:compact_finite_complete}), and a complete subset of the complete space
    $\mathbb{R}^n$ is closed (\cref{rem:complete_vs_closed}).
  \item \textbf{($\Leftarrow$)} is where the real content sits, and it is exactly
    \textbf{Bolzano--Weierstrass}: a bounded sequence in $\mathbb{R}^n$ has a convergent
    subsequence, obtained by bisecting a box $n$ times over, or by extracting coordinatewise
    (\cref{prop:Rn_coordinatewise_complete}) $n$ times in succession. Closedness then puts the
    limit back inside $K$, giving sequential compactness, which is compactness by
    \cref{thm:sequential_topological_compactness_equivalence}.
\end{itemize}
The reason the theorem is special to $\mathbb{R}^n$ with the standard metric is that both steps
use the completeness of $\mathbb{R}$ and the finiteness of the dimension.
\Cref{ex:heine_borel_fails} below shows how thoroughly it fails otherwise.
\end{ainote}

% Source: Lukas Krause/Week 3 March 2-6 Notes Lukas Krause.pdf, p. 2
% Generator: Gemini 3.6 Flash (High)
\begin{proposition}[Topologically compact subsets of $\mathbb{R}^n$ are bounded]
\label{prop:topologically_compact_is_bounded}
If a subset $S \subseteq \mathbb{R}^n$ is topologically compact, then $S$ is bounded.
\end{proposition}

\begin{proof}
Consider the collection of open balls $\mathcal{U} := \{B_k(0) \mid k \in \mathbb{N}_{\ge 1}\}$ centered at the origin. Since every point $x \in \mathbb{R}^n$ has $|x| < k$ for some large integer $k$, $\mathcal{U}$ forms an open cover of $\mathbb{R}^n$, and hence an open cover of $S$:
\[
S \;\subseteq\; \mathbb{R}^n \;=\; \bigcup_{k=1}^\infty B_k(0).
\]
By topological compactness of $S$, there exists a finite subcover $\{B_{k_1}(0), B_{k_2}(0), \dots, B_{k_m}(0)\}$. Setting $K := \max\{k_1, k_2, \dots, k_m\}$, we obtain
\[
S \;\subseteq\; \bigcup_{j=1}^m B_{k_j}(0) \;=\; B_K(0).
\]
Therefore, $S$ is contained in the ball of radius $K$ centered at the origin, which means $S$ is bounded.
\end{proof}

% Supplement: Sascha Brack/Class Notes/Week_02_Notes_Friday_Updated.pdf, p. 7
\begin{example}[Compactness via Heine-Borel and continuity]
\label{ex:compactness_heine_borel_examples}
Consider the following subsets and determine whether they are compact under the standard Euclidean metric:
\begin{enumerate}[label=\textbf{(\alph*)}]
  \item \textbf{Is $\mathbb{N} \subset \mathbb{R}$ compact?} No. While it is closed in $\mathbb{R}$, it is not bounded. By Heine-Borel, it cannot be compact.
  \item \textbf{Is $[0,1] \times [0,1] \subset \mathbb{R}^2$ complete?} Yes. The square $[0,1]^2$ is closed and bounded in $\mathbb{R}^2$, hence compact by Heine-Borel. As every compact metric space is complete (\cref{ex:compact_finite_complete}), it is complete.
  \item Let $f: \mathbb{R} \to \mathbb{R}^3, x \mapsto (\sin x, \cos x, 1)$. \textbf{Is $f([0,\pi/2])$ compact in $\mathbb{R}^3$?} Yes. We do not need to show this manually using sequences or open covers. Since the domain $[0,\pi/2]$ is compact in $\mathbb{R}$ (closed and bounded) and the function $f$ is continuous, its image $f([0,\pi/2])$ is compact in $\mathbb{R}^3$ (by \cref{item:continuity_preserves_compact}).
\end{enumerate}
\end{example}

% Supplement: Sascha Brack/Class Notes/Week_03_Notes_Monday.pdf, p. 3
\begin{exercise}[Topological properties of various subsets]
\label{ex:topological_properties_table}
Decide whether the following subsets are open, closed, compact, or complete (under the standard Euclidean metric):
\begin{enumerate}[label=\textbf{(\alph*)}]
  \item $[0,1] \subseteq \mathbb{R}$
  \item $\mathbb{Q} \subseteq \mathbb{R}$
  \item $B_1(0) \subseteq \mathbb{R}^3$
  \item $B_1(0) \cap B_1(1) \subseteq \mathbb{R}^2$
  \item $\{0\} \cup \{1/n \mid n \in \mathbb{N}_{\ge 1}\} \subseteq \mathbb{R}$
  \item $\bigcup_{n=1}^\infty B_{1/n}(n) \subseteq \mathbb{R}$
  \item $(-\infty, 1] \subseteq \mathbb{R}$
  \item $f([0,1])$ where $f : \mathbb{R} \to \mathbb{R}^n$ is continuous
  \item $f^{-1}([0,1])$ where $f : \mathbb{R}^n \to \mathbb{R}$ is continuous
  \item $f^{-1}((0,1))$ where $f : \mathbb{R}^n \to \mathbb{R}$ is continuous
  \item $S := \operatorname{span}\{(1,0,0)\transp, (3,1,0)\transp\} \subseteq \mathbb{R}^3$
  \item $\{x \in \mathbb{R}^3 \mid \exists s \in S \text{ s.t.\ } d(x, s) < 1\}$ for a subset $S \subseteq \mathbb{R}^3$
\end{enumerate}
\end{exercise}

\begin{exercisesolution}[Proof of \cref{ex:topological_properties_table}]
\begin{enumerate}[label=\textbf{(\alph*)}]
  \item $[0,1] \subseteq \mathbb{R}$: \textbf{Closed}, \textbf{compact}, \textbf{complete}. (Not open).
  \item $\mathbb{Q} \subseteq \mathbb{R}$: \textbf{None}. It is neither open nor closed, not bounded (so not compact), and not complete.
  \item $B_1(0) \subseteq \mathbb{R}^3$: \textbf{Open}.
  \item $B_1(0) \cap B_1(1) \subseteq \mathbb{R}^2$: \textbf{Open} (finite intersection of open sets).
  \item $\{0\} \cup \{1/n \mid n \in \mathbb{N}_{\ge 1}\} \subseteq \mathbb{R}$: \textbf{Closed}, \textbf{compact}, \textbf{complete} (bounded and contains its only limit point $0$).
  \item $\bigcup_{n=1}^\infty B_{1/n}(n) \subseteq \mathbb{R}$: \textbf{Open} (arbitrary union of open sets).
  \item $(-\infty, 1] \subseteq \mathbb{R}$: \textbf{Closed}, \textbf{complete}.
  \item $f([0,1])$: \textbf{Closed}, \textbf{compact}, \textbf{complete} (continuous image of a compact set).
  \item $f^{-1}([0,1])$: \textbf{Closed}, \textbf{complete} (continuous preimage of a closed set).
  \item $f^{-1}((0,1))$: \textbf{Open} (continuous preimage of an open set).
  \item $S := \operatorname{span}\{(1,0,0)\transp, (3,1,0)\transp\} \subseteq \mathbb{R}^3$: \textbf{Closed}, \textbf{complete} (it is a $2$-dimensional subspace, hence closed in $\mathbb{R}^3$).
  \item $\{x \in \mathbb{R}^3 \mid \exists s \in S \text{ s.t.\ } d(x, s) < 1\}$: \textbf{Open}. This set is equivalently written as $\{x \in \mathbb{R}^3 \mid d(x, S) < 1\}$ or $\bigcup_{s \in S} B_1(s)$, which is an arbitrary union of open balls.
\end{enumerate}
\end{exercisesolution}

This captures our intuition of a \qt{compact} set. We will briefly show that Heine--Borel is
\textbf{false} for $\mathbb{R}^n$ with an \textbf{arbitrary metric}.

\begin{exercise}[Heine--Borel fails for arbitrary metrics]
\label{ex:heine_borel_fails}
Let $(X,d)$ be a metric space.
\begin{enumerate}[label=\textbf{(\alph*)}]
  \item \textit{(Optional)} Show that $\tilde d(x,y) := \dfrac{d(x,y)}{1 + d(x,y)}$ is a metric on
    $X$.
  \item Show that, if $U \subseteq X$ is open with respect to $d$, then $U$ is open with respect
    to $\tilde d$.
  \item Show that, for $X = \mathbb{R}$, $d(x,y) = \dfrac{|x-y|}{1+|x-y|}$,
    $\mathbb{N} \subseteq X$ is closed and bounded but not compact.
\end{enumerate}
\end{exercise}

% Correction: Claude Opus 5 (High) --- part (c) read $X = \mathbb{R}^n$ in the source. Flagged by
% Gemini 3.1 Pro (2026-08-11) as confusing to leave standing in the statement, and it is: the
% ainote below diagnosed the defect but the reader still met the wrong hypothesis first.
\begin{ainote}
Part \textbf{(c)} is stated in the source for $X = \mathbb{R}^n$, but it then works with
$\mathbb{N} \subseteq X$ and the absolute value $|x-y|$, which only parses for $n = 1$. It is
corrected to $X = \mathbb{R}$ above. Nothing is lost: the same argument runs verbatim in
$\mathbb{R}^n$ once $|x-y|$ is read as the Euclidean norm and $\mathbb{N}$ as
$\mathbb{N} \times \{0\}^{n-1}$.
\end{ainote}

\begin{aiexample}[Open Cover of $(0,1)$ with No Finite Subcover]
% Generator: Gemini 3.6 Flash (Medium)
Consider the open interval $X = (0,1)$ under the standard Euclidean metric. The family of open sets
\[
U_n := \left(\frac{1}{n}, 1\right) \quad \text{for } n \in \mathbb{N}_{\ge 2}
\]
forms an open cover of $(0,1)$ because for any $x \in (0,1)$, there exists $n \in \mathbb{N}$ with $1/n < x$. However, any finite subcollection $\{U_{n_1}, \dots, U_{n_k}\}$ has a maximum index $N := \max\{n_1, \dots, n_k\}$, so its union is $U_N = (1/N, 1)$, which fails to cover points in $(0, 1/N]$. Thus $(0,1)$ is bounded but not compact.
\end{aiexample}

% Supplement: Analysis_II_Script_v1.pdf, p. 22
\begin{example}[Heine--Borel needs $\mathbb{R}^n$ itself, not just the standard-looking metric]
\label{ex:rationals_not_compact}
$\mathbb{Q} \cap [0,2]$, with the metric inherited from the standard metric on $\mathbb{R}$, is
closed and bounded \emph{as a subset of $\mathbb{Q}$}, and yet it is not compact. The family
\[\mathbb{Q}\cap[0,2] \ =\ \big(\mathbb{Q}\cap[0,\sqrt2)\big) \ \cup\!\!
  \bigcup_{p\in\mathbb{Q},\,p>\sqrt2}\!\big(\mathbb{Q}\cap(p,2]\big)\]
is an open cover of $\mathbb{Q}\cap[0,2]$ (each set is the trace of an open subset of
$\mathbb{R}$, hence open in the subspace $\mathbb{Q}\cap[0,2]$): every rational in $[0,2]$ is
either $<\sqrt2$ or $>\sqrt2$, since $\sqrt2 \notin \mathbb{Q}$. Any finite subcover uses only
finitely many of the sets $\mathbb{Q}\cap(p,2]$, hence has a largest such $p$, call it
$p_{\max}$, and then misses every rational in $(\sqrt2,\,p_{\max}]$, of which there are infinitely
many. This is sharper than \cref{ex:compactness_heine_borel_examples}\textbf{(a)}'s $\mathbb{N}$.
There, compactness fails because the set is unbounded. Here it fails on a closed \emph{and}
bounded set, purely because the ambient space $\mathbb{Q}$ is not $\mathbb{R}^n$. The
irrationality of $\sqrt2$ is doing all the work, leaving a \qt{hole} exactly where the cover would
need to close up.
\end{example}

% Kept inline: solved directly in Corsin Nick/Class Notes/Week 3.pdf, pp. 1--2.
\begin{exercisesolution}[Proof of \cref{ex:heine_borel_fails}]
\begin{enumerate}[label=\textbf{(\alph*)}]
  \item Non-negativity, definiteness and symmetry are inherited directly from the
    corresponding properties of $d$, since $t \mapsto \tfrac{t}{1+t}$ is non-negative on
    $[0,\infty)$ and vanishes only at $t=0$.

    \textit{Triangle inequality.} We want to show
    \[
    \frac{d(x,y)}{1+d(x,y)} \leq \frac{d(x,z)}{1+d(x,z)} + \frac{d(z,y)}{1+d(z,y)}.
    \]
    Denote $d(x,y) = d_1$, $d(x,z) = d_2$, $d(z,y) = d_3$ and multiply out:
    \[
    \begin{aligned}
    &d_1(1+d_2)(1+d_3) \leq d_2(1+d_1)(1+d_3) + d_3(1+d_1)(1+d_2) \\
    \iff\ & (d_1 + d_1d_2)(1+d_3) \leq (d_2 + d_2d_1)(1+d_3) + (d_3 + d_1d_3)(1+d_2) \\
    \iff\ & d_1 + d_1d_3 + d_1d_2d_3 + d_1d_2 \leq d_2 + d_3d_2 + d_2d_1d_3 + d_2d_1 \\
    & \qquad\qquad\qquad\qquad\qquad\quad + d_3 + d_3d_2 + d_1d_3 + d_1d_3d_2 \\
    \iff\ & d_1 \leq d_2 + 2d_2d_3 + d_3 + d_1d_2d_3 \\
    \Longleftarrow\ & 0 \leq 2d_2d_3 + d_1d_2d_3 \quad \text{together with } d_1 \leq d_2+d_3
      \text{ (the triangle inequality for } d\text{)}
    \end{aligned}
    \]
    which is true by non-negativity.
  \item Let $U \subseteq X$ be open. Then, given $x \in U$, we find $r > 0$ such that
    $B_r(x) := \{y \in X : d(x,y) < r\} \subseteq U$. Let $\varepsilon := \dfrac{r}{1+r}$. Then
    $\tilde B_\varepsilon(x) \subseteq B_r(x) \subseteq U$, where
    $\tilde B_\varepsilon(x) := \{y \in X : \tilde d(x,y) < \varepsilon\}$.
  \item Note that for all $x, y \in \mathbb{R}$, $\dfrac{|x-y|}{1+|x-y|} < 1$. So
    $\mathbb{N} \subseteq B_1(0)$, thus it is \textbf{bounded} with respect to this metric. Also,
    $\mathbb{N}$ is \textbf{closed}, as follows from \textbf{(b)}. But $x_n = n$ defines a
    sequence in $\mathbb{N}$ with no convergent subsequence, thus $\mathbb{N}$ is \textbf{not
    compact}.
\end{enumerate}
\end{exercisesolution}

\begin{ainote}
The cancellation in the derivation above originally contained a stray $d_2d_1d_2$ where the
expansion gives $d_2d_1d_3$; the cancelled terms and the conclusion $0 \leq 2d_2d_3 + d_1d_2d_3$
are unaffected. Re-derived cleanly above.
\end{ainote}

% Generator: Claude Opus 5 (Medium)
\begin{remark}[The moral: boundedness is not a topological property]
\label{rem:boundedness_not_topological}
It is worth saying out loud what \cref{ex:heine_borel_fails} has just demonstrated, because the
three parts are more striking together than separately.

Part \textbf{(b)} shows that $d$ and $\tilde d$ have the \emph{same open sets}. (Only one
direction is asked for, but the converse follows the same way: given
$\tilde B_\varepsilon(x)$, the choice $r := \tfrac{\varepsilon}{1-\varepsilon}$ gives
$B_r(x) \subseteq \tilde B_\varepsilon(x)$.) Two metrics with the same open sets have the same
convergent sequences, the same continuous functions, and, crucially, the same \textbf{compact}
sets, since compactness is defined purely by open covers (\cref{def:compact_metric_space}).

Part \textbf{(c)} shows that they do \emph{not} have the same bounded sets. Indeed
$\tilde d < 1$ always, so under $\tilde d$ \emph{every} subset of \emph{every} metric space is
bounded. The notion collapses entirely.

Putting those together: \textbf{compactness is topological, boundedness is not.} Boundedness is
a property of the chosen metric, not of the open sets it generates, so no theorem of the form
\qt{compact if and only if closed and bounded} can hold for a general metric. That is precisely why
\Cref{thm:heine_borel} carries the hypothesis \textbf{standard metric}. That hypothesis is not a
technicality to be skimmed; it is the entire content of the theorem.

% Reclassified from ainote to remark: it introduces terminology with \newterm and separates two
% notions by an explicit example, so it is mathematics.
% Feedback: Gemini 3.1 Pro (2026-08-11) --- singled out for placing total boundedness against
% boundedness at exactly the point where the reader has just watched boundedness collapse. Keep
% it nested here; moving it to the definition of total boundedness would lose the contrast.
\begin{remark}[What boundedness should have been]
The right general replacement is already available: compact if and only if complete \emph{and}
\newterm{totally bounded} (\cref{def:totally_bounded},
\cref{rem:compact_iff_complete_totally_bounded}). Total boundedness is the notion that boundedness
was trying and failing to be, and $\mathbb{N}$ under $\tilde d$ separates the two cleanly. For
$\varepsilon < \tfrac12$ the ball $\tilde B_\varepsilon(n)$ contains no integer other than $n$
itself, so no finite family of $\varepsilon$-balls covers $\mathbb{N}$. The set is therefore
bounded but not totally bounded, and that is exactly the gap through which compactness escapes.
\end{remark}
\end{remark}

% Source: old_exams/august2025.pdf (Prof. Laura Kobel-Keller), Analysis II examination,
% August 2025, Frage 10, printed p. 9
% Extractor: Claude Opus 5 (High)
\begin{exercise}[True or False \textnormal{(what actually characterises compactness)}]
\label{ex:kobel_compactness_characterisation}
Let $E \subseteq X$ be a subset of a metric space. For each of the following completions of the
statement \qt{$E$ is compact if \dots}, decide whether it is true or false. More than one may be
correct, and so may none.
\begin{enumerate}[label=\textbf{(\alph*)}]
  \item $E$ is closed and bounded.
  \item $E$ is simultaneously open and closed.
  \item $E$ is complete and totally bounded.
  \item $E$ is open and totally bounded.
\end{enumerate}
\exinfo{This exercise is taken from the Analysis II examination of August 2025
(Prof.\ Laura Kobel-Keller), Question 10, the opening question of the multiple-choice section.}
\end{exercise}

% Originally: content/week-03.tex, \subsection{Why do we care about compactness?}

\section{Why compactness matters}
\label{sec:why_compactness_matters}

% Sonnet 5 (Medium): added \label on each item so these three facts (heavily cited by name in
% later chapters, e.g. the extreme value theorem) can be \cref-ed instead of just named in prose.
\begin{enumerate}[label=\textbf{(\alph*)}]
  \item \label{item:continuity_preserves_compact}
    \textbf{Preserved by continuous functions:} $X$, $Y$ metric spaces, $K \subseteq X$
    compact, $f : X \to Y$ continuous. Then $f(K) \subseteq Y$ is compact.
  \item \label{item:extreme_value_theorem}
    \textbf{Weierstrass theorem / Extreme value theorem:} $X$ metric space,
    $K \subseteq X$ compact, $f : K \to \mathbb{R}$ continuous. Then $f$ attains its maximum and
    minimum in $K$. In particular, $f$ is bounded. This is proved as
    \cref{lem:sequential_compactness_extreme_values}.
    \begin{ainote}
    The official script's statement of this theorem (Cor.\ 9.79) writes the conclusion as
    \qt{there exist $\bar x \in X$ such that $f(\bar x) = \sup_K f$}, with $\bar x$ ranging over
    the whole ambient space $X$ rather than over $K$; the proof given there repeats the slip,
    writing $f^{-1}(\sup f(X))$ where it must mean $f^{-1}(\sup f(K))$. As stated, the claim is
    false as soon as $f$ is unbounded off $K$: e.g.\ $K:=\{0\}\subset X:=\mathbb{R}$,
    $f(x):=x$. The maximiser $\bar x$ must be sought in $K$, as in the statement above.
    \end{ainote}
  \item \label{item:uniform_continuity_compact}
    \textbf{Uniform continuity:} $X$, $Y$ metric spaces, $f : X \to Y$ continuous. If $X$ is
    compact, then $f$ is uniformly continuous (\cref{def:uniformly_continuous}). This is proved
    as \cref{prop:compact_implies_uniformly_continuous}, via the Lebesgue number of the cover
    $\{f^{-1}(B_{\varepsilon/2}(y))\}_{y}$.
\end{enumerate}

% Sonnet 5 (Medium): illustrates item (b), the extreme value theorem.
\begin{center}
\begin{tikzpicture}[>=Latex, scale=1.2]
  \draw[->] (-0.3,0) -- (3.5,0) node[right] {$x$};
  \draw[->] (0,-0.3) -- (0,2.8) node[above] {$y$};

  \draw[green!60!black, thick] plot[smooth, tension=0.8] coordinates {(0.3,0.8) (0.9,2.3) (2.1,0.4) (3.0,1.8)};

  \draw (0.3,0.05) -- (0.3,-0.05) node[below] {$a$};
  \draw (0.9,0.05) -- (0.9,-0.05) node[below] {$c$};
  \draw (2.1,0.05) -- (2.1,-0.05) node[below] {$d$};
  \draw (3.0,0.05) -- (3.0,-0.05) node[below] {$b$};

  \draw[dashed, red!80!black] (0,2.3) node[left, font=\small] {$f(c)$} -- (0.9,2.3) -- (0.9,0);
  \fill[red!80!black] (0.9,2.3) circle (2pt);

  \draw[dashed, blue!80!black] (0,0.4) node[left, font=\small] {$f(d)$} -- (2.1,0.4) -- (2.1,0);
  \fill[blue!80!black] (2.1,0.4) circle (2pt);

  \node[below, font=\small, align=center] at (1.7,-0.55)
    {\textbf{(b)} on the compact $K=[a,b]$, $f$ attains its maximum at $c$ and its minimum at $d$};
\end{tikzpicture}
\end{center}

Both extrema are attained \emph{inside} $K$. That is the content of \textbf{(b)}, and it is
exactly what fails on a non-compact domain: on the open interval $(a,b)$ the same $f$ may
approach a supremum it never reaches.

% Quelle: Diego Torres Tejeda/Notes - 02.03 - Diego Torres Tejeda.pdf, pp. 8--10
\begin{remark}[What compactness and connectedness are \emph{for}: local to global]
\label{rem:local_to_global}
The three facts above look like a list of unrelated services compactness happens to provide.
Diego Torres Tejeda's notes give the idea that unifies them, and it is worth having early
because it also explains why connectedness sits in the same chapter.

\textbf{Local-to-global arguments} are a recurring theme across mathematics. It is often easy to
check a property \emph{locally}, on some ball $B_\varepsilon(x)$ around each point, while what one
actually wants is the property \emph{globally}, on all of $X$. Compactness and
connectedness are the two hypotheses that license that jump. Schematically, a local-to-global
statement reads:
\begin{quote}
\itshape
If for every $x \in X$ there is $\varepsilon>0$ such that $B_\varepsilon(x)$ has property $P$,
then $X$ has property $P$.
\end{quote}
Each of the two hypotheses powers its own version, and the proofs are three lines each.

\textbf{Connectedness turns \qt{locally constant} into \qt{constant}.} Let $(X,d)$ be connected
and $f : X \to \mathbb{R}$ locally constant. Every level set $f^{-1}(a)$ is then open. Fix
$a \in f(X)$ and split
\[X = \underbrace{f^{-1}(a)}_{=:U} \ \sqcup \underbrace{\bigcup_{b \neq a}f^{-1}(b)}_{=:V},\]
a disjoint union of two open sets. Since $U \neq \emptyset$, connectedness forces $V=\emptyset$,
i.e.\ $f \equiv a$.

\textbf{Compactness turns \qt{locally bounded} into \qt{bounded}.} Let $X$ be compact and
suppose each $x$ has an open $U_x \ni x$ on which $f$ is bounded. The $\{U_x\}_{x\in X}$ cover
$X$, so finitely many suffice, say $U_1,\dots,U_n$, and then
\[\sup_{x \in X}|f(x)| \;=\; \max_{1\leq i\leq n}\ \sup_{x \in U_i}|f(x)| \;<\; \infty,\]
because a \emph{maximum of finitely many} finite numbers is finite. A continuous $f$ is locally
bounded, since $|f|<|f(x)|+1$ near each $x$, so this recovers the boundedness half of
\cref{item:extreme_value_theorem}.

Read this way the two proofs are the same proof twice, with different notions of \qt{finitely
many}: connectedness says a partition into two open pieces must be trivial, compactness says a
cover can be taken finite. Both convert local information into global. Every application in
\Cref{item:continuity_preserves_compact,item:extreme_value_theorem,item:uniform_continuity_compact}
is an instance. Uniform continuity is the most visible of them, since there the whole difficulty
is that the $\delta$ supplied by continuity is local and one needs a single global one
(\cref{rem:continuity_vs_uniform_continuity}).
\end{remark}

% The exercise below overlaps heavily with ex:topological_properties_table above. Both are Sascha
% Brack's, from two different files (his Week 2 Friday and Week 3 Monday notes), and he evidently
% reused the self-test. Both are kept: deciding that one is redundant would be a judgement about
% his notes rather than about this transcription.
% Transition: Claude Opus 5 (High)
The exercise below covers the same ground as \cref{ex:topological_properties_table}, and is worth
doing anyway. It adds the explicit \textbf{?} answers for the cases where the data is
insufficient, a trap in the final item involving two \emph{parallel} spanning vectors, and a
four-column solution table which is where the lesson actually lands. Read
\cref{ex:topological_properties_table} as the warm-up and this one as the real test.

% Supplement: Sascha Brack/Class Notes/Week_02_Notes_Friday_Updated.pdf, p. 8
\begin{exercise}[Classify these sets]
\label{ex:classification_table}
Sascha Brack sets this as a self-test at the end of his compactness lecture, and it is the
fastest way to find out whether the four notions have actually separated in your head. For each
subset of $\mathbb{R}^n$ (standard metric), decide whether it is \textbf{open}, \textbf{closed},
\textbf{compact}, \textbf{complete}. Throughout, $f$ denotes a continuous function and
\textbf{?} means \qt{not decidable from the information given}.
\begin{enumerate}[label=\textbf{\arabic*.}]
  \item $[0,1] \subseteq \mathbb{R}$
  \item $\mathbb{Q} \subseteq \mathbb{R}$
  \item $B_1(0) \subseteq \mathbb{R}^3$
  \item $\left\{\tfrac1n : n \geq 1\right\} \subseteq \mathbb{R}$, and separately
    $\{0\} \cup \left\{\tfrac1n : n \geq 1\right\}$
  \item $\bigcup_{n\geq1}B_{1/n}(n) \subseteq \mathbb{R}$
  \item $(-\infty,1] \subseteq \mathbb{R}$
  \item $f([0,1])$, \quad $f^{-1}([0,1])$, \quad $f^{-1}\big((0,1)\big)$
  \item $S := \operatorname{span}\!\left\{\begin{pmatrix}1\\0\end{pmatrix},
    \begin{pmatrix}3\\0\end{pmatrix}\right\} \subseteq \mathbb{R}^2$
\end{enumerate}
\end{exercise}

% Quelle: Diego Torres Tejeda/Notes - 16.03 - Diego Torres Tejeda.pdf, p. 1 -- Sonnet 5 (Medium)
\begin{remark}[What continuous maps do and do not preserve]
Diego's exercise-class notes collect \cref{item:continuity_preserves_compact} and the analogous
fact for connectedness (\cref{thm:continuity_preserves_connected}) into one table, together with
the dual statement for preimages, and a caution worth stating explicitly:
\begin{center}
\begin{tabular}{lll}
& \textbf{Preserved by $f^{-1}$} & \textbf{Preserved by $f$} \\
\hline
open & yes (\cref{def:continuous}) & not in general \\
closed & yes & not in general \\
compact & not in general & yes (\cref{item:continuity_preserves_compact}) \\
connected & not in general & yes (\cref{thm:continuity_preserves_connected}) \\
\end{tabular}
\end{center}
\textbf{Caution:} completeness and boundedness are \emph{not} topological properties, so in
general they are preserved neither by continuous maps nor by their inverses. For instance,
$\tan : (-\tfrac\pi2,\tfrac\pi2) \to \mathbb{R}$ is a continuous bijection carrying a bounded set
onto an unbounded one.
\end{remark}


% --- exercises relocated here from the dissolved sheet 3 ---

% Originally: content/exercise-sheets/week-03.tex, problem 3.3
% Source: exercises/Ex3_Analysis2_eng.pdf

\begin{exercise}[Open, closed, complete and compact]
\label{ex:3.3}
For each of the following subsets of $\mathbb{R}^N$ say whether they are open / closed / none /
compact (with respect to the standard Euclidean structure). Try to prove your assertions
``efficiently''.
\begin{enumerate}[label=\textbf{\arabic*.}]
  \item $E_1 := \{x \in \mathbb{R}^3 : 0 < x_2 \leq 2\}$
  \item $E_2 := \{x \in \mathbb{R}^n : \sin(|x|) \geq \tfrac{1}{4}\}$
  \item $E_3 := \bigcup_{n \geq 1}\{x \in \mathbb{R}^3 : x_1^4 - \tfrac{1}{n}x_3^2 - x_2^2 > \tfrac{1}{n},\ x_1 + x_2 < 6\}$
  \item $E_4 := \{(x,y) \in (\mathbb{R}^n \times \mathbb{R}^n) : x \cdot y > \tfrac{1}{2}|x||y|\}$
  \item $E_5 := \{X \in \mathbb{R}^{n\times n} : \text{the matrix } X \text{ is invertible}\}$
  \item $E_6 := \{X \in \mathbb{R}^{n\times n} : \text{the matrix } X \text{ is symmetric}\}$
  \item $E_7 := \{X \in \mathbb{R}^{n\times n} : \text{the entries of } X \text{ are either } 0 \text{ or } 1\}$
  \item $E_8 := \{x \in \mathbb{R}^n : |x_i| \leq 6 \text{ for all } 1 \leq i \leq n\}$
  \item $E_9 := E_8 \times E_3 \subset \mathbb{R}^{n+3}$
  \item $E_{10} := E_2 \cap E_8 \subset \mathbb{R}^n$
\end{enumerate}
\exinfo{This exercise is Problem 3.3 of Problem Sheet 3, Analysis II, Spring Semester 2026.}
\end{exercise}

% Originally: content/exercise-sheets/week-03.tex, problem 3.4
% Source: exercises/Ex3_Analysis2_eng.pdf

\begin{exercise}[Multiple choice (open, complete, compact)]
\label{ex:3.4}
Select all the statements below that are true.
\begin{enumerate}[label=\textbf{(\alph*)}]
  \item A nonempty open strict subset of $\mathbb{R}^n$ cannot be compact.
  \item A nonempty open strict subset of $\mathbb{R}^n$ cannot be complete.
  \item A complete subset of $\mathbb{R}^n$ contains all its accumulation points.
  \item Countable union of complete subsets of $\mathbb{R}^n$ is complete.
  \item A subset is closed in $\mathbb{R}^n$ if and only if it is complete as a metric space
    itself (with the distance inherited from $\mathbb{R}^n$).
\end{enumerate}
\exinfo{This exercise is Problem 3.4 of Problem Sheet 3, Analysis II, Spring Semester 2026.}
\end{exercise}

% Source: old_exams/examFS24.pdf (21 August 2024), Wahr/Falsch-Fragen, Aufgabe 1, p. 2
% Extractor: Gemini 3.6 Flash (High)
% Correction: Claude Opus 5 (High) --- \subset -> \subseteq, and part (c) added back: the original
% also asks whether U is *simply* connected, which is the one item on that true/false block with a
% trap in it, and Flash dropped it. Citation corrected: the source is the true/false section, not a
% "Part A".
\begin{exercise}[Continuous image of an annulus]
\label{ex:annulus_continuous_image}
Consider the closed annulus $U := \{(x,y) \in \mathbb{R}^2 \mid 1 \leq x^2+y^2 \leq 4\}$ and the continuous function $f : U \to \mathbb{R}$, $f(x,y) := \sin(xy) - y^4$.
\begin{enumerate}[label=\textbf{(\alph*)}]
  \item Determine whether $U$ is compact and whether $U$ is connected.
  \item Prove that the image $f(U) \subseteq \mathbb{R}$ is a compact, connected interval $[a,b]$ for some real numbers $a < b$.
  \item Is $U$ simply connected?
\end{enumerate}
\exinfo{This exercise is taken from the exam of 21 August 2024, where it appears as the first block
of true/false questions. It is reformulated here as a proof exercise.}
\end{exercise}

% Source: old_exams/fs2023/Prüfung.pdf (9 August 2021), Teil B, Aufgabe 3, p. 15
% Extractor: Claude Opus 5 (High)
\begin{exercise}[Nearest point to the origin]
\label{ex:nearest_point_to_origin}
Let $A \subseteq \mathbb{R}^n$ be non-empty and let $f : A \to \mathbb{R}_{\ge 0}$, $f(x) :=
\lvert x\rvert$, be the restriction of the Euclidean norm to $A$. Prove the following.
\begin{enumerate}[label=\textbf{(\alph*)}]
  \item If $A$ is closed, then $f$ attains a minimum on $A$: there is $P_0 \in A$ with
    $f(P_0) \le f(P)$ for all $P \in A$.
  \item Statement \textbf{(a)} is false in general when $A$ is not assumed closed.
  \item If $A$ is convex, then the minimum is attained at no more than one point of $A$.
  \item If $0 \notin A$, then every minimiser $P_0 \in A$ lies in $\partial A$.
\end{enumerate}
\exinfo{This exercise is taken from the exam of 9 August 2021, Part B, Problem 3.}
\end{exercise}


% Originally: new section (no predecessor in the week-based skeleton).
% Quelle: Linus Lüchinger/Slides/Slides-05-28.pdf, slides 2--3.

\section{Compactness in spaces of functions: Ascoli--Arzelà}
\label{sec:ascoli_arzela}

% Supplement: Linus Lüchinger/Slides/Slides-05-28.pdf, slides 2--3
\begin{ainote}
Corsin's notes stop at \cref{thm:heine_borel}, which settles compactness in $\mathbb{R}^n$ and
says nothing about spaces of \emph{functions}. Linus devotes his final session to the missing
theorem (Theorem 15.47 of the lecture notes), presenting it as \qt{the Bolzano--Weierstrass
theorem, one level up}. It is worth having here for a concrete reason: it is the compactness
input that \cref{rem:peano_vs_picard} appeals to without naming, when it says Peano's theorem
extracts a convergent subsequence from a family of approximate solutions.
\end{ainote}

\subsection{The space of continuous functions on a compact set}

\begin{definition}[Supremum norm on $C(K,\mathbb{R}^m)$]
\label{def:sup_norm_function_space}
Let $K \subseteq \mathbb{R}^n$ be compact. Write $C(K,\mathbb{R}^m)$ for the set of continuous
maps $f : K \to \mathbb{R}^m$, a real vector space under pointwise operations, and equip it with
the \newterm{supremum norm} \germanfor{Supremumsnorm}
\[\lVert f\rVert_\infty := \sup_{x \in K} \lvert f(x)\rvert.\]
\end{definition}

The supremum is finite because $\lvert f\rvert$ is continuous on a compact set, hence bounded
(\cref{sec:why_compactness_matters}), which is the same point already made for the supremum metric
in \cref{ex:sup_metric_is_metric}. Convergence in this norm is exactly \emph{uniform} convergence.

\begin{proposition}[$C(K,\mathbb{R}^m)$ is complete]
\label{prop:function_space_complete}
For $K \subseteq \mathbb{R}^n$ compact, $\big(C(K,\mathbb{R}^m), \lVert\cdot\rVert_\infty\big)$
is a complete normed vector space.
\end{proposition}

\begin{proof}
Let $(f_k)_{k\in\mathbb{N}}$ be Cauchy for $\lVert\cdot\rVert_\infty$. For each fixed $x \in K$
we have $\lvert f_k(x)-f_\ell(x)\rvert \leq \lVert f_k-f_\ell\rVert_\infty$, so $(f_k(x))_k$ is
Cauchy in $\mathbb{R}^m$ and converges, by completeness of $\mathbb{R}^m$, to a limit we call
$f(x)$. This defines a function $f : K \to \mathbb{R}^m$.

The convergence is uniform. Given $\varepsilon>0$ choose $N$ with
$\lVert f_k-f_\ell\rVert_\infty < \varepsilon$ for $k,\ell \geq N$; fixing $k \geq N$ and letting
$\ell \to \infty$ in $\lvert f_k(x)-f_\ell(x)\rvert < \varepsilon$ gives
$\lvert f_k(x)-f(x)\rvert \leq \varepsilon$ for \emph{every} $x$, i.e.\
$\lVert f_k - f\rVert_\infty \leq \varepsilon$.

Finally $f$ is continuous, being a uniform limit of continuous functions, by the standard
$\varepsilon/3$ argument from \emph{Analysis I}. Hence $f \in C(K,\mathbb{R}^m)$ and
$f_k \to f$ in the norm.
\end{proof}

\subsection{Why Bolzano--Weierstrass fails one level up}

In $\mathbb{R}^n$, bounded sequences have convergent subsequences. The naive transfer of that
statement to $C(K,\mathbb{R}^m)$ is \textbf{false}, and it is worth seeing exactly how.

% Generator: Claude Opus 5 (High)
\begin{aiexample}[A bounded sequence of functions with no convergent subsequence]
\label{ex:bounded_no_convergent_subsequence}
Take $K := [0,1]$, $m := 1$ and $f_k(x) := x^k$. Every $f_k$ is continuous with
$\lVert f_k\rVert_\infty = 1$, so the sequence is bounded, and indeed lies on the unit sphere.

Yet no subsequence converges in $C([0,1])$. Pointwise,
\[f_k(x) \longrightarrow g(x) := \begin{cases} 0, & 0 \leq x < 1, \\ 1, & x = 1,\end{cases}\]
and the same holds for every subsequence. If some $f_{k_j}$ converged uniformly to a limit
$f \in C([0,1])$, then it would converge pointwise to $f$ as well, forcing $f = g$; but $g$ is
not continuous. So no subsequence can converge.

\begin{remark}[The diagnosis]
Boundedness was never the issue. What goes wrong is that the $f_k$ become arbitrarily
\emph{steep} near $x=1$: to keep $\lvert f_k(x)-f_k(y)\rvert$ below a fixed $\varepsilon$ there,
one needs $\lvert x-y\rvert$ smaller and smaller as $k$ grows. Each individual $f_k$ is
uniformly continuous (\cref{def:uniformly_continuous}), being continuous on a compact interval.
What fails is that no single $\delta$ works for all of them at once. Ruling that out is precisely
what the next definition does.
\end{remark}
\end{aiexample}

\subsection{Equicontinuity, and the theorem}

\begin{definition}[Equicontinuity]
\label{def:equicontinuous}
A family $(f_k)_{k\in\mathbb{N}}$ in $C(K,\mathbb{R}^m)$ is \newterm{equicontinuous}
\germanfor{gleichgradig stetig} if
\[\forall \varepsilon>0\ \ \exists \delta>0\ \ \forall k \in \mathbb{N}\ \ \forall x,y \in K :
  \quad \lvert x-y\rvert < \delta \implies \lvert f_k(x)-f_k(y)\rvert < \varepsilon.\]
\end{definition}

% Feedback: Gemini 3.1 Pro (2026-08-11) --- called didactically outstanding: the three-way
% quantifier comparison packs the chapter's most abstract notion into three bullets. Keep the
% three-item list and the closing sentence that names the gap.
\begin{remark}[Reading the quantifiers]
\label{rem:equicontinuity_quantifiers}
Compare the three conditions, which differ only in what $\delta$ is allowed to depend on:
\begin{itemize}
  \item \textbf{Continuity} of one $f$: $\delta$ may depend on $\varepsilon$ \emph{and on the
    point} $x$.
  \item \textbf{Uniform continuity} of one $f$: $\delta$ may depend on $\varepsilon$ only, so
    there is one $\delta$ for all points.
  \item \textbf{Equicontinuity} of a family: $\delta$ may depend on $\varepsilon$ only, so there
    is one $\delta$ for all points \emph{and all members of the family}.
\end{itemize}
So equicontinuity is uniform continuity made uniform in one further variable, the index. Each
$x^k$ of \cref{ex:bounded_no_convergent_subsequence} is uniformly continuous, and the family is
not equicontinuous; that is the whole gap.
\end{remark}

% Quelle: Linus Lüchinger/Slides/Slides-05-28.pdf, slide 3
\begin{theorem}[Ascoli--Arzelà]
\label{thm:ascoli_arzela}
Let $K \subseteq \mathbb{R}^n$ be compact and let $(f_k)_{k\in\mathbb{N}}$ be a sequence in
$C(K,\mathbb{R}^m)$ that is
\begin{enumerate}[label=\textbf{(\alph*)}]
  \item \textbf{bounded}: $\sup_k \lVert f_k\rVert_\infty < \infty$; and
  \item \textbf{equicontinuous} in the sense of \cref{def:equicontinuous}.
\end{enumerate}
Then $(f_k)$ has a subsequence converging uniformly to some $f \in C(K,\mathbb{R}^m)$.
\end{theorem}

% Linus states the theorem without proof, remarking that it is "mostly just used in proofs".
\begin{ainote}
None of the transcribed notes prove this theorem; Linus states it and remarks that it is
\qt{mostly just used in proofs}. The sketch below is included because its mechanism, a
diagonal argument, explains why \emph{both} hypotheses are needed, which the statement alone does
not. It is a sketch and not a proof: the details are omitted.
\end{ainote}

\begin{remark}[Proof sketch: the diagonal argument]
\label{rem:ascoli_diagonal_sketch}
Fix a countable dense subset $D = \{q_1,q_2,\dots\} \subseteq K$ (the points of $K$ with rational
coordinates will do).

\textbf{Step 1: converge on $D$, using boundedness.} The values $(f_k(q_1))_k$ form a bounded
sequence in $\mathbb{R}^m$, so by Bolzano--Weierstrass some subsequence converges at $q_1$. Pass
to a further subsequence converging also at $q_2$, then at $q_3$, and so on. Each refinement is
legitimate, but no single one of them handles all of $D$. Taking the \emph{diagonal} sequence,
whose $j$-th member is the $j$-th member of the $j$-th subsequence, produces one subsequence
converging at every point of $D$ simultaneously.

\textbf{Step 2: upgrade to uniform convergence, using equicontinuity.} Convergence on a dense
set says nothing on its own about the rest of $K$; a sequence can converge at every rational and
oscillate wildly in between. Equicontinuity forbids that: given $\varepsilon$, pick the $\delta$
it supplies, cover the compact $K$ by finitely many balls of radius $\delta$ centred at points
of $D$, and control any $x$ by the nearby dense point. This makes the diagonal subsequence
uniformly Cauchy, hence convergent by \cref{prop:function_space_complete}.
\end{remark}

\begin{remark}[Where each hypothesis does its work]
The two hypotheses are not interchangeable, and the sketch shows which does what:
\textbf{boundedness} is what lets Bolzano--Weierstrass run at each individual point, and
\textbf{equicontinuity} is what converts pointwise control on a dense set into uniform control
everywhere. Delete the second and \cref{ex:bounded_no_convergent_subsequence} is a
counterexample; delete the first and the constant sequences $f_k \equiv k$ are equicontinuous
(they are all constant, so any $\delta$ works) with no convergent subsequence at all.
\end{remark}

% Generator: Claude Opus 5 (High)
\begin{aiexercise}[A uniform derivative bound gives compactness]
\label{ex:uniform_lipschitz_family}
Let $(f_k)_{k\in\mathbb{N}}$ be a sequence in $C^1([0,1],\mathbb{R})$ with
\[\lvert f_k'(x)\rvert \leq 1 \quad \text{for all } x \in [0,1] \text{ and all } k,
  \qquad\text{and}\qquad f_k(0) = 0 \text{ for all } k.\]
\begin{enumerate}[label=\textbf{(\alph*)}]
  \item Show that the family is equicontinuous.
  \item Show that it is bounded in $\lVert\cdot\rVert_\infty$.
  \item Conclude that some subsequence converges uniformly, and explain why the argument would
    break down if the bound $\lvert f_k'\rvert \leq 1$ were replaced by \qt{each $f_k$ is
    differentiable}.
\end{enumerate}
\exsol{sol:uniform_lipschitz_family}
\end{aiexercise}

% Transition: Claude Opus 5 (High)
That closes the account of compactness. What follows returns to $\mathbb{R}^n$ and to
differentiation, where compactness reappears as a hypothesis rather than a subject.

% Originally: the \section{Solutions} of content/week-03.tex

\section{Solutions}
\label{sec:compactness_solutions}

\begin{exercisesolution}[Solution to \cref{ex:classification_table}]
\label{sol:classification_table}
% Generator: Claude Opus 5 (Medium)
\begin{center}
\begin{tabular}{@{}l cccc@{}}
\textbf{Set} & \textbf{open} & \textbf{closed} & \textbf{compact} & \textbf{complete} \\
\hline
$[0,1]$ & $\times$ & $\checkmark$ & $\checkmark$ & $\checkmark$ \\
$\mathbb{Q}$ & $\times$ & $\times$ & $\times$ & $\times$ \\
$B_1(0) \subseteq \mathbb{R}^3$ & $\checkmark$ & $\times$ & $\times$ & $\times$ \\
$\left\{\tfrac1n\right\}$ & $\times$ & $\times$ & $\times$ & $\times$ \\
$\{0\}\cup\left\{\tfrac1n\right\}$ & $\times$ & $\checkmark$ & $\checkmark$ & $\checkmark$ \\
$\bigcup_{n\geq1}B_{1/n}(n)$ & $\checkmark$ & $\times$ & $\times$ & $\times$ \\
$(-\infty,1]$ & $\times$ & $\checkmark$ & $\times$ & $\checkmark$ \\
$f([0,1])$ & $\times$ & $\checkmark$ & $\checkmark$ & $\checkmark$ \\
$f^{-1}([0,1])$ & \textbf{?} & $\checkmark$ & \textbf{?} & $\checkmark$ \\
$f^{-1}\big((0,1)\big)$ & $\checkmark$ & \textbf{?} & \textbf{?} & \textbf{?} \\
$S = \operatorname{span}\{(1,0),(3,0)\}$ & $\times$ & $\checkmark$ & $\times$ & $\checkmark$ \\
\end{tabular}
\end{center}

The rows that carry the lesson:

\textbf{$\left\{\tfrac1n\right\}$ versus $\{0\}\cup\left\{\tfrac1n\right\}$.} Adding the single
point $0$ changes all four answers. Without it the set misses its own limit point, so it is not
closed, and the Cauchy sequence $\tfrac1n$ has nowhere to converge inside the set. With it, the
set is closed and bounded, hence compact by \cref{thm:heine_borel}.

\textbf{$\bigcup_{n\geq1}B_{1/n}(n)$.} Open as a union of open sets
(\cref{item:arbitrary_union_open}); not closed, since each endpoint $n\pm\tfrac1n$ is a limit
point outside the set; unbounded, so certainly not compact.

\textbf{$S$ is a trap.} The two spanning vectors are parallel, so $S$ is not $\mathbb{R}^2$ but
the $x$-axis, a line, which is closed, unbounded and not compact.

\textbf{The $f$ rows show exactly how far continuity gets you.} Preimages behave: $f^{-1}$ of a
closed set is closed, of an open set is open (\cref{def:continuous}). Images do not, and only
compactness transfers forward (\cref{item:continuity_preserves_compact}), which is why
$f([0,1])$ is compact but $f^{-1}([0,1])$ need not be: take $f \equiv 0$, giving
$f^{-1}([0,1]) = \mathbb{R}^n$. That same $f$ makes $f^{-1}([0,1])$ open, while $f = \id$ makes
it $[0,1]$, which is not, hence the \textbf{?}. And $f([0,1])$ is never open: it is compact,
hence closed and bounded, and a non-empty proper subset of $\mathbb{R}^m$ cannot be both open
and closed, since $\mathbb{R}^m$ is connected (\cref{def:connectedness}).

% Reclassified from ainote to remark: an observation about the table's mathematical content.
\begin{remark}[Why the last two columns agree]
The punchline of the table is that its last two columns are \textbf{identical}. That is not a
coincidence of this particular list. Inside a \emph{complete} ambient space a subset is complete
exactly when it is closed (\cref{rem:complete_vs_closed}), and $\mathbb{R}^n$ is complete. Build
the same table inside $\mathbb{Q}$ instead and the two columns come apart immediately.
\end{remark}
\end{exercisesolution}

% Transition: Claude Opus 5 (High)
% Correction: Claude Opus 5 (High) --- replaces an ainote that was wrong three ways. It named
% ex:connectedness_tf as being "in this chapter" when that exercise is in ch. 8; it claimed only
% aiexercise problems are solved in this section, though ex:open_closed_complete_compact and ex:classification_table are
% here and are neither; and it listed two of the six exercises actually solved inline.
Not every solution in this chapter is collected here. \Cref{ex:totally_bounded_implies_bounded}
(\cpageref{ex:totally_bounded_implies_bounded}), \cref{ex:compact_finite_complete}
(\cpageref{ex:compact_finite_complete}), \cref{ex:convergent_sequence_compact}
(\cpageref{ex:convergent_sequence_compact}), \cref{ex:finite_set_compact}
(\cpageref{ex:finite_set_compact}), \cref{ex:topological_properties_table}
(\cpageref{ex:topological_properties_table}) and \cref{ex:heine_borel_fails}
(\cpageref{ex:heine_borel_fails}) are solved where they appear.

\begin{exercisesolution}[Solution to \cref{ex:open_closed_complete_compact}]
\label{sol:open_closed_complete_compact}
% Generator: Claude Opus 5 (High)
The efficient route is to write each set as a preimage of an open or closed set under a
continuous map, and then read openness or closedness off \cref{def:continuous}\textbf{(c)}
rather than chasing balls. Compactness is then Heine--Borel (\cref{thm:heine_borel}): closed
\emph{and} bounded.
\begin{center}
\begin{tabular}{@{}c l c c c@{}}
& \textbf{Set} & \textbf{open} & \textbf{closed} & \textbf{compact} \\
\hline
1 & $E_1 = \{0 < x_2 \leq 2\} \subseteq \mathbb{R}^3$ & $\times$ & $\times$ & $\times$ \\
2 & $E_2 = \{\sin|x| \geq \tfrac14\}$ & $\times$ & $\checkmark$ & $\times$ \\
3 & $E_3$ (union of open sets) & $\checkmark$ & $\times$ & $\times$ \\
4 & $E_4 = \{x\cdot y > \tfrac12|x||y|\}$ & $\checkmark$ & $\times$ & $\times$ \\
5 & $E_5 = \{X \text{ invertible}\}$ & $\checkmark$ & $\times$ & $\times$ \\
6 & $E_6 = \{X \text{ symmetric}\}$ & $\times$ & $\checkmark$ & $\times$ \\
7 & $E_7 = \{X \text{ with entries } 0,1\}$ & $\times$ & $\checkmark$ & $\checkmark$ \\
8 & $E_8 = [-6,6]^n$ & $\times$ & $\checkmark$ & $\checkmark$ \\
9 & $E_9 = E_8 \times E_3$ & $\times$ & $\times$ & $\times$ \\
10 & $E_{10} = E_2 \cap E_8$ & $\times$ & $\checkmark$ & $\checkmark$ \\
\end{tabular}
\end{center}

\begin{enumerate}[label=\textbf{(\alph*)}]
  \item Half-open in the $x_2$ direction, so \textbf{neither}: no ball around a point with
    $x_2 = 2$ stays inside, and the sequence $(0,\tfrac1k,0) \to 0 \notin E_1$ shows it is not
    closed. Unbounded, hence not compact.
  \item \textbf{Closed:} $E_2 = h^{-1}\big([\tfrac14,\infty)\big)$ for the continuous
    $h(x) := \sin(|x|)$. Not open, since a point with $\sin|x| = \tfrac14$
    has points of $E_2^c$ arbitrarily close. Not compact: $\sin$ is periodic, so $E_2$ contains
    infinitely many concentric annuli and is unbounded.
  \item \textbf{Open:} each set in the union is defined by two \emph{strict} inequalities between
    continuous functions, hence open, and an arbitrary union of open sets is open
    (\cref{item:arbitrary_union_open}). It is non-empty (take $x = (2,0,0)$ and $n=1$) and
    unbounded (take $x = (-M,0,0)$ for large $M$), so it is not compact; and a non-empty proper
    open subset of the connected space $\mathbb{R}^3$ cannot also be closed
    (\cref{def:connectedness}).
  \item \textbf{Open:} $E_4 = F^{-1}\big((0,\infty)\big)$ for the continuous
    $F(x,y) := x\cdot y - \tfrac12|x||y|$. Unbounded (scale any element),
    hence not compact, and not closed by the connectedness argument of part \textbf{(c)}.
    Geometrically $E_4$ is the set of pairs enclosing an angle strictly less than $60^\circ$
    (\cref{rem:inner_product_angle_geometry}).
  \item \textbf{Open:} $E_5 = \det^{-1}\big(\mathbb{R}\setminus\{0\}\big)$ and $\det$ is a
    polynomial in the entries, hence continuous. Unbounded, so not compact.
  \item \textbf{Closed:} $E_6$ is the solution set of the linear equations $X_{ij} - X_{ji} = 0$,
    i.e.\ the preimage of the closed set $\{0\}$ under a continuous (indeed linear) map. Being a
    linear subspace it is unbounded, hence not compact, and for $n \geq 2$ it is a proper
    subspace, hence has empty interior and is not open.
  \item \textbf{Closed and compact:} $E_7$ is \emph{finite}, with $2^{n^2}$ elements. Finite sets
    are compact (\cref{ex:compact_finite_complete}\textbf{(a)}) and closed. Not open, since no
    ball around a matrix of $0$s and $1$s consists only of such matrices.
  \item \textbf{Closed and compact:} $E_8 = [-6,6]^n$ is closed (a finite intersection of the
    closed sets $\{|x_i| \leq 6\}$) and bounded, so \cref{thm:heine_borel} applies. Not open:
    no ball around a point with some $|x_i| = 6$ fits inside.
  \item \textbf{Neither, and not compact.} For non-empty $A, B$, the product $A \times B$ is open
    exactly when both factors are, closed exactly when both are, and compact exactly when both
    are. Here $E_8$ is not open and $E_3$ is neither closed nor compact, so $E_9$ fails all
    three.

    \begin{ainote}
    The official solution opens this part with \qt{Since $E_8$ is unbounded so is $E_9$}. That
    is a slip, since $E_8 = [-6,6]^n$ is bounded. It is $E_3$ that is unbounded, as the official
    solution itself established two parts earlier. The conclusion (that $E_9$ is not compact) is
    correct; only the reason is misattributed.
    \end{ainote}
  \item \textbf{Closed and compact:} an intersection of the closed set $E_2$ with the closed
    set $E_8$ is closed (\cref{prop:open_closed_operations}\textbf{(c)}), and it is bounded
    because $E_8$ is. Heine--Borel again. Alternatively, quote
    \cref{prop:closed_subset_of_compact}: a closed subset of the compact $E_8$ is compact.
\end{enumerate}

\begin{remark}[The efficient method, extracted]
Nine of the ten sets were settled without producing a single $\varepsilon$. The recipe is:
\emph{strict} inequalities between continuous functions give open sets, \emph{weak} ones give
closed sets, and mixing the two gives neither. Compactness then reduces to checking boundedness,
by Heine--Borel. This is what the problem means by proving assertions \qt{efficiently}, and it is
the practical payoff of \cref{def:continuous}\textbf{(c)} over the $\varepsilon$--$\delta$
formulation.
\end{remark}
\end{exercisesolution}

\begin{exercisesolution}[Solution to \cref{ex:open_complete_compact_multiple_choice}]
\label{sol:open_complete_compact_multiple_choice}
% Generator: Claude Opus 5 (High)
\textbf{(a), (b), (c) and (e) are true; (d) is false.}

The engine behind \textbf{(a)}, \textbf{(b)}, \textbf{(c)} and \textbf{(e)} is a single fact,
\cref{rem:complete_vs_closed}: inside a \emph{complete} ambient space, a subset is complete if and
only if it is closed. Here the ambient space is $\mathbb{R}^n$, which is complete by
\cref{prop:Rn_coordinatewise_complete}.

\textbf{(a) True.} A compact subset of $\mathbb{R}^n$ is closed and bounded
(\cref{thm:heine_borel}). If it were also open and non-empty, it would be a non-empty
\emph{clopen} proper subset of $\mathbb{R}^n$, which is impossible because $\mathbb{R}^n$ is
connected. Indeed $\mathbb{R}^n$ is path-connected, any two points being joined by a segment, and
hence connected by \cref{lem:path_connected_implies_connected}.

\textbf{(b) True,} by the same argument: complete $\Rightarrow$ closed, and non-empty, open,
closed and proper is again impossible in the connected space $\mathbb{R}^n$.

\textbf{(c) True.} Complete $\Rightarrow$ closed, and a closed set contains all its accumulation
points. Indeed, if $x$ is an accumulation point of $A$ then some sequence in $A$ converges to $x$,
and \cref{prop:open_closed_via_sequences}\textbf{(b)} puts $x$ in $A$.

\textbf{(d) False.} Every singleton $\{q\}$ is complete (its only Cauchy sequences are
eventually constant), yet
\[\mathbb{Q} = \bigcup_{q \in \mathbb{Q}}\{q\}\]
is a countable union of complete sets which is famously \emph{not} complete
(\cref{rem:complete_vs_closed}). Equivalently, in terms of closedness: a countable union of
closed sets need not be closed, exactly as in \cref{rem:finiteness_is_necessary}.

\textbf{(e) True,} this being \cref{rem:complete_vs_closed} stated for the ambient space
$\mathbb{R}^n$, and it is worth stressing that the \qt{if} direction is where completeness of the
ambient space is used. Inside $\mathbb{Q}$ instead of $\mathbb{R}$ the statement collapses:
$\mathbb{Q}$ is closed in itself but not complete.
\end{exercisesolution}

% Reclassified from ainote to remark: this is a second proof, which is content.
\begin{remark}[A second route to \textbf{(b)}, avoiding connectedness]
Part \textbf{(b)} is also script Exercise 9.76, and the script proves it without invoking
connectedness at all. The argument is more elementary and worth recording. Fix $x \in U$ and
$y \notin U$, and let $\gamma(t) := x+t(y-x)$, $t\in[0,1]$, be the segment between them. Put
$t_* := \sup\{t\in[0,1] : \gamma(t)\in U\}$. The point $\gamma(t_*)$ cannot lie in $U$, since
openness of $U$ would let $t_*$ be pushed further out, and yet it is a limit of points
$\gamma(t_k) \in U$ with $t_k \nearrow t_*$. That is a convergent, hence Cauchy, sequence in $U$
whose limit escapes $U$.

This avoids \cref{lem:path_connected_implies_connected} entirely, at the cost of only working in
$\mathbb{R}^n$, since the segment needs the vector space structure. The connectedness proof above
generalises to any connected, locally connected ambient space.
\end{remark}

\begin{exercisesolution}[Solution to \cref{ex:uniform_lipschitz_family}]
\label{sol:uniform_lipschitz_family}
% Generator: Claude Opus 5 (High)
\begin{enumerate}[label=\textbf{(\alph*)}]
  \item By the mean value theorem, for any $x,y \in [0,1]$ there is $\xi$ between them with
    \[\lvert f_k(x)-f_k(y)\rvert = \lvert f_k'(\xi)\rvert\,\lvert x-y\rvert \leq \lvert x-y\rvert,\]
    so every $f_k$ is $1$-Lipschitz \emph{with the same constant}. Given $\varepsilon>0$, the
    choice $\delta := \varepsilon$ then satisfies \cref{def:equicontinuous} for every $k$ and
    every pair $x,y$ at once. A uniformly Lipschitz family is always equicontinuous, and this is
    by far the most common way the hypothesis is verified in practice.
  \item Apply the same bound with $y := 0$ and use $f_k(0)=0$:
    \[\lvert f_k(x)\rvert = \lvert f_k(x)-f_k(0)\rvert \leq \lvert x\rvert \leq 1
      \quad\text{for all } x \in [0,1],\]
    so $\lVert f_k\rVert_\infty \leq 1$ for every $k$. Note that the normalisation $f_k(0)=0$ is
    doing real work: without it the family $f_k :\equiv k$ satisfies the derivative bound and is
    unbounded.
  \item Both hypotheses of \cref{thm:ascoli_arzela} hold on the compact set $[0,1]$, so some
    subsequence converges uniformly to a continuous limit.

    Replacing the uniform bound by mere differentiability destroys this, and
    \cref{ex:bounded_no_convergent_subsequence} is already the counterexample: $f_k(x) := x^k$
    is $C^1$ on $[0,1]$ with $f_k(0)=0$, so it satisfies every hypothesis \emph{except} the
    uniform one. Its derivative $f_k'(x) = kx^{k-1}$ has $\lVert f_k'\rVert_\infty = k$, which
    grows without bound. And we saw that no subsequence of $(x^k)$ converges. What matters is
    therefore not that each $f_k$ is well behaved, but that they are well behaved \emph{by a
    common margin}.
\end{enumerate}
\end{exercisesolution}

\begin{exercisesolution}[Proof of \cref{ex:annulus_continuous_image}]
\label{sol:annulus_continuous_image}
% Generator: Gemini 3.6 Flash (High)
% Correction: Claude Opus 5 (High) --- solution for part (c) supplied; U is bounded by
% \overline{B_2(0)}, not B_2(0), since (0,2) lies in U.
\begin{enumerate}[label=\textbf{(\alph*)}]
  \item The map $(x,y) \mapsto x^2+y^2$ is continuous, and $U$ is the preimage of the closed set
    $[1,4]$ under it, so $U$ is closed. It is also bounded, being contained in
    $\overline{B_2(0)}$. By Heine--Borel (\cref{thm:heine_borel}), $U$ is compact. It is moreover
    path-connected --- any two of its points can be joined by moving along a circle and then
    radially --- and therefore connected.

  \item Since $f$ is continuous and $U$ is compact, $f(U)$ is compact
    (\cref{item:continuity_preserves_compact}); since $U$ is connected, $f(U)$ is connected
    (\cref{thm:continuity_preserves_connected}). The connected subsets of $\mathbb{R}$ are exactly
    the intervals, and a compact interval is one of the form $[a,b]$. It remains to rule out
    $a = b$, that is, to check $f$ is non-constant: $(1,0)$ and $(0,2)$ both lie in $U$, and
    $f(1,0) = \sin 0 - 0 = 0$ while $f(0,2) = \sin 0 - 16 = -16$. Hence $f(U) = [a,b]$ with $a < b$.

  \item \textbf{No.} A loop running once around the inner hole cannot be contracted to a point
    within $U$ without crossing the removed disc $\{x^2+y^2 < 1\}$. The annulus is the standard
    example separating \qt{connected} from \qt{simply connected}, and it is worth keeping the two
    apart: part \textbf{(a)} is what \textbf{(b)} needs, and simple connectedness plays no role
    there at all.
\end{enumerate}
\end{exercisesolution}

\begin{exercisesolution}[Proof of \cref{ex:nearest_point_to_origin}]
\label{sol:nearest_point_to_origin}
% Generator: Claude Opus 5 (High)
\begin{enumerate}[label=\textbf{(\alph*)}]
  \item The set $A$ need not be bounded, so Weierstrass does not apply to it directly. The remedy is
    to cut it down to a compact piece without losing the minimum. Fix any $a_0 \in A$ and set
    $R := \lvert a_0\rvert$. The set
\[K := A \cap \overline{B_R(0)}\]
    is closed (an intersection of two closed sets) and bounded, hence compact by Heine--Borel
    (\cref{thm:heine_borel}), and it is non-empty because $a_0 \in K$. Since $f$ is continuous, it
    attains a minimum on $K$ (\cref{item:extreme_value_theorem}), say at $P_0 \in K \subseteq A$.
    This $P_0$ minimises over all of $A$: any $P \in A \setminus K$ has $\lvert P\rvert > R =
    \lvert a_0\rvert \ge f(P_0)$, because $a_0 \in K$.

  \item Take $A := \{x \in \mathbb{R}^n \mid 0 < \lvert x\rvert \le 1\}$, the closed unit ball with
    its centre removed. Then $\inf_A f = 0$, since $\lvert x\rvert$ can be made as small as we like,
    but $f(x) = 0$ would force $x = 0 \notin A$. So the infimum is not attained. The set is bounded,
    which shows that closedness --- not boundedness --- is the hypothesis doing the work in
    \textbf{(a)}.

  \item Let $A$ be convex and suppose $P_0, P_1 \in A$ both minimise, with common value
    $m := \lvert P_0\rvert = \lvert P_1\rvert$. If $m = 0$ then $P_0 = P_1 = 0$ and there is nothing
    to prove, so assume $m > 0$. By convexity the midpoint $M := \tfrac12(P_0+P_1)$ lies in $A$, and
    the triangle inequality gives
\[\lvert M\rvert \le \tfrac12\big(\lvert P_0\rvert + \lvert P_1\rvert\big) = m.\]
    Equality in the triangle inequality for the Euclidean norm forces $P_0$ and $P_1$ to be
    non-negative multiples of one another; having equal non-zero norms, they would then be equal.
    So if $P_0 \neq P_1$ the inequality is strict, $\lvert M\rvert < m$, contradicting minimality of
    $m$. Hence $P_0 = P_1$.

  \item Let $P_0 \in A$ be a minimiser and put $m := \lvert P_0\rvert$, which is positive because
    $0 \notin A$. Suppose, for contradiction, that $P_0$ were an interior point of $A$, so that
    $B_\varepsilon(P_0) \subseteq A$ for some $\varepsilon > 0$. Note first that
    $\varepsilon \le m$: otherwise $\lvert 0 - P_0\rvert = m < \varepsilon$ would put the origin
    inside $B_\varepsilon(P_0) \subseteq A$, contrary to $0 \notin A$. Now move from $P_0$ straight
    towards the origin: the point
\[P_1 := \left(1 - \frac{\varepsilon}{2m}\right) P_0\]
    satisfies $\lvert P_1 - P_0\rvert = \frac{\varepsilon}{2m}\lvert P_0\rvert = \frac{\varepsilon}{2}
    < \varepsilon$, so $P_1 \in A$. Since $\varepsilon \le m$, the scalar $1 - \frac{\varepsilon}{2m}$
    lies in $[\tfrac12, 1)$ and in particular is positive, so
    $\lvert P_1\rvert = \left(1 - \frac{\varepsilon}{2m}\right)m = m - \frac{\varepsilon}{2} < m$.
    This contradicts the minimality of $\lvert P_0\rvert$. Therefore $P_0$ is not interior, and since
    $P_0 \in A \subseteq \overline{A}$ it must lie in $\partial A$.
% Correction: Claude Opus 5 (High) --- the bound epsilon <= m was implicit in the first draft. It is
% needed: without it the scalar could be negative and |P_1| would be epsilon/2 - m, not m -
% epsilon/2. The official solution secures the same point by taking epsilon := min{epsilon', |P_0|}.
\end{enumerate}
\begin{remark}
The four parts isolate four different hypotheses and show what each one buys: closedness gives
\emph{existence}, convexity gives \emph{uniqueness}, and $0 \notin A$ gives \emph{location}. Part
\textbf{(b)} is what stops the list reading as decoration. It is worth noticing that boundedness is
never assumed: the trick in \textbf{(a)} of intersecting with a ball large enough to catch one
point of $A$ is the standard way to run a Weierstrass argument on an unbounded closed set, and it
recurs whenever a coercive function is minimised.
\end{remark}
\end{exercisesolution}

% Generator: Claude Opus 5 (High)
% No solution key ships with the Jossen papers, so this is an independent derivation.
\begin{exercisesolution}[\cref{ex:pseudocompactness_implications}]
\label{sol:pseudocompactness_implications}
\begin{enumerate}[label=\textbf{(\alph*)}]
  \item Assume \textbf{(ii)} and let $f : X \to \mathbb{R}$ be continuous. It attains a maximum
    at some $p$ and a minimum at some $q$, so $f(X) \subseteq [f(q),f(p)]$ and $f$ is bounded.

    Conversely, assume \textbf{(i)} and let $f : X \to \mathbb{R}$ be continuous. Then
    $M := \sup f(X)$ is a real number, since $f(X)$ is non-empty and bounded. Suppose $f$ never
    took the value $M$. Then $M - f(x) > 0$ for every $x$, so
    \[g : X \to \mathbb{R}, \qquad g(x) := \frac{1}{M-f(x)}\]
    is well defined and continuous. By definition of the supremum there is a sequence $(x_k)$ in
    $X$ with $f(x_k) \to M$, whence $g(x_k) \to +\infty$. Therefore $g$ is an unbounded continuous
    function, contradicting \textbf{(i)}. Consequently $f$ attains its maximum, and applying this
    to $-f$ gives the minimum.

  \item Assume \textbf{(iii)} and let $f : X \to \mathbb{R}$ be continuous. The preimages
    $f^{-1}\big((-k,k)\big)$, $k \in \mathbb{N}$, are open (\cref{def:continuous}) and cover $X$,
    because every $f(x)$ is a real number and so lies in $(-k,k)$ for $k$ large. Compactness
    yields finitely many indices $k_1 < \dots < k_r$ whose preimages already cover $X$, and these
    are nested, so $X = f^{-1}\big((-k_r,k_r)\big)$. Hence $\lvert f\rvert < k_r$ throughout.

  \item Assume \textbf{(iv)} and suppose some continuous $f : X \to \mathbb{R}$ were unbounded.
    Choose $x_k \in X$ with $\lvert f(x_k)\rvert \ge k$ for every $k$. By \textbf{(iv)} the
    sequence $(x_k)$ has an accumulation point $x \in X$, so some subsequence satisfies
    $x_{k_j} \to x$. Continuity gives $f(x_{k_j}) \to f(x)$, so the numbers $f(x_{k_j})$ are
    bounded. But $\lvert f(x_{k_j})\rvert \ge k_j \to \infty$, a contradiction.
\end{enumerate}
\begin{remark}[Weierstrass, split into its two halves]
Part \textbf{(a)} is the striking one, and it is worth seeing what it says: the two conclusions of
the extreme value theorem, \emph{bounded} and \emph{attained}, are not independent facts about the
space. Either one forces the other, with no compactness assumed anywhere. What compactness (or
\textbf{(iv)}) supplies is only the entry ticket \textbf{(i)}.

The trick in \textbf{(a)} is worth keeping: to force a supremum to be attained, divide by the gap
$M - f$ and let the failure to attain it manufacture an unbounded function. It reappears whenever a
statement about suprema has to be converted into a statement about boundedness.

None of the reverse implications hold. A space satisfying \textbf{(i)} is called
\newterm{pseudocompact}, and there are pseudocompact metric spaces that are not compact. Inside
$\mathbb{R}^n$ the gap closes, by \cref{thm:heine_borel}. The chain proved here is
\textbf{(iii)} $\Rightarrow$ \textbf{(i)} $\Leftrightarrow$ \textbf{(ii)} and \textbf{(iv)}
$\Rightarrow$ \textbf{(i)}; that \textbf{(iii)} and \textbf{(iv)} are themselves equivalent for
metric spaces is \cref{thm:sequential_topological_compactness_equivalence}, which is a much deeper
statement than either implication above.
\end{remark}
\end{exercisesolution}

% Generator: Claude Opus 5 (High)
\begin{exercisesolution}[\cref{ex:sumset_open_closed_compact}]
\label{sol:sumset_open_closed_compact}
\begin{enumerate}[label=\textbf{(\alph*)}]
  \item Suppose $A$ is open. For fixed $b \in \mathbb{R}^2$ the translate
    $A + b := \{a+b \mid a \in A\}$ is open, since translation is a homeomorphism of
    $\mathbb{R}^2$; concretely, $B_r(a) \subseteq A$ gives $B_r(a+b) \subseteq A+b$. Now
    \[A + B = \bigcup_{b \in B} (A+b)\]
    is a union of open sets and therefore open. The case of $B$ open is symmetric, since
    $A+B = B+A$.

  \item Let $c \in \overline{A+B}$ and pick $c_k = a_k + b_k \in A+B$ with $c_k \to c$, where
    $a_k \in A$ and $b_k \in B$. Since $B$ is compact, it is sequentially compact
    (\cref{thm:sequential_topological_compactness_equivalence}), so a subsequence
    $b_{k_j} \to b \in B$ converges. Then
    \[a_{k_j} = c_{k_j} - b_{k_j} \longrightarrow c - b,\]
    and $A$ is closed, so $c - b =: a$ lies in $A$. Hence $c = a + b \in A+B$, and $A+B$ is closed.

  \item The product $A \times B \subseteq \mathbb{R}^4$ is compact: it is closed and bounded,
    because $A$ and $B$ are, so \cref{thm:heine_borel} applies. Addition
    \[s : \mathbb{R}^2\times\mathbb{R}^2 \to \mathbb{R}^2, \qquad s(a,b) := a+b\]
    is continuous, being linear, and $A+B = s(A\times B)$ is the continuous image of a compact set,
    hence compact.
\end{enumerate}
\begin{remark}[Compactness of $B$ is not decoration]
Two closed sets can have a non-closed sum, so \textbf{(b)} genuinely needs more than closedness.
Take
\[A := \{(n,0) \mid n \in \mathbb{N}\}, \qquad B := \{(-n,\tfrac1n) \mid n \in \mathbb{N}\}.\]
Both are closed, since each is a set of isolated points running off to infinity and so has no
accumulation point at all. Their sum contains $(0,\tfrac1n)$ for every $n$, and those converge to
the origin. But $(0,0) \in A+B$ would need $n = m$ and $\tfrac1m = 0$ at once, so the origin is
missing and $A+B$ is not closed.

What compactness of $B$ supplies in the proof is precisely the convergent subsequence of
$(b_k)$: closedness of $B$ alone would let $(b_k)$ escape to infinity, which is exactly what the
counterexample arranges. Note also that nothing in any of the three arguments used the dimension,
so all three hold verbatim in $\mathbb{R}^n$, and \textbf{(a)} and \textbf{(b)} hold in any normed
space.
\end{remark}
\end{exercisesolution}

\begin{exercisesolution}[Solution to \cref{ex:kobel_compactness_characterisation}]
\label{sol:kobel_compactness_characterisation}
% Generator: Claude Opus 5 (High)
Only \textbf{(c)} is true.
\begin{enumerate}[label=\textbf{(\alph*)}]
  \item \textbf{False.} This is Heine--Borel (\cref{thm:heine_borel}), and it is a theorem about
    $\mathbb{R}^n$ with the standard metric, not about metric spaces. \Cref{ex:heine_borel_fails}
    breaks it with the bounded metric $\tilde d$, under which $\mathbb{N}$ is closed and bounded
    and not compact; an infinite set with the discrete metric does the same.
  \item \textbf{False.} Being open and closed is a statement about how $E$ sits in $X$, and it has
    no bearing on compactness. Any space is open and closed in itself, so $E = X = \mathbb{R}$ is
    a counterexample.
  \item \textbf{True}, and it is the honest general replacement for \textbf{(a)}: a metric space is
    compact if and only if it is complete and totally bounded
    (\cref{rem:compact_iff_complete_totally_bounded}, built on
    \cref{lem:sequentially_compact_totally_bounded}).
  \item \textbf{False.} Total boundedness alone is not enough --- completeness is the other half of
    \textbf{(c)} --- and openness does not supply it. The interval $(0,1) \subseteq \mathbb{R}$ is
    open and totally bounded and not compact.
\end{enumerate}
\begin{remark}[Why \textbf{(a)} is the tempting one]
\Cref{rem:boundedness_not_topological} already contains the reason \textbf{(a)} fails and
\textbf{(c)} succeeds, and it is worth restating in one line, because the question is designed to
catch a reader who remembers Heine--Borel without its hypothesis. Compactness is a topological
property; boundedness is not, since two metrics with the same open sets can disagree about which
sets are bounded. Total boundedness, unlike boundedness, is not so easily destroyed by rescaling
the metric, and that is exactly why it is the notion that survives into the general statement.
\end{remark}
\end{exercisesolution}

\begin{exercisesolution}[Solution to \cref{ex:box_preimage_of_compact_not_compact}]
\label{sol:box_preimage_of_compact_not_compact}
% Generator: Claude Opus 5 (High)
Take the projection onto the first coordinate,
\[g(x,y) := (x,\,0), \qquad K := \{(0,0)\}.\]
Then $g$ is linear and therefore continuous, and it is nonconstant since $g(1,0) = (1,0)$ while
$g(0,0) = (0,0)$. The set $K$ is a single point, hence compact. But
\[g^{-1}(K) = \{(x,y) \mid x = 0\} = \{0\}\times\mathbb{R}\]
is the whole $y$-axis, which is closed but unbounded and so not compact by
\cref{thm:heine_borel}.

\begin{remark}[Continuity preserves compactness in one direction only]
\Cref{item:continuity_preserves_compact} says that $g(C)$ is compact whenever $C$ is, and it is easy
to carry that away as a statement about $g$ and compactness rather than about images. It is not:
the preimage of a compact set under a continuous map can be as large as one likes, and the
example above makes it a whole line by collapsing one direction. A map that \emph{does} have the
property is called \newterm{proper}, and properness is a genuine extra hypothesis, not a
consequence of continuity. \Cref{ex:box_proper_map_not_diffeomorphism} constructs one, and shows
that properness in turn buys much less than one might hope.

The official key gives a different answer with the same moral, $g(x,y) := (\arctan x, \arctan y)$
together with any compact $K$ containing the open square $(-\tfrac\pi2,\tfrac\pi2)^2$, for which
$g^{-1}(K) = \mathbb{R}^2$. There the map is injective; it is the \emph{boundedness} of its
image that does the work, rather than a collapsed direction.
\end{remark}
\end{exercisesolution}




% ============================================================
% Chapter 7 --- Connectedness
% Originally: content/week-03.tex, \section{Connectedness} (Corsin Week 3, Monday) and
%             \section{Path connectedness} (Corsin Week 3, Friday)
% ============================================================
\chapter{Connectedness}
\label{ch:connectedness}

% Transition: Claude Opus 5 (Medium)
Compactness told us when a space has \qt{enough finiteness} to guarantee extrema and uniform
behaviour. This chapter turns to a completely different, qualitative kind of structure: whether a
space can be split into separate pieces at all. Connectedness and path connectedness are the two
answers, they are not equivalent, and the topologist's sine curve is the standard witness.

% Originally: content/week-03.tex, \section{Connectedness} (Corsin Week 3, Monday)

\section{Connectedness}
\label{sec:connectedness}

\begin{definition}[Connected and disconnected metric spaces]
\label{def:connectedness}
A metric space $(X,d)$ is \newterm{disconnected} (\germanterm{unzusammenh\"angend}) if there
exist $U_1, U_2 \subseteq X$ subsets which are
\begin{enumerate}[label=\textbf{(\alph*)}]
  \item \textbf{open},
  \item \textbf{non-empty}, $U_1 \neq \emptyset \neq U_2$,
  \item \textbf{disjoint}, $U_1 \cap U_2 = \emptyset$,
\end{enumerate}
such that $X = U_1 \cup U_2$. Similarly, for a subset $A \subseteq X$, it is disconnected if it is
disconnected as a metric space with the induced metric $d|_{A \times A}$. \newterm{Connected}
(\germanterm{zusammenh\"angend}) is the negation of disconnected.
\end{definition}

Typically, proving that a set or metric space is connected is done \textbf{by contradiction},
which is why one should remember the definition of \textit{disconnected}. Intuitively, connected
means: \qt{the space cannot be separated by open sets.}

\begin{center}
\begin{tikzpicture}[scale=1.0]
  % Left panel: Disconnected
  \begin{scope}[shift={(0,0)}]
    \draw[thick] plot[smooth cycle, tension=0.7] coordinates {(-1.5,-0.8) (-1.8,0.6) (-0.8,1.2) (-0.6,-0.2)};
    \draw[thick] plot[smooth cycle, tension=0.7] coordinates {(0.4,-0.6) (0.2,0.8) (1.2,1.0) (1.0,-0.4)};

    \draw[blue!80!black, thick, dashed] plot[smooth cycle, tension=0.7] coordinates {(-1.7,-1.0) (-2.0,0.8) (-0.6,1.4) (-0.4,-0.4)};
    \node[blue!80!black] at (-1.2,1.2) {$U_1$};

    \draw[blue!80!black, thick, dashed] plot[smooth cycle, tension=0.7] coordinates {(0.2,-0.8) (0.0,1.0) (1.4,1.2) (1.2,-0.6)};
    \node[blue!80!black] at (0.8,1.0) {$U_2$};

    \node[below, font=\small] at (-0.3,-1.3) {\textbf{Disconnected}};
  \end{scope}

  % Right panel: Connected
  \begin{scope}[shift={(5.0,0)}]
    \draw[thick] plot[smooth cycle, tension=0.7] coordinates {(-1.4,-0.6) (-1.6,1.2) (0.2,0.6) (1.2,1.0) (1.0,-0.8) (-0.2,-1.0)};

    \draw[blue!80!black, thick, dashed] (-1.7,-0.8).. controls (-1.9,1.4) and (0.0,1.0).. (0.1,0.2).. controls (0.0,-0.8) and (-0.5,-1.2).. (-1.7,-0.8);
    \node[blue!80!black] at (-1.2,1.1) {$U_1$};

    \draw[blue!80!black, thick, dashed] (1.3,-0.9).. controls (1.5,1.2) and (-0.1,0.8).. (0.2,0.0).. controls (0.1,-0.7) and (0.6,-1.1).. (1.3,-0.9);
    \node[blue!80!black] at (1.0,0.9) {$U_2$};

    \draw[orange!90!black, <->, thick] (0.1,0.2) -- node[above, font=\tiny] {gap} (0.2,0.0);

    \node[below, font=\small] at (0,-1.3) {\textbf{Connected}};
  \end{scope}
\end{tikzpicture}
\end{center}

% Generator: Claude Opus 5 (High) --- the base case of connectedness, used three times in this
% chapter (path-connectedness, the topologist's sine curve, the intermediate value theorem)
% but stated nowhere in the source notes.
\begin{proposition}[Intervals are connected]
\label{prop:intervals_connected}
Every interval $I \subseteq \mathbb{R}$ --- open, closed, half-open, bounded or not --- is
connected (\cref{def:connectedness}). Conversely, a connected subset of $\mathbb{R}$ is an
interval.
\end{proposition}

\begin{proof}
Suppose $I = U_1 \sqcup U_2$ with $U_1, U_2$ open in $I$, disjoint and both non-empty. Pick
$a \in U_1$ and $b \in U_2$; without loss of generality $a<b$, and $[a,b] \subseteq I$ because
$I$ is an interval. Set
\[c := \sup\{\,t \in [a,b] \ :\ [a,t] \subseteq U_1\,\},\]
which is well defined since $a$ belongs to the set. Now $c \in [a,b] \subseteq I$, so $c$ lies
in $U_1$ or in $U_2$, and both lead to a contradiction:
\begin{itemize}
  \item If $c \in U_1$, openness gives $\varepsilon>0$ with $(c-\varepsilon,c+\varepsilon)\cap I
    \subseteq U_1$; together with $[a,c) \subseteq U_1$ this yields $[a,t] \subseteq U_1$ for
    some $t>c$ unless $c=b$ --- contradicting either the supremum or $b \in U_2$.
  \item If $c \in U_2$, openness gives $\varepsilon>0$ with $(c-\varepsilon,c] \cap I \subseteq
    U_2$. But $c>a$ (as $a \in U_1$ is open in $I$), so by definition of the supremum there is
    $t \in (c-\varepsilon,c]$ with $[a,t] \subseteq U_1$ --- contradicting disjointness.
\end{itemize}
So no such splitting exists, and $I$ is connected. For the converse, a set $A \subseteq
\mathbb{R}$ that is not an interval admits $a<c<b$ with $a,b \in A$ and $c \notin A$; then
$A \cap (-\infty,c)$ and $A \cap (c,\infty)$ disconnect it.
\end{proof}

\begin{remark}[Why this is the only example one ever needs]
\label{rem:intervals_are_the_base_case}
Every connectedness proof in this chapter reduces to \cref{prop:intervals_connected}: a path is
a continuous image of $[0,1]$ (\cref{def:path}), so path-connected sets are connected
(\cref{lem:path_connected_implies_connected}) purely because the \emph{interval} is; and the
intermediate value theorem, invoked in the topologist's sine curve below, is
\cref{thm:continuity_preserves_connected} applied to an interval, since the connected subsets of
$\mathbb{R}$ are exactly the intervals. In other words, the single fact that $\mathbb{R}$ has no
gaps --- the completeness axiom, which is what powers the supremum above --- is the source of
all of it.
\end{remark}

\begin{aiexample}[Disconnectedness of Rational Numbers]
% Generator: Gemini 3.6 Flash (Medium)
The set of rational numbers $\mathbb{Q} \subset \mathbb{R}$ equipped with the subspace topology is disconnected. To see this, define
\[
U := \mathbb{Q} \cap (-\infty, \sqrt{2}) \quad \text{and} \quad V := \mathbb{Q} \cap (\sqrt{2}, \infty).
\]
Since $\sqrt{2} \notin \mathbb{Q}$, we have $U \cup V = \mathbb{Q}$ and $U \cap V = \emptyset$. Both $U$ and $V$ are non-empty and open in the subspace topology of $\mathbb{Q}$ (as intersections of $\mathbb{Q}$ with open intervals in $\mathbb{R}$). Hence $U$ and $V$ form a separation of $\mathbb{Q}$.
\end{aiexample}

% Supplement: Sascha Brack/Class Notes/Week_02_Notes_Friday_Updated.pdf, p. 9
\begin{example}[A disconnected set in $\mathbb{R}^2$]
Let $X := ([0,1] \times [0,1]) \cup ([2,3] \times \{1\}) \subseteq \mathbb{R}^2$. Is $X$ connected?
No. Let $U := \{(x,y) \in X \mid x < 1.5\}$ and $V := \{(x,y) \in X \mid x > 1.5\}$. We can write these as $U = X \cap \{(x,y) \in \mathbb{R}^2 \mid x < 1.5\}$ and $V = X \cap \{(x,y) \in \mathbb{R}^2 \mid x > 1.5\}$. Both $U$ and $V$ are open in the subspace topology of $X$ because they are intersections of $X$ with open half-planes in $\mathbb{R}^2$. Furthermore, $X = U \cup V$, and $U, V$ are disjoint and non-empty. Thus, $X$ is disconnected.
\begin{center}
\begin{tikzpicture}[scale=1.5]
  \draw[->, thick] (-0.2,0) -- (3.5,0);
  \draw[->, thick] (0,-0.2) -- (0,1.5);
  
  \draw[very thick, fill=purple!20, pattern=north east lines, pattern color=purple] (0,0) rectangle (1,1);
  \draw[very thick, purple] (2,1) -- (3,1);
  \node[purple, left] at (-0.1, 0.5) {$X$};
  
  \draw[green!60!black, dashed] (0.5, 0.5) circle (0.8);
  \node[green!60!black] at (1.2, 0.6) {$U$};
  
  \draw[green!60!black, dashed] (2.5, 1) ellipse (0.7 and 0.3);
  \node[green!60!black] at (2.5, 1.4) {$V$};
\end{tikzpicture}
\end{center}
\end{example}

\begin{theorem}[Continuous functions preserve connectedness]
\label{thm:continuity_preserves_connected}
Let $X$, $Y$ be metric spaces, $f : X \to Y$ continuous. If $A \subseteq X$ is connected, then
$f(A)$ is connected.
\end{theorem}

% Generator: Claude Opus 5 (High) --- stated without proof in the source; the argument is three
% lines and is the exact dual of \cref{item:continuity_preserves_compact}.
\begin{proof}
We may assume $A = X$, replacing $X$ by $A$ with the induced metric and $Y$ by $f(A)$. Suppose
$f(X)$ were disconnected, say $f(X) = V_1 \sqcup V_2$ with $V_1,V_2$ open in $f(X)$, disjoint
and non-empty. Then $U_i := f^{-1}(V_i)$ is open by \cref{def:continuous}\textbf{(c)},
non-empty because $V_i$ meets the image, and $U_1 \cap U_2 = \emptyset$ with
$U_1 \cup U_2 = X$ because the $V_i$ partition $f(X)$. So $U_1,U_2$ disconnect $X$, contrary to
hypothesis.
\end{proof}

\begin{corollary}[Intermediate value theorem]
\label{cor:intermediate_value_theorem}
Let $f : [a,b] \to \mathbb{R}$ be continuous. Then $f$ attains every value between $f(a)$ and
$f(b)$.
\end{corollary}

\begin{proof}
By \cref{prop:intervals_connected}, $[a,b]$ is connected, so $f([a,b])$ is connected by
\cref{thm:continuity_preserves_connected}, hence is an interval (\cref{prop:intervals_connected}
again). An interval containing $f(a)$ and $f(b)$ contains everything between them.
\end{proof}

\begin{ainote}
The intermediate value theorem is a result from \emph{Analysis I}, but it is worth re-deriving
it here in one line: it is the first place where connectedness earns its keep, and the proof
makes visible that the theorem is really a statement about $\mathbb{R}$ having no gaps rather
than about $f$. It is used in the topologist's sine curve
(\cpageref{ex:topologists_sine_curve}) below.
\end{ainote}

% Supplement: Sascha Brack/Class Notes/Week_03_Notes_Monday.pdf, p. 4
\begin{corollary}[Linear maps preserve the unit cube]
\label{cor:linear_maps_unit_cube}
Given the $n$-dimensional unit cube $[0,1]^n \subset \mathbb{R}^n$, any linear map $L : \mathbb{R}^n \to \mathbb{R}^n$ maps $[0,1]^n$ to an $n$-dimensional parallelepiped. Since linear maps in Euclidean space are continuous, the image $L([0,1]^n)$ is connected. Furthermore, since $[0,1]^n$ is compact, $L([0,1]^n)$ is also compact.
\end{corollary}

\begin{exercise}[Properties of connectedness]
\label{ex:connectedness_tf}
Are the following true or false?
\begin{enumerate}[label=\textbf{(\alph*)}]
  \item The union of connected subsets with a common point is connected.
  \item The intersection of connected subsets is connected.
  % Not a first introduction -- \newterm/\germanterm for "closure" belong to
  % \cref{def:interior_closure_boundary}, so this is a plain back-reference.
  \item The closure (\cref{def:interior_closure_boundary}) of a connected subset is connected.
\end{enumerate}
\end{exercise}

% Kept inline: solved directly in Corsin Nick/Class Notes/Week 3.pdf, pp. 5--6.
\begin{exercisesolution}[Proof of \cref{ex:connectedness_tf}]
\textbf{(a) True.} \textit{Proof.} Let $X$ be a metric space, $V_1, V_2 \subseteq X$ connected,
and $x_0 \in V_1 \cap V_2$. Suppose $V_1 \cup V_2 \subseteq U_1 \cup U_2$ with
$U_1, U_2 \subseteq X$ open and disjoint. Assume without loss of generality $x_0 \in U_1$. Then
$U_1 \cap V_1 \neq \emptyset$ and $U_1 \cap V_2 \neq \emptyset$. So, since $V_1, V_2$ are
connected, $V_1 \cap U_2 = V_2 \cap U_2 = \emptyset$. Therefore
$(V_1 \cup V_2) \cap U_2 = \emptyset$ and $V_1 \cup V_2$ is connected.

\textbf{(b) False.} \textit{Counterexample:}
\[
X = \mathbb{R}^2, \quad U = [-1,1] \times \{0\}, \quad V = S^1 = \{x \in \mathbb{R}^2 : \|x\| = 1\},
\]
\[
U \cap V = \{(-1,0), (1,0)\}.
\]

\textbf{(c) True.} Suppose $X$ is a metric space, $A \subseteq X$ a subset with closure
$\overline{A}$. Suppose $U_1, U_2 \subseteq X$ are disjoint open sets covering $\overline{A}$, and
assume without loss of generality $\overline{A} \cap U_1 \neq \emptyset$. Then, since $U_1$ is
open, $A \cap U_1 \neq \emptyset$. By connectedness of $A$, $U_2 \cap A = \emptyset$, and since
$U_2$ is open, also $U_2 \cap \overline{A} = \emptyset$.
\end{exercisesolution}

\begin{ainote}
Corsin writes $\|x\| = R^2$ in the definition of $S^1$ above; it must be $\|x\| = 1$, as the
stated intersection $\{(\pm 1, 0)\}$ confirms. Corrected here.
\end{ainote}

\begin{center}
\begin{tikzpicture}[scale=1.1]
  \draw[orange!90!black, thick, fill=orange!10] plot[smooth cycle, tension=0.7] coordinates {(-1.6,-0.6) (-1.8,0.8) (-0.3,0.6) (-0.5,-0.6)};
  \node[orange!90!black, font=\large] at (-1.2,0) {$V_1$};

  \draw[orange!90!black, thick, fill=orange!10] plot[smooth cycle, tension=0.7] coordinates {(-0.3,-0.6) (-0.5,0.6) (1.0,0.8) (0.8,-0.6)};
  \node[orange!90!black, font=\large] at (0.5,0) {$V_2$};

  \fill[purple] (-0.4,0) circle (2.5pt) node[below, font=\small] {$x_0$};

  \draw[red!80!black, thick, dotted] plot[smooth cycle, tension=0.7] coordinates {(-2.0,-0.9) (-2.1,1.1) (1.3,1.1) (1.1,-0.9)};
  \node[red!80!black, font=\large] at (-1.6,-0.7) {$U_1$};

  \draw[red!80!black, thick, dotted] plot[smooth cycle, tension=0.7] coordinates {(2.0,-0.8) (1.8,0.9) (3.0,0.9) (2.8,-0.8)};
  \node[red!80!black, font=\large] at (2.4,0) {$U_2$};
  \node[red!80!black, font=\footnotesize, align=center] at (2.4,-1.1) {no\\intersection};
\end{tikzpicture}
\end{center}

\begin{center}
\begin{tikzpicture}[>=Latex, scale=1.3]
  \draw[->] (-1.6,0) -- (1.6,0) node[right] {$x$};
  \draw[->] (0,-1.6) -- (0,1.6) node[above] {$y$};

  \draw[red!80!black, thick] (0,0) circle (1.0);
  \node[red!80!black, above right] at (0.7,0.7) {$S^1$};

  \draw[orange!90!black, ultra thick] (-1,0) -- (1,0);
  \node[orange!90!black, below] at (0,-0.1) {$U = [-1,1]\times\{0\}$};

  \fill[purple] (-1,0) circle (2pt) node[above left, font=\footnotesize] {$(-1,0)$};
  \fill[purple] (1,0) circle (2pt) node[above right, font=\footnotesize] {$(1,0)$};
\end{tikzpicture}
\end{center}

% Originally: content/week-03.tex, \section{Path connectedness} (Corsin Week 3, Friday)

\section{Path connectedness}
\label{sec:path_connectedness}

% Transition: Claude Opus 5 (High)
\Cref{def:connectedness} defines connectedness negatively. A space is connected when it
\emph{cannot} be split, which is why proofs about it run by contradiction. That is mathematically
convenient and intuitively unsatisfying, since it says nothing whatever about what holds a
connected space together.

Path connectedness is the positive version, and it says the obvious thing: any two points can be
joined by a curve running inside the space. It implies connectedness
(\cref{lem:path_connected_implies_connected}), and on open subsets of $\mathbb{R}^n$ the two
notions agree (\cref{prop:open_euclidean_connected_path_connected}). In general they come apart,
and the topologist's sine curve at the end of this section is the standard witness.

% Source: Corsin Nick/Class Notes/Week 3.pdf, pp. 1--14
\begin{definition}[Path]
\label{def:path}
Let $(X,d)$ be a metric space. A continuous function $\gamma : [0,1] \to X$, where $[0,1]$
carries the standard metric, is called a \newterm{path} (\germanterm{Weg}). More generally, we
can replace $[0,1]$ by $[a,b]$ for $a < b$.
\end{definition}

\begin{definition}[Path-connected]
\label{def:path_connected}
A metric space $X$ (or subset thereof) is \newterm{path-connected}
(\germanterm{wegzusammenh\"angend}) if for all $x, y \in X$ there exists a path
(\cref{def:path}) $\gamma : [0,1] \to X$ with $\gamma(0) = x$, $\gamma(1) = y$.
\begin{center}
\begin{tikzpicture}
  \draw[very thick, fill=gray!10] plot[smooth cycle, tension=0.7] coordinates {(-2, -1) (-2.5, 0.5) (-1, 1.5) (1.5, 1) (2.5, -0.5) (1, -1.5)};
  \node[font=\large] at (2, 0.8) {$X$};
  
  \fill (-1.5, -0.2) circle (1.5pt) node[below] {$x$};
  \fill (1.8, -0.1) circle (1.5pt) node[below] {$y$};
  
  \draw[blue, thick, ->] (-1.5, -0.2) .. controls (-1, 1) and (1, -1) .. (1.8, -0.1);
  
  \node[blue, above] at (-1.5, -0.2) {$\gamma(0)$};
  \node[blue, above] at (1.8, -0.1) {$\gamma(1)$};
  \node[blue, below] at (0, -0.1) {$\gamma([0,1])$};
\end{tikzpicture}
\end{center}
\end{definition}

\begin{lemma}[Path-connected implies connected]
\label{lem:path_connected_implies_connected}
Let $X$ be a metric space. Then $X$ path-connected (\cref{def:path_connected}) implies $X$
connected (\cref{def:connectedness}).
\end{lemma}

% Generator: Claude Opus 5 (High) --- stated without proof in the source.
\begin{proof}
Suppose $X$ were path-connected but disconnected, say $X = U_1 \sqcup U_2$ with $U_1,U_2$ open,
disjoint and non-empty. Pick $x \in U_1$, $y \in U_2$ and a path $\gamma : [0,1] \to X$ with
$\gamma(0)=x$, $\gamma(1)=y$. Then $\gamma^{-1}(U_1)$ and $\gamma^{-1}(U_2)$ are open in $[0,1]$
by \cref{def:continuous}\textbf{(c)}, disjoint, cover $[0,1]$, and are non-empty (they contain
$0$ and $1$ respectively). This disconnects $[0,1]$, contradicting
\cref{prop:intervals_connected}.
\end{proof}

% Source: Jérôme Paschoud/Notizen/Woche 3.1 Banach und Kompaktheit.pdf, p. 6
% Generator: Gemini 3.6 Flash (High)
\begin{exercise}[The unit disk is connected]
\label{ex:disk_connected}
Prove that the set $S^1 := \{ (x,y) \in \mathbb{R}^2 \mid x^2 + y^2 \leq 1 \}$ is connected.
\end{exercise}

\begin{exercisesolution}[Proof of \cref{ex:disk_connected}]
Instead of proving connectedness directly via open sets, we can show that $S^1$ is path-connected, which implies connectedness (\cref{lem:path_connected_implies_connected}). Given any two points $a, b \in S^1$, we can connect them via a straight line segment. The segment can be parameterized as $\gamma(t) := (1-t)a + tb$ for $t \in [0,1]$. Since $S^1$ (the unit disk in $\mathbb{R}^2$) is convex, for any $t \in [0,1]$ we have $\lVert \gamma(t) \rVert \leq (1-t)\lVert a \rVert + t \lVert b \rVert \leq 1-t + t = 1$, so the entire path lies in $S^1$. Because any two points can be connected by a continuous path within $S^1$, the set is path-connected, and therefore connected.

(Note: The standard notation for the unit disk is $D^2$ or $B_1(0)$, whereas $S^1$ usually denotes the boundary circle $\{x^2+y^2=1\}$, but we reproduce the notation $S^1$ here as it appears in the source notes.)
\end{exercisesolution}

\begin{proposition}[Connectedness and path-connectedness in open Euclidean subsets]
\label{prop:open_euclidean_connected_path_connected}
Let $U \subseteq \mathbb{R}^n$ be open (\cref{def:open_closed}). Then
\[
U \text{ path-connected} \iff U \text{ connected}.
\]
\end{proposition}

% Generator: Claude Opus 5 (High) --- stated without proof in the source; the "(<=)" direction
% is the standard "the set of reachable points is clopen" argument and is worth having, since it
% is the only place in this chapter where openness is genuinely used.
\begin{proof}
\textbf{($\Rightarrow$)} is \cref{lem:path_connected_implies_connected}, which needs neither
openness nor $\mathbb{R}^n$.

\textbf{($\Leftarrow$)} Assume $U$ is connected and non-empty (the empty set is vacuously
path-connected), and fix $x_0 \in U$. Let
\[A := \{x \in U \ :\ \text{some path in } U \text{ joins } x_0 \text{ to } x\}.\]
We show $A$ and $U \setminus A$ are both open; since $x_0 \in A$ (take the constant path),
connectedness then forces $U\setminus A = \emptyset$, i.e.\ $A = U$, which is the claim.

Let $x \in U$ and choose $r>0$ with $B_r(x) \subseteq U$, using that $U$ is open. The ball is
convex, so every $y \in B_r(x)$ is joined to $x$ by the straight segment
$t \mapsto (1-t)x+ty$, which stays inside $B_r(x) \subseteq U$. Consequently:
\begin{itemize}
  \item If $x \in A$, concatenating a path from $x_0$ to $x$ with that segment joins $x_0$ to
    any $y \in B_r(x)$, so $B_r(x) \subseteq A$ and $A$ is open.
  \item If $x \notin A$, then no $y \in B_r(x)$ can lie in $A$, since otherwise running the
    segment backwards from $y$ to $x$ would join $x_0$ to $x$. So $B_r(x) \subseteq U\setminus A$ and
    $U\setminus A$ is open.
\end{itemize}
(Concatenation of two paths is again a path: traverse the first on $[0,\tfrac12]$ and the second
on $[\tfrac12,1]$, each at double speed; the result is continuous because the two pieces agree
at $t=\tfrac12$.)
\end{proof}

\begin{importantremark}[Only for open subsets]
\label{rem:path_connected_needs_open}
The implication \qt{connected $\Rightarrow$ path-connected} is \emph{false} without openness,
and the proof shows exactly why: the ball $B_r(x)$ around a point of $U$ was needed to supply
the straight segment. For a set that is not open there may be points with no room around them at
all, and the topologist's sine curve (\cref{ex:topologists_sine_curve}) below is precisely such
a case. There the origin has no neighbourhood in $S$ that is a nice piece of curve.
\end{importantremark}

% Source: Corsin Nick/Class Notes/Week 3.pdf, pp. 1--14
% (Re-cited: the two proofs and the important remark above are additions, not transcription.)
\begin{example}[The topologist's sine curve]
\label{ex:topologists_sine_curve}
\[S := \left\{\left(t, \sin\left(\tfrac{1}{t}\right)\right) : t > 0\right\} \cup \{(0,0)\} \subseteq \mathbb{R}^2\]
$S$ is \textbf{connected} but \textbf{not path-connected}!
\end{example}

\begin{ainote}
Corsin writes the second piece as $\cup\,\{0\}$; since $S \subseteq \mathbb{R}^2$, the point
being adjoined is the origin $(0,0) \in \mathbb{R}^2$, not the number $0$. Written as
$\cup\,\{(0,0)\}$ above, which is also what his own figure and the proof below use.
\end{ainote}

\begin{center}
\begin{tikzpicture}[>=Latex, scale=1.3]
  \draw[->] (-0.2,0) -- (3.2,0) node[right] {$t$};
  \draw[->] (0,-1.4) -- (0,1.4) node[above] {$y$};
  \draw (0,1.0) -- (-0.05,1.0) node[left, font=\footnotesize] {$1$};
  \draw (0,-1.0) -- (-0.05,-1.0) node[left, font=\footnotesize] {$-1$};
  % Horizontal figure-coordinate \x corresponds to t = \x/10, since the plot is sin(10/\x).
  \draw (1.0,0.05) -- (1.0,-0.05) node[below, font=\footnotesize] {$0.1$};
  \draw (2.0,0.05) -- (2.0,-0.05) node[below, font=\footnotesize] {$0.2$};

  \draw[blue!80!black, thick] plot[domain=0.08:2.8, samples=250] (\x, {sin((1.0/(\x*0.1)) r)});
  \fill[blue!80!black] (0,0) circle (1.5pt) node[left, font=\footnotesize] {$(0,0)$};
  \node[above, font=\small] at (1.5,1.3) {$S = \{(t, {\sin(1/t})) : t > 0\} \cup \{(0,0)\}$};
\end{tikzpicture}
\end{center}

\begin{proof}[$S$ is connected]
Write $G := S \setminus \{(0,0)\} = \{(t,\sin(1/t)) : t>0\}$ for the graph part, and suppose for
contradiction that $S$ is disconnected, say $S = V_1 \sqcup V_2$ with $V_1,V_2$ non-empty,
disjoint and open in $S$. Say $(0,0) \in V_1$.

$G$ is the continuous image of the interval $(0,\infty)$ under
$t \mapsto (t,\sin(1/t))$, hence path-connected, hence connected
(\cref{lem:path_connected_implies_connected}). A connected set meeting both parts of a
separation would itself be separated, so $G$ lies entirely in $V_1$ or entirely in $V_2$.

It cannot lie in $V_1$: then $V_2 \subseteq S \setminus (G \cup \{(0,0)\}) = \emptyset$,
contradicting $V_2 \neq \emptyset$. So $G \subseteq V_2$, and therefore
$V_1 = \{(0,0)\}$.

But $V_1$ is open in $S$, so some ball $B_\varepsilon\big((0,0)\big) \cap S$ equals
$\{(0,0)\}$, and that is false. For $t := \tfrac{1}{k\pi}$ with $k$ large, the point
$(t,\sin(1/t)) = (t,0) \in G$ lies within $\varepsilon$ of the origin. This is the required
contradiction, so $S$ is connected.
\end{proof}

\begin{proof}[$S$ is not path-connected]
Assume by contradiction there is a path from $(0,0)$ to $(\tfrac{1}{\pi}, 0)$,
\[
\gamma : [0,1] \to [0, \tfrac{1}{\pi}] \times [-1,1], \qquad t \mapsto (x(t), y(t)).
\]
Since $\gamma$ is continuous, $x$ and $y$ must also be continuous. By the \newterm{intermediate
value theorem} (\cref{cor:intermediate_value_theorem}), $x : [0,1] \to [0,\tfrac{1}{\pi}]$
attains every value in $[0,\tfrac{1}{\pi}]$, since $x(0) = 0$ and $x(1) = \tfrac{1}{\pi}$.
Therefore, we find a sequence $(a_n)_{n=1}^{\infty}$ such that
\[
x(a_n) = \frac{1}{2\pi n + \tfrac{\pi}{2}}.
\]
We may moreover choose the $a_n$ \textbf{decreasing}: having found $a_{n-1}$, apply the
intermediate value theorem on $[0,a_{n-1}]$ rather than on $[0,1]$, which is legitimate because
$x(0) = 0$ and $x(a_{n-1}) = \tfrac{1}{2\pi(n-1)+\pi/2}$ already straddle the next target value.
This yields $a_n \in (0,a_{n-1})$, so $(a_n)$ is decreasing and bounded below, hence convergent,
say $a_n \to \ell \in [0,1]$.

Now evaluate $\gamma$ along this sequence. By construction $x(a_n) \to 0$, so continuity gives
$x(\ell) = 0$; and $y(a_n) = \sin\!\big(2\pi n + \tfrac\pi2\big) = 1$ for every $n$, so
$y(\ell) = 1$. Hence $\gamma(\ell) = (0,1)$. But $(0,1) \notin S$: the only point of $S$ whose
first coordinate vanishes is the origin. This contradicts $\gamma([0,1]) \subseteq S$, so no
such path exists.
\end{proof}

\begin{ainote}
Corsin's argument produces the $a_n$ from the intermediate value theorem and then asserts
$a_n \to 0$. That does not follow from the choice as stated. The intermediate value theorem
returns \emph{some} preimage, and nothing so far prevents all of them from sitting near $1$.
The repair is the standard one and costs one sentence: apply the theorem on the shrinking
intervals $[0,a_{n-1}]$ rather than on $[0,1]$ each time, which forces the $a_n$ down. Inserted
above; the rest of his proof is unchanged.
\end{ainote}

\begin{center}
\begin{tikzpicture}[>=Latex, scale=1.3]
  \draw[->] (-0.3,0) -- (3.2,0) node[right] {$t$};
  \draw[->] (0,-1.4) -- (0,1.4) node[above] {$y$};

  \draw[blue!80!black, thick] plot[domain=0.08:2.8, samples=250] (\x, {sin((1.0/(\x*0.1)) r)});

  \fill[red!80!black] (0,0) circle (2pt) node[below left, font=\small] {$\gamma(0)$};

  % The peaks of the plotted sin(10/x) sit at x = 20/(pi(1+4k)), k = 1,2,3,...
  % -- these must land ON the curve, since they are the points gamma(a_n).
  \fill[red!80!black] (1.273,1.0) circle (2pt) node[above right, font=\footnotesize] {$\gamma(a_1)$};
  \fill[red!80!black] (0.707,1.0) circle (1.7pt);
  \fill[red!80!black] (0.490,1.0) circle (1.4pt);
  \fill[red!80!black] (0.374,1.0) circle (1.1pt);
  \fill[red!80!black] (0.303,1.0) circle (0.9pt);

  % The limit (0,1) is NOT a point of S -- hence the hollow circle.
  \draw[red!80!black, thick] (0,1.0) circle (2.5pt);
  \node[red!80!black, font=\footnotesize, left] at (-0.08,1.0) {$(0,1)$};
  \draw[red!80!black, ->, dashed] (1.15,1.19) -- (0.13,1.07);

  \node[red!80!black, font=\small] at (1.75,-1.15) {$\gamma(a_n) \to (0,1) \neq (0,0) = \gamma(0)$};
\end{tikzpicture}
\end{center}

% Sonnet 5 (Medium)
We close the week with a change of subject: away from topology and back to the linear-algebraic
structure (\cref{sec:structured_spaces}) that lets us talk about angles, not just distances.



% Originally: content/week-03.tex, end of \subsection{Contraction mappings}

\begin{aiexercise}[Continuous image of a path-connected set]
\label{ex:ai_connected_image}
% Generator: Claude Opus 5 (High). Replaces an earlier exercise asking for
% \cref{thm:continuity_preserves_connected}, which now carries its own proof.
Let $f \colon X \to Y$ be a continuous map between metric spaces.
\begin{enumerate}[label=\textbf{(\alph*)}]
  \item Prove that if $X$ is path-connected (\cref{def:path_connected}), then so is $f(X)$.
  \item Conclude that the graph $\Gamma_g := \{(x,g(x)) : x \in U\}$ of a continuous
    $g : U \to \mathbb{R}^m$ is path-connected whenever $U \subseteq \mathbb{R}^n$ is convex.
    (Compare \cref{ex:4.2}, which asks the same question with \qt{connected} in place
    of \qt{path-connected}.)
\end{enumerate}
\end{aiexercise}

% --- exercise relocated here from the dissolved sheet 4 ---

% Originally: content/week-04.tex, \exercisesheet{4}
% Source: exercises/Ex4_Analysis2_eng.pdf

\begin{exercise}[4.2 --- Connected graphs \textnormal{(important)}]
\label{ex:4.2}
Let $U \subset \mathbb{R}^n$ be open and connected and let $f \in C^1(U, \mathbb{R}^m)$. Show that
its graph
\[\Gamma_f := \{(x, f(x)) : x \in U\}\]
is a connected subset of $\mathbb{R}^n \times \mathbb{R}^m$.
\end{exercise}

% Originally: the \section{Solutions} of content/week-03.tex

\section{Solutions}
\label{sec:connectedness_solutions}

% Transition: Claude Opus 5 (High)
Not every solution in this chapter is collected here. \Cref{ex:connectedness_tf}
(\cpageref{ex:connectedness_tf}) and \cref{ex:disk_connected} (\cpageref{ex:disk_connected}) are
solved where they appear.

\begin{exercisesolution}[Proof of \cref{ex:ai_connected_image}]
% Generator: Claude Opus 5 (High)
\begin{enumerate}[label=\textbf{(\alph*)}]
  \item Let $p, q \in f(X)$, say $p = f(x)$ and $q = f(y)$. Since $X$ is path-connected there is
    a path $\gamma : [0,1] \to X$ with $\gamma(0)=x$ and $\gamma(1)=y$. Then
    $f \circ \gamma : [0,1] \to f(X)$ is continuous as a composition of continuous maps, and it
    satisfies $(f\circ\gamma)(0) = p$, $(f\circ\gamma)(1) = q$. So it is a path in $f(X)$ from
    $p$ to $q$, and $f(X)$ is path-connected.
  \item A convex set $U$ is path-connected: the straight segment $t \mapsto (1-t)x + ty$ lies in
    $U$ and is continuous. The map $\varphi : U \to \mathbb{R}^n\times\mathbb{R}^m$,
    $\varphi(x) := (x,g(x))$, is continuous because each component is, and
    $\Gamma_g = \varphi(U)$. Apply \textbf{(a)}.
\end{enumerate}

\begin{remark}
Part \textbf{(a)} is strictly easier than its connectedness counterpart
(\cref{thm:continuity_preserves_connected}), because path-connectedness is a statement one
verifies by \emph{exhibiting} something, namely a path, and continuous maps push paths forward
directly. Connectedness has no such witness to push forward, which is why its proof has to run
by contradiction through preimages instead. This asymmetry is worth noticing: it is the reason
path-connectedness is usually the easier of the two to prove, and connectedness the easier of
the two to \emph{use}.
\end{remark}
\end{exercisesolution}

% ===========================================================================
% Official sheet 3 (3.1--3.11). Derived independently, then checked against
% exercises/Sol3_Analysis2_eng.pdf.
% ===========================================================================

\begin{exercisesolution}[Proof of \cref{ex:4.2}]
% Generator: Claude Sonnet 5 (Medium)
Consider the map
\[
\varphi : U \to \mathbb{R}^n \times \mathbb{R}^m, \qquad x \mapsto (x, f(x)).
\]
Since $f \in C^1(U,\mathbb{R}^m)$ is in particular continuous, and the identity on $U$ is
continuous, $\varphi$ is continuous. By definition $\Gamma_f = \varphi(U)$. Since $U$ is connected
and $\varphi$ is continuous, \cref{thm:continuity_preserves_connected} gives that $\varphi(U) =
\Gamma_f$ is connected.
\end{exercisesolution}

\begin{exercisesolution}[Solution of \cref{ex:closed_unbounded_not_simply_connected}]
% Generator: Claude Opus 5 (High)
Take the closed exterior of the unit disc,
\[U := \{x \in \mathbb{R}^2 \mid \lvert x\rvert \ge 1\}.\]
\begin{itemize}
  \item \textbf{Closed:} $U$ is the preimage of the closed set $[1,\infty)$ under the continuous map
    $x \mapsto \lvert x\rvert$.
  \item \textbf{Unbounded:} it contains $(k,0)$ for every $k \in \mathbb{N}$.
  \item \textbf{Connected:} it is in fact path-connected. Given $p,q \in U$, travel radially outwards
    from $p$ to the circle of radius $R := \max\{\lvert p\rvert, \lvert q\rvert\}$, then along that
    circle, then radially inwards to $q$. Every point on this path has norm between $1$ and $R$, so
    the path stays in $U$.
  \item \textbf{Not simply connected:} the loop $\gamma(t) := (2\cos 2\pi t,\, 2\sin 2\pi t)$ winds
    once around the origin, and no loop with winding number $1$ about a point outside $U$ can be
    contracted to a point within $U$. Contracting it would force the loop to sweep across the
    removed open disc $\{\lvert x\rvert < 1\}$.
\end{itemize}
\begin{ainote}
The exercise is really a test of whether the three notions have been kept apart. Closed-and-unbounded
rules out the reflex answer of an annulus $\{1 \le \lvert x\rvert \le 2\}$, and \qt{connected but not
simply connected} rules out a half-plane or all of $\mathbb{R}^2$. Deleting an \emph{open} disc
rather than a point is what keeps the result closed while still puncturing it.
\end{ainote}
\end{exercisesolution}

\begin{exercisesolution}[Proof of \cref{ex:sphere_path_connected}]
% Generator: Claude Opus 5 (High)
Both directions are proved by looking at what $S^{n-1}$ actually is in the two cases.

\textbf{The case $n = 1$.} Here $S^0 = \{x \in \mathbb{R} \mid x^2 = 1\} = \{-1, +1\}$, a two-point
space. Suppose $\gamma : [0,1] \to S^0$ were a path with $\gamma(0) = -1$ and $\gamma(1) = 1$.
Regarded as a map into $\mathbb{R}$, it is continuous on an interval, so by the intermediate value
theorem it takes the value $0$ somewhere. But $0 \notin S^0$, a contradiction. Hence $S^0$ is not
path-connected.

\textbf{The case $n \ge 2$.} Fix $p, q \in S^{n-1}$, and suppose first that $q \neq -p$. Then the
straight segment between them misses the origin: if $(1-t)p + tq = 0$ for some $t \in [0,1]$, taking
norms gives $(1-t) = t$, so $t = \tfrac12$ and $p = -q$, which we excluded. We may therefore
normalise the segment and set
\[\gamma(t) := \frac{(1-t)p + tq}{\lvert (1-t)p + tq\rvert}, \qquad t \in [0,1].\]
This is continuous, takes values in $S^{n-1}$, and satisfies $\gamma(0) = p$, $\gamma(1) = q$.

It remains to handle $q = -p$, and this is exactly where $n \ge 2$ enters: the sphere $S^{n-1}$ then
contains a point $r$ with $r \neq \pm p$. (For $n \ge 2$ the sphere is infinite --- pick any unit
vector $v$ orthogonal to $p$, which exists because $p^\perp$ has dimension $n - 1 \ge 1$.) Neither
$r$ nor $q$ is the antipode of the other's partner, so the previous construction gives a path from
$p$ to $r$ and a path from $r$ to $q$; concatenating them joins $p$ to $q$. Hence $S^{n-1}$ is
path-connected.
\begin{remark}
The construction in the main case is the radial projection of the segment, and it is the standard way
to move around a sphere: it works verbatim on any $S^{n-1}$ and explains why the antipodal pair is
the only obstruction. Note also that $n \ge 2$ is used only once, to produce a third point --- which
is precisely what $S^0$ does not have.
\end{remark}
\begin{ainote}
The official solution introduces the same construction with the words \qt{now let $n \ge 1$}, and
then disposes of the antipodal case $x = -y$ by composing two paths through \qt{an arbitrary other
point $w \in S^{n-1}$}. For $n = 1$ there is no such $w$: $S^0 = \{\pm1\}$ consists of $x$ and $-x$
and nothing else. The bound is therefore $n \ge 2$, exactly as the statement being proved requires
--- and as the official solution's own preceding paragraph, which shows $S^0$ is \emph{not}
path-connected, confirms. The step is written out explicitly above rather than left to the reader,
since it is the only place the hypothesis is used.
\end{ainote}
\end{exercisesolution}

\begin{exercisesolution}[Proof of \cref{ex:examHS24_TF_topology}]
% Generator: Gemini 3.1 Pro (High)
\begin{enumerate}[label=\textbf{(\alph*)}]
  \item \textbf{True:} $U$ is defined by non-strict inequalities involving continuous functions (specifically, the pre-image of $[0, \infty)$ under $g_1(x,y) = y+x^2$ and $g_2(x,y) = x^2-y$). Consequently, $U$ is closed.
  \item \textbf{True:} We can contract $U$ to the origin via the homotopy $H((x,y), t) = (tx, t^2y)$ for $t \in [0,1]$. Note that $(tx)^2 = t^2x^2$, and since $-x^2 \le y \le x^2$, multiplying by $t^2 \ge 0$ gives $-(tx)^2 \le t^2y \le (tx)^2$, which ensures the contraction stays within $U$.
  \item \textbf{False:} $U$ is not bounded (for example, the sequence $(n,0)$ is in $U$ for all $n \in \mathbb{N}$), and therefore not compact.
  \item \textbf{True:} Since $U$ is path-connected (from the homotopy to the origin) and therefore connected, its continuous image $f(U)$ is connected in $\mathbb{R}$.
  \item \textbf{False:} The function is not bounded on $U$. Consider the path $(x,0)$ inside $U$; here $f(x,0) = x^2$, which grows without bound as $x \to \infty$.
  \item \textbf{True:} Since $f(x,y) \ge 0$ everywhere on $U$ and $f(0,0) = 0$ with $(0,0) \in U$, the infimum is $0$ and it is attained at the origin.
\end{enumerate}
\end{exercisesolution}

\begin{exercisesolution}[Solution to \cref{ex:box_intersection_disconnected_union_connected}]
% Generator: Claude Opus 5 (High) --- the construction below is the official key's, with the
% strip widened from half-width 1/10 to 1/4 so that the figure is legible; every verification is
% written for these notes.
Take an open annulus and an open vertical strip through its hole:
\[U_1 := B_1(0) \setminus \overline{B_{1/2}(0)}
  = \big\{p \in \mathbb{R}^2 \ \big|\ \tfrac12 < \lvert p\rvert < 1\big\}, \qquad
  U_2 := \big(-\tfrac14,\tfrac14\big)\times\mathbb{R}.\]

% Generator: Claude Opus 5 (High) --- FIG-MOCK-19
% Geometry: two concentric circles of radii 1 and 1/2 (drawn at 2.2 cm and 1.1 cm per unit),
% and a vertical strip |x| < 1/4 (drawn at +-0.55 cm) running past both of them. All four
% boundaries are DASHED, because both sets are open and neither contains its boundary. The
% annulus is shaded blue, the strip orange, and the two pieces of the intersection purple on
% top of both. The purple pieces are two lenses, one above and one below the x-axis, and the
% whole point of the picture is that they do not touch.
% Arithmetic: on U_1 \cap U_2 we have |x| < 1/4 and |p| > 1/2, so
%   y^2 > 1/4 - x^2 >= 1/4 - 1/16 = 3/16,  i.e. |y| > sqrt(3)/4 = 0.433.
% At the drawing scale that is 0.433 * 2.2 = 0.953 cm, which is where each lens must stop:
% it is sqrt(1.1^2 - 0.55^2) = 1.1 * sqrt(3)/2 = 0.953 cm, the height at which the strip edge
% x = 0.55 crosses the inner circle. The two agree, as they must.
\begin{center}
\begin{tikzpicture}[>=Latex]
  % --- U_1, the annulus: fill the outer disc, then punch out the inner one ---
  \fill[blue!12] (0,0) circle (2.2);
  \fill[white] (0,0) circle (1.1);
  % --- U_2, the strip (covers the annulus shading where they meet) ---
  \fill[orange!14] (-0.55,-2.9) rectangle (0.55,2.9);
  % --- the intersection, painted back on top of both: purple over the outer disc inside the
  %     strip, then orange restored over the hole, which belongs to U_2 alone ---
  \begin{scope}
    \clip (-0.55,-2.9) rectangle (0.55,2.9);
    \fill[purple!40] (0,0) circle (2.2);
    \fill[orange!14] (0,0) circle (1.1);
  \end{scope}
  % --- boundaries, all dashed: both sets are open ---
  \draw[blue!70!black, dashed, thick] (0,0) circle (2.2);
  \draw[blue!70!black, dashed, thick] (0,0) circle (1.1);
  \draw[orange!85!black, dashed, thick] (-0.55,-2.9) -- (-0.55,2.9);
  \draw[orange!85!black, dashed, thick] (0.55,-2.9) -- (0.55,2.9);
  % --- labels ---
  % label placed inside the blue band: radius 1.65 at 135 degrees, between 1.1 and 2.2
  \node[blue!70!black, font=\small] at (-1.17,1.17) {$U_1$};
  \node[orange!85!black, font=\small, anchor=south] at (0,2.95) {$U_2$};
  \node[purple!75!black, font=\small, align=left, anchor=west] at (2.7,0)
    {$U_1\cap U_2$:\\two components};
  \draw[purple!75!black, ->] (2.65,0.35) to[out=150,in=20] (0.62,1.6);
  \draw[purple!75!black, ->] (2.65,-0.35) to[out=210,in=340] (0.62,-1.6);
  \draw[gray, dotted] (-2.6,0) -- (2.4,0);
\end{tikzpicture}
\end{center}

Four things have to be checked, and the figure does none of them by itself.

\begin{enumerate}[label=\textbf{(\alph*)}]
  \item \textbf{$U_1$ is open and connected.} It is the intersection of the open set $B_1(0)$
    with the complement of the closed set $\overline{B_{1/2}(0)}$, hence open. It is
    path-connected, since any two of its points can be joined by moving radially to the circle of
    radius $\tfrac34$ and then along that circle, and a path-connected set is connected.
  \item \textbf{$U_2$ is open and connected.} It is a product of an open interval with
    $\mathbb{R}$, hence open, and it is convex, hence path-connected and so connected.
  \item \textbf{$U_1\cap U_2$ is not connected.} If $(x,y)$ lies in it then $\lvert x\rvert <
    \tfrac14$ and $x^2+y^2 > \tfrac14$, so
    \[y^2 > \tfrac14 - x^2 > \tfrac14 - \tfrac1{16} = \tfrac3{16} > 0 .\]
    In particular $y \neq 0$ throughout, so the two open half-planes $\{y>0\}$ and $\{y<0\}$
    separate $U_1\cap U_2$. Both pieces are non-empty, since $(0,\tfrac34)$ and $(0,-\tfrac34)$
    both lie in $U_1\cap U_2$. That is a separation, so the intersection is disconnected.
  \item \textbf{$U_1\cup U_2$ is connected.} Both sets are connected and they share the point
    $(0,\tfrac34)$, so \cref{cor:union_connected_common_point} applies to the two-element family
    $\{U_1,U_2\}$.
\end{enumerate}

\begin{remark}[Why the intersection is the odd one out]
\Cref{cor:union_connected_common_point} makes unions easy, and nothing in this chapter says
anything at all about intersections. That silence is not an oversight. A union of connected sets
sharing a point cannot be split, because any separation would have to split one of the pieces;
an intersection carries no such argument, and the picture above shows why: the strip enters the
annulus twice, and what it takes away in between is exactly what would have joined the two
pieces.

Note also that the construction needs the hole. If $U_1$ were a disc, then $U_1 \cap U_2$ would
be convex and the example would collapse; the failure of connectedness in \textbf{(c)} is the
annulus's non-simple-connectedness showing up in a second guise.
\end{remark}
\end{exercisesolution}



\part{Differential Calculus}

% Generator: Claude Opus 5 (High) --- part-level framing prose, new content.
In one variable the derivative is a number, which makes it tempting to expect the several-variable
derivative to be a list of numbers. It is not. The right object is a linear map, the one that best
approximates $f$ near a point, and the partial derivatives are merely its matrix in the standard
basis. Getting this order of ideas right is what turns the chain rule into a statement about
composing linear maps rather than a formula to be memorised. This part sets up that differential,
works out how to compute with it, and closes with Taylor's theorem, which sharpens the linear
approximation to any order and supplies the quantitative estimates the optimization theory of the
next part runs on.

% ============================================================
% Chapter 9 --- The differential and the Jacobi matrix
% Originally: content/week-04.tex, \section{The differential} and
%             \section{The Jacobi matrix} (Corsin Week 4, Monday)
% ============================================================
\chapter{The Differential and the Jacobi Matrix}
\label{ch:differential}

% Transition: Claude Opus 5 (Medium)
The previous part treated continuous functions as structure-preserving maps between metric
spaces. This one asks for more: that a map be well approximated, near each point, by a
\emph{linear} one. That is the differential, and this chapter sets out the hierarchy it sits in
--- Fréchet differentiability, Gâteaux derivatives, partial derivatives --- together with the
matrix that represents it.

% Originally: content/week-04.tex, \section{The differential} (Corsin Week 4, Monday)

% Source: Corsin Nick/Class Notes/Week 4.pdf, p. 2
\section{The differential}
\label{sec:the_differential}

For $f : \mathbb{R}\to\mathbb{R}$ differentiable at $x_0 \in \mathbb{R}$:
\[f(x_0+h) = f(x_0) + f'(x_0)h + o(h)\]
\[\Updownarrow\]
\[\lim_{h\to 0}\frac{f(x_0+h)-f(x_0)}{h} = f'(x_0)\]
\[\Updownarrow\]
\[\lim_{h\to 0}\left|\frac{f(x_0+h)-f(x_0)-f'(x_0)h}{h}\right| = 0\]

For $F : \mathbb{R}^n \to \mathbb{R}^m$ differentiable at $x_0 \in \mathbb{R}^n$:
\[F(x_0+h) = F(x_0) + \underbrace{DF_{x_0}(h)}_{\text{linear map}} + o(|h|)\]
\[\Updownarrow\]
\[\lim_{h\to 0}\frac{|F(x_0+h)-F(x_0)-DF_{x_0}(h)|}{|h|} = 0\]
\[\Downarrow\]
\[\lim_{s\to 0}\frac{F(x_0+sv)-F(x_0)}{s} = DF_{x_0}(v) = \partial_vF(x_0) \quad \text{for each fixed } v \in \mathbb{R}^n\]
\[\text{(the \newterm{directional derivative} of } F \text{ at } x_0 \text{ in direction } v\text{)}\]

Note that the last arrow points \textbf{one way only}. The first two lines say the same thing;
the third is a strictly weaker consequence, and recovering differentiability from it is exactly
what fails in several variables --- see \cref{ex:partial_derivatives_not_continuous} below.

% Transition: Claude Opus 5 (High)
It is worth pausing on why the definition had to be restated rather than merely copied. The
one-variable derivative is defined as a \emph{quotient}, and that is precisely what does not
survive: for $h \in \mathbb{R}^n$ with $n \geq 2$ the expression
$\big(F(x_0+h)-F(x_0)\big)/h$ is meaningless, because one cannot divide by a vector. The first
line above is the version that does survive, and it is the one to remember: \textbf{$F$ is
differentiable at $x_0$ if it is approximated near $x_0$ by an affine map, to an error smaller
than $|h|$.} Everything in this chapter --- the Jacobian, the chain rule, Taylor expansions ---
is bookkeeping for that single idea, and the reason the derivative becomes a \emph{linear map}
rather than a number is simply that a linear map is what an approximation in $n$ variables looks
like.

\begin{ainote}
Corsin's middle line reads
$\lim_{h\to0}\frac{F(x_0+h)-F(x_0)}{|h|} = DF_{x_0}\!\left(\frac{h}{|h|}\right)$, inside a chain
of $\Updownarrow$. As written that limit does not exist: $h$ is being sent to $0$ on the left
while still appearing on the right, and the quotient genuinely depends on the direction of
approach. The intended statement is the one typeset above --- fix a direction $v$ first, then
let the scalar $s \to 0$ --- and it is only an \emph{implication}, not an equivalence. The
distinction is not pedantic: it is the entire content of
\Cref{ex:partial_derivatives_not_continuous} and of
\Cref{rem:directional_derivative_caveat}.
\end{ainote}

% Source: Corsin Nick/Class Notes/Week 4.pdf, p. 2
% (Re-cited: the paragraph and AI-Note above are additions, not transcription.)
\begin{definition}[Differentiability]
\label{def:differentiable}
We say that $F : \mathbb{R}^n \to \mathbb{R}^m$ is \newterm{differentiable}
(\germanterm{differenzierbar}) at $x_0$ if a linear map $DF_{x_0}$ as above exists, where
$x_0, h \in \mathbb{R}^n$. The same definition applies verbatim to $F : U \to \mathbb{R}^m$
for any open subset $U \subseteq \mathbb{R}^n$, and we use it in that generality throughout.
\end{definition}

% Sonnet 5 (Medium)
\begin{remark}
The linear map $DF_{x_0}$ in \cref{def:differentiable}, when it exists, is automatically unique:
if $L_1, L_2$ both satisfy the defining limit, then
$(L_1-L_2)(h)/|h| \to 0$ as $h \to 0$ along every direction, which forces the linear map
$L_1 - L_2$ to vanish identically. Concretely: write $L := L_1-L_2$ and suppose $L(h) \neq 0$
for some $h$. Along the ray $th$, $t > 0$, linearity gives
\[\frac{|L(th)|}{|th|} = \frac{t\,|L(h)|}{t\,|h|} = \frac{|L(h)|}{|h|} \neq 0,\]
a nonzero \emph{constant} --- so the quotient cannot tend to $0$ as $t \to 0$. Hence $L(h)=0$
for every $h$. This
is the multivariable analogue of the uniqueness-of-limits argument from
\Cref{prop:uniqueness_of_limits}.
\end{remark}

\begin{proposition}[Differentiability implies continuity]
\label{prop:differentiable_implies_continuous}
If $F : U \to \mathbb{R}^m$, $U \subseteq \mathbb{R}^n$ open, is differentiable
(\cref{def:differentiable}) at $x_0 \in U$, then $F$ is continuous at $x_0$.
\end{proposition}

\begin{proof}
By \cref{def:differentiable}, $F(x_0+h) - F(x_0) = DF_{x_0}(h) + o(|h|)$ as $h \to 0$. Since
$DF_{x_0}$ is linear on the finite-di\-men\-sional space $\mathbb{R}^n$, it is bounded, i.e.\ there
exists $C \geq 0$ with $|DF_{x_0}(h)| \leq C|h|$ for all $h$. Hence
\[|F(x_0+h) - F(x_0)| \leq |DF_{x_0}(h)| + o(|h|) \leq C|h| + o(|h|) \xrightarrow{h\to 0} 0,\]
so $F(x_0+h) \to F(x_0)$ as $h \to 0$, i.e.\ $F$ is continuous at $x_0$.
\end{proof}

\begin{remark}
The converse fails: continuity does not imply differentiability, exactly as in one variable
(e.g.\ $x \mapsto |x|$ at $0$). \Cref{ex:partial_derivatives_not_continuous} below shows the
gap is even wider in several variables --- \emph{all} partial derivatives can exist at a point
without $F$ even being continuous there, let alone differentiable.
\end{remark}

% Originally: content/week-04.tex, \section{The Jacobi matrix} (Corsin Week 4, Monday)

% Source: Corsin Nick/Class Notes/Week 4.pdf, p. 3
\section{The Jacobi matrix}
\label{sec:jacobi_matrix}

Since $DF_{x_0} : \mathbb{R}^n \to \mathbb{R}^m$ (\cref{def:differentiable}) is a linear map, there is a \textbf{1--1
correspondence with a matrix once a basis is fixed}! In the standard basis of $\mathbb{R}^n$,
$(e_1,\dots,e_n)$, this is the \newterm{Jacobi matrix} \germanfor{Jacobi-Matrix}:
\[DF_x \;\cong\; \Jac F(x) = \begin{pmatrix} \dfrac{\partial F_1}{\partial x_1} & \cdots & \dfrac{\partial F_1}{\partial x_n} \\[1ex] \vdots & \ddots & \vdots \\[1ex] \dfrac{\partial F_m}{\partial x_1} & \cdots & \dfrac{\partial F_m}{\partial x_n}\end{pmatrix}\]
($m$ rows, $n$ columns), where we wrote $F(x) = (F_1, \dots, F_m)\transp$ and
\[\frac{\partial F}{\partial x_i} =: \partial_i F(x)\]
is the \newterm{$i$-th partial derivative} \germanfor{partielle Ableitung} at $x$:
\[\partial_i F(x) = \lim_{s\to 0}\frac{F(x+se_i)-F(x)}{s}.\]
This defines a function $\partial_i F : U \to \mathbb{R}^m$. If $F$ is differentiable, then
\[\partial_i F(x) = DF_x(e_i).\]

\begin{importantremark}
The partial derivatives can exist even if $F : U \to \mathbb{R}^m$ is not differentiable (recall \cref{ex:partial_derivatives_not_continuous})!
\end{importantremark}

% Sonnet 5 (Medium)
\begin{aiexample}[All partial derivatives exist, yet $f$ is not even continuous]
\label{ex:partial_derivatives_not_continuous}
% Generator: Claude Sonnet 5 (Medium)
Let $f : \mathbb{R}^2 \to \mathbb{R}$ be defined by
\[f(x,y) := \begin{cases} \dfrac{xy}{x^2+y^2} & (x,y) \neq (0,0), \\ 0 & (x,y) = (0,0). \end{cases}\]
Along each axis, $f$ vanishes identically ($f(x,0) = f(0,y) = 0$), so both partial derivatives at
the origin exist and equal $0$:
\[\partial_x f(0,0) = \lim_{s\to 0}\frac{f(s,0)-f(0,0)}{s} = 0, \qquad \partial_y f(0,0) = 0.\]
Yet $f$ is not even continuous at $(0,0)$: along the diagonal $y=x$, $f(x,x) = \tfrac{x^2}{2x^2} =
\tfrac12$ for all $x \neq 0$, so $f(x,x) \to \tfrac12 \neq 0 = f(0,0)$ as $x \to 0$. By
\Cref{prop:differentiable_implies_continuous}, $f$ is therefore not differentiable at $(0,0)$,
even though both partial derivatives exist there.
\end{aiexample}

% Supplement: Sascha Brack/Class Notes/Week_03_Notes_Monday.pdf, p. 7
\begin{remark}[The same expression, squared, behaves completely differently]
\label{rem:continuity_contrast_pair}
Sascha Brack pairs the function above with one that looks almost identical, as a test of whether
you can tell them apart:
\[f_1(x,y) := \frac{xy}{x^2+y^2}, \qquad\qquad f_2(x,y) := \frac{x^2y^2}{x^2+y^2},\]
both extended by $0$ at the origin. Substituting $x = r\cos\theta$, $y = r\sin\theta$ settles
both at once:
\[f_1 = \frac{r^2\cos\theta\sin\theta}{r^2} = \cos\theta\sin\theta,
\qquad\qquad
f_2 = \frac{r^4\cos^2\theta\sin^2\theta}{r^2} = r^2\cos^2\theta\sin^2\theta.\]
In $f_1$ every $r$ cancels: the function is constant along each ray and takes a different
constant on different rays, so no limit exists and $f_1$ is discontinuous at the origin. In
$f_2$ a factor $r^2$ survives, and since $|\cos^2\theta\sin^2\theta| \leq 1$ we get
$|f_2| \leq r^2 \to 0$ \emph{uniformly in $\theta$}, so $f_2$ is continuous at the origin.

The pair is worth doing once because it shows what the polar substitution is actually testing.
The question is never \qt{does the expression go to zero along each fixed ray?}, which $f_1$ does
not even manage. The question is \textbf{does the $\theta$-dependence get killed by a positive
power of $r$?} If the $r$'s cancel completely, the direction survives into the limit and there is
none.
This is the same test used in \cref{q:quiz_q4} on the repetition quiz, and the reason
step 3 of \cref{rem:differentiability_recipe} insists on \qt{uniformly in $\theta$}.
\end{remark}

% Sonnet 5 (Medium)
\begin{center}
\begin{tikzpicture}[>=Latex, scale=1.0]
  \draw[->] (-1.6,0) -- (1.6,0) node[right, font=\small] {$x$};
  \draw[->] (0,-1.6) -- (0,1.6) node[above, font=\small] {$y$};
  \draw[orange!90!black, thick, ->] (0,0) -- (1.3,0) node[below right, font=\scriptsize] {$f\equiv0$};
  \draw[orange!90!black, thick, ->] (0,0) -- (0,1.3) node[above left, font=\scriptsize] {$f\equiv0$};
  \draw[purple, thick, ->] (0,0) -- (1.1,1.1) node[above right, font=\scriptsize] {$f\equiv\tfrac12$};
  % Generator: Claude Opus 5 (High) --- caption added; the picture previously stood with no
  % text of its own, several paragraphs away from the example it illustrates.
  \node[below, font=\footnotesize, align=center] at (0,-1.75)
    {$f = \tfrac{xy}{x^2+y^2}$ near the origin: constant along every ray,\\
     but a \emph{different} constant on the diagonal than on the axes};
\end{tikzpicture}
\end{center}

% Transition: Claude Opus 5 (High) --- the diagram below arrived with no text at all around
% it, between a figure and a \subsection heading.
Sascha Brack's notes collect these relationships into one picture. Read it as a hierarchy: the
strongest condition is at the bottom, and every arrow points at something weaker.

% Supplement: Sascha Brack/Class Notes/Week_04_Notes_Monday_Updated.pdf, p. 4
\begin{center}
\begin{tikzcd}[row sep=3.5em, column sep=1.5em]
  & \begin{tabular}{c}\textbf{Differentiable}\\($Df_{x_0}$ exists)\end{tabular} \arrow[dl, Rightarrow] \arrow[dr, Rightarrow] & \\
  \begin{tabular}{c}\textbf{Partial derivatives}\\exist\end{tabular} & & \begin{tabular}{c}\textbf{Directional derivatives}\\exist\end{tabular} \\
  \begin{tabular}{c}\textbf{Partial derivatives}\\exist\\and are continuous\end{tabular} \arrow[u, Rightarrow] \arrow[uur, Rightarrow] \arrow[rr, Leftrightarrow] & & \begin{tabular}{c}\textbf{Directional derivatives}\\exist\\and are continuous\end{tabular} \arrow[u, Rightarrow] \arrow[uul, Rightarrow]
\end{tikzcd}
\end{center}

% Generator: Claude Opus 5 (High)
\begin{importantremark}[Not one of these arrows reverses]
The diagram is easy to misread as a set of equivalences, and it is nothing of the kind: every
$\Rightarrow$ in it is strict. \Cref{ex:partial_derivatives_not_continuous} already shows the
bottom-left arrow failing backwards, since there both partial derivatives exist at the origin
while $f$ is not even continuous. The other two failures, with their counterexamples, are collected in
\cref{sec:sufficient_condition_differentiability}, which redraws this same landscape as a single
vertical chain and attaches a witness to each broken link.

The one genuine equivalence is the horizontal arrow along the bottom, and it is worth seeing
why. Left to right: continuous partials make $f$ continuously differentiable
(\cref{thm:continuous_partials_imply_differentiable}), so every directional derivative exists
and equals $Df_x(v)$, which depends continuously on $x$. Right to left is immediate: take
$v = e_i$.
\end{importantremark}

\subsection{Examples}

% Transition: Claude Opus 5 (High)
The three examples that follow are the three shapes the Jacobian takes, and it is worth naming
them, because the layout of the matrix changes completely between them while the underlying
object, the linear map $DF_{x_0}$, does not.
\begin{itemize}
  \item \textbf{A curve} ($n=1$, $m$ arbitrary): one column, so $\Jac\gamma(t)$ is a
    \emph{column vector}, the velocity.
  \item \textbf{A scalar field} ($n$ arbitrary, $m=1$): one row, so $\Jac f(x)$ is a
    \emph{row vector}, namely the transposed gradient $(\nabla f)\transp$.
  \item \textbf{A vector field} (both $>1$): the full $m \times n$ matrix, whose columns are the
    partial derivatives and whose rows are the gradients of the components.
\end{itemize}
That the Jacobian of a scalar field is a row while the Jacobian of a curve is a column is not a
convention to memorise: it follows from \qt{$m$ rows, $n$ columns} and nothing else. Note the
transpose in the second bullet, which is the one place a real convention enters: the
\emph{gradient} $\nabla f$ is a column vector, so it is $\Jac f = (\nabla f)\transp$ that is the
row. Confusing the two is the commonest source of shape errors when the chain rule is applied
later in this chapter.

% Source: Corsin Nick/Class Notes/Week 4.pdf, p. 4
\begin{example}[Derivative of a curve]
\[\gamma : \mathbb{R}\to\mathbb{R}^2, \qquad t \mapsto \begin{pmatrix}\cos t \\ \sin t\end{pmatrix}\]
What is $\gamma'(\tfrac{\pi}{4}) := D\gamma_{\pi/4}$?
\[\gamma'\!\left(\tfrac{\pi}{4}\right) = \begin{pmatrix}\partial_t\gamma_1(\pi/4) \\ \partial_t\gamma_2(\pi/4)\end{pmatrix} = \begin{pmatrix}-1/\sqrt{2} \\ 1/\sqrt{2}\end{pmatrix}\]
\end{example}

\begin{aiexample}[Jacobian Matrix of Polar Coordinates]
% Generator: Gemini 3.6 Flash (Medium)
Let $f \colon (0,\infty) \times (0,2\pi) \to \mathbb{R}^2$ be the polar coordinate map $f(r,\theta) := (r\cos\theta, r\sin\theta)\transp$. The Jacobian matrix $\Jac f(r,\theta)$ is:
\[
\Jac f(r,\theta) = \begin{pmatrix} \dfrac{\partial f_1}{\partial r} & \dfrac{\partial f_1}{\partial \theta} \\[1ex] \dfrac{\partial f_2}{\partial r} & \dfrac{\partial f_2}{\partial \theta} \end{pmatrix} = \begin{pmatrix} \cos\theta & -r\sin\theta \\ \sin\theta & r\cos\theta \end{pmatrix}.
\]
Its determinant is $\det \Jac f(r,\theta) = r\cos^2\theta + r\sin^2\theta = r > 0$.
\end{aiexample}

\begin{center}
\begin{tikzpicture}[>=Latex, scale=1.3]
  \draw[->] (-1.4,0) -- (1.4,0) node[right] {$x$};
  \draw[->] (0,-1.4) -- (0,1.4) node[above] {$y$};
  \draw[thick] (0,0) circle (1.0);
  \coordinate (P) at (0.707, 0.707);
  \draw[dotted, orange!90!black, thick] (0,0) -- (P);
  \fill[orange!90!black] (P) circle (2pt) node[above right, font=\footnotesize] {$\gamma(\pi/4)$};
  \draw[orange!90!black, ultra thick, ->] (P) -- ++(-0.6, 0.6) node[above left, font=\small] {$\gamma'(\pi/4)$};
\end{tikzpicture}
\end{center}

This vector is the \textbf{velocity of the curve} \germanfor{Geschwindigkeit} at
$t = \tfrac{\pi}{4}$; its length $|\gamma'(t)|$ is the \textbf{speed}. We get that
\[\gamma\!\left(\tfrac{\pi}{4}+h\right) = \frac{1}{\sqrt{2}}\begin{pmatrix}1\\1\end{pmatrix} + \frac{h}{\sqrt{2}}\begin{pmatrix}-1\\1\end{pmatrix} + o(|h|).\]

\begin{ainote}
Corsin calls $\gamma'(\tfrac{\pi}{4})$ the \qt{speed} of the curve. Strictly, $\gamma'(t)$ is the
\emph{velocity}, a vector carrying both a direction and a magnitude, whereas the speed is the
scalar $|\gamma'(t)|$, here $|(-1,1)|/\sqrt2 = 1$. Since the display below uses the vector, the
distinction matters. Corrected above.
\end{ainote}

% Source: Corsin Nick/Class Notes/Week 4.pdf, p. 4
From this example, we see that the \textbf{column} vectors of the Jacobian are the partial
derivatives:
\[\Jac F(x) = \left(\frac{\partial F}{\partial x_1}\ \Big|\ \cdots\ \Big|\ \frac{\partial F}{\partial x_n}\right)\]
and therefore, in the canonical basis, with $v := (v_1,\dots,v_n)\transp$:
\[\partial_v F(x) = DF_x(v) = \Jac F(x)\,v = \sum_{i=1}^{n} v_i\,\frac{\partial F}{\partial x_i}.\]

% Source: Corsin Nick/Class Notes/Week 4.pdf, p. 5
\begin{example}[Scalar field gradient]
\[f : \mathbb{R}^3\to\mathbb{R}, \qquad (x,y,z) \mapsto 2x^2 + y^2 + 3z^2 - 2xyz\]
What is $\Jac f(x,y,z)$?
\[\Jac f(x,y,z) = (4x - 2yz,\ 2y - 2xz,\ 6z - 2xy) = \left(\frac{\partial f}{\partial x}, \frac{\partial f}{\partial y}, \frac{\partial f}{\partial z}\right) = (\nabla f)\transp\]
\end{example}

% Supplement: Sascha Brack/Class Notes/Week_04_Notes_Friday_Updated.pdf, p. 3
\begin{example}[Jacobian matrix of a vector field]
\label{ex:jacobian_vector_field}
Consider $f : \mathbb{R}^2 \to \mathbb{R}^3$ defined by
\[f(x,y) := \begin{pmatrix} x^2+y^2 \\ 2xy \\ e^x\sin(y) \end{pmatrix}.\]
The partial derivatives with respect to $x$ and $y$ are the column vectors:
\[\partial_1 f(x,y) = \begin{pmatrix} 2x \\ 2y \\ e^x\sin(y) \end{pmatrix}, \qquad \partial_2 f(x,y) = \begin{pmatrix} 2y \\ 2x \\ e^x\cos(y) \end{pmatrix}.\]
Because these partial derivatives exist everywhere and are continuous, $f$ is differentiable on all of $\mathbb{R}^2$. The differential $Df_{(x_0,y_0)}$ is represented by the $3 \times 2$ Jacobian matrix:
\[\Jac f(x_0,y_0) = \begin{pmatrix} 2x_0 & 2y_0 \\ 2y_0 & 2x_0 \\ e^{x_0}\sin(y_0) & e^{x_0}\cos(y_0) \end{pmatrix}.\]
\end{example}

From this, we see that the \textbf{rows} of the Jacobian are the \newterm{gradients}
\germanfor{Gradient}:
\[\Jac F(x) = \begin{pmatrix} \nabla F_1 \\ \hline \vdots \\ \hline \nabla F_m \end{pmatrix}\]

\begin{ainote}
Corsin's sentence reads ``the columns of the Jacobian are the gradients'', but the displayed
matrix stacks $\nabla F_1, \dots, \nabla F_m$ as rows. The display is the correct one, and it is
consistent with $(\nabla f)\transp$ just above. Corrected to ``rows''.
\end{ainote}

% Supplement: Sascha Brack/Class Notes/Week_04_Notes_Monday_Updated.pdf, pp. 10--11
\begin{example}[A function that is differentiable but not $C^1$]
\label{ex:differentiable_not_c1}
Let $f : \mathbb{R}^2 \to \mathbb{R}$ be defined by
\[
f(x,y) := \begin{cases}
  (x^2+y^2)\sin\left(\dfrac{1}{\sqrt{x^2+y^2}}\right), & \text{if } (x,y) \neq (0,0), \\
  0, & \text{if } (x,y) = (0,0).
\end{cases}
\]
We claim that $f$ is differentiable at $(0,0)$ with $Df_{(0,0)} = (0, 0)$, but the partial derivatives are not continuous at $(0,0)$. 

\textbf{Step 1: Differentiability at $(0,0)$.} We use the definition of the differential. Let $L = (0, 0)$. We need to show that
\[
  \lim_{h \to (0,0)} \frac{f(h_1, h_2) - f(0,0) - L\begin{pmatrix}h_1\\h_2\end{pmatrix}}{|h|} = 0.
\]
Plugging in the definitions and changing to polar coordinates $(h_1, h_2) = (r\cos\theta, r\sin\theta)$ with $r = |h| = \sqrt{h_1^2+h_2^2}$:
\[
  \lim_{r \to 0} \frac{r^2 \sin(1/r) - 0 - 0}{r} = \lim_{r \to 0} r \sin(1/r) = 0.
\]
The sine term is bounded between $-1$ and $1$, so the limit evaluates to $0$. Thus, $f$ is differentiable at $(0,0)$ and $Df_{(0,0)} = (0, 0)$, which implies $\partial_1 f(0,0) = 0$ and $\partial_2 f(0,0) = 0$.

\textbf{Step 2: Discontinuity of partial derivatives.} Let us compute $\partial_1 f(x,y)$ for $(x,y) \neq (0,0)$. Using the product and chain rules:
\[
  \partial_1 f(x,y) = 2x \sin\left(\frac{1}{\sqrt{x^2+y^2}}\right) - \frac{x}{\sqrt{x^2+y^2}} \cos\left(\frac{1}{\sqrt{x^2+y^2}}\right).
\]
Let us choose a sequence $(x_n, y_n) := (\frac{1}{n}, 0)$ which converges to $(0,0)$ as $n \to \infty$. Then
\[
  \lim_{n \to \infty} \partial_1 f\left(\frac{1}{n}, 0\right) = \lim_{n \to \infty} \left[ \frac{2}{n} \sin(n) - \cos(n) \right].
\]
The first term vanishes, but $\lim_{n \to \infty} \cos(n)$ does not exist. Thus, $\partial_1 f$ is not continuous at $(0,0)$, and so $f$ is not continuously differentiable ($C^1$) at $(0,0)$. (By symmetry, the same holds for $\partial_2 f$).
\end{example}

% Sonnet 5 (Medium)


% --- exercises relocated here from the dissolved sheet 4 ---

% Originally: content/week-04.tex, end of the chapter

\begin{aiexercise}[Differential and Linearization at $(1,0)$]
\label{ex:ai_differential_linearization}
% Generator: Gemini 3.6 Flash (Medium)
Let $f(x,y) := x^2 y + \sin(xy)$ defined on $\mathbb{R}^2$.
\begin{enumerate}[label=\textbf{(\alph*)}]
  \item Compute the gradient $\nabla f(x,y)$ and the total differential $Df(1,0)$.
  \item Find the linear approximation $L(x,y)$ of $f$ around the point $(1,0)$.
\end{enumerate}
\end{aiexercise}

% Quelle: Linus Lüchinger/Slides/Slides-03-19.pdf, slides 9--11
% Generator: Claude Opus 5 (High)
\subsection{Differentiating with respect to a matrix}
\label{sec:differentiating_wrt_matrix}

Nothing in \cref{def:differentiable} requires the domain to be $\mathbb{R}^n$ with $n$ small, or
even to be written as a column of numbers. A function of an $n\times m$ matrix is a function of
$nm$ real variables, and its partial derivatives $\partial f/\partial A_{ij}$ are defined exactly
as before, by varying one entry and freezing the rest.

% Reclassified from ainote to remark: motivation for the mathematics, not commentary on the
% transcription. Linus spends part of a session making this argument.
\begin{remark}[Why this is worth a subsection]
Taking derivatives with respect to a matrix argument can look like a formal exercise with no
purpose. It is not, and the reason is worth recording: \textbf{it is how neural networks are
trained.} The parameters of such a network are matrices. The function being minimised sends those
matrices to a single number measuring how badly the network fits the data, and training consists
of computing the gradient with respect to those matrices and then stepping a little way against
it, because $-\nabla f$ points in the direction of steepest decrease. Every step of that loop is
the material of this chapter.
\end{remark}

The simplest instance is already worth doing by hand, and it is exactly linear regression.

% Generator: Claude Opus 5 (High)
\begin{aiexample}[The gradient of a least-squares fit]
\label{ex:least_squares_matrix_gradient}
Fix data matrices $X \in \mathbb{R}^{m\times k}$ (the $k$ inputs, each a column in
$\mathbb{R}^m$) and $Y \in \mathbb{R}^{n\times k}$ (the $k$ desired outputs). The model applies
an unknown matrix $A \in \mathbb{R}^{n\times m}$ and adds an unknown offset
$b \in \mathbb{R}^n$ to every column, so its total error is
\[f(A,b) := \big\lVert \underbrace{AX + b\mathbf{1}\transp - Y}_{=:\,R} \big\rVert^2,
  \qquad \mathbf{1} := (1,\dots,1)\transp \in \mathbb{R}^k,\]
where $\lVert\cdot\rVert$ is the Hilbert--Schmidt norm of \cref{sec:structured_spaces}, so that
$f(A,b) = \sum_{p,q} R_{pq}^2$ is the total squared error over all $k$ data points. Writing out
one entry of the residual,
\[R_{pq} = \sum_{s=1}^{m} A_{ps}X_{sq} + b_p - Y_{pq},\]
so that $\partial R_{pq}/\partial A_{ij} = \delta_{pi}X_{jq}$ and
$\partial R_{pq}/\partial b_i = \delta_{pi}$. The chain rule then gives
\begin{align}
  \frac{\partial f}{\partial A_{ij}} &= 2\sum_{p,q} R_{pq}\,\delta_{pi}X_{jq}
    = 2\sum_{q} R_{iq}X_{jq} = 2\big(RX\transp\big)_{ij}, \nonumber \\
  \frac{\partial f}{\partial b_i} &= 2\sum_{p,q} R_{pq}\,\delta_{pi}
    = 2\sum_{q} R_{iq} = 2\big(R\mathbf{1}\big)_i. \nonumber
\end{align}
Collecting the entries back into arrays of the same shape as the variables,
\[\nabla_{\!A} f = 2\big(AX + b\mathbf{1}\transp - Y\big)X\transp, \qquad
  \nabla_{\!b} f = 2\big(AX + b\mathbf{1}\transp - Y\big)\mathbf{1}.\]
The gradient with respect to $b$ is twice the vector of row sums of the residual, and the
gradient with respect to $A$ is the residual correlated against the inputs. Both are entirely
explicit, and both were computed with nothing beyond partial differentiation and the chain rule.
\end{aiexample}

% Reclassified from ainote to remark: a caution about how the identity is read is mathematics.
\begin{remark}[A notational caution]
One point trips people up here. Since $f$ above has domain
$\mathbb{R}^{n\times m}\times\mathbb{R}^n$ of dimension $nm+n$, its Jacobian is formally a
$1\times(nm+n)$ row, and $\partial_v f(A,b) = \Jac f(A,b)\,v$ would require flattening the matrix
direction $v$ into a column of length $nm$. That flattening is pure bookkeeping and is almost
never performed: one keeps $\nabla_{\!A}f$ shaped like $A$, as above. The identity
$\partial_vf = \Jac f\cdot v$ of \cref{sec:jacobi_matrix} still holds, and is simply being read
in a reshaped coordinate system.
\end{remark}

% Originally: content/week-04.tex, \exercisesheet{4}
% Source: exercises/Ex4_Analysis2_eng.pdf

\begin{exercise}[Derivative computation \textnormal{(important)}]
\label{ex:4.1}
Consider the function $u : (x,y) \mapsto x^{\sin(y)}$, defined for
$(x,y) \in (0,\infty)\times\mathbb{R} \subset \mathbb{R}^2$. Compute $\partial_x u$ and
$\partial_y u$.
\exinfo{This exercise is Problem 4.1 of Problem Sheet 4, Analysis II, Spring Semester 2026.}
\end{exercise}

% Source: old_exams/examHS24.pdf (13 February 2025), Wahr/Falsch-Fragen, Aufgabe 4, p. 3
% Extractor: Gemini 3.1 Pro (High)
\begin{exercise}[True or False \textnormal{(piecewise differentiability)}]
\label{ex:examHS24_TF_diff_piecewise}
Consider the function $f \colon \mathbb{R}^2 \to \mathbb{R}$ given by
\[
f(x,y) = \begin{cases}
y^2 \cos x & \text{if } y \ge 0, \\
0 & \text{if } y < 0.
\end{cases}
\]
Decide whether each of the following statements is true or false.
\begin{enumerate}[label=\textbf{(\alph*)}]
  \item $f \in C^0(\mathbb{R}^2)$.
  \item $\partial_y f \in C^0(\mathbb{R}^2)$.
  \item $\partial_y f \in C^1(\mathbb{R}^2)$.
  \item $f \in C^2(\mathbb{R}^2)$.
  \item $\frac{\partial^2 f}{\partial x^2}(0,0)$ exists.
\end{enumerate}
\exinfo{This exercise is taken from the exam of 13 February 2025 (Prof.\ Joaquim Serra), where it appears as the fourth
block of true/false questions.}
\end{exercise}

% Originally: content/week-04.tex, \section{Solutions}

\section{Solutions}
\label{sec:differential_solutions}

\begin{exercisesolution}[Solution to \cref{ex:4.1}]
% Generator: Claude Sonnet 5 (Medium)
Write $u(x,y) = x^{\sin(y)} = e^{\sin(y)\ln x}$ for $x > 0$. Differentiating with the chain rule:
\[
\partial_x u = \sin(y)\,x^{\sin(y)-1}, \qquad \partial_y u = \cos(y)\ln(x)\,x^{\sin(y)} = \cos(y)\ln(x)\,e^{\sin(y)\ln x}.
\]
\end{exercisesolution}

\begin{exercisesolution}[Proof of \cref{ex:ai_differential_linearization}]
% Generator: Gemini 3.6 Flash (Medium)
\begin{enumerate}[label=\textbf{(\alph*)}]
  \item Partial derivatives are $\frac{\partial f}{\partial x} = 2xy + y\cos(xy)$ and $\frac{\partial f}{\partial y} = x^2 + x\cos(xy)$. Evaluating at $(1,0)$ yields $\frac{\partial f}{\partial x}(1,0) = 0$ and $\frac{\partial f}{\partial y}(1,0) = 1 + 1 = 2$. Thus $\nabla f(1,0) = (0, 2)\transp$, and the total differential $Df(1,0) \colon \mathbb{R}^2 \to \mathbb{R}$ is $Df(1,0)(h_1, h_2) = 2h_2$.
  \item Since $f(1,0) = 0$, the linearization at $(1,0)$ is $L(x,y) = f(1,0) + Df(1,0)(x-1, y-0) = 0 + 0(x-1) + 2y = 2y$.
\end{enumerate}
\end{exercisesolution}

\begin{exercisesolution}[Proof of \cref{ex:examHS24_TF_diff_piecewise}]
% Generator: Gemini 3.1 Pro (High)
\begin{enumerate}[label=\textbf{(\alph*)}]
  \item \textbf{True:} For $y > 0$ and $y < 0$ the function is given by continuous formulas. On the interface $y=0$, we have $\lim_{y \to 0^+} y^2 \cos x = 0 = \lim_{y \to 0^-} 0$, which matches $f(x,0) = 0$.
  \item \textbf{True:} For $y > 0$, $\partial_y f = 2y \cos x$; for $y < 0$, $\partial_y f = 0$. At $y=0$, the one-sided derivatives are $\lim_{h \to 0^+} \frac{h^2 \cos x - 0}{h} = 0$ and $\lim_{h \to 0^-} \frac{0 - 0}{h} = 0$. So $\partial_y f$ exists everywhere and evaluates to $2y \cos x$ for $y \ge 0$ and $0$ for $y < 0$. These two pieces patch continuously at $y=0$.
  \item \textbf{False:} For $\partial_y f$ to be $C^1$, its own partial derivatives must be continuous. However, for $y>0$, $\partial_y(\partial_y f) = 2 \cos x$, while for $y<0$, it is $0$. As $y \to 0$, the upper half approaches $2 \cos x$, which does not match the lower half limit of $0$ unless $\cos x = 0$. Thus $\partial_{yy} f$ has jump discontinuities along the $x$-axis.
  \item \textbf{False:} Since $\partial_{yy}f$ is not continuous (from the previous part), $f$ is not twice continuously differentiable on $\mathbb{R}^2$.
  \item \textbf{True:} Along the $x$-axis ($y=0$), the function is $f(x,0) = 0$ everywhere. Therefore, the function restricted to the $x$-axis is identically zero, meaning $\partial_x f(x,0) = 0$ everywhere, and consequently its derivative with respect to $x$ again, $\partial_{xx} f(0,0)$, is also $0$.
\end{enumerate}
\end{exercisesolution}


% ============================================================
% Chapter 10 --- The chain rule and directional derivatives
% Originally: content/week-04.tex, \section{Chain rule} with its subsections
%             \subsection{Directional derivatives} and
%             \subsection{What does imply differentiability} (Corsin Week 4, Monday),
%             plus \section{The chain rule along a path} (Corsin Week 4, Friday).
%             The two subsections are promoted to sections here.
% ============================================================
\chapter{The Chain Rule and Directional Derivatives}
\label{ch:chain_rule}

% Transition: Claude Opus 5 (High)
Having encoded the differential as a matrix, we now ask how it interacts with composition. The
answer is the chain rule, and in this setting it says something cleaner than its one-variable
version does: the differential of a composition is the \emph{composition of the differentials},
which is to say the product of their matrices. Reading it that way makes it a statement of linear
algebra rather than a formula to be memorised, and explains where the sums of products in the
coordinate form come from.

The same machinery yields directional derivatives, and with them the standard trap of
several-variable calculus. A function can possess every directional derivative at a point and
still fail to be differentiable there, or even continuous.
\Cref{sec:sufficient_condition_differentiability} says what does suffice, and it is the criterion
used in practice for the rest of the course.

% Originally: content/week-04.tex, \section{Chain rule} (Corsin Week 4, Monday)

% Source: Corsin Nick/Class Notes/Week 4.pdf, p. 5
\section{Chain rule}
\label{sec:chain_rule}

\begin{theorem}[Chain rule]
\label{thm:chain_rule}
The \newterm{chain rule} \germanfor{Kettenregel} states: for
$f : U \subseteq \mathbb{R}^n \to \mathbb{R}^m$, $g : V \subseteq \mathbb{R}^m \to \mathbb{R}^d$
differentiable, it holds:
\[D(g\circ f)(x) = Dg(f(x))\,Df(x).\]
The same identity holds for the corresponding Jacobi matrices, where the composition on the
right becomes an ordinary matrix product.
\end{theorem}

% Generator: Claude Opus 5 (High)
\begin{ainote}
Corsin states the chain rule without proof, as does the lecture at this point. The one-line
reason it is true is worth carrying, though: differentiability says $f$ and $g$ are each well
approximated by a linear map near the relevant point, and composing two approximations composes
the linear maps, with the error terms of the two compositions still $o(|h|)$. Making that
precise is the whole proof, and it is the same argument as in one variable.
\end{ainote}

% Supplement: Sascha Brack/Class Notes/Week_04_Notes_Friday_Updated.pdf, pp. 7--10
\begin{remark}[The chain rule as a sum over paths]
\label{rem:chain_rule_paths}
Stated as $D(g\circ f) = Dg\cdot Df$ the rule is compact but, as Sascha Brack puts it in his
notes, \qt{pretty abstract}. It says nothing about \emph{why} a sum appears when the rule is
written out in coordinates. His device is to draw the variables as a graph and read the formula
off it.

Write $x \in \mathbb{R}^n \xmapsto{\ f\ } y \in \mathbb{R}^m \xmapsto{\ g\ } z \in \mathbb{R}^k$
and join each variable to those it feeds into. Then $x_i$ influences $z_j$ only by first moving
some intermediate $y_\ell$, and the contribution of that route is the product of the two rates
along it. Summing over the routes gives the chain rule:

% Sonnet 5 (Medium) --- FIG-W04-03, after the variable graph on Sascha Brack's p. 10.
\begin{center}
\begin{tikzpicture}[>=Latex, scale=1.0]
  \node (xi) at (0,2.4) {$x_i$};
  \node (y1) at (-2.2,1.1) {$y_1$};
  \node (y2) at (-0.7,1.1) {$y_2$};
  \node (yd) at (0.75,1.1) {$\cdots$};
  \node (ym) at (2.2,1.1) {$y_m$};
  \node (zj) at (0,-0.2) {$z_j$};

  \foreach \t in {y1,y2,ym}{\draw[->, blue!70!black] (xi) -- (\t);}
  \foreach \s in {y1,y2,ym}{\draw[->, orange!90!black] (\s) -- (zj);}

  \node[blue!70!black, font=\scriptsize, above left] at (-1.25,1.75)
    {$\dfrac{\partial y_\ell}{\partial x_i}$};
  \node[orange!90!black, font=\scriptsize, below left] at (-1.25,0.55)
    {$\dfrac{\partial z_j}{\partial y_\ell}$};
\end{tikzpicture}
\end{center}

\[\frac{\partial z_j}{\partial x_i}
  \;=\; \sum_{\ell=1}^{m}
  \underbrace{\frac{\partial z_j}{\partial y_\ell}}_{(\Jac g)_{j\ell}}
  \underbrace{\frac{\partial y_\ell}{\partial x_i}}_{(\Jac f)_{\ell i}}
  \;=\; \big(\Jac g(f(x))\,\Jac f(x)\big)_{ji}.\]

The right-hand equality is just the definition of matrix multiplication, and that is precisely the
point: \textbf{matrix multiplication is a sum over paths}, and the chain rule looks like a product of
matrices precisely because composing functions composes those graphs. Sascha's summary is that
computing $\partial z_j/\partial x_i$ this way is nothing more than \qt{computing single
components of the chain rule}.

The reading is worth having because it makes the coordinate form reconstructible instead of
memorised: count the intermediate variables, one term each, each term a product of two factors,
and the index that is summed over is the one belonging to the middle layer. It is also the
form in which the rule is actually used in physics. Compare
\Cref{sec:chain_rule_along_path}, which is this picture with a single intermediate layer
$\gamma_1,\dots,\gamma_n$ and a single source variable $t$.
\end{remark}

% Generator: Claude Opus 5 (High) --- not in any source read so far; added because it is the
% cleanest one-line payoff of the chain rule along a path and is quoted constantly in physics.
\begin{theorem}[Euler's homogeneous function theorem]
\label{thm:eulers_homogeneous_function_theorem}
Let $f \in C^1(\mathbb{R}^n\setminus\{0\},\mathbb{R})$ be \newterm{positively homogeneous of
degree $k$} for some $k \in \mathbb{R}$, meaning that
\[f(tx) = t^kf(x) \qquad \text{for all } x \in \mathbb{R}^n\setminus\{0\} \text{ and all } t>0.\]
Then $\langle x, \nabla f(x)\rangle = k\,f(x)$ for all $x \in \mathbb{R}^n\setminus\{0\}$.
\end{theorem}

\begin{proof}
Fix $x \neq 0$ and read both sides of the homogeneity identity as functions of $t>0$ alone. The
left-hand side is $f$ composed with the path $\gamma(t) := tx$, whose velocity is
$\gamma'(t) = x$, so the chain rule along a path (\cref{sec:chain_rule_along_path}) gives
\[\frac{d}{dt}f(tx) = \big\langle \nabla f(tx),\ x\big\rangle.\]
On the right-hand side $f(x)$ is a constant in $t$, and therefore
$\frac{d}{dt}\big(t^kf(x)\big) = kt^{k-1}f(x)$. The two derivatives agree for every $t>0$;
evaluating at $t=1$ leaves $\langle \nabla f(x), x\rangle = k f(x)$, which is the claim by
symmetry of the inner product.
\end{proof}

\begin{aiexample}[Chain Rule for Vector Fields]
% Generator: Gemini 3.6 Flash (Medium)
Let $f(u,v) := u^2 + v^2$ and let $g(r,\theta) := (r\cos\theta, r\sin\theta)\transp$. The composite function is $(f \circ g)(r,\theta) = r^2\cos^2\theta + r^2\sin^2\theta = r^2$.
By the multi-variable Chain Rule, the Jacobian matrix of $f \circ g$ satisfies:
\[
D(f \circ g)(r,\theta) = Df(g(r,\theta)) \cdot Dg(r,\theta) = \begin{pmatrix} 2r\cos\theta & 2r\sin\theta \end{pmatrix} \begin{pmatrix} \cos\theta & -r\sin\theta \\ \sin\theta & r\cos\theta \end{pmatrix} = \begin{pmatrix} 2r & 0 \end{pmatrix},
\]
which matches the direct calculation $\frac{\partial}{\partial r}(r^2) = 2r$ and $\frac{\partial}{\partial \theta}(r^2) = 0$.
\end{aiexample}

% Source: Corsin Nick/Class Notes/Week 4.pdf, pp. 5--6
\begin{exercise}[Jacobian of a composite map]
\label{ex:chain_rule_computation}
Compute $\Jac(g\circ f)$ for
\[f : \mathbb{R}^3\to\mathbb{R}^2,\ x \mapsto (x_1x_2,\ x_2+x_3), \qquad g : \mathbb{R}^2\to\mathbb{R}^2,\ x \mapsto (e^{x_1},\ x_1x_2).\]
\end{exercise}

% Kept inline: solved directly in Corsin Nick/Class Notes/Week 4.pdf, pp. 5--6.
\begin{exercisesolution}[\cref{ex:chain_rule_computation}]
We have
\[\Jac f(x) = \begin{pmatrix} x_2 & x_1 & 0 \\ 0 & 1 & 1\end{pmatrix}, \qquad \Jac g(x) = \begin{pmatrix} e^{x_1} & 0 \\ x_2 & x_1 \end{pmatrix}.\]
So:
\[
\begin{aligned}
\Jac(g\circ f) &= \begin{pmatrix} e^{x_1x_2} & 0 \\ x_2+x_3 & x_1x_2\end{pmatrix}\begin{pmatrix} x_2 & x_1 & 0 \\ 0 & 1 & 1\end{pmatrix} \\
&= \begin{pmatrix} x_2e^{x_1x_2} & x_1e^{x_1x_2} & 0 \\ x_2^2 + x_2x_3 & 2x_1x_2 + x_1x_3 & x_1x_2 \end{pmatrix}
\end{aligned}
\]
\end{exercisesolution}

% Source: Jérôme Paschoud/Notizen/Woche 4.2 Kettenregel.pdf, p. 1
\begin{exercise}[Computing differentials with and without the chain rule]
\label{ex:chain_rule_computation_paschoud}
Let $f: \mathbb{R}^2 \to \mathbb{R}^3$ and $g: \mathbb{R}^3 \to \mathbb{R}$ be defined by
\[f(x,y) := \begin{pmatrix} x+y \\ x-y \\ x^2+y^2-1 \end{pmatrix}, \qquad g(x,y,z) := x^2+y^2+z^2.\]
Compute the differential $D(g \circ f)$ in two ways:
\begin{enumerate}[label=\textbf{(\alph*)}]
  \item By directly computing the composition $g \circ f$ and then differentiating.
  \item By computing $Df$ and $Dg$ and using the chain rule.
\end{enumerate}
\end{exercise}

\begin{exercisesolution}[\cref{ex:chain_rule_computation_paschoud}]
\begin{enumerate}[label=\textbf{(\alph*)}]
  \item \textbf{Directly:} First compute $g(f(x,y))$:
    \begin{align*}
    (g \circ f)(x,y) &= (x+y)^2 + (x-y)^2 + (x^2+y^2-1)^2 \\
    &= (x^2+2xy+y^2) + (x^2-2xy+y^2) + (x^4+y^4+1+2x^2y^2-2x^2-2y^2) \\
    &= 2x^2 + 2y^2 + x^4 + y^4 + 1 + 2x^2y^2 - 2x^2 - 2y^2 \\
    &= x^4 + y^4 + 2x^2y^2 + 1 \\
    &= (x^2+y^2)^2 + 1.
    \end{align*}
    Now differentiate this scalar function with respect to $x$ and $y$:
    \[D(g \circ f)(x,y) = \begin{pmatrix} 2(x^2+y^2)\cdot 2x & 2(x^2+y^2)\cdot 2y \end{pmatrix} = \begin{pmatrix} 4x(x^2+y^2) & 4y(x^2+y^2) \end{pmatrix}.\]

  \item \textbf{Via chain rule:} Compute the differentials of $f$ and $g$ individually.
    \[Df(x,y) = \begin{pmatrix} 1 & 1 \\ 1 & -1 \\ 2x & 2y \end{pmatrix}, \qquad Dg(u,v,w) = \begin{pmatrix} 2u & 2v & 2w \end{pmatrix}.\]
    Evaluate $Dg$ at $f(x,y)$:
    \[Dg(f(x,y)) = \begin{pmatrix} 2(x+y) & 2(x-y) & 2(x^2+y^2-1) \end{pmatrix}.\]
    Multiply the matrices:
    \begin{align*}
    D(g \circ f)(x,y) &= Dg(f(x,y)) Df(x,y) \\
    &= \begin{pmatrix} 2(x+y) & 2(x-y) & 2(x^2+y^2-1) \end{pmatrix} \begin{pmatrix} 1 & 1 \\ 1 & -1 \\ 2x & 2y \end{pmatrix} \\
    &= \begin{pmatrix} 4x + 4x(x^2+y^2-1) & 4y + 4y(x^2+y^2-1) \end{pmatrix} \\
    &= \begin{pmatrix} 4x(x^2+y^2) & 4y(x^2+y^2) \end{pmatrix}.
    \end{align*}
    Both methods yield the same result.
\end{enumerate}
\end{exercisesolution}

% Source: old_exams/examHS24.pdf (13 February 2025), Wahr/Falsch-Fragen, Aufgabe 3, p. 2
% Extractor: Gemini 3.1 Pro (High)
\begin{exercise}[True or False \textnormal{(the higher-order chain rule)}]
\label{ex:examHS24_TF_chain_rule}
Let $u \in C^2(\mathbb{R}^n)$ and $F = (F_1, \dots, F_n) \in C^2(\mathbb{R}^n, \mathbb{R}^n)$, with
$n \ge 2$. Decide whether each of the following identities is true or false.
\begin{enumerate}[label=\textbf{(\alph*)}]
  \item The identity $\partial_i(u \circ F) = \sum_{j=1}^n \partial_ju \circ F \, \partial_iF_j$ holds for all $i \in \{1, \dots, n\}$.
  \item The identity $\partial_{ii}(u \circ F) = \sum_{j=1}^n \partial_{jj}u \circ F \, (\partial_iF_j)^2 + \sum_{j=1}^n \partial_ju \circ F \, \partial_{ii}F_j$ holds for all $i \in \{1, \dots, n\}$.
  \item The identity $\partial_i\partial_j(u \circ F) = \partial_j\partial_i(u \circ F)$ holds for all $i, j \in \{1, \dots, n\}$.
\end{enumerate}
\exinfo{This exercise is taken from the exam of 13 February 2025 (Prof.\ Joaquim Serra), where it appears as the third
block of true/false questions.}
\exsol{sol:chain_rule_higher_order_tf}
\end{exercise}

% Source: old_exams/examFS24.pdf (21 August 2024), Wahr/Falsch-Fragen, Aufgabe 3, pp. 2--3
% Extractor: Claude Opus 5 (High)
% Presented as a worked example rather than as an exercise: the block is three one-line verdicts,
% and its value is entirely in seeing which hypothesis each one consumes. Its companion above,
% ex:examHS24_TF_chain_rule, is the same theme set as a problem, and assumes F in C^2 where this
% one assumes only C^1 --- which is the whole of part (c).
\begin{example}[Which identities are forced, and by what]
\label{ex:forced_identities_chain_rule}
% Generator: Claude Opus 5 (High) --- the three statements are the examiner's; the discussion is
% written for these notes.
Let $u \in C^2(\mathbb{R}^n)$ and let $F = (F_1,\dots,F_n) \in C^1(\mathbb{R}^n,\mathbb{R}^n)$,
with $n \geq 2$. Exactly one of the following three identities is forced by those hypotheses, and
the two that fail do so for quite different reasons.
\begin{enumerate}[label=\textbf{(\alph*)}]
  \item \textbf{$\partial_iF_j = \partial_jF_i$ for all $i,j$: not forced.} This asks the Jacobi
    matrix $\Jac F$ to be symmetric at every point, which is the integrability condition of
    \cref{lem:integrability_conditions}: it holds precisely when $F$ is a candidate for a gradient,
    and nothing in the hypotheses makes it one. The rotation field
    $F(x) := (-x_2,\, x_1,\, 0,\dots,0)$ is as smooth as one likes and has $\partial_1F_2 = 1$
    while $\partial_2F_1 = -1$.
  \item \textbf{$\partial_j\partial_iu = \partial_i\partial_ju$ for all $i,j$: forced.} This is
    Schwarz's lemma (\cref{thm:schwarz_mixed_partials}), and $u \in C^2$ is exactly its hypothesis.
    Note how differently the two statements read once they are set side by side: \textbf{(a)}
    constrains an arbitrary vector field, \textbf{(b)} constrains the gradient of a function. Only
    the second is automatic.
  \item \textbf{$u \circ F \in C^2$: not forced.} The chain rule gives
    $\partial_i(u\circ F) = \sum_{j}(\partial_ju\circ F)\,\partial_iF_j$, and the factors
    $\partial_iF_j$ are merely continuous, so this says only that $u \circ F \in C^1$. To
    differentiate once more one would have to differentiate $\partial_iF_j$, which the hypotheses
    do not permit. Take $n = 2$, $u(y_1,y_2) := y_1$ and $F(x_1,x_2) := (x_1\lvert x_1\rvert,\, x_2)$.
    Then $F \in C^1$, since $\partial_1F_1 = 2\lvert x_1\rvert$ is continuous, and
    $(u\circ F)(x) = x_1\lvert x_1\rvert$ has the same non-differentiable first derivative at
    $x_1 = 0$. Consequently $u \circ F \notin C^2$.
\end{enumerate}
The general rule that \textbf{(c)} illustrates is that a composition inherits the \emph{weaker} of
the two regularities. Raise $F$ to $C^2$ and $u \circ F$ becomes $C^2$, which is what
\cref{ex:examHS24_TF_chain_rule}\textbf{(c)} is then free to assume.
\exinfo{This example is taken from the exam of 21 August 2024 (Prof.\ Joaquim Serra), where it appears as the third block
of true/false questions.}
\end{example}

% Originally: content/week-04.tex, \subsection{Directional derivatives} --- promoted to a section

% Sonnet 5 (Medium)
\section{Directional derivatives}
\label{sec:directional_derivatives}

% Source: Corsin Nick/Class Notes/Week 4.pdf, p. 6
Sometimes the \newterm{differential} (\germanterm{das Differential}) is not easily expressed as a
Jacobian. Then it can be convenient to use the formula for \newterm{directional derivatives}
(\germanterm{Richtungsableitung}):
\begin{equation*}
DF_{x_0}(v) = \partial_v F(x_0) = \left.\frac{d}{ds}\right|_{s=0} F(x_0 + sv) \tag{$*$}
\end{equation*}

% Generator: Claude Opus 5 (Medium) --- the geometry behind (*): pick a ray, restrict to it,
% and you are back to a one-variable derivative from Analysis I.
Formula $(*)$ describes a two-step recipe, and the picture below is worth carrying around.
\textbf{Step one:} pick a direction $v$ and walk away from $x_0$ along the straight ray
$s \mapsto x_0+sv$, ignoring the rest of the domain. \textbf{Step two:} record the values of $F$
along that ray. What you get is an ordinary function of \emph{one} variable $s$, and
$\partial_vF(x_0)$ is simply its derivative at $s=0$ --- an Analysis I object. The partial
derivative $\partial_iF$ is the special case $v = e_i$: walk parallel to the $i$-th axis.

\begin{center}
\begin{tikzpicture}[>=Latex, scale=1.0]
  % --- Panel 1: the domain, and the rays we are allowed to walk along ---
  \begin{scope}
    \draw[->] (-0.3,0) -- (2.9,0) node[right, font=\small] {$x_1$};
    \draw[->] (0,-0.3) -- (0,2.7) node[above, font=\small] {$x_2$};

    \coordinate (P) at (1.15,1.0);
    % the rays, extended in grey
    \draw[gray!55, dashed] (-0.15,1.0) -- (2.75,1.0);
    \draw[gray!55, dashed] (0.15,0.15) -- (2.35,2.35);

    \draw[orange!90!black, very thick, ->] (P) -- ++(0.8,0) node[below, font=\scriptsize] {$e_1$};
    \draw[green!55!black, very thick, ->] (P) -- ++(0,0.8) node[left, font=\scriptsize] {$e_2$};
    \draw[purple, very thick, ->] (P) -- ++(0.62,0.62) node[above left, font=\scriptsize] {$v$};

    \fill (P) circle (1.8pt) node[below left, font=\scriptsize] {$x_0$};
    \node[below, font=\small, align=center] at (1.3,-0.75)
      {pick a direction and\\ walk along that ray};
  \end{scope}

  % --- Panel 2: the profile along each ray, and its slope at s=0 ---
  \begin{scope}[shift={(6.6,0)}]
    \draw[->] (-1.1,0) -- (2.2,0) node[right, font=\small] {$s$};
    \draw[->] (0,-0.3) -- (0,2.8) node[above right, font=\small] {$F(x_0+sv)$};
    \draw (0,0.05) -- (0,-0.05) node[below right, font=\scriptsize] {$0$};

    % profile along e_1 (gentle positive slope) and along v (negative slope)
    \draw[orange!90!black, very thick, domain=-0.55:1.95, samples=50, smooth]
      plot ({\x}, {1.4 + 0.35*\x + 0.18*\x*\x});
    \draw[purple, very thick, domain=-0.55:1.55, samples=50, smooth]
      plot ({\x}, {1.4 - 0.75*\x + 0.30*\x*\x});

    % their tangents at s = 0
    \draw[orange!90!black, dashed] (-0.5,1.225) -- (1.95,2.0825);
    \draw[purple, dashed] (-0.5,1.775) -- (1.55,0.2375);

    \fill (0,1.4) circle (1.8pt);
    \draw[dotted, gray] (-0.62,1.4) -- (0,1.4);
    \node[font=\scriptsize, left] at (-0.65,1.4) {$F(x_0)$};

    % anchored just off the right end of each tangent line, clear of both curves
    \node[orange!90!black, font=\scriptsize, right] at (2.0,2.20) {slope $=\partial_1F(x_0)$};
    \node[purple, font=\scriptsize, right] at (1.62,0.22) {slope $=\partial_vF(x_0)$};

    \node[below, font=\small, align=center] at (0.5,-0.75)
      {each ray gives a one-variable function;\\ its slope at $0$ is $\partial_vF(x_0)$};
  \end{scope}
\end{tikzpicture}
\end{center}

The right-hand panel is also the cleanest way to see what differentiability adds. Every direction
produces \emph{some} slope; nothing so far forces those slopes to be related to one another. Being
differentiable is precisely the demand that all of them are read off a \emph{single} linear map
via $\partial_vF(x_0) = DF_{x_0}(v)$ --- so that knowing the $n$ partial derivatives already
determines every other direction. \Cref{ex:all_directional_derivatives_exist} below is the
cautionary case where all the slopes exist but no such linear map governs them.

\begin{importantremark}
\label{rem:directional_derivative_caveat}
The directional derivative $\partial_v F(x_0)$ can exist (depending on $v$ and $x_0$) even if $F$
is not differentiable at $x_0$, i.e.\ $DF_{x_0}$ does not exist. However, if $DF_{x_0}$ exists,
the formula $(*)$ holds.
\end{importantremark}

% Generator: Claude Opus 5 (Medium)
\begin{aiexample}[Every direction, and still not continuous]
\label{ex:all_directional_derivatives_exist}
\Cref{ex:partial_derivatives_not_continuous} showed that the two \emph{partial} derivatives are
too little. One might hope that asking for \emph{all} directions is enough. It is not. Let
\[f(x,y) := \begin{cases} \dfrac{x^2y}{x^4+y^2} & (x,y) \neq (0,0), \\[1ex] 0 & (x,y)=(0,0).\end{cases}\]
Fix a direction $v = (a,b) \neq 0$ and compute along the ray $t \mapsto tv$:
\[\frac{f(ta,tb)}{t} = \frac{1}{t}\cdot\frac{t^2a^2\cdot tb}{t^4a^4+t^2b^2}
   = \frac{a^2b}{t^2a^4+b^2} \;\xrightarrow{\ t\to0\ }\;
   \begin{cases} \dfrac{a^2}{b} & b \neq 0, \\[1ex] 0 & b = 0.\end{cases}\]
So $\partial_vf(0,0)$ exists for \textbf{every} direction $v$. Yet $f$ is not even continuous at
the origin: approaching along the \emph{parabola} $y=x^2$ instead of along a straight line,
\[f(x,x^2) = \frac{x^2\cdot x^2}{x^4+x^4} = \frac12 \quad\text{for all } x \neq 0.\]
By \cref{prop:differentiable_implies_continuous}, $f$ is therefore not differentiable at the
origin.

Two lessons worth keeping.
\begin{itemize}
  \item \textbf{Straight lines cannot see everything.} Every directional derivative is computed
    along a line, and this $f$ is well behaved along every line while misbehaving along a
    parabola. Approaching a point in $\mathbb{R}^n$ is genuinely harder than in $\mathbb{R}$,
    where there are only two directions.
  \item \textbf{Watch whether $v \mapsto \partial_vf(x_0)$ is linear.} If $f$ were differentiable
    we would need $\partial_vf(0,0) = DF_0(v)$, a \emph{linear} function of $v$; but $a^2/b$ is
    not linear in $(a,b)$. That failure alone rules out differentiability, without any mention of
    continuity --- and it is precisely the test that settles \cref{ex:5.1}\textbf{(b)} on the
    \cref{ex:5.3}.
\end{itemize}
\end{aiexample}

% Originally: content/week-04.tex, \subsection{What does imply differentiability} --- promoted

% Supplement: Sascha Brack/Class Notes/Week_04_Notes_Monday_Updated.pdf, pp. 3--5
\section{What \emph{does} imply differentiability}
\label{sec:sufficient_condition_differentiability}

\begin{ainote}
Corsin's notes give the definition and the warnings, but never state a \emph{sufficient}
condition --- so far we have two counterexamples and no positive test. Sascha Brack's notes
(Lecture 8) state one as Theorem 2.9 and organise the whole picture into a chain of
implications; both are reproduced below, since without them the practical question
\qt{how do I actually check differentiability?} is left unanswered.
\end{ainote}

\begin{theorem}[Continuous partial derivatives suffice]
\label{thm:continuous_partials_imply_differentiable}
Let $U \subseteq \mathbb{R}^n$ be open and $F : U \to \mathbb{R}^m$. If all partial derivatives
$\partial_iF$ exist on $U$ \textbf{and are continuous} at every $x \in U$, then $F$ is
differentiable at every $x \in U$, and
\[DF_x = \big(\partial_1F(x),\ \dots,\ \partial_nF(x)\big).\]
Such an $F$ is called \newterm{continuously differentiable}, written $F \in C^1(U)$.
\end{theorem}

This is the theorem that makes the subject workable: in practice one almost never checks the
definition directly. Compute the partials, observe that they are built from continuous pieces,
and differentiability follows --- which is exactly why the course can write $f \in C^1$, $C^2$,
$C^k$ so freely from here onward.

Together with the two counterexamples, the whole landscape is a single chain, each link
one-directional:

% Sonnet 5 (Medium) --- FIG-W04-02, after the diagram on Sascha Brack's p. 3.
\begin{center}
\begin{tikzpicture}[>=Latex, node distance=0pt]
  \tikzset{
    box/.style={draw, rounded corners=2pt, align=center, font=\small,
                inner sep=4pt, minimum height=8mm},
    imp/.style={->, very thick, green!45!black},
    nimp/.style={->, thick, red!75!black, dashed}}

  \node[box, fill=green!8] (c1) at (0,4.5) {$F \in C^1$ near $x_0$\\[-2pt]
    {\scriptsize (all $\partial_iF$ exist \emph{and} are continuous)}};
  \node[box, fill=blue!8] (df) at (0,2.9) {$F$ differentiable at $x_0$};
  \node[box] (dir) at (0,1.4) {all directional derivatives $\partial_vF(x_0)$ exist};
  \node[box] (par) at (0,-0.1) {all partial derivatives $\partial_iF(x_0)$ exist};

  \draw[imp] (c1) -- (df);
  \draw[imp] (df) -- (dir);
  \draw[imp] (dir) -- (par);

  \node[font=\scriptsize, green!45!black, right] at (0.15,3.7)
    {\cref{thm:continuous_partials_imply_differentiable}};
  \node[font=\scriptsize, green!45!black, right] at (0.15,2.15) {formula $(*)$};
  \node[font=\scriptsize, green!45!black, right] at (0.15,0.65) {take $v = e_i$};

  % The reverse arrows, all false. Each carries its counterexample as a midway node on the
  % arc itself, filled white so it interrupts the dashed line cleanly. Previously the three
  % labels were free-standing nodes anchored at x = -2.6 while the arcs bulged out past
  % x = -3, so every arrow was drawn straight through its own label.
  \tikzset{noLbl/.style={midway, font=\scriptsize, red!75!black, align=right,
                         fill=white, inner sep=1.5pt}}
  \draw[nimp] (df.west)  to[out=160, in=200, looseness=1.7]
    node[noLbl] {\textbf{no:} \cref{ex:differentiable_not_c1}} (c1.west);
  \draw[nimp] (dir.west) to[out=160, in=200, looseness=1.7]
    node[noLbl] {\textbf{no:} \cref{ex:all_directional_derivatives_exist}} (df.west);
  \draw[nimp] (par.west) to[out=160, in=200, looseness=1.7]
    node[noLbl] {\textbf{no:} \cref{ex:partial_derivatives_not_continuous}} (dir.west);
\end{tikzpicture}
\end{center}

Every downward arrow holds, and \textbf{not one of them reverses} --- each failure is witnessed
by an example already in this chapter, so the diagram doubles as a map of where those
counterexamples belong:
\begin{itemize}
  \item \textbf{Differentiable $\not\Rightarrow$ $C^1$:} \cref{ex:differentiable_not_c1}, whose
    partial derivatives oscillate without limit at the origin even though the differential
    exists there.
  \item \textbf{All directional derivatives $\not\Rightarrow$ differentiable:}
    \cref{ex:all_directional_derivatives_exist}, where every ray gives a slope but
    $v \mapsto \partial_vf(0,0)$ is not linear.
  \item \textbf{Partial derivatives $\not\Rightarrow$ all directional derivatives:}
    \cref{ex:partial_derivatives_not_continuous} again. For $f = \tfrac{xy}{x^2+y^2}$ and a
    direction $v = (a,b)$ with $ab \neq 0$,
    \[\frac{f(ta,tb)}{t} = \frac{1}{t}\cdot\frac{t^2ab}{t^2(a^2+b^2)}
      = \frac{ab}{t\,(a^2+b^2)} \ \xrightarrow{\ t \to 0\ }\ \pm\infty,\]
    so $\partial_vf(0,0)$ exists only along the two axes --- the two directions the partial
    derivatives happen to look at.
\end{itemize}
The three counterexamples are not variations on one theme; they fail at three different
heights, which is precisely why the chain has three separate links.

% Supplement: Sascha Brack/Class Notes/Week_04_Notes_Monday_Updated.pdf, p. 5
\begin{remark}[A recipe for checking differentiability at a point]
\label{rem:differentiability_recipe}
Sascha's notes distil the above into three steps, worth following in this order.
\begin{enumerate}[label=\textbf{\arabic*.}]
  \item \textbf{Compute the partial derivatives $\partial_iF(x_0)$.} If one of them fails to
    exist, $F$ is not differentiable and you are done. If they all exist \emph{and are
    continuous near $x_0$}, $F$ \emph{is} differentiable by
    \Cref{thm:continuous_partials_imply_differentiable} --- also done. Most exercises end here.
  \item \textbf{Otherwise, guess the candidate.} If $F$ is differentiable at all, the differential
    can only be $L := \big(\partial_1F(x_0),\dots,\partial_nF(x_0)\big)$, since the partials are
    forced (\cref{sec:jacobi_matrix}).
  \item \textbf{Verify the candidate against the definition:}
    \[\lim_{h\to0}\frac{\big|F(x_0+h)-F(x_0)-L(h)\big|}{\lVert h\rVert} \overset{?}{=} 0.\]
    In $\mathbb{R}^2$, substituting $h = (r\cos\theta,\ r\sin\theta)$ and checking that the
    result tends to $0$ \emph{uniformly in $\theta$} is usually the quickest route --- the same
    polar trick used in \cref{ex:all_directional_derivatives_exist} and, later, in
    \Cref{q:quiz_q4}.
\end{enumerate}
Step 1 is the one students skip; it settles the overwhelming majority of cases without any
$\varepsilon$--$\delta$ work at all.
\end{remark}

% Supplement: Sascha Brack/Class Notes/Week_03_Notes_Friday.pdf, p. 7
\begin{example}[Checking differentiability at the origin]
\label{ex:differentiability_recipe_application}
Let $f(x,y) := \frac{x^2 y}{x^2+y^2}$ for $(x,y) \neq (0,0)$ and $f(0,0) := 0$. We check if $f$ is differentiable at the origin.
\begin{enumerate}[label=\textbf{\arabic*.}]
  \item \textbf{Partial derivatives:}
    \[\partial_1 f(0,0) = \lim_{t\to 0}\frac{f(t,0)-f(0,0)}{t} = \lim_{t\to 0}\frac{0}{t} = 0, \qquad \partial_2 f(0,0) = \lim_{t\to 0}\frac{f(0,t)-f(0,0)}{t} = 0.\]
    They exist, but to use \cref{thm:continuous_partials_imply_differentiable} we would have to compute them everywhere and show they are continuous at $(0,0)$, which is tedious and (as we will see) false. We proceed to step 2.
  \item \textbf{Guess the candidate:} The only possible differential is $L(h_1,h_2) = \partial_1 f(0,0)h_1 + \partial_2 f(0,0)h_2 = 0$.
  \item \textbf{Verify against the definition:} We check whether the following limit is zero:
    \[\lim_{h\to 0} \frac{|f(h_1,h_2) - f(0,0) - L(h)|}{\lVert h\rVert} = \lim_{h\to 0} \frac{\frac{h_1^2 h_2}{h_1^2+h_2^2}}{\sqrt{h_1^2+h_2^2}} = \lim_{h\to 0} \frac{h_1^2 h_2}{(h_1^2+h_2^2)^{3/2}}.\]
    Passing to polar coordinates $h = (r\cos\theta, r\sin\theta)$, the expression becomes
    \[\frac{r^3\cos^2\theta\sin\theta}{r^3} = \cos^2\theta\sin\theta.\]
    This depends on the angle of approach $\theta$ (for instance, it is $0$ along the axes but $\frac{1}{2\sqrt{2}} \neq 0$ along $h_1=h_2$), so the limit as $r \to 0$ is not uniformly $0$. Hence $f$ is \textbf{not differentiable} at $(0,0)$.
\end{enumerate}
\end{example}

% Supplement: Sascha Brack/Class Notes/Week_04_Notes_Friday_Updated.pdf, pp. 4--5
\begin{example}[A function that is differentiable at the origin]
\label{ex:differentiability_recipe_positive}
Let $f(x,y) := \frac{x^2 y^2}{\sqrt{x^2+y^2}}$ for $(x,y) \neq (0,0)$ and $f(0,0) := 0$. We check differentiability at the origin following the recipe.
\begin{enumerate}[label=\textbf{\arabic*.}]
  \item \textbf{Partial derivatives:} For $(x,y) \neq (0,0)$, the quotient and chain rules give
    \[\partial_1 f(x,y) = xy^2\frac{x^2+2y^2}{(x^2+y^2)^{3/2}}, \qquad \partial_2 f(x,y) = x^2y\frac{2x^2+y^2}{(x^2+y^2)^{3/2}}.\]
    At the origin, computing directly from the definition yields $\partial_1 f(0,0) = 0$ and $\partial_2 f(0,0) = 0$.
    To see if they are continuous at $(0,0)$, we substitute $x = r\cos\theta$ and $y = r\sin\theta$:
    \[\partial_1 f(r\cos\theta, r\sin\theta) = \frac{r\cos\theta\,r^2\sin^2\theta \cdot r^2(1+\sin^2\theta)}{r^3} = r^2\cos\theta\sin^2\theta(1+\sin^2\theta).\]
    As $r \to 0$, this tends to $0$ uniformly in $\theta$, so $\partial_1 f$ is continuous at $(0,0)$. The same holds for $\partial_2 f$. By \cref{thm:continuous_partials_imply_differentiable}, $f$ is differentiable at the origin.
  \item \textbf{Guess the candidate:} We already know $f$ is differentiable and $L = (0,0)$. But for practice, let's verify it with the definition directly.
  \item \textbf{Verify against the definition:} We check the limit
    \[\lim_{h\to 0} \frac{|f(h_1,h_2) - f(0,0) - L(h)|}{\lVert h\rVert} = \lim_{h\to 0} \frac{\frac{h_1^2 h_2^2}{\sqrt{h_1^2+h_2^2}}}{\sqrt{h_1^2+h_2^2}} = \lim_{h\to 0} \frac{h_1^2 h_2^2}{h_1^2+h_2^2}.\]
    In polar coordinates, this is $\frac{r^4\cos^2\theta\sin^2\theta}{r^2} = r^2\cos^2\theta\sin^2\theta$. As $r \to 0$, the limit is clearly $0$ uniformly in $\theta$, confirming $f$ is differentiable with $Df_0 = 0$.
\end{enumerate}
\end{example}

% Originally: content/week-04.tex, \section{The chain rule along a path} (Corsin Week 4, Friday)

% Source: Corsin Nick/Class Notes/Week 4.pdf, p. 8
\section{The chain rule along a path}
\label{sec:chain_rule_along_path}

An important special case of the chain rule is that of a composition of a scalar field with a
path. Suppose $f : \mathbb{R}^n \supseteq U \to \mathbb{R}$ and
$\gamma : [0,1]\to\mathbb{R}^n$ are a $C^1$ \newterm{scalar field} (\germanterm{Skalarfeld}) and
path, respectively. Then
\[
\begin{aligned}
D(f\circ\gamma)(t) &= Df(\gamma(t))\cdot D\gamma(t) \\
&= \langle \nabla f(\gamma(t)),\ \gamma'(t)\rangle \\
&= \sum_{i=1}^{n} \frac{\partial f}{\partial \gamma_i}\frac{\partial \gamma_i}{\partial t}
\end{aligned}
\]
where $\gamma = (\gamma_1,\dots,\gamma_n)$ and
$\dfrac{\partial f}{\partial \gamma_i} := \partial_i f(\gamma(t))$.

This operation is so common in physics that it has its own notation: for $f$ as above, we look at
each coordinate as a function of \textbf{time} and differentiate:
\[\frac{d}{dt}f(x_1,\dots,x_n) := Df(x_1(t),\dots,x_n(t)) = \sum_{i=1}^{n}\frac{\partial f}{\partial x_i}\frac{\partial x_i}{\partial t}.\]

% Generator: Claude Opus 5 (High)
The middle line, $D(f\circ\gamma)(t) = \langle\nabla f(\gamma(t)),\gamma'(t)\rangle$, is worth
dwelling on, because it is the formula that gives the gradient its meaning. It says the rate at
which $f$ changes along a path depends on the velocity only through its inner product with
$\nabla f$ --- so by Cauchy--Schwarz (\cref{thm:cauchy_schwarz}) the change is fastest when
$\gamma'$ points along $\nabla f$, and zero when $\gamma'$ is perpendicular to it. Two standard
facts fall straight out:

\begin{itemize}
  \item \textbf{$\nabla f$ points in the direction of steepest ascent}, with
    $|\nabla f|$ the rate of ascent in that direction.
  \item \textbf{$\nabla f$ is perpendicular to the level sets of $f$.} If $\gamma$ stays inside
    a level set $\{f = c\}$, then $f\circ\gamma$ is constant, so
    $\langle\nabla f(\gamma(t)),\gamma'(t)\rangle = 0$ for every $t$ --- and $\gamma'$ runs
    through the directions tangent to the level set.
\end{itemize}

\begin{center}
% Generator: Claude Opus 5 (High) --- FIG-W04-04: level sets of f(x,y) = x^2 + 2y^2, a path
% crossing them, and the decomposition of gamma' at one point.
% Level sets are the ellipses x^2 + 2y^2 = c, i.e. semi-axes sqrt(c) and sqrt(c/2).
% At P = (1.10, 0.55): grad f = (2x, 4y) = (2.20, 2.20), i.e. the 45-degree direction,
% which is indeed normal to the ellipse there (the ellipse tangent has slope -x/(2y) = -1).
\begin{tikzpicture}[>=Latex, scale=1.5]
  \draw[->] (-2.35,0) -- (2.5,0) node[right, font=\footnotesize] {$x$};
  \draw[->] (0,-1.75) -- (0,1.95) node[above, font=\footnotesize] {$y$};

  % Three level sets of x^2 + 2y^2. The middle one, c = 2.142, is chosen to pass through the
  % marked point P = (1.10, 0.6826): 1.10^2 + 2*0.6826^2 = 1.21 + 0.932 = 2.142.
  \foreach \c in {0.9, 2.142, 4.0}{
    \draw[gray!70, thin] (0,0) ellipse ({sqrt(\c)} and {sqrt(\c/2)});
  }
  % Kept in the upper LEFT: the gradient arrow and its label occupy the upper right.
  \node[gray!75!black, font=\scriptsize, anchor=west] at (-2.25,1.45) {level sets $\{f=c\}$};

  % a path crossing the level sets
  \draw[blue!75!black, thick, domain=0.15:2.05, samples=60, smooth]
    plot (\x, {0.28 + 0.30*\x + 0.06*\x*\x});
  \node[blue!75!black, font=\scriptsize, anchor=south east] at (0.55,0.42) {$\gamma$};

  % the point P on the middle ellipse (c = 1.815) and on the path:
  % path at x = 1.10 gives y = 0.28 + 0.33 + 0.0726 = 0.6826 -- adjust c to match the point.
  \coordinate (P) at (1.10,0.6826);
  \fill[black] (P) circle (0.8pt);
  \node[font=\scriptsize, below right] at (1.13,0.63) {$\gamma(t)$};

  % gradient of f at P: (2x, 4y) = (2.20, 2.73), normalised and scaled to length 0.62
  \draw[red!80!black, very thick, ->] (P) -- ++(0.39,0.48)
    node[anchor=south west, font=\scriptsize] {$\nabla f$};
  % velocity: gamma'(x) = 0.30 + 0.12x = 0.432, direction (1, 0.432), scaled
  \draw[blue!75!black, very thick, ->] (P) -- ++(0.57,0.246)
    node[anchor=north west, font=\scriptsize] {$\gamma'$};

  \node[font=\scriptsize, align=center, anchor=north] at (0,-1.85)
    {$\dfrac{d}{dt}f(\gamma(t)) = \langle \nabla f, \gamma'\rangle$: only the component of
     $\gamma'$ along $\nabla f$ contributes,\\
     so $f$ is constant along a level set and grows fastest along $\nabla f$};
\end{tikzpicture}
\end{center}

% Sonnet 5 (Medium)


% --- exercises relocated here from the dissolved sheet 4 ---

% Originally: content/week-04.tex, end of the chapter

\begin{aiexercise}[Directional Derivative along a Unit Vector]
\label{ex:ai_directional_derivative}
% Generator: Gemini 3.6 Flash (Medium)
Let $f(x,y) := x^2 y$ defined on $\mathbb{R}^2$. Compute the directional derivative $D_v f(1,2)$ of $f$ at the point $(1,2)$ along the unit vector $v := \frac{1}{\sqrt{2}}(1,1)\transp$.
\end{aiexercise}

% Originally: content/week-04.tex, \exercisesheet{4}
% Source: exercises/Ex4_Analysis2_eng.pdf

\begin{exercise}[4.6 --- A directional derivative vanishes \textnormal{(important)}]
\label{ex:4.6}
Let $u \in C^1(\mathbb{R}^n)$ and $\nu \in \mathbb{R}^n$. Show that
\begin{enumerate}[label=\textbf{\arabic*.}]
  \item If $\partial_1 u \equiv 0$ then ``$u$ does not depend on $x_1$''; more rigorously: there
    exists a unique function $v \in C^1(\mathbb{R}^{n-1})$ such that
    \begin{equation*}
    u(x_1,\dots,x_n) = v(x_2,\dots,x_n) \quad \text{for all } x \in \mathbb{R}^n. \tag{2}
    \end{equation*}
  \item If $\partial_\nu u \equiv 0$ and $\nu\cdot e_1 \neq 0$ then ``$u$ is a function of $n-1$
    variables''; more rigorously: there exists a unique function $w \in C^1(\mathbb{R}^{n-1})$
    such that
    \[u(x_1,\dots,x_n) = w\!\left(x_2 - \frac{x_1\nu_2}{\nu_1}, \dots, x_n - \frac{x_1\nu_n}{\nu_1}\right) \quad \text{for all } x \in \mathbb{R}^n.\]
  \item \textbf{(*)} What can we conclude if we assume only that $\partial_1 u = 0$ in an open
    connected subset $U \subset \mathbb{R}^n$?
\end{enumerate}
\textit{Official hints:} \textbf{4.6.2} --- apply the mean value theorem to $u(x + t\nu)$.
\textbf{4.6.3} --- consider the domain $U := \{(x,y) : x^2 < y < x^2+1\}$ and the function
\[u(x,y) := \begin{cases} \max\{0, y-2\}^2 & \text{if } x \geq 0, \\ 0 & \text{if } x \leq 0.\end{cases}\]
\end{exercise}

% Originally: content/week-04.tex, \section{Solutions}

\section{Solutions}
\label{sec:chain_rule_solutions}

\begin{exercisesolution}[Solution to \cref{ex:4.6}]
% Generator: Claude Sonnet 5 (Medium)
\begin{enumerate}[label=\textbf{\arabic*.}]
  \item Define $v(x_2,\dots,x_n) := u(0,x_2,\dots,x_n)$, which is $C^1$ as the restriction of $u$
    to the hyperplane $x_1 = 0$. For fixed $(x_2,\dots,x_n)$, the function
    $t \mapsto u(t,x_2,\dots,x_n)$ has derivative $\partial_1 u \equiv 0$ everywhere, hence is
    constant in $t$; evaluating at $t=0$ gives $u(x_1,\dots,x_n) = u(0,x_2,\dots,x_n) =
    v(x_2,\dots,x_n)$ for all $x_1$. Uniqueness: if some $\tilde v$ also satisfies
    $u(x) = \tilde v(x_2,\dots,x_n)$, then setting $x_1 = 0$ gives $\tilde v = v$.
  \item Since $\nu\cdot e_1 = \nu_1 \neq 0$, consider the invertible linear change of coordinates
    $y_i := x_i - \dfrac{x_1\nu_i}{\nu_1}$ for $i = 2,\dots,n$, together with $x_1$ itself.
    Moving $x_1$ by $t$ while holding $y_2,\dots,y_n$ fixed forces $x_i$ to move by
    $t\nu_i/\nu_1$, i.e.\ moves $x$ in the direction $\tfrac{1}{\nu_1}\nu$. By the hint, applying
    the mean value theorem to $u(x+t\nu)$ shows $t \mapsto u(x+t\nu)$ is constant since
    $\partial_\nu u \equiv 0$; consequently $u$, expressed in the coordinates $(x_1, y_2,\dots,
    y_n)$, has vanishing $x_1$-derivative and is therefore (by part 1) independent of $x_1$ in
    these coordinates. Hence, there is a (unique, by the same argument as in part 1) function
    $w \in C^1(\mathbb{R}^{n-1})$ with
    \[
    u(x_1,\dots,x_n) = w\!\left(x_2 - \frac{x_1\nu_2}{\nu_1},\dots,x_n - \frac{x_1\nu_n}{\nu_1}\right).
    \]
  \item In general, only \textbf{locally}: on any open subset where every fiber
    $\{x_1 : (x_1,x_2,\dots,x_n) \in U\}$ (for fixed $(x_2,\dots,x_n)$) is connected, the same
    argument as in part 1 shows $u$ is locally of the form $v(x_2,\dots,x_n)$. But this need not
    hold \textbf{globally} on all of $U$, even though $U$ itself is connected: with
    $U := \{(x,y) : x^2 < y < x^2+1\}$ and $u(x,y) := \max\{0,y-2\}^2$ for $x \geq 0$, $u := 0$
    for $x \leq 0$, one checks $\partial_1 u \equiv 0$ throughout $U$ (near $x=0$ both branches
    agree, since $y < 2$ there), yet $u$ is not globally independent of $x$: for $y > 2$, the
    fiber of $U$ over $y$ splits into two disconnected intervals (one with $x>0$, one with
    $x<0$), and $u$ takes different values on each. So, connectedness of $U$ alone does not
    suffice; what is needed is connectedness of each vertical fiber of $U$.
\end{enumerate}
\end{exercisesolution}

\begin{exercisesolution}[Proof of \cref{ex:ai_directional_derivative}]
% Generator: Gemini 3.6 Flash (Medium)
The gradient of $f(x,y)$ is $\nabla f(x,y) = (2xy, x^2)\transp$. Evaluating at $(1,2)$ gives $\nabla f(1,2) = (4,1)\transp$. Since $f$ is continuously differentiable ($C^\infty$), the directional derivative along the unit vector $v$ is given by the inner product:
\[
D_v f(1,2) = \nabla f(1,2) \cdot v = \begin{pmatrix} 4 \\ 1 \end{pmatrix} \cdot \frac{1}{\sqrt{2}}\begin{pmatrix} 1 \\ 1 \end{pmatrix} = \frac{4 + 1}{\sqrt{2}} = \frac{5}{\sqrt{2}}.
\]
\end{exercisesolution}


% ============================================================
% Chapter 11 --- The mean value theorem and Taylor expansions
% Originally: content/week-04.tex, \subsection{The mean value theorem} (promoted to a
%             section) and \section{Taylor expansions} (Corsin Week 4)
% ============================================================
\chapter{The Mean Value Theorem and Taylor Expansions}
\label{ch:taylor}

% Transition: Claude Opus 5 (Medium)
Having linearized a map once, we ask how good that linear approximation is, and how to improve it
order by order --- the several-variable analogue of the Taylor series from \emph{Analysis I}. The
mean value theorem comes first, because it is what the error terms are estimated with, and
because it is the first result in this course that becomes an \emph{inequality} rather than an
equality once the target is $\mathbb{R}^m$ with $m > 1$ (\cref{ex:4.5}).

% Originally: content/week-04.tex, \subsection{The mean value theorem} --- promoted to a section

% Quelle: Diego Torres Tejeda/Notes - 09.03 - Diego Torres Tejeda.pdf, p. 8
\section{The mean value theorem}
\label{sec:mean_value_theorem}

\begin{ainote}
The mean value theorem is used twice in this chapter --- in the official hint to
\Cref{ex:4.6} and in its solution --- but Corsin never states it. Diego Torres Tejeda's notes
state it, together with the warning that follows, which is what \cref{ex:4.6} actually turns on.
\end{ainote}

\begin{theorem}[Mean value theorem]
\label{thm:mean_value_theorem}
Let $f : U \to \mathbb{R}$ be differentiable, $x \in U$ and $h \in \mathbb{R}^n$ with
$x+th \in U$ for all $t \in [0,1]$. Then there is $t_0 \in (0,1)$ such that, writing
$\xi := x+t_0h$,
\[f(x+h)-f(x) = \partial_hf(\xi).\]
\end{theorem}

\begin{proof}
Apply the one-variable mean value theorem to $\varphi(t) := f(x+th)$ on $[0,1]$, which is
differentiable with $\varphi'(t) = \partial_hf(x+th)$ by the chain rule along a path
(\cref{sec:chain_rule_along_path}). It gives $t_0 \in (0,1)$ with
$\varphi(1)-\varphi(0) = \varphi'(t_0)$, which is the claim.
\end{proof}

\begin{importantremark}[It fails for vector-valued $f$]
\label{rem:mvt_fails_vector_valued}
The hypothesis \qt{$f : U \to \mathbb{R}$} is essential --- \textbf{the theorem is false the
moment the target is $\mathbb{R}^m$ with $m \geq 2$}, and this is exactly what
\cref{ex:4.5} below probes.
% Corrected: this pointed at \cref{ex:4.6}, which applies the *scalar* theorem to $u(x+t\nu)$
% and never leaves $\mathbb{R}$. The problem that probes the vector-valued failure is 4.5 ---
% which was not in the built document when the sentence was written.

Applying the scalar result to each component $f_1,\dots,f_m$ separately does produce points
$\xi_1,\dots,\xi_m$ with
\[f(x+h)-f(x) = \begin{pmatrix}\partial_hf_1(\xi_1) \\ \vdots \\ \partial_hf_m(\xi_m)\end{pmatrix},\]
but the $\xi_i$ are \emph{different points} in general, and no single $\xi$ need work for all
components at once. The standard witness is the \emph{circle}
$\gamma(t) := (\cos t, \sin t)$ on $[0,2\pi]$: here $\gamma(2\pi)-\gamma(0) = 0$, while
$\gamma'(t) = (-\sin t,\cos t)$ has $|\gamma'(t)| = 1$ and so never vanishes --- there is no
$\xi$ at all. Geometrically this is entirely reasonable: a walker who returns to the starting
point must have had zero \emph{average} velocity, but need never have stood still. In one
dimension \qt{returning} forces a turning point; in two it does not.

What survives for vector-valued $f$ is the \emph{mean value inequality},
$|f(x+h)-f(x)| \leq \sup_{t\in[0,1]}\lVert Df_{x+th}\rVert\,|h|$, which is what one actually
uses in practice.
\end{importantremark}

\begin{center}
% Generator: Claude Opus 5 (High) --- FIG-W04-06: the scalar theorem and its vector-valued
% failure, side by side.
% Left: f(x) = 0.35 + 1.5x - 0.55x^2 on [0.15, 2.4]. Chord slope between a=0.15 and b=2.4 is
% f'(xi) with f'(x) = 1.5 - 1.1x, and (f(b)-f(a))/(b-a) = 1.5 - 1.1*(a+b)/2 = 1.5 - 1.4025
% = 0.0975, so xi = (1.5 - 0.0975)/1.1 = 1.275 -- the midpoint (a+b)/2, as it must be for a
% parabola. Marked there.
% Right: the circle t -> (cos t, sin t); start = end, yet the velocity has length 1 throughout.
\begin{tikzpicture}[>=Latex, scale=1.0]

  % ---------- Left: the scalar mean value theorem ----------
  \begin{scope}
    \draw[->] (-0.2,0) -- (3.0,0) node[right, font=\footnotesize] {$x$};
    \draw[->] (0,-0.2) -- (0,2.5) node[above, font=\footnotesize] {$y$};

    \draw[blue!75!black, thick, domain=0.15:2.55, samples=70, smooth]
      plot (\x, {0.35 + 1.5*\x - 0.55*\x*\x});

    % Chord from a = 0.35 (y = 0.8076) to b = 2.4 (y = 0.7820); slope -0.0125.
    % For a parabola the mean value point is the midpoint, xi = (0.35+2.4)/2 = 1.375,
    % where f = 1.3727 and f' = 1.5 - 1.1*1.375 = -0.0125 -- equal to the chord slope.
    \draw[red!80!black, thick] (0.35,0.8076) -- (2.4,0.7820);
    \fill (0.35,0.8076) circle (1.2pt) node[below, font=\scriptsize] {$x$};
    \fill (2.4,0.7820) circle (1.2pt) node[below right, font=\scriptsize] {$x+h$};

    \draw[green!50!black, thick] (0.70,1.3811) -- (2.05,1.3643);
    \fill[green!50!black] (1.375,1.3727) circle (1.4pt);
    \node[green!50!black, font=\scriptsize, above] at (1.375,1.44) {$\xi$};

    \node[font=\scriptsize, align=center, anchor=north] at (1.4,-0.42)
      {\textbf{scalar $f$:} some $\xi$ has\\ tangent parallel to the chord};
  \end{scope}

  % ---------- Right: the vector-valued failure ----------
  \begin{scope}[shift={(5.6,1.15)}]
    \draw[->] (-1.5,0) -- (1.55,0) node[above right, font=\footnotesize] {$x$};
    \draw[->] (0,-1.5) -- (0,1.6) node[above, font=\footnotesize] {$y$};

    \draw[blue!75!black, thick] (0,0) circle (1);
    \fill[red!80!black] (1,0) circle (1.6pt);
    % Label placed below the axis: the velocity arrow at t = 0 points straight up from (1,0),
    % and the x-axis label sits just above the axis, so both sides are occupied.
    \draw[red!80!black, thin] (1,0) -- (1.18,-0.55);
    \node[red!80!black, font=\scriptsize, anchor=north west] at (1.10,-0.52)
      {$\gamma(0)=\gamma(2\pi)$};

    % velocity vectors, always of length 1, never zero
    \foreach \a in {0,72,144,216,288}{
      \draw[green!50!black, very thick, ->]
        ({cos(\a)},{sin(\a)}) -- ++({-0.52*sin(\a)},{0.52*cos(\a)});
    }
    \node[green!50!black, font=\scriptsize, anchor=east] at (-1.1,0.95)
      {$|\gamma'|=1$};

    \node[font=\scriptsize, align=center, anchor=north] at (0,-1.57)
      {\textbf{vector-valued $\gamma$:} $\gamma(2\pi)-\gamma(0)=0$,\\ yet $\gamma'$ never
       vanishes --- no $\xi$ exists};
  \end{scope}
\end{tikzpicture}
\end{center}

% Source: Corsin Nick/Class Notes/Week 4.pdf, p. 7
\begin{exercise}[Differential of $A \mapsto A\transp A$]
\label{ex:transpose_diff}
We identify the inner product space of real-valued matrices $\mathbb{R}^{n\times n}$ with inner
product
\[\langle A, B\rangle = \sum_{i,j=1}^{n} A_{ij}B_{ij} = \Tr(A\transp B)\]
for $A, B \in \mathbb{R}^{n\times n}$, with the euclidean space
$\mathbb{R}^{n^2}\cong\mathbb{R}^{n\times n}$, with the standard inner product. Then, consider
the function
\[F : \mathbb{R}^{n\times n}\to\mathbb{R}^{n\times n}, \qquad A \mapsto A\transp A.\]
Compute the differential at the identity matrix $\id \in \mathbb{R}^{n\times n}$ applied to
$X \in \mathbb{R}^{n\times n}$:
\[DF_{\id}(X) = \partial_X F(\id).\]
\end{exercise}

% Kept inline: solved directly in Corsin Nick/Class Notes/Week 4.pdf, p. 7.
\begin{exercisesolution}[\cref{ex:transpose_diff}]
We use $(*)$:
\[
\begin{aligned}
DF_{\id}(X) &= \left.\frac{d}{ds}\right|_{s=0} F(\id+sX) \\
&= \left.\frac{d}{ds}\,(\id+sX)\transp(\id+sX)\right|_{s=0} \\
&= \left.\frac{d}{ds}\left[\id + s(X\transp+X) + s^2X\transp X\right]\right|_{s=0} \\
&= X\transp + X
\end{aligned}
\]
\end{exercisesolution}


% --- exercise relocated here from the dissolved sheet 4 ---

% Originally: content/exercise-sheets/week-04.tex, problem 4.5 --- never previously built
% Source: exercises/Ex4_Analysis2_eng.pdf

\begin{exercise}[4.5 --- Mean value for vector-valued functions \textnormal{(optional)}]
\label{ex:4.5}
Let $f \in C^1(\mathbb{R},\mathbb{R}^m)$ for $m > 1$. Is it true that there is $t \in [0,1]$ such
that
$$f(1) - f(0) = Df_t(1) = \begin{pmatrix} f_1'(t) \\ \vdots \\ f_m'(t) \end{pmatrix}?$$
Prove it or provide a counterexample.

\textit{Official hint:} try $f(t) = (\sin(2\pi t), \cos(2\pi t))$.
\end{exercise}

% Originally: content/week-04.tex, \section{Taylor expansions} (Corsin Week 4, Friday)

% Source: Corsin Nick/Class Notes/Week 4.pdf, p. 9
\section{Taylor expansions}
\label{sec:taylor_expansions}

% Sonnet 5 (Medium)
\subsection{The general formula}

In the lecture, you have seen the following formula. For
$f \in C^k(U \subseteq \mathbb{R}^n, \mathbb{R})$, $\bar x \in U$:
\[f(\bar x + h) = \sum_{|\alpha| = 0}^{k} \partial^\alpha f(\bar x)\,\frac{h^\alpha}{\alpha!} + o(|h|^k)\]
where $\alpha$ is a \newterm{multi-index} \germanfor{Multiindex} in $\mathbb{N}^n$.

In practice, we never use this formula directly, but it is worth unpacking once (e.g., as in \cref{ex:taylor_maximal_expansion}).

% Generator: Claude Opus 5 (High) --- the multi-index notation above is only well defined
% because mixed partials commute; the fact is used twice below (in the expansion and in
% \cref{rem:second_order_taylor_hessian}) but stated nowhere in the source notes.
% Neither Corsin's notes nor the working formula below mention Schwarz's theorem, yet both depend
% on it. Reclassified from ainote to plain prose: this is a mathematical prerequisite.
% Transition: Claude Opus 5 (High)
Before unpacking the formula, one hypothesis hidden in the notation deserves to be made explicit.
Writing $\partial^\alpha$ for a multi-index $\alpha$ presumes that it does not matter \emph{in
which order} the derivatives are taken, since otherwise $\partial^{(1,1)}$ would be ambiguous
between $\partial_1\partial_2$ and $\partial_2\partial_1$. That the presumption is legitimate is
itself a theorem, and it is stated next.

\begin{theorem}[Schwarz / Clairaut --- mixed partial derivatives commute]
\label{thm:schwarz_mixed_partials}
Let $U \subseteq \mathbb{R}^n$ be open and $f \in C^2(U,\mathbb{R})$. Then for all
$i,j \in \{1,\dots,n\}$ and all $x \in U$,
\[\partial_i\partial_j f(x) = \partial_j\partial_i f(x).\]
More generally, for $f \in C^k(U,\mathbb{R})$ any $k$ partial derivatives may be permuted
freely, so $\partial^\alpha f$ depends only on the multi-index $\alpha$ and not on the order of
differentiation. In particular the Hessian $\Hess f(x)$ of
\cref{rem:second_order_taylor_hessian} is a \textbf{symmetric} matrix.
\end{theorem}

\begin{remark}[The hypothesis is not decoration]
\label{rem:schwarz_needs_c2}
Continuity of the second derivatives is genuinely needed. The standard counterexample is
\[f(x,y) := \begin{cases} xy\,\dfrac{x^2-y^2}{x^2+y^2} & (x,y)\neq(0,0), \\[1ex]
  0 & (x,y)=(0,0),\end{cases}\]
for which a direct computation from the definition gives
$\partial_1\partial_2 f(0,0) = 1$ but $\partial_2\partial_1 f(0,0) = -1$. This $f$ has both
mixed partials everywhere, but they are not continuous at the origin, so
\cref{thm:schwarz_mixed_partials} does not apply, and indeed fails. Since everything in
\cref{sec:taylor_expansions} assumes $f \in C^k$, the issue never arises here. It is worth
knowing that it \emph{could}.
\end{remark}

% Supplement: Sascha Brack/Class Notes/Week_05_Notes_Monday.pdf, p. 7
\begin{exercise}[Applying Schwarz's theorem]
\label{ex:schwarz_application}
Let $f : \mathbb{R}^2 \to \mathbb{R}$ be defined by
\[f(x,y) := x^2 e^{y^2} + e^{-x^2} \sin\left(e^{x^2} \left(x^5 + 3x^2 + 100x + \frac{1}{2+\sin(x)}\right)\right).\]
\begin{enumerate}[label=\textbf{(\alph*)}]
  \item Show that $f \in C^2(\mathbb{R}^2, \mathbb{R})$.
  \item Compute $\partial_2 \partial_1 f(x,y)$. (Make it as easy as possible for yourself.)
\end{enumerate}
\end{exercise}

\begin{exercisesolution}[\cref{ex:schwarz_application}]
\begin{enumerate}[label=\textbf{(\alph*)}]
  \item Every term in $f$ is a composition of $C^2$ functions (polynomials, exponentials, sine),
    and the denominator $2+\sin(x)$ is bounded away from zero. Thus $f$ is $C^2$.
  \item Since $f \in C^2$, by Schwarz's theorem (\cref{thm:schwarz_mixed_partials}) we have
    $\partial_2 \partial_1 f = \partial_1 \partial_2 f$. Computing $\partial_1$ of the second
    term would be tedious. However, notice that the second term depends only on $x$, so its
    derivative with respect to $y$ is $0$. Thus, we first compute $\partial_2 f$:
    \[\partial_2 f(x,y) = x^2 (2y e^{y^2}) + 0 = 2x^2 y e^{y^2}.\]
    Then we compute $\partial_1$ of this result:
    \[\partial_1 \partial_2 f(x,y) = \partial_1 (2x^2 y e^{y^2}) = 4xy e^{y^2}.\]
    By Schwarz's theorem, $\partial_2 \partial_1 f(x,y) = 4xy e^{y^2}$.
\end{enumerate}
\end{exercisesolution}

% Source: Corsin Nick/Class Notes/Week 4.pdf, pp. 9--10
\begin{exercise}[Maximal Taylor expansion in $\mathbb{R}^2$]
\label{ex:taylor_maximal_expansion}
Let $f \in C^3(\mathbb{R}^2,\mathbb{R})$. Write out the maximal Taylor approximation.
\end{exercise}

% Kept inline: solved directly in Corsin Nick/Class Notes/Week 4.pdf, pp. 9--10.
\begin{exercisesolution}[\cref{ex:taylor_maximal_expansion}]
The multi-indices are:
\begin{itemize}
  \item $|\alpha| = 0$: $\alpha = (0,0)$
  \item $|\beta| = 1$: $\beta_1 = (1,0)$, $\beta_2 = (0,1)$
  \item $|\gamma| = 2$: $\gamma_1 = (1,1)$, $\gamma_2 = (2,0)$, $\gamma_3 = (0,2)$
  \item $|\delta| = 3$: $\delta_1 = (3,0)$, $\delta_2 = (2,1)$, $\delta_3 = (1,2)$, $\delta_4 = (0,3)$
\end{itemize}
Recall that
\[\partial^\alpha = \partial_1^{\alpha_1}\partial_2^{\alpha_2}\cdots\partial_n^{\alpha_n}, \qquad \text{e.g.\ } \partial^{\delta_2} = \partial^{(2,1)} = \partial_1^2\partial_2,\]
and
\[\alpha! = \alpha_1!\cdots\alpha_n!, \qquad \text{e.g.\ } \gamma_2! = (2,0)! = 2\cdot 1 = 2.\]
So:
\[
\begin{aligned}
f(\bar x + h) = &\ f(\bar x) && (\alpha) \\
&+ \partial_1 f(\bar x)h_1 + \partial_2 f(\bar x)h_2 && (\beta_1, \beta_2)\\
&+ \partial_1\partial_2 f(\bar x)h_1h_2 + \partial_1^2 f(\bar x)\frac{h_1^2}{2} + \partial_2^2 f(\bar x)\frac{h_2^2}{2} && (\gamma_1, \gamma_2, \gamma_3)\\
&+ \partial_1^3 f(\bar x)\frac{h_1^3}{6} + \partial_1^2\partial_2 f(\bar x)\frac{h_1^2h_2}{2} && (\delta_1, \delta_2)\\
&+ \partial_1\partial_2^2 f(\bar x)\frac{h_1h_2^2}{2} + \partial_2^3 f(\bar x)\frac{h_2^3}{6} && (\delta_3, \delta_4)\\
&+ o(|h|^3)
\end{aligned}
\]
\end{exercisesolution}

\begin{ainote}
Corsin writes the recall line as
$\partial^\alpha = \partial_{\alpha_1}\partial_{\alpha_2}\cdots\partial_{\alpha_n}$ with the
example $\partial^{(2,1)} = \partial_2\partial_1$. Read literally that is the wrong object, since
a multi-index $\alpha$ records \textit{how many times} each $\partial_i$ is applied, so
$\partial^{(2,1)} = \partial_1^2\partial_2$. His own expansion above uses
$\partial_1^2\partial_2 f(\bar x)\,\tfrac{h_1^2h_2}{2}$ for $\delta_2$, confirming the intended
meaning. Corrected above.
\end{ainote}

% Sonnet 5 (Medium)
\subsection{The working formula}

% Source: Corsin Nick/Class Notes/Week 4.pdf, p. 11
We can rewrite this as
\[f(\bar x + h) = \sum_{\ell=0}^{k}\frac{1}{\ell!}\big(h_1\partial_1 + h_2\partial_2\big)^\ell f(\bar x).\]

\begin{theorem}[Taylor working formula]
\label{thm:taylor_working_formula}
For $f \in C^k(U \subseteq \mathbb{R}^n, \mathbb{R})$ with $\bar x + th \in U$ for all
$t \in [0,1]$:
\[f(\bar x + h) = \sum_{\ell=0}^{k}\frac{1}{\ell!}\left(\sum_{m=1}^{n} h_m\partial_m\right)^{\!\ell} f(\bar x) + o(|h|^k)\]
\end{theorem}

\begin{ainote}
Both sums in \cref{thm:taylor_working_formula} are written with lower limit $\ell = 1$ (resp.\
starting index) in the source. The $\ell = 0$ term is $f(\bar x)$ itself, which the explicit
expansion above does include, so the lower limit must be $0$. Corrected above; compare the
$|\alpha| = 0$ line of \cref{ex:taylor_maximal_expansion}, which is the same term written in
multi-index form.
\end{ainote}

% Generator: Claude Sonnet 5 (High) --- new content: \cref{thm:taylor_working_formula} only
% gives the asymptotic (Peano) remainder $o(|h|^k)$. The exact (Lagrange) remainder is genuinely
% different information --- an equality valid for that specific $h$, not just eventually small
% --- and was not stated anywhere in this document.
\begin{theorem}[Taylor's theorem with Lagrange remainder]
\label{thm:taylor_lagrange_remainder}
Let $U \subseteq \mathbb{R}^n$ be open, $f \in C^{k+1}(U,\mathbb{R})$, $\bar x \in U$, and
$h \in \mathbb{R}^n$ with $\bar x+th \in U$ for all $t \in [0,1]$. Then there exists
$\theta \in (0,1)$ such that
\[f(\bar x+h) = \sum_{\ell=0}^{k}\frac{1}{\ell!}\Big(\sum_{m=1}^n h_m\partial_m\Big)^{\!\ell}
  f(\bar x)
  \ +\ \frac{1}{(k+1)!}\Big(\sum_{m=1}^n h_m\partial_m\Big)^{\!k+1}f(\bar x+\theta h).\]
Equivalently, in multi-index notation,
\[f(\bar x+h) = \sum_{|\alpha|\leq k}\partial^\alpha f(\bar x)\,\frac{h^\alpha}{\alpha!}
  \ +\ \sum_{|\alpha|=k+1}\partial^\alpha f(\bar x+\theta h)\,\frac{h^\alpha}{\alpha!}.\]
\end{theorem}

\begin{proof}
Let $\varphi(t) := f(\bar x+th)$ for $t \in [0,1]$, well defined and $C^{k+1}$ since $f \in
C^{k+1}$ and the segment $\bar x+th$ stays in $U$. By induction on $\ell$ via the chain rule,
which is exactly the computation underlying \cref{thm:taylor_working_formula} itself,
\[\varphi^{(\ell)}(t) = \Big(\sum_{m=1}^n h_m\partial_m\Big)^{\!\ell}f(\bar x+th)
  \qquad \text{for } 0 \leq \ell \leq k+1.\]
Apply the one-variable Taylor theorem with Lagrange remainder (\emph{Analysis I}) to $\varphi$
on $[0,1]$: there is $\theta \in (0,1)$ with
\[\varphi(1) = \sum_{\ell=0}^k\frac{\varphi^{(\ell)}(0)}{\ell!} + \frac{\varphi^{(k+1)}(\theta)}{(k+1)!}.\]
Substituting $\varphi(1) = f(\bar x+h)$, $\varphi^{(\ell)}(0) = \big(\sum_mh_m\partial_m\big)^\ell
f(\bar x)$, and $\varphi^{(k+1)}(\theta) = \big(\sum_mh_m\partial_m\big)^{k+1}f(\bar x+\theta h)$
gives the working-formula version. The multi-index version follows by expanding
$\big(\sum_mh_m\partial_m\big)^\ell$ via the multinomial theorem, exactly as in
\cref{ex:taylor_maximal_expansion}.
\end{proof}

\begin{remark}
Unlike \cref{thm:taylor_working_formula}'s $o(|h|^k)$ remainder, this is an \emph{exact}
equality for the specific $h$ in question, not merely a statement about what happens as
$h \to 0$. That is genuinely different information, and it is what is needed whenever a hard,
uniform error bound is wanted rather than an asymptotic one (e.g.\ bounding
$|f(\bar x+h) - T_k(h)|$ over an entire region at once, given a bound on the $(k{+}1)$-th
derivatives there). The price is that $\theta$ depends on $\bar x$, $h$ \emph{and} $k$ together,
and is almost never computable in practice. That is exactly why the rest of this chapter uses the
asymptotic form for actually computing expansions, and reserves this one for error estimates.
\end{remark}

% Quelle: Linus Lüchinger/Slides/Slides-03-16.pdf, slide 12;
%         Linus Lüchinger/Slides/Slides-03-19.pdf, slides 4--6
% Generator: Claude Opus 5 (High)
\begin{corollary}[A polynomial that fits to order $k$ \emph{is} the Taylor polynomial]
\label{cor:taylor_expansion_unique}
Let $U \subseteq \mathbb{R}^n$ be open, $x_0 \in U$, $f \in C^{k+1}(U)$, and let $P$ be a
polynomial of degree at most $k$ with
\[\lvert f(x_0+h) - P(h)\rvert = o\big(\lvert h\rvert^{k}\big) \qquad \text{as } h \to 0.\]
Then $\partial^\alpha f(x_0) = \partial^\alpha P(0)$ for every multi-index with
$\lvert\alpha\rvert \leq k$. Equivalently, $P$ is \emph{the} Taylor polynomial of $f$ of degree
$k$ at $x_0$.
\end{corollary}

\begin{proof}
By \cref{thm:taylor_working_formula}, the Taylor polynomial $T$ of degree $k$ also satisfies
$\lvert f(x_0+h)-T(h)\rvert = o(\lvert h\rvert^k)$. Subtracting, the polynomial $Q := P-T$ has
degree at most $k$ and satisfies $\lvert Q(h)\rvert = o(\lvert h\rvert^k)$. Suppose
$Q \not\equiv 0$ and let $d \leq k$ be the degree of its lowest non-vanishing homogeneous part
$Q_d$. Choosing $v$ with $Q_d(v) \neq 0$ and restricting to the ray $h := tv$,
\[Q(tv) = t^dQ_d(v) + O(t^{d+1}), \qquad\text{so}\qquad
  \frac{\lvert Q(tv)\rvert}{t^{k}} \sim \lvert Q_d(v)\rvert\, t^{d-k} \not\to 0\]
as $t \downarrow 0$, since $d \leq k$. That contradicts $\lvert Q(h)\rvert = o(\lvert h\rvert^k)$,
so $Q \equiv 0$ and $P = T$.
\end{proof}

\begin{remark}[Reading derivatives off an asymptotic expansion]
\label{rem:reading_derivatives_off_expansions}
This is more useful than it looks, because it runs \emph{backwards}. Normally one computes
derivatives in order to produce an expansion; the corollary says that an expansion, however it
was obtained, hands back the derivatives. Given
\[f(x_0+h) = c + a\transp h + \tfrac12 h\transp Bh + o(\lvert h\rvert^2)\]
one may read off $f(x_0) = c$, $\nabla f(x_0) = a$ and $\Hess f(x_0) = B$ without differentiating
anything. This is the same move by which \cref{def:differentiable} recognises a differential from
an $o(\lvert h\rvert)$ estimate, now carried to higher order. It is the standard route to a Hessian
for functions built out of known series.

Two cautions on the bookkeeping, both of which the exercise below turns on. The estimate must be
$o(\lvert h\rvert^{k})$, not $O(\lvert h\rvert^{k})$: an $O(\lvert h\rvert^{k})$ error can
contain a genuine degree-$k$ term and so corrupts the order-$k$ coefficients, leaving only the
derivatives of order $< k$ determined. And the expansion determines derivatives only up to
order $k$, and nothing beyond it.
\end{remark}

% Supplement: Linus Lüchinger/Slides/Slides-03-19.pdf, slide 6
\begin{exercise}[What does an expansion actually tell you?]
\label{ex:reading_off_expansions}
For each of the following, state precisely which values of $f$ and of its derivatives at the
indicated point are determined by the given information, and which are not.
\begin{enumerate}[label=\textbf{(\alph*)}]
  \item $f : \mathbb{R}^3 \to \mathbb{R}$ with
    $f(x_1,x_2,x_3) - x_1x_2 \in O\big(\lvert x\rvert\big)$, at $x = 0$.
  \item $f : \mathbb{R}^2 \to \mathbb{R}$ with
    $f(1+x_1,\,-1+x_2) - 2 - x_1 \in o\big(\lvert x\rvert\big)$, at $(1,-1)$.
  \item $f : \mathbb{R}^2 \to \mathbb{R}$ with
    $f(-1+x_1,\,x_2) + 2 - 2x_1x_2 + 2x_2^2 - 3x_2^3 \in o\big(\lvert x\rvert^3\big)$,
    at $(-1,0)$.
\end{enumerate}
Assume in each case that $f$ is as smooth as the relevant statement requires.
\end{exercise}

% Supplement: Toby Lane/class-document.pdf, §3.4.4
\begin{remark}[Two ways to weight the same expansion]
\Cref{ex:taylor_maximal_expansion} above and \cref{thm:taylor_working_formula} appear to disagree
on how heavily mixed partial derivatives are weighted. Take the coefficient of $h_1^2h_2$, so
$\alpha = (2,1)$ and $\ell = |\alpha| = 3$. The explicit multi-index expansion gives $\tfrac12$
directly. The working formula instead offers $\tfrac{1}{\ell!}(h_1\partial_1+h_2\partial_2)^\ell$,
and expanding that by the multinomial theorem produces
\[
  \frac{\ell!}{\alpha!}\cdot\frac{1}{\ell!} \;=\; \frac{1}{\alpha!}
  \;=\; \frac{1}{2!\,1!} \;=\; \frac12
\]
for the term $h^\alpha\partial^\alpha f$. The two agree, but they arrive by different routes, and
seeing why is worth the paragraph.

The working formula's $\tfrac{1}{\ell!}$ weights \emph{all} $\ell$-fold products of partials
globally, including repeated ones such as $\partial_1\partial_2$, which appears in more than one
order in the multinomial expansion of $(h_1\partial_1+h_2\partial_2)^\ell$. By Schwarz's lemma
(equality of mixed partial derivatives, valid here because $f \in C^k$) all those repeats collapse
onto the same $\partial^\alpha$, and there are exactly $\ell!/\alpha!$ of them. That collapsing is
what turns the global $1/\ell!$ into the per-term $1/\alpha!$ seen in
\Cref{ex:taylor_maximal_expansion}.

In short, $1/\ell!$ weights \emph{paths} through the mixed partial derivatives, whereas
$1/\alpha!$ weights the \emph{distinct} derivatives themselves, once those paths have been
identified with one another.
\end{remark}

% Sonnet 5 (Medium)
\subsection{A worked example}

% Source: Corsin Nick/Class Notes/Week 4.pdf, pp. 11--12
\begin{exercise}[Taylor expansion of $(x+y^2)e^{-x^2-y^2}$]
\label{ex:taylor_explicit_example}
Compute the Taylor expansion of
\[f : \mathbb{R}^2\to\mathbb{R}, \qquad (x,y) \mapsto (x+y^2)e^{-x^2-y^2}\]
around $(0,0)$ up to third order.
\end{exercise}

% Kept inline: solved directly in Corsin Nick/Class Notes/Week 4.pdf, pp. 11--12.
\begin{exercisesolution}[\cref{ex:taylor_explicit_example}, cumbersome solution]
Calculate all partial derivatives up to third order and use the definition. (Please never do
this.)
\end{exercisesolution}

\begin{exercisesolution}[\cref{ex:taylor_explicit_example}, the easy solution]
The substitution and multiplication rules we used in Analysis 1 still apply:
\[
\begin{aligned}
(x+y^2)e^{-x^2-y^2} &= (x+y^2)\left(1 - x^2 - y^2 + o(|x^2+y^2|)\right) \\
&= x + y^2 - x^3 - xy^2 + o\!\left(|(x,y)|^3\right)
\end{aligned}
\]
\end{exercisesolution}

% Supplement: Sascha Brack/Class Notes/Week_05_Notes_Monday.pdf, pp. 9--11
\begin{example}[More Taylor expansions via Analysis 1]
\label{ex:more_taylor_via_analysis1}
The substitution and multiplication rules from Analysis 1 often make expansions much faster than using the formula directly, provided the expansion is around $(0,0)$.
\begin{enumerate}[label=\textbf{(\alph*)}]
  \item Expand $f(x,y) := \cos(x)e^y$ up to second order around $(0,0)$:
  \[f(x,y) = \left(1 - \frac{x^2}{2} + o(|(x,y)|^2)\right)\left(1 + y + \frac{y^2}{2} + o(|(x,y)|^2)\right) = 1 + y - \frac{1}{2}x^2 + \frac{1}{2}y^2 + o(|(x,y)|^2).\]
  \item Expand $f(x,y) := \cos(x+y^2)$ up to second order around $(0,0)$. Setting $z := x+y^2$:
  \[\cos(z) = 1 - \frac{z^2}{2} + o(|z|^2) \implies \cos(x+y^2) = 1 - \frac{(x+y^2)^2}{2} + o(|(x,y)|^2) = 1 - \frac{1}{2}x^2 + o(|(x,y)|^2),\]
  where terms like $xy^2$ and $y^4$ are absorbed into $o(|(x,y)|^2)$ since their total degree is strictly greater than $2$.
  \item Expand $f(x,y) := \sqrt{y} e^{\sqrt{x}}$ to first order around $(1,1)$. Here the formula is better since the expansion point is not the origin.
  We have $f(1,1) = e$.
  $\partial_1 f(x,y) = \sqrt{y} e^{\sqrt{x}} \frac{1}{2\sqrt{x}} \implies \partial_1 f(1,1) = \frac{e}{2}$.
  $\partial_2 f(x,y) = \frac{1}{2\sqrt{y}} e^{\sqrt{x}} \implies \partial_2 f(1,1) = \frac{e}{2}$.
  Writing $h = (h_1,h_2) = (x-1, y-1)$:
  \[f(1+h_1, 1+h_2) = f(1,1) + \partial_1 f(1,1)h_1 + \partial_2 f(1,1)h_2 + o(|h|) = e + \frac{e}{2}h_1 + \frac{e}{2}h_2 + o(|h|).\]
\end{enumerate}
\end{example}

% Sonnet 5 (Medium)
\subsection{Second-order form and the Hessian}

% Source: Corsin Nick/Class Notes/Week 4.pdf, p. 12
\begin{remark}[Second-order Taylor formula with Hessian]
\label{rem:second_order_taylor_hessian}
The second-order Taylor approximation of $f \in C^2(\mathbb{R}^n,\mathbb{R})$ around
$x_0 \in \mathbb{R}^n$ can be expressed compactly as
\[f(x_0+x) = f(x_0) + \nabla f(x_0)\cdot x + \tfrac{1}{2}x\transp \cdot \Hess f(x_0)\cdot x + o(|x|^2)\]
where
\[\Hess f(x_0) = \big(\partial_i\partial_j f(x_0)\big)_{i,j=1,\dots,n}\]
is the \newterm{Hessian matrix} \germanfor{Hesse-Matrix} of $f$.
\end{remark}

% Generator: Claude Opus 5 (High)
This is the form the rest of the course actually uses, and it is worth reading term by term as
a sequence of ever-better approximations to $f$ near $x_0$:
\begin{itemize}
  \item the \textbf{constant} term $f(x_0)$, the crudest guess, exact only at $x_0$;
  \item the \textbf{linear} term $\nabla f(x_0)\cdot x$, which is the tangent plane and hence
    exactly the differential of \cref{def:differentiable};
  \item the \textbf{quadratic} term $\tfrac12 x\transp\Hess f(x_0)\,x$, the first term that sees
    curvature, and the one that decides in \cref{sec:hessian_test} whether a critical point is a
    minimum, a maximum, or a saddle.
\end{itemize}
By \cref{thm:schwarz_mixed_partials} the Hessian is symmetric, which is what makes that
last decision a question about the \emph{signs of its eigenvalues} rather than anything harder.

\begin{center}
% Generator: Claude Opus 5 (High) --- FIG-W04-05: the three approximations, in one variable,
% where they can actually be drawn. f(x) = 1 + 0.55x - 0.5x^2 + 0.28x^3 about x_0 = 0, so
% f(0) = 1, f'(0) = 0.55, f''(0) = -1: constant 1, tangent 1 + 0.55x, parabola 1 + 0.55x - 0.5x^2.
\begin{tikzpicture}[>=Latex, scale=1.25]
  \draw[->] (-1.7,0) -- (2.1,0) node[right, font=\footnotesize] {$x$};
  \draw[->] (0,-0.3) -- (0,2.6) node[above, font=\footnotesize] {$y$};
  \draw (0,0.05) -- (0,-0.05) node[below, font=\scriptsize] {$x_0$};

  % order 0
  \draw[gray!70, thick, dashed] (-1.6,1) -- (2.0,1);
  \node[gray!70!black, font=\scriptsize, anchor=west] at (2.02,1.0) {order $0$};

  % order 1: tangent
  \draw[green!50!black, thick, domain=-1.6:2.0, samples=2]
    plot (\x, {1 + 0.55*\x});
  \node[green!50!black, font=\scriptsize, anchor=west] at (2.02,2.1) {order $1$};

  % order 2: parabola
  \draw[orange!90!black, thick, domain=-1.3:1.85, samples=60, smooth]
    plot (\x, {1 + 0.55*\x - 0.5*\x*\x});
  \node[orange!90!black, font=\scriptsize, anchor=west] at (1.88,0.40) {order $2$};

  % the function itself
  \draw[blue!75!black, very thick, domain=-1.35:1.75, samples=90, smooth]
    plot (\x, {1 + 0.55*\x - 0.5*\x*\x + 0.28*\x*\x*\x});
  \node[blue!75!black, font=\scriptsize, anchor=south east] at (-0.85,1.05) {$f$};

  \fill (0,1) circle (1.1pt);

  \node[font=\scriptsize, align=center, anchor=north] at (0.3,-0.62)
    {each order matches one more derivative at $x_0$,\\
     and the error drops from $O(|x|)$ to $O(|x|^2)$ to $O(|x|^3)$};
\end{tikzpicture}
\end{center}

In several variables the same picture holds with \qt{tangent line} replaced by \emph{tangent
plane} and \qt{parabola} by a paraboloid, bowl-shaped or saddle-shaped according to the
signs of the eigenvalues of $\Hess f(x_0)$. That is exactly the classification carried out in
\cref{sec:hessian_test}.

% Source: Corsin Nick/Class Notes/Week 4.pdf, p. 12
% (Re-cited: the reading of the three terms and the figure above are additions.)
\begin{question}
For $f \in C^k(U, \mathbb{R})$, do the Taylor expansions of $t \mapsto f(x_0+th)$ in $t$ and of
$(th) \mapsto f(x_0+th)$ in $(th)$ around $0$ coincide?
\end{question}

\begin{answer}
\textbf{Yes.} This is, in fact, how we prove the multi-dimensional Taylor approximation, and why
we need the assumption $x_0 + th \in U$ for all $t \in [0,1]$.
\end{answer}

% Source: Corsin Nick/Class Notes/Week 4.pdf, p. 13
\begin{exercise}[Proof of second-order Taylor formula]
\label{ex:hessian_taylor_proof}
Prove the formula from above:
\[f(x_0+h) = f(x_0) + \nabla f(x_0)h + \tfrac{1}{2}h\transp \Hess f(x_0)h + \dots\]
for $f \in C^k(U,\mathbb{R})$ with $x_0 + th \in U$ for all $t \in [0,1]$.
\end{exercise}

% Kept inline: solved directly in Corsin Nick/Class Notes/Week 4.pdf, p. 13.
\begin{exercisesolution}[Proof of \cref{ex:hessian_taylor_proof}]
We expand $t \mapsto f(x_0+th)$ in $t$:
\[
\begin{aligned}
f(x_0+th) &= f(x_0) + t\left.\frac{d}{dt}\right|_{t=0}f(x_0+th) + \frac{t^2}{2}\left.\frac{d^2}{dt^2}\right|_{t=0}f(x_0+th) + o(|th|^2) \\
&= f(x_0) + t\sum_{i=1}^{n}\partial_i f(x_0)h_i + \frac{t^2}{2}\frac{d}{dt}\left.\sum_{i=1}^{n}\partial_i f(x_0+th)h_i\right|_{t=0} + o(|h|^2) \\
&= f(x_0) + t\nabla f(x_0)h + \frac{t^2}{2}\sum_{i,j=1}^{n}\partial_i\partial_j f(x_0)h_ih_j + o(|th|^2) \\
&= f(x_0) + t\nabla f(x_0)h + \frac{t^2}{2}h\transp \Hess f(x_0)h + o(|th|^2)
\end{aligned}
\]
Setting $t = 1$ gives the result.
\end{exercisesolution}


% --- exercises relocated here from the dissolved sheets ---

% Originally: content/week-05.tex, \exercisesheet{5}, problem 5.3
% Source: exercises/Ex5_Analysis2_eng.pdf

\begin{exercise}[Multiple choice (Landau notation) \textnormal{(important)}]
\label{ex:5.3}
Mark all and only the statements which are true.
\begin{enumerate}[label=\textbf{(\alph*)}]
  \item As $|x| \to 0$, if $f(x) = x_2x_1^2 + O(x_1^2)$, then $f(x) = O(|x|^2)$.
  \item As $x \to 0$, if $f(x) = x_2x_1^2 + O(x_1^3)$, then $f(x) = O(x_1^3)$.
  \item As $|x| \to 0$, if $f(x) = x_2x_1 + O(|x|^2)$, then $f$ is differentiable at $x = 0$.
  \item As $x \to 0$, if $f(x) = x_2x_1 + O(|x|^2)$, then $f$ is twice differentiable around
    $x = 0$.
  \item As $|x| \to 0$, if $f(x) = x_2x_1 + O(|x|^3)$, then $f$ is twice differentiable around
    $x = 0$ and $\partial_{11}f(0) = 0$.
\end{enumerate}
\exhint[Official hint]{Revise \textit{Corollary 10.45}.}
\exinfo{This exercise is Problem 5.3 of Problem Sheet 5, Analysis II, Spring Semester 2026.}
\end{exercise}

% Originally: content/exercise-sheets/week-05.tex, problem 5.5 --- never previously built
% Source: exercises/Ex5_Analysis2_eng.pdf

\begin{exercise}[Multiple choice (cost of a Taylor expansion) \textnormal{(optional)}]
\label{ex:5.5}
Assume $g \in C^\infty(\mathbb{R})$ and $f \in C^\infty(\mathbb{R}^n)$ is such that
$$f(x) = x_1 + x_2x_1 + O(|x|^4) \quad \text{as } x \to 0.$$
We want to compute the Taylor expansion of $(g\circ f)$ at $x = 0$ of the highest possible order
with the information we have. We can ask for the value $g^{(k)}(t)$ for arbitrary
$k \in \mathbb{N}$, $t \in \mathbb{R}$, but we have to pay 5 CHF each time.

For a general $g$, which is the degree of the best (i.e.\ highest-order) Taylor polynomial we can
compute? How expensive will it be to compute it?
\begin{enumerate}[label=\textbf{(\alph*)}]
  \item degree 3 and 20 CHF
  \item degree 2 and 15 CHF
  \item degree 2 and 10 CHF
  \item degree 3 and 15 CHF
\end{enumerate}
\exinfo{This exercise is Problem 5.5 of Problem Sheet 5, Analysis II, Spring Semester 2026.}
\end{exercise}

% Originally: content/exercise-sheets/week-05.tex, problem 5.6 --- never previously built
% Source: exercises/Ex5_Analysis2_eng.pdf

\begin{exercise}[Polynomials \textnormal{(optional --- ``interesting, but irrelevant for the course'')}]
\label{ex:5.6}
Let $P, Q \in \mathbb{R}[x_1,\dots,x_n]$ be polynomials of $n$ variables; assume that there are
positive numbers $M, \sigma$ such that
$$|P(x) - Q(x)| \leq M|x|^\sigma \quad \text{for all } |x| \leq 1.$$
Show that $P$ and $Q$ have the same coefficients of order smaller than $\sigma$. That is, if we
write $P(X) = \sum_{\alpha\in\mathbb{N}^n} a_\alpha X^\alpha$,
$Q(X) = \sum_{\alpha\in\mathbb{N}^n} b_\alpha X^\alpha$, then $a_\alpha = b_\alpha$ for all
$|\alpha| < \sigma$.
\exinfo{This exercise is Problem 5.6 of Problem Sheet 5, Analysis II, Spring Semester 2026.}
\end{exercise}

% Source: Jérôme Paschoud/Notizen/Woche 5.1 Taylor.pdf, p. 3
% Generator: Gemini 3.6 Flash (High)
\begin{exercise}[Evaluating partial derivatives using multi-index notation]
\label{ex:multi_index_evaluation}
Let $f(x,y,z) := x^3+2y+e^{2x}\cos(z)$. Calculate the partial derivative $\partial^\alpha f$ for the multi-index $\alpha = (2,0,1)$.
\end{exercise}

\begin{exercisesolution}[\cref{ex:multi_index_evaluation}]
The multi-index $\alpha = (2,0,1)$ corresponds to differentiating twice with respect to $x$, zero times with respect to $y$, and once with respect to $z$. Since $f$ is smooth, the order of differentiation does not matter.
\begin{align*}
\partial_x f &= 3x^2 + 2e^{2x}\cos(z) \\
\partial_{xx} f &= 6x + 4e^{2x}\cos(z) \\
\partial_{xxz} f &= -4e^{2x}\sin(z).
\end{align*}
Hence $\partial^{(2,0,1)} f(x,y,z) = -4e^{2x}\sin(z)$.
\end{exercisesolution}

% Source: Jérôme Paschoud/Notizen/Woche 5.1 Taylor.pdf, pp. 3--4
% Generator: Gemini 3.6 Flash (High)
\begin{exercise}[Taylor expansion of $\sin(xy)$]
\label{ex:taylor_expansion_sin_xy}
Compute the Taylor expansion of $f(x,y) := \sin(xy)$ to order $3$ around $(0,0)$. 
\end{exercise}

\begin{exercisesolution}[\cref{ex:taylor_expansion_sin_xy}]
We can solve this in two ways:
\begin{enumerate}[label=\textbf{(\alph*)}]
  \item \textbf{Using the general formula:}
    We list the multi-indices $\alpha$ with $|\alpha| \leq 3$ and compute the partial derivatives at $(0,0)$:
    \begin{itemize}
      \item $|\alpha|=0$: $f(0,0) = \sin(0) = 0$.
      \item $|\alpha|=1$: $\partial_x f = y\cos(xy) \implies 0$. $\partial_y f = x\cos(xy) \implies 0$.
      \item $|\alpha|=2$: $\partial_{xx} f = -y^2\sin(xy) \implies 0$. $\partial_{yy} f = -x^2\sin(xy) \implies 0$. $\partial_{xy} f = \cos(xy) - xy\sin(xy) \implies 1$.
      \item $|\alpha|=3$: $\partial_{xxx} f = -y^3\cos(xy) \implies 0$. $\partial_{yyy} f = -x^3\cos(xy) \implies 0$. $\partial_{xxy} f = -2y\sin(xy) - xy^2\cos(xy) \implies 0$. $\partial_{xyy} f = -2x\sin(xy) - x^2y\cos(xy) \implies 0$.
    \end{itemize}
    The only non-zero term is $\partial^{(1,1)} f(0,0) = 1$. The corresponding factor in the formula is $\frac{1}{\alpha!}h^\alpha = \frac{1}{(1!)(1!)}xy = xy$. Thus,
    \[f(x,y) = xy + o(|(x,y)|^3).\]
  \item \textbf{Using 1D expansions (the clever way):}
    We know that $\sin(t) = t - \frac{t^3}{6} + O(t^5)$ as $t \to 0$. Substituting $t = xy$ yields
    \[\sin(xy) = xy - \frac{(xy)^3}{6} + O((xy)^5) = xy + o(|(x,y)|^3).\]
\end{enumerate}
\end{exercisesolution}

% Source: Jérôme Paschoud/Notizen/Woche 5.1 Taylor.pdf, p. 4
% Generator: Gemini 3.6 Flash (High)
\begin{exercise}[Taylor expansion of $\sqrt{1+x+y^2}$]
\label{ex:taylor_expansion_sqrt_1_plus_x_plus_y2}
Compute the Taylor expansion of $f(x,y) := \sqrt{1+x+y^2}$ to order $2$ around $(0,0)$.
\end{exercise}

\begin{exercisesolution}[\cref{ex:taylor_expansion_sqrt_1_plus_x_plus_y2}]
Using the 1D expansion $\sqrt{1+t} = 1 + \frac{1}{2}t - \frac{1}{8}t^2 + O(t^3)$ with $t = x+y^2$, we obtain:
\begin{align*}
\sqrt{1+x+y^2} &= 1 + \frac{1}{2}(x+y^2) - \frac{1}{8}(x+y^2)^2 + O(|(x,y)|^3) \\
&= 1 + \frac{1}{2}x + \frac{1}{2}y^2 - \frac{1}{8}(x^2+2xy^2+y^4) + o(|(x,y)|^2) \\
&= 1 + \frac{1}{2}x - \frac{1}{8}x^2 + \frac{1}{2}y^2 + o(|(x,y)|^2).
\end{align*}
\end{exercisesolution}

% Supplement: Analysis_II_Script_v1.pdf, p. 59
\begin{exercise}[Taylor expansion of $\exp(\arctan(x-y))$]
\label{ex:taylor_expansion_exp_arctan}
Compute the Taylor expansion of $f(x,y) := \exp(\arctan(x-y))$ to order $2$ around $(0,0)$.
\end{exercise}

\begin{exercisesolution}[\cref{ex:taylor_expansion_exp_arctan}]
% Generator: Claude Sonnet 5 (High)
Write $s := x-y$, itself of order $|h|$. Since $\arctan$ is odd, $\arctan(t) = t + O(t^3)$, so
\[\arctan(s) = s + O(|h|^3) = s + o(|h|^2).\]
Since $\exp$ is smooth and bounded near $0$, an $o(|h|^2)$ correction inside $\exp$ contributes
only an $o(|h|^2)$ correction to the output. Consequently, to order $2$, $\exp(\arctan(s))$ and
$\exp(s)$ agree:
\[\exp(\arctan(x-y)) = \exp(x-y) + o(|(x,y)|^2).\]
Now $\exp(s) = 1+s+\tfrac12s^2+O(s^3)$ with $s^2 = x^2-2xy+y^2$, so
\[\exp(\arctan(x-y)) = 1 + x - y + \tfrac12x^2 - xy + \tfrac12y^2 + o(|(x,y)|^2).\]
\end{exercisesolution}

% Supplement: Analysis_II_Script_v1.pdf, p. 59
\begin{exercise}[Taylor expansion of $1/(1-x^2-y^2)$]
\label{ex:taylor_expansion_one_over_1_minus_x2_y2}
Compute the Taylor expansion of $f(x,y) := \dfrac{1}{1-x^2-y^2}$ to order $2$ around $(0,0)$.
\end{exercise}

\begin{exercisesolution}[\cref{ex:taylor_expansion_one_over_1_minus_x2_y2}]
% Generator: Claude Sonnet 5 (High)
Using $\tfrac{1}{1-t} = 1+t+O(t^2)$ with $t := x^2+y^2$, itself already $O(|h|^2)$, so
$t^2 = O(|h|^4) = o(|h|^2)$:
\[\frac{1}{1-x^2-y^2} = 1 + (x^2+y^2) + o(|(x,y)|^2).\]
\end{exercisesolution}

\begin{remark}
Together with \cref{ex:taylor_expansion_sin_xy,ex:taylor_expansion_sqrt_1_plus_x_plus_y2} above,
these two exercises are exactly script Exercise 10.36: \qt{compute the Taylor polynomials up
to the quadratic order of $\sin(xy)$, $\sqrt{1+x+y^2}$, $\exp(\arctan(x-y))$,
$1/(1-x^2-y^2)$ --- don't use the general formula!} All four are worked by the same \qt{1D
substitution} technique this section uses throughout.
\end{remark}

% Source: Jérôme Paschoud/Notizen/Woche 5.1 Taylor.pdf, p. 4
% Generator: Gemini 3.6 Flash (High)
\begin{exercise}[Taylor expansion of $\frac{\cos x}{1-y^2}$]
\label{ex:taylor_expansion_cos_over_1_minus_y2}
Compute the Taylor expansion of $f(x,y) := \frac{\cos x}{1-y^2}$ to order $4$ around $(0,0)$.
\end{exercise}

\begin{exercisesolution}[\cref{ex:taylor_expansion_cos_over_1_minus_y2}]
We multiply the 1D expansions for $\cos x$ and $\frac{1}{1-y^2}$.
As $x \to 0$, $\cos x = 1 - \frac{x^2}{2} + \frac{x^4}{24} + O(x^6)$.
As $y \to 0$, $\frac{1}{1-y^2} = 1 + y^2 + y^4 + O(y^6)$.
Multiplying these together and keeping terms up to order $4$:
\begin{align*}
\frac{\cos x}{1-y^2} &= \left(1 - \frac{x^2}{2} + \frac{x^4}{24} + O(x^6)\right)\left(1 + y^2 + y^4 + O(y^6)\right) \\
&= 1 + y^2 + y^4 - \frac{x^2}{2} - \frac{x^2y^2}{2} + \frac{x^4}{24} + o(|(x,y)|^4) \\
&= 1 - \frac{1}{2}x^2 + y^2 + \frac{1}{24}x^4 - \frac{1}{2}x^2y^2 + y^4 + o(|(x,y)|^4).
\end{align*}
\end{exercisesolution}

% Source: old_exams/fs2023/Prüfung.pdf (9 August 2021), Teil B, Aufgabe 1, p. 13
% Extractor: Gemini 3.6 Flash (High)
% Correction: Claude Opus 5 (High) --- part (c) restored to the source's "all critical points";
% Flash had weakened it to "in a neighbourhood of x = 0", and then answered neither version. Part
% (d) restored: the source asks about x = 1 as well, and that is the interesting one.
\begin{exercise}[Taylor expansion of a parametric integral]
\label{ex:taylor_parametric_integral}
Consider the function $f : \mathbb{R} \to \mathbb{R}$ defined by the parametric integral
\[f(x) := \int_{x-1}^{x+1} \cos(\pi t^2) \, dt.\]
The integral need not be evaluated explicitly.
\begin{enumerate}[label=\textbf{(\alph*)}]
  \item Compute $f'(x)$, and simplify it to a product of two sines.
  \item Find the coefficients of the fourth-order Taylor polynomial of $f$ around $x = 0$,
  \[f(x) = f(0) + a_1 x + a_2 x^2 + a_3 x^3 + a_4 x^4 + O(x^5).\]
  \item Determine \emph{all} critical points of $f$.
  \item Decide, for each of the critical points $x = 0$ and $x = 1$, whether it is a local maximum,
    a local minimum, or neither.
\end{enumerate}
\exinfo{This exercise is taken from the exam of 9 August 2021, Part B, Problem 1. The original
opens with a part asking whether $f$ is continuously differentiable, bounded, and even.}
\end{exercise}


% Originally: new. Problem 4.5 was transcribed into content/exercise-sheets/week-04.tex but
% that file was never \input by anything, so the problem had never appeared in the built
% document and carried no solution.
\section{Solutions}
\label{sec:taylor_solutions}

\begin{exercisesolution}[Solution to \cref{ex:4.5}]
% Generator: Claude Opus 5 (High)
\textbf{No} --- and the official hint's function is the cleanest witness. Take
\[f(t) := \big(\sin(2\pi t),\ \cos(2\pi t)\big), \qquad f \in C^1(\mathbb{R},\mathbb{R}^2).\]
Then $f(1) = (0,1) = f(0)$, so the left-hand side is
\[f(1) - f(0) = (0,0).\]
On the other hand $f'(t) = 2\pi\big(\cos(2\pi t),\ -\sin(2\pi t)\big)$, whose norm is
\[\lVert f'(t)\rVert = 2\pi\sqrt{\cos^2(2\pi t) + \sin^2(2\pi t)} = 2\pi \neq 0
  \qquad \text{for every } t.\]
So $f'(t)$ never equals $(0,0)$, and no $t \in [0,1]$ can satisfy the claimed identity.

The reason is the one given in \cref{rem:mvt_fails_vector_valued}: applying the scalar mean value
theorem componentwise produces one intermediate point \emph{per component},
\[f_1(1)-f_1(0) = f_1'(\xi_1), \qquad f_2(1)-f_2(0) = f_2'(\xi_2),\]
and nothing forces $\xi_1 = \xi_2$. Here $\xi_1 = \tfrac14$ and $\xi_2 = \tfrac12$ work
componentwise --- both components do vanish, just not at the same time.

\begin{ainote}
Note what does \emph{not} fail. The \emph{mean value inequality}
$\lVert f(1)-f(0)\rVert \leq \sup_{t \in [0,1]}\lVert f'(t)\rVert$ holds here as
$0 \leq 2\pi$, comfortably. That inequality is what survives into $\mathbb{R}^m$ and what every
later estimate in this document actually uses; the \emph{equality} is a one-dimensional accident.
Geometrically the counterexample is a walker going once round a circle: they return to the start,
so the total displacement is zero, but they never once stand still.
\end{ainote}
\end{exercisesolution}



\part{Optimization and Convexity}

% Generator: Claude Opus 5 (High) --- part-level framing prose, new content.
Differentiability turns the search for extrema into algebra. An interior extremum forces the
gradient to vanish, so the candidates can be found by solving equations, and the sign behaviour of
the Hessian then decides which kind of point each candidate is. Two complications occupy the rest
of the part. When the variables are constrained to a level set, the gradient need not vanish at an
extremum at all, and the method of Lagrange multipliers says what condition replaces it. When the
Hessian test comes out inconclusive, or when the question asked is global rather than local,
convexity is the hypothesis that settles it: for a convex function a critical point is not a
candidate but an answer.

% ============================================================
% Chapter 12 --- Extrema and the Hessian test
% Originally: content/week-05.tex, \section{Optimization} and \section{The Hessian test}
%             (Corsin Week 5), with problems 6.2, 6.8, 6.9 and 6.12 from sheet 6.
% ============================================================
\chapter{Extrema and the Hessian Test}
\label{ch:extrema}

% Transition: Claude Opus 5 (Medium)
The Taylor expansion of the previous chapter stops being a formal exercise the moment one asks
what it is \emph{for}. Its first-order term locates the candidates for a maximum or a minimum ---
the points where the gradient vanishes --- and its second-order term decides which of them are
which. This chapter sets up both halves: the vocabulary of local and global extrema, and the
Hessian test that classifies a critical point from the sign of a quadratic form, together with
the practical question of how to read that sign off a matrix.

% Originally: content/week-05.tex, \section{Optimization} (Corsin Week 5)

% Source: Corsin Nick/Class Notes/Week 5.pdf, p. 2
\section{Optimization}
\label{sec:optimization}

\begin{center}
\begin{tikzpicture}[scale=1.2]
  \draw[very thick] (0, 1) .. controls (1, 3) and (2, -1) .. (3, 0)
                    .. controls (4, 1) and (4.5, 4) .. (5, 3.5)
                    .. controls (5.5, 3) and (6, 0.5) .. (6.5, 0.5);
                    
  \node[above, font=\footnotesize] at (0.8, 2.1) {local max};
  \draw[->, thick, shorten >= 2pt] (0.8, 2.1) -- (0.65, 1.85);
  
  \node[below, font=\footnotesize] at (2.2, -0.6) {global min};
  \draw[->, thick, shorten >= 2pt] (2.2, -0.6) -- (1.9, -0.3);
  
  \node[above, font=\footnotesize] at (4.6, 4.1) {global max};
  \draw[->, thick, shorten >= 2pt] (4.6, 4.1) -- (4.65, 3.8);
  
  \node[below, font=\footnotesize] at (6.2, 0.1) {local min};
  \draw[->, thick, shorten >= 2pt] (6.2, 0.1) -- (6.3, 0.45);
\end{tikzpicture}
\end{center}

\begin{definition}[Local extrema and critical points]
\label{def:local_extrema_critical_points}
\begin{itemize}
  \item \newterm{Local minimum} (\germanterm{lokales Minimum}) of
    $f : U \subseteq \mathbb{R}^n \to \mathbb{R}$: $f$ has a local minimum at $x_0 \in U$ if
    \[\exists\,\varepsilon > 0 \text{ such that } \forall x \in B_\varepsilon(x_0)\cap U : f(x) \geq f(x_0).\]
  \item \textbf{Local maxima} are defined analogously with $f(x) \leq f(x_0)$.
  \item The extremum is \textbf{strict} if the inequality is strict for $x \neq x_0$:
    $f(x) > f(x_0)$ for a strict minimum, $f(x) < f(x_0)$ for a strict maximum.
  \item A \newterm{critical point} (\germanterm{kritischer Punkt}) of $f \in C^1(U,\mathbb{R})$,
    $U \subseteq \mathbb{R}^n$ open, is a point $x_0 \in U$ where $\nabla f(x_0) = 0$.
\end{itemize}
\end{definition}

\begin{proposition}[Local extremum $\implies$ critical point]
\label{prop:extremum_implies_critical}
If $f \in C^1(U,\mathbb{R})$ and $x_0 \in \operatorname{int}(U)$ is a local extremum, then
$\nabla f(x_0) = 0$.
\end{proposition}

% Generator: Claude Opus 5 (Medium) --- Corsin states this without proof; the one-variable
% reduction is two lines and is the argument every later Lagrange computation leans on.
\begin{proof}
Fix a direction $v \in \mathbb{R}^n$ and consider the one-variable function
$\varphi(t) := f(x_0 + tv)$, defined for $|t|$ small enough that $x_0 + tv$ stays in $U$ --- this
is where $x_0 \in \operatorname{int}(U)$ is used, and it is the only place. If $f$ has a local
extremum at $x_0$, then $\varphi$ has one at the \emph{interior} point $t=0$, so by the
one-variable result from Analysis I, $\varphi'(0) = 0$. By the chain rule along a path
(\cref{sec:chain_rule_along_path}),
\[0 = \varphi'(0) = \langle \nabla f(x_0), v\rangle.\]
Since $v$ was arbitrary, $\nabla f(x_0)$ is orthogonal to every vector of $\mathbb{R}^n$,
including itself: $\langle\nabla f(x_0),\nabla f(x_0)\rangle = |\nabla f(x_0)|^2 = 0$, hence
$\nabla f(x_0) = 0$.
\end{proof}

\begin{remark}[Intuition, and two standing warnings]
The proof says: at an interior extremum every direction is a \qt{flat} direction, and the only
vector orthogonal to everything is the zero vector,
\[\langle\nabla f(x_0), v\rangle = 0 \text{ for all } v \in \mathbb{R}^n
  \iff \nabla f(x_0) = 0.\]
Two things this proposition does \textbf{not} say, both of which the rest of the chapter is
built around.
\begin{itemize}
  \item \textbf{The converse is false.} A critical point need not be an extremum:
    $f(x,y) := x^2-y^2$ has $\nabla f(0,0)=0$ but a saddle there, and even in one variable
    $f(x) := x^3$ has $f'(0)=0$ with no extremum. Critical points are \emph{candidates};
    separating them is exactly what \cref{thm:hessian_test} is for.
  \item \textbf{Interiority is essential.} On $U := [0,1]$ the function $f(x) := x$ attains both
    its minimum and its maximum, at $x=0$ and $x=1$, and $f' \equiv 1$ never vanishes. The
    hypothesis $x_0 \in \operatorname{int}(U)$ is what fails. This is why every optimization over
    a closed region in this chapter is done in two steps --- interior critical points first, then
    the boundary separately, the latter being what Lagrange multipliers handle.
\end{itemize}
\end{remark}

\begin{center}
\begin{tikzpicture}[>=Latex, scale=1.3]
  % Grid on xy-plane
  \begin{scope}[x={(0.866cm,-0.3cm)}, y={(0.5cm,0.4cm)}]
    % Plane axes
    \draw[thick, ->] (-1.5,-1.5) -- (2.2,-1.5) node[right] {$x$};
    \draw[thick, ->] (-1.5,-1.5) -- (-1.5,2.2) node[above] {$y$};
    \draw[gray!40, thin] (-1.5,-0.5) -- (1.8,-0.5);
    \draw[gray!40, thin] (-1.5,0.5) -- (1.8,0.5);
    \draw[gray!40, thin] (-1.5,1.5) -- (1.8,1.5);
    \draw[gray!40, thin] (-0.5,-1.5) -- (-0.5,1.8);
    \draw[gray!40, thin] (0.5,-1.5) -- (0.5,1.8);
    \draw[gray!40, thin] (1.5,-1.5) -- (1.5,1.8);
    % Projected level set
    \draw[purple, thick] (0,0) circle (1.0);
    \fill[green!60!black] (0.707,-0.707) circle (2pt) node[below right, font=\tiny] {minimum};
    \draw[orange!90!black, thick, ->] (0.707,-0.707) -- (0.1,-0.707) node[below, font=\tiny] {gradient};
  \end{scope}

  % Paraboloid surface
  \draw[blue!80!black, thick] (0,0.8) ellipse (1.3 and 0.5);
  \draw[blue!80!black, thick] (-1.3,0.8).. controls (-1.0,-0.6) and (1.0,-0.6).. (1.3,0.8);
  \draw[purple, thick, dashed] (0,0.1) ellipse (0.85 and 0.32);

  % Connecting dotted line & steepest ascent arrow
  \draw[orange!90!black, thick, dotted] (0.55,-0.55) -- (0.55,0.1);
  \draw[red, thick, ->] (0.55,0.1).. controls (0.75,0.3).. (0.95,0.75) node[above left, font=\tiny] {steepest ascent};
  \node[blue!80!black, right] at (1.3,0.8) {$f(x,y)$};
  \node[purple, right, font=\tiny] at (0.85,0.1) {level sets};
\end{tikzpicture}
\end{center}

% Sonnet 5 (Medium)
Unconstrained extrema are found by \cref{prop:extremum_implies_critical}. We now ask the same
question when $x$ is restricted to a lower-dimensional constraint set, where the gradient
condition $\nabla f(x_0) = 0$ is generally too strong to hold anywhere.

% Originally: content/week-05.tex, \section{The Hessian test} (Corsin Week 5)

% Source: Corsin Nick/Class Notes/Week 5.pdf, p. 7
\section{The Hessian test}
\label{sec:hessian_test}

% Sonnet 5 (Medium): "Hessian" was already introduced in \cref{rem:second_order_taylor_hessian}
% (\cref{ch:taylor}); referencing it here instead of re-introducing the term with \newterm.
The \textbf{Hessian} (\cref{rem:second_order_taylor_hessian}) of $f \in C^2(U,\mathbb{R})$,
$U \subseteq \mathbb{R}^n$, is given by
\[\Hess f(x) = (\partial_i\partial_j f)_{i,j=1,\dots,n}\]
with respect to the standard basis of $\mathbb{R}^n$. By \newterm{Schwarz's lemma}
(\germanterm{Satz von Schwarz}), stated as \cref{thm:schwarz_mixed_partials},
$\Hess f(x)$ is symmetric.

% Supplement: Lukas Krause/Week 5 March 16-20 Notes Lukas Krause.pdf, pp. 2--3
\begin{ainote}
Symmetry of the Hessian is the hypothesis every statement in this section rests on --- the
spectral theorem in \cref{rem:diagonalizing_hessian}, the eigenvalue criteria, the Hurwitz
minors --- and neither Corsin nor \cref{thm:schwarz_mixed_partials} proves it. Lukas Krause's notes give a
complete proof, and unusually he leaves his \emph{failed} first attempt in place, with the
comment that it is instructive to see why the obvious calculation does not work. Both are
reproduced below, because the failure is genuinely the more informative half.
\end{ainote}

\begin{proof}[Proof of \cref{thm:schwarz_mixed_partials}]
It suffices to treat $m=1$ (apply the result to each component), $i<j$ (there is nothing to
prove when $i=j$), and $n=2$ with $i=1$, $j=2$ (freeze all the other variables). So let
$f \in C^2$ near $(x,y)$ and write $\partial_1,\partial_2$ for the two partial derivatives.

\textit{The attempt that fails.} Unwinding both mixed derivatives as iterated limits gives
\[\partial_2\partial_1f(x,y)
  = \lim_{h_2\to0}\lim_{h_1\to0}
    \frac{f(x+h_1,y+h_2) - f(x+h_1,y) - f(x,y+h_2) + f(x,y)}{h_1h_2},\]
and one would like to apply the mean value theorem in each variable and be done. The obstruction
is that the two intermediate points the mean value theorem returns depend on $h_1$ and $h_2$
\emph{separately} and need not agree, so forcing them together leaves a stray factor
$h_2/h_1$, which blows up when $h_1 \to 0$ at fixed $h_2$. The symmetric-looking expression
above is not enough on its own; one has to keep the two increments coupled.

\textit{The proof.} Couple them by fixing a single $h \neq 0$ and setting
\[F(h) := f(x+h,y+h) - f(x+h,y) - \big(f(x,y+h) - f(x,y)\big),\]
which is the second difference of $f$ over a square of side $h$, and
$\varphi(t) := f(x+th, y+h) - f(x+th, y)$, so that $F(h) = \varphi(1)-\varphi(0)$. The mean
value theorem gives $\xi_1 \in (0,1)$ with
\[F(h) = \varphi'(\xi_1) = h\,\partial_1f(x+\xi_1h,\ y+h) - h\,\partial_1f(x+\xi_1h,\ y).\]
Now apply the mean value theorem a second time, to $\psi(t) := \partial_1f(x+\xi_1h,\ y+th)$ ---
legitimate precisely because $f \in C^2$, so $\partial_1 f$ is itself continuously
differentiable in the second variable. Since $F(h) = h\big(\psi(1)-\psi(0)\big)$, there is
$\xi_2 \in (0,1)$ with
\[F(h) = h\,\psi'(\xi_2) = h^2\,\partial_2\partial_1f(x+\xi_1h,\ y+\xi_2h).\]
Grouping the four terms of $F(h)$ the other way --- first and third together, second and fourth
together --- runs the identical argument with the roles of the variables exchanged and produces
$\tilde\xi_1,\tilde\xi_2 \in (0,1)$ with
$F(h) = h^2\,\partial_1\partial_2f(x+\tilde\xi_1h,\ y+\tilde\xi_2h)$. Dividing both expressions
by $h^2 \neq 0$,
\[\partial_2\partial_1f(x+\xi_1h,\ y+\xi_2h)
  \ =\ \partial_1\partial_2f(x+\tilde\xi_1h,\ y+\tilde\xi_2h).\]
All four numbers $\xi_1,\xi_2,\tilde\xi_1,\tilde\xi_2$ lie in $(0,1)$ --- bounded, though they
depend on $h$ --- so both arguments converge to $(x,y)$ as $h \to 0$. Continuity of the mixed
second partials (again $f \in C^2$) lets us pass to the limit on both sides, giving
$\partial_2\partial_1f(x,y) = \partial_1\partial_2f(x,y)$.
\end{proof}

\begin{remark}[What the failed attempt shows]
\label{rem:why_schwarz_needs_care}
The two attempts differ in exactly one respect: the failed one lets $h_1$ and $h_2$ go to zero
\emph{independently}, the successful one ties them together as a single $h$ and only takes a
limit at the very end. That is why the second difference $F(h)$ over a \emph{square} is the
right object --- it is symmetric in the two variables by construction, so the two groupings of
its four terms compute the two mixed partials, and the symmetry of the answer is inherited from
the symmetry of the square rather than proved by force.

It also locates precisely where continuity is used: twice, once to apply the mean value theorem
to $\partial_1f$ and once to pass to the limit at the end. Drop it and the conclusion genuinely
fails --- see \cref{rem:schwarz_needs_c2} for the standard counterexample, which is
also the function Lukas sets as an exercise alongside this proof.
\end{remark}

\begin{definition}[Sign of a symmetric matrix]
\label{def:sign_symmetric_matrix}
A symmetric matrix $A \in \mathbb{R}^{n\times n}$ is:
\begin{itemize}
  \item \newterm{positive definite} (\germanterm{positiv definit}) if $v\transp Av > 0$ for all
    $v \in \mathbb{R}^n\setminus\{0\}$;
  \item \newterm{negative definite} (\germanterm{negativ definit}) if $v\transp Av < 0$ for all
    $v \in \mathbb{R}^n\setminus\{0\}$;
  \item \textbf{indefinite} if there exist $v, u \in \mathbb{R}^n\setminus\{0\}$ such that
    $v\transp Av > 0$ and $u\transp Au < 0$;
  \item \textbf{degenerate} (\germanterm{entartet}) if there exists
    $v \in \mathbb{R}^n\setminus\{0\}$ such that $Av = 0$; equivalently, if $\det A = 0$, or
    if $0$ is an eigenvalue of $A$.
\end{itemize}
\end{definition}

\begin{ainote}
Corsin defines degenerate as \qt{there exists $v \neq 0$ with $v\transp Av = 0$}. That
condition is much weaker than intended: if $A$ is indefinite, with $v\transp Av > 0$ and
$u\transp Au < 0$, then $t \mapsto (tv+(1-t)u)\transp A(tv+(1-t)u)$ is continuous and changes
sign, so it vanishes at some $w \neq 0$ --- making \emph{every} indefinite matrix degenerate,
and the third bullet of \cref{thm:hessian_test} vacuous. The intended notion is a nontrivial
kernel ($\det A = 0$), which is also what \cref{ex:6.2} states. Corrected above.

\end{ainote}

% Supplement: Sascha Brack/Class Notes/Week_05_Notes_Friday.pdf, pp. 3--4
\begin{remark}[Linear algebra interlude: diagonalizing the Hessian]
\label{rem:diagonalizing_hessian}
Since the Hessian $\Hess f(x_0)$ is a symmetric matrix, it is diagonalizable over $\mathbb{R}$ by the spectral theorem. There exist $n$ real eigenvalues $\lambda_1, \dots, \lambda_n \in \mathbb{R}$ and a corresponding orthonormal basis of eigenvectors $v_1, \dots, v_n \in \mathbb{R}^n$ such that $\Hess f(x_0) v_i = \lambda_i v_i$ for all $i$.

Every vector $x \in \mathbb{R}^n$ can be written in this basis as $x = \sum_{i=1}^n a_i v_i$ with $a_i \in \mathbb{R}$. Applying the Hessian yields:
\[
  \Hess f(x_0) x = \Hess f(x_0) \sum_{i=1}^n a_i v_i = \sum_{i=1}^n a_i \Hess f(x_0) v_i = \sum_{i=1}^n \lambda_i a_i v_i.
\]
In this eigenbasis, the Hessian matrix becomes diagonal. Consequently, the quadratic form evaluates to $x\transp \Hess f(x_0) x = \sum_{i=1}^n \lambda_i a_i^2$, which immediately links the definiteness of the Hessian to the signs of its eigenvalues.
\end{remark}

\begin{theorem}[The Hessian test]
\label{thm:hessian_test}
If $f \in C^2(U,\mathbb{R})$ has a critical point $x_0 \in \operatorname{int}(U)$, then:
\begin{itemize}
  \item $\Hess f(x_0)$ \textbf{positive definite} $\Rightarrow$ local \textbf{min} at $x_0$;
  \item $\Hess f(x_0)$ \textbf{negative definite} $\Rightarrow$ local \textbf{max} at $x_0$;
  \item $\Hess f(x_0)$ \textbf{indefinite} $\Rightarrow$ \newterm{saddle point}
    (\germanterm{Sattelpunkt}) at $x_0$.
\end{itemize}
If $\Hess f(x_0)$ is degenerate but not indefinite --- that is, positive or negative
semidefinite with $0$ an eigenvalue --- the test gives \textbf{no information}, and all three
behaviours are possible.
\end{theorem}

\begin{ainote}
Corsin's third bullet reads \qt{indefinite and not degenerate}. With degeneracy corrected to
$\det A = 0$ the extra hypothesis is unnecessary: indefiniteness alone already forces a saddle,
since directions $v$ with $v\transp\Hess f(x_0)v > 0$ and $u$ with $u\transp\Hess f(x_0)u < 0$
make $t\mapsto f(x_0+tv)$ increase and $t\mapsto f(x_0+tu)$ decrease away from $x_0$. Dropped
above, and the genuinely inconclusive case --- semidefinite \emph{and} degenerate --- stated
explicitly instead, since that is the one students actually run into
(\cref{ex:6.8} is exactly this trap).
\end{ainote}

% Moved here from the Lagrange-multiplier section (review item D2): it uses the Hessian,
% "indefinite" and the eigenvalue criterion, none of which were available at its former
% position 260 lines earlier.
\begin{aiexample}[Classifying saddle points via the Hessian]
\label{ex:ai_saddle_via_hessian}
% Generator: Gemini 3.6 Flash (Medium)
Consider $f(x,y) := x^2-y^2$ on $\mathbb{R}^2$. The gradient
$\nabla f(x,y) = (2x,\,-2y)\transp$ vanishes at $(0,0)$, and the Hessian is constant:
\[
\Hess f(0,0) = \begin{pmatrix} 2 & 0 \\ 0 & -2 \end{pmatrix}.
\]
Its eigenvalues are $\lambda_1 = 2 > 0$ and $\lambda_2 = -2 < 0$, so $\Hess f(0,0)$ is
indefinite and $(0,0)$ is a saddle point by \cref{thm:hessian_test} --- the model case the word
\qt{saddle} is named after, since $f$ rises along the $x$-axis and falls along the $y$-axis.
\end{aiexample}

\begin{remark}[Mnemonic]
positive $= \smile \dots$ min; negative $= \frown \dots$ max; indefinite $=$ saddle.
\end{remark}

\begin{center}
% Generator: Claude Opus 5 (High) --- FIG-W05-01: the three conclusive cases of
% \cref{thm:hessian_test}, drawn as the model quadratic each one is.
% Left:  x^2 + y^2, Hess = diag(2,2), both eigenvalues > 0 -- a bowl.
% Mid:  -x^2 - y^2, Hess = diag(-2,-2), both < 0 -- a dome.
% Right: x^2 - y^2, Hess = diag(2,-2), opposite signs -- a saddle (\cref{ex:ai_saddle_via_hessian}).
% In each panel the two coloured curves are the profiles along the two eigendirections; the
% signs of the eigenvalues are exactly the up/down curvature of those two profiles.
\begin{tikzpicture}[>=Latex, scale=1.0]

  % ---------- positive definite: bowl ----------
  \begin{scope}
    \draw[gray!55] (-1.5,0) -- (1.5,0);
    \draw[green!50!black, very thick, domain=-1.25:1.25, samples=60, smooth]
      plot (\x, {0.62*\x*\x});
    \draw[orange!90!black, very thick, domain=-1.25:1.25, samples=60, smooth]
      plot (\x, {0.30*\x*\x});
    \fill (0,0) circle (1.5pt);
    \node[font=\scriptsize, anchor=south west] at (0.06,0.02) {$x_0$};
    % Caption baseline shared with the other two panels so the row reads straight.
    \node[font=\small, align=center, anchor=north] at (0,-1.15)
      {\textbf{positive definite}\\ $\lambda_1,\lambda_2>0$\\ local \textbf{min}};
  \end{scope}

  % ---------- negative definite: dome ----------
  \begin{scope}[shift={(4.2,0)}]
    \draw[gray!55] (-1.5,0) -- (1.5,0);
    \draw[green!50!black, very thick, domain=-1.25:1.25, samples=60, smooth]
      plot (\x, {-0.62*\x*\x});
    \draw[orange!90!black, very thick, domain=-1.25:1.25, samples=60, smooth]
      plot (\x, {-0.30*\x*\x});
    \fill (0,0) circle (1.5pt);
    \node[font=\scriptsize, anchor=south west] at (0.06,0.06) {$x_0$};
    \node[font=\small, align=center, anchor=north] at (0,-1.15)
      {\textbf{negative definite}\\ $\lambda_1,\lambda_2<0$\\ local \textbf{max}};
  \end{scope}

  % ---------- indefinite: saddle ----------
  \begin{scope}[shift={(8.4,0)}]
    \draw[gray!55] (-1.5,0) -- (1.5,0);
    \draw[green!50!black, very thick, domain=-1.25:1.25, samples=60, smooth]
      plot (\x, {0.62*\x*\x});
    \draw[orange!90!black, very thick, domain=-1.25:1.25, samples=60, smooth]
      plot (\x, {-0.62*\x*\x});
    \fill (0,0) circle (1.5pt);
    \node[font=\scriptsize, anchor=south west] at (0.08,0.04) {$x_0$};
    \node[green!50!black, font=\scriptsize, anchor=west] at (1.28,0.90) {up};
    \node[orange!90!black, font=\scriptsize, anchor=west] at (1.28,-0.90) {down};
    \node[font=\small, align=center, anchor=north] at (0,-1.15)
      {\textbf{indefinite}\\ opposite signs\\ \textbf{saddle}};
  \end{scope}
\end{tikzpicture}
\end{center}

% Generator: Claude Opus 5 (High)
The picture is the whole theorem. By \cref{rem:diagonalizing_hessian} the second-order Taylor
term is $\tfrac12\sum_i\lambda_ia_i^2$ in the eigenbasis, so along the $i$-th eigendirection $f$
behaves like the parabola $\tfrac12\lambda_it^2$: it curves \emph{up} when $\lambda_i>0$ and
\emph{down} when $\lambda_i<0$. All up gives a minimum, all down a maximum, and one of each
gives a saddle --- there is a direction of ascent and a direction of descent through the same
point, so it is neither. And when some $\lambda_i = 0$ that direction is \emph{flat} to second
order, the quadratic term says nothing about it, and the test must fall silent --- which is
exactly the degenerate case, illustrated by $x^4+y^4$ against $x^4-y^4$ above.

\begin{ainote}
The regularity in \cref{thm:hessian_test} is stated by Corsin as $C^3$; the test needs only
$C^2$ (which is what the Hessian's definition assumes). Relaxed to $C^2$ above.
\end{ainote}

\begin{aiexample}[Inconclusive Hessian Test at a Degenerate Critical Point]
% Generator: Gemini 3.6 Flash (Medium)
Consider $f(x,y) := x^4 + y^4$ and $g(x,y) := x^4 - y^4$ at the critical point $(0,0)$.
\begin{itemize}
  \item For both functions, all second-order partial derivatives at $(0,0)$ vanish, so $\Hess f(0,0) = \Hess g(0,0) = \left(\begin{smallmatrix} 0 & 0 \\ 0 & 0 \end{smallmatrix}\right)$.
  \item The Hessian test is \textbf{inconclusive} since $\det \Hess = 0$.
  \item However, direct inspection reveals that $f(x,y) \geq 0 = f(0,0)$ for all $(x,y)$, so $(0,0)$ is a strict local minimum for $f$, whereas $g(x,0) > 0$ and $g(0,y) < 0$ for non-zero $x,y$, so $(0,0)$ is a saddle point for $g$.
\end{itemize}
\end{aiexample}

% Source: Corsin Nick/Class Notes/Week 5.pdf, pp. 8--9
\subsection{How do I determine positive/negative definiteness?}

Let $A \in \mathbb{R}^{n\times n}$ be symmetric with eigenvalues $\lambda_1,\dots,\lambda_n$.
Note: $\det(A) = \lambda_1\cdots\lambda_n$ and $\Tr(A) = \lambda_1+\dots+\lambda_n$.

\textbf{For $n = 2$:}
\begin{itemize}
  \item $\det(A) = \lambda_1\lambda_2 < 0 \implies$ \textbf{indefinite}
  \item $\det(A) > 0$, $\Tr(A) > 0 \implies$ \textbf{positive definite}
  \item $\det(A) > 0$, $\Tr(A) < 0 \implies$ \textbf{negative definite}
\end{itemize}

\textbf{For $n \geq 3$:}
\begin{itemize}
  \item Solve the characteristic equation $\det(A - \lambda\id) = 0$ for $\lambda$ and check all
    eigenvalues.
  \item \newterm{Hurwitz criterion} (\germanterm{Hurwitz-Kriterium}): for
    \[A = \begin{pmatrix} a_{11} & a_{12} & \cdots & a_{1n} \\ a_{21} & \ddots & & \vdots \\ \vdots & & \ddots & \\ a_{n1} & \cdots & & a_{nn}\end{pmatrix}\]
    define the leading principal minors
    \[A_1 = \det(a_{11}), \quad A_2 = \det\begin{pmatrix} a_{11} & a_{12} \\ a_{21} & a_{22}\end{pmatrix}, \quad \dots, \quad A_n = \det(A).\]
    Then:
    \begin{itemize}
      \item $A_1 > 0$, $A_2 > 0$, $A_3 > 0$, \ldots\ if and only if the matrix is \textbf{positive definite}
      \item $A_1 < 0$, $A_2 > 0$, $A_3 < 0$, \ldots\ if and only if the matrix is \textbf{negative definite}
      \item \textbf{indefinite} if neither $A_1 \geq 0, A_2 \geq 0, \dots$ nor
        $A_1 \leq 0, A_2 \geq 0, A_3 \leq 0, \dots$
    \end{itemize}
\end{itemize}

\begin{ainote}
Corsin writes the characteristic equation as ``$A - \lambda\id = 0$'', omitting the determinant.
Corrected above.
\end{ainote}

% Source: Corsin Nick/Class Notes/Week 5.pdf, pp. 9--10
\begin{example}[Definiteness calculations]
\begin{center}
\begin{tabular}{lll}
\textbf{Matrix} & \textbf{Computation} & \textbf{Conclusion} \\
\hline
$\left(\begin{smallmatrix}2&2\\2&5\end{smallmatrix}\right)$ & $\det = 6 > 0$, $\Tr = 7 > 0$ & positive definite \\
$\left(\begin{smallmatrix}-4&2\\2&-4\end{smallmatrix}\right)$ & $\det = 12 > 0$, $\Tr = -8 < 0$ & negative definite \\
$\left(\begin{smallmatrix}4&2\\2&-4\end{smallmatrix}\right)$ & $\det = -20 < 0$ & indefinite \\
$\left(\begin{smallmatrix}-5&0&0\\0&-4&-2\\0&-2&-4\end{smallmatrix}\right)$ & $A_1 = -5 < 0$, $A_2 = 20 > 0$, $A_3 = -60 < 0$ & negative definite \\
$\left(\begin{smallmatrix}-1&0&-4\\0&2&-2\\-4&-2&22\end{smallmatrix}\right)$ & $A_1 = -1 < 0$, $A_2 = -2 < 0$, $A_3 = -72 < 0$ & indefinite \\
\end{tabular}
\end{center}
\end{example}

% Source: Corsin Nick/Class Notes/Week 5.pdf, pp. 10--11
\begin{exercise}[Critical values of $x^3 - y^3 + 3\alpha xy$]
\label{ex:cubic_saddle_family}
Find the critical values of
\[f_\alpha : \mathbb{R}^2 \to \mathbb{R}, \qquad (x,y) \mapsto x^3 - y^3 + 3\alpha xy\]
for $\alpha \in \mathbb{R}\setminus\{0\}$ and determine min/max/saddle points.
\end{exercise}

% Kept inline: solved directly in Corsin Nick/Class Notes/Week 5.pdf, pp. 10--11.
\begin{exercisesolution}[Solution to \cref{ex:cubic_saddle_family}]
\[
\begin{aligned}
\frac{\partial f_\alpha}{\partial x} &= 3x^2 + 3\alpha y \overset{!}{=} 0 &&\implies y = -\frac{x^2}{\alpha} \\
\frac{\partial f_\alpha}{\partial y} &= -3y^2 + 3\alpha x \overset{!}{=} 0 &&\implies -\frac{x^4}{\alpha^2} + \alpha x = 0
\end{aligned}
\]
\begin{enumerate}[label=\textbf{\arabic*.}]
  \item $x = y = 0$
  \item $x \neq 0 \implies x^3 = \alpha^3 \implies x = \alpha$, $y = -\alpha$
\end{enumerate}
So, we have critical points $(0,0)$ and $(\alpha, -\alpha)$.

Compute the Hessian:
\[\begin{pmatrix} \partial_x^2 f & \partial_x\partial_y f \\ \partial_x\partial_y f & \partial_y^2 f\end{pmatrix} = \begin{pmatrix} 6x & 3\alpha \\ 3\alpha & -6y \end{pmatrix} = \Hess f(x,y)\]
\[\Hess f(0,0) = \begin{pmatrix} 0 & 3\alpha \\ 3\alpha & 0\end{pmatrix}, \qquad \det = -9\alpha^2 < 0 \implies \text{indefinite}\]
\begin{align*}
\Hess f(\alpha,-\alpha) &= \begin{pmatrix} 6\alpha & 3\alpha \\ 3\alpha & 6\alpha\end{pmatrix}
  = 3\alpha\begin{pmatrix} 2 & 1 \\ 1 & 2\end{pmatrix}, \\
\det &= 27\alpha^2 > 0, \qquad \Tr = 12\alpha.
\end{align*}
For $\alpha > 0$ this is positive definite; for $\alpha < 0$ it is negative definite. So:
\begin{itemize}
  \item $(0,0) \longrightarrow$ \textbf{saddle point}
  \item $(\alpha,-\alpha) \longrightarrow$ \textbf{local minimum} for $\alpha > 0$; \textbf{local
    maximum} for $\alpha < 0$
\end{itemize}
\end{exercisesolution}

% Originally: content/week-05.tex, end of the chapter

\begin{aiexercise}[Global Extrema on a Compact Triangle]
\label{ex:ai_extrema_triangle}
% Generator: Gemini 3.6 Flash (Medium)
Find the global maximum and minimum values of $f(x,y) := xy$ on the compact triangular region $T := \{(x,y) \in \mathbb{R}^2 \mid x \ge 0, y \ge 0, x+y \le 2\}$.
\end{aiexercise}


% --- exercises relocated here from the dissolved sheets ---

% Originally: content/week-06.tex, \exercisesheet{6}, problem 6.2
% Source: exercises/Ex6_Analysis2_eng.pdf

\begin{exercise}[6.2 --- The signature of a $2\times 2$ matrix \textnormal{(important)}]
\label{ex:6.2}
Despite the definition, it is not necessary to compute the eigenvalues of a matrix to find its
signature. Prove that for a symmetric $2\times 2$ matrix $M$ we have the following simple rule to
determine the signature in terms of $\det M$ and $\Tr M$:
\begin{itemize}
  \item If $\det M > 0$, $\Tr M > 0$ then $M$ is positive definite,
  \item If $\det M > 0$, $\Tr M < 0$ then $M$ is negative definite,
  \item If $\det M < 0$ then $M$ is indefinite,
  \item If $\det M = 0$, then $M$ is degenerate.
\end{itemize}
\end{exercise}

% Originally: content/week-06.tex, \exercisesheet{6}, problem 6.8
% Source: exercises/Ex6_Analysis2_eng.pdf

\begin{exercise}[6.8 --- Multiple choice \textnormal{(important)}]
\label{ex:6.8}
The Hessian matrix of $f \in C^2(\mathbb{R}^n)$ is positive semidefinite at a critical point $x_0$
of $f$, i.e.,
\[\langle v, \Hess f(x_0) v\rangle \geq 0 \quad \text{for all } v \in \mathbb{R}^n.\]
Which of the following statements necessarily hold? (There may be more than one).
\begin{enumerate}[label=\textbf{(\alph*)}]
  \item $x_0$ is a strict local minimum of $f$.
  \item $x_0$ is a local minimum of $f$.
  \item $x_0$ is not a local maximum of $f$.
  \item None of the above statements.
\end{enumerate}
\end{exercise}

% Originally: content/week-06.tex, \exercisesheet{6}, problem 6.9
% Source: exercises/Ex6_Analysis2_eng.pdf

\begin{exercise}[6.9 --- Minimization \textnormal{(important)}]
\label{ex:6.9}
The function $f : \mathbb{R}^2 \to \mathbb{R}$ is given by $f(x,y) = 2x^2+y^2-x$. Determine the
extrema of $f$ on\dots
\begin{enumerate}[label=\textbf{(\alph*)}]
  \item \dots the unit circle $S^1 = \{(x,y) \in \mathbb{R}^2 \mid x^2+y^2 = 1\}$;
  \item \dots the closed unit disk $D = \{(x,y) \in \mathbb{R}^2 \mid x^2+y^2 \leq 1\}$.
\end{enumerate}
\end{exercise}

% Originally: content/week-06.tex, \exercisesheet{6}, problem 6.12
% Source: exercises/Ex6_Analysis2_eng.pdf

\begin{exercise}[6.12 --- Critical points \textnormal{(important)}]
\label{ex:6.12}
Let $f : \mathbb{R}^2 \to \mathbb{R}$ be the function $f(x,y) = (ax^2+by^2)e^{-x^2-y^2}$ with real
parameters $a, b \in \mathbb{R}$. Find all critical points and determine their nature with the
Hessian test, depending on $a, b$.
\end{exercise}

% Originally: the \section{Solutions} of content/week-05.tex and content/week-06.tex

\section{Solutions}
\label{sec:extrema_solutions}

\begin{exercisesolution}[Proof of \cref{ex:ai_extrema_triangle}]
% Generator: Gemini 3.6 Flash (Medium)
Since $T$ is compact and $f$ is continuous, global extrema exist by the extreme value theorem
(\cref{item:extreme_value_theorem}).
\begin{itemize}
  \item \textbf{Interior of $T$:} Critical points satisfy $\nabla f(x,y) = (y,x)\transp = (0,0)\transp$, and therefore $(x,y) = (0,0)$, which lies on the boundary. Thus, there are no critical points in $\operatorname{int}(T)$.
  \item \textbf{Boundary of $T$:}
    \begin{itemize}
      \item Along the leg $x = 0$ ($0 \le y \le 2$), $f(0,y) = 0$.
      \item Along the leg $y = 0$ ($0 \le x \le 2$), $f(x,0) = 0$.
      \item Along the hypotenuse $x+y = 2$ ($0 \le x \le 2$), substitute $y = 2-x$:
        \[h(x) := x(2-x) = 2x-x^2, \qquad h'(x) = 2-2x,\]
        which vanishes at $x = 1$. This gives the candidate point $(1,1)$ with $f(1,1) = 1$. The endpoints $(2,0)$ and $(0,2)$ yield $f = 0$.
    \end{itemize}
\end{itemize}
Comparing values, the \textbf{global maximum} is $1$ at $(1,1)$, and the \textbf{global minimum} is $0$, attained everywhere on the coordinate axes legs of $T$.
\end{exercisesolution}

\begin{exercisesolution}[Proof of \cref{ex:6.2}]
% Generator: Claude Sonnet 5 (Medium)
By the spectral theorem, since $M$ is symmetric there exists an orthogonal matrix $O$
(so $OO\transp = \id$, hence $\det O = \pm 1$) and $\lambda_1, \lambda_2 \in \mathbb{R}$ such that
\[OMO\transp = \begin{pmatrix} \lambda_1 & 0 \\ 0 & \lambda_2 \end{pmatrix}.\]
Using $\det(AB) = \det(A)\det(B)$ and $\Tr(AB) = \Tr(BA)$,
\[\det M = \det(OMO\transp) = \lambda_1\lambda_2, \qquad \Tr M = \Tr(MO\transp O) = \Tr(OMO\transp) = \lambda_1 + \lambda_2.\]
So, $\det M$ and $\Tr M$ are exactly the product and sum of the eigenvalues, and the four cases
follow immediately:
\begin{itemize}
  \item $\det M > 0$: $\lambda_1, \lambda_2$ have the same sign. If additionally $\Tr M > 0$,
    both are positive, so $M$ is \textbf{positive definite}; if $\Tr M < 0$, both are negative,
    so $M$ is \textbf{negative definite}.
  \item $\det M < 0$: $\lambda_1, \lambda_2$ have opposite signs, so $M$ is \textbf{indefinite}.
  \item $\det M = 0$: one of $\lambda_1, \lambda_2$ is zero, so $M$ is \textbf{degenerate}.
\end{itemize}
\end{exercisesolution}

\begin{exercisesolution}[Solution to \cref{ex:6.8}]
% Generator: Claude Sonnet 5 (Medium)
None of \textbf{(a)}--\textbf{(c)} hold in general:
\begin{itemize}
  \item \textbf{(a) False:} $f \equiv 0$ has $\Hess f(x_0) = 0$ positive semidefinite at every
    $x_0$, but $x_0$ is not a \emph{strict} local minimum (every nearby value equals $f(x_0)$).
  \item \textbf{(b) False:} $f(x) = x_1^3$ at $x_0 = 0$ has $\Hess f(0) = 0$ (positive
    semidefinite), but $f$ takes negative values arbitrarily close to $0$, so $x_0$ is not a
    local minimum at all.
  \item \textbf{(c) False:} $f(x) = -x_1^4$ at $x_0 = 0$ has $\Hess f(0) = 0$ (positive
    semidefinite, vacuously, since it is the zero matrix), yet $x_0$ \emph{is} a (strict) local
    maximum of $f$.
\end{itemize}
Since \textbf{(a)}--\textbf{(c)} all fail, \textbf{(d) is the correct answer}: none of the above
statements necessarily hold.
\end{exercisesolution}

\begin{exercisesolution}[Solution to \cref{ex:6.9}]
% Generator: Claude Sonnet 5 (Medium)
\begin{enumerate}[label=\textbf{(\alph*)}]
  \item With $g(x,y) := x^2+y^2-1$, the Lagrangian is $L(x,y,\lambda) := f(x,y) - \lambda g(x,y)$,
    giving the system
    \[0 = \partial_x L = 2x(2-\lambda) - 1, \qquad 0 = \partial_y L = 2y(2-2\lambda), \qquad 0 = \partial_\lambda L = -(x^2+y^2-1).\]
    From the middle equation, either $y = 0$ or $\lambda = 1$.

    \textbf{If $y = 0$:} then $x^2 = 1$, so $x = \pm 1$, giving the points $(1,0)$ and $(-1,0)$
    with $f(1,0) = 1$ and $f(-1,0) = 3$.

    \textbf{If $\lambda = 1$:} the first equation gives $2x(1) = 1$, i.e.\ $x = \tfrac12$, and the
    constraint gives $y = \pm\tfrac{\sqrt3}{2}$, with $f\big(\tfrac12,\pm\tfrac{\sqrt3}{2}\big) =
    \tfrac12 + \tfrac34 - \tfrac12 = \tfrac34$.

    Comparing all four values, $f$ attains a \textbf{global maximum} at $(-1,0)$ (value $3$) and
    \textbf{global minima} at $\big(\tfrac12,\pm\tfrac{\sqrt3}{2}\big)$ (value $\tfrac34$); $(1,0)$
    is a local (non-global) extremum with value $1$.
  \end{enumerate}
\end{exercisesolution}

\begin{ainote}
The official solution lists $(1,0)$ alongside $(-1,0)$ as \qt{local extrema} from the Lagrange
system, then only mentions it again implicitly by omission from the final global-extrema summary.
It is worth flagging explicitly that \textbf{Lagrange's condition finds all \emph{critical} points
of $f|_{S^1}$, not automatically the global extrema}. Here $(1,0)$ is a genuine local maximum along
the circle (value $1$) that is simply not the largest one; only comparing all candidate values
at the end distinguishes it from the global maximum at $(-1,0)$ (value $3$).
\end{ainote}

\begin{exercisesolution}[Solution to \cref{ex:6.9}, continued]
% Generator: Claude Sonnet 5 (Medium)
\begin{enumerate}[label=\textbf{(\alph*)}, start=2]
  \item Having handled the boundary $S^1 = \partial D$ in \textbf{(a)}, we look for interior
    critical points: $Df(x,y) = (4x-1, 2y) = 0$ gives $(x,y) = \big(\tfrac14, 0\big) \in
    \operatorname{int}(D)$. The Hessian $\Hess f = \left(\begin{smallmatrix} 4 & 0 \\ 0 & 2\end{smallmatrix}\right)$ is
    positive definite, so this is a local minimum, with value $f\big(\tfrac14,0\big) =
    \tfrac18 - \tfrac14 = -\tfrac18$. Since $-\tfrac18$ is smaller than every boundary value found
    in \textbf{(a)}, $f$ on $D$ attains a \textbf{global minimum} at $\big(\tfrac14,0\big)$ and a
    \textbf{global maximum} at $(-1,0)$.
\end{enumerate}
\end{exercisesolution}

\begin{exercisesolution}[Solution to \cref{ex:6.12}]
% Generator: Claude Sonnet 5 (Medium)
Writing $r^2 := x^2+y^2$, the gradient is
\[\nabla f(x,y) = \big(2x(a-ax^2-by^2),\ 2y(b-ax^2-by^2)\big)e^{-r^2}.\]

\textbf{Case $a=b=0$:} $f \equiv 0$, so every point of $\mathbb{R}^2$ is critical, and (being the
constant function) each is simultaneously a global maximum and minimum.

\textbf{Case $a=0, b\neq 0$ (and symmetrically $a \neq 0, b=0$):} the gradient equations reduce to
$-2xby^2=0$ and $2by(1-y^2)=0$, whose common solutions are the entire line $\{(x,0): x\in
\mathbb{R}\}$ together with the two points $(0,\pm1)$. On the critical line the Hessian is
singular (degenerate, as expected since a whole line is critical); at $(0,\pm1)$,
$\Hess f(0,\pm1) = \operatorname{diag}(-2b,-4b)\,e^{-1}$, giving \textbf{minima} if $b<0$ and
\textbf{maxima} if $b>0$. The case $a\neq0,b=0$ is symmetric, with critical line $\{(0,y)\}$ and
extra points $(\pm1,0)$: minima if $a<0$, maxima if $a>0$.

\textbf{Case $a,b \neq 0$, $a \neq b$:} setting $x=0$ or $y=0$ in the gradient equations gives the
five critical points $(0,0)$, $(0,\pm1)$, $(\pm1,0)$ (checking $x,y\neq0$ simultaneously forces
$a=b$, excluded here). The Hessians are
\begin{align*}
\Hess f(0,0) &= \begin{pmatrix}2a&0\\0&2b\end{pmatrix}, \\
\Hess f(0,\pm1) &= e^{-1}\begin{pmatrix}2(a-b)&0\\0&-4b\end{pmatrix}, \\
\Hess f(\pm1,0) &= e^{-1}\begin{pmatrix}-4a&0\\0&2(b-a)\end{pmatrix}.
\end{align*}
Reading off signs for each ordering of $a,b$ against $0$:
\begin{itemize}
  \item $a>b>0$: min at $(0,0)$, saddle at $(0,\pm1)$, max at $(\pm1,0)$.
  \item $b>a>0$: min at $(0,0)$, max at $(0,\pm1)$, saddle at $(\pm1,0)$.
  \item $a>0>b$: saddle at $(0,0)$, min at $(0,\pm1)$, max at $(\pm1,0)$.
  \item $b>0>a$: saddle at $(0,0)$, max at $(0,\pm1)$, min at $(\pm1,0)$.
  \item $0>a>b$: max at $(0,0)$, min at $(0,\pm1)$, saddle at $(\pm1,0)$.
  \item $0>b>a$: max at $(0,0)$, saddle at $(0,\pm1)$, min at $(\pm1,0)$.
\end{itemize}

\textbf{Case $a=b\neq0$:} beyond $(0,0)$, $(0,\pm1)$, $(\pm1,0)$, setting $x,y\neq0$ now gives
$a-ax^2-ay^2=0$, that is, $x^2+y^2=1$, so the entire unit circle is critical. On that circle,
\[\Hess f(x,y) = -4a\,e^{-1}\begin{pmatrix} x^2 & xy \\ xy & y^2 \end{pmatrix},\]
whose determinant is $16a^2e^{-2}(x^2y^2 - x^2y^2) = 0$. The Hessian is therefore degenerate
(\cref{def:sign_symmetric_matrix}) at every such point, and the test is inconclusive, as it must
be, since a whole curve of critical points cannot consist of isolated extrema.
\end{exercisesolution}

\begin{exercisesolution}[Solution to \cref{ex:ai_fta_hypotheses}]
% Generator: Claude Opus 5 (High)
\begin{enumerate}[label=\textbf{(\alph*)}]
  \item The restriction is what makes $\lvert z\rvert^k \leq \lvert z\rvert^{n-1}$ for every
    $k \leq n-1$, which is the step that collapses the whole sum into a single power. For
    $\lvert z\rvert < 1$ the inequality reverses, the low powers dominate, and the estimate is
    simply false: at $z = 0$ it would read $\lvert a_0\rvert \geq 0 \cdot (0 - C)$ only by
    accident of the right-hand side being negative, and it gives nothing usable.

    Nothing is lost, because the estimate is only ever applied for $\lvert z\rvert \geq R$, and
    $R \geq 1$ automatically: $C = \sum_{k=0}^n\lvert a_k\rvert$ includes the term
    $\lvert a_n\rvert = 1$, so $C \geq 1$ and hence $R \geq C+1 \geq 2$. The hypothesis is free.
  \item Here $f(z) = \lvert z\rvert^2+1 \geq 1$, so $f$ has no root, and indeed
    $\lvert f\rvert$ attains its minimum $1$ at $z_0 = 0$. Both halves of the first part of the
    proof therefore go through unchanged: $\lvert f\rvert$ grows without bound, the infimum is
    attained, and the minimiser is interior.

    The step that breaks is the expansion. The proof writes $g(z) := f(z_0+z) = \sum_k b_kz^k$,
    and $f$ here is not a polynomial in $z$: it is $z\overline z + 1$, and $\overline z$ is not
    expressible in powers of $z$. Concretely, at $z_0 = 0$ and $z = re^{i\varphi}$,
    \[g\big(re^{i\varphi}\big) = 1 + r^2,\]
    in which the perturbation carries \emph{no} factor $e^{i\ell\varphi}$. The direction
    $\varphi$ has vanished from the expression, so there is no direction to aim, and the term
    $r^2$ increases $\lvert g\rvert$ whichever way one leaves the origin.

    That is the sharpest way to say what the proof actually uses: not that $f$ is smooth, nor
    that it grows, but that the lowest-order correction to $b_0$ has a \emph{phase} that the free
    parameter $\varphi$ can rotate onto $-b_0$. The function $z\overline z$ is the standard
    example of a perturbation with no phase to rotate.
  \item The requirement is $e^{i(\ell\varphi+\psi)} = -1$, i.e.\
    $\ell\varphi + \psi \equiv \pi \pmod{2\pi}$, so
    \[\varphi_j := \frac{\pi - \psi + 2\pi j}{\ell}, \qquad j = 0,1,\ldots,\ell-1,\]
    which are $\ell$ distinct directions in $[0,2\pi)$, equally spaced by $2\pi/\ell$. The
    argument does not care which is used, because it needs only \emph{one} ray along which
    $\lvert g\rvert$ falls below $\lvert b_0\rvert$; the contradiction is reached from a single
    such ray. That $\ell$ of them exist rather than one is a bonus of working over $\mathbb{C}$,
    and it is exactly the freedom that \textbf{(b)} shows to be indispensable.
\end{enumerate}
\end{exercisesolution}

\begin{exercisesolution}[Proof of \cref{ex:hessian_parameterized_exp}]
% Generator: Gemini 3.6 Flash (High)
% Correction: Claude Opus 5 (High) --- Hf -> \Hess f, \text{Tr} -> \Tr, ^T -> \transp, and the
% \implies signs standing in for verbs replaced by prose. The range argument in (c) also spelled
% out: "by continuity" was doing a lot of unstated work, and it is the connectedness of R^2 plus
% the intermediate value theorem that closes the gaps.
\begin{enumerate}[label=\textbf{(\alph*)}]
  \item The gradient of $f_\alpha$ is
\[\nabla f_\alpha(x,y) = \begin{pmatrix} e^x(x^2 + \alpha y^2 + 2x) \\ 2\alpha y e^x \end{pmatrix},\]
and we look for the points where it vanishes. Since $e^x > 0$ everywhere, the second component
vanishes when $\alpha y = 0$, which splits the problem in two.
\begin{itemize}
  \item \textbf{If $\alpha = 0$:} the second component vanishes identically, and the first reduces
    to $x^2 + 2x = 0$, so $x = 0$ or $x = -2$. The critical set is the pair of vertical lines
    $\{(0,y) \mid y \in \mathbb{R}\}$ and $\{(-2,y) \mid y \in \mathbb{R}\}$.
  \item \textbf{If $\alpha \neq 0$:} the second component forces $y = 0$, and the first then reads
    $x^2 + 2x = 0$ again. The critical points are $(0,0)$ and $(-2,0)$.
\end{itemize}

  \item For $\alpha \neq 0$, differentiating once more gives
\[\Hess f_\alpha(x,y) = e^x \begin{pmatrix} x^2+4x+2+\alpha y^2 & 2\alpha y \\ 2\alpha y & 2\alpha \end{pmatrix}.\]
\begin{itemize}
  \item \textbf{At $(0,0)$:} $\Hess f_\alpha(0,0) = \left(\begin{smallmatrix} 2 & 0 \\ 0 & 2\alpha \end{smallmatrix}\right)$, with determinant $4\alpha$. For $\alpha > 0$ the determinant is positive and $\Tr = 2+2\alpha > 0$, so the matrix is positive definite and $(0,0)$ is a strict local minimum. For $\alpha < 0$ the determinant is negative, so $(0,0)$ is a saddle point.
  \item \textbf{At $(-2,0)$:} $\Hess f_\alpha(-2,0) = e^{-2}\left(\begin{smallmatrix} -2 & 0 \\ 0 & 2\alpha \end{smallmatrix}\right)$, with determinant $-4\alpha e^{-4}$. For $\alpha > 0$ this is negative, so $(-2,0)$ is a saddle point. For $\alpha < 0$ it is positive, and the entry $\partial_1^2 f_\alpha(-2,0) = -2e^{-2} < 0$ makes the matrix negative definite, so $(-2,0)$ is a strict local maximum.
\end{itemize}
Notice that the two critical points swap roles as $\alpha$ changes sign, and that neither test is
inconclusive: $\det \Hess$ vanishes only at $\alpha = 0$, the case excluded here.

  \item Each range is determined by exhibiting the extreme values and then filling in between by the
    intermediate value theorem, which applies because $\mathbb{R}^2$ is connected and $f_\alpha$ is
    continuous.
\begin{itemize}
  \item \textbf{If $\alpha > 0$:} both $x^2$ and $\alpha y^2$ are non-negative, so $f_\alpha \geq 0$,
    with $f_\alpha(0,0) = 0$ attained. Along the $x$-axis, $f_\alpha(x,0) = x^2 e^x \to \infty$ as
    $x \to \infty$, so values grow without bound. Hence $f_\alpha(\mathbb{R}^2) = [0,\infty)$.
  \item \textbf{If $\alpha < 0$:} along the $x$-axis $f_\alpha(x,0) = x^2 e^x$ sweeps out
    $[0,\infty)$ as above, while along the $y$-axis $f_\alpha(0,y) = \alpha y^2$ sweeps out
    $(-\infty,0]$. Together these give $f_\alpha(\mathbb{R}^2) = \mathbb{R}$.
  \item \textbf{If $\alpha = 0$:} $f_0(x,y) = x^2 e^x$ does not depend on $y$ and is non-negative,
    with value $0$ at $x = 0$ and unbounded above. Hence $f_0(\mathbb{R}^2) = [0,\infty)$.
\end{itemize}
\end{enumerate}
\begin{ainote}
It is worth pausing on the $\alpha > 0$ case: $(0,0)$ is a strict local minimum \emph{and} the global
minimum, but $(-2,0)$ is only a saddle, and $f_\alpha$ has no global maximum at all. Part
\textbf{(c)} is not a corollary of part \textbf{(b)} --- the Hessian test is purely local, and
reading a range off it is a standard way to go wrong.
\end{ainote}
\end{exercisesolution}




% ============================================================
% Chapter 13 --- Lagrange multipliers
% Originally: content/week-05.tex, \section{Lagrange multipliers} (Corsin Week 5),
%             with problems 5.2, 6.10 and 6.11 from sheets 5 and 6.
% ============================================================
\chapter{Lagrange Multipliers}
\label{ch:lagrange}

% Transition: Claude Opus 5 (Medium)
The previous chapter looked for extrema of $f$ on an open set, where the condition
$\nabla f(x_0) = 0$ is available. Restrict $f$ to a lower-dimensional constraint set and that
condition is generally too strong to hold anywhere; what replaces it is the requirement that
$\nabla f$ be a linear combination of the gradients of the constraints. This chapter develops
that criterion, the geometry behind it, and what goes wrong when the constraint gradients fail to
be independent.

% Originally: content/week-05.tex, \section{Lagrange multipliers} (Corsin Week 5)

% Source: Corsin Nick/Class Notes/Week 5.pdf, pp. 3--4
\section{Lagrange multipliers}
\label{sec:lagrange_multipliers}

Suppose we want to find the hottest place on the surface of the earth. In theory, temperature is
defined everywhere, giving some function
\[T \in C^\infty(\mathbb{R}^3, (0,\infty)) \qquad \text{(absolute zero is not attainable)}\]
If we model the earth as $S^2 := \{x \in \mathbb{R}^3 : |x| = 1\}$, then we want to find the local
maxima of $T|_{S^2}$, which \textbf{are not necessarily critical points of $T$}!

\begin{center}
\begin{tikzpicture}[>=Latex, scale=1.3]
  \draw[blue!80!black, thick] (0,0) circle (1.3);
  \draw[blue!80!black, dashed] (1.3,0) arc (0:180:1.3 and 0.4);
  \draw[blue!80!black] (-1.3,0) arc (180:360:1.3 and 0.4);
  \node[blue!80!black, font=\large] at (-0.6,0.5) {$S^2$};

  % Point q (radial gradient = extremum)
  \coordinate (Q) at (-0.55, 0.95);
  \fill[red!80!black] (Q) circle (1.8pt) node[left, font=\footnotesize] {$q$};
  \draw[red!80!black, thick, ->] (Q) -- ++(-0.4, 0.7) node[above right, font=\small] {$\nabla T(q)$};
  \node[red!80!black, font=\tiny, right] at (-0.1, 1.4) {local extremum of $T|_{S^2}$};

  % Point p (oblique gradient = not extremum)
  \coordinate (P) at (0.35, -0.45);
  \fill[orange!90!black] (P) circle (1.8pt) node[left, font=\footnotesize] {$p$};
  \draw[orange!90!black, thick, ->] (P) -- ++(0.75, -0.2) node[below right, font=\small] {$\nabla T(p)$};
  \draw[orange!90!black, dotted, thick] (-0.3, -0.3) -- (1.0, -0.6);
  \node[orange!90!black, font=\tiny, align=left] at (1.1, -0.1) {has tangent component $\implies$\\not a local extremum};
\end{tikzpicture}
\end{center}

% Generator: Claude Opus 5 (Medium) --- Corsin's sentence reads "T has a local extremum if it is
% perpendicular", which is both missing its subject and the wrong direction of implication.
Read the picture as follows. At $p$ the gradient $\nabla T(p)$ has a non-zero component
\emph{along} the surface, so walking on $S^2$ in that direction increases $T$: the point $p$
cannot be a local extremum of $T|_{S^2}$. So if $T|_{S^2}$ \emph{does} have a local extremum at a
point, then $\nabla T$ there must have no tangential component left --- it must be
\textbf{perpendicular to the surface}, as at $q$.

Note the direction of that statement carefully: perpendicularity is \textbf{necessary, not
sufficient}. Plenty of points have a perpendicular gradient without being extrema at all --- on
the earth, the coldest place has a perpendicular gradient just as the hottest does, and so does
many a saddle. What the condition delivers is a shortlist of \emph{candidates}; deciding which
of them are maxima, which minima, and which neither is a separate step, and in practice is done
by simply evaluating $T$ at each and comparing.

Note also that $S^2$ is the \newterm{level set} (\germanterm{Niveaumenge}) $g(x) = 0$ for
$g(x) := |x|^2 - 1$ (i.e., $S^2 = g^{-1}\{0\}$).

\begin{remark}[Gradient perpendicular to level sets]
For $g(x) := |x|^2-1$ we have $\nabla g(x) = 2x$, which at a point $x \in S^2$ is the outward
radial direction --- perpendicular to $S^2$. In general, the gradient is perpendicular to the
level sets.
\end{remark}

\begin{ainote}
Corsin writes $\nabla g(x) = \tfrac{x}{|x|}$ here. That is the gradient of $x \mapsto |x|$, not
of the $g(x) = |x|^2-1$ just fixed above, whose gradient is $2x$. The two differ only by the
positive factor $2|x|$, so they point the same way and the geometric conclusion is unaffected
--- but the constraint actually differentiated in the Lagrangian below is the squared one. See \cref{ex:xyz_on_ball} for the same $|x|$ versus $|x|^2$ slip in the exercise.
\end{ainote}

% Source: Corsin Nick/Class Notes/Week 5.pdf, pp. 4--5
This intuition is made rigorous in the following.

\begin{proposition}[Constrained optimization problems]
\label{prop:constrained_optimization}
Suppose $f : U \to \mathbb{R}$, $g : U \to \mathbb{R}^k$ are $C^1$ and $U \subseteq \mathbb{R}^n$
is open, and write $g := (g_1,\dots,g_k)$. If $f|_{g^{-1}\{0\}}$ has a local extremum at
$x_0 \in g^{-1}\{0\}$, then
\[\{\nabla f(x_0), \nabla g_1(x_0), \dots, \nabla g_k(x_0)\} \text{ are linearly dependent.}\]
\end{proposition}

\begin{remark}
Once $(\nabla g_1,\dots,\nabla g_k)$ are linearly independent everywhere on $g^{-1}\{0\}$ (the
standing assumption below), the constraint set $g^{-1}\{0\}$ is in fact a submanifold of
$\mathbb{R}^n$ --- this is made precise by \cref{thm:regular_value_theorem}, which also
explains \cref{prop:constrained_optimization} geometrically: $\nabla f(x_0)$ must be normal to the
tangent space of the constraint set at a constrained extremum.
\end{remark}

For $k = 1$, this corresponds to our intuition above of $\nabla f(x_0)$ \textbf{parallel to}
$\nabla g(x_0)$.

\textit{In practice}, we assume $(\nabla g_1, \dots, \nabla g_k)$ nowhere linearly dependent,
since otherwise $\lambda_0 = 0$ is possible (which does not give us anything).

% Generator: Claude Opus 5 (Medium)
\begin{aiexample}[What goes wrong when $\nabla g$ vanishes]
\label{ex:constraint_qualification_fails}
The standing assumption is not bookkeeping --- drop it and the method can miss an extremum
outright. Let
\[f(x,y) := x, \qquad M := \{(x,y) \in \mathbb{R}^2 : y^2 = x^3\},\]
the \newterm{cuspidal cubic} (\germanterm{Neilsche Parabel}). Since $y^2 \geq 0$ forces
$x^3 \geq 0$, every point of $M$ has $x \geq 0$, and $(0,0) \in M$. So $f$ attains a clear
\textbf{global minimum on $M$ at the origin}.

Now run the method with $g(x,y) := y^2-x^3$. The Lagrange condition $\nabla f = \lambda\nabla g$
reads
\[(1,0) = \lambda\,(-3x^2,\ 2y),\]
whose first component demands $1 = -3\lambda x^2$. At the origin this says $1 = 0$: no
multiplier $\lambda$ exists. The method returns \textbf{no candidates at all}, while the minimum
sits there in plain sight.

The culprit is visible in $\nabla g(0,0) = (0,0)$: the standing assumption fails at exactly the
point that matters. Geometrically, $M$ has a \emph{cusp} at the origin --- it is not a smooth
curve there, so it has no tangent line, and \qt{$\nabla f$ perpendicular to the constraint} has
nothing to be perpendicular to.

\textbf{The practical moral:} before trusting a Lagrange computation, check that
$\nabla g \neq 0$ (or, for several constraints, that the $\nabla g_j$ are independent) at every
point of the constraint set. Where it fails, add those points to your candidate list by hand.
\end{aiexample}

% Supplement: Adrien Martelli/Notes/20.03-Notes_Chap11_Lagrangemultipliers_geometric_explanation.pdf
\begin{remark}[Where the coefficient $\lambda_0$ went]
\Cref{prop:constrained_optimization} asserts linear \emph{dependence} of
$\{\nabla f(x_0), \nabla g_1(x_0), \dots, \nabla g_k(x_0)\}$, i.e.\ scalars
$\lambda_0, \lambda_1, \dots, \lambda_k$, not all zero, with
$\lambda_0 \nabla f(x_0) + \lambda_1 \nabla g_1(x_0) + \dots + \lambda_k \nabla g_k(x_0) = 0$.
The standing assumption above --- $(\nabla g_1,\dots,\nabla g_k)$ linearly independent --- is
exactly what rules out $\lambda_0 = 0$ (else the $\nabla g_j$'s alone would already be dependent),
so $\lambda_0 \neq 0$ and we may divide by it, silently renaming $\lambda_j/\lambda_0$ as
$\lambda_j$; this is why the Lagrangian below carries no separate $\lambda_0$. Adrien Martelli's
notes keep $\lambda_0$ explicit instead of dividing it away, pairing the dependence relation with
a normalization constraint $\lambda_0^2 + \dots + \lambda_k^2 = 1$ (any nonzero vector of
multipliers can be rescaled to satisfy this). The two formulations are equivalent whenever
$\lambda_0 \neq 0$; keeping $\lambda_0$ explicit only matters in degenerate settings where the
independence assumption on $(\nabla g_1,\dots,\nabla g_k)$ fails and $\lambda_0 = 0$ is genuinely
forced, e.g.\ if $x_0$ is itself a critical point of the constraint map $g$.
\end{remark}

% Generator: Claude Opus 5 (High)
\begin{remark}[What the multiplier itself measures]
\label{rem:lagrange_multiplier_sensitivity}
So far $\lambda$ has been bookkeeping --- a number produced by the method and then discarded once
the candidate points are found. It carries information of its own. Consider the constrained
problem with the constraint level left free,
\[m(c) := \min\{f(x) : g(x) = c\},\]
and suppose the minimiser $x^*(c)$ depends differentiably on $c$. Differentiating
$g(x^*(c)) = c$ gives $\langle \nabla g(x^*(c)), (x^*)'(c)\rangle = 1$, while the Lagrange
condition $\nabla f = \lambda\nabla g$ at $x^*(c)$ turns the chain rule for $m$ into
\[m'(c) = \big\langle \nabla f(x^*(c)),\ (x^*)'(c)\big\rangle
        = \lambda\big\langle \nabla g(x^*(c)),\ (x^*)'(c)\big\rangle = \lambda.\]
So $\lambda$ is the rate at which the optimal value responds to relaxing the constraint by one
unit --- in economics this reading is standard enough to have its own name, the
\qt{shadow price} of the constraint. A large $\lvert\lambda\rvert$ says the constraint is
expensive; $\lambda = 0$ says it is not binding at all.
\end{remark}

% Source: Corsin Nick/Class Notes/Week 5.pdf, pp. 4--5
% Supplement: Sascha Brack/Class Notes/Week_05_Notes_Friday.pdf, p. 5
\begin{center}
\begin{tikzpicture}[>=Latex, scale=1.3]
  \draw[very thick] (0,0) circle (1.2);
  \node[font=\large] at (-0.5,0.7) {$M$};
  \fill (0,0) circle (1.2pt) node[below right, font=\scriptsize] {$0$};

  \coordinate (P1) at ({1.2*cos(20)}, {1.2*sin(20)});
  \draw[magenta!90!black, thick, ->] (P1) -- ++({0.8*cos(20)}, {0.8*sin(20)}) node[above, font=\scriptsize] {$\nabla g(x_1, y_1)$};
  \fill (P1) circle (1.5pt) node[below left, font=\scriptsize, magenta!90!black] {$(x_1, y_1)$};

  \coordinate (P2) at ({1.2*cos(60)}, {1.2*sin(60)});
  \draw[magenta!90!black, thick, ->] (P2) -- ++({0.8*cos(60)}, {0.8*sin(60)}) node[right, font=\scriptsize] {$\nabla g(x_2, y_2)$};
  \fill (P2) circle (1.5pt) node[below left, font=\scriptsize, magenta!90!black] {$(x_2, y_2)$};

  \coordinate (P3) at ({1.2*cos(150)}, {1.2*sin(150)});
  \draw[magenta!90!black, thick, ->] (P3) -- ++({0.8*cos(150)}, {0.8*sin(150)}) node[above left, font=\scriptsize] {$\nabla g(x_3, y_3)$};
  \fill (P3) circle (1.5pt) node[below right, font=\scriptsize, magenta!90!black] {$(x_3, y_3)$};
\end{tikzpicture}
\end{center}

As the figure suggests, $\nabla g(x)$ is perpendicular to the constraint set $M := g^{-1}\{0\}$. To find extrema along $M$, we basically only care about the tangential part of $\nabla f$. We can then find a multiplier $\lambda$ to \qt{erase} the perpendicular component of $\nabla f(x_0)$. We do this by forming the \newterm{Lagrangian}
(\germanterm{Lagrange-Funktion}):
\[
\begin{aligned}
L : U \times \mathbb{R}^k &\to \mathbb{R} \\
(x,\lambda) &\mapsto f(x) - \lambda\cdot g(x) = f(x) - \lambda_1g_1(x) - \dots - \lambda_kg_k(x)
\end{aligned}
\]
We then find local extrema by forcing $\nabla L(x,\lambda) = 0$, i.e.
\[\underbrace{\frac{\partial L}{\partial x_i} = 0 \ \ \forall i = 1,\dots,n}_{\nabla f(x) - \lambda\cdot\nabla g(x) = 0} \qquad \text{and} \qquad \underbrace{\frac{\partial L}{\partial \lambda_j} = 0 \ \ \forall j = 1,\dots,k}_{g_j(x) = 0}\]

% Generator: Claude Opus 5 (Medium)
\begin{remark}[Why the Lagrangian is built this way]
\label{rem:why_the_lagrangian}
The construction can look like a trick pulled out of nowhere. It is not: $L$ is assembled
precisely so that setting \emph{one} gradient to zero encodes \emph{both} halves of the problem
at once.

Read the two underbraces. Differentiating in the $x_i$ reproduces the geometric condition ---
$\nabla f$ parallel to the $\nabla g_j$, i.e.\ no tangential component left. Differentiating in
$\lambda_j$ returns $g_j(x)=0$, the constraint itself, because $\lambda_j$ appears in $L$ only
in the term $-\lambda_jg_j(x)$, linearly. So the multipliers are not auxiliary clutter: they are
exactly the bookkeeping that lets a \emph{constrained} problem in $n$ variables be written as an
\emph{unconstrained} critical-point problem in $n+k$ variables, which is the one kind of problem
we already know how to solve (\cref{prop:extremum_implies_critical}).

That also explains the minus sign, which is arbitrary: writing $f + \lambda\cdot g$ merely flips
the sign of every $\lambda_j$ and changes nothing. And it explains why $\lambda$ is discarded at
the end --- it is scaffolding, not part of the answer.

\begin{ainote}
Though $\lambda$ is not meaningless. In applications it is the \newterm{shadow price}
(\germanterm{Schattenpreis}): if the constraint is relaxed from $g=0$ to $g=c$, the optimal value
changes at rate $\lambda$ to first order, $\tfrac{d}{dc}f_{\text{opt}}(c) = \lambda$. This is
where the name comes from in economics, and it is why $\lambda$ carries units. Not needed for
this course, but it is the usual answer to \qt{what \emph{is} $\lambda$?}
\end{ainote}
\end{remark}

% Source: Corsin Nick/Class Notes/Week 5.pdf, pp. 5--6
\begin{exercise}[Extrema of $xyz$ on the unit ball]
\label{ex:xyz_on_ball}
Find the local extrema of
\[f : K \to \mathbb{R}, \quad (x,y,z) \mapsto xyz, \qquad K = \{(x,y,z) \in \mathbb{R}^3 : x^2+y^2+z^2 \leq 1\}.\]
\end{exercise}

% Kept inline: solved directly in Corsin Nick/Class Notes/Week 5.pdf, pp. 5--6.
\begin{exercisesolution}[Proof of \cref{ex:xyz_on_ball}]
We first find the \textbf{interior critical points}:
\[\nabla f(x,y,z) = (yz,\ xz,\ xy) \overset{!}{=} 0 \implies x = y = 0 \ \text{ or } \ x = z = 0 \ \text{ or } \ y = z = 0.\]
Then we look for local extrema \textbf{on the boundary}. We need to find the local extrema of
$f|_{g^{-1}\{0\}}$, where $g : \mathbb{R}^3\to\mathbb{R}$ is given by $g(x) = |x|^2 - 1$.

The Lagrangian is
\[L(x,y,z,\lambda) := f(x,y,z) - \lambda g(x,y,z) = xyz - \lambda(x^2+y^2+z^2-1),\]
giving the system of equations
\[
\begin{aligned}
\frac{\partial L}{\partial x} &= yz - 2\lambda x = 0, & \frac{\partial L}{\partial \lambda} &= 1 - (x^2+y^2+z^2) = 0. \\
\frac{\partial L}{\partial y} &= xz - 2\lambda y = 0, \\
\frac{\partial L}{\partial z} &= xy - 2\lambda z = 0,
\end{aligned}
\]
We then have
\begin{equation*}
2\lambda = \frac{yz}{x} = \frac{xz}{y} = \frac{xy}{z} \tag{$*$}
\end{equation*}
\[\implies \underbrace{y^2z^2 = x^2z^2}_{y^2 = x^2} = \underbrace{x^2y^2}_{y^2 = z^2} \qquad \text{(if any coordinate is 0, so are the others)}\]
\[\implies x^2 = y^2 = z^2.\]
Then $x^2+y^2+z^2 = 1 \implies 3x^2 = 1 \implies x = \pm\frac{1}{\sqrt{3}}$, giving
\textbf{8 solutions}:
\[(x,y,z) \in \left\{\tfrac{1}{\sqrt{3}}(\alpha,\beta,\gamma) : \alpha,\beta,\gamma \in \{\pm 1\}\right\}\]
plus the solutions $\{(x_1,x_2,x_3) \in K : x_i = x_j = 0,\ i \neq j\}$.

\textbf{To justify division by $z$, $x$ or $y$ in $(*)$:} if $z = 0$, then
$-\lambda x = -\lambda y = xy = 0$, giving $x = 0$ or $y = 0$, and $\lambda = 0$, since
$x = y = z = 0$ does not satisfy $x^2+y^2+z^2 = 1$. Then, since $\lambda = 0$,
$\nabla f(x,y,z) = 0$ and we have already found that point in the first step. Similarly for
$y = 0$ or $x = 0$.

% Generator: Claude Opus 5 (Medium) --- Corsin's solution stops at the candidate list. The
% exercise asks for the local extrema, so the classification step is completed here.
\textbf{Classifying the candidates.} The work so far produced candidates, not answers
(\cref{prop:extremum_implies_critical} and the warning after it). Since $K$ is compact and $f$
is continuous, $f$ does attain a global maximum and minimum on $K$, so it is enough to evaluate.

At a boundary candidate $\tfrac{1}{\sqrt3}(\alpha,\beta,\gamma)$ with
$\alpha,\beta,\gamma \in \{\pm1\}$,
\[f\Big(\tfrac{1}{\sqrt3}(\alpha,\beta,\gamma)\Big)
  = \frac{\alpha\beta\gamma}{3\sqrt3},\]
so the value depends only on the sign $\alpha\beta\gamma$. Four of the eight points have
$\alpha\beta\gamma = +1$ and give the maximal value $\tfrac{1}{3\sqrt3}$; the other four have
$\alpha\beta\gamma = -1$ and give $-\tfrac{1}{3\sqrt3}$. On the coordinate axes found in the
first step, $f \equiv 0$.

Hence $f$ attains its \textbf{global maximum} $\tfrac{1}{3\sqrt3}$ at the four points with an
even number of minus signs, and its \textbf{global minimum} $-\tfrac{1}{3\sqrt3}$ at the four
with an odd number. The interior candidates are \textbf{not extrema at all}: at any point of a
coordinate axis $f=0$, yet every neighbourhood contains points where $xyz>0$ and points where
$xyz<0$ --- flip the sign of one coordinate. They are the three-variable analogue of a saddle.
\end{exercisesolution}

\begin{ainote}
Corsin writes $g(x) = |x| - 1$ for the constraint above, but the Lagrangian uses
$x^2+y^2+z^2-1$, i.e.\ $|x|^2-1$. Both cut out the same sphere; the squared version is the one
actually differentiated, so it is used throughout.
\end{ainote}

% Source: Jérôme Paschoud/Notizen/Woche 5.2 Optimierung I.pdf, p. 3
% Generator: Gemini 3.6 Flash (High)
\begin{exercise}[Extrema on a circle]
\label{ex:extrema_on_circle}
Find the maximum and minimum of $f(x,y) = xy$ under the condition that $(x,y)$ lies on the circle with radius $\sqrt{2}$. (cf.\ Michaels Bsp 7.2.3)
\end{exercise}

\begin{exercisesolution}[\cref{ex:extrema_on_circle}]
We want to optimize $f(x,y) = xy$ subject to the constraint $g(x,y) = x^2+y^2-2 = 0$.
The Lagrangian is
\[L(x,y,\lambda) := f(x,y) - \lambda g(x,y) = xy - \lambda(x^2+y^2-2).\]
Setting the gradient of $L$ to zero yields the system:
\begin{align*}
  \partial_x L &= y - 2\lambda x = 0 \\
  \partial_y L &= x - 2\lambda y = 0 \\
  \partial_\lambda L &= 2 - (x^2+y^2) = 0.
\end{align*}
From the first two equations, $y = 2\lambda x$ and $x = 2\lambda y$. Substituting the second into the first gives $y = 4\lambda^2 y$, or $y(1-4\lambda^2) = 0$.
If $y=0$, then $x=0$, which contradicts $x^2+y^2=2$. Hence $y \neq 0$ and $4\lambda^2 = 1 \implies \lambda = \pm 1/2$.
If $\lambda = 1/2$, then $y = x$. The constraint gives $2x^2 = 2 \implies x = \pm 1$, giving points $(1,1)$ and $(-1,-1)$. At these points, $f(\pm 1, \pm 1) = 1$.
If $\lambda = -1/2$, then $y = -x$. The constraint gives $2x^2 = 2 \implies x = \pm 1$, giving points $(1,-1)$ and $(-1,1)$. At these points, $f(\pm 1, \mp 1) = -1$.
Since the circle is compact, these must be the global extrema. The maximum is $1$ and the minimum is $-1$.
\end{exercisesolution}

% Source: Jérôme Paschoud/Notizen/Woche 5.2 Optimierung I.pdf, p. 3
% Generator: Gemini 3.6 Flash (High)
\begin{exercise}[Largest rectangle in an ellipse]
\label{ex:largest_rectangle_ellipse}
Given an ellipse $E = \{ (x,y) \in \mathbb{R}^2 \mid x^2/A^2 + y^2/B^2 = 1 \}$, find the rectangle (with sides parallel to the axes) whose vertices lie on $E$ and has the largest area. (cf.\ Michaels Bsp 7.2.7)
\end{exercise}

\begin{exercisesolution}[\cref{ex:largest_rectangle_ellipse}]
Let the vertices of the rectangle in the four quadrants be $(\pm x, \pm y)$ with $x, y \geq 0$. The area of the rectangle is $f(x,y) = 4xy$.
We maximize $f(x,y)$ subject to the constraint $g(x,y) = \frac{x^2}{A^2} + \frac{y^2}{B^2} - 1 = 0$ in the region $x > 0, y > 0$.
The Lagrangian is
\[L(x,y,\lambda) := 4xy - \lambda\left(\frac{x^2}{A^2} + \frac{y^2}{B^2} - 1\right).\]
The critical point equations are:
\begin{align*}
  \partial_x L &= 4y - \frac{2\lambda x}{A^2} = 0 \implies \lambda = \frac{2yA^2}{x} \\
  \partial_y L &= 4x - \frac{2\lambda y}{B^2} = 0 \implies \lambda = \frac{2xB^2}{y} \\
  \partial_\lambda L &= 1 - \frac{x^2}{A^2} - \frac{y^2}{B^2} = 0.
\end{align*}
Equating the expressions for $\lambda$ yields $\frac{2yA^2}{x} = \frac{2xB^2}{y} \implies y^2 A^2 = x^2 B^2$, or $\frac{x^2}{A^2} = \frac{y^2}{B^2}$.
Substituting this into the constraint gives $2\frac{x^2}{A^2} = 1 \implies x = \frac{A}{\sqrt{2}}$.
Similarly, $y = \frac{B}{\sqrt{2}}$.
The maximal area is $f\left(\frac{A}{\sqrt{2}}, \frac{B}{\sqrt{2}}\right) = 4 \cdot \frac{A}{\sqrt{2}} \cdot \frac{B}{\sqrt{2}} = 2AB$.
\end{exercisesolution}

% Source: Corsin Nick/Class Notes/Week 5.pdf, p. 6
\begin{center}
\begin{tikzpicture}[>=Latex, scale=1.4]
  % Sphere
  \draw[blue!80!black, thick] (0,0) circle (1.4);
  \draw[blue!80!black, dashed] (1.4,0) arc (0:180:1.4 and 0.45);
  \draw[blue!80!black] (-1.4,0) arc (180:360:1.4 and 0.45);

  % Axes
  \draw[green!60!black, thick, ->] (-1.6,0) -- (1.6,0) node[right, font=\footnotesize] {$x$};
  \draw[green!60!black, thick, ->] (0,-1.6) -- (0,1.6) node[above, font=\footnotesize] {$z$};
  \draw[green!60!black, thick, ->] (0.8,-1.0) -- (-0.8,1.0) node[above left, font=\footnotesize] {$y$};

  % Inscribed cube vertices (+-1/sqrt(3), +-1/sqrt(3), +-1/sqrt(3))
  \coordinate (A1) at (-0.6, -0.5);
  \coordinate (A2) at (0.6, -0.5);
  \coordinate (A3) at (0.8, 0.2);
  \coordinate (A4) at (-0.4, 0.2);

  \coordinate (B1) at (-0.6, 0.7);
  \coordinate (B2) at (0.6, 0.7);
  \coordinate (B3) at (0.8, 1.4);
  \coordinate (B4) at (-0.4, 1.4);

  % Cube edges
  \draw[red!80!black, thick, dotted] (A1)--(A2)--(A3)--(A4)--cycle;
  \draw[red!80!black, thick, dotted] (B1)--(B2)--(B3)--(B4)--cycle;
  \draw[red!80!black, thick, dotted] (A1)--(B1) (A2)--(B2) (A3)--(B3) (A4)--(B4);

  % Vertex dots
  \foreach \p in {A1,A2,A3,A4,B1,B2,B3,B4} {
    \fill[red!80!black] (\p) circle (1.8pt);
  }
  \node[below, font=\small] at (0,-1.7) {8 critical points $\tfrac{1}{\sqrt{3}}(\pm 1, \pm 1, \pm 1)$ on $S^2$};
\end{tikzpicture}
\end{center}

% Generator: Claude Opus 5 (High) --- not in Corsin's notes; included because it is the one
% place in the course where an analytic tool repays a debt to linear algebra, and because the
% Rayleigh-quotient argument is a genuine use of \cref{prop:constrained_optimization} with more
% than one constraint.
\begin{theorem}[Spectral theorem, via constrained optimization]
\label{thm:spectral_theorem_optimization}
Every symmetric $A \in \mathbb{R}^{n\times n}$ admits an orthonormal basis
$\{v_1,\dots,v_n\}$ of $\mathbb{R}^n$ consisting of eigenvectors of $A$, with real eigenvalues.
\end{theorem}

\begin{proof}
Let $f(x) := \langle x, Ax\rangle$, a polynomial and hence continuous, and let
$g(x) := \lvert x\rvert^2 - 1$, so that the unit sphere is $S^{n-1} = g^{-1}\{0\}$. Since $A$ is
symmetric, $\nabla f(x) = 2Ax$; also $\nabla g(x) = 2x$.

We construct the $v_j$ one at a time. Suppose $v_1,\dots,v_k$ are already orthonormal
eigenvectors of $A$, with eigenvalues $\lambda_1,\dots,\lambda_k$ (for $k=0$ there is nothing to
suppose). Put
\[K_k := S^{n-1} \cap \{x : \langle x,v_j\rangle = 0 \text{ for } j=1,\dots,k\},\]
the unit sphere of the orthogonal complement of $\operatorname{span}\{v_1,\dots,v_k\}$. This is
closed and bounded, hence compact by \cref{thm:heine_borel}, and non-empty as long as $k<n$, so
$f$ attains a maximum on it at some $v_{k+1} \in K_k$.

The constraints active at $v_{k+1}$ are $g$ together with the $k$ linear functions
$h_j(x) := \langle x,v_j\rangle$, whose gradients are $\nabla g(v_{k+1}) = 2v_{k+1}$ and
$\nabla h_j(v_{k+1}) = v_j$. These $k+1$ vectors are orthogonal and non-zero, hence linearly
independent, so \cref{prop:constrained_optimization} applies and yields multipliers with
\[2Av_{k+1} = 2\lambda_{k+1}v_{k+1} + \sum_{j=1}^{k}\mu_jv_j.\]
Pair both sides with $v_i$ for a fixed $i \leq k$. On the right, $\langle v_{k+1},v_i\rangle = 0$
and $\langle v_j,v_i\rangle = \delta_{ji}$, leaving $\mu_i$. On the left, symmetry of $A$ moves it
across the inner product,
\[\langle Av_{k+1}, v_i\rangle = \langle v_{k+1}, Av_i\rangle
  = \lambda_i\langle v_{k+1},v_i\rangle = 0,\]
using that $v_i$ is already an eigenvector. Hence $\mu_i = 0$ for every $i \leq k$, and the
displayed identity collapses to $Av_{k+1} = \lambda_{k+1}v_{k+1}$. So $v_{k+1}$ is a unit
eigenvector orthogonal to all its predecessors, and after $n$ steps the $v_j$ form the required
basis. Each eigenvalue is real because $\lambda_j = \langle v_j, Av_j\rangle$ is a real number by
construction.
\end{proof}

% Generator: Claude Opus 5 (High)
\begin{corollary}[Extreme eigenvalues as Rayleigh quotients]
\label{cor:courant_fischer_extreme}
For symmetric $A \in \mathbb{R}^{n\times n}$,
\[\lambda_{\min}(A) = \min_{x \in S^{n-1}}\langle x,Ax\rangle, \qquad
  \lambda_{\max}(A) = \max_{x \in S^{n-1}}\langle x,Ax\rangle.\]
\end{corollary}

\begin{proof}
Write $x \in S^{n-1}$ in the orthonormal eigenbasis of \cref{thm:spectral_theorem_optimization}
as $x = \sum_j c_jv_j$ with $\sum_j c_j^2 = 1$. Then
$\langle x,Ax\rangle = \sum_j \lambda_jc_j^2$ is a convex combination of the eigenvalues, so it
lies between $\lambda_{\min}$ and $\lambda_{\max}$; both bounds are attained by taking $x$ to be
the corresponding eigenvector.
\end{proof}

% Originally: content/week-05.tex, end of the chapter

\begin{aiexercise}[Extrema on a Circle via Lagrange Multipliers]
\label{ex:ai_lagrange_circle}
% Generator: Gemini 3.6 Flash (Medium)
Find the maximum and minimum values of $f(x,y) := x + 2y$ subject to the constraint $g(x,y) := x^2 + y^2 - 5 = 0$.
\end{aiexercise}


% --- exercises relocated here from the dissolved sheets ---

% Originally: content/week-05.tex, \exercisesheet{5}, problem 5.2
% Source: exercises/Ex5_Analysis2_eng.pdf

\begin{exercise}[5.2 --- Lagrange multipliers \textnormal{(important)}]
\label{ex:5.2}
Consider the ball $U := \{(x,y,z) : x^2+y^2+z^2 < 1\}$ and the function
\[f(x,y,z) := 1 - x^2 - y^2 + 4x.\]
\begin{enumerate}[label=\textbf{\arabic*.}]
  \item Motivate rigorously why $f$ attains its absolute maximum and minimum in $\overline{U}$.
  \item Find all the critical points of $f$ which lie in the interior of $U$. Say whether
    $\sup_U f$ (or $\inf_U f$) is attained at some point in $U$.
  \item Say whether $\max_{\partial U} f$ and $\min_{\partial U}$ exist.
  \item With the method of Lagrange multipliers, find the possible critical points for the
    constrained problem $\max\{f(x,y,z) : (x,y,z) \in \partial U\}$, and same for min.
  \item Among all the points you have found, clarify at which points the maximum/minimum of $f$
    is attained.
\end{enumerate}
\end{exercise}

% Originally: content/week-06.tex, \exercisesheet{6}, problem 6.10
% Source: exercises/Ex6_Analysis2_eng.pdf

\begin{exercise}[6.10 --- Lagrange multipliers \textnormal{(important)}]
\label{ex:6.10}
Consider the function $f(x,y,z) = 3x-y+2z$ and the set
\[M = \{(x,y,z) \in \mathbb{R}^3 \mid x^2+y^2+z^2 = 1,\ x+y = 0\}.\]
Determine the extrema of $f$ on $M$ and their nature.
\end{exercise}

% Originally: content/week-06.tex, \exercisesheet{6}, problem 6.11
% Source: exercises/Ex6_Analysis2_eng.pdf

\begin{exercise}[6.11 --- Closest point on a hyperbolic paraboloid \textnormal{(important)}]
\label{ex:6.11}
Find the points on the submanifold
\[M = \{(x,y,z) \in \mathbb{R}^3 : z = x^2-y^2\}\]
that are closest to the point $(0,0,1)$.
\end{exercise}

% Originally: the \section{Solutions} of content/week-05.tex and content/week-06.tex

\section{Solutions}
\label{sec:lagrange_solutions}

\begin{exercisesolution}[Solution to \cref{ex:5.2}]
% Generator: Claude Sonnet 5 (Medium)
\begin{enumerate}[label=\textbf{\arabic*.}]
  \item $\overline{U}$ is closed and bounded in $\mathbb{R}^3$, hence compact by
    \Cref{thm:heine_borel}, and $f$ is a polynomial, hence continuous. By the extreme value
    theorem (\cref{item:extreme_value_theorem}), $f$ attains its absolute maximum and minimum on
    $\overline{U}$.
  \item $\nabla f(x,y,z) = (-2x+4,\ -2y,\ 0)$ vanishes only at $x=2$, $y=0$ (any $z$), and no such
    point lies in $U$ (since $x=2$ already violates $x^2+y^2+z^2<1$). So $f$ has \textbf{no
    critical point in the interior of $U$}; consequently, by
    \Cref{prop:extremum_implies_critical}, $\sup_U f$ and $\inf_U f$ are \textbf{not attained} at
    any interior point.
  \item $\partial U$ is compact (a sphere) and $f$ is continuous, so $\max_{\partial U} f$ and
    $\min_{\partial U} f$ \textbf{both exist}, again by the extreme value theorem
    (\cref{item:extreme_value_theorem}).
  \item With $g(x,y,z) := x^2+y^2+z^2-1$, the Lagrange condition $\nabla f = \lambda \nabla g$
    reads $(-2x+4,-2y,0) = \lambda(2x,2y,2z)$. The last coordinate gives $\lambda z = 0$.
    \textbf{If $z \neq 0$:} then $\lambda = 0$, so $-2x+4=0 \Rightarrow x=2$, impossible on
    $\partial U$. So $z = 0$. Then $x^2+y^2=1$, and the remaining equations are
    $4-2x = 2\lambda x$ and $-2y = 2\lambda y$, i.e.\ $-2y(1+\lambda) = 0$.
    \textbf{If $\lambda = -1$:} the first equation becomes $4-2x=-2x$, i.e.\ $4=0$, impossible.
    So $y = 0$, and $x^2=1$ gives $x = \pm 1$. The Lagrange candidates are
    $(1,0,0)$ (with $\lambda = 1$) and $(-1,0,0)$ (with $\lambda = -3$).
  \item Evaluating, $f(1,0,0) = 1-1+4 = 4$ and $f(-1,0,0) = 1-1-4 = -4$. Since there are no
    interior critical points, these are the global extrema on $\overline{U}$: the
    \textbf{maximum} $4$ is attained at $(1,0,0)$, and the \textbf{minimum} $-4$ at $(-1,0,0)$.
\end{enumerate}
\end{exercisesolution}

\begin{exercisesolution}[Proof of \cref{ex:ai_lagrange_circle}]
% Generator: Gemini 3.6 Flash (Medium)
The constraint set $S := \{(x,y) \in \mathbb{R}^2 \mid x^2 + y^2 = 5\}$ is compact and $\nabla g(x,y) = (2x, 2y)\transp \neq (0,0)\transp$ on $S$. By Lagrange Multipliers, candidates satisfy $\nabla f = \lambda \nabla g$, i.e.:
\[
1 = 2\lambda x \quad \text{and} \quad 2 = 2\lambda y \implies y = 2x.
\]
Substituting $y = 2x$ into $x^2 + y^2 = 5$ gives $5x^2 = 5 \implies x = \pm 1$. The candidate points are $(1,2)$ and $(-1,-2)$. Evaluating $f$:
\[
f(1,2) = 1 + 4 = 5 \quad (\text{maximum}), \qquad f(-1,-2) = -1 - 4 = -5 \quad (\text{minimum}).
\]
\end{exercisesolution}

\begin{exercisesolution}[Solution to \cref{ex:6.10}]
% Generator: Claude Sonnet 5 (Medium)
With $g_1(x,y,z) := x^2+y^2+z^2-1$ and $g_2(x,y,z) := x+y$, the Lagrangian
$L := f - \lambda_1 g_1 - \lambda_2 g_2$ gives the system
\[3 - 2\lambda_1 x - \lambda_2 = 0, \qquad -1-2\lambda_1 y - \lambda_2 = 0, \qquad 2 - 2\lambda_1 z = 0, \qquad x+y=0, \qquad x^2+y^2+z^2=1.\]
From $x+y=0$, $y=-x$. Subtracting the second equation from the first eliminates $\lambda_2$:
\[4 - 2\lambda_1 x - 2\lambda_1(-x) \cdot(-1)= 4-2\lambda_1x-2\lambda_1x=4-4\lambda_1 x = 0 \implies \lambda_1 x = 1.\]
From $2-2\lambda_1 z=0$, $\lambda_1 = 1/z$ (so $z \neq 0$), hence $x = z$ (from $\lambda_1 x = 1 =
\lambda_1 z$). With $y=-x=-z$, the constraint becomes $3x^2 = 1$, i.e.\ $x = \pm\tfrac{1}{\sqrt3}$,
giving the two points $\tfrac{1}{\sqrt3}(1,-1,1)$ and $-\tfrac{1}{\sqrt3}(1,-1,1)$. Evaluating,
\[f\Big(\tfrac{1}{\sqrt3}(1,-1,1)\Big) = \frac{3+1+2}{\sqrt3} = 2\sqrt3, \qquad f\Big({-\tfrac{1}{\sqrt3}}(1,-1,1)\Big) = -2\sqrt3.\]
Since $M$ is compact (closed and bounded in $\mathbb{R}^3$) and $f$ is continuous, these are the
only Lagrange candidates, so $f$ attains its \textbf{maximum} $2\sqrt3$ at $\tfrac{1}{\sqrt3}(1,-1,1)$
and its \textbf{minimum} $-2\sqrt3$ at $-\tfrac{1}{\sqrt3}(1,-1,1)$.
\end{exercisesolution}

\begin{exercisesolution}[Solution to \cref{ex:6.11}]
% Generator: Claude Sonnet 5 (Medium)
Minimizing the distance to $(0,0,1)$ is the same as minimizing
$f(x,y,z) := x^2+y^2+(z-1)^2$ subject to $g(x,y,z) := z-x^2+y^2 = 0$. Since $f$ is a distance
squared to a fixed point, its minimum over the (closed) constraint set is attained. Lagrange's
condition $\nabla f = \lambda \nabla g$ reads
\[2x = -2\lambda x, \qquad 2y = 2\lambda y, \qquad 2(z-1) = \lambda,\]
i.e.\ $x(1+\lambda) = 0$ and $y(1-\lambda) = 0$.

\textbf{Case $x=y=0$:} the constraint gives $z=0$, so $(0,0,0)$ with $d^2 = 1$.

\textbf{Case $x \neq 0$ (so $\lambda=-1$):} then $2(z-1)=-1 \implies z = \tfrac12$, and $y(1-\lambda)=2y=0
\implies y=0$. The constraint $z=x^2-y^2$ gives $x^2=\tfrac12$, i.e.\ $x=\pm\tfrac{1}{\sqrt2}$,
giving $\big(\pm\tfrac{1}{\sqrt2},0,\tfrac12\big)$ with $d^2 = \tfrac12+0+\tfrac14 = \tfrac34$.

\textbf{Case $y \neq 0$ (so $\lambda=1$):} then $2(z-1)=1 \implies z=\tfrac32$, and $x(1+\lambda)=2x=0
\implies x=0$; the constraint becomes $-y^2 = \tfrac32$, which has no real solution.

Comparing $d^2=1$ against $d^2=\tfrac34$, the \textbf{closest points} are
$\big(\pm\tfrac{1}{\sqrt2},0,\tfrac12\big)$, at minimal distance $\tfrac{\sqrt3}{2}$.
\end{exercisesolution}


% ============================================================
% Chapter 14 --- Convexity
% Originally: content/week-06.tex, \section{Convexity} (Corsin Week 6),
%             with problem 6.7 from sheet 6.
% ============================================================
\chapter{Convexity}
\label{ch:convexity}

% Transition: Claude Opus 5 (Medium)
Convexity is the hypothesis under which the local information of the previous two chapters
becomes global: for a convex function a critical point is not merely a local minimum but the
minimum, and the Hessian test stops being a test at all. This chapter gives the definition, the
equivalent characterizations in terms of the gradient and the Hessian, and proves that they agree.

% Originally: content/week-06.tex, \section{Convexity} (Corsin Week 6)

% Source: Corsin Nick/Class Notes/Week 6.pdf, p. 2
\section{Convexity}
\label{sec:convexity}

% Sonnet 5 (Medium)
\subsection{Definitions and characterizations}

\begin{definition}[Convex set and convex function]
\label{def:convex_set_function}
A subset $A \subseteq \mathbb{R}^n$ is \newterm{convex} (\germanterm{konvex}) if for all
$x,y \in A$ and every $t \in [0,1]$, the point $p := (1-t)x+ty$ is an element of $A$.

A function $f : A \to \mathbb{R}$ is convex if $A$ is convex and, for all $x,y \in A$ and
$t \in [0,1]$,
\[f\big((1-t)x+ty\big) \leq (1-t)f(x)+tf(y).\]
\end{definition}

% Sonnet 5 (Medium)
\begin{center}
\begin{tikzpicture}[>=Latex, scale=1.0]
  % Curve is 0.18(r-2.6)^2 + 0.4; the chord endpoints must SIT on it:
  %   f(0.7) = 1.0498,  f(4.6) = 1.12.
  % Interior point at t = 0.35: r = 2.065, chord height 1.0744, curve height 0.4515.
  \draw[->] (-0.3,0) -- (5.6,0);
  \draw[->] (0,-0.3) -- (0,2.0);
  \draw[blue!70!black, thick, domain=0.2:5, samples=60, smooth]
    plot ({\x}, {0.18*(\x-2.6)*(\x-2.6)+0.4});
  \node[blue!70!black, font=\small] at (5.3,1.55) {$f$};

  \draw[orange!90!black, thick] (0.7,1.0498) -- (4.6,1.12);
  \node[orange!90!black, font=\scriptsize] at (3.5,1.32) {chord};
  \fill (0.7,1.0498) circle (1.5pt) node[left, font=\small] {$f(x)$};
  \fill (4.6,1.12) circle (1.5pt) node[below right, font=\small] {$f(y)$};

  % The convexity inequality, read off vertically at the interior point.
  \draw[gray, dashed] (2.065,0) -- (2.065,1.0744);
  \fill[orange!90!black] (2.065,1.0744) circle (1.5pt);
  \fill[blue!70!black] (2.065,0.4515) circle (1.5pt);
  \draw[red!80!black, <->, thick] (2.065,0.4515) -- (2.065,1.0744);

  \draw (0.7,0.05) -- (0.7,-0.05) node[below, font=\small] {$x$};
  \draw (4.6,0.05) -- (4.6,-0.05) node[below, font=\small] {$y$};
  \draw (2.065,0.05) -- (2.065,-0.05);
  \node[font=\scriptsize, below] at (2.065,-0.12) {$(1-t)x+ty$};
  \node[red!80!black, font=\small] at (2.9,-0.85)
    {$f\big((1-t)x+ty\big) \leq (1-t)f(x)+tf(y)$};
\end{tikzpicture}
\end{center}

\begin{definition}[Semidefinite matrices]
\label{def:semidefinite_matrix}
Let $A \in \mathbb{R}^{n\times n}$ be symmetric. Then $A$ is
\begin{itemize}
  \item \newterm{positive semidefinite} (\germanterm{positiv semidefinit}), also called
    \textbf{nonnegative definite}, if $v\transp Av \geq 0$ for all $v \in \mathbb{R}^n$;
  \item \newterm{negative semidefinite} (\germanterm{negativ semidefinit}), or
    \textbf{nonpositive definite}, if $v\transp Av \leq 0$ for all $v \in \mathbb{R}^n$.
\end{itemize}
These are the boundary cases of \cref{def:sign_symmetric_matrix}: positive definite implies
positive semidefinite, and the two differ exactly on the degenerate matrices, where
$v\transp Av = 0$ is allowed for some $v \neq 0$.
\end{definition}

\begin{ainote}
\Cref{def:semidefinite_matrix} is not in Corsin's notes. The terms are used freely from
\Cref{prop:convexity_characterizations} and \cref{ex:6.8} onward but were never named when
\Cref{def:sign_symmetric_matrix} is introduced later, so they are defined here instead of
being left implicit.
\end{ainote}

\begin{proposition}[Characterizations of convexity]
\label{prop:convexity_characterizations}
Let $U \subseteq \mathbb{R}^n$ be open and convex, and let $f \in C^2(U)$. Then the following are
equivalent:
\begin{enumerate}[label=\textbf{(\alph*)}]
  \item $f$ is convex;
  \item $\Hess f(x)$ is positive semidefinite for all $x \in U$;
  \item $f(y)-f(x) \geq Df_x(y-x)$ for all $x,y \in U$.
\end{enumerate}
\end{proposition}

\begin{example}[The parabola]
The parabola $p : \mathbb{R}^n \to \mathbb{R}$, $x \mapsto |x|^2 = \sum_{i=1}^n x_i^2$, is convex:
\[\partial_i\partial_j p(x) = 2\delta_{ij} \implies \Hess p(x) = 2\id_n,\]
and every eigenvalue is $2>0$, so $\Hess p(x)$ is positive definite (hence positive semidefinite)
for every $x$, and $p$ is convex by \cref{prop:convexity_characterizations}.
\end{example}

% Sonnet 5 (Medium)
\subsection{Proof of the characterizations}

% Source: Corsin Nick/Class Notes/Week 6.pdf, pp. 3--6
\begin{exercise}[Equivalent characterizations of convexity]
\label{ex:convexity_characterizations}
Let $U \subseteq \mathbb{R}^n$ be open and convex.
\begin{enumerate}[label=\textbf{(\alph*)}]
  \item Show that $f : U \to \mathbb{R}$ is convex if and only if the functions
    \[f_{x,y} : [0,1] \to \mathbb{R}, \qquad s \mapsto f(x+s(y-x))\]
    are convex for all $x,y \in U$.
  \item Show that
    \[f \text{ convex} \iff f(y)-f(x) \geq Df_x(y-x) \quad \text{for all } x,y \in U.\]
    This also holds for $f \in C^1(U)$, $U \subseteq \mathbb{R}^n$ open and convex.
  \item Conclude that if $f \in C^1(U)$ is convex and has a critical point at $x_0 \in U$, then $f$
    has a minimum at $x_0$.
\end{enumerate}
\end{exercise}

% Kept inline: solved directly in Corsin Nick/Class Notes/Week 6.pdf, pp. 3--6.
\begin{exercisesolution}[Solution to \cref{ex:convexity_characterizations}]
\begin{enumerate}[label=\textbf{(\alph*)}]
  \item By definition, $f_{x,y}$ is convex if and only if
    \[f_{x,y}\big((1-t)s_0+ts_1\big) \leq (1-t)f_{x,y}(s_0)+tf_{x,y}(s_1) \quad \text{for all } s_0,s_1,t \in [0,1],\]
    which, unwinding the definition of $f_{x,y}$, reads
    \[f\big(x+[(1-t)s_0+ts_1](y-x)\big) \leq (1-t)f(x+s_0(y-x))+tf(x+s_1(y-x)),\]
    that is,
    \[f\big((1-t)(x+s_0(y-x))+t(x+s_1(y-x))\big) \leq (1-t)f(x+s_0(y-x))+tf(x+s_1(y-x)).\]
    Setting $s_0=0$, $s_1=1$ recovers exactly the convexity inequality for $f$ at the points $x$
    and $y$, so if this holds for all $x,y \in U$, then $f$ is convex. Conversely, if $f$ is
    convex, then convexity of $f_{x,y}$ for each fixed $x,y \in U$ follows by applying convexity
    of $f$ to the two points $x+s_0(y-x)$ and $x+s_1(y-x)$, which both lie in $U$ since $U$ is
    convex.
  \item Recall that $g \in C^1(I)$, for $I \subseteq \mathbb{R}$ an interval, is convex if and
    only if $g'$ is increasing. Combined with \textbf{(a)}, this gives
    \[f \text{ convex} \iff f'_{x,y} \text{ is increasing, for all } x,y \in U.\]
    By the chain rule, for $t \in [0,1]$,
    \[f'_{x,y}(t) = Df_{x+t(y-x)}(y-x),\]
    and by the mean value theorem there exists $s \in (0,1)$ such that
    \[f(y)-f(x) = f_{x,y}(1)-f_{x,y}(0) = f'_{x,y}(s) = Df_{x+s(y-x)}(y-x).\]

    \textbf{($\Rightarrow$)} If $f$ is convex, then $f_{x,y}$ is convex for all $x,y \in U$ by
    \textbf{(a)}, so $f'_{x,y}$ is increasing, and in particular $f'_{x,y}(s) \geq f'_{x,y}(0)$.
    Combined with the mean value theorem identity above, this yields
    \[f(y)-f(x) \geq Df_x(y-x).\]

    \textbf{($\Leftarrow$)} Suppose $f(y)-f(x) \geq Df_x(y-x)$ for all $x,y \in U$. We must show
    that $0 \leq s < t \leq 1 \implies f'_{x,y}(s) \leq f'_{x,y}(t)$. Set $u := x+s(y-x)$ and
    $v := x+t(y-x)$, so that $u-v = (s-t)(y-x)$. Applying the hypothesis to the pairs $(u,v)$ and
    $(v,u)$ gives
    \[Df_u(v-u) \leq f(v)-f(u), \qquad Df_v(u-v) \leq f(u)-f(v).\]
    Therefore,
    \[(t-s)f'_{x,y}(s) = Df_u(v-u) \leq f(v)-f(u) \leq -Df_v(u-v) = Df_v(v-u) = (t-s)f'_{x,y}(t),\]
    and since $t-s>0$, this gives $f'_{x,y}(s) \leq f'_{x,y}(t)$, i.e., $f'_{x,y}$ is increasing,
    hence $f_{x,y}$ is convex for all $x,y \in U$, and so $f$ is convex by \textbf{(a)}.

    \begin{remark}
    The argument above only ever used differentiability of $f$, never twice-differentiability, so
    the equivalence in \textbf{(b)} already holds for $f \in C^1(U)$ with $U \subseteq
    \mathbb{R}^n$ open and convex --- one order of regularity weaker than
    \Cref{prop:convexity_characterizations}, which is stated for $f \in C^2(U)$. This is exactly
    the fact used in \cref{ex:6.7}\textbf{(b)}.
    \end{remark}
  \item Since $f$ is convex and $x_0 \in U$ is a critical point, $Df_{x_0} = 0$, so \textbf{(b)}
    gives
    \[0 = Df_{x_0}(y-x_0) \leq f(y)-f(x_0) \quad \text{for all } y \in U,\]
    i.e., $f(y) \geq f(x_0)$ for all $y \in U$. Hence $f$ has a minimum at $x_0$ --- in fact a
    global minimum on $U$, sharpening \cref{prop:extremum_implies_critical} (which only guarantees
    that critical points are the sole \emph{candidates} for extrema) in the convex case.
\end{enumerate}
\end{exercisesolution}

% Sonnet 5 (Medium) --- FIG-W06-01, drawn from the page image seen while transcribing
\begin{center}
\begin{tikzpicture}[>=Latex, scale=1.15]
  \draw[->] (-0.3,0) -- (5.2,0) node[right, font=\small] {$r$};
  \draw[->] (0,-0.3) -- (0,3.0);

  \draw[green!60!black, thick, domain=0.2:5.0, samples=60, smooth]
    plot ({\x}, {0.20*(\x-2.6)*(\x-2.6) + 1.0});
  \node[green!60!black, font=\small] at (5.35,2.2) {$f_{x,y}(r)$};

  % True tangents to f_{x,y}(r) = 0.20(r-2.6)^2 + 1.0:
  %   at s = 1.2: value 1.392, slope -0.56;   at t = 3.6: value 1.200, slope +0.40.
  % They must touch the curve and lie below it -- that is what part (b) says.
  \draw[orange!90!black, thick, domain=0.3:3.0]
    plot ({\x}, {-0.56*(\x-1.2) + 1.392});
  \draw[purple, thick, domain=1.8:5.0]
    plot ({\x}, {0.40*(\x-3.6) + 1.200});

  \fill[orange!90!black] (1.2,1.392) circle (1.6pt);
  \fill[purple] (3.6,1.200) circle (1.6pt);

  \node[orange!90!black, font=\scriptsize, align=center] (LabU) at (1.45,2.45)
    {slope $f'_{x,y}(s) = Df_u(y-x) < 0$};
  \draw[orange!90!black, thin, ->] (LabU) -- (1.2,1.47);
  \node[purple, font=\scriptsize, align=center] (LabV) at (4.15,0.45)
    {slope $f'_{x,y}(t) = Df_v(y-x) > 0$};
  \draw[purple, thin, ->] (LabV) -- (3.6,1.12);

  \draw (0,0.05) -- (0,-0.05) node[below, font=\small] {$0$};
  \draw (1.2,0.05) -- (1.2,-0.05) node[below, font=\small] {$s$};
  \draw (3.6,0.05) -- (3.6,-0.05) node[below, font=\small] {$t$};
  \draw (4.9,0.05) -- (4.9,-0.05) node[below, font=\small] {$1$};
  \draw[dashed, gray] (1.2,0) -- (1.2,1.392);
  \draw[dashed, gray] (3.6,0) -- (3.6,1.200);
\end{tikzpicture}
\end{center}

The picture records both halves of \textbf{(b)} at once: each tangent line lies weakly
\emph{below} the graph --- which is the inequality $f(y)-f(x) \geq Df_x(y-x)$ read along the
segment --- and the slope grows from $s$ to $t$, which is $f'_{x,y}$ increasing.

% Sonnet 5 (Medium)
We now change subject entirely: from convexity to the question of when a map is locally
invertible --- the tool that makes the implicit function theorem, and later submanifolds
(\cref{sec:submanifolds}), possible.


% --- exercises relocated here from the dissolved sheets ---

% Originally: content/week-06.tex, \exercisesheet{6}, problem 6.7
% Source: exercises/Ex6_Analysis2_eng.pdf

\begin{exercise}[6.7 --- Multiple choice \textnormal{(important)}]
\label{ex:6.7}
Among the following statements about convex functions mark those (and only those) which are
always true.
\begin{enumerate}[label=\textbf{(\alph*)}]
  \item If $f \in C^1(U)$ is convex in some open convex set $U \subset \mathbb{R}^n$ and $f$ has a
    local maximum at $z \in U$, then $\nabla f \equiv 0$ in $U$.
  \item If $f \in C^1(U)$ is convex in some open convex set $U \subset \mathbb{R}^n$ and $f$ has a
    global maximum at $z \in U$, then $\nabla f \equiv 0$ in $U$.
  \item Assume $f_n \in C^2(\mathbb{R})$ is a sequence of convex functions that converge
    pointwise to some $f : \mathbb{R} \to \mathbb{R}$. Is $f$ necessarily convex?
  \item There exists a convex function $f \in C^\infty(\mathbb{R}^2)$ such that
    \[f(x) = 1 - 2x_1 + x_2^3 + O(|x|^4) \quad \text{as } |x| \to 0.\]
  \item There exists a convex function $f \in C^\infty(\mathbb{R}^2)$ such that
    \[f(x) = 1 - 2x_1 + x_2^4 + O(|x|^4) \quad \text{as } |x| \to 0.\]
  \item A convex set is not necessarily connected.
\end{enumerate}
\end{exercise}

% Originally: the \section{Solutions} of content/week-06.tex

\section{Solutions}
\label{sec:convexity_solutions}

\begin{exercisesolution}[Solution to \cref{ex:6.7}]
% Generator: Claude Sonnet 5 (Medium)
\begin{enumerate}[label=\textbf{(\alph*)}]
  \item \textbf{False.} Let $f(x) := \max\{x_1, \tfrac12\}^4$ on $U := B_1 \subset \mathbb{R}^n$.
    The inner function $x \mapsto \max\{x_1,\tfrac12\}$ is convex and takes values in
    $[\tfrac12,\infty)$, and $t \mapsto t^4$ is convex and increasing there. A convex
    increasing function composed with a convex function is convex, so $f$ is convex. For all $x \in U$ with $x_1 \leq \tfrac12$,
    $f(x) = (\tfrac12)^4$ identically, so $f$ has a (non-strict) local maximum at $x=0$, yet
    $\nabla f$ does not vanish once $x_1 > \tfrac12$.

    \begin{remark}
    Worth noting: $x=0$ here is simultaneously a (non-strict) local \emph{minimum} too, since $f$
    is locally constant on the whole region $\{x_1 \leq \tfrac12\} \cap U$. This is exactly why
    the counterexample doesn't contradict \textbf{(b)}: $x=0$ is not a \emph{global} maximum of
    $f$ on $U$ (as $f$ grows without bound as $x_1 \to 1$), so \textbf{(b)}'s stronger hypothesis
    simply isn't triggered.
    \end{remark}
  \item \textbf{True.} Since $z$ is a local maximum, $\nabla f(z) = 0$
    (\cref{prop:extremum_implies_critical}). Convexity gives $f(x) \geq f(z) + \nabla f(z)\cdot(x-z)
    = f(z)$ for all $x \in U$; but $z$ being a global maximum also gives $f(x) \leq f(z)$ for all
    $x \in U$. Hence $f \equiv f(z)$ on $U$, so $\nabla f \equiv 0$ in $U$.
  \item \textbf{True.} Fix $x, y \in U$ and $t \in [0,1]$. Convexity of each $f_n$ gives
    $f_n(tx+(1-t)y) \leq tf_n(x) + (1-t)f_n(y)$ for every $n$; letting $n \to \infty$ (with
    $x,y,t$ fixed) and using pointwise convergence yields the same inequality for $f$. Since
    $x,y,t$ were arbitrary, $f$ is convex.
  \item \textbf{False.} A convex function lies above its tangent plane, so from the Taylor data at
    $0$, any such $f$ would satisfy $f(x) \geq 1-2x_1$ for all $x \in \mathbb{R}^n$. Combined with
    $f(x) = 1-2x_1+x_2^3+O(|x|^4)$, this forces $x_2^3 \geq -M|x|^4$ near $0$ for some $M>0$,
    which fails along $x = (0,-r,0,\dots,0)$ as $r \downarrow 0$ (left side $\sim -r^3$, right
    side $\sim -Mr^4$, and $-r^3 < -Mr^4$ for small $r$). Contradiction, so no such convex $f$
    exists.
  \item \textbf{True.} Simply take $f(x) := 1-2x_1+x_2^4$ itself (with the $O(|x|^4)$ remainder
    identically zero): it is convex, being a sum of the affine (hence convex) term $1-2x_1$ and
    the convex term $x_2^4$.
  \item \textbf{False.} A convex set is always path-connected, since for any two points the
    straight segment joining them lies inside the set by definition of convexity; path-connected
    sets are connected.
\end{enumerate}
\end{exercisesolution}



\part{Inverse and Implicit Functions; Submanifolds}

% Generator: Claude Opus 5 (High) --- part-level framing prose, new content.
Two questions run through this part. When can an equation $f(x,y) = 0$ be solved for $y$ as a
function of $x$, and when is a map invertible near a point? Neither has a workable global answer,
but both have clean local ones, and both reduce to the same linear-algebra hypothesis:
invertibility of a Jacobian. That is the content of the implicit and inverse function theorems,
which turn out to be equivalent to each other and are proved here from the contraction principle
of Part II. Their deeper use is definitional. A level set on which the differential is surjective
is locally a graph, and being locally a graph is precisely what it means to be a submanifold, a
subset of $\mathbb{R}^n$ that looks like $\mathbb{R}^m$ from close up. The closing chapters
develop that notion, since submanifolds are the objects the last two parts integrate over.

% ============================================================
% Chapter 15 --- The inverse function theorem
% Originally: content/week-06.tex, \section{The inverse function theorem} (Corsin Week 6),
%             with problem 7.1 from sheet 7.
% ============================================================
\chapter{The Inverse Function Theorem}
\label{ch:inverse_function_theorem}

% Transition: Claude Opus 5 (Medium)
So far the differential has been used to \emph{describe} a map near a point. The inverse function
theorem turns that around. If the differential at a point is invertible, then so is the map
itself near that point, and the inverse is exactly as differentiable as the map was. Everything
in the rest of this part rests on that one statement.

% Originally: content/week-06.tex, \section{The inverse function theorem} (Corsin Week 6)

% Source: Corsin Nick/Class Notes/Week 6.pdf, pp. 8--9
\section{The inverse function theorem}
\label{sec:inverse_function_theorem}

\begin{theorem}[Inverse function theorem]
\label{thm:inverse_function_theorem}
Let $U \subseteq \mathbb{R}^n$ be open, let $f \in C^k(U,\mathbb{R}^n)$, and let $x_0 \in U$ be
such that $Df_{x_0}$ is invertible. Then there exists an open neighbourhood $U_0 \subseteq U$ of
$x_0$ such that
\[f|_{U_0} : U_0 \to f(U_0)\]
is a $C^k$-\newterm{diffeomorphism} (\germanterm{Diffeomorphismus}) and
\[D\big(f|_{U_0}^{-1}\big)_{f(x)} = D\big(f|_{U_0}\big)_x^{-1} \quad \text{for all } x \in U_0.\]
\end{theorem}

% Supplement: Sascha Brack/Class Notes/Week_06_Notes_Friday.pdf, p. 2
\begin{remark}[Checking for a global diffeomorphism]
\label{rem:global_diffeomorphism_check}
To verify that a map $f \in C^k(U, V)$ between open sets is a global $C^k$-diffeomorphism, it is best to check the following two conditions:
\begin{enumerate}
  \item \textbf{Global bijectivity:} Check that $f$ is bijective. This assures that a global inverse $f^{-1} : V \to U$ exists, and any local inverse will coincide with it.
  \item \textbf{Local diffeomorphism everywhere:} Check that the differential $Df_x$ (or $\Jac f(x)$) is invertible at every point $x \in U$. By the inverse function theorem, this ensures $f$ is a local $C^k$-diffeomorphism everywhere.
\end{enumerate}
If both conditions hold, the global inverse $f^{-1}$ inherits the $C^k$ regularity of the local inverses, making $f$ a global $C^k$-diffeomorphism.
\end{remark}

% Supplement: Sascha Brack/Class Notes/Week_06_Notes_Monday.pdf, p. 4
\begin{example}[Polar and spherical coordinates]
\label{ex:polar_spherical_coordinates}
The most important examples of diffeomorphisms are transformations between coordinate systems. For instance, the map assigning polar coordinates to Cartesian coordinates on $\mathbb{R}^2 \setminus ( (-\infty, 0] \times \{0\} )$, and similarly for spherical coordinates on $\mathbb{R}^3$, are global $C^\infty$-diffeomorphisms onto their images.
\end{example}

\begin{question}
\begin{enumerate}[label=\textbf{Q\arabic*.}]
  \item If $f : \mathbb{R}^n \to \mathbb{R}^n$ is $C^k$ and a \newterm{homeomorphism}
    (\germanterm{Hom\"oomorphismus}) --- that is, bijective with continuous inverse --- is it
    necessarily a $C^k$-diffeomorphism?
  \item If $f \in C^\infty(U,\mathbb{R}^n)$ has an everywhere invertible differential, i.e.,
    $Df_x$ is invertible for all $x \in U$ with $U \subseteq \mathbb{R}^n$ open, then $f$ is a
    $C^\infty$-diffeomorphism onto its image. True or false?
\end{enumerate}
\end{question}

\begin{answer}
\begin{enumerate}[label=\textbf{A\arabic*.}]
  \item \textbf{False.} Consider $f : \mathbb{R} \to \mathbb{R}$, $x \mapsto x^3$. Then
    $f^{-1} : \mathbb{R} \to \mathbb{R}$, $x \mapsto \sqrt[3]{x}$, is not continuously
    differentiable at $0 \in \mathbb{R}$.
  \item \textbf{False.} Consider
    \[\cos|_{\mathbb{R}\setminus\pi\mathbb{Z}} : \mathbb{R}\setminus\pi\mathbb{Z} \to (-1,1), \qquad x \mapsto \cos(x).\]
    For any $x \in \mathbb{R}\setminus\pi\mathbb{Z}$, $\cos'(x) = -\sin(x) \neq 0$, so the
    differential is everywhere invertible, but the cosine is not injective on this domain, so it
    fails to be a diffeomorphism onto its image.
\end{enumerate}
\end{answer}

\begin{ainote}
Corsin illustrates \textbf{A2} with a second sketch: one $C^\infty$ curve carrying all three
possibilities at once. The figure below reproduces it, with the three regions labelled by
\emph{which hypothesis of the inverse function theorem fails}.
\end{ainote}

% Sonnet 5 (Medium) --- FIG-W06-03, drawn from the page image seen while transcribing.
% The first region must have a genuinely HORIZONTAL tangent (f'=0) -- drawn as an explicit
% cubic 1.95 + 0.15(x-2.15)^3, which is strictly increasing yet has f'(2.15) = 0. A merely
% shallow slope would not do: where f' != 0 the IFT applies and f IS a local diffeomorphism.
\begin{center}
\begin{tikzpicture}[>=Latex, scale=1.05]
  \draw[->] (-0.3,0) -- (8.4,0) node[right, font=\small] {$x$};
  \draw[->] (0,-0.3) -- (0,5.6) node[above, font=\small] {$y$};

  % Region 1: strictly increasing, but with f'(2.15) = 0 (an x^3-type inflection).
  \draw[blue!80!black, thick, domain=0.25:3.1, samples=60, smooth]
    plot ({\x}, {1.95 + 0.15*(\x-2.15)*(\x-2.15)*(\x-2.15)});
  % Regions 2 and 3: a hump (three preimages), then a strictly monotone stretch.
  \draw[blue!80!black, thick] plot[smooth] coordinates {
    (3.1,2.079) (3.5,2.40) (4.0,2.62) (4.5,2.45) (4.95,2.20)
    (5.4,2.50) (6.1,3.20) (6.9,4.00) (7.6,4.70)};

  % Region 1 marker: the horizontal tangent itself.
  \draw[red!80!black, dashed, thick] (1.45,1.95) -- (2.85,1.95);
  \fill[red!80!black] (2.15,1.95) circle (1.7pt);
  \draw (2.15,0.05) -- (2.15,-0.05) node[below, font=\scriptsize] {$x_0$};
  \node[red!80!black, font=\scriptsize, align=center] (LabFlat) at (1.15,3.9)
    {$f'(x_0)=0$:\\ injective, but $f^{-1}$\\ not differentiable};
  \draw[red!80!black, thin, ->] (LabFlat) -- (2.10,2.10);

  % Region 2 marker: one value, three preimages.
  \draw[orange!90!black, dashed, thick] (3.2,2.35) -- (5.65,2.35);
  \fill[orange!90!black] (3.45,2.35) circle (1.5pt);
  \fill[orange!90!black] (4.61,2.35) circle (1.5pt);
  \fill[orange!90!black] (5.24,2.35) circle (1.5pt);
  \node[orange!90!black, font=\scriptsize, align=center] (LabHump) at (4.4,1.15)
    {not injective:\\ one value, three preimages};
  \draw[orange!90!black, thin, ->] (LabHump) -- (4.55,2.24);

  % Region 3 marker: the IFT applies.
  \node[green!60!black, font=\scriptsize, align=center] (LabIFT) at (5.75,4.75)
    {$f'\neq0$ and injective:\\ diffeomorphism by the IFT};
  \draw[green!60!black, thin, ->] (LabIFT) -- (6.85,4.00);
\end{tikzpicture}
\end{center}

Reading the three regions left to right makes the two answers above concrete.

\begin{itemize}
  \item At $x_0$ the curve is still \emph{strictly increasing}, so $f$ is injective there and
    $f^{-1}$ exists --- but the tangent is horizontal, so the tangent to $f^{-1}$ at $f(x_0)$ is
    \emph{vertical} and $f^{-1}$ fails to be differentiable. Invertible, yet not a
    diffeomorphism: this is exactly $x \mapsto x^3$ from \textbf{A1}, and it is why
    \Cref{thm:inverse_function_theorem} assumes $Df_{x_0}$ invertible rather than merely
    assuming $f$ bijective.
  \item On the hump, $f' \neq 0$ at every individual point, so $f$ \emph{is} a local
    diffeomorphism near each of them --- yet one value has three preimages, so $f$ is not
    globally injective. This is \textbf{A2}: an everywhere-invertible differential buys locality
    and nothing more.
  \item On the final stretch both hypotheses hold, and the inverse function theorem delivers a
    genuine $C^\infty$-diffeomorphism onto the image.
\end{itemize}

% Sonnet 5 (Medium) --- FIG-W06-02, drawn from the page image seen while transcribing
\begin{center}
\begin{tikzpicture}[>=Latex, scale=1.1]
  \draw[->] (-0.3,0) -- (7.3,0) node[right, font=\small] {$x$};
  \draw[->] (0,-1.6) -- (0,1.6) node[above, font=\small] {$y$};
  \draw (0,1) -- (-0.05,1) node[left, font=\footnotesize] {$1$};
  \draw (0,-1) -- (-0.05,-1) node[left, font=\footnotesize] {$-1$};

  \draw[blue!80!black, thick, domain=0:7, samples=120, smooth]
    plot ({\x}, {cos(\x r)}) node[right, font=\small] {$\cos(x)$};

  \draw[red!80!black, dashed, thick] (0,0.4) -- (7,0.4);
  \node[red!80!black, font=\scriptsize, right] at (7,0.4) {not injective};
\end{tikzpicture}
\end{center}

% Sonnet 5 (Medium)
The inverse function theorem asks when $f(x)=y$ can be solved for $x$ as a function of $y$. The
implicit function theorem generalizes this: it asks when an equation $f(x,y)=0$ can be solved for
\emph{some} of the variables in terms of the others, without $f$ itself being invertible.

% Originally: content/week-07.tex, \section{Solutions}

\section{Solutions}
\label{sec:ift_solutions}

% Transition: Claude Opus 5 (High)
Not every solution in this chapter is collected here. \Cref{ex:open_mapping_grad_nonzero} is solved
where it appears, on \cpageref{ex:open_mapping_grad_nonzero}.

\begin{exercisesolution}[Solution to \cref{ex:7.1}]
% Generator: Claude Sonnet 5 (Medium)
\begin{enumerate}[label=\textbf{(\alph*)}]
  \item First build a diffeomorphism $\mathbb{R} \to (0,1)$: the function
    \[\psi : \mathbb{R} \to (0,1), \qquad \psi(t) := \frac12+\frac1\pi\arctan(t)\]
    is $C^\infty$, strictly increasing (since $\arctan' > 0$), and bijective onto $(0,1)$ (as
    $\arctan$ is bijective onto $(-\tfrac\pi2,\tfrac\pi2)$), with $C^\infty$ inverse
    $\psi^{-1}(s) = \tan\big(\pi(s-\tfrac12)\big)$. Applying $\psi$ in each coordinate gives the
    diffeomorphism
    \[\phi : \mathbb{R}^2 \to (0,1)\times(0,1), \qquad \phi(x,y) := \Big(\frac12+\frac1\pi\arctan(x),\ \frac12+\frac1\pi\arctan(y)\Big).\]
  \item \textbf{No.} $f$ is a $C^\infty$ bijection of $\mathbb{R}$ with $f^{-1}(x) =
    \sqrt[5]{x}$, but $f'(0) = 5\cdot0^4 = 0$. If $f$ were a diffeomorphism, then by the
    chain rule applied to $f^{-1}\circ f = \id$,
    \[1 = (f^{-1})'(f(x))\cdot f'(x) \quad \text{for all } x,\]
    which is impossible at $x=0$ since $f'(0)=0$. Equivalently, $(f^{-1})'(x) =
    \tfrac15x^{-4/5} \to \infty$ as $x \to 0$, so $f^{-1}$ is not even differentiable at $0$,
    let alone continuously so.
\end{enumerate}
\end{exercisesolution}

\begin{exercisesolution}[Solution of \cref{ex:hessian_gradient_diffeomorphism}]
% Generator: Claude Opus 5 (High)
Write $H := \Hess f(0)$. Everything follows from its spectrum, so compute that first. The matrix is
block diagonal, with the $1\times1$ block $(1)$ and the $2\times2$ block
$\left(\begin{smallmatrix} 0 & 1 \\ 1 & 0\end{smallmatrix}\right)$, whose eigenvalues are $\pm 1$.
Hence $H$ has eigenvalues $1, 1, -1$, and in particular $\det H = -1 \neq 0$: the critical point is
non-degenerate, and $H$ is indefinite.

Since $f$ is smooth with $f(0) = 0$ and $\nabla f(0) = 0$, Taylor's theorem gives
\begin{equation}
\label{eq:degenerate_critical_taylor}
f(x) = \tfrac12 \langle Hx, x\rangle + o(\lvert x\rvert^2)
     = \tfrac12 x_1^2 + x_2x_3 + o(\lvert x\rvert^2)
     \qquad \text{as } x \to 0 .
\end{equation}
\begin{enumerate}[label=\textbf{(\alph*)}]
  \item \textbf{False.} Along the $x_1$-axis, \eqref{eq:degenerate_critical_taylor} gives
    $f(t,0,0) = \tfrac12 t^2 + o(t^2) > 0 = f(0)$ for small $t \neq 0$, so $f$ takes values larger
    than $f(0)$ arbitrarily close to the origin.
  \item \textbf{True.} We have just seen a direction along which $f$ increases. In the plane
    $x_1 = 0$ the quadratic part is $x_2x_3$, so along $x_2 = -x_3 = t$ it reads $-t^2$, and
    $f(0,t,-t) = -t^2 + o(t^2) < 0 = f(0)$. Taking both values above and below $f(0)$ in every
    neighbourhood of the origin is exactly what a saddle point does.
  \item \textbf{False.} The quadratic form is correct only up to a change of coordinates, and the
    statement is made in the original ones. By \eqref{eq:degenerate_critical_taylor} the second-order
    part of $f$ is $\tfrac12 x_1^2 + x_2 x_3$, which is not $x_1^2 - x_2^2 - x_3^2$: evaluating both
    at $(0,1,1)$ gives $1$ and $-2$. What \emph{is} true is that in suitable coordinates the form
    becomes $\tfrac12(u_1^2 + u_2^2 - u_3^2)$, matching the eigenvalues $1,1,-1$ --- note in
    passing that this has signature $(2,1)$, not the $(1,2)$ the stated expression suggests.
  \item \textbf{True.} This is the same Taylor expansion read with a weaker conclusion: the
    right-hand side of \eqref{eq:degenerate_critical_taylor} is bounded by a constant times
    $\lvert x\rvert^2$ near $0$. Vanishing to first order at a point is precisely what $O(\lvert
    x\rvert^2)$ records.
  \item \textbf{True}, and this is where the chapter's theorem enters. The map $\nabla f$ is smooth,
    being the gradient of a smooth function, and its Jacobian at the origin is
    $\Jac(\nabla f)(0) = \Hess f(0) = H$, which we computed to be invertible. By the inverse function
    theorem (\cref{thm:inverse_function_theorem}), $\nabla f$ is therefore a diffeomorphism from some
    ball $B_\varepsilon(0)$ onto an open neighbourhood of $\nabla f(0) = 0$.
\end{enumerate}
\begin{remark}
Item \textbf{(e)} is the useful one to remember: \qt{non-degenerate critical point} means exactly
\qt{$\nabla f$ is a local diffeomorphism there}, so such points are isolated. Item \textbf{(c)} is
the standard trap --- diagonalising the Hessian is a statement about \emph{some} coordinates, and
carrying its conclusion back to the original ones without the change of variables is where the
argument breaks.
\end{remark}
\end{exercisesolution}

\begin{exercisesolution}[Solution of \cref{ex:local_diffeo_parameter_alpha}]
% Generator: Claude Opus 5 (High)
Everything turns on one number:
\[\det\Jac\Phi(0,0) = (\alpha+5)(\alpha-5) - (-3)\cdot 3 = \alpha^2 - 25 + 9 = \alpha^2 - 16 .\]
\begin{enumerate}[label=\textbf{(\alph*)}]
  \item \textbf{False.} The determinant vanishes exactly when $\alpha^2 = 16$, so $\alpha^2 \neq 25$
    is the wrong condition. Take $\alpha = 4$, which satisfies $\alpha^2 = 16 \neq 25$ and so is
    covered by the claim. There $\det\Jac\Phi(0,0) = 0$, and no restriction of $\Phi$ to a
    neighbourhood of the origin can be a diffeomorphism: if $\Phi|_V$ had a $C^1$ inverse $\Psi$,
    the chain rule applied to $\Psi\circ\Phi = \id$ would give
    $\Jac\Psi(\Phi(0,0))\,\Jac\Phi(0,0) = I_2$, forcing $\Jac\Phi(0,0)$ to be invertible.
  \item \textbf{True.} If $\alpha^2 \neq 16$, then $\det\Jac\Phi(0,0) \neq 0$, so
    $D\Phi_{(0,0)}$ is invertible, and $\Phi \in C^1$. The inverse function theorem
    (\cref{thm:inverse_function_theorem}) supplies open sets $V \ni (0,0)$ and $W \ni \Phi(0,0)$
    with $\Phi|_V : V \to W$ a $C^1$-diffeomorphism. Any ball $B_r(0,0) \subseteq V$ then works,
    since a restriction of a diffeomorphism to an open subset is again one onto its image.
  \item \textbf{True.} For $\alpha = 0$ the determinant is $-16 \neq 0$, so \textbf{(b)} applies.
    Fix $V, W$ as there and $r > 0$ with $B_r(0,0) \subseteq V$. A diffeomorphism is in particular
    a homeomorphism onto the open set $W$, so it maps open subsets of $V$ to open subsets of $W$,
    and open subsets of an open set are open in $\mathbb{R}^2$. Hence $\Phi(B_r(0,0))$ is open.
\end{enumerate}
\begin{remark}[Where $25$ comes from, and why it is wrong]
The two thresholds are not a random pair. A reader who computes the determinant as a product of
the diagonal entries and forgets the off-diagonal contribution gets $(\alpha+5)(\alpha-5)$ and
lands on $\alpha^2 = 25$, which is exactly the false condition of \textbf{(a)}. The off-diagonal
pair $-3, 3$ contributes $+9$ and moves the singular set to $\alpha^2 = 16$, that is
$\alpha = \pm4$. Item \textbf{(a)} and item \textbf{(b)} are therefore the same sentence with one
arithmetic slip between them, and reading them as a pair is faster than checking either alone.
\end{remark}
\end{exercisesolution}

\begin{exercisesolution}[Solution of \cref{ex:kobel_ift_single_choice}]
% Generator: Claude Opus 5 (High)
The correct answer is \textbf{(b)}.
\begin{enumerate}[label=\textbf{(\alph*)}]
  \item \textbf{False}, and it fails before any hypothesis on the derivative is examined.
    \Cref{thm:inverse_function_theorem} needs domain and codomain of the \emph{same} dimension,
    because a local inverse has to be defined on a neighbourhood of the image point. For
    $n > 2$ the image of $f$ has empty interior in $\mathbb{R}^n$ --- it is at best a
    $2$-dimensional surface --- so no neighbourhood of $f(0,0)$ in $\mathbb{R}^n$ lies in the
    image, and there is nothing for an inverse to be defined on. Being locally \emph{injective},
    which is what an immersion gives, is a strictly weaker conclusion.
  \item \textbf{True.} For a map $\mathbb{R}^2 \to \mathbb{R}^2$, maximal rank means rank $2$,
    which for a square matrix is invertibility, that is $\det\Jac f(0,0) \neq 0$. That is exactly
    the hypothesis of \cref{thm:inverse_function_theorem}, and $f \in C^2 \subseteq C^1$ supplies
    the regularity.
  \item \textbf{False.} One non-vanishing partial derivative is one column of the Jacobian, and a
    matrix with one non-zero column can still be singular. Take $f(x,y) := (x,0)\transp$: then
    $\partial f/\partial x = (1,0)\transp \neq 0$ everywhere, while
    $\Jac f = \left(\begin{smallmatrix}1&0\\0&0\end{smallmatrix}\right)$ has determinant $0$, and
    $f$ is not injective on any neighbourhood of the origin --- it collapses every vertical
    segment to a point.
  \item \textbf{False}, by the same example with the roles of the variables exchanged: for
    $f(x,y) := (y,0)\transp$ the partial $\partial f/\partial y$ is non-zero and the Jacobian is
    again singular.
\end{enumerate}
\begin{remark}[The same answer if the partials are read component-wise]
Items \textbf{(c)} and \textbf{(d)} are printed with $f$ itself, so
$\partial f/\partial x(0,0) \neq 0$ says a whole column of the Jacobian is non-zero. A reader who
takes $f = (f_1,f_2)\transp$ and reads them as $\partial f_1/\partial x(0,0) \neq 0$ gets a weaker
hypothesis --- a single \emph{entry} is non-zero --- and the same counterexample refutes it. The
question is robust to the ambiguity, which is presumably why it went unnoticed.

What both readings test is the difference between \qt{some derivative does not vanish} and
\qt{the derivative is invertible}. Only the second is a statement about the linear map $Df_{(0,0)}$
as a whole, and only the second is what \cref{rem:inverse_function_theorem_intuition} argues is the
right hypothesis.
\end{remark}
\end{exercisesolution}


% ============================================================
% Chapter 16 --- The implicit function theorem
% Originally: content/week-06.tex, \section{The implicit function theorem} (Corsin Week 6),
%             with problem 7.5 from sheet 7 (added in the next stage).
% ============================================================
\chapter{The Implicit Function Theorem}
\label{ch:implicit_function_theorem}

% Transition: Claude Opus 5 (Medium)
The inverse function theorem asks when a map can be undone. The implicit function theorem
generalizes the question: when can an equation $f(x,y) = 0$ be solved for \emph{some} of the
variables in terms of the others, without $f$ itself being invertible? The answer is the same
condition applied to a submatrix of the Jacobian, and it is what turns level sets into graphs.

% Originally: content/week-06.tex, \section{The implicit function theorem} (Corsin Week 6)

% Source: Corsin Nick/Class Notes/Week 6.pdf, pp. 10--12
\section{The implicit function theorem}
\label{sec:implicit_function_theorem}

\begin{theorem}[Implicit function theorem]
\label{thm:implicit_function_theorem}
Let $1 \leq n-d < n$, let $f \in C^k(\mathbb{R}^d\times\mathbb{R}^{n-d}, \mathbb{R}^{n-d})$ (or an
open subset of $\mathbb{R}^d\times\mathbb{R}^{n-d}$), and suppose $f(x_0,y_0) = 0$ for some
$(x_0,y_0) \in \mathbb{R}^d\times\mathbb{R}^{n-d}$ such that the differential $D_yf(x_0,y_0)$ of
$y \mapsto f(x_0,y)$ at $y_0$ is bijective (so $\Jac_yf(x_0,y_0)$ is invertible). Then there exist
balls $B_r(x_0) \subseteq \mathbb{R}^d$ and $B_s(y_0) \subseteq \mathbb{R}^{n-d}$ and
$g \in C^k(B_r(x_0), B_s(y_0))$ such that, for all $(x,y) \in B_r(x_0)\times B_s(y_0)$,
\[f(x,y) = 0 \iff y = g(x).\]
Furthermore,
\[\Jac g(x) = -\big(\Jac_yf(x,g(x))\big)^{-1}\Jac_xf(x,g(x)).\]
\end{theorem}

% Sonnet 5 (Medium)
In words: the equation $f(x,y)=0$ ties $x$ and $y$ together without telling us what $y$ is. The
theorem says that near a point where the tie is \emph{non-degenerate in $y$}, the equation can be
untangled --- there is an honest function $g$ producing the right $y$ for each nearby $x$, it is
as smooth as $f$ was, and its derivative is computable without ever writing $g$ down.

Three things are worth reading off the statement before the pictures.
\begin{itemize}
  \item \textbf{Which variables get solved for is a choice.} Splitting the $n$ coordinates into
    $d$ \qt{free} ones and $n-d$ \qt{solved} ones is ours to make; the theorem only demands that
    the block $\Jac_yf$ belonging to the solved variables be invertible. Choose a different
    split and a different block has to be invertible.
  \item \textbf{The conclusion is an equivalence, not just an inclusion.} Inside the box,
    $f(x,y)=0 \iff y=g(x)$. So $g$ does not merely produce \emph{some} solutions: within the
    box there are no others.
  \item \textbf{Nothing is claimed globally.} The level set $\{f(x,y)=0\}$ is a graph over $x$
    only near the chosen point --- \cref{sec:submanifolds} will call such a set a
    submanifold, and this local-graph property is exactly what that definition abstracts.
\end{itemize}

% Supplement: Sascha Brack/Class Notes/Week_06_Notes_Friday.pdf, pp. 3--5
\begin{example}[A concrete example in $\mathbb{R}^2$]
\label{ex:implicit_function_circle}
Consider $f : \mathbb{R}^2 \to \mathbb{R}$ defined by $f(x,y) := x^2 + y^2 - 1$.
Here $n=2$ and $d=1$, so $x \in \mathbb{R}$ and $y \in \mathbb{R}$. We want to see whether we can write $y$ locally as a function of $x$, i.e., find $g(x)$ such that $f(x,g(x)) = 0$.

First, let us find a point $(x_0,y_0)$ where $f(x_0,y_0) = 0$. We guess $f(0,1) = 0^2+1^2-1 = 0$, so $(0,1)$ works.
Next, we check whether the $y$-derivative $\Jac_y f(0,1)$ is invertible. Since $y$ is 1-dimensional, this is just the partial derivative $\partial_y f(0,1)$:
\[\partial_y f(0,1) = \partial_y(x^2+y^2-1)\big|_{(0,1)} = 2y\big|_{y=1} = 2 \neq 0.\]
Since it is non-zero, it is invertible as a $1\times 1$ matrix.
The theorem then guarantees that for small $r, s > 0$, there exists $g : B_r(0) \to B_s(1)$ such that for all $(x,y) \in B_r(0)\times B_s(1)$,
\[f(x,y) = 0 \iff y = g(x).\]
Since $f \in C^\infty(\mathbb{R}^2)$, we also have $g \in C^\infty(B_r(0))$.
Moreover, we can compute $g'(x) = \Jac g(x)$ using the formula:
\[g'(x) = - \big(\partial_y f(x,g(x))\big)^{-1} \partial_x f(x,g(x)) = - (2g(x))^{-1} (2x) = -\frac{x}{g(x)}.\]
(Indeed, since $y=g(x) = \sqrt{1-x^2}$ near $(0,1)$, its derivative is exactly $-x/\sqrt{1-x^2}$.)
\end{example}

\begin{remark}[Why $\Jac g$ has this form]
The formula is just implicit differentiation, done in blocks. By construction $g$ satisfies
$f(x,g(x)) \equiv 0$ for every $x \in B_r(x_0)$: the left-hand side is identically zero as a
function of $x$, so its derivative is too. Differentiating with the chain rule
(\cref{thm:chain_rule}), and remembering that $x$ enters both directly and through $g$,
\[0 = \underbrace{D_xf(x,g(x))}_{\text{direct}} + \underbrace{D_yf(x,g(x))\cdot Dg(x)}_{\text{through } g},\]
and since $D_yf(x,g(x))$ is invertible near $(x_0,y_0)$ we may move it to the other side:
\[Dg(x) = -\big(D_yf(x,g(x))\big)^{-1} D_xf(x,g(x)),\]
which is the formula in \cref{thm:implicit_function_theorem}.

Read it as an exchange rate. Nudging $x$ changes $f$ by $D_xf$; to keep the equation satisfied,
$y$ must move by whatever compensates, and $\big(D_yf\big)^{-1}$ converts \qt{how much $f$ moved}
back into \qt{how far $y$ must travel}. The minus sign is the compensation. This is also why
invertibility of $D_yf$ is the hypothesis: if $D_yf$ were singular, some change in $f$ could not
be absorbed by moving $y$ at all, and there would be nothing for $g$ to do.
\end{remark}

\begin{ainote}
Corsin remarks in his own notes that ``the geometrical meaning of this one eludes me, sorry''.
The exchange-rate reading above and the two figures below are this document's attempt at the
geometry he was after; the algebra itself is his chain-rule derivation.
\end{ainote}

% Sonnet 5 (Medium) --- FIG-W06-04, replacing the original arrow diagram.
% The hypothesis "D_y f invertible" is a statement about the DIRECTION of the level set at
% the point, so the figure shows the level set and its tangent, not an abstract arrow.
\subsubsection*{What the hypothesis means}

The condition on $D_yf(x_0,y_0)$ is a statement about the \emph{direction} in which the level
set $\{f=0\}$ runs at the point. Moving off $(x_0,y_0)$ straight up, in the $y$-direction alone,
changes $f$ to first order by $D_yf(x_0,y_0)$. If that map is invertible, no vertical motion
leaves $f$ unchanged, so the level set cannot run vertically at $(x_0,y_0)$ --- and a curve that
is not vertical is locally a graph over $x$. If instead $D_yf(x_0,y_0)=0$, the level set is free
to turn vertical, and above a single $x$ there may sit two values of $y$, or none.

\begin{center}
\begin{tikzpicture}[>=Latex, scale=1.05]
  % --- Left panel: D_y f invertible, level set transverse to the vertical ---
  \begin{scope}
    \draw[green!60!black, thick, ->] (-0.3,0) -- (3.5,0) node[right, font=\small] {$x$};
    \draw[orange!90!black, thick, ->] (0,-0.3) -- (0,3.2) node[above, font=\small] {$y$};

    \draw[blue!80!black, very thick, domain=0.35:3.1, samples=60, smooth]
      plot ({\x}, {1.55 + 0.55*(\x-1.6) - 0.16*(\x-1.6)*(\x-1.6)});
    \node[blue!80!black, font=\scriptsize, below right] at (2.75,1.95) {$\{f=0\}$};

    \coordinate (P) at (1.6,1.55);
    \draw[dotted, gray] (1.6,0) -- (P) -- (0,1.55);
    \node[below, font=\scriptsize] at (1.6,-0.05) {$x_0$};
    \node[left, font=\scriptsize] at (-0.05,1.55) {$y_0$};
    \fill (P) circle (1.7pt);

    % vertical probe: moving in y alone leaves the level set
    \draw[red!80!black, very thick, ->] (P) -- ++(0,0.85);
    \node[red!80!black, font=\scriptsize, align=center] at (1.6,2.78)
      {move in $y$ only:\\ $f$ changes};

    \node[font=\scriptsize, align=center] at (1.55,-0.85)
      {$D_yf(x_0,y_0)$ invertible\\ \textbf{solvable for $y$}};
  \end{scope}

  % --- Right panel: D_y f = 0, level set vertical, no graph over x ---
  \begin{scope}[shift={(5.4,0)}]
    \draw[green!60!black, thick, ->] (-0.3,0) -- (3.5,0) node[right, font=\small] {$x$};
    \draw[orange!90!black, thick, ->] (0,-0.3) -- (0,3.2) node[above, font=\small] {$y$};

    \draw[blue!80!black, very thick, domain=0.35:2.75, samples=60, smooth]
      plot ({1.6 - 1.05*(\x-1.55)*(\x-1.55)}, {\x});
    \node[blue!80!black, font=\scriptsize, right] at (0.30,0.38) {$\{f=0\}$};

    \coordinate (Q) at (1.6,1.55);
    \draw[dotted, gray] (1.6,0) -- (Q) -- (0,1.55);
    \node[below, font=\scriptsize] at (1.6,-0.05) {$x_0$};
    \node[left, font=\scriptsize] at (-0.05,1.55) {$y_0$};
    \fill (Q) circle (1.7pt);

    \draw[red!80!black, very thick, ->] (Q) -- ++(0,0.85);
    \node[red!80!black, font=\scriptsize, align=center] at (1.75,2.95)
      {move in $y$ only: $f$ unchanged to first order};

    % two y's over one x, just left of x_0; none just right of it
    \draw[gray, dashed] (1.15,0) -- (1.15,2.75);
    \fill[orange!90!black] (1.15,0.89) circle (1.6pt);
    \fill[orange!90!black] (1.15,2.21) circle (1.6pt);
    \node[orange!90!black, font=\scriptsize, align=center] at (0.62,1.55)
      {two $y$'s\\ here};
    \draw[gray, dashed] (2.15,0) -- (2.15,2.75);
    \node[orange!90!black, font=\scriptsize, align=center] at (2.75,1.15)
      {none\\ here};

    \node[font=\scriptsize, align=center] at (1.55,-0.85)
      {$D_yf(x_0,y_0)=0$\\ \textbf{not solvable for $y$}};
  \end{scope}
\end{tikzpicture}
\end{center}

% Sonnet 5 (Medium) --- FIG-W06-05, replacing the original box-and-arrow diagram.
\subsubsection*{What the conclusion gives}

The theorem does not say the whole level set is a graph --- only that it is one \emph{inside a
small enough box}. On $B_r(x_0)\times B_s(y_0)$ the two descriptions coincide exactly: the
zero set of $f$ and the graph of $g$ are the same set. Outside the box the level set may do
anything at all, including doubling back or carrying further branches, and that is precisely
why the statement is local.

\begin{center}
\begin{tikzpicture}[>=Latex, scale=1.2]
  \draw[green!60!black, thick, ->] (-0.4,0) -- (5.3,0) node[right, font=\small] {$x \in \mathbb{R}^d$};
  \draw[orange!90!black, thick, ->] (0,-0.4) -- (0,3.4) node[above, font=\small] {$y \in \mathbb{R}^{n-d}$};

  % Level set: inside the box a clean graph; outside it doubles back and branches.
  \draw[gray!55, very thick, smooth]
    plot coordinates {(0.45,2.95) (0.8,2.35) (1.1,1.98) (1.4,1.751)};
  \draw[gray!55, very thick, smooth]
    plot coordinates {(3.6,1.531) (4.05,1.33) (4.35,0.78) (4.15,0.28)};
  \node[gray!55!black, font=\scriptsize, right] at (4.45,0.80) {further branches};

  % The box, then the graph redrawn on top so it reads as the box's content.
  \draw[green!60!black, thick, fill=green!7] (1.4,0.9) rectangle (3.6,2.4);
  \draw[blue!80!black, very thick, domain=1.4:3.6, samples=60, smooth]
    plot ({\x}, {1.72 - 0.30*(\x-1.5) + 0.10*(\x-1.5)*(\x-1.5)});

  \coordinate (P) at (2.5,1.52);
  \fill (P) circle (1.7pt) node[above left, font=\scriptsize] {$(x_0,y_0)$};

  % Box extents.
  \draw[green!60!black, <->] (1.4,0.62) -- node[below, font=\scriptsize] {$B_r(x_0)$} (3.6,0.62);
  \draw[orange!90!black, <->] (1.13,0.9) -- node[left, font=\scriptsize] {$B_s(y_0)$} (1.13,2.4);
  \draw[dotted, gray] (2.5,0) -- (P);
  \node[below, font=\scriptsize] at (2.5,-0.05) {$x_0$};

  \node[blue!80!black, font=\scriptsize, right, align=left] at (3.72,2.55)
    {inside the box:\\ $f(x,y)=0 \iff y=g(x)$};
  \draw[blue!80!black, thin, ->] (3.72,2.35) -- (3.05,1.62);
\end{tikzpicture}
\end{center}

\begin{example}[The circle $S^1$]
The circle $S^1 = \{(x,y) \in \mathbb{R}^2 : x^2+y^2-1=0\}$ is the level set at $0$ of
$F(x,y) := x^2+y^2-1$, i.e., $S^1 = F^{-1}\{0\}$. Locally, there are explicit functions
\[x \mapsto \pm\sqrt{1-x^2} = y(x),\]
defined on open intervals around $x_0$ whenever $y(x_0)\neq0$. And indeed,
\[D_yF(x,y) = \frac{\partial F}{\partial y} = 2y\]
is bijective whenever $y \neq 0$, so the implicit function theorem holds there.

\begin{remark}[The circle realizes both panels]
The circle exhibits both cases of the figure above, and it is worth checking them against each
other.

At a point with $y_0 \neq 0$ --- say $(0,1)$ --- we have $D_yF = 2y_0 \neq 0$, the circle meets
that point at a non-vertical angle, and near it the circle really is the graph
$y = \sqrt{1-x^2}$. That is the left panel.

At $\bar x = 1$, where $y = 0$, we get $D_yF(1,0) = 0$. The circle is \emph{vertical} there, and
no graph $y(x)$ can exist around it: for $x$ slightly less than $1$ there are two points of the
circle above $x$, namely $\pm\sqrt{1-x^2}$, and for $x$ slightly greater than $1$ there are none.
That is the right panel exactly --- two $y$'s on one side, none on the other.

Note what has \emph{not} gone wrong at $\bar x = 1$: the circle is a perfectly smooth curve
there, with no corner and no break. What fails is only the attempt to describe it as a graph
\emph{over $x$}. Swap the roles of the two variables and everything works again, since
$D_xF(1,0) = 2 \neq 0$: near $(1,0)$ the circle is the graph $x = \sqrt{1-y^2}$ of a function
of $y$. This is the \qt{which variables get solved for is a choice} point above, in its
smallest possible instance.
\end{remark}
\end{example}

% Sonnet 5 (Medium) --- FIG-W06-06, drawn from the page image seen while transcribing
\begin{center}
\begin{tikzpicture}[>=Latex, scale=1.3]
  \draw[->] (-1.6,0) -- (1.9,0) node[right, font=\small] {$x$};
  \draw[->] (0,-1.6) -- (0,1.6) node[above, font=\small] {$y$};
  \draw[green!60!black, thick] (0,0) circle (1);

  \draw[red!80!black, dashed, thick, domain=-0.9:0.9, samples=40, smooth]
    plot ({\x}, {sqrt(1-\x*\x)}) node[right, font=\scriptsize] {$y=\sqrt{1-x^2}$};

  \coordinate (X0) at (0.4,0);
  \draw[dotted, gray] (X0) -- (0.4,0.917);
  \node[below, font=\scriptsize] at (X0) {$x_0$};

  \fill[red!80!black] (1,0) circle (1.5pt) node[below right, font=\scriptsize] {$\bar x$};
  \node[red!80!black, font=\scriptsize, below right, align=left] at (1.0,-0.5)
    {not possible\\around $y=0$};

  \draw[->, thick] (2.6,-1.0) -- (2.6,1.3) node[above, font=\small] {$\mathbb{R}$};
  \fill (2.6,0) circle (1.5pt) node[right, font=\small] {$0$};
  \draw[->, thick, blue!80!black] (1.15,0.6).. controls (2.0,1.0).. (2.6,0)
    node[midway, above, font=\small] {$F$};
\end{tikzpicture}
\end{center}

\begin{example}[An implicit quintic]
Let $G(x,y) := y^5+y^3+y+x$, and consider its level set at $0$,
$G^{-1}\{0\} = \{(x,y)\in\mathbb{R}^2 : G(x,y)=0\}$. Clearly $(0,0)\in G^{-1}\{0\}$, since
$G(0,0)=0^5+0^3+0+0=0$. Now both
\[\partial_xG(0,0)=1\neq0, \qquad \partial_yG(0,0)=\big[5y^4+3y^2+1\big]_{y=0}=1\neq0,\]
so by the implicit function theorem there are open intervals (balls) $(-r,r)$, $(-s,s)$ and
functions
\[(-r,r)\to(-s,s), \quad y\mapsto x(y) = -y^5-y^3-y,\]
\[(-s,s)\to(-r,r), \quad x\mapsto y(x),\]
satisfying
\[y^5+y^3+y+x(y)=0, \qquad y(x)^5+y(x)^3+y(x)+x=0.\]

\begin{remark}
As Galois and Abel showed, no closed-form expression for $y(x)$ in radicals exists in general,
since $y^5+y^3+y+x=0$ is a genuine quintic in $y$. Nevertheless, the implicit function theorem
guarantees that $y(x)$ exists (and is $C^\infty$, since $G$ is a polynomial), and it can be
computed and plotted numerically.
\end{remark}
\end{example}

% Sonnet 5 (Medium) --- FIG-W06-07, drawn from the page image seen while transcribing
\begin{center}
\begin{tikzpicture}[>=Latex, scale=0.62]
  \draw[->] (-11,0) -- (11,0) node[right, font=\small] {$x$};
  \draw[->] (0,-2.0) -- (0,2.0) node[above, font=\small] {$y(x)$};
  \foreach \x in {-10,-7.5,-5,-2.5,2.5,5,7.5,10}
    \draw (\x,0.08) -- (\x,-0.08) node[below, font=\tiny] {$\x$};
  \foreach \y in {-1.5,-1.0,-0.5,0.5,1.0,1.5}
    \draw (0.15,\y) -- (-0.15,\y) node[left, font=\tiny] {$\y$};

  \draw[blue!80!black, thick, variable=\t, domain=-1.42:1.42, samples=120, smooth]
    plot ({-\t*\t*\t*\t*\t - \t*\t*\t - \t}, {\t});
\end{tikzpicture}
\end{center}


% --- exercise relocated here from the dissolved sheet 7 ---

% Source: Jérôme Paschoud/Notizen/Woche 7.1 Satz über implizite Funktionen.pdf, p. 5
% Generator: Gemini 3.6 Flash (High)
\begin{exercise}[Implicitly defined functions from a system]
\label{ex:implicit_system_yz_of_x}
Show that one can solve the system of equations
\[
\begin{cases}
-2x^2 + y^2 + z^2 = 0 \\
x^2 + e^{y-1} - 2y = 0
\end{cases}
\]
for $y, z$ as functions of $x$ (i.e.\ $y(x), z(x)$) close enough to the point $(1,1,1)$. Calculate $\begin{pmatrix} y'(1) \\ z'(1) \end{pmatrix}$. (cf.\ Michaels Bsp 8.3.5)
\end{exercise}

\begin{exercisesolution}[\cref{ex:implicit_system_yz_of_x}]
Let $F(x,y,z) = \begin{pmatrix} -2x^2 + y^2 + z^2 \\ x^2 + e^{y-1} - 2y \end{pmatrix}$. We have $F(1,1,1) = \begin{pmatrix} -2+1+1 \\ 1+1-2 \end{pmatrix} = \begin{pmatrix} 0 \\ 0 \end{pmatrix}$.
The Jacobian of $F$ with respect to $(y,z)$ is
\[
D_{y,z} F(x,y,z) = \begin{pmatrix} \partial_y F_1 & \partial_z F_1 \\ \partial_y F_2 & \partial_z F_2 \end{pmatrix} = \begin{pmatrix} 2y & 2z \\ e^{y-1} - 2 & 0 \end{pmatrix}.
\]
At $(x,y,z) = (1,1,1)$, we have
\[
D_{y,z} F(1,1,1) = \begin{pmatrix} 2 & 2 \\ e^0 - 2 & 0 \end{pmatrix} = \begin{pmatrix} 2 & 2 \\ -1 & 0 \end{pmatrix}.
\]
The determinant is $\det(D_{y,z} F(1,1,1)) = 2(0) - 2(-1) = 2 \neq 0$. Thus the matrix is invertible, and by the Implicit Function Theorem, there exist functions $y(x), z(x)$ solving the system in a neighbourhood of $x=1$ such that $y(1)=1$ and $z(1)=1$.

To find the derivatives at $x=1$, we use the formula
\[
\begin{pmatrix} y'(1) \\ z'(1) \end{pmatrix} = - (D_{y,z} F(1,1,1))^{-1} D_x F(1,1,1).
\]
We calculate $D_x F(x,y,z) = \begin{pmatrix} -4x \\ 2x \end{pmatrix}$, so $D_x F(1,1,1) = \begin{pmatrix} -4 \\ 2 \end{pmatrix}$.
The inverse of $D_{y,z} F(1,1,1)$ is
\[
\begin{pmatrix} 2 & 2 \\ -1 & 0 \end{pmatrix}^{-1} = \frac{1}{2} \begin{pmatrix} 0 & -2 \\ 1 & 2 \end{pmatrix}.
\]
Then
\[
\begin{pmatrix} y'(1) \\ z'(1) \end{pmatrix} = - \frac{1}{2} \begin{pmatrix} 0 & -2 \\ 1 & 2 \end{pmatrix} \begin{pmatrix} -4 \\ 2 \end{pmatrix} = - \frac{1}{2} \begin{pmatrix} -4 \\ 0 \end{pmatrix} = \begin{pmatrix} 2 \\ 0 \end{pmatrix}.
\]
Therefore, $y'(1) = 2$ and $z'(1) = 0$.
\end{exercisesolution}

% Source: Jérôme Paschoud/Notizen/Woche 7.1 Satz über implizite Funktionen.pdf, p. 6
% Generator: Gemini 3.6 Flash (High)
\begin{exercise}[Implicit solving in different variables]
\label{ex:implicit_xyz_equation}
Consider the equation $xyz - yz^4 + x^2y = 2$ near the point $(0,-2,1)$. Can one solve this equation for:
\begin{enumerate}[label=\textbf{(\alph*)}]
  \item $x$?
  \item $y$?
  \item $z$?
\end{enumerate}
Calculate the derivative of the solution at the given point, if possible. (cf.\ Michaels Bsp 8.3.3)
\end{exercise}

\begin{exercisesolution}[\cref{ex:implicit_xyz_equation}]
Let $F(x,y,z) = xyz - yz^4 + x^2y - 2$. At the point $P = (0,-2,1)$, we have $F(0,-2,1) = 0 - (-2)(1) + 0 - 2 = 0$.
We compute the partial derivatives of $F$:
\begin{align*}
\partial_x F(x,y,z) &= yz + 2xy \\
\partial_y F(x,y,z) &= xz - z^4 + x^2 \\
\partial_z F(x,y,z) &= xy - 4yz^3
\end{align*}
Evaluating at $P = (0,-2,1)$:
\begin{align*}
\partial_x F(0,-2,1) &= (-2)(1) + 0 = -2 \neq 0 \\
\partial_y F(0,-2,1) &= 0 - 1 + 0 = -1 \neq 0 \\
\partial_z F(0,-2,1) &= 0 - 4(-2)(1) = 8 \neq 0
\end{align*}

\textbf{(a)} Since $\partial_x F(P) = -2 \neq 0$, the Implicit Function Theorem guarantees we can solve for $x = x(y,z)$ locally. The derivative (gradient) is
\[
\nabla x(y,z) \Big|_{(-2,1)} = \begin{pmatrix} \partial_y x \\ \partial_z x \end{pmatrix} = -\frac{1}{\partial_x F(P)} \begin{pmatrix} \partial_y F(P) \\ \partial_z F(P) \end{pmatrix} = -\frac{1}{-2} \begin{pmatrix} -1 \\ 8 \end{pmatrix} = \begin{pmatrix} -1/2 \\ 4 \end{pmatrix}.
\]

\textbf{(b)} Since $\partial_y F(P) = -1 \neq 0$, we can solve for $y = y(x,z)$ locally. The derivative is
\[
\nabla y(x,z) \Big|_{(0,1)} = \begin{pmatrix} \partial_x y \\ \partial_z y \end{pmatrix} = -\frac{1}{\partial_y F(P)} \begin{pmatrix} \partial_x F(P) \\ \partial_z F(P) \end{pmatrix} = -\frac{1}{-1} \begin{pmatrix} -2 \\ 8 \end{pmatrix} = \begin{pmatrix} -2 \\ 8 \end{pmatrix}.
\]

\textbf{(c)} Since $\partial_z F(P) = 8 \neq 0$, we can solve for $z = z(x,y)$ locally. The derivative is
\[
\nabla z(x,y) \Big|_{(0,-2)} = \begin{pmatrix} \partial_x z \\ \partial_y z \end{pmatrix} = -\frac{1}{\partial_z F(P)} \begin{pmatrix} \partial_x F(P) \\ \partial_y F(P) \end{pmatrix} = -\frac{1}{8} \begin{pmatrix} -2 \\ -1 \end{pmatrix} = \begin{pmatrix} 1/4 \\ 1/8 \end{pmatrix}.
\]
\end{exercisesolution}

% Originally: content/week-07.tex, \exercisesheet{7}, problem 7.5
% Source: exercises/Ex7_Analysis2_eng.pdf

\begin{exercise}[7.5 --- Implicit function II \textnormal{(important)}]
\label{ex:7.5}
Show that the system of equations
\[\begin{cases} xy^5+yu^5+zv^5=1, \\ x^5y+y^5u+z^5v=1, \end{cases}\]
is solvable for the variables $u$ and $v$ in a neighborhood of the point
$(x_0,y_0,z_0,u_0,v_0) = (0,1,1,1,0)$ and determine the derivatives $D_{(0,1,1)}u$ and
$D_{(0,1,1)}v$ of the implicitly defined functions $u=u(x,y,z)$ and $v=v(x,y,z)$.
\end{exercise}


% ============================================================
% Chapter 17 --- Submanifolds, tangent and normal spaces
% Originally: content/week-07.tex, \section{Submanifolds} (with the regular value and local
%             parametrization theorems) and \section{Tangent and normal spaces}
%             (Corsin Week 7), plus problem 7.6 from sheet 7.
% ============================================================
\chapter{Submanifolds, Tangent and Normal Spaces}
\label{ch:submanifolds}

% Transition: Claude Opus 5 (Medium)
The implicit function theorem said that a level set is locally a graph. This chapter takes that
as a definition: a submanifold of $\mathbb{R}^n$ is a set that looks, near each of its points,
like a piece of $\mathbb{R}^k$. Three descriptions turn out to be equivalent --- as a level set of
a submersion, as a local graph, and as a local parametrization --- and each is the convenient one
for a different purpose. The chapter closes with the linear algebra attached to each point: the
tangent space of directions along the submanifold, and the normal space orthogonal to it.

% Originally: content/week-07.tex, \section{Submanifolds} (Corsin Week 7)

% Source: Corsin Nick/Class Notes/Week 7.pdf, p. 2
\section{Submanifolds}
\label{sec:submanifolds}

Recall that a \newterm{$C^k$-diffeomorphism} (\cref{thm:inverse_function_theorem}) is a $C^k$
bijective function whose inverse is also $C^k$.

\begin{definition}[Submanifold]
\label{def:submanifold}
$M^m \subseteq \mathbb{R}^n$ is a $C^k$, $m$-dimensional \newterm{submanifold}
\germanfor{Untermannigfaltigkeit} of $\mathbb{R}^n$ if for any point $p \in M$ there are an
open neighbourhood $U \subseteq \mathbb{R}^n$ of $p$, an open neighbourhood $V \subseteq
\mathbb{R}^n$ of $0$, and a $C^k$-diffeomorphism $\Phi : U \to V$ with $\Phi(p) = 0$ (a
\newterm{submanifold chart}) such that
\[\Phi(U \cap M) = (\mathbb{R}^m\times\{0\}) \cap V = \{x \in V : x_{m+1}=\dots=x_n=0\}.\]
\end{definition}

% Sonnet 5 (Medium) --- FIG-W07-01
\begin{center}
\begin{tikzpicture}[>=Latex, scale=1.0]
  % Left panel: the curved submanifold M with a chart neighbourhood U around p
  \begin{scope}[shift={(0,0)}]
    \draw[blue!70!black, thick] plot[smooth] coordinates
      {(-1.6,-0.4) (-0.9,0.5) (0.0,0.1) (0.9,0.7) (1.7,-0.1)};
    \node[blue!70!black, font=\small] at (1.9,-0.4) {$M$};
    \draw[orange!90!black, thick, dashed] (0.0,0.35) ellipse (0.85 and 0.75);
    \node[orange!90!black, font=\small] at (0.0,1.25) {$U$};
    \fill (0.0,0.1) circle (1.5pt) node[below, font=\small] {$p$};
    \node[below, font=\small] at (0,-1.3) {$\mathbb{R}^n$};
  \end{scope}

  \draw[->, thick] (2.5,0) -- (3.7,0) node[midway, above, font=\small] {$\Phi$};

  % Right panel: the flattened image, R^m x {0} as a horizontal line inside V
  \begin{scope}[shift={(4.4,0)}]
    \draw[orange!90!black, thick] (-1.1,-0.9) rectangle (1.9,1.1);
    \node[orange!90!black, font=\small] at (0.4,1.35) {$V$};
    \draw[green!60!black, thick, ->] (-1.05,0) -- (1.85,0)
      node[right, font=\scriptsize] {$\mathbb{R}^m\times\{0\}$};
    \fill (0,0) circle (1.5pt) node[below, font=\small] {$0 = \Phi(p)$};
    \node[below, font=\small] at (0.4,-1.3) {$\mathbb{R}^n$};
  \end{scope}
\end{tikzpicture}
\end{center}

% Corsin's own observation; the note it used to sit in was an ainote until 2026-08-09, but by
% Test 1 in style.md it is mathematics rather than commentary about the transcription.
\begin{remark}[Why the definition is almost never used directly]
Producing such a $\Phi$ by hand is difficult even for simple $M$, so the definition is rarely
what one works with. Two equivalent descriptions do the job instead, the \textbf{implicit} one
(the regular value theorem below) and the \textbf{explicit or graphical} one (the local
parametrization theorem below). \Cref{rem:three_submanifold_descriptions} compares all three
once both are available.
\end{remark}

\subsection{The regular value theorem}

\begin{theorem}[Regular value theorem]
\label{thm:regular_value_theorem}
Suppose $U \subseteq \mathbb{R}^n$ is open, $F \in C^k(U,\mathbb{R}^\ell)$, and $y \in F(U)$. If
\[DF_x : \mathbb{R}^n \to \mathbb{R}^\ell\]
is surjective (i.e., $\Jac F(x)$ has full rank) for all $x \in F^{-1}\{y\} = \{\bar x \in U :
F(\bar x) = y\}$, then $F^{-1}\{y\}$ is a $C^k$ submanifold of dimension $n-\ell$.
\end{theorem}

% Sonnet 5 (Medium)
\begin{importantremark}[Full rank means two opposite things]
\label{rem:full_rank_two_meanings}
\qt{Full rank} is decided by the \textbf{shape} of the matrix, and the two representations of a
submanifold sit on opposite sides of it:
\begin{itemize}
  \item An \textbf{implicit} description $F : \mathbb{R}^n \to \mathbb{R}^\ell$ cuts $M$ out by
    $\ell$ equations, so $\Jac F$ is \textbf{wide} ($\ell \times n$ with $\ell \leq n$). Full
    rank $=\ell$ means $DF_x$ is \textbf{surjective}, so the $\ell$ constraints are independent,
    each one genuinely cutting a dimension. Here $\dim M = n - \ell$.
  \item A \textbf{parametrization} $f : \mathbb{R}^m \to \mathbb{R}^n$ sweeps $M$ out by $m$
    parameters, so $\Jac f$ is \textbf{tall} ($n \times m$ with $m \leq n$). Full rank $=m$
    means $Df_x$ is \textbf{injective}, so the $m$ parameters move in independent directions,
    none of them wasted. Here $\dim M = m$.
\end{itemize}
Same phrase, opposite condition. A wide matrix is never injective and a tall one is never
surjective, so there is no ambiguity in practice. But it is worth checking which side you are on
before writing down what \qt{full rank} is supposed to give you.
\end{importantremark}

% Source: Corsin Nick/Class Notes/Week 7.pdf, p. 2
\begin{example}[Circle of radius $R$]
Let $F : \mathbb{R}^2 \to \mathbb{R}$, $(x,y) \mapsto x^2+y^2$. Then
\[\Jac F(x,y) = (2x,\ 2y)\]
has full rank at every $(x,y) \in \mathbb{R}^2\setminus\{0\}$. Here $\ell = 1$, so by
\cref{rem:full_rank_two_meanings} full rank means that $DF_{(x,y)} : \mathbb{R}^2 \to \mathbb{R}$
is \textbf{surjective}, which for a single row is simply the requirement that $(2x,2y)$ be
nonzero. For $R>0$, $(0,0) \notin F^{-1}\{R^2\}$, so
\[F^{-1}\{R^2\} = \{(x,y) \in \mathbb{R}^2 : x^2+y^2=R^2\}\]
is a submanifold of dimension $2-1=1$ by \cref{thm:regular_value_theorem}. Similarly,
$\{x \in \mathbb{R}^n : |x|^2=R^2\}$ is a submanifold of dimension $n-1$ for $n \geq 1$ (the
\newterm{$n$-sphere} of radius $R$).
\end{example}

% Sonnet 5 (Medium) --- FIG-W07-02
\begin{center}
\begin{tikzpicture}[>=Latex, scale=1.1]
  \draw[->] (-1.7,0) -- (1.7,0) node[right, font=\small] {$x$};
  \draw[->] (0,-1.7) -- (0,1.7) node[above, font=\small] {$y$};
  \draw[blue!70!black, thick] (0,0) circle (1.3);
  \node[blue!70!black, font=\small] at (1.05,1.05) {$F^{-1}\{R^2\}$};
  \coordinate (Q) at (0.92,0.92);
  \fill (Q) circle (1.5pt);
  \draw[orange!90!black, thick, ->] (Q) -- ++(0.55,0.55)
    node[above right, font=\scriptsize] {$\nabla F(x,y) = (2x,2y)$};
\end{tikzpicture}
\end{center}

% Generator: Claude Opus 5 (Medium)
\begin{aiexample}[A set that is \emph{not} a submanifold, and how to prove it]
\label{ex:cross_not_submanifold}
The regular value theorem is a one-way street: if $DF_x$ is surjective everywhere on the level
set, we get a submanifold. When it is not, we get \emph{no information}. The level set may still
be a submanifold, or it may not. It is worth seeing a case where it is not.

Let $F(x,y) := xy$ and consider $C := F^{-1}\{0\}$, the union of the two coordinate axes.
Here $\Jac F(x,y) = (y,\ x)$, which is nonzero, and hence surjective, at every point of $C$
\emph{except} the origin. So, \cref{thm:regular_value_theorem} does apply on
$C \setminus \{(0,0)\}$, telling us that the four open half-axes are $1$-dimensional
submanifolds, and says nothing at the origin.

And at the origin $C$ genuinely fails to be a submanifold. The argument is a connectedness trick
worth remembering, since it is the standard way to prove a negative here.

Suppose $C$ were a $1$-dimensional submanifold. By \cref{def:submanifold} there would be a
neighbourhood $U$ of the origin and a diffeomorphism $\Phi : U \to V$ carrying $C \cap U$ onto a
piece of a \emph{line}, $(\mathbb{R}\times\{0\})\cap V$. Now delete the centre point from each
side and count connected components (\cref{def:connectedness}):
\begin{itemize}
  \item Removing a point from a small piece of a line leaves \textbf{two} pieces.
  \item Removing the origin from $C \cap U$ leaves \textbf{four} pieces, one along each
    half-axis.
\end{itemize}
But $\Phi$ is a homeomorphism, and homeomorphisms preserve the number of connected components
(\cref{thm:continuity_preserves_connected}, applied to $\Phi$ and to $\Phi^{-1}$). Four cannot
equal two, so no such $\Phi$ exists and $C$ is not a submanifold.

\begin{remark}
The moral is that submanifolds are forbidden from having \emph{crossings, corners or cusps}: near
each of its points a submanifold must look like a flat piece of $\mathbb{R}^m$, and a crossing
does not. Compare \cref{ex:constraint_qualification_fails}, where the constraint set
$\{y^2=x^3\}$ had a cusp at exactly the point where $\nabla g$ vanished. That is the same
phenomenon, seen through the Lagrange-multiplier lens back in \cref{ch:lagrange}.
\end{remark}
\end{aiexample}

% Source: Corsin Nick/Class Notes/Week 7.pdf, p. 4
\subsubsection{Connection to Lagrange multipliers}

\begin{remark}[Constraint sets are submanifolds]
Recall the Lagrangian from \cref{prop:constrained_optimization}: when optimizing $f|_M$ for
$f \in C^k(\mathbb{R}^n,\mathbb{R})$, where $g_1,\dots,g_k : U \to \mathbb{R}$,
\[M := \{x \in \mathbb{R}^n : g_1(x)=\dots=g_k(x)=0\}, \qquad (\nabla g_1,\dots,\nabla g_k) \text{ linearly independent for all } x \in M\]
(hence $k \leq n$). This means precisely that, for $g : \mathbb{R}^n \to \mathbb{R}^k$,
$x \mapsto (g_1(x),\dots,g_k(x))$, the Jacobian
\[\Jac g(x) = \begin{pmatrix} \nabla g_1(x) \longrightarrow \\ \vdots \\ \nabla g_k(x) \longrightarrow \end{pmatrix}\]
has full rank. Therefore, $M = g^{-1}\{0\}$ is an $(n-k)$-dimensional submanifold by the regular
value theorem: the constraint sets of the Lagrange-multiplier problems in \cref{ch:lagrange}
were submanifolds all
along.
\end{remark}

% Source: Corsin Nick/Class Notes/Week 7.pdf, p. 5
\subsection{The local parametrization theorem}

\begin{theorem}[Local parametrization]
\label{thm:local_parametrization}
Let $M \subseteq \mathbb{R}^n$. Then $M$ is a $C^k$ submanifold with $\dim M = m$ if and only if
for every $p \in M$ there exist an open neighbourhood $V$ of $p$, an open set $U \subseteq
\mathbb{R}^m$, and a $C^k$ map $f : U \to V \cap M$ such that
\begin{itemize}
  \item $Df_x : \mathbb{R}^m \to \mathbb{R}^n$ is injective for all $x \in U$;
  \item $f$ is a homeomorphism (bijective, continuous, with continuous inverse $f^{-1}$).
\end{itemize}
\end{theorem}

% Generator: Claude Opus 5 (Medium)
\begin{remark}[Three descriptions, and which one to reach for]
\label{rem:three_submanifold_descriptions}
There are now three ways to certify that $M$ is an $m$-dimensional submanifold, and they are
equivalent. They are not equally convenient, though, and choosing badly is what makes these
exercises feel hard.

\begin{itemize}
  \item \textbf{Chart} (\cref{def:submanifold}) is the definition itself. Straighten $M$ out
    with a diffeomorphism, so that the first $m$ coordinates of the chart carry the submanifold
    and the rest vanish. It is almost never used directly on a concrete $M$, since producing
    $\Phi$ by hand is painful. Its role is theoretical, and it is what the other two are proved
    from.
  \item \textbf{Implicit, via the regular value theorem} (\cref{thm:regular_value_theorem}) is
    the one to reach for when $M$ arrives as \emph{equations}, such as $\{x^2+y^2=R^2\}$ or
    $\{g_1=\dots=g_k=0\}$. It rests on the same idea as the chart definition, with the
    diffeomorphism produced implicitly rather than by hand. Check a rank condition on $\Jac F$
    and you are done, with no parametrization ever written. This is normally the cheapest route.
  \item \textbf{Parametrization} (\cref{thm:local_parametrization}) is the one for an $M$ that
    arrives as a \emph{picture you can sweep out}, such as a graph, a surface of revolution, or
    a curve given by a formula in a parameter. It is often the most intuitive, since it builds
    the submanifold out of parameters explicitly. It is also the one to use when the tangent
    space is wanted next, because $\partial_1f,\dots,\partial_mf$ hand you a basis directly
    (\cref{def:tangent_space}).
\end{itemize}

The rough rule: \textbf{equations $\to$ regular value theorem; a formula with parameters $\to$
local parametrization.} A graph $\{(x,f(x))\}$ satisfies both, which is why graphs are the
easiest case of all.

Two cautions carried over from \cref{rem:full_rank_two_meanings}. \qt{Full rank} means
\emph{surjective} in the first case and \emph{injective} in the second, and the shape of the
matrix decides which. And neither theorem yields a negative. If the rank condition fails you
have learned nothing, and must argue separately (see \cref{ex:cross_not_submanifold}).
\end{remark}

\begin{example}[The $n$-parabola]
Let $P := \{(x,t) \in \mathbb{R}^n\times\mathbb{R} : t=|x|^2\}$. One could use the regular value
theorem, but instead take
\[f : \mathbb{R}^n \to \mathbb{R}^n\times\mathbb{R}, \qquad x \mapsto (x,|x|^2).\]
\begin{enumerate}[label=\textbf{(\alph*)}]
  \item The Jacobian matrix has full rank $n$ at every $x$:
\[ \Jac f(x) = \begin{pmatrix} \id_{n\times n} \\ 2x\transp \end{pmatrix}. \]
    This one is tall, so by \cref{rem:full_rank_two_meanings} full rank means that $Df_x$ is
    \textbf{injective}, which is clear from the $\id_{n\times n}$ block alone.
  \item $f$ is bijective onto $P$ and continuous, with inverse $\pi : f(U) \to \mathbb{R}^n$,
    $(x,t) \mapsto x$ (the canonical projection), which is also continuous.
  \item $f(\mathbb{R}^n) = P$.
\end{enumerate}
So, $P$ is an $n$-dimensional submanifold by \cref{thm:local_parametrization}.

\begin{remark}
The graphical representation of a submanifold (as the graph of a function, here $x \mapsto |x|^2$)
is therefore a common, and equivalent, special case of a local parametrization: it is always
injective on its domain by construction, and its inverse is simply the projection forgetting the
last coordinate.
\end{remark}
\end{example}

% Sonnet 5 (Medium) --- FIG-W07-03
\begin{center}
\begin{tikzpicture}[>=Latex, scale=1.1]
  \draw[->] (-1.9,0) -- (1.9,0) node[right, font=\small] {$x$};
  \draw[->] (0,-0.3) -- (0,2.7) node[above, font=\small] {$t$};
  \draw[blue!70!black, thick, domain=-1.6:1.6, samples=60, smooth]
    plot ({\x}, {\x*\x}) node[above right, font=\small] {$P = \{(x,t) : t=|x|^2\}$};
\end{tikzpicture}
\end{center}

% Source: Corsin Nick/Class Notes/Week 7.pdf, p. 6
\subsubsection{An alternative optimization technique}

\begin{remark}[Parametrized optimization, revisiting \cref{ex:extrema_circle_disk}]
Local parametrizations give a second route to constrained-extremum problems, alongside Lagrange
multipliers: parametrize the constraint set and optimize the resulting unconstrained function of
fewer variables. Corsin illustrates this by revisiting $f(x,y) := 2x^2+y^2-x$ on
$S^1 = \{(x,y)\in\mathbb{R}^2 : x^2+y^2=1\}$, which is exactly \cref{ex:extrema_circle_disk}\textbf{(a)}, solved
there via Lagrange multipliers.

Parametrize $S^1$ by $(x(\theta),y(\theta)) := (\cos\theta,\sin\theta)$. Then
\[f(\theta) := f(x(\theta),y(\theta)) = 2\cos^2\theta+\sin^2\theta-\cos\theta = \cos^2\theta-\cos\theta+1,\]
using $\sin^2\theta = 1-\cos^2\theta$. Differentiating,
\[f'(\theta) = -2\cos\theta\sin\theta+\sin\theta \overset{!}{=} 0.\]
\begin{enumerate}[label=\textbf{(\alph*)}]
  \item $\sin\theta = 0 \implies \theta_1 \in \{0,\pi\}$.
  \item $\sin\theta \neq 0 \implies -2\cos\theta+1=0 \implies \cos\theta = \tfrac12 \implies
    \theta_2 \in \{\tfrac{\pi}{3},\tfrac{5\pi}{3}\}$.
\end{enumerate}
So, the extrema are at $(\cos\theta_1,\sin\theta_1) = (\pm1,0)$ and
$(\cos\theta_2,\sin\theta_2) = \big(\tfrac12,\pm\tfrac{\sqrt3}{2}\big)$. These are exactly the
four points found via Lagrange multipliers in the solution to \cref{ex:extrema_circle_disk}\textbf{(a)}, so the
two methods agree.
\end{remark}

% Source: Jérôme Paschoud/Notizen/Woche 7.2 Untermannigfaltigkeiten.pdf, p. 4
% Generator: Gemini 3.6 Flash (High)
\begin{exercise}[Dimension of a submanifold]
\label{ex:dimension_submanifold_system}
Show that $M := \{(x,y,z) \in \mathbb{R}^3 \mid 2x+z=1, x^2+y^2-z^2=1\}$ is a submanifold of $\mathbb{R}^3$. What is its dimension? (cf.\ Michaels Bsp 9.1.2 d)
\end{exercise}

\begin{exercisesolution}[\cref{ex:dimension_submanifold_system}]
Let $F(x,y,z) = \begin{pmatrix} 2x+z-1 \\ x^2+y^2-z^2-1 \end{pmatrix}$. Then $M = F^{-1}\{0\}$.
We compute the Jacobian matrix of $F$:
\[ \Jac F(x,y,z) = \begin{pmatrix} 2 & 0 & 1 \\ 2x & 2y & -2z \end{pmatrix}. \]
We need to check that $\rank(\Jac F(x,y,z)) = 2$ for all $(x,y,z) \in M$.
The rank is less than 2 if and only if all $2 \times 2$ subdeterminants are zero:
\begin{align*}
\det \begin{pmatrix} 2 & 0 \\ 2x & 2y \end{pmatrix} &= 4y = 0 \implies y = 0 \\
\det \begin{pmatrix} 2 & 1 \\ 2x & -2z \end{pmatrix} &= -4z - 2x = 0 \implies x = -2z \\
\det \begin{pmatrix} 0 & 1 \\ 2y & -2z \end{pmatrix} &= -2y = 0 \implies y = 0.
\end{align*}
So, the rank is not maximal only when $y=0$ and $x=-2z$.
Let's check if such a point can belong to $M$. Substituting $x=-2z$ into the first equation of $M$ gives $2(-2z) + z = -3z = 1 \implies z = -1/3$.
Then $x = 2/3$. Thus the point is $(2/3, 0, -1/3)$.
Now check the second equation of $M$: $(2/3)^2 + 0^2 - (-1/3)^2 = 4/9 - 1/9 = 3/9 = 1/3 \neq 1$.
Therefore, this point is not in $M$.
Since $\Jac F(x,y,z)$ has maximal rank 2 for all $(x,y,z) \in M$, $M$ is a submanifold of $\mathbb{R}^3$.
Its dimension is $n - k = 3 - 2 = 1$.
\end{exercisesolution}

% Source: Jérôme Paschoud/Notizen/Woche 7.2 Untermannigfaltigkeiten.pdf, p. 4
% Generator: Gemini 3.6 Flash (High)
\begin{exercise}[Submanifold of matrices]
\label{ex:submanifold_matrices_rank1}
Show that $M := \{ A \in \operatorname{Mat}_{2\times 2}(\mathbb{R}) \mid \rank(A) = 1 \}$ is a submanifold of $\operatorname{Mat}_{2\times 2}(\mathbb{R})$. What is the dimension of $M$?
\end{exercise}

\begin{exercisesolution}[\cref{ex:submanifold_matrices_rank1}]
Let $A = \begin{pmatrix} a & b \\ c & d \end{pmatrix}$. The condition $\rank(A) \leq 1$ is equivalent to $\det(A) = ad - bc = 0$.
The condition $\rank(A) = 1$ means $\det(A) = 0$ but $A \neq 0$.
So, $M$ is the zero set of $F : \mathbb{R}^4 \setminus \{0\} \to \mathbb{R}$, $F(a,b,c,d) = ad - bc$.
The space $\operatorname{Mat}_{2\times 2}(\mathbb{R})$ is isomorphic to $\mathbb{R}^4$.
The Jacobian of $F$ is
\[ \Jac F(a,b,c,d) = \begin{pmatrix} \partial_a F & \partial_b F & \partial_c F & \partial_d F \end{pmatrix} = \begin{pmatrix} d & -c & -a & b \end{pmatrix}. \]
For $\Jac F$ to not have maximal rank (which is 1), we must have $d = 0$, $c = 0$, $a = 0$, and $b = 0$, which means $A = 0$.
But the zero matrix is excluded from $M$ because its rank is 0, not 1.
Thus $\Jac F(A) \neq 0$ for all $A \in M$. By the regular value theorem, $M$ is a submanifold.
The ambient dimension is $4$ and the number of independent constraints is $1$.
So, the dimension of $M$ is $4 - 1 = 3$.
\end{exercisesolution}

% Source: old_exams/HS18.pdf (Analysis I & II, winter 2018/2019, Prof. Manfred Einsiedler),
% Aufgabe 7, p. 2
% Extractor: Claude Opus 5 (High)
% Kept as a worked example: it is the half of thm:local_parametrization that the theorem's own
% statement hides, namely that the homeomorphism clause is free once one shrinks the domain, and
% that is a piece of theory rather than a drill.
\begin{example}[An immersion is locally a parametrization]
\label{ex:immersion_locally_submanifold}
\Cref{thm:local_parametrization} asks a candidate parametrization $f$ for two things: that $Df_x$
be injective everywhere, and that $f$ be a homeomorphism onto its image. The first is a condition
on a matrix and is easy to check; the second is a global condition and is the one that fails in
practice, as the figure-eight and the lollipop show. The following says that, locally, the second
condition comes for free.

\textbf{Claim.} Let $U \subseteq \mathbb{R}^n$ be open and let $f : U \to \mathbb{R}^{n+1}$ be
smooth with $D_{x_0}f$ injective for every $x_0 \in U$ (such an $f$ is called an
\newterm{immersion} \germanfor{Immersion}). Then every $x_0 \in U$ has an open neighbourhood
$V \subseteq U$ such that $f(V)$ is an $n$-dimensional submanifold of $\mathbb{R}^{n+1}$.

Fix $x_0 \in U$. The matrix $\Jac f(x_0)$ has $n+1$ rows and $n$ columns and is injective, so it
has rank $n$: some $n$ of its rows are linearly independent. Permuting the coordinates of
$\mathbb{R}^{n+1}$ is a linear diffeomorphism of the ambient space, and it carries submanifolds to
submanifolds, so we may assume those are the first $n$ rows. Write accordingly
\[f = \begin{pmatrix} g \\ h \end{pmatrix}, \qquad g : U \to \mathbb{R}^n,\quad
h : U \to \mathbb{R},\]
so that $\Jac g(x_0) \in \mathbb{R}^{n\times n}$ is invertible.

By the inverse function theorem (\cref{thm:inverse_function_theorem}) there is an open
$V \ni x_0$, $V \subseteq U$, such that $g|_V : V \to W := g(V)$ is a diffeomorphism onto an open
set $W \subseteq \mathbb{R}^n$. Let $\varphi := (g|_V)^{-1} : W \to V$ and set
\[\psi := h \circ \varphi : W \to \mathbb{R},\]
which is smooth as a composition of smooth maps. For $x \in V$ we have $g(x) \in W$ and
\[f(x) = \big(g(x),\, h(x)\big) = \big(g(x),\, \psi(g(x))\big),\]
using $\varphi(g(x)) = x$. As $x$ runs over $V$, the point $g(x)$ runs over all of $W$. Therefore
\[f(V) = \{(w,\psi(w)) \mid w \in W\}\]
is exactly the graph of the smooth function $\psi$ over the open set $W$, and a graph is an
$n$-dimensional submanifold of $\mathbb{R}^{n+1}$. To see that last point, apply the regular value
theorem to $(w,s) \mapsto s - \psi(w)$, whose Jacobian $(-\nabla\psi(w)\transp,\ 1)$ never
vanishes.

\begin{remark}[What is local here, and what is not]
The injectivity of $D_{x_0}f$ is a pointwise, linear-algebraic condition, and yet it forces $f$
itself to be injective near $x_0$: on the shrunken $V$ the map $f$ is injective because its first
block $g$ already is. That is the whole trick, and it is why an immersion cannot fail to be a
parametrization for a local reason.

What it cannot repair is a global collision. The figure-eight curve is an immersion of an open
interval whose image meets itself at the crossing point, and no shrinking of the domain changes
that the \emph{image} is not a submanifold there. The claim promises a neighbourhood $V$ in the
\emph{domain} and says nothing about a neighbourhood of $f(x_0)$ in $\mathbb{R}^{n+1}$; the two
are different requests, and only the first is granted. Compare
\cref{ex:lollipop_tangent_field}, where the failure is of exactly this global kind.
\end{remark}
\exinfo{This example is taken from the Analysis I \& II examination of winter 2018/2019
(Prof.\ Manfred Einsiedler), Problem 7, where it is set as a proof exercise. The framing against
\cref{thm:local_parametrization} and the closing remark are added here.}
\end{example}

% Source: old_exams/HS19.pdf (Analysis I & II, 22 January 2020, Prof. Peter S. Jossen),
% Teil B, Aufgabe 4, p. 21
% Extractor: Claude Opus 5 (High)
% Filed next to ex:submanifold_matrices_rank1 deliberately: both live in Mat_2(R) = R^4, and the
% contrast is the point. There the regular value theorem applies as it stands; here the naive
% application fails, because F lands in the symmetric matrices and so is nowhere submersive.
% Parts (e) and (f) of the original invoke the constant rank theorem, which this document does not
% state; the exercise is re-routed through the regular value theorem instead. See the ainote below.
\begin{exercise}[The orthogonal group as a submanifold]
\label{ex:orthogonal_group_submanifold}
Let $\operatorname{Mat}_2(\mathbb{R})$ denote the $2\times 2$ real matrices, identified with
$\mathbb{R}^4$ by listing the entries $a,b,c,d$ of
$\left(\begin{smallmatrix} a & b \\ c & d\end{smallmatrix}\right)$ as a column. Define
\[F : \operatorname{Mat}_2(\mathbb{R}) \to \operatorname{Mat}_2(\mathbb{R}), \qquad
F(A) := AA\transp,\]
and let $\Orth_2(\mathbb{R}) := \{A \in \operatorname{Mat}_2(\mathbb{R}) \mid F(A) = \id\}$.
\begin{enumerate}[label=\textbf{(\alph*)}]
  \item Write $F$ out in the coordinates $a,b,c,d$ and explain why it is smooth.
  \item Show that $F(A)$ is symmetric for every $A$, so that $F$ takes values in the
    three-dimensional subspace $\operatorname{Sym}_2(\mathbb{R}) \subseteq
    \operatorname{Mat}_2(\mathbb{R})$ of symmetric matrices.
  \item Compute the Jacobian matrix of $F$ at $A = \left(\begin{smallmatrix} a & b \\ c &
    d\end{smallmatrix}\right)$. It is a $4\times 4$ matrix.
  \item Show that this Jacobian has rank $3$ at every $A \in \Orth_2(\mathbb{R})$. Why does this
    already rule out applying the regular value theorem to $F$ as it stands?
  \item Regard $F$ instead as a map $\widetilde F : \mathbb{R}^4 \to \mathbb{R}^3$ into
    $\operatorname{Sym}_2(\mathbb{R}) \cong \mathbb{R}^3$ and deduce from
    \cref{thm:regular_value_theorem} that $\Orth_2(\mathbb{R})$ is a submanifold of
    $\operatorname{Mat}_2(\mathbb{R})$. What is its dimension?
  \item Determine the tangent space $T_{\id}\Orth_2(\mathbb{R})$ from the Jacobian computed in
    \textbf{(c)}, and give a basis for it.
\end{enumerate}
\exinfo{This exercise is taken from the Analysis I \& II examination of 22 January 2020
(Prof.\ Peter S.\ Jossen), Part B, Problem 4. The original reaches the submanifold property
through the constant rank theorem, which it also asks the candidate to state; parts \textbf{(d)}
and \textbf{(e)} are rephrased here to use the regular value theorem instead.}
\exsol{sol:orthogonal_group_submanifold}
\end{exercise}

\begin{ainote}
The re-routing in \textbf{(e)} is not a simplification of Jossen's problem, it is the reason the
problem is interesting. The constant rank theorem states that a $C^k$ map whose differential has
the \emph{same} rank $r$ throughout a neighbourhood looks, in suitable charts, like the linear
projection $(x_1,\dots,x_n) \mapsto (x_1,\dots,x_r,0,\dots,0)$; its level sets are then
submanifolds of dimension $n-r$. This document does not state that theorem, and it does not need
to here: restricting the codomain of $F$ to the subspace $\operatorname{Sym}_2(\mathbb{R})$ that
its image already lies in turns the constant rank $3$ into a \emph{full} rank $3$, and the regular
value theorem takes over. That manoeuvre, shrinking the target until the map becomes submersive,
is the standard way to handle matrix groups without the heavier theorem, and it is worth
knowing because $\Orth_n(\mathbb{R})$, $\operatorname{SL}_n(\mathbb{R})$ and their relatives all
fall to it.
\end{ainote}

% Transition: Claude Opus 5 (High)
Three equivalent descriptions are on the table, each characterising the same class of sets. What
they do not yet come with is practice, both in choosing which description to write down for a
surface handed to you and, harder, in arguing that a set is \emph{not} a submanifold at all.
None of the three theorems helps with the second, since each of them only ever yields a positive.

% Originally: new section (no predecessor in the week-based skeleton).
% Quelle: Linus Lüchinger/Slides/Slides-03-30.pdf, slides 12--13 (examples and
%         non-examples), 16, 23, 26 (one surface in three representations);
%         Linus Lüchinger/Slides/Slides-04-30.pdf, slide 7 (tangent-field test).

\section{Recognising a submanifold in practice}
\label{sec:recognising_submanifolds}

% Supplement: Linus Lüchinger/Slides/Slides-03-30.pdf, slides 16, 23, 26
% Why this section exists at all: rem:three_submanifold_descriptions says which description to
% reach for, but Corsin's notes never run one surface through all three, and give no worked
% non-example beyond the crossing in ex:cross_not_submanifold. Linus spends most of one session
% on exactly those two gaps, and this section follows his treatment. (Was an ainote until
% 2026-08-09; it is source-mining logistics, which Test 2 in style.md sends to a % comment.)

\subsection{One surface, three descriptions}

% Supplement: Linus Lüchinger/Slides/Slides-03-30.pdf, slides 16, 23, 26
The three characterisations of a submanifold are easiest to keep apart when they are applied to
the same object at the same point. Take the open upper half-sphere
\[\mathbb{S}^2_+ := \{(x,y,z) \in \mathbb{R}^3 : x^2+y^2+z^2 = 1,\ z > 0\},\]
a $2$-dimensional submanifold of $\mathbb{R}^3$, and throughout fix the point
\[p := \Big(0,\ \tfrac{1}{\sqrt2},\ \tfrac{1}{\sqrt2}\Big) \in \mathbb{S}^2_+.\]

\begin{example}[The upper half-sphere, three ways]
\label{ex:half_sphere_three_ways}
\begin{enumerate}[label=\textbf{(\alph*)}]
  \item \textbf{Implicit.} Put $F(x,y,z) := x^2+y^2+z^2-1$ on the open set
    $U := \{z>0\}$. Then $\mathbb{S}^2_+ = F^{-1}\{0\}$ and
    \[\nabla F(x,y,z) = (2x,\ 2y,\ 2z), \qquad \nabla F(p) = (0,\ \sqrt2,\ \sqrt2) \neq 0,\]
    so $\Jac F$ has full rank $\ell = 1$ everywhere on $U$ (indeed $2z > 0$ already forces the
    row to be nonzero). \cref{thm:regular_value_theorem} gives a submanifold of dimension
    $3-1 = 2$.
  \item \textbf{Graph.} Over the open unit disc $V := \{(x,y) : x^2+y^2 < 1\}$ put
    $f(x,y) := \sqrt{1-x^2-y^2}$. Then $\mathbb{S}^2_+ = \{(x,y,f(x,y)) : (x,y) \in V\}$ is the
    graph of $f$, and $f(0,1/\sqrt2) = \sqrt{1-1/2} = 1/\sqrt2$ recovers the third coordinate of
    $p$. Here the half-sphere is globally a graph, which is exactly why the \emph{upper half} is
    a friendlier example than the whole sphere.
  \item \textbf{Parametrization.} With $W := (0,\pi/2) \times (0,2\pi)$ and
    \[\phi(\theta,\psi) := (\sin\theta\cos\psi,\ \sin\theta\sin\psi,\ \cos\theta),\]
    we get $p = \phi(\pi/4,\ \pi/2)$. The Jacobian
    \[\Jac\phi(\theta,\psi) = \begin{pmatrix}
        \cos\theta\cos\psi & -\sin\theta\sin\psi \\
        \cos\theta\sin\psi & \sin\theta\cos\psi \\
        -\sin\theta & 0
      \end{pmatrix}\]
    is tall ($3\times2$), and full rank $=2$ means $D\phi$ is \emph{injective}
    (\cref{rem:full_rank_two_meanings}): the bottom row forces the first column to be nonzero
    whenever $\sin\theta \neq 0$, and the second column is then independent of it because
    $\sin\theta \neq 0$ makes $(-\sin\theta\sin\psi,\ \sin\theta\cos\psi)$ nonzero. On $W$ we
    have $\theta \in (0,\pi/2)$, so $\sin\theta > 0$ throughout.
\end{enumerate}
\end{example}

% Generator: Claude Opus 5 (High) --- FIG-17-04
% Geometry check. Screen projection of a 3D point (X,Y,Z) is (Y - 0.5X, Z - 0.3X), then
% scaled by the sphere radius R = 1.25. The marked point is p = (0, 1/sqrt2, 1/sqrt2),
% which has X = 0, so it projects to R*(0.7071, 0.7071) = (0.884, 0.884) -- a point at
% screen distance sqrt(0.884^2 + 0.884^2) = 1.250 = R from the origin, i.e. exactly on the
% dome's silhouette circle, at 45 degrees. That is why the dot sits ON the arc in all three
% panels, and why the gradient arrow leaves it along the outward radial direction (45 deg):
% nabla F(p) = (0, sqrt2, sqrt2) is radial, and radial directions in the y-z plane are
% undistorted by this projection.
\begin{center}
\begin{tikzpicture}[>=Latex, scale=0.95,
    dome/.style={blue!65!black, thick},
    equator/.style={blue!45!black, thin},
    panellabel/.style={font=\small, align=center}]

  \def\R{1.25}
  % All three captions sit on the SAME baseline y = -2.75, below every drawn element.
  \def\capy{-2.75}
  % --- reusable dome: outline arc + equator ellipse (back dashed, front solid) ---
  \newcommand{\halfsphere}{%
    \draw[dome] (-\R,0) arc[start angle=180, end angle=0, radius=\R];
    \draw[equator, dashed] (-\R,0) arc[start angle=180, end angle=0,
        x radius=\R, y radius={0.40*\R}];
    \draw[equator] (\R,0) arc[start angle=0, end angle=-180,
        x radius=\R, y radius={0.40*\R}];
  }

  % ============ Panel (a): implicit ============
  \begin{scope}[shift={(0,0)}]
    \halfsphere
    \coordinate (pa) at (0.884,0.884);
    \fill[black] (pa) circle (1.6pt);
    \node[font=\scriptsize, below left=-1pt and 0pt] at (pa) {$p$};
    \draw[orange!90!black, very thick, ->] (pa) -- ++(0.60,0.60);
    \node[orange!90!black, font=\scriptsize, above] at (1.30,1.46) {$\nabla F(p)$};
    \node[panellabel] at (0,\capy) {\textbf{(a)} implicit\\[1pt]
      {\footnotesize$\mathbb{S}^2_+ = F^{-1}\{0\}$}};
  \end{scope}

  % ============ Panel (b): graph ============
  \begin{scope}[shift={(4.6,0)}]
    \halfsphere
    \coordinate (pb) at (0.884,0.884);
    % The disc V, drawn detached below the dome so the drop line stays readable.
    \draw[green!45!black, thick, fill=green!12] (0,-1.30) ellipse ({\R} and {0.40*\R});
    \node[green!45!black, font=\scriptsize] at (-1.62,-1.30) {$V$};
    \fill[black] (pb) circle (1.6pt);
    \node[font=\scriptsize, below left=-1pt and 0pt] at (pb) {$p$};
    % Drop line from p to its shadow. p has (x,y) = (0, 1/sqrt2), whose screen abscissa is
    % R*0.7071 = 0.884 -- strictly inside the disc, since |(0,1/sqrt2)| = 0.707 < 1.
    \draw[gray, dashed] (0.884,0.884) -- (0.884,-1.30);
    \fill[green!45!black] (0.884,-1.30) circle (1.4pt);
    \node[font=\scriptsize, below right=-1pt and -2pt] at (0.884,-1.42) {$(x,y)$};
    \node[panellabel] at (0,\capy) {\textbf{(b)} graph\\[1pt]
      {\footnotesize$z = \sqrt{1-x^2-y^2}$}};
  \end{scope}

  % ============ Panel (c): parametrization ============
  \begin{scope}[shift={(9.2,0)}]
    \halfsphere
    \coordinate (pc) at (0.884,0.884);
    \fill[black] (pc) circle (1.6pt);
    \node[font=\scriptsize, below left=-1pt and 0pt] at (pc) {$p$};
    % Parameter rectangle W, drawn detached below the dome.
    \draw[orange!90!black, thick, fill=orange!12] (-1.15,-1.75) rectangle (0.15,-0.95);
    \node[orange!90!black, font=\scriptsize] at (-1.50,-1.35) {$W$};
    \fill[orange!90!black] (-0.50,-1.35) circle (1.4pt);
    \node[font=\scriptsize, below right=-1pt and -4pt] at (-0.50,-1.86)
      {$(\tfrac{\pi}{4},\tfrac{\pi}{2})$};
    \draw[->, thick, orange!90!black] (0.30,-1.15) to[out=15, in=255] (0.82,0.70);
    \node[orange!90!black, font=\scriptsize] at (1.15,-0.45) {$\phi$};
    \node[panellabel] at (0,\capy) {\textbf{(c)} parametrization\\[1pt]
      {\footnotesize$p = \phi(\tfrac{\pi}{4},\tfrac{\pi}{2})$}};
  \end{scope}
\end{tikzpicture}
\end{center}

% Generator: Claude Opus 5 (High)
\begin{remark}[Why all three theorems are stated \emph{locally}]
\label{rem:why_submanifolds_are_local}
The half-sphere is unusually obliging: one graph and one parametrization cover all of it. That is
not the general situation, and each description fails to globalise for its own reason.
\begin{itemize}
  \item \textbf{A parametrization need not be global.} The circle $\mathbb{S}^1$ is a
    $1$-dimensional submanifold, yet no single parametrization covers it. Suppose
    $\gamma : I \to \mathbb{S}^1$ were a homeomorphism from an open interval, as
    \cref{thm:local_parametrization} demands. Delete one point from each side: $I \setminus \{t\}$
    falls into \emph{two} pieces, while $\mathbb{S}^1 \setminus \{\gamma(t)\}$ is still an arc, so
    \emph{one} piece. A homeomorphism preserves the number of connected components
    (\cref{thm:continuity_preserves_connected}), so no such $\gamma$ exists. That is the counting
    argument of \cref{ex:cross_not_submanifold}, run in the opposite direction. Two overlapping
    arcs do cover the circle, and two is the minimum.
  \item \textbf{An implicit description need not be global either.} The Möbius strip is a
    $2$-dimensional submanifold of $\mathbb{R}^3$, but it is not $g^{-1}\{0\}$ for any single
    $g$ with $\nabla g \neq 0$: such a $g$ would supply a globally consistent normal direction
    $\nabla g/\lvert\nabla g\rvert$, and the Möbius strip has none
    (\cref{sec:orientable_manifolds}).
\end{itemize}
Requiring the descriptions only near each point costs nothing, since being a submanifold is
itself a local condition, and it buys all the examples that matter.
\end{remark}

\subsection{A gallery of non-examples}

% Supplement: Linus Lüchinger/Slides/Slides-03-30.pdf, slides 12--13
Linus prefaces the formal definitions with a list of sets that are \emph{not} submanifolds, which
is worth reproducing because each entry fails for a different reason. Being a $d$-dimensional
submanifold means looking, near every one of its points, like a flat piece of $\mathbb{R}^d$, and
there are several distinct ways to violate that.

\begin{itemize}
  \item \textbf{The figure eight} $\{(x,y) : y^2 = x^2 - x^4\}$ has a \emph{crossing}. Away from
    the origin it is a perfectly good curve. At the origin four branches meet. This is the same
    failure as \cref{ex:cross_not_submanifold}, and the same connectedness count settles it.
  \item \textbf{The lollipop} is a circle with a segment attached, or any \qt{line sticking out
    of a sphere}. Here the trouble is not a crossing but a \emph{change of dimension} at the
    junction: near the attachment point the set is neither $1$- nor $2$-dimensional.
  \item \textbf{The closed square} $[0,1]^2 \subseteq \mathbb{R}^2$ has a \emph{boundary}. Its
    interior is a fine $2$-dimensional submanifold, but no neighbourhood of a corner or edge
    point looks like an open piece of $\mathbb{R}^2$, because points of $[0,1]^2$ near an edge
    fill only a half-plane. (Sets like this are the subject of
    \cref{sec:submanifolds_with_boundary}; a submanifold \emph{with boundary} is a genuinely
    weaker notion, not a submanifold that happens to have edges.)
  \item \textbf{$\mathbb{Q}^n \subseteq \mathbb{R}^n$} has \emph{no dimension at all}. It is not
    $0$-dimensional, since a $0$-dimensional submanifold is a set of isolated points and no
    rational point is isolated. And it is not $n$-dimensional, since it contains no open ball.
\end{itemize}

% Generator: Claude Opus 5 (High) --- FIG-17-05
% Geometry check. Panel (a) is the lemniscate y^2 = x^2 - x^4, parametrized as
% (cos t, sin t cos t) for t in [0, 2pi); at t = pi/2 and t = 3pi/2 the curve passes through
% (0,0), which is why the self-crossing is at the origin and nowhere else. Panel (b): the
% circle has centre (-0.45,0) radius 0.62, so its rightmost point is (0.17, 0); the stick
% runs from (0.17,0) to (1.15,0), i.e. it starts exactly ON the circle -- the junction dot
% is at (0.17,0). Panel (c): the square is [0,1]^2 scaled to side 1.3; the marked edge point
% (0.65, 0) is the midpoint of the bottom edge, and the half-disc above it is the part of a
% small ball that actually lies in the square.
\begin{center}
\begin{tikzpicture}[>=Latex, scale=1.15,
    bad/.style={red!80!black},
    leader/.style={gray!70, very thin}]
  % All captions share the baseline y = -1.70; all annotations sit at y = -1.10,
  % clear of every drawn curve (the lemniscate reaches only |y| = 0.625).
  \def\capy{-1.70}
  \def\anny{-1.10}

  % ---- (a) figure eight: (x,y) = 1.25*(cos t, sin t cos t), which satisfies
  %      y^2 = x^2 - x^4 after undoing the scale. It passes through the origin
  %      exactly twice, at t = pi/2 and t = 3pi/2.
  \begin{scope}[shift={(0,0)}]
    \draw[blue!70!black, thick, domain=0:360, samples=200, smooth, variable=\t]
      plot ({1.25*cos(\t)}, {1.25*sin(\t)*cos(\t)});
    \fill[bad] (0,0) circle (1.8pt);
    \draw[leader] (0,\anny+0.16) -- (0,-0.14);
    \node[bad, font=\scriptsize] at (0,\anny) {crossing};
    \node[font=\small] at (0,\capy) {\textbf{(a)} figure eight};
  \end{scope}

  % ---- (b) lollipop: circle centred (-0.50,0) of radius 0.62, so its rightmost
  %      point is (0.12,0); the stick runs from there to (1.25,0), so the junction
  %      dot at (0.12,0) really does lie on both pieces.
  \begin{scope}[shift={(4.0,0)}]
    \draw[blue!70!black, thick] (-0.50,0) circle (0.62);
    \draw[blue!70!black, thick] (0.12,0) -- (1.25,0);
    \fill[bad] (0.12,0) circle (1.8pt);
    \draw[leader] (0.42,\anny+0.16) -- (0.16,-0.10);
    \node[bad, font=\scriptsize] at (0.42,\anny) {junction};
    \node[font=\small] at (0.30,\capy) {\textbf{(b)} lollipop};
  \end{scope}

  % ---- (c) closed square: [0,1]^2 drawn with side 1.3. The marked point is the
  %      midpoint of the bottom edge; the orange region is the part of a small ball
  %      about it that actually lies in the square -- the upper half only.
  \begin{scope}[shift={(7.9,0)}]
    \fill[blue!10] (0,-0.65) rectangle (1.3,0.65);
    \draw[blue!70!black, thick] (0,-0.65) rectangle (1.3,0.65);
    \fill[orange!30] (0.65,-0.65) -- ++(0.42,0)
      arc[start angle=0, end angle=180, radius=0.42] -- cycle;
    \draw[orange!90!black, thick] (0.65-0.42,-0.65)
      arc[start angle=180, end angle=0, radius=0.42];
    \fill[bad] (0.65,-0.65) circle (1.8pt);
    \draw[leader] (0.65,\anny+0.16) -- (0.65,-0.75);
    \node[bad, font=\scriptsize] at (0.65,\anny) {only half a ball};
    \node[font=\small] at (0.65,\capy) {\textbf{(c)} closed square};
  \end{scope}
\end{tikzpicture}
\end{center}

\subsection{Proving a negative: the tangent field test}

% Supplement: Linus Lüchinger/Slides/Slides-04-30.pdf, slide 7
\cref{ex:cross_not_submanifold} disproved a submanifold claim by counting connected components
after deleting a point. That argument is sharp but narrow: it needs the bad point to disconnect
the set. Linus gives a second, complementary criterion, which catches cases the counting argument
misses.

\begin{proposition}[Continuous tangent field along a curve]
\label{prop:continuous_tangent_field}
Let $M \subseteq \mathbb{R}^n$ be a $1$-dimensional $C^1$ submanifold. Then every point of $M$
has a neighbourhood on which there is a continuous map $T$ assigning to each $q$ a nonzero vector
$T(q) \in T_qM$.
\end{proposition}

\begin{proof}
Fix $p \in M$. By \cref{thm:local_parametrization} there are an open interval $I \subseteq
\mathbb{R}$, a neighbourhood $V$ of $p$ in $\mathbb{R}^n$, and a $C^1$ homeomorphism
$\gamma : I \to V \cap M$ with $D\gamma_t$ injective for every $t$. Injectivity of a linear map
$\mathbb{R} \to \mathbb{R}^n$ says exactly that $\gamma'(t) \neq 0$. Set
\[T(q) := \gamma'\big(\gamma^{-1}(q)\big), \qquad q \in V \cap M.\]
This is nonzero by the above, lies in $T_qM$ by \cref{prop:tangent_space_via_curves}, and is
continuous as a composition of the continuous maps $\gamma^{-1}$ and $\gamma'$.
\end{proof}

The test is used in contrapositive form: \textbf{if no continuous nonvanishing tangent field
exists near a point, the set is not a $1$-dimensional submanifold there.}

% Generator: Claude Opus 5 (High)
\begin{aiexample}[The lollipop, disproved by tangent fields]
\label{ex:lollipop_tangent_field}
Let $L$ be the lollipop of the gallery above and let $q$ be its junction point. Deleting $q$
leaves \emph{three} pieces where a curve would leave two, so the counting argument of
\cref{ex:cross_not_submanifold} already applies. But it is instructive to see the tangent-field
test reach the same verdict, since it localises the contradiction differently.

Suppose $L$ were a $1$-dimensional submanifold near $q$, and let $T$ be the tangent field
supplied by \cref{prop:continuous_tangent_field}. Approach $q$ along the stick: the tangent
direction there is constant, so $T$ tends to a vector along the stick. Approach $q$ along the
circle instead: the circle meets the stick at a right angle at $q$, so $T$ tends to a vector
perpendicular to the stick. A continuous map cannot have two different limits at the same point,
so no such $T$ exists.

\begin{remark}[When to reach for which]
Both tests disprove the same statement, and neither subsumes the other.
Counting components is the tool when the offending point \emph{separates} the set, and it needs
no differentiability. The tangent-field test is the tool when the set stays connected but the
direction of travel jumps. A cusp such as $\{y^2 = x^3\}$ at the origin is the standard case,
where removing the origin leaves the expected two pieces and only the tangent direction gives the
failure away.
\end{remark}
\end{aiexample}

\subsection{Transporting a submanifold by a diffeomorphism}

% Generator: Claude Opus 5 (High)
% Gap closed 2026-08-10: def:submanifold is stated through charts, and the fact that a
% diffeomorphism carries submanifolds to submanifolds of the same dimension follows from it in
% four lines, but it was nowhere in the document. ex:image_of_segment_submanifold below needs it,
% and so does any argument that moves a submanifold into better coordinates.
Each of the three descriptions above answers the question \qt{is this particular set a
submanifold?}. A different question arises as soon as a map is in play: given that $M$ is a
submanifold and $\Phi$ a diffeomorphism, what is $\Phi(M)$? Here the chart definition is the easy
one to use, and it is about the only place where that is so.

\begin{proposition}[Diffeomorphisms carry submanifolds to submanifolds]
\label{prop:diffeomorphism_preserves_submanifold}
Let $U, W \subseteq \mathbb{R}^n$ be open, let $\Phi : U \to W$ be a $C^k$-diffeomorphism, and let
$M \subseteq U$ be an $m$-dimensional $C^k$ submanifold of $\mathbb{R}^n$. Then $\Phi(M)$ is an
$m$-dimensional $C^k$ submanifold of $\mathbb{R}^n$. In particular the dimension is unchanged.
\end{proposition}

\begin{proof}
Fix $q \in \Phi(M)$ and write $q = \Phi(p)$ with $p \in M$. By \cref{def:submanifold} there are an
open neighbourhood $U_p \subseteq \mathbb{R}^n$ of $p$, an open $V \ni 0$, and a
$C^k$-diffeomorphism $\Psi : U_p \to V$ with $\Psi(p) = 0$ and
$\Psi(U_p \cap M) = (\mathbb{R}^m\times\{0\}) \cap V$. Replacing $U_p$ by $U_p \cap U$, we may
assume $U_p \subseteq U$.

Then $\Phi(U_p)$ is an open neighbourhood of $q$, and
\[\Psi \circ \Phi^{-1} : \Phi(U_p) \longrightarrow V\]
is a composition of $C^k$-diffeomorphisms, hence itself one, and it sends $q$ to $\Psi(p) = 0$.
Since $\Phi$ is injective, $\Phi(U_p) \cap \Phi(M) = \Phi(U_p \cap M)$, and therefore
\[(\Psi\circ\Phi^{-1})\big(\Phi(U_p) \cap \Phi(M)\big) = \Psi(U_p \cap M)
  = (\mathbb{R}^m\times\{0\}) \cap V .\]
That is a submanifold chart for $\Phi(M)$ at $q$, of the same dimension $m$. As $q$ was arbitrary,
$\Phi(M)$ is an $m$-dimensional $C^k$ submanifold.
\end{proof}

% Source: old_exams/fs2023/Probeprfg1.pdf (Probeprüfung Analysis II, undated), Aufgabe 7, p. 2
% Extractor: Claude Opus 5 (High)
% The paper names no lecturer; see old-exam-mining.md.
% Placed directly after prop:diffeomorphism_preserves_submanifold because it is the same question
% -- which operations preserve the property -- asked about four more of them. Parts (c) and (d)
% add products and transversal intersections, neither of which appears anywhere in this chapter,
% and (d) is the only place in the document where transversality is used.
\begin{exercise}[True or False \textnormal{(which operations preserve submanifolds?)}]
\label{ex:operations_preserving_submanifolds}
Decide whether each of the following is true or false, with justification.
\begin{enumerate}[label=\textbf{(\alph*)}]
  \item If $M$ is a $k$-dimensional submanifold of $\mathbb{R}^n$ and $N \subseteq M$ is open in
    $M$, then $N$ is itself a $k$-dimensional submanifold of $\mathbb{R}^n$.
  \item The union of two disjoint $k$-dimensional submanifolds of $\mathbb{R}^n$ is a
    $k$-dimensional submanifold of $\mathbb{R}^n$.
  \item The Cartesian product of a $k$-dimensional submanifold of $\mathbb{R}^m$ and an
    $l$-dimensional submanifold of $\mathbb{R}^n$ is a $(k+l)$-dimensional submanifold of
    $\mathbb{R}^{m+n}$.
  \item If $M$ and $N$ are $(n-1)$-dimensional submanifolds of $\mathbb{R}^n$ for $n \geq 2$, and
    $T_pM \neq T_pN$ for every $p \in M \cap N$, then $M \cap N$ is an $(n-2)$-dimensional
    submanifold of $\mathbb{R}^n$.
\end{enumerate}
\exinfo{This exercise is taken from the mock examination \emph{Probepr\"ufung Analysis II}
(undated) in \texttt{old\_exams/fs2023/}, Problem 7. That paper names no lecturer.}
\exsol{sol:operations_preserving_submanifolds}
\end{exercise}

\begin{remark}[Why the dimension cannot move]
\label{rem:dimension_is_a_diffeomorphism_invariant}
The proof transports one chart into another and never touches $m$, which is the whole content: a
diffeomorphism can bend a submanifold arbitrarily, but it cannot change how many independent
directions there are along it. This is worth stating because the tempting error runs the other
way, imagining that a map into $\mathbb{R}^2$ whose image \qt{spreads out} raises the dimension.
It cannot. A curve stays a curve.
\end{remark}

% Source: old_exams/examFS24.pdf (21 August 2024), Wahr/Falsch-Fragen, Aufgabe 6, p. 4
% Extractor: Claude Opus 5 (High)
% Split: parts 6.1-6.3 of this block are ex:local_diffeo_parameter_alpha in
% content/15-inverse-function-theorem/01-inverse-function-theorem.tex. The enumeration is continued
% at (d) with \setcounter so that the two halves read as one exam question, which is what they are.
\begin{exercise}[True or False \textnormal{(the image of a segment)}]
\label{ex:image_of_segment_submanifold}
Continuing \cref{ex:local_diffeo_parameter_alpha}: let $\Phi \in C^1(\mathbb{R}^2,\mathbb{R}^2)$
satisfy $\Phi(0,0) = (0,0)\transp$ and
$\Jac\Phi(0,0) = \left(\begin{smallmatrix} \alpha+5 & -3 \\ 3 & \alpha-5\end{smallmatrix}\right)$,
and take $\alpha = 0$. Decide whether each of the following is true or false, with justification.
\begin{enumerate}[label=\textbf{(\alph*)}]
  \setcounter{enumi}{3}
  \item For $r > 0$ small enough, the image of $\{(x,y) \in B_r(0,0) \mid y = 0\}$ under $\Phi$ is
    a $C^1$ submanifold of dimension $1$.
  \item For $r > 0$ small enough, the image of $\{(x,y) \in B_r(0,0) \mid y = 0\}$ under $\Phi$ is
    a $C^1$ submanifold of dimension $2$.
\end{enumerate}
\exinfo{This exercise is taken from the exam of 21 August 2024 (Prof.\ Joaquim Serra), where it appears as the fourth and
fifth parts of the sixth block of true/false questions. Its first three parts are
\cref{ex:local_diffeo_parameter_alpha}.}
\exsol{sol:image_of_segment_submanifold}
\end{exercise}

% Source: exercises_2024/mock.pdf (Analysis II mock examination, Spring Semester 2024,
% Prof. Joaquim Serra), Exercise 9, p. 2
% Extractor: Claude Opus 5 (High)
% The R^3 -> R^2 counterpart of ex:local_diffeo_parameter_alpha / ex:image_of_segment_submanifold
% above, and not a duplicate of them: there the map is square and the question is invertibility,
% here it is wide and the question is what a surjective differential buys. Parts (a) and (d) are
% the ones with teeth --- the curve gamma passes through the point but does not lie in M, which
% both of them turn on.
% Part (4) of the printed block carries no "True / False" prompt, unlike its five siblings;
% transcribed as a statement to be decided, in line with the rest.
\begin{exercise}[True or False \textnormal{(a level set in $\mathbb{R}^3$, and a curve through it)}]
\label{ex:level_set_of_a_wide_jacobian}
Let $F \in C^1(\mathbb{R}^3,\mathbb{R}^2)$ satisfy
\[F(0,0,1) = \begin{pmatrix} 2 \\ -2\end{pmatrix} \qquad\text{and}\qquad
  \Jac F(0,0,1) = \begin{pmatrix} 1 & 1 & -2 \\ 0 & 1 & -1\end{pmatrix},\]
and set $M := F^{-1}\big(\{(2,-2)\transp\}\big)$. For $t < 1$ let
\[\gamma(t) := \Big(\sin(2t),\ t+\tfrac12t^2,\ \tfrac{\cos t}{1-t}\Big).\]
Decide whether each of the following is true or false, with justification.
\begin{enumerate}[label=\textbf{(\alph*)}]
  \item $(F\circ\gamma)'(0) = (0,0)\transp$.
  \item In a small neighbourhood of $(0,0,1)$, the set $M$ is a smooth submanifold of
    dimension $2$.
  \item In a small neighbourhood of $(0,0,1)$, the set $M$ is a smooth submanifold of
    dimension $1$.
  \item The vector $\gamma'(0)$ is normal to $M$ at the point $(0,0,1)$.
  \item There exist $\delta > 0$ and $\varphi,\psi \in C^1((-\delta,\delta),\mathbb{R})$ such
    that $F(\varphi(y),\,y,\,\psi(y)) = (2,-2)\transp$ for all $y \in (-\delta,\delta)$.
  \item For all sufficiently small $\varepsilon > 0$, the set
    $F\big(B_\varepsilon((0,0,1))\big)$ is open in $\mathbb{R}^2$.
\end{enumerate}
\exinfo{This exercise is taken from the mock examination for Analysis II, Spring Semester 2024
(Prof.\ Joaquim Serra), where it appears as the ninth block of multiple-choice questions.}
\exsol{sol:level_set_of_a_wide_jacobian}
\end{exercise}

% Transition: Claude Opus 5 (High)
Every argument in this section leaned on knowing which vectors point \emph{along} $M$ at a given
point. The next section makes that space of directions an object in its own right.

% Originally: content/week-07.tex, \section{Tangent and normal spaces} (Corsin Week 7)

% Source: Corsin Nick/Class Notes/Week 7.pdf, pp. 7--8
\section{Tangent and normal spaces}
\label{sec:tangent_normal_spaces}

\subsection{Tangent space}

Suppose $M^m \subseteq \mathbb{R}^n$ is a submanifold and $f : U \to f(U) \subseteq M$ is a local
parametrization around some point $p \in f(U)$ with $f(q) = p$. The $n\times m$ matrix
\[\Jac f(q) = \begin{pmatrix} \dfrac{\partial f}{\partial q_1} & \cdots & \dfrac{\partial f}{\partial q_m} \end{pmatrix}\]
has full rank, that is, the column vectors $\partial_if(q)$ are linearly independent. They form a
basis of $T_pM$:

% Sonnet 5 (Medium) --- FIG-W07-04
\begin{center}
\begin{tikzpicture}[>=Latex, scale=1.1]
  % A wireframe surface patch (parallel curve families) with a marked point p
  \foreach \i in {0,0.6,...,2.4} {
    \draw[blue!40!black, thin] plot[smooth] coordinates
      {(-1.4+\i*0.15,-1.1+\i*0.25) (-0.5+\i*0.15,-0.3+\i*0.25) (0.6+\i*0.15,0.1+\i*0.25) (1.6+\i*0.15,-0.2+\i*0.25)};
  }
  \foreach \j in {0,0.55,...,2.2} {
    \draw[blue!40!black, thin] plot[smooth] coordinates
      {(-1.4+\j*0.15,-1.1+\j*0.05) (-0.6+\j*0.35,-0.55+\j*0.15) (0.3+\j*0.55,-0.1+\j*0.2) (1.3+\j*0.7,0.05+\j*0.2)};
  }
  \coordinate (P) at (0.1,-0.05);
  \fill (P) circle (1.7pt) node[below, font=\small] {$p = f(q)$};
  \draw[orange!90!black, thick, ->] (P) -- ++(1.1,0.35)
    node[right, font=\scriptsize] {$\partial f/\partial x_1$};
  \draw[purple, thick, ->] (P) -- ++(0.35,0.95)
    node[above, font=\scriptsize] {$\partial f/\partial x_2$};
\end{tikzpicture}
\end{center}

\begin{definition}[Tangent space]
\label{def:tangent_space}
For $M$, $f$ as above, the \newterm{tangent space} \germanfor{Tangentialraum} of $M$ at
$p \in M$ is
\[T_pM := \operatorname{span}\{\partial_1f(q),\dots,\partial_mf(q)\} = Df_q(\mathbb{R}^m);\]
it is independent of the choice of $f$.
\end{definition}

\begin{example}[Tangent space to the upper hemisphere]
Parametrize the upper hemisphere $S^2_+ := \{(x,y,z)\in\mathbb{R}^3 : x^2+y^2+z^2=1,\ z>0\}$
graphically, i.e.,
\[f(x,y) = \begin{pmatrix} x \\ y \\ \sqrt{1-x^2-y^2} \end{pmatrix}.\]
We determine $T_{e_3}S^2_+$, where $e_3 = (0,0,1)$:
\begin{align*}
\left.\frac{\partial f}{\partial x}\right|_{(x,y)=(0,0)}
  &= \left.\begin{pmatrix} 1 \\ 0 \\ \dfrac{-x}{\sqrt{1-x^2-y^2}}\end{pmatrix}\right|_{(x,y)=(0,0)}
   = \begin{pmatrix}1\\0\\0\end{pmatrix} = e_1, \\[1ex]
\left.\frac{\partial f}{\partial y}\right|_{(x,y)=(0,0)}
  &= \left.\begin{pmatrix} 0 \\ 1 \\ \dfrac{-y}{\sqrt{1-x^2-y^2}}\end{pmatrix}\right|_{(x,y)=(0,0)}
   = \begin{pmatrix}0\\1\\0\end{pmatrix} = e_2.
\end{align*}
So, $T_{e_3}S^2_+ = T_{e_3}S^2 = \operatorname{span}\{e_1,e_2\} = \mathbb{R}^2\times\{0\}$.
\end{example}

% Source: Diego Torres Tejeda/Notes - 13.04 - Diego Torres Tejeda.pdf, p. 2
% Generator: Gemini 3.6 Flash (High)
\begin{example}[Tangent space to the 2-Torus]
\label{ex:torus_tangent_space}
Consider the 2-torus $\mathbb{T}^2 \subset \mathbb{R}^3$ parametrized by $\phi(t,\theta) : \mathbb{R}^2 \to \mathbb{R}^3$,
\[
\phi(t,\theta) := \begin{pmatrix} (R_1 + R_2 \cos\theta)\cos t \\ (R_1 + R_2 \cos\theta)\sin t \\ R_2 \sin\theta \end{pmatrix},
\]
where $R_1 > R_2 > 0$ are the major and minor radii. The partial derivatives of $\phi$ with respect to the angles $t$ and $\theta$ yield two linearly independent column vectors:
\[
\frac{\partial \phi}{\partial t} = \begin{pmatrix} -(R_1 + R_2\cos\theta)\sin t \\ (R_1 + R_2\cos\theta)\cos t \\ 0 \end{pmatrix}, \qquad
\frac{\partial \phi}{\partial \theta} = \begin{pmatrix} -R_2\sin\theta\cos t \\ -R_2\sin\theta\sin t \\ R_2\cos\theta \end{pmatrix}.
\]
By \cref{def:tangent_space}, at any point $p = \phi(t,\theta) \in \mathbb{T}^2$, the tangent space is the two-dimensional vector subspace
\[
T_p \mathbb{T}^2 = \operatorname{span}\left\{ \frac{\partial \phi}{\partial t}(t,\theta), \; \frac{\partial \phi}{\partial \theta}(t,\theta) \right\}.
\]
In particular, at the outer point $p_0 := (R_1+R_2, 0, 0)\transp$ (where $t = 0, \theta = 0$), we obtain $\frac{\partial \phi}{\partial t}(0,0) = (0, R_1+R_2, 0)\transp$ and $\frac{\partial \phi}{\partial \theta}(0,0) = (0, 0, R_2)\transp$, so $T_{p_0}\mathbb{T}^2 = \{0\} \times \mathbb{R}^2$.
\end{example}

% Sonnet 5 (Medium) --- FIG-W07-05
\begin{center}
\begin{tikzpicture}[scale=1.1]
  \draw[blue!70!black, thick] (0,0) circle (1.3);
  \draw[blue!70!black, dashed] (1.3,0) arc (0:180:1.3 and 0.4);
  \draw[blue!70!black] (-1.3,0) arc (180:360:1.3 and 0.4);
  \node[blue!70!black, font=\small] at (-1.0,0.5) {$S^2$};

  \fill (0,1.3) circle (1.7pt) node[right, font=\small] {$e_3$};
  \draw[orange!90!black, thick, fill=orange!10, fill opacity=0.5]
    (-1.7,1.3) -- (1.7,1.3) -- (2.0,1.8) -- (-1.4,1.8) -- cycle;
  \node[orange!90!black, font=\small] at (1.5,1.65) {$T_{e_3}S^2$};
\end{tikzpicture}
\end{center}

% Quelle: Diego Torres Tejeda/Notes - 30.03 - Diego Torres Tejeda.pdf, p. 2 -- Sonnet 5 (Medium)
% Corsin defines T_pM only via a parametrization; the curve description below is Diego's. (Was
% an ainote until 2026-08-09. The mathematics is the proposition, and who-wrote-which-definition
% is editor-facing by Test 2 in style.md.)
\Cref{def:tangent_space} builds $T_pM$ out of a parametrization, so the space is pinned down by a
formula rather than by what its elements are. There is an equivalent description that says what
they are. A vector is tangent to $M$ at $p$ exactly when some curve running inside $M$ passes
through $p$ with that velocity, which is both how tangent vectors are usually pictured and the
quicker test in practice, since no parametrization need be written down.

\begin{proposition}[Tangent space via curves]
\label{prop:tangent_space_via_curves}
Let $M^m \subseteq \mathbb{R}^n$ be a $C^k$ submanifold ($k\geq1$) and $p \in M$. Then
\[T_pM = \big\{\gamma'(0) : \gamma \in C^1\big((-\varepsilon,\varepsilon), \mathbb{R}^n\big),\ \varepsilon>0,\ \gamma(t) \in M \text{ for all } t,\ \gamma(0)=p\big\}.\]
\end{proposition}

\begin{proof}
Let $\Phi : U \to V$ be a submanifold chart at $p$ (\cref{def:submanifold}), so
$\Phi(U\cap M) = (\mathbb{R}^m\times\{0\})\cap V$ and $\Phi(p)=0$; equivalently, let
$f := \Phi^{-1}|_{(\mathbb{R}^m\times\{0\})\cap V}$ be the corresponding local parametrization
(\cref{thm:local_parametrization}), $f(q)=p$.

\textbf{($\subseteq$)} Let $v \in T_pM = Df_q(\mathbb{R}^m)$, say $v = Df_q(w)$. The curve
$\gamma(t) := f(q+tw)$, defined for $t$ small enough that $q+tw$ stays in $f$'s domain, satisfies
$\gamma(t) \in M$ for all such $t$, $\gamma(0) = p$, and by the chain rule
$\gamma'(0) = Df_q(w) = v$.

\textbf{($\supseteq$)} Let $\gamma$ be a curve as in the statement. Since $\Phi$ is a
diffeomorphism with $\Phi(\gamma(t)) \in \mathbb{R}^m\times\{0\}$ for all $t$ near $0$, write
$\Phi\circ\gamma =: (\eta,0)$ with $\eta \in C^1\big((-\varepsilon,\varepsilon),\mathbb{R}^m\big)$.
Then $\gamma = \Phi^{-1}(\eta,0) = f(\eta(\cdot))$ near $0$ (shrinking $U$ if necessary so that
$f$ is defined on all of $\eta$'s image), so by the chain rule,
\[\gamma'(0) = Df_{\eta(0)}\big(\eta'(0)\big) \in Df_q(\mathbb{R}^m) = T_pM\]
(using $\eta(0)=q$, since $\Phi(p)=0=\Phi(f(q))$ and $\Phi$ is injective).
\end{proof}

% Sonnet 5 (Medium)
\begin{center}
\begin{tikzpicture}[>=Latex, scale=1.1]
  \draw[blue!70!black, thick] plot[smooth] coordinates
    {(-1.8,-0.8) (-0.8,0.6) (0.4,-0.2) (1.4,0.9) (2.2,0.1)};
  \coordinate (P) at (0.4,-0.2);
  \fill (P) circle (1.7pt) node[below, font=\small] {$p=\gamma(0)$};
  \draw[orange!90!black, thick, ->] (P) -- ++(0.9,0.75) node[right, font=\small] {$\gamma'(0)$};
\end{tikzpicture}
\end{center}

% Source: Jérôme Paschoud/Notizen/Woche 8.1 Tangentialräume und Jordan Mass.pdf, p. 1
% Generator: Gemini 3.1 Pro (High)
\begin{remark}[Four representations of the tangent space]
\label{rem:four_representations_tangent_space}
Depending on how a $m$-dimensional submanifold $M \subseteq \mathbb{R}^n$ is locally described near $p$, its tangent space $T_pM$ can be computed algebraically in four equivalent ways:
\begin{itemize}
  \item \textbf{Diffeomorphism} (from the definition, $\Psi(p)=0$): $T_pM = (D\Psi_p)^{-1}(\mathbb{R}^m \times \{0\})$.
  \item \textbf{Implicit representation} ($M = F^{-1}\{0\}$): $T_pM = \ker(DF_p)$.
  \item \textbf{Parametric representation} (local parametrization $g$ with $g(q)=p$): $T_pM = Dg_q(\mathbb{R}^m)$.
  \item \textbf{Graph representation} (graph of $h$, $p=(z, h(z))$): $T_pM = \{(w, Dh_z(w)) \mid w \in \mathbb{R}^m\}$.
\end{itemize}
These algebraic descriptions are all equivalent to the geometric description via velocities of curves (\cref{prop:tangent_space_via_curves}).
\end{remark}

\subsection{Normal space}

% Source: Corsin Nick/Class Notes/Week 7.pdf, pp. 9--10
Suppose $F : U \to \mathbb{R}^{n-m}$ is an implicit representation of $M^m = F^{-1}\{0\}$ (via the
regular value theorem) and $f : V^m \to M$ is a local parametrization, $M$ a submanifold. Then, by
definition,
\begin{enumerate}[label=\textbf{\arabic*.}]
  \item The Jacobian matrix:
\[ \Jac F(x) = \begin{pmatrix} \nabla F_1(x) \longrightarrow \\ \vdots \\ \nabla F_{n-m}(x) \longrightarrow \end{pmatrix} \]
    has full rank, i.e., $\{\nabla F_i(x)\}_{i=1,\dots,n-m}$ are linearly independent;
  \item $T_{f(q)}M = \operatorname{span}\{\partial_if(q)\}$, as above;
  \item $F(f(x)) = 0$ for all $x \in V$, since $f(x) \in M = F^{-1}\{0\}$.
\end{enumerate}
From \textbf{3}, it follows that
\begin{align*}
\Jac(F\circ f)(q)
  &= \Jac F(f(q))\cdot\Jac f(q) \\
  &= \begin{pmatrix} \nabla F_1(f(q)) \longrightarrow \\ \vdots \\ \nabla F_{n-m}(f(q)) \longrightarrow \end{pmatrix}
     \begin{pmatrix} \dfrac{\partial f}{\partial q_1} & \cdots & \dfrac{\partial f}{\partial q_m}\end{pmatrix}
   \;=\; 0.
\end{align*}
Now, since $\big(\Jac F(f(q))\cdot\Jac f(q)\big)_{ij} = \langle \nabla F_i(f(q)), \partial_jf(q)\rangle = 0$
for all $i=1,\dots,n-m$, $j=1,\dots,m$, the vectors $\nabla F_i(f(q))$ are perpendicular, that is,
\newterm{normal}, to $T_{f(q)}M$. Since by \textbf{1} we have $n-m$ linearly independent vectors
normal to the $m$-dimensional plane $T_{f(q)}M$, we can define:

\begin{definition}[Normal space]
\label{def:normal_space}
Let $M^m \subseteq \mathbb{R}^n$ be a submanifold and $F : U^n \to \mathbb{R}^{n-m}$ an implicit
representation, $M = F^{-1}\{0\}$. Then the \newterm{normal space} \germanfor{Normalraum} of
$M$ at $p \in M$ is
\[N_pM := \operatorname{span}\{\nabla F_1(p),\dots,\nabla F_{n-m}(p)\}.\]
\end{definition}

\begin{example}[Normal space to the sphere]
Represent $S^2$ implicitly: $S^2 = F^{-1}\{1\} = \{(x,y,z)\in\mathbb{R}^3 : F(x,y,z)=1\}$ where
$F(x,y,z) = x^2+y^2+z^2$. We determine $N_{e_3}S^2$:
\[\nabla F(0,0,1) = (2x,2y,2z)\big|_{(x,y,z)=(0,0,1)} = 2e_3.\]
So, $N_{e_3}S^2 = \operatorname{span}\{e_3\} = \{0\}\times\mathbb{R}$.
\end{example}

% Generator: Claude Opus 5 (High) --- the example above is the $n=3$, $F = \lvert x\rvert^2$
% instance of a general fact worth stating once, since it is what makes the Lagrange condition
% of \cref{prop:constrained_optimization} geometric rather than algebraic.
\begin{proposition}[The gradient is normal to the level set]
\label{prop:gradient_orthogonality_level_sets}
Let $U \subseteq \mathbb{R}^n$ be open, $g \in C^1(U,\mathbb{R})$, and let
$S := \{x \in U : g(x) = c\}$ for some $c \in \mathbb{R}$. If $\nabla g(x_0) \neq 0$ at a point
$x_0 \in S$, then near $x_0$ the set $S$ is an $(n-1)$-dimensional submanifold and
\[N_{x_0}S = \operatorname{span}\{\nabla g(x_0)\}.\]
In particular the normal space is a \emph{line}, and the affine \emph{tangent} hyperplane to $S$
at $x_0$ is $\{x \in \mathbb{R}^n : \langle x - x_0,\ \nabla g(x_0)\rangle = 0\}$.
\end{proposition}

\begin{proof}
That $S$ is a submanifold of dimension $n-1$ near $x_0$ is \cref{thm:regular_value_theorem},
since $\nabla g(x_0) \neq 0$ says exactly that $Dg_{x_0} : \mathbb{R}^n \to \mathbb{R}$ is
surjective. For the normal space, let $v \in T_{x_0}S$ and pick a $C^1$ curve
$\gamma : (-\varepsilon,\varepsilon) \to S$ with $\gamma(0) = x_0$ and $\gamma'(0) = v$. Then
$g \circ \gamma \equiv c$, so differentiating at $t=0$ with the chain rule along a path
(\cref{sec:chain_rule_along_path}) gives
\[0 = \frac{d}{dt}\Big|_{t=0}g(\gamma(t)) = \big\langle \nabla g(x_0),\ v\big\rangle.\]
Hence $\nabla g(x_0) \perp T_{x_0}S$, that is,
$\operatorname{span}\{\nabla g(x_0)\} \subseteq N_{x_0}S$. Both are subspaces of dimension
$n - \dim T_{x_0}S = n-(n-1) = 1$, so they are equal.
\end{proof}

% Supplement: Sascha Brack/Class Notes/Week_08_Notes_Monday.pdf, pp. 3--4
\begin{example}[Tangent and normal vectors via parametrization]
\label{ex:tangent_normal_parametrized_sphere}
For a parametrized surface $f : V \to \mathbb{R}^3$, the tangent space at $p = f(x)$ is $T_pM = \operatorname{span}\{\partial_1 f(x), \partial_2 f(x)\}$. A normal vector can be found via the cross product: $\vec{n} = \partial_1 f(x) \times \partial_2 f(x)$.

Let us apply this to the sphere $S^2$ using spherical coordinates:
\[\Phi(\theta,\varphi) = \begin{pmatrix} \sin\theta\cos\varphi \\ \sin\theta\sin\varphi \\ \cos\theta \end{pmatrix}, \qquad (\theta,\varphi) \in (0,\pi) \times (-\pi,\pi).\]
(Note that the image of $\Phi$ is the unit sphere minus a \qt{thin slice}, namely the meridian where $\varphi = \pm\pi$ together with the two poles where $\theta \in \{0,\pi\}$.)
The tangent space $T_{\Phi(\theta,\varphi)} S^2$ is spanned by the partial derivatives:
\[
\partial_\theta \Phi(\theta,\varphi) = \begin{pmatrix} \cos\theta\cos\varphi \\ \cos\theta\sin\varphi \\ -\sin\theta \end{pmatrix}, \qquad
\partial_\varphi \Phi(\theta,\varphi) = \begin{pmatrix} -\sin\theta\sin\varphi \\ \sin\theta\cos\varphi \\ 0 \end{pmatrix}.
\]
A normal vector is given by their cross product:
% Split across three lines: as one display this overflowed the text block by 70pt, which the
% house style forbids for chains of matrix blocks.
\begin{align*}
\vec{n} = \partial_\theta \Phi \times \partial_\varphi \Phi
  &= \begin{pmatrix} \cos\theta\cos\varphi \\ \cos\theta\sin\varphi \\ -\sin\theta \end{pmatrix}
     \times
     \begin{pmatrix} -\sin\theta\sin\varphi \\ \sin\theta\cos\varphi \\ 0 \end{pmatrix} \\
  &= \begin{pmatrix} \sin^2\theta\cos\varphi \\ \sin^2\theta\sin\varphi \\
       \sin\theta\cos\theta\cos^2\varphi + \sin\theta\cos\theta\sin^2\varphi \end{pmatrix} \\
  &= \sin\theta \begin{pmatrix} \sin\theta\cos\varphi \\ \sin\theta\sin\varphi \\ \cos\theta \end{pmatrix}.
\end{align*}
Its length is $|\vec{n}| = \sin\theta$ (since $\sin\theta > 0$ for $\theta \in (0,\pi)$).
The unit normal vector is therefore
\[
\nu(\theta,\varphi) := \frac{\vec{n}}{|\vec{n}|} = \begin{pmatrix} \sin\theta\cos\varphi \\ \sin\theta\sin\varphi \\ \cos\theta \end{pmatrix} = \Phi(\theta,\varphi).
\]
This confirms that the outward normal vector at a point on the unit sphere is precisely the position vector of that point.
\end{example}

% Sonnet 5 (Medium) --- FIG-W07-06
\begin{center}
\begin{tikzpicture}[scale=1.1]
  \draw[blue!70!black, thick] (0,0) circle (1.3);
  \draw[blue!70!black, dashed] (1.3,0) arc (0:180:1.3 and 0.4);
  \draw[blue!70!black] (-1.3,0) arc (180:360:1.3 and 0.4);
  \node[blue!70!black, font=\small] at (-1.0,0.5) {$S^2$};

  \draw[purple, thick, ->] (0,-2.0) -- (0,2.1) node[above, font=\small] {$N_{e_3}S^2$};
  \fill (0,1.3) circle (1.7pt) node[right, font=\small] {$e_3$};
  \fill (0,-1.3) circle (1.7pt) node[right, font=\small] {$-e_3$};
\end{tikzpicture}
\end{center}


% --- exercise relocated here from the dissolved sheet 7 ---

% Originally: content/week-07.tex, \exercisesheet{7}, problem 7.6
% Source: exercises/Ex7_Analysis2_eng.pdf

\begin{exercise}[7.6 --- Spherical coordinates \textnormal{(important)}]
\label{ex:7.6}
The mapping $\Phi : (0,\infty)\times(0,\pi)\times(-\pi,\pi) \to \mathbb{R}^3$ defined by
\[\Phi(r,\theta,\varphi) = \begin{pmatrix} r\sin\theta\cos\varphi \\ r\sin\theta\sin\varphi \\ r\cos\theta \end{pmatrix}\]
is called \textit{spherical coordinates}.
\begin{enumerate}[label=\textbf{\arabic*.}]
  \item Sketch the images of $r \mapsto \Phi(r,\theta_0,\varphi_0)$, $\theta \mapsto
    \Phi(r_0,\theta,\varphi_0)$ and $\varphi \mapsto \Phi(r_0,\theta_0,\varphi)$ for some fixed
    $r_0 \in (0,\infty)$, $\theta_0 \in (0,\pi)$, $\varphi_0 \in (-\pi,\pi)$.
  \item What is the image of $\Phi$?
  \item Show that $\det(D_{(r,\theta,\varphi)}\Phi) = r^2\sin\theta$ holds.
  \item Conclude that the mapping $\Phi$ is a diffeomorphism onto its image.
\end{enumerate}
\end{exercise}

% Originally: content/week-07.tex, \section{Solutions}

\section{Solutions}
\label{sec:submanifolds_solutions}

% Transition: Claude Opus 5 (High)
Not every solution in this chapter is collected here. \Cref{ex:dimension_submanifold_system}
(\cpageref{ex:dimension_submanifold_system}) and \cref{ex:submanifold_matrices_rank1}
(\cpageref{ex:submanifold_matrices_rank1}) are solved where they appear.

\begin{ainote}
Corsin's notes for this material cover only the theory (submanifolds, the regular value and local
parametrization theorems, tangent and normal spaces), with no page working out any of the
official-sheet exercises, so all solutions below are original. Those answering a problem from
Sheet 7 have been checked against \texttt{exercises/Sol7\_Analysis2\_eng.pdf}, and no divergence
from the official solutions was found. The exam problems mined into this chapter come from papers
that ship no key at all, so those solutions rest on independent derivation alone.
\end{ainote}

\begin{exercisesolution}[Solution to \cref{ex:7.6}]
% Generator: Claude Sonnet 5 (Medium)
\begin{enumerate}[label=\textbf{(\alph*)}]
  \item Fixing $\theta_0,\varphi_0$, the image of $r \mapsto \Phi(r,\theta_0,\varphi_0)$ is the
    radial half-line through the point $p_0 := (\sin\theta_0\cos\varphi_0,
    \sin\theta_0\sin\varphi_0, \cos\theta_0)$ on the unit sphere, i.e., $\{rp_0 : r \in
    (0,\infty)\}$. Fixing $r_0,\varphi_0$, the image of $\theta \mapsto \Phi(r_0,\theta,\varphi_0)$
    is a semicircle of radius $r_0$ from the north pole to the south pole (excluding the poles),
    at longitude $\varphi_0$. Fixing $r_0,\theta_0$, the image of $\varphi \mapsto
    \Phi(r_0,\theta_0,\varphi)$ is a circle of latitude of radius $r_0\sin\theta_0$ at height
    $r_0\cos\theta_0$ (excluding the point at $\varphi = \pm\pi$).
  \item For each fixed $r$, the image sweeps out the sphere of radius $r$ minus the half-plane
    at $\varphi = \pm\pi$ (the \emph{anti}-meridian, since the prime meridian is $\varphi = 0$).
    Overall,
    \[\operatorname{im}(\Phi) = \mathbb{R}^3\setminus\big((-\infty,0]\times\{0\}\times\mathbb{R}\big).\]
  \item Computing,
    \[D_{(r,\theta,\varphi)}\Phi = \begin{pmatrix} \sin\theta\cos\varphi & r\cos\theta\cos\varphi & -r\sin\theta\sin\varphi \\ \sin\theta\sin\varphi & r\cos\theta\sin\varphi & r\sin\theta\cos\varphi \\ \cos\theta & -r\sin\theta & 0 \end{pmatrix},\]
    hence, expanding along the last row,
    \[\begin{aligned}
    \det(D_{(r,\theta,\varphi)}\Phi) &= r^2\sin\theta\cos^2\theta\cos^2\varphi+r^2\sin^3\theta\sin^2\varphi+r^2\sin\theta\cos^2\theta\sin^2\varphi+r^2\sin^3\theta\cos^2\varphi \\
    &= r^2\big(\sin\theta\cos^2\theta+\sin^3\theta\big) = r^2\sin\theta.
    \end{aligned}\]
  \item Since $r \in (0,\infty)$ and $\theta \in (0,\pi)$, $r^2\sin\theta$ is never $0$, so
    $\Jac\Phi(r,\theta,\varphi)$ is everywhere invertible and $\Phi$ is a local diffeomorphism by
    \Cref{thm:inverse_function_theorem}. To see it is a \textbf{global} diffeomorphism onto its
    image, it remains to check $\Phi$ is injective: if $\Phi(r_1,\theta_1,\varphi_1) =
    \Phi(r_2,\theta_2,\varphi_2)$, then $r_1=r_2$ follows by taking the norm of both sides; since
    $\cos : (0,\pi) \to (-1,1)$ is injective, the third coordinate gives $\theta_1=\theta_2$;
    finally $(\cos\varphi,\sin\varphi)$ with $\varphi \in (-\pi,\pi)$ determines a unique point on
    the unit circle, so $\varphi_1=\varphi_2$. Hence $\Phi$ is injective, and (being a local
    diffeomorphism everywhere on its domain) a diffeomorphism onto its image.
\end{enumerate}
\end{exercisesolution}

\begin{exercisesolution}[Proof of \cref{ex:submanifold_quadric_tangent_space}]
% Generator: Gemini 3.6 Flash (High)
% Correction: Claude Opus 5 (High) --- ^T -> \transp throughout, \implies replaced by prose, and the
% check that p lies on M added: the tangent space computation presupposes it.
\begin{enumerate}[label=\textbf{(\alph*)}]
  \item Define the smooth function $F : \mathbb{R}^3 \to \mathbb{R}$,
    $F(x_1,x_2,x_3) := x_1 x_2 + x_2 x_3 + x_1 x_3 - 1$, so that $M = F^{-1}(0)$. Its gradient is
\[\nabla F(x) = (x_2+x_3,\; x_1+x_3,\; x_1+x_2)\transp .\]
    Suppose $\nabla F(x) = 0$. Adding the three equations $x_2+x_3 = x_1+x_3 = x_1+x_2 = 0$ gives
    $2(x_1+x_2+x_3) = 0$, and subtracting each equation in turn then forces $x_1 = x_2 = x_3 = 0$.
    But $F(0,0,0) = -1 \neq 0$, so the origin is not in $M$. Hence $\nabla F$ is non-zero at every
    point of $M$, which is to say $0$ is a regular value of $F$. By the regular value theorem
    (\cref{thm:regular_value_theorem}), $M$ is a smooth submanifold of $\mathbb{R}^3$ of dimension
    $3 - 1 = 2$.

  \item First, $p = (-1,2,3)$ does lie on $M$: indeed $(-1)(2) + (2)(3) + (-1)(3) = -2 + 6 - 3 = 1$.
    At this point,
\[\nabla F(-1,2,3) = (2+3,\; -1+3,\; -1+2)\transp = (5, 2, 1)\transp .\]
    The tangent space is the kernel of $DF(p)$, equivalently the orthogonal complement of
    $\nabla F(p)$:
\[T_p M = \left\{(v_1, v_2, v_3)\transp \in \mathbb{R}^3 \;\middle|\; 5v_1 + 2v_2 + v_3 = 0\right\}.\]
    Reading $v_3$ off from $v_1$ and $v_2$ gives two independent solutions: taking $(v_1,v_2) =
    (1,0)$ yields $v_3 = -5$, and taking $(v_1,v_2) = (0,1)$ yields $v_3 = -2$. Hence
\[\left\{\begin{pmatrix} 1 \\ 0 \\ -5 \end{pmatrix}, \begin{pmatrix} 0 \\ 1 \\ -2 \end{pmatrix}\right\}\]
    is a basis of $T_p M$; the two vectors are independent because their first two components
    already are, and there are two of them, matching $\dim T_p M = 2$.
\end{enumerate}
\begin{ainote}
The official solution gives the basis $\left\{(2,-5,0)\transp,\ (1,0,-5)\transp\right\}$, which
looks different from the one above but is not a disagreement: both are bases of the same plane
$\ker(5,2,1)$, and each vector of either list satisfies $5v_1+2v_2+v_3 = 0$. \qt{Give a basis} has
infinitely many correct answers, and no canonical one --- so when checking work against a master
solution here, verify that each vector is annihilated by $\nabla F(p)$ and that the two are
independent, rather than comparing entries. The two lists even share the vector $(1,0,-5)\transp$.
The official solution also invokes the constant rank theorem via surjectivity of $D_xF$ where we
cite the regular value theorem (\cref{thm:regular_value_theorem}); for a real-valued $F$ these say
the same thing, since a linear map $\mathbb{R}^3 \to \mathbb{R}$ is surjective precisely when its
gradient is non-zero.
\end{ainote}
\end{exercisesolution}

\begin{exercisesolution}[Solution of \cref{ex:image_of_segment_submanifold}]
% Generator: Claude Opus 5 (High)
For $\alpha = 0$ the determinant computed in \cref{ex:local_diffeo_parameter_alpha} is
$-16 \neq 0$, so fix $r > 0$ small enough that $\Phi$ restricted to $B_r(0,0)$ is a
$C^1$-diffeomorphism onto the open set $\Phi(B_r(0,0))$. Write
\[S := \{(x,y) \in B_r(0,0) \mid y = 0\} = (-r,r) \times \{0\}\]
for the segment in question. It is itself a $1$-dimensional $C^1$ submanifold of $\mathbb{R}^2$:
given $p = (a,0) \in S$, the translation $\Psi(x,y) := (x-a,\, y)$ is a diffeomorphism from
$U_p := B_r(0,0)$ onto the open set $V := \Psi(U_p)$, it sends $p$ to $0$, and
$\Psi(U_p \cap S) = (\mathbb{R}\times\{0\}) \cap V$.
\begin{enumerate}[label=\textbf{(\alph*)}]
  \setcounter{enumi}{3}
  \item \textbf{True.} Apply \cref{prop:diffeomorphism_preserves_submanifold} to the
    $C^1$-diffeomorphism $\Phi|_{B_r(0,0)} : B_r(0,0) \to \Phi(B_r(0,0))$ and to $S$. The image
    $\Phi(S)$ is a $C^1$ submanifold of $\mathbb{R}^2$ of dimension $1$.
  \item \textbf{False.} By \textbf{(d)} the set is $1$-dimensional, and no non-empty set is a
    submanifold of two different dimensions at once. Independently of that: a $2$-dimensional
    submanifold of $\mathbb{R}^2$ is an \emph{open} subset of $\mathbb{R}^2$, because a chart with
    $m = n = 2$ reads $\Psi(U \cap M) = V = \Psi(U)$, so $U \cap M = U$ and $M$ contains the open
    neighbourhood $U$ of each of its points. But $\Phi(S)$ is not open. Its complement in
    $\Phi(B_r(0,0))$ is the image of $B_r(0,0) \setminus S$, which is dense in $B_r(0,0)$, and a
    homeomorphism carries dense sets to dense sets, so every point of $\Phi(S)$ is a limit of
    points outside it.
\end{enumerate}
\begin{remark}[The pair is a dimension check, not a submanifold check]
The two items ask about the same set, so at most one can be true, and the real question is which
number a diffeomorphism is allowed to change. It changes none:
\cref{rem:dimension_is_a_diffeomorphism_invariant} is the whole content of the block. The picture
that makes \textbf{(e)} tempting is that $\Phi$ maps a two-dimensional ball onto a two-dimensional
region, so its image \qt{fills out} the plane. It does, but the segment does not, and it is the
segment being transported.
\end{remark}
\end{exercisesolution}

% Generator: Claude Opus 5 (High)
% No solution key ships with the Jossen papers, so this is an independent derivation.
\begin{exercisesolution}[\cref{ex:orthogonal_group_submanifold}]
\begin{enumerate}[label=\textbf{(\alph*)}]
  \item With $A = \left(\begin{smallmatrix} a & b \\ c & d\end{smallmatrix}\right)$,
\[AA\transp = \begin{pmatrix} a & b \\ c & d \end{pmatrix}
              \begin{pmatrix} a & c \\ b & d \end{pmatrix}
            = \begin{pmatrix} a^2+b^2 & ac+bd \\ ac+bd & c^2+d^2 \end{pmatrix},\]
    so in the coordinates of the statement
\[F(a,b,c,d) = \big(a^2+b^2,\ ac+bd,\ ac+bd,\ c^2+d^2\big).\]
    Every component is a polynomial, hence $F$ is smooth.

  \item $(AA\transp)\transp = (A\transp)\transp A\transp = AA\transp$, so $F(A)$ is symmetric for
    every $A$. This is already visible in the coordinate formula: the second and third components
    of $F$ are the same function.

  \item Differentiating the four components with respect to $a,b,c,d$ in turn,
\[\Jac F(a,b,c,d) = \begin{pmatrix}
  2a & 2b & 0 & 0 \\
  c & d & a & b \\
  c & d & a & b \\
  0 & 0 & 2c & 2d
\end{pmatrix}.\]

  \item The second and third rows coincide for every $A$, so $\rank \Jac F(A) \le 3$ always. To see
    that equality holds on $\Orth_2(\mathbb{R})$, note that $AA\transp = \id$ says exactly that the
    two rows $u := (a,b)$ and $v := (c,d)$ of $A$ are orthonormal. Suppose
\[\alpha(2a,2b,0,0) + \beta(c,d,a,b) + \gamma(0,0,2c,2d) = 0.\]
    Its first two coordinates read $2\alpha u + \beta v = 0$, and $u,v$ are linearly independent,
    so $\alpha = \beta = 0$. The last two then read $2\gamma v = 0$, so $\gamma = 0$ as well. The
    three distinct rows are therefore independent and the rank is exactly $3$.

    This is why \cref{thm:regular_value_theorem} cannot be applied to $F$ as a map
    $\mathbb{R}^4 \to \mathbb{R}^4$: the theorem needs $D_AF$ to be \emph{surjective}, that is of
    rank $4$, and the rank is $3$ at every point of $\Orth_2(\mathbb{R})$. Nor is this a defect of
    the particular level: by \textbf{(b)} the image of $F$ lies in a three-dimensional subspace, so
    $D_AF$ is nowhere surjective, for any $A$ whatsoever.

  \item Identify $\operatorname{Sym}_2(\mathbb{R}) \cong \mathbb{R}^3$ by
    $\left(\begin{smallmatrix} p & q \\ q & s\end{smallmatrix}\right) \mapsto (p,q,s)$ and let
\[\widetilde F : \mathbb{R}^4 \to \mathbb{R}^3, \qquad
  \widetilde F(a,b,c,d) := \big(a^2+b^2,\ ac+bd,\ c^2+d^2\big),\]
    which is $F$ with the duplicated component dropped. Then
    $\Orth_2(\mathbb{R}) = \widetilde F^{-1}\{(1,0,1)\}$, and
\[\Jac \widetilde F(a,b,c,d) = \begin{pmatrix}
  2a & 2b & 0 & 0 \\
  c & d & a & b \\
  0 & 0 & 2c & 2d
\end{pmatrix}\]
    has rank $3$ at every point of $\Orth_2(\mathbb{R})$, by the computation in \textbf{(d)}. So
    $(1,0,1)$ is a regular value of $\widetilde F$, and the regular value theorem makes
    $\Orth_2(\mathbb{R})$ a smooth submanifold of $\mathbb{R}^4 = \operatorname{Mat}_2(\mathbb{R})$
    of dimension $4 - 3 = 1$.

  \item By \cref{prop:gradient_orthogonality_level_sets} applied to each of the three defining
    equations, or directly from the description of the tangent space to a regular level set,
    $T_{\id}\Orth_2(\mathbb{R}) = \ker \Jac \widetilde F(\id)$. At
    $(a,b,c,d) = (1,0,0,1)$,
\[\Jac \widetilde F(1,0,0,1) = \begin{pmatrix}
  2 & 0 & 0 & 0 \\
  0 & 1 & 1 & 0 \\
  0 & 0 & 0 & 2
\end{pmatrix},\]
    so a vector $(v_1,v_2,v_3,v_4)$ lies in the kernel precisely when $v_1 = 0$, $v_4 = 0$ and
    $v_2 + v_3 = 0$. Hence
\[T_{\id}\Orth_2(\mathbb{R}) = \operatorname{span}\left\{
  \begin{pmatrix} 0 \\ 1 \\ -1 \\ 0\end{pmatrix}\right\}
  = \operatorname{span}\left\{\begin{pmatrix} 0 & 1 \\ -1 & 0\end{pmatrix}\right\},\]
    a line, in agreement with $\dim \Orth_2(\mathbb{R}) = 1$ from \textbf{(e)}. Using the full
    $4\times4$ Jacobian of \textbf{(c)} gives the same kernel, since the repeated row imposes no
    new condition.
\end{enumerate}
\begin{remark}[The tangent space names the answer]
The line found in \textbf{(f)} consists of the \emph{antisymmetric} $2\times2$ matrices, and that
is not a coincidence: differentiating $A(t)A(t)\transp = \id$ along any curve in
$\Orth_2(\mathbb{R})$ through $\id$ gives $\dot A(0) + \dot A(0)\transp = 0$. The same computation
in $n$ dimensions gives $\dim\Orth_n(\mathbb{R}) = \binom{n}{2}$, the dimension of the
antisymmetric matrices, which for $n=2$ is $1$.

That $\Orth_2(\mathbb{R})$ is one-dimensional is worth pausing on, because it is easy to expect
more. Writing $u = (\cos\theta,\sin\theta)$ for the first row, orthonormality leaves exactly two
choices of second row, so $\Orth_2(\mathbb{R})$ is a disjoint union of two circles: the rotations
$\det A = 1$ and the reflections $\det A = -1$. A submanifold need not be connected, and the
regular value theorem never claims it is.

Set against \cref{ex:submanifold_matrices_rank1}, which lives in the same $\mathbb{R}^4$, the
contrast is instructive. There a single scalar equation cut out a three-dimensional submanifold and
the regular value theorem applied as it stands. Here the naive count would suggest
$4 - 4 = 0$ dimensions, and it is wrong precisely because the four equations $AA\transp = \id$ are
not independent: one of them is a duplicate. Counting equations is only a shortcut for computing a
rank, and when the shortcut and the rank disagree it is the rank that is right.
\end{remark}
\end{exercisesolution}

\begin{exercisesolution}[Solution to \cref{ex:operations_preserving_submanifolds}]
% Generator: Claude Opus 5 (High) --- the four statements are the examiner's; the working, the
% counterexample in (b) and the remark are written for these notes. Probeprfg1 ships no key.
Throughout, \qt{submanifold} is used in the sense of \cref{def:submanifold}: locally
straightenable by a diffeomorphism of the ambient space.
\begin{enumerate}[label=\textbf{(\alph*)}]
  \item \textbf{True}, and there is nothing to prove beyond intersecting with one more open set.
    Being open in $M$ means $N = M \cap W$ for some open $W \subseteq \mathbb{R}^n$. Fix $p \in
    N$ and take a submanifold chart $\Psi : U_p \to V$ for $M$ at $p$. Replacing $U_p$ by
    $U_p \cap W$, which is still an open neighbourhood of $p$, turns it into a submanifold chart
    for $N$ at $p$ of the same dimension. The definition is local, so shrinking the neighbourhood
    costs nothing.
  \item \textbf{False}, and this is the one that has to be met rather than reasoned about.
    In $\mathbb{R}^2$ take
    \[M := \{(x,0) \mid x > 0\}, \qquad N := \{(0,y) \mid y \in \mathbb{R}\}.\]
    They are disjoint, since every point of $M$ has $x > 0$ and every point of $N$ has $x = 0$,
    and each is a $1$-dimensional submanifold: $M$ is an open subset of the $x$-axis, so
    part \textbf{(a)} applies, and $N$ is a line.

    But $M \cup N$ is not a submanifold at the origin. Every ball $B_\varepsilon(0,0)$ meets it
    in a segment of the $y$-axis together with a half-open segment running right, and deleting
    the origin from that set leaves \emph{three} connected pieces, whereas deleting a point from
    a straightened $(\mathbb{R}\times\{0\})\cap V$ leaves two. No diffeomorphism can change the
    number of pieces, so no submanifold chart exists at $(0,0)$.

    What goes wrong is that disjointness is not enough: the two pieces are disjoint but one of
    them accumulates on the other. Requiring instead that $M$ and $N$ have disjoint
    \emph{closures} makes the statement true, and then each point of the union has a
    neighbourhood meeting only one of the two.
  \item \textbf{True}, and it is the chart definition again. Fix $(p,q) \in M \times N$ and take
    submanifold charts $\Psi : U_p \to V$ for $M$ and $\Xi : U_q \to W$ for $N$, with
    $\Psi(U_p \cap M) = (\mathbb{R}^k\times\{0\})\cap V$ and
    $\Xi(U_q \cap N) = (\mathbb{R}^l\times\{0\})\cap W$. Then $\Psi\times\Xi$ is a
    diffeomorphism of the open set $U_p\times U_q \subseteq \mathbb{R}^{m+n}$ onto $V\times W$,
    and it carries $(U_p\times U_q) \cap (M\times N)$ onto
    \[\big((\mathbb{R}^k\times\{0\})\cap V\big) \times \big((\mathbb{R}^l\times\{0\})\cap W\big).\]
    Composing with the permutation of coordinates that gathers the $k$ and the $l$ live
    directions together, which is linear and invertible and hence a diffeomorphism, turns this
    into $(\mathbb{R}^{k+l}\times\{0\}) \cap (V\times W)$. That is a submanifold chart of
    dimension $k+l$.
  \item \textbf{True}, and the hypothesis on the tangent spaces is exactly what is needed. Fix
    $p \in M \cap N$. Each of $M$ and $N$ is locally a level set, and \cref{def:submanifold}
    supplies this directly: a submanifold chart $\Psi : U \to V$ for $M$ at $p$ satisfies
    $\Psi(U \cap M) = (\mathbb{R}^{n-1}\times\{0\}) \cap V$, so its last component
    $F := \Psi_n$ is $C^1$ with $M \cap U = F^{-1}\{0\}$, and $\nabla F(p) \neq 0$ because it is
    a row of the invertible matrix $\Jac\Psi(p)$. Do the same for $N$, shrink to a common
    neighbourhood, and write $G$ for its function. So
    \[M \cap U = F^{-1}\{0\}, \quad \nabla F(p) \neq 0, \qquad
      N \cap U = G^{-1}\{0\}, \quad \nabla G(p) \neq 0 .\]
    By \cref{prop:gradient_orthogonality_level_sets}, $T_pM$ has $\nabla F(p)$ as its
    normal line, and likewise for $N$. Two hyperplanes in $\mathbb{R}^n$ coincide precisely when
    their normal lines do, so $T_pM \neq T_pN$ says exactly that $\nabla F(p)$ and $\nabla G(p)$
    are linearly independent.

    That makes $H := (F,G) : U \to \mathbb{R}^2$ a map whose Jacobian at $p$ has the two
    independent rows $\nabla F(p)\transp$ and $\nabla G(p)\transp$, hence rank $2$, so $DH_p$ is
    surjective. Full rank is an open condition, so after shrinking $U$ it holds throughout, and
    \cref{thm:regular_value_theorem} makes
    \[H^{-1}\{0\} = (M\cap U) \cap (N \cap U) = (M\cap N)\cap U\]
    a submanifold of dimension $n-2$. As $p$ was arbitrary, $M \cap N$ is one.
\end{enumerate}

\begin{remark}[Transversality is the name of the hypothesis in \textbf{(d)}]
Two submanifolds are said to meet \newterm{transversally} at $p$ if their tangent spaces together
span the ambient space, $T_pM + T_pN = \mathbb{R}^n$. For two hyperplanes that is the same as
being distinct, which is why part \textbf{(d)} can state the condition the way it does, and the
general theorem says that a transversal intersection is a submanifold of dimension
$\dim M + \dim N - n$. Here $(n-1)+(n-1)-n = n-2$, as found.

The hypothesis is doing real work, and dropping it breaks the conclusion immediately: two spheres
in $\mathbb{R}^3$ meeting at a single point of tangency have the same tangent plane there, and
their intersection is a point, of dimension $0$ rather than $1$. Two coincident planes are worse
still, intersecting in a plane. Transversality is what rules both out, and the reason it is the
right condition rather than a convenient one is visible in the proof: it is precisely the
statement that stacking the two constraints gives a map of full rank.
\end{remark}
\end{exercisesolution}

\begin{exercisesolution}[Solution to \cref{ex:level_set_of_a_wide_jacobian}]
% Generator: Claude Opus 5 (High) --- the six statements are the examiner's; the working is
% written for these notes. Verdicts checked afterwards against exercises_2024/mocksol.pdf and
% agree on all six.
Everything rests on two computations. First, $\gamma(0) = (0,0,1)$, and differentiating
componentwise,
\[\gamma'(t) = \left(2\cos(2t),\ 1+t,\
  \frac{-\sin t\,(1-t) + \cos t}{(1-t)^2}\right), \qquad\text{so}\qquad
  \gamma'(0) = \begin{pmatrix} 2 \\ 1 \\ 1 \end{pmatrix}.\]
Second, $\Jac F(0,0,1)$ has rank $2$: its leftmost $2\times2$ minor is
$\left(\begin{smallmatrix} 1 & 1 \\ 0 & 1\end{smallmatrix}\right)$, of determinant $1$. Full rank
is an open condition, since that minor is a continuous function of the point, so $\Jac F$ has
rank $2$ on a whole neighbourhood $U$ of $(0,0,1)$, and \cref{thm:regular_value_theorem} applied
to $F|_U$ makes $M \cap U$ a submanifold of dimension $3-2 = 1$. Two independent equations cut
two dimensions out of three and leave a curve; parts \textbf{(b)} and \textbf{(c)} follow at
once.
\begin{enumerate}[label=\textbf{(\alph*)}]
  \item \textbf{False.} By the chain rule,
    \[(F\circ\gamma)'(0) = \Jac F(\gamma(0))\,\gamma'(0)
      = \begin{pmatrix} 1 & 1 & -2 \\ 0 & 1 & -1\end{pmatrix}
        \begin{pmatrix} 2 \\ 1 \\ 1 \end{pmatrix}
      = \begin{pmatrix} 2+1-2 \\ 1-1 \end{pmatrix}
      = \begin{pmatrix} 1 \\ 0 \end{pmatrix} \neq \begin{pmatrix} 0 \\ 0\end{pmatrix}.\]
    This is worth reading as information rather than as a failed computation: had $\gamma$ run
    inside $M$, then $F\circ\gamma$ would be constant and the derivative would vanish. So
    $\gamma$ meets $M$ at $t=0$ and immediately leaves it.
  \item \textbf{False:} the dimension is $1$, not $2$.
  \item \textbf{True.}
  \item \textbf{False.} By \cref{def:normal_space} the normal space is
    $N_{(0,0,1)}M = \operatorname{span}\{(1,1,-2)\transp,\ (0,1,-1)\transp\}$, the span of the
    rows of $\Jac F(0,0,1)$. If $\gamma'(0) = a(1,1,-2)\transp + b(0,1,-1)\transp$, the first
    coordinate forces $a = 2$ and the second then forces $b = -1$; but the third coordinate of
    that combination is $-2\cdot2 - (-1) = -3$, not $1$. Note that $\gamma'(0)$ is not tangent
    either: part \textbf{(a)} says exactly that it is not in $\ker\Jac F(0,0,1) = T_{(0,0,1)}M$.
    It is simply a vector in neither subspace, which is what a curve leaving $M$ produces.
  \item \textbf{True.} To solve for $x$ and $z$ in terms of $y$, the block of $\Jac F(0,0,1)$ in
    the $x$- and $z$-columns must be invertible, and
    \[\det\begin{pmatrix} 1 & -2 \\ 0 & -1\end{pmatrix} = -1 \neq 0 .\]
    \Cref{thm:implicit_function_theorem} therefore supplies $\delta > 0$ and $C^1$ functions
    $\varphi,\psi$ on $(-\delta,\delta)$ with $\varphi(0) = 0$, $\psi(0) = 1$ and
    $F(\varphi(y),y,\psi(y)) = (2,-2)\transp$ throughout. This exhibits the curve of part
    \textbf{(c)} concretely, as a graph over the $y$-axis.
  \item \textbf{True.} Complete $F$ to a map of $\mathbb{R}^3$ to itself by appending the
    coordinate that part \textbf{(e)} solved \emph{for}, namely $y$:
    \[\Phi(x,y,z) := \big(F_1(x,y,z),\ F_2(x,y,z),\ y\big), \qquad
      \Jac\Phi(0,0,1) = \begin{pmatrix} 1 & 1 & -2 \\ 0 & 1 & -1 \\ 0 & 1 & 0\end{pmatrix},\]
    whose determinant is $1 \cdot (1\cdot 0 - (-1)\cdot 1) = 1 \neq 0$. By
    \cref{thm:inverse_function_theorem}, $\Phi$ is a $C^1$ diffeomorphism of some
    $B_\varepsilon((0,0,1))$ onto an open subset of $\mathbb{R}^3$. Writing $\pi(u,v,w) := (u,v)$
    for the projection, we have $F = \pi\circ\Phi$, and $\pi$ maps open sets to open sets, since
    $\pi(B_r(p)) = B_r(\pi(p))$. Hence $F(B_\varepsilon) = \pi(\Phi(B_\varepsilon))$ is open.
\end{enumerate}

\begin{ainote}
The examination calls $M$ a \qt{smooth manifold} in parts \textbf{(b)} and \textbf{(c)}, but it
assumes only $F \in C^1$, so what the regular value theorem delivers is a $C^1$ submanifold. The
distinction changes neither verdict, since both parts are about the dimension.
\end{ainote}
\end{exercisesolution}

\begin{exercisesolution}[Solution to \cref{ex:box_normal_to_a_graph_in_r4}]
% Generator: Claude Opus 5 (High)
Write points of $\mathbb{R}^4$ as $(x,x_4)$ with $x \in \mathbb{R}^3$, and present the graph as a
level set: it is $G^{-1}(\{0\})$ for
\[G(x,x_4) := x_4 - \phi(x), \qquad
  \nabla G(x,x_4) = \begin{pmatrix} -\nabla\phi(x) \\ 1 \end{pmatrix} \in \mathbb{R}^4 .\]
This gradient never vanishes, its last entry being $1$, so
\cref{prop:gradient_orthogonality_level_sets} applies: the graph is a $3$-dimensional submanifold
of $\mathbb{R}^4$ and its normal space at $X$ is the line spanned by $\nabla G(X)$. Normalising,
\[N(X) = \frac{1}{\sqrt{1+\lvert\nabla\phi(x)\rvert^2}}
  \begin{pmatrix} -\partial_1\phi(x) \\ -\partial_2\phi(x) \\ -\partial_3\phi(x) \\ 1
  \end{pmatrix},\]
determined up to sign; this is the choice pointing in the direction of increasing $x_4$.

The formula is worth recognising rather than deriving. For $\phi : \mathbb{R}^2 \to \mathbb{R}$
it is the familiar unit normal $(-\partial_1\phi,-\partial_2\phi,1)/\sqrt{1+\lvert\nabla\phi
\rvert^2}$ of a surface in $\mathbb{R}^3$, and for $\phi : \mathbb{R} \to \mathbb{R}$ it is
$(-\phi',1)/\sqrt{1+\phi'^2}$, perpendicular to the tangent $(1,\phi')$ as it should be. Nothing
about the derivation used $n = 3$; presenting a graph as a level set of $x_{n+1} - \phi(x)$ works
in every dimension, which is the point of setting the question one dimension past the picture.

\begin{ainote}[A slip in the official key]
The key gives two displays for this answer. The second, $(-\nabla\phi(x),1)\transp$ divided by
$\sqrt{1+\lvert\nabla\phi(x)\rvert^2}$, is the formula above. The first is not: it lists a fourth
component $-\partial_4\phi(x)$ and normalises by
$\sqrt{1+\partial_1\phi(x)^2+\partial_2\phi(x)^2+\partial_3\phi(x)^2+\partial_4\phi(x)^2}$. Since
$\phi$ is a function of three variables there is no $\partial_4\phi$, and a vector with those
four components plus the entry $1$ would have five. Only the second display is correct.
\end{ainote}
\end{exercisesolution}




% ============================================================
% Chapter 18 --- Geodesics
% Originally: content/week-10.tex, \section{A short introduction to geodesics}
%             (Corsin Week 10). Moved forward from the vector-calculus material it sat
%             in front of: a geodesic is a curve in a submanifold, so it reads far better
%             next to tangent spaces than next to flux integrals.
% ============================================================
\chapter{Geodesics}
\label{ch:geodesics}

% Transition: Claude Opus 5 (Medium)
With submanifolds and their tangent spaces in place, a natural question is what a straight line
becomes once space is curved. This short chapter answers it: a geodesic is a curve of locally
minimal length inside a submanifold, and the condition characterising one is the first genuine
calculus of variations argument in the course --- differentiate the length functional along a
perturbation and require the derivative to vanish for every perturbation.

% Originally: content/week-10.tex, \section{A short introduction to geodesics} (Corsin Week 10)


\section{A short introduction to geodesics}
\label{sec:geodesics}

% Checked 2026-08-10 against lec_notes.pdf (239 pp, 27 July 2026): "geodesic" occurs exactly
% once in the whole file, in Example 9.14 on printed p. 5, and the surrounding pages were
% rendered and read rather than grepped. There is no geodesics section, no length functional
% and no variational argument anywhere in the lecture notes.
\begin{ainote}[This chapter goes beyond the course]
\label{ainote:geodesics_beyond_the_course}
Two warnings before starting, and they point the same way.

Corsin marks this section \qt{not exam relevant}. And the official lecture notes have no
counterpart to it at all: the word \newterm{geodesic} appears in them exactly once, in
Example 9.14 on printed p.\ 5, where the great-circle distance
$d_1(x,y) = R\arccos\big(x\cdot y/R^2\big)$ is offered as a more natural metric on the sphere
than the Euclidean one, with no indication of why that arc is the shortest. Nothing in the
course computes a geodesic, and the calculus of variations does not appear in it.

So this chapter is not lecture material, and nothing in it can be assumed known. It is here
because the calculation is short and because the idea behind it, treating length as a functional
and applying optimization to it, is worth meeting once; and it is here in particular because it
answers the question Example 9.14 raises and then drops.
\end{ainote}

We will try to find the shortest path (i.e., a geodesic) from $0 \in \mathbb{R}^n$ to $x \in \mathbb{R}^n$. This
is, of course, the straight line, but the essentially same calculation can be used to find the
shortest path between points on a sphere or any other submanifold.

\textbf{Idea.} Apply \qt{optimization} to the map
\[L : \mathcal A \to [0,\infty), \qquad c \mapsto L(c) = \int_0^1 \lvert c'(t)\rvert\,dt,
  \qquad \mathcal A := \{\gamma \in C^1([0,1],\mathbb{R}^n) : \gamma(0)=0,\ \gamma(1)=x\}.\]

\begin{lemma}[A few differentiation identities]
\label{lem:inner_product_derivative_identities}
Let $\alpha,\beta : (0,1) \to \mathbb{R}^n$ be $C^1$. Then
\[\frac{d}{dt}\langle\alpha(t),\beta(t)\rangle = \langle\alpha'(t),\beta(t)\rangle
  + \langle\alpha(t),\beta'(t)\rangle, \qquad
  \frac{d}{dt}\lvert\alpha(t)\rvert = \frac{1}{\lvert\alpha(t)\rvert}\langle\alpha(t),\alpha'(t)\rangle.\]
\end{lemma}

% Source: Corsin Nick/Class Notes/Week 10.pdf, pp. 4--5
\begin{proposition}[The shortest path is a straight line]
\label{prop:geodesic_straight_line}
Assume that $c \in C^\infty([0,1],\mathbb{R}^n)$ with $c(0)=0$, $c(1)=x$ and
$\lvert c'(t)\rvert \equiv \lambda$ is the optimal path we are looking for. Then $c(t) = xt$
and $\lambda = \lvert x\rvert$.
\end{proposition}

\begin{proof}
Let $g \in C^\infty([0,1],\mathbb{R}^n)$ with $g(0)=g(1)=0$, and define, for $\varepsilon \in
\mathbb{R}$, the family of paths $c_\varepsilon(t) := c(t) + \varepsilon g(t)$, so that
$c_0 = c$. The boundary conditions on $g$ are what keep the whole family admissible:
$c_\varepsilon(0) = c(0) = 0$ and $c_\varepsilon(1) = c(1) = x$ for every $\varepsilon$, so each
$c_\varepsilon$ is a competitor in $\mathcal{A}$ and not merely a nearby curve.

% Sonnet 5 (Medium) --- FIG-W10-02
\begin{center}
\begin{tikzpicture}[>=Latex, scale=1.2]
  \draw[red, thick] (0,0) -- (3.6,1.1) node[midway, below, font=\scriptsize, sloped] {$c$};
  \foreach \k in {1,2,3} {
    \draw[green!45!black, thick, opacity={0.5+0.15*\k}]
      plot[smooth] coordinates {(0,0) (1.1,{0.55+0.13*\k}) (2.3,{0.75+0.10*\k}) (3.6,1.1)};
  }
  \node[font=\scriptsize] at (2.0,1.55) {$c_\varepsilon$};
  \fill (0,0) circle (1.3pt) node[left, font=\scriptsize] {$0$};
  \fill (3.6,1.1) circle (1.3pt) node[right, font=\scriptsize] {$x$};
\end{tikzpicture}
\end{center}

By assumption, $c$ is a minimizer of $L$, so it is a critical point of the family, i.e.
\[\left.\frac{d}{d\varepsilon}\right|_{\varepsilon=0} L(c_\varepsilon) = 0.\]
(This equality itself requires a proof, but the idea is exactly the same as that of
optimization in $\mathbb{R}^n$.) First,
\[\left.\frac{\partial}{\partial\varepsilon}\right|_{\varepsilon=0} c_\varepsilon(t)
  = \left.\frac{\partial}{\partial\varepsilon}\right|_{\varepsilon=0}
    \big(c(t)+\varepsilon g(t)\big) = g(t).\]
Recalling $\lvert c'\rvert \equiv \lambda$ and using
\Cref{lem:inner_product_derivative_identities},
\[
\begin{aligned}
\left.\frac{d}{d\varepsilon}\right|_{\varepsilon=0} L(c_\varepsilon)
  &= \left.\frac{d}{d\varepsilon}\int_0^1 \lvert c_\varepsilon'(t)\rvert\,dt\right|_{\varepsilon=0}
   = \int_0^1 \left.\frac{\partial}{\partial\varepsilon}\lvert c_\varepsilon'(t)\rvert\,dt
     \right|_{\varepsilon=0} \\
  &= \int_0^1 \frac{\big\langle c_\varepsilon'(t),\ \partial_\varepsilon c_\varepsilon'(t)
       \big\rangle}{\lvert c_\varepsilon'(t)\rvert}\,dt\bigg|_{\varepsilon=0}
   = \frac1\lambda \int_0^1 \Big\langle c'(t),\ \partial_\varepsilon\big|_{\varepsilon=0}
       c_\varepsilon'(t)\Big\rangle\,dt \\
  &= \frac1\lambda \int_0^1 \frac{d}{dt}\Big\langle c'(t), g(t)\Big\rangle
       - \big\langle c''(t), g(t)\big\rangle\,dt
   = \frac1\lambda\left[\big\langle c'(t),g(t)\big\rangle\Big|_0^1
       - \int_0^1 \big\langle c''(t),g(t)\big\rangle\,dt\right] \\
  &= -\frac1\lambda \int_0^1 \big\langle c''(t),g(t)\big\rangle\,dt \ \overset{!}{=}\ 0,
\end{aligned}
\]
where the boundary term vanished since $g(0)=g(1)=0$. Since $g$ is arbitrary, it follows that
$c''(t) \equiv 0$. This last step is the \newterm{fundamental lemma of the calculus of
variations}: a continuous function integrating to zero against \emph{every} smooth $g$ vanishing
at the endpoints must itself vanish, since otherwise one could choose $g$ to be a bump
concentrated where $c''$ has a fixed sign and make the integral non-zero. Together with
$c(0)=0$ and $c(1)=x$ this gives $c(t)=xt$ and $\lambda = \lvert x\rvert$.
\end{proof}

% Corsin flags both of these steps in his own proof with a warning triangle but does not resolve
% either. (Was an ainote until 2026-08-09; by Test 1 in style.md it is mathematics, and the
% duplicate \newterm of the fundamental lemma it carried is now left only at its first
% occurrence, inside the proof above.)
\begin{remark}[The two steps that hide real work]
\label{rem:geodesic_hidden_steps}
The proof moves quickly over $\left.\frac{d}{d\varepsilon}\right|_{\varepsilon=0}L(c_\varepsilon)
= 0$ and over the final $\overset{!}{=}\,0$, and those are precisely the two places where
something is being assumed.

The first is the claim that a minimizer of $L$ over the infinite-dimensional set $\mathcal A$ is
a critical point of every one-parameter family through it. That is the same fact used for
optimization on $\mathbb{R}^n$ in \cref{prop:extremum_implies_critical}, and it wants $\mathcal A$
to be a vector space. It is not one, since $\mathcal A$ is an affine subspace. What rescues the
argument is that the \emph{differences} $c_\varepsilon - c$ range over
$\{g \in C^\infty([0,1],\mathbb{R}^n) : g(0)=g(1)=0\}$, which is a genuine vector space, so the
perturbation directions form one even though the curves themselves do not.

The second is the fundamental lemma named in the proof, which upgrades \qt{the integral vanishes
for every $g$} to \qt{$c''$ vanishes at every $t$}, rather than merely on average. Corsin gives
neither its name nor its proof.
\end{remark}

% Generator: Claude Opus 5 (High) --- the hypothesis $\lvert c'\rvert \equiv \lambda$ is used
% twice in the computation (to pull $1/\lambda$ out of the integral) and is never commented on.
\begin{remark}[Why constant speed is assumed, and why it costs nothing]
\label{rem:geodesic_constant_speed}
\Cref{prop:geodesic_straight_line} assumes $\lvert c'(t)\rvert \equiv \lambda$, which looks like
a restriction on the competitor curves and is not one. Length is invariant under
reparametrization, since substituting $t = \varphi(s)$ for an increasing $C^1$ bijection
$\varphi$ leaves $\int\lvert c'\rvert\,dt$ unchanged. Every curve with nowhere-vanishing speed
therefore has a constant-speed representative tracing the same image, obtained by running along
it proportionally to arc length. The hypothesis chooses that representative.

It is worth choosing, because the functional $L$ cannot see the difference and the equation can.
Without it the computation still goes through, but $1/\lvert c'(t)\rvert$ stays inside the
integral and the conclusion is that $\big(c'/\lvert c'\rvert\big)' \equiv 0$, which says the
\emph{direction} of travel is constant while the speed is free. Constant speed converts that
into the cleaner $c'' \equiv 0$. \Cref{ex:ai_variation_without_constant_speed} carries this out.
\end{remark}

\begin{aiexercise}[The fundamental lemma of the calculus of variations]
% Generator: Claude Opus 5 (High)
\label{ex:ai_fundamental_lemma}
The proof of \cref{prop:geodesic_straight_line} passes from \qt{the integral vanishes for every
$g$} to \qt{$c''$ vanishes at every $t$} in one step. Prove that step. Precisely: let
$h \in C^0\big([0,1],\mathbb{R}^n\big)$ satisfy
\[\int_0^1 \big\langle h(t),\,g(t)\big\rangle\,dt = 0
  \qquad\text{for every } g \in C^\infty\big([0,1],\mathbb{R}^n\big) \text{ with } g(0)=g(1)=0.\]
Show that $h \equiv 0$.

\exhint{Argue by contradiction, and use a bump function of the kind constructed in
\cref{def:smooth_bump_function} to concentrate $g$ where $h$ is known not to vanish.}
\end{aiexercise}

\begin{aiexercise}[The variation without the constant-speed hypothesis]
% Generator: Claude Opus 5 (High)
\label{ex:ai_variation_without_constant_speed}
Repeat the computation in the proof of \cref{prop:geodesic_straight_line} for a minimiser
$c \in C^2\big([0,1],\mathbb{R}^n\big)$ that is \emph{not} assumed to have constant speed, only
$\lvert c'(t)\rvert > 0$ for every $t$.
\begin{enumerate}[label=\textbf{(\alph*)}]
  \item Show that $\left.\frac{d}{d\varepsilon}\right|_{\varepsilon=0}L(c_\varepsilon)
    = \int_0^1 \big\langle c'(t)/\lvert c'(t)\rvert,\ g'(t)\big\rangle\,dt$, and integrate by
    parts to bring it into the form $-\int_0^1\langle \cdot,\,g(t)\rangle\,dt$.
  \item Conclude, using \cref{ex:ai_fundamental_lemma}, that
    $\big(c'/\lvert c'\rvert\big)' \equiv 0$.
  \item Deduce that the image of $c$ is still the straight segment from $0$ to $x$, and that
    $L(c) = \lvert x\rvert$. Which feature of $c$ does the argument now fail to pin down, and
    why is that unavoidable?
\end{enumerate}
\end{aiexercise}



\part{Integration}

% Generator: Claude Opus 5 (High) --- part-level framing prose, new content.
Integration in several variables raises a question that one variable never had to face. Before
asking what $\int_A f$ is, one has to say which sets $A$ possess a volume at all, and Jordan
measure is the answer this course uses. With that settled, Fubini's theorem reduces a multiple
integral to iterated one-dimensional ones, and substitution generalizes to the change of variables
formula, in which the factor that appears is $\lvert\det \Jac\phi\rvert$, the local volume
distortion of the substitution. The final chapter leaves flat regions behind altogether. The Gram
determinant measures how much a parametrization stretches length, area or volume, and with it one
can integrate over a curve or a surface sitting inside $\mathbb{R}^n$ rather than over a piece of
$\mathbb{R}^n$ itself.

% ============================================================
% Chapter 19 --- Jordan measure and the Riemann integral
% Originally: content/week-08.tex, \section{Multivariate integration} (Corsin Week 8),
%             with problems 8.2 and 8.4 from sheet 8.
% ============================================================
\chapter{Jordan Measure and the Riemann Integral}
\label{ch:jordan_measure}

% Transition: Claude Opus 5 (Medium)
Integration in several variables needs a notion of the volume of a domain before it can define an
integral over one. This chapter sets up Jordan measure and the multidimensional Riemann integral
built on it, and then spends most of its effort on the sets that are \emph{not} Jordan
measurable. That is where every counterexample on the problem sheet lives, and the reason they
exist is a single one: a Jordan cover may use only finitely many cubes, all of one size. A short
section extends the integral to unbounded domains, which every later chapter uses and none of them
defines. The chapter closes with a digression on the Lebesgue theory, which repairs that single
defect properly; it is stated without proofs and used nowhere afterwards.

% Originally: content/week-08.tex, \section{Multivariate integration} (Corsin Week 8)

% Source: Corsin Nick/Class Notes/Week 8.pdf, p. 7
\section{Multivariate integration}
\label{sec:multivariate_integration}

% Supplement: Sascha Brack/Class Notes/Week_08_Notes_Monday.pdf, pp. 5--7
\subsection{Jordan measure and the Riemann integral}
The rigorous definition of the multidimensional integral relies on the \newterm{Jordan measure} (\germanterm{Jordan-Maß}). The intuition, as presented in Sascha Brack's notes, involves \qt{dyadic cubes} (pixels) that approximate a region from the inside and outside.

\begin{center}
\begin{tikzpicture}[scale=0.8]
  % first grid
  \draw[gray, dotted] (0,0) grid (4,2);
  \draw[very thick, blue!80!black] (0,0) rectangle (2,2);
  \node[below, font=\scriptsize] at (1,0) {dyadic cube, side length $1=2^{-0}$};
  
  \begin{scope}[shift={(6,0)}]
    \draw[gray, dotted, step=0.5] (0,0) grid (4,2);
    \draw[very thick, blue!80!black] (0,0) rectangle (0.5,0.5);
    \draw[very thick, blue!80!black] (1,0.5) rectangle (1.5,1.0);
    \draw[very thick, blue!80!black] (2,0) rectangle (2.5,0.5);
    \draw[very thick, blue!80!black] (2.5,0) rectangle (3.0,0.5);
    \draw[very thick, blue!80!black] (3,1.5) rectangle (3.5,2.0);
    \node[below, font=\scriptsize, align=center] at (2,-0.2) {dyadic set\\(union of same size cubes)};
  \end{scope}
\end{tikzpicture}
\end{center}

% Generator: Claude Opus 5 (High) --- the notion was used throughout this chapter and in
% \cref{ch:change_of_variables} but had no numbered definition to point at.
\begin{definition}[Dyadic cube and dyadic set]
\label{def:dyadic_set}
For $p \in \mathbb{N}$ and $a \in \mathbb{Z}^n$, the \newterm{dyadic cube}
(\germanterm{dyadischer Würfel}) of generation $p$ at $a$ is
\[Q_p(a) := 2^{-p}\big(a + [0,1)^n\big),\]
a half-open cube of side length $2^{-p}$. A \newterm{dyadic set} (\germanterm{dyadische Menge})
is a finite union $E = \bigcup_{\ell=1}^N Q_p(a(\ell))$ of distinct dyadic cubes \emph{of one
common generation} $p$, and its volume is
\[\mu(E) := 2^{-pn}N.\]
\end{definition}

The requirement that all cubes share a generation is not cosmetic: it is what separates the
Jordan theory from the Lebesgue theory, as \cref{rem:jordan_vs_lebesgue_sizes} makes precise.

% Supplement: Sascha Brack/Class Notes/Week_08_Notes_Monday.pdf, pp. 5--7
Similar to the 1D Riemann integral, we approximate the volume of a set $E \subseteq \mathbb{R}^n$ using an outer sequence of dyadic sets $G_i \supseteq E$ and an inner sequence $F_i \subseteq E$. If the limits of their volumes coincide as the pixel size $2^{-p} \to 0$, $E$ is \newterm{Jordan measurable}.

Furthermore, the multidimensional integral of a non-negative function $f : A \to \mathbb{R}$ is exactly the $(n+1)$-dimensional Jordan measure (volume) of its hypograph $\Gamma_A(f)$:

\begin{center}
\begin{tikzpicture}[scale=1.2]
  % axes
  \draw[->, thick] (0,0,0) -- (3,0,0) node[right] {$y$};
  \draw[->, thick] (0,0,0) -- (0,2.5,0) node[above] {$z = f(x,y)$};
  \draw[->, thick] (0,0,0) -- (0,0,3) node[below left] {$x$};
  
  % base domain A
  \draw[thick, dashed, fill=yellow!20] (0.5,0,0.5) -- (2.5,0,0.5) -- (2.0,0,2.5) -- (0.2,0,2.0) -- cycle;
  \node[font=\small] at (1.5, 0, 1.5) {$A$};
  
  % top surface graph(f)
  \draw[thick, fill=yellow!40, opacity=0.8] (0.5,1.2,0.5) .. controls (1.5,1.8,0.5) .. (2.5,1.5,0.5)
    -- (2.0,0.8,2.5) .. controls (1.0,1.2,2.5) .. (0.2,1.0,2.0) -- cycle;
  \node[font=\small] at (1.5, 1.4, 1.5) {$\Gamma_A(f)$};
  
  % vertical lines
  \draw[dashed] (0.5,0,0.5) -- (0.5,1.2,0.5);
  \draw[dashed] (2.5,0,0.5) -- (2.5,1.5,0.5);
  \draw[dashed] (2.0,0,2.5) -- (2.0,0.8,2.5);
  \draw[dashed] (0.2,0,2.0) -- (0.2,1.0,2.0);
  
  \node[right, font=\small, align=center] at (2.5, 0.7, 0.5) {volume is\\$\int_A f$};
\end{tikzpicture}
\end{center}

% Source: Jérôme Paschoud/Notizen/Woche 8.2 Riemann-Integral.pdf, p. 1
% Generator: Gemini 3.1 Pro (High)
\begin{definition}[Riemann integral in $\mathbb{R}^n$]
\label{def:riemann_integral_rn}
Let $A \subseteq \mathbb{R}^n$ and $f: A \to \mathbb{R}$. Analogous to Analysis I, the integral is defined via dyadic step functions (\germanterm{dyadische Treppenfunktionen}). The lower and upper integrals are defined as
\begin{align*}
  \underline{\int}_A f &:= \sup \left\{ \int_A g \;\middle|\; g: \mathbb{R}^n \to \mathbb{R} \text{ dyadic step function with } g \leq f \right\}, \\
  \overline{\int}_A f &:= \inf \left\{ \int_A h \;\middle|\; h: \mathbb{R}^n \to \mathbb{R} \text{ dyadic step function with } h \geq f \right\}.
\end{align*}
The function $f$ is \newterm{Riemann integrable} if $\underline{\int}_A f = \overline{\int}_A f$. In this case, their common value is denoted by $\int_A f$.

\begin{center}
\begin{tikzpicture}[scale=0.9]
  % axes
  \draw[->, gray] (0,0,0) -- (3.5,0,0) node[right, text=black] {$y$};
  \draw[->, gray] (0,0,0) -- (0,2.5,0) node[above, text=black] {$z = g(x,y)$};
  \draw[->, gray] (0,0,0) -- (0,0,3.5) node[below left, text=black] {$x$};
  
  % Base squares
  \draw[gray, dashed] (0.5,0,0.5) grid (2.5,0,2.5);
  
  % Column 1
  \draw[thick, blue!80!black, fill=blue!10, opacity=0.8] (0.5,0,1.5) -- (0.5,1.2,1.5) -- (1.5,1.2,1.5) -- (1.5,0,1.5) -- cycle;
  \draw[thick, blue!80!black, fill=blue!20, opacity=0.8] (1.5,0,1.5) -- (1.5,1.2,1.5) -- (1.5,1.2,2.5) -- (1.5,0,2.5) -- cycle;
  \draw[thick, blue!80!black, fill=blue!30, opacity=0.8] (0.5,1.2,1.5) -- (1.5,1.2,1.5) -- (1.5,1.2,2.5) -- (0.5,1.2,2.5) -- cycle;
  
  % Column 2
  \draw[thick, red!80!black, fill=red!10, opacity=0.8] (1.5,0,0.5) -- (1.5,0.8,0.5) -- (2.5,0.8,0.5) -- (2.5,0,0.5) -- cycle;
  \draw[thick, red!80!black, fill=red!20, opacity=0.8] (2.5,0,0.5) -- (2.5,0.8,0.5) -- (2.5,0.8,1.5) -- (2.5,0,1.5) -- cycle;
  \draw[thick, red!80!black, fill=red!30, opacity=0.8] (1.5,0.8,0.5) -- (2.5,0.8,0.5) -- (2.5,0.8,1.5) -- (1.5,0.8,1.5) -- cycle;

  \node[font=\small, align=left] at (5, 1, 1) {Integral of a\\dyadic step function:\\sum of column volumes};
\end{tikzpicture}
\end{center}

This integral satisfies all the \qt{good} properties from the one-dimensional case, such as linearity. Moreover, if $E \subseteq \mathbb{R}^n$ is Jordan measurable and $f: \overline{E} \to [0, \infty)$ is continuous, then $f$ is Riemann integrable on $E$.
\end{definition}

% Source: Diego Torres Tejeda/Notes - 17.04 - Diego Torres Tejeda.pdf, p. 1
% Generator: Gemini 3.6 Flash (High)
\begin{theorem}[Jordan measurability criterion via boundary]
\label{thm:jordan_measurability_boundary_criterion}
A bounded subset $E \subseteq \mathbb{R}^n$ is Jordan measurable if and only if its boundary $\partial E$ is a Jordan null set, i.e.\ $\mu_{\mathrm{out}}(\partial E) = 0$.
\end{theorem}

\begin{lemma}[Compact sets: Jordan null versus Lebesgue null]
\label{lem:compact_jordan_vs_lebesgue_null}
Let $K \subseteq \mathbb{R}^n$ be compact. Then $K$ is a Jordan null set if and only if $K$ is a Lebesgue null set.
\end{lemma}

\begin{proof}
If $K$ is Jordan null, then for every $\varepsilon > 0$, $K$ is covered by finitely many dyadic cubes $Q_1, \dots, Q_m$ with total volume $\sum_{i=1}^m \mu(Q_i) < \varepsilon$, so $K$ is a Lebesgue null set. Conversely, if $K$ is Lebesgue null, there is a countable cover by dyadic cubes $\{Q_j\}_{j \ge 1}$ with $\sum_{j=1}^\infty \mu(Q_j) < \varepsilon$. Taking the interiors of slightly enlarged cubes yields an open cover of the compact set $K$, which admits a finite subcover of total volume $< 2\varepsilon$, proving $K$ is Jordan null.
\end{proof}

% Source: Diego Torres Tejeda/Notes - 17.04 - Diego Torres Tejeda.pdf, pp. 3--4
% Generator: Gemini 3.6 Flash (High)
\begin{example}[The Cantor set is a Jordan null set]
\label{ex:cantor_set_jordan_null}
Consider the middle-thirds \newterm{Cantor set} $C \subseteq [0,1]$, defined inductively by $C_0 := [0,1]$ and
\[
C_{n+1} := \tfrac{1}{3} C_n \cup \left( \tfrac{2}{3} + \tfrac{1}{3} C_n \right), \qquad C := \bigcap_{n=0}^\infty C_n.
\]
Each $C_n$ consists of $2^n$ closed intervals of length $3^{-n}$, so $\mu(C_n) = (2/3)^n$. Since $C \subseteq C_n$ for all $n$, we have
\[
\mu_{\mathrm{out}}(C) \leq \mu(C_n) = \left(\tfrac{2}{3}\right)^n \xrightarrow{n\to\infty} 0,
\]
proving that the Cantor set $C$ is a Jordan null set.

Moreover, as Diego Torres Tejeda highlights in his notes, $C$ consists precisely of those points $x \in [0,1]$ possessing a ternary (base-3) expansion $x = \sum_{n=1}^\infty \frac{a_n}{3^n}$ composed entirely of digits $a_n \in \{0, 2\}$ without the digit $1$.
\end{example}

% Source: Diego Torres Tejeda/Notes - 17.04 - Diego Torres Tejeda.pdf, p. 5
% Generator: Gemini 3.6 Flash (High)
\begin{example}[An open set that is not Jordan measurable]
\label{ex:non_jordan_measurable_open_set}
Jordan measurability can fail even for open sets. Enumerate the rational numbers in $[0,1]$ as $\mathbb{Q} \cap [0,1] = \{q_1, q_2, \dots\}$, and define the open set
\[
U := \bigcup_{k=1}^\infty B_{2^{-(k+2)}}(q_k) \;\cap\; (0,1).
\]
Since $\mathbb{Q} \cap [0,1]$ is dense in $[0,1]$, the closure of $U$ is $\overline{U} = [0,1]$. If $U$ were Jordan measurable, its measure would satisfy $\mu(U) = \mu(\overline{U}) = 1$. However, any finite union of intervals covering a dyadic inner approximation of $U$ has measure bounded by $\sum_{k=1}^\infty 2 \cdot 2^{-(k+2)} = \frac{1}{2}$. Thus $\mu_{\mathrm{in}}(U) \leq \frac{1}{2} < 1 = \mu_{\mathrm{out}}(U)$, showing that $U$ is \textbf{not} Jordan measurable. Its boundary $\partial U = [0,1] \setminus U$ is a \qt{fat} boundary of positive outer measure.
\end{example}

% Source: Jérôme Paschoud/Notizen/Woche 8.1 Tangentialräume und Jordan Mass.pdf, p. 3
% Generator: Gemini 3.1 Pro (High)
\begin{proposition}[Algebra of Jordan measurable sets]
\label{prop:algebra_jordan_measurable}
If $E, F \subset \mathbb{R}^n$ are Jordan measurable, then their union $E \cup F$, intersection $E \cap F$, and differences $E \setminus F$, $F \setminus E$ are also Jordan measurable.
Furthermore, if $E \cap F = \emptyset$ (or if they only intersect in a Jordan null set), then
\[
\mu(E \cup F) = \mu(E) + \mu(F).
\]
\end{proposition}

% Source: Jérôme Paschoud/Notizen/Woche 8.1 Tangentialräume und Jordan Mass.pdf, p. 4
% Generator: Gemini 3.1 Pro (High)
\begin{remark}[Lebesgue null vs.\ Jordan null]
\label{rem:lebesgue_vs_jordan_null}
A set $E \subset \mathbb{R}^n$ is \newterm{Lebesgue null} if for every $\varepsilon > 0$, there exist \emph{countably} many dyadic cubes $Q_i$ covering $E$ with $\sum \mu(Q_i) < \varepsilon$. For a \newterm{Jordan null} set, the cover must be formed by \emph{finitely} many cubes (which implicitly requires $E$ to be bounded).
Because a countable cover is allowed, the set $\mathbb{Q} \cap [0,1]$ is Lebesgue null. However, it is not Jordan null (in fact, it is not even Jordan measurable, as seen in \cref{ex:non_jordan_measurable_open_set}), since any finite union of intervals covering it must also cover its closure $[0,1]$, thus having measure at least $1$.
\end{remark}

% Quelle: Linus Lüchinger/Slides/Slides-04-16.pdf, slide 6
% Generator: Claude Opus 5 (High)
\begin{remark}[The relaxation is really about \emph{sizes}, not just counts]
\label{rem:jordan_vs_lebesgue_sizes}
Stating the difference as \qt{finitely many} against \qt{countably many} is correct but hides
where the extra power actually comes from. A Jordan cover is a \emph{dyadic set}, and every cube
in a dyadic set of generation $p$ has the same side length $2^{-p}$
(\cref{def:dyadic_set}). So the two conditions really read:
\begin{itemize}
  \item \textbf{Jordan null:} for each $\varepsilon>0$, cover $E$ by cubes that all have
    \emph{one common side length}, of total volume $<\varepsilon$. Finiteness is then automatic:
    only finitely many cubes of a fixed size fit into a bounded region.
  \item \textbf{Lebesgue null:} for each $\varepsilon>0$, cover $E$ by cubes of \emph{whatever
    side lengths you like}, of total volume $<\varepsilon$. Allowing the sizes to vary is exactly
    what lets the cover be infinite while the total volume stays finite.
\end{itemize}
The two relaxations are therefore one relaxation. It is permission to shrink the cubes as you go
that buys the countable cover, and it is the countable cover that makes the notions differ.

$\mathbb{N} \subseteq \mathbb{R}$ shows this in its purest form. Cover the integer $k$ by an
interval of length $\varepsilon\,2^{-k-1}$; the total length is $\varepsilon$, so $\mathbb{N}$ is
Lebesgue null. No Jordan cover can do this: cubes of one fixed length $2^{-p}$ need at least one
cube per integer, infinitely many of them, of infinite total length. (Indeed $\mathbb{N}$ is
unbounded, so it is not Jordan measurable at all.) The point is that the geometric series is
available only when the sizes may vary.
\end{remark}

% Source: Jérôme Paschoud/Notizen/Woche 8.1 Tangentialräume und Jordan Mass.pdf, p. 4
% Generator: Gemini 3.1 Pro (High)
\begin{theorem}[Linear transformations of Jordan measurable sets]
\label{thm:jordan_measure_linear_transformation}
Let $L: \mathbb{R}^n \to \mathbb{R}^n$ be a linear, invertible map. If $E \subset \mathbb{R}^n$ is Jordan measurable, then the image $L(E)$ is also Jordan measurable, and its measure scales by the absolute value of the determinant of $L$:
\[
\mu(L(E)) = |\det L| \, \mu(E).
\]
\end{theorem}


% --- exercises relocated here from the dissolved sheet 8 ---

% Originally: content/week-08.tex, \exercisesheet{8}, problem 8.2
% Source: exercises/Ex8_Analysis2_eng.pdf

% Source: exercises/Ex8_Analysis2_eng.pdf, p. 1
\begin{exercise}[8.2 --- True or False \textnormal{(important)}]
\label{ex:8.2}
\begin{enumerate}[label=\textbf{\arabic*.}]
  \item A bounded countable set is always Jordan-null.
  \item A countable set is always Lebesgue-null.
  \item Let $D \subset [0,1]$ be a dense set (i.e.\ $\overline D = [0,1]$). Then
    $\mu_{out}(D) = 1$. ($\mu_{out}$ was defined in Definition 13.10.)
  \item Let $X, Y \subset [0,1]$ Jordan measurable sets such that $\mu(X) > 1/2$ and
    $\mu(Y) > 1/2$. Then $X \cap Y \neq \emptyset$.
  \item Let $X, Y \subset [0,1]$ such that $\mu_{out}(X) > 1/2$ and $\mu_{out}(Y) > 1/2$. Then
    $X \cap Y \neq \emptyset$.
\end{enumerate}
\textit{Official hint:} \textbf{8.2.5} --- $\mu_{out}$ can be positive and large on very sparse
sets\dots

% Supplement: Sascha Brack/Ex Sheet Hints/Ex8_Analysis2_hints.pdf, p. 1
\textit{Sascha Brack's hints:} a toolkit of sets worth testing every part against ---
$\mathbb{Q}$, the irrationals $\mathbb{R}\setminus\mathbb{Q}$, $\mathbb{Q}\cap[0,1]$, and
$(\mathbb{R}\setminus\mathbb{Q})\cap[0,1]$. For \textbf{8.2.3} specifically: monotonicity already
gives $\mu_{out}(D) \leq \mu_{out}([0,1]) = 1$, so the work is showing
$\mu_{out}(D) \geq 1$ --- essentially, that the \qt{smallest} dyadic set containing $D$ also
contains $[0,1]$.
\end{exercise}

% Originally: content/week-08.tex, \exercisesheet{8}, problem 8.4
% Source: exercises/Ex8_Analysis2_eng.pdf

% Source: exercises/Ex8_Analysis2_eng.pdf, p. 1
\begin{exercise}[8.4 --- Multiple Choice \textnormal{(important)}]
\label{ex:8.4}
Let $U \subset \mathbb{R}^n$ be a nonempty, open subset, $f : U \to \mathbb{R}^m$ a function,
and $N \subset U$ a Jordan null set. In which of the following cases is the image
$f(N) \subset \mathbb{R}^m$ necessarily a null set? \textit{Attention: Only one answer is
correct!}
\begin{enumerate}[label=\textbf{\arabic*.}]
  \item If $f$ is uniformly continuous.
  \item If $f$ is uniformly continuous and $m \geq n$.
  \item If $f$ is locally Lipschitz continuous.
  \item If $f$ is locally Lipschitz continuous and $m \geq n$.
\end{enumerate}
\end{exercise}

% Originally: new section (no predecessor in the week-based skeleton, and none in any tutor's
% notes). Added 2026-08-10 after a pass over lec_notes.pdf found that this document uses improper
% multidimensional integrals in at least four places without ever defining one:
%   * ex:improper_integral_polar (ch. 20) states outright that "its Riemann improper integral is
%     well-defined", which is a claim about a definition we did not have;
%   * ex:gaussian_integral_polar (ch. 20) integrates over all of R^2;
%   * ex:ai_frullani (ch. 20) integrates over the half-line;
%   * ex:10.3 (ch. 21), Gabriel's horn, integrates over [1, infinity).
% The lecture notes devote two pages to it and no tutor covers it, so this is a genuine gap
% rather than a deliberate skip. Corsin never reaches it either.

% Supplement: lec_notes.pdf, pp. 124--125
\section{Improper integrals}
\label{sec:improper_integrals}

Everything so far has integrated a continuous function over a \emph{bounded} Jordan measurable
set. Two situations fall outside that and both occur constantly in practice: the domain may be
unbounded, and the integrand may be unbounded near the edge of its domain. What should
$\int_{\mathbb{R}^2}f$ mean?

The answer avoids limits entirely, at the cost of one hypothesis.

\begin{definition}[Improper integral]
\label{def:improper_integral}
Let $U \subseteq \mathbb{R}^n$ be open and let $f \in C^0(U)$ be \textbf{non-negative}. The
\newterm{improper integral} \germanfor{uneigentliches Integral} of $f$ over $U$, still written
$\int_U f$, is
\[\int_U f := \sup\left\{\int_K f \ \middle|\ K \subseteq U \text{ compact and Jordan
  measurable}\right\} \ \in\ [0,+\infty].\]
\end{definition}

Each $\int_K f$ on the right is an ordinary integral: $K$ is compact and Jordan measurable and
$f$ is continuous on it, so \cref{def:riemann_integral_rn} applies. The supremum of a set of
real numbers bounded below by $0$ always exists in $[0,+\infty]$, so $\int_U f$ is always
defined. It may be $+\infty$, and in that case one says the improper integral \newterm{diverges}.

\begin{remark}[Why non-negativity, and what it buys]
\label{rem:improper_needs_nonnegativity}
The hypothesis $f \geq 0$ is not a technicality that a more careful definition would remove. It
is what makes the supremum monotone, so that enlarging $K$ can only increase $\int_K f$, and
that monotonicity is the whole of \cref{lem:improper_integral_as_limit} below.

Drop it and the value stops being well defined. For a sign-changing $f$ one can arrange
exhaustions that pick up the positive part faster or slower than the negative part and obtain
different answers, exactly as a conditionally convergent series can be rearranged to any sum.
\Cref{ex:feynman_dirichlet_integral} is the case in point: $\int_0^\infty (\sin x)/x\,dx = \pi/2$
is \emph{not} an improper integral in the sense above, since $(\sin x)/x$ changes sign infinitely
often and $\int_0^\infty \lvert\sin x\rvert/x\,dx = +\infty$. Its value depends on integrating in
the order $\int_0^R$ and then letting $R \to \infty$, which is why that example needs a
justification the others do not.
\end{remark}

The definition is clean but useless for computing, since one cannot take a supremum over all
compact subsets by hand. The next result says that any single sensible exhaustion suffices, and
that it does not matter which one is chosen.

\begin{lemma}[Improper integrals as limits]
\label{lem:improper_integral_as_limit}
Let $U \subseteq \mathbb{R}^n$ be open and $f \in C^0(U)$ non-negative. Let
$K_1 \subseteq K_2 \subseteq K_3 \subseteq \cdots$ be compact Jordan measurable sets with
\[\bigcup_{\ell=1}^\infty \operatorname{int}(K_\ell) = U.\]
Then
\[\int_U f = \lim_{\ell\to\infty}\int_{K_\ell} f,\]
and in particular the limit exists in $[0,+\infty]$.
\end{lemma}

\begin{proof}
Since the $K_\ell$ are nested and $f \geq 0$, the sequence $\int_{K_\ell}f$ is non-decreasing, so
it converges in $[0,+\infty]$; call the limit $\Lambda$.

\textbf{$\int_U f \leq \Lambda$.} Let $K \subseteq U$ be compact and Jordan measurable. The sets
$\operatorname{int}(K_\ell)$ are open and cover $U$, hence cover $K$, so by compactness finitely
many of them suffice; since they are nested, a single one does, say
$K \subseteq \operatorname{int}(K_{\ell_0}) \subseteq K_{\ell_0}$. Non-negativity then gives
$\int_K f \leq \int_{K_{\ell_0}} f \leq \Lambda$. Taking the supremum over all such $K$ yields
$\int_U f \leq \Lambda$.

\textbf{$\Lambda \leq \int_U f$.} Each $K_\ell$ is itself compact and Jordan measurable and
contained in $U$, so it competes in the supremum defining $\int_U f$. Consequently
$\int_{K_\ell} f \leq \int_U f$ for every $\ell$, and letting $\ell \to \infty$ gives
$\Lambda \leq \int_U f$.
\end{proof}

\begin{remark}[What the lemma licenses]
\label{rem:improper_exhaustion_free_choice}
Read the two halves together: the value does not depend on the exhaustion, so one may choose
whichever exhaustion makes the computation easy. That is the licence
\cref{ex:improper_integral_polar} and \cref{ex:gaussian_integral_polar} were already using when
they integrated over $\mathbb{R}^2$ in polar coordinates, and it is what
\cref{ex:ai_gaussian_two_exhaustions} makes explicit by running two different exhaustions against
each other and comparing the answers.
\end{remark}

\begin{aiexercise}[The Gaussian integral, two exhaustions at once]
% Generator: Claude Opus 5 (High)
\label{ex:ai_gaussian_two_exhaustions}
Let $f(x,y) := e^{-(x^2+y^2)}$ on $U := \mathbb{R}^2$, which is continuous and positive, and put
$I := \int_{-\infty}^{\infty} e^{-t^2}\,dt$.
\begin{enumerate}[label=\textbf{(\alph*)}]
  \item Check that the closed disks $D_\ell := \overline{B_\ell(0)}$ and the closed squares
    $Q_\ell := [-\ell,\ell]^2$ each form an exhaustion of $\mathbb{R}^2$ in the sense of
    \cref{lem:improper_integral_as_limit}.
  \item Compute $\lim_{\ell\to\infty}\int_{D_\ell} f$ in polar coordinates.
  \item Compute $\lim_{\ell\to\infty}\int_{Q_\ell} f$ using \cref{thm:fubini}, leaving the answer
    in terms of $I$.
  \item Conclude the value of $I$, and say which step would collapse if $f$ were allowed to
    change sign.
\end{enumerate}
\end{aiexercise}

% Originally: written for this chapter; no predecessor week file.
% Generator: Claude Opus 5 (High)
% Added 2026-08-11 on the user's explicit instruction, after they asked whether Jordan measure and
% Lebesgue measure are the same thing and then asked for a short digression section on Lebesgue
% theory. This is a deliberate, user-authorised exception to the "enrich existing sections only"
% scope rule in gemini.md, which is recorded there. Nothing here is proved, nothing here is used
% by any later chapter, and the section says both of those things to the reader.

\section{A digression: a glimpse of the Lebesgue integral}
\label{sec:lebesgue_digression}

This course builds one theory of integration: the Riemann integral over Jordan measurable sets. It
is enough for everything in the chapters that follow, and no result later in this document depends
on a single line of this section. Read it as an answer to a question the previous section provokes
and cannot settle.

The question is what \qt{Lebesgue-null} means, given that this course never defines a Lebesgue
integral, and yet the term appears on the problem sheets (\cref{ex:8.2}\textbf{(b)}) and in the
comparison remarks above. It comes from a second theory of integration, and the sanest way to place
it is to see what that theory fixes. \textbf{Nothing below is proved.} The statements are given so
that the vocabulary has somewhere to attach, and so that a reader meeting measure theory later
recognises what it was built for.

\subsection{Two things the Riemann theory cannot do}

\begin{example}[The measurable sets are too few]
\label{ex:lebesgue_defect_sets}
\Cref{prop:algebra_jordan_measurable} says the Jordan measurable sets are closed under finite
unions, intersections and differences. They are \emph{not} closed under countable unions. Each
singleton $\{q\} \subseteq \mathbb{R}$ is Jordan measurable with measure $0$, and yet
\[\mathbb{Q}\cap[0,1] = \bigcup_{k\ge1}\{q_k\}\]
is not Jordan measurable at all: any finite cover of it by cubes also covers its closure $[0,1]$,
so $\mu_{\mathrm{out}} = 1$ while $\mu_{\mathrm{in}} = 0$. \Cref{ex:non_jordan_measurable_open_set}
is worse still, since the set there is \emph{open} and the volume of an open set ought not to be a
question a theory declines to answer.
\end{example}

\begin{example}[Limits escape the class of integrable functions]
\label{ex:lebesgue_defect_limits}
Enumerate $\mathbb{Q}\cap[0,1] = \{q_1,q_2,\dots\}$ and let $f_k$ be the indicator function of
$\{q_1,\dots,q_k\}$. Each $f_k$ is Riemann integrable, having only finitely many discontinuities,
and $\int_0^1 f_k = 0$. The sequence increases pointwise to the Dirichlet function
\[f(x) = \begin{cases} 1, & x \in \mathbb{Q}, \\ 0, & x \notin \mathbb{Q},\end{cases}\]
which is not Riemann integrable, its upper and lower sums being $1$ and $0$ for every partition.

A bounded increasing sequence of integrable functions therefore need not have an integrable limit.
The only exchange theorem available in this course needs \emph{uniform} convergence, and uniform
convergence is a strong hypothesis that the sequences arising in practice routinely fail.
\end{example}

Both defects have the same shape. The Riemann--Jordan theory is built on finite covers by cubes of
one size, and it is stable under finite operations and unstable under countable ones. Lebesgue's
construction changes exactly that.

\subsection{Lebesgue measure}

\begin{definition}[Lebesgue outer measure]
\label{def:lebesgue_outer_measure}
For $E \subseteq \mathbb{R}^n$ the \newterm{Lebesgue outer measure}
\germanfor{äusseres Lebesgue-Mass} is
\[\lambda^*(E) := \inf\left\{ \sum_{i=1}^\infty \vol(Q_i) \;\middle|\;
E \subseteq \bigcup_{i=1}^\infty Q_i,\ Q_i \text{ cubes} \right\},\]
the infimum running over all \emph{countable} covers of $E$ by cubes of arbitrary sizes.
\end{definition}

This is the notion already in use above: $E$ is Lebesgue null exactly when $\lambda^*(E) = 0$, which
is \cref{rem:lebesgue_vs_jordan_null} restated. The defect of $\lambda^*$ is that it is only
subadditive, not additive, so one restricts it to a good class of sets.

\begin{definition}[Lebesgue measurable set]
\label{def:lebesgue_measurable}
A set $E \subseteq \mathbb{R}^n$ is \newterm{Lebesgue measurable}
\germanfor{Lebesgue-messbar} if it splits every test set additively:
\[\lambda^*(A) = \lambda^*(A\cap E) + \lambda^*(A\setminus E) \qquad
\text{for every } A \subseteq \mathbb{R}^n.\]
On this class one writes $\lambda := \lambda^*$ and calls it \newterm{Lebesgue measure}
\germanfor{Lebesgue-Mass}.
\end{definition}

\begin{theorem}[What the construction delivers, without proof]
\label{thm:lebesgue_measure_properties}
The Lebesgue measurable subsets of $\mathbb{R}^n$ form a \newterm{$\sigma$-algebra}
\germanfor{$\sigma$-Algebra}: a family containing $\emptyset$ and closed under complements and
under \emph{countable} unions. It contains every open set, every closed set and every Lebesgue null
set. On it, $\lambda$ is countably additive: for pairwise disjoint measurable $E_1,E_2,\dots$,
\[\lambda\left(\bigcup_{k\ge1}E_k\right) = \sum_{k\ge1}\lambda(E_k).\]
Moreover every Jordan measurable set $E$ is Lebesgue measurable with $\lambda(E) = \mu(E)$, and the
inclusion is strict.
\end{theorem}

\begin{remark}[The answer to the question, in one line]
\label{rem:jordan_is_not_lebesgue}
Jordan measure and Lebesgue measure are \emph{not} the same, and the difference is not a matter of
convention. Lebesgue measure is a strict extension: it agrees with Jordan measure wherever the
latter is defined, so no volume computed anywhere in this document is put at risk, and it is
defined on many more sets. The single structural upgrade is from an algebra to a
$\sigma$-algebra, that is, from finite to countable stability, and both defects of
\cref{ex:lebesgue_defect_sets,ex:lebesgue_defect_limits} dissolve in it.
\end{remark}

\subsection{The integral: partition the range, not the domain}

The Riemann integral chops the \emph{domain} into small pieces and hopes the function is nearly
constant on each. Lebesgue chops the \emph{range} instead, and asks how large the set is on which
the function takes each value. The standard image: to count a heap of coins, a Riemann integrator
picks them up in the order they lie and adds as they come, while a Lebesgue integrator first sorts
them into denominations and multiplies. The second survives a much messier heap.

Concretely, a \newterm{simple function} is a finite combination $s = \sum_{i=1}^N c_i \chi_{E_i}$
with $c_i \ge 0$ and $E_i$ measurable, where $\chi_E$ is the indicator function of $E$; its
integral is declared to be $\int s \,d\lambda := \sum_i c_i \lambda(E_i)$. For measurable
$f \ge 0$ one then sets
\[\int f \,d\lambda := \sup\left\{ \int s\,d\lambda \;\middle|\; s \text{ simple},\ 0 \le s \le f
\right\},\]
and for general $f$ one splits $f = f^+ - f^-$ into positive and negative parts, calling $f$
\newterm{Lebesgue integrable} if $\int\lvert f\rvert\,d\lambda < \infty$.

The Dirichlet function of \cref{ex:lebesgue_defect_limits} is now trivial: it is the indicator
function of $\mathbb{Q}\cap[0,1]$, a null set, so its integral is $0$. Its Riemann sums never
settle because they are forced to look at intervals; its Lebesgue integral is immediate because the
set carrying the value $1$ is negligible.

\subsection{What the extra work buys}

\begin{theorem}[The convergence theorems, without proof]
\label{thm:lebesgue_convergence_theorems}
Let $f_k$ be measurable functions on a measurable $E \subseteq \mathbb{R}^n$.
\begin{enumerate}[label=\textbf{(\alph*)}]
  \item \textbf{Monotone convergence.} If $0 \le f_1 \le f_2 \le \dots$ pointwise and
    $f_k \to f$ pointwise, then $\int_E f_k\,d\lambda \to \int_E f\,d\lambda$.
  \item \textbf{Fatou's lemma.} If $f_k \ge 0$, then
    $\int_E \liminf_k f_k \,d\lambda \le \liminf_k \int_E f_k\,d\lambda$.
  \item \textbf{Dominated convergence.} If $f_k \to f$ pointwise and there is a single integrable
    $g$ with $\lvert f_k\rvert \le g$ for all $k$, then $f$ is integrable and
    $\int_E f_k\,d\lambda \to \int_E f\,d\lambda$.
\end{enumerate}
\end{theorem}

\begin{remark}[Why these are the point of the whole construction]
\label{rem:why_dominated_convergence_matters}
Compare the hypotheses. Ours is uniform convergence, a statement about $\sup_x\lvert f_k(x)-f(x)
\rvert$, which fails as soon as the functions develop a moving spike. Dominated convergence asks
only for pointwise convergence plus one integrable envelope, and the envelope is usually the easy
part. Monotone convergence asks for nothing but monotonicity.

This is where the cost of staying inside the Riemann theory is actually paid in this document.
Uniform convergence is assumed by hand in \cref{sec:ascoli_arzela}, in the treatment of improper
integrals in \cref{sec:improper_integrals}, and again when differentiating under the integral sign
in \cref{sec:feynmans_trick}, where a domination hypothesis is what the argument morally wants.
Each of those becomes shorter and more general in the Lebesgue setting.
\end{remark}

\subsection{How the two integrals compare}

\begin{theorem}[Lebesgue's criterion, and the comparison, without proof]
\label{thm:lebesgue_criterion_riemann}
Let $Q \subseteq \mathbb{R}^n$ be a closed box and $f : Q \to \mathbb{R}$ bounded.
\begin{enumerate}[label=\textbf{(\alph*)}]
  \item $f$ is Riemann integrable if and only if its set of discontinuities is Lebesgue null.
  \item If $f$ is Riemann integrable, then it is Lebesgue integrable and the two integrals agree.
\end{enumerate}
\end{theorem}

Part \textbf{(a)} is the reason a course with no Lebesgue integral still has to say
\qt{Lebesgue null}. It is the sharp characterisation of Riemann integrability, and no notion built
from finite covers replaces it: the Dirichlet function is discontinuous everywhere, a set of
measure $1$, and it fails; a function with countably many jumps is discontinuous on a null set, and
it succeeds.

\begin{ainote}
It is tempting to conclude that the Lebesgue integral simply subsumes the Riemann integral. For
\emph{proper} integrals on a box, \cref{thm:lebesgue_criterion_riemann}\textbf{(b)} says exactly
that. For \emph{improper} ones it is false in both directions, and the standard witness is
\[\int_0^\infty \frac{\sin x}{x}\,dx = \frac{\pi}{2},\]
which exists as an improper Riemann integral (\cref{def:improper_integral}) but not as a Lebesgue
integral, because $\int_0^\infty \lvert \sin x\rvert/x \,dx = \infty$ and the Lebesgue theory
integrates only absolutely. The improper Riemann integral is a limit of integrals and can exploit
cancellation between the humps; the Lebesgue integral is defined through $\lvert f\rvert$ and
cannot. So neither theory contains the other outright, and \qt{Lebesgue is strictly stronger} is a
half-truth worth not repeating.
\end{ainote}

\begin{remark}[What to take away]
\label{rem:lebesgue_digression_summary}
Three sentences are enough to carry. Jordan measure is Lebesgue measure restricted to a smaller,
finitely-stable class of sets, and where both are defined they agree. The Lebesgue theory exists
because that class is too small to be closed under limits, and its payoff is the convergence
theorems rather than any new volume. And the one place its vocabulary genuinely reaches into this
course is Lebesgue's criterion, which is why \qt{Lebesgue null} had to be defined here at all.
\end{remark}

% Generator: Claude Opus 5 (High)
% First use of the `aside' environment, added to main.tex in the same pass. This is the intended
% shape: a genuinely startling fact that no later argument touches.
\begin{aside}[Some sets have no volume at all, and a ball can be doubled]
\label{aside:vitali_banach_tarski}
\Cref{thm:lebesgue_measure_properties} carefully says which sets $\lambda$ is defined on, and never
claims that is all of them. It cannot: there exist subsets of $\mathbb{R}$ that are not Lebesgue
measurable. Vitali's construction takes the quotient $\mathbb{R}/\mathbb{Q}$ and picks one
representative from each class inside $[0,1]$. The resulting set $V$ has countably many disjoint
rational translates that together fill up an interval, so countable additivity would force
$\lambda(V)$ to be simultaneously $0$ and positive. No value works, and $V$ is unmeasurable.

The construction needs the axiom of choice, and this is not a technicality. It is consistent with
the other axioms of set theory that \emph{every} subset of $\mathbb{R}$ is Lebesgue measurable, so
the pathology is imported by the choice principle rather than discovered in the reals.

Pushed into $\mathbb{R}^3$, the same idea gives the Banach--Tarski paradox: a solid ball can be cut
into five pieces which, moved only by rotations and translations, reassemble into two solid balls
of the original size. Nothing is stretched, and no volume is created, because the five pieces are
unmeasurable and simply do not have volumes to add up. The paradox is a statement about the axiom
of choice, not about geometry.
\end{aside}

\subsection{Exercises}

% Generator: Claude Opus 5 (High)
\begin{aiexercise}[Null sets from countability]
\label{ex:ai_lebesgue_null_countable}
Work directly from \cref{def:lebesgue_outer_measure}.
\begin{enumerate}[label=\textbf{(\alph*)}]
  \item Show that a single point of $\mathbb{R}^n$ is Lebesgue null.
  \item Show that every countable subset of $\mathbb{R}^n$ is Lebesgue null, and deduce that
    $\mathbb{Q}^n$ is.
  \item Show that a Lebesgue null set has empty interior. (Compare
    \cref{ex:lebesgue_null_sets_properties}\textbf{(a)}, which is the same statement contraposed.)
\end{enumerate}
\end{aiexercise}

% Generator: Claude Opus 5 (High)
\begin{aiexercise}[An integral the Riemann theory cannot see]
\label{ex:ai_lebesgue_three_valued}
Define $f : [0,1] \to \mathbb{R}$ by
\[f(x) := \begin{cases}
1, & x \in \mathbb{Q}, \\
2, & x \notin \mathbb{Q} \text{ and } x < \tfrac12, \\
3, & x \notin \mathbb{Q} \text{ and } x \ge \tfrac12.
\end{cases}\]
\begin{enumerate}[label=\textbf{(\alph*)}]
  \item Compute $\int_{[0,1]} f \,d\lambda$.
  \item Determine the set of points at which $f$ is continuous, and conclude from
    \cref{thm:lebesgue_criterion_riemann} that $f$ is not Riemann integrable.
\end{enumerate}
\end{aiexercise}

% Generator: Claude Opus 5 (High)
\begin{aiexercise}[Dominated convergence where uniform convergence fails]
\label{ex:ai_dominated_beats_uniform}
For $k \ge 1$ let $f_k : [0,1] \to \mathbb{R}$, $f_k(x) := x^k$.
\begin{enumerate}[label=\textbf{(\alph*)}]
  \item Determine the pointwise limit $f$ of $(f_k)$ and show that the convergence is not uniform.
  \item Exhibit an integrable dominating function and apply
    \cref{thm:lebesgue_convergence_theorems}\textbf{(c)} to compute
    $\lim_{k\to\infty}\int_{[0,1]} f_k\,d\lambda$.
  \item Confirm the answer by evaluating $\int_0^1 x^k\,dx$ directly.
\end{enumerate}
\end{aiexercise}

% Generator: Claude Opus 5 (High)
\begin{aiexercise}[Improper Riemann but not Lebesgue]
\label{ex:ai_improper_not_lebesgue}
Define $f : [1,\infty) \to \mathbb{R}$ by $f(x) := (-1)^{k+1}/k$ for $x \in [k,k+1)$,
$k \in \mathbb{N}_{\ge1}$.
\begin{enumerate}[label=\textbf{(\alph*)}]
  \item Show that the improper integral $\int_1^\infty f(x)\,dx$ converges, and identify its value.
  \item Show that $\int_1^\infty \lvert f(x)\rvert\,dx = \infty$, so that $f$ is not Lebesgue
    integrable.
  \item Why does this not contradict \cref{thm:lebesgue_criterion_riemann}\textbf{(b)}?
\end{enumerate}
\end{aiexercise}

% Originally: content/week-08.tex had a \section{Solutions} holding nothing but a placeholder
% note, which had already gone stale, since the statements were in fact transcribed. So 8.2 and
% 8.4 stood in the document unsolved. Written here for the first time.
\section{Solutions}
\label{sec:jordan_measure_solutions}

% Transition: Claude Opus 5 (High)
Not every solution in this chapter is collected here. \Cref{ex:inclusion_exclusion_dyadic} is
solved where it appears, on \cpageref{ex:inclusion_exclusion_dyadic}.

\begin{exercisesolution}[Solution to \cref{ex:8.2}]
% Generator: Claude Opus 5 (High)
The whole problem turns on one distinction: \textbf{Jordan} outer measure is an infimum over
\emph{finite} covers, \textbf{Lebesgue} outer measure over \emph{countable} ones. Parts
\textbf{(a)}--\textbf{(c)} are that distinction three times; parts \textbf{(d)} and \textbf{(e)}
are the difference between a measure and a merely sub-additive outer measure.
\begin{enumerate}[label=\textbf{(\alph*)}]
  \item \textbf{False.} Take $X := \mathbb{Q}\cap[0,1]$: bounded and countable, yet
    $\mu_{out}(X) = 1$ by part \textbf{(c)}, since $X$ is dense in $[0,1]$. Countability buys
    nothing here, since only finitely many intervals are allowed at a time, and finitely many
    intervals covering a dense set must cover the whole of $[0,1]$.
  \item \textbf{True}, and this is exactly where countable covers do buy something. Enumerate
    $X = \{x_1,x_2,\dots\}$ and, given $\varepsilon > 0$, cover $x_k$ by an interval of length
    $\varepsilon 2^{-k}$. The total length is $\sum_{k\geq1}\varepsilon2^{-k} = \varepsilon$, and
    $\varepsilon$ was arbitrary, so the Lebesgue outer measure is $0$.
  \item \textbf{True.} Monotonicity gives $\mu_{out}(D) \leq \mu_{out}([0,1]) = 1$, so only the
    lower bound needs an argument. Let $I_1,\dots,I_N$ be finitely many intervals covering $D$.
    Enlarging each to its closure changes no length, so we may assume they are closed; then
    $\bigcup_{k=1}^N I_k$ is \emph{closed}, which is where finiteness enters, and contains
    $D$, hence contains $\overline D = [0,1]$. Therefore
    $\sum_{k=1}^N |I_k| \geq \mu([0,1]) = 1$. Taking the infimum over such covers,
    $\mu_{out}(D) \geq 1$.
  \item \textbf{True.} Suppose $X \cap Y = \emptyset$. Jordan measure is additive on disjoint
    Jordan measurable sets, so
    \[\mu(X \cup Y) = \mu(X) + \mu(Y) > \tfrac12 + \tfrac12 = 1.\]
    But $X \cup Y \subseteq [0,1]$, so monotonicity forces $\mu(X\cup Y) \leq 1$, a
    contradiction. Hence $X \cap Y \neq \emptyset$.
  \item \textbf{False}, and this is the official hint's point. Take
    \[X := \mathbb{Q}\cap[0,1], \qquad Y := (\mathbb{R}\setminus\mathbb{Q})\cap[0,1].\]
    Both are dense in $[0,1]$, so part \textbf{(c)} gives $\mu_{out}(X) = \mu_{out}(Y) = 1 >
    \tfrac12$; yet $X \cap Y = \emptyset$ by construction. The argument of part \textbf{(d)}
    breaks at its first step:
    $\mu_{out}$ is only \emph{sub}additive, so $\mu_{out}(X\cup Y) \leq \mu_{out}(X)+\mu_{out}(Y)$,
    an inequality pointing the wrong way to produce a contradiction. Neither set is Jordan
    measurable, which is precisely why part \textbf{(d)} does not apply.
\end{enumerate}

\begin{ainote}
Parts \textbf{(a)} and \textbf{(e)} are the same counterexample used twice, and part \textbf{(c)}
is the lemma both rest on. It is
worth extracting what makes it work: a \emph{finite} union of closed sets is closed, so a finite
cover of $D$ automatically covers $\overline D$. Every difference between Jordan and Lebesgue
theory visible on this sheet traces back to that one sentence.
\end{ainote}
\end{exercisesolution}

\begin{exercisesolution}[Solution to \cref{ex:8.4}]
% Generator: Claude Opus 5 (High)
\textbf{The correct answer is 4: $f$ locally Lipschitz continuous and $m \geq n$.} Both
hypotheses are needed, and dropping either one admits a counterexample.

\textbf{Why 1 and 2 fail:} uniform continuity is not enough. Let $c : [0,1] \to [0,1]$ be the
Cantor function, which is continuous on a compact interval and hence uniformly continuous, and
let $N \subset [0,1]$ be the Cantor set. Then $N$ is Jordan-null (it is covered by $2^k$
intervals of total length $(2/3)^k \to 0$), while $c(N) = [0,1]$ is not null. Here $n = m = 1$, so
this single example kills option 1 \emph{and} option 2. The extra hypothesis $m \geq n$ does
nothing to rescue mere continuity. Intuitively, a uniformly continuous map has no bound on how
much it may stretch, only on how abruptly.

\textbf{Why 3 fails:} $m \geq n$ is not decoration. Let $f : \mathbb{R}^2 \to \mathbb{R}$,
$f(x,y) := y$, which is Lipschitz with constant $1$, and let
$N := \{0\}\times[0,1]$, a segment and so Jordan-null in $\mathbb{R}^2$. Then $f(N) = [0,1]$,
which is not null in $\mathbb{R}^1$. Here $m = 1 < 2 = n$: a projection loses a dimension, and
what was negligible in the source stops being negligible in the smaller target.

\textbf{Why 4 holds.} Let $N$ be Jordan-null. Being Jordan-null it is bounded, and we may work on
a compact neighbourhood $K \subset U$ of $\overline N$ on which $f$ is Lipschitz with some
constant $L$ (local Lipschitz continuity plus compactness gives a single constant). Fix
$\varepsilon > 0$ and cover $N$ by finitely many cubes $Q_1,\dots,Q_M \subset K$ of sides
$s_i \leq 1$ with $\sum_i s_i^{\,n} < \varepsilon$. Each $Q_i$ has diameter $s_i\sqrt n$, so
$f(Q_i \cap N)$ has diameter at most $L s_i\sqrt n$ and therefore fits inside a cube of side
$L s_i \sqrt n$ in $\mathbb{R}^m$, of volume $(L\sqrt n)^m s_i^{\,m}$. Since $s_i \leq 1$ and
$m \geq n$ we have $s_i^{\,m} \leq s_i^{\,n}$, so the images are covered with total volume at
most
\[(L\sqrt n)^m \sum_{i=1}^M s_i^{\,n} \;<\; (L\sqrt n)^m\,\varepsilon.\]
As $\varepsilon$ was arbitrary, $f(N)$ is null.

\begin{ainote}
The proof shows exactly where each hypothesis is spent. \emph{Lipschitz} is what converts a bound
on the side of a cube into a bound on the diameter of its image. Continuity alone gives no
such conversion. \emph{$m \geq n$} is what makes $s_i^{\,m} \leq s_i^{\,n}$, i.e.\ what stops the
volume from being measured with too few powers of $s_i$. The two counterexamples attack precisely
these two steps.
\end{ainote}
\end{exercisesolution}

\begin{exercisesolution}[Solution to \cref{ex:ai_gaussian_two_exhaustions}]
% Generator: Claude Opus 5 (High)
\begin{enumerate}[label=\textbf{(\alph*)}]
  \item Both families are nested and consist of compact sets whose boundaries are a circle
    and the four sides of a square respectively, each a Jordan null set, so both are Jordan
    measurable by \cref{thm:jordan_measurability_boundary_criterion}. For the covering condition,
    $\operatorname{int}(D_\ell) = B_\ell(0)$ and $\operatorname{int}(Q_\ell) = (-\ell,\ell)^2$,
    and every point of $\mathbb{R}^2$ lies in both once $\ell$ exceeds its distance from the
    origin. Hence $\bigcup_\ell\operatorname{int}(D_\ell) =
    \bigcup_\ell\operatorname{int}(Q_\ell) = \mathbb{R}^2$, and
    \cref{lem:improper_integral_as_limit} applies to each.
  \item In polar coordinates, using \cref{thm:change_of_variables} with Jacobian $r$,
    \[\int_{D_\ell} e^{-(x^2+y^2)}\,dx\,dy
      = \int_0^{2\pi}\!\!\int_0^{\ell} e^{-r^2}\,r\,dr\,d\varphi
      = 2\pi\left[-\frac{e^{-r^2}}{2}\right]_0^{\ell} = \pi\big(1 - e^{-\ell^2}\big)
      \ \xrightarrow[\ell\to\infty]{}\ \pi.\]
  \item On a square the integrand factorises, so \cref{thm:fubini} splits it:
    \[\int_{Q_\ell} e^{-(x^2+y^2)}\,dx\,dy
      = \left(\int_{-\ell}^{\ell} e^{-x^2}dx\right)\left(\int_{-\ell}^{\ell} e^{-y^2}dy\right)
      = \left(\int_{-\ell}^{\ell} e^{-x^2}dx\right)^{\!2}
      \ \xrightarrow[\ell\to\infty]{}\ I^2.\]
  \item Both limits compute the same number $\int_{\mathbb{R}^2}f$, by
    \cref{lem:improper_integral_as_limit}, so $I^2 = \pi$ and $I = \sqrt\pi$, the positive root
    since the integrand is positive.
\end{enumerate}

The step that collapses without $f \geq 0$ is the identification in \textbf{(d)}. Non-negativity
is what makes $\ell \mapsto \int_{K_\ell}f$ monotone, and monotonicity is what forces every
exhaustion to the same limit. For a sign-changing integrand the disks and the squares are free to
disagree, and \cref{rem:improper_needs_nonnegativity} gives the standard example of an integral
whose value genuinely depends on how the domain is exhausted.

\begin{remark}[This is the argument \cref{ex:gaussian_integral_polar} was already making]
\label{rem:gaussian_two_exhaustions_retro}
The classical derivation of $\int_{\mathbb{R}}e^{-x^2}dx = \sqrt\pi$ squares the integral, reads
the product as an integral over $\mathbb{R}^2$, and switches to polar coordinates. Written out,
that is precisely the comparison above: the squaring step is the square exhaustion and the polar
step is the disk exhaustion, and the derivation works because the two agree. What
\cref{lem:improper_integral_as_limit} adds is the permission to move between them, which the
classical presentation uses without mentioning.
\end{remark}
\end{exercisesolution}

\begin{exercisesolution}[Proof of \cref{ex:examHS24_TF_improper}]
% Generator: Gemini 3.1 Pro (High)
Both improper integrals can be evaluated by switching to polar coordinates, where $f(r\cos\theta, r\sin\theta) = r^{-2\alpha}$ and the area element is $r\,dr\,d\theta$.
\begin{enumerate}[label=\textbf{(\alph*)}]
  \item \textbf{False:} Over the region $A = \{r \ge 1\}$, the integral becomes $2\pi \int_1^\infty r^{1-2\alpha} \, dr$. The integral of $r^p$ on $[1,\infty)$ converges if and only if $p < -1$. Thus, we need $1-2\alpha < -1$, which simplifies to $2\alpha > 2$, or $\alpha > 1$. The statement claims it converges if and only if $\alpha \le 1$, which is the exact opposite.
  \item \textbf{False:} Over the region $B = \{r \le 1\}$, the integral becomes $2\pi \int_0^1 r^{1-2\alpha} \, dr$. The integral of $r^p$ on $(0,1]$ converges if and only if $p > -1$. This requires $1-2\alpha > -1$, meaning $\alpha < 1$. For $\alpha = 10/9 > 1$, the integral diverges.
  \item \textbf{False:} As established above, for $\alpha < 1$, the integral evaluates to
  \[ \int_B f = 2\pi \left[ \frac{r^{2-2\alpha}}{2-2\alpha} \right]_0^1 = \frac{2\pi}{2-2\alpha} = \frac{\pi}{1-\alpha}. \]
  Taking the limit as $\alpha \to 0^+$, we get $\lim_{\alpha \to 0^+} \frac{\pi}{1-\alpha} = \pi$. The statement incorrectly claims the limit is $\frac{\pi^2}{6}$.
\end{enumerate}
\end{exercisesolution}

\begin{exercisesolution}[Solution of \cref{ex:improper_cos_over_r4}]
% Generator: Claude Opus 5 (High)
The integrand is undefined at the origin, so $\lvert f\rvert$ is a continuous non-negative function
on the open set $\mathbb{R}^2\setminus\{(0,0)\}$ and \cref{def:improper_integral} applies there.
The two boundary circles are Jordan null sets, so it makes no difference whether $A$ and $B$ are
read as closed or open. Throughout, write $r := \sqrt{x^2+y^2}$, so that $\lvert f\rvert =
\lvert\cos(xy)\rvert/r^4$, and recall from \cref{rem:polar_as_shells} that a radial bound
integrates as $2\pi\int g(r)\,r\,dr$.

Nothing here can be computed exactly, because $\cos(xy)$ has no antiderivative in polar
coordinates. Both verdicts come from bounding that numerator, once from above and once from below.
\begin{enumerate}[label=\textbf{(\alph*)}]
  \item \textbf{True.} On all of $\mathbb{R}^2$, $\lvert\cos(xy)\rvert \leq 1$, so on $A$
    \[\int_A \lvert f\rvert \;\leq\; \int_A \frac{1}{r^4}
      \;=\; 2\pi\int_1^\infty \frac{1}{r^4}\,r\,dr
      \;=\; 2\pi\int_1^\infty r^{-3}\,dr
      \;=\; 2\pi\left[\frac{-1}{2r^2}\right]_1^\infty = \pi \;<\; \infty .\]
  \item \textbf{False.} Here the numerator must be bounded \emph{below}, and on $B$ it is bounded
    below by a positive constant. Indeed $\lvert xy\rvert \leq \tfrac12(x^2+y^2) \leq \tfrac12$ for
    $(x,y) \in B$ by the inequality of arithmetic and geometric means, and $\cos$ is positive and
    decreasing on $[0,\tfrac12] \subseteq [0,\pi/2)$, so $\cos(xy) \geq \cos\tfrac12 > 0$
    throughout $B$. Consequently
    \[\int_B \lvert f\rvert \;\geq\; \cos\tfrac12 \int_B \frac{1}{r^4}
      \;=\; 2\pi\cos\tfrac12 \int_0^1 r^{-3}\,dr \;=\; +\infty,\]
    since $\int_0^1 r^p\,dr$ diverges for $p \leq -1$ and here $p = -3$.
  \item \textbf{False.} $B \setminus\{(0,0)\}$ is contained in $\mathbb{R}^2\setminus\{(0,0)\}$, so
    monotonicity of the improper integral in the domain gives
    $\int_{\mathbb{R}^2}\lvert f\rvert \geq \int_B \lvert f\rvert = +\infty$ by \textbf{(b)}. One
    divergent piece is enough; the convergence found in \textbf{(a)} cannot repair it.
\end{enumerate}
\begin{remark}[The singularity is at the origin, not at infinity]
Set against \cref{ex:examHS24_TF_improper}, which asks the same two questions for
$r^{-2\alpha}$, this block is the case $\alpha = 2$: the radial integrand is $r^{-4}\cdot r =
r^{-3}$, which converges at infinity and diverges at $0$, and no exponent does both. That is the
general picture in the plane. On $A$ convergence needs $r^{-2\alpha+1}$ integrable at infinity,
that is $\alpha > 1$; on $B$ it needs the same expression integrable at $0$, that is $\alpha < 1$.
The two conditions are complementary, so $\int_{\mathbb{R}^2}r^{-2\alpha}$ diverges for every
$\alpha$, and the bounded oscillating factor $\cos(xy)$ changes none of it.
\end{remark}
\end{exercisesolution}


% ============================================================
% Chapter 20 --- Change of variables, Fubini and Feynman's trick
% Originally: content/week-08.tex --- \subsection{Change of variables} and
%             \subsection{Fubini's theorem} (both promoted to sections) together with
%             \section{Feynman's trick} (Corsin Week 8), plus problem 8.5 from sheet 8.
% ============================================================
\chapter{Change of Variables, Fubini and Feynman's Trick}
\label{ch:change_of_variables}

% Transition: Claude Opus 5 (Medium)
With the integral defined, the problem becomes computing one, and three tools do nearly all of
that work. Change of variables trades an awkward domain for a convenient one and pays a Jacobian
factor for the trade. Fubini reduces an $n$-dimensional integral to $n$ one-dimensional ones.
Differentiating under the integral sign reaches integrals with no elementary antiderivative at
all, by turning the integral into an ODE in a parameter it never had. A short closing section puts
all three to work at once on the centroid of a region, which is the first quantity in the course
that is defined by an integral rather than computed by one.

% Originally: content/week-08.tex, \subsection{Change of variables} --- promoted to a section

\section{Change of variables}
\label{sec:change_of_variables}

\begin{theorem}[Change of variables]
\label{thm:change_of_variables}
Suppose $U,V \subseteq \mathbb{R}^n$ are open sets and $\Phi : U \to V$ is a
$C^1$-diffeomorphism. If $A \subseteq U$ is Jordan measurable with $\overline{A} \subseteq U$,
and $f : A \to \mathbb{R}$ is continuous, then
\[\int_A f(x)\,dx
  = \int_{\Phi(A)} \frac{f\circ\Phi^{-1}(y)}{\big|\det \Jac\Phi\big(\Phi^{-1}(y)\big)\big|}\,dy.\]
\end{theorem}

\begin{ainote}
Corsin uses \newterm{Jordan measurable} \germanfor{Jordan-messbar} here without defining it;
the definition is in the lecture (ch.~13). For this chapter it is enough to know that boxes,
and the bounded regions cut out by finitely many smooth curves that appear in the examples, are
all Jordan measurable.

On the shape of the formula: it is often quoted the other way round, as
$\int_{\Phi(A)} f = \int_A (f\circ\Phi)\,\lvert\det\Jac\Phi\rvert$. The version above is the same
statement with $\Phi$ pushed to the other side, which is why the determinant appears
\emph{divided} rather than multiplied. Which form is convenient depends on whether it is $\Phi$
or $\Phi^{-1}$ that you can write down. In \cref{ex:change_of_variables_domain} below it is
$\Phi$, so the dividing form is the one that saves work.

% Added 2026-08-13 (Claude Opus 5, Extra): the two halves of the proof are worked out at the very
% end of this section, six hundred lines below, and nothing here said so. A reader has no way to
% find them.
The theorem is not proved here, but the two steps that carry its proof are, and both sit at the
end of this section. \Cref{ex:linear_substitution_elementary_matrices} settles the linear case,
where the determinant enters and everything else is bookkeeping.
\Cref{ex:cube_covering_lemma_substitution} settles the geometric step that a Lipschitz estimate
cannot reach, namely that a diffeomorphism close to the identity cannot let a cube collapse.
\end{ainote}

% Source: Jérôme Paschoud/Notizen/Woche 8.2 Riemann-Integral.pdf, p. 2
% Generator: Gemini 3.1 Pro (High)
\begin{remark}[Intuition for the Jacobian determinant]
From linear algebra, we know that for a linear map given by a matrix $L$, the absolute value of its determinant $\lvert\det L\rvert$ describes how much the map scales volumes. For example, the map $L = \left(\begin{smallmatrix} 2 & 3 \\ 1 & -1 \end{smallmatrix}\right)$ expands areas by a factor of $\lvert\det L\rvert = \lvert-2 - 3\rvert = 5$.

\begin{center}
\begin{tikzpicture}[>=Latex, scale=0.8]
  % Left side: Unit square
  \draw[->, gray] (-0.5, 0) -- (2, 0);
  \draw[->, gray] (0, -0.5) -- (0, 2);
  \draw[thick, blue, fill=blue!10] (0,0) -- (1,0) -- (1,1) -- (0,1) -- cycle;
  \draw[->, thick, blue] (0,0) -- (1,0) node[below right, font=\scriptsize, text=black] {$\left(\begin{smallmatrix} 1 \\ 0 \end{smallmatrix}\right)$};
  \draw[->, thick, blue] (0,0) -- (0,1) node[above left, font=\scriptsize, text=black] {$\left(\begin{smallmatrix} 0 \\ 1 \end{smallmatrix}\right)$};
  
  % Arrow in middle
  \draw[->, thick] (2.5, 0.5) -- (4, 0.5) node[midway, above] {$L$};
  
  % Right side: Parallelogram
  \begin{scope}[shift={(6, 1)}]
    \draw[->, gray] (-0.5, 0) -- (5.5, 0);
    \draw[->, gray] (0, -2.5) -- (0, 2.5);
    
    \coordinate (O) at (0,0);
    \coordinate (V1) at (2,1);
    \coordinate (V2) at (3,-1);
    \coordinate (V3) at (5,0);
    
    \draw[thick, blue, fill=blue!10] (O) -- (V1) -- (V3) -- (V2) -- cycle;
    \draw[->, thick, blue] (O) -- (V1) node[above left, font=\scriptsize, text=black] {$\left(\begin{smallmatrix} 2 \\ 1 \end{smallmatrix}\right)$};
    \draw[->, thick, blue] (O) -- (V2) node[below left, font=\scriptsize, text=black] {$\left(\begin{smallmatrix} 3 \\ -1 \end{smallmatrix}\right)$};
    
    % Correction: Claude Opus 5 (High) --- this caption sat at (2.5,-1.5), which is where the
    % `below left' label of V2 = (3,-1) also lands, so the column vector printed on top of
    % "Area is |det L| = 5". Dropped clear of it; nothing else lies below the parallelogram.
    \node[blue, align=center, font=\small] at (2.5, -2.1) {Area is $\lvert\det L\rvert = 5$\\ times larger};
  \end{scope}
\end{tikzpicture}
\end{center}

The change of variables formula generalizes this to non-linear maps: the Jacobian determinant acts as the local volume scaling factor.
\end{remark}

% Source: Jérôme Paschoud/Notizen/Woche 8.2 Riemann-Integral.pdf, pp. 2--3
% Generator: Gemini 3.1 Pro (High)
\begin{example}[Gaussian integral via polar coordinates]
\label{ex:gaussian_integral_polar}
A classic application of the change of variables formula is computing the Gaussian integral $I := \int_{-\infty}^\infty e^{-x^2}\,dx$. We consider its square:
\[
I^2 = \left( \int_{-\infty}^\infty e^{-x^2}\,dx \right) \left( \int_{-\infty}^\infty e^{-y^2}\,dy \right) = \int_{\mathbb{R}^2} e^{-(x^2+y^2)} \mathop{d\text{vol}}(x,y).
\]
We use polar coordinates, but strictly following \cref{thm:change_of_variables}, we define the forward map $\Phi: \mathbb{R}^2 \setminus \{(0,0)\} \to (0, \infty) \times [-\pi, \pi)$ by
\[
\Phi(x,y) := \left( \sqrt{x^2+y^2}, \operatorname{atan2}(y,x) \right) =: (r, \varphi).
\]

\begin{center}
% Correction: Claude Opus 5 (Max) --- the picture exhibited no size difference at all, under a
% caption reading "not equal in size". The two source rectangles were $[0.5,1.5]\times[0,0.5]$
% and $[0.5,1.5]\times[0.5,1]$: congruent, area $0.5$ each. Their images were the sectors
% $r\in[0.5,1.5]$ over $0$--$30^\circ$ and over $30$--$60^\circ$: also congruent, area $0.5236$
% each. Stacking the two rectangles in $\varphi$ was the cause, because $\det\Jac\Phi^{-1} = r$
% depends on $r$ alone: shifting a rectangle in $\varphi$ cannot change its image's area, and
% shifting it in $r$ is the only thing that can.
%
% The two rectangles are now side by side in $r$, both $0.5 \times 0.6$, at the same
% $\varphi$-band. Their images are then sectors of the same angular width but different radii:
%   red   $\tfrac12(45^\circ-10^\circ)(0.9^2-0.4^2) = \tfrac12(0.6109)(0.65) = 0.199$,
%   blue  $\tfrac12(45^\circ-10^\circ)(2.0^2-1.5^2) = \tfrac12(0.6109)(1.75) = 0.535$,
% a factor of $2.69$, which is the ratio of the mean radii $1.75/0.65$ as $\det\Jac = r$ predicts.
\begin{tikzpicture}[>=Latex, scale=1]
  % Left side: (r, phi) plane. Two CONGRUENT rectangles, 0.5 wide and 0.6 tall, differing only
  % in where they sit along r -- so r_2 - r_1 = r_4 - r_3 = 0.5 by construction.
  \draw[->, gray] (-0.2, 0) -- (2.6, 0) node[right, font=\scriptsize, text=black] {$r$};
  \draw[->, gray] (0, -0.2) -- (0, 1.8) node[above, font=\scriptsize, text=black] {$\varphi$};

  \draw[thick, fill=red!20]  (0.4, 0.5) rectangle (0.9, 1.1);
  \draw[thick, fill=blue!20] (1.5, 0.5) rectangle (2.0, 1.1);

  \node[left, font=\scriptsize] at (0, 1.1) {$\varphi_2$};
  \node[left, font=\scriptsize] at (0, 0.5) {$\varphi_1$};
  \node[below, font=\scriptsize] at (0.4, 0) {$r_1$};
  \node[below, font=\scriptsize] at (0.9, 0) {$r_2$};
  \node[below, font=\scriptsize] at (1.5, 0) {$r_3$};
  \node[below, font=\scriptsize] at (2.0, 0) {$r_4$};
  % Anchored south WEST, not south: centred at x = 1.2 this line is about 2.6 units wide, so it
  % reached back past x = 0 and the $\varphi$ axis printed through its opening word. Rendering
  % p.278 showed it; anchoring the left edge is the only placement that cannot drift back.
  \node[font=\scriptsize, anchor=south west] at (0.15, 1.25) {same $\Delta r$, same $\Delta\varphi$};

  % Arrow
  \draw[->, thick] (3, 0.5) .. controls (3.3, 0.8) and (3.7, 0.8) .. (4, 0.5) node[midway, above] {$\Phi^{-1}$};

  % Right side: (x,y) plane. Same angular band for both, so only the radius differs.
  \begin{scope}[shift={(6.5, 0.5)}]
    \draw[->, gray] (-0.6, 0) -- (2.6, 0) node[right, font=\scriptsize, text=black] {$x$};
    \draw[->, gray] (0, -0.6) -- (0, 2.6) node[above, font=\scriptsize, text=black] {$y$};

    \foreach \r in {0.4, 0.9, 1.5, 2.0} { \draw[gray!45, dashed] (0,0) circle (\r); }

    % Red sector: the image of the rectangle near the origin.
    \draw[thick, fill=red!20]
      (10:0.4) -- (10:0.9) arc (10:45:0.9) -- (45:0.4) arc (45:10:0.4) -- cycle;
    % Blue sector: the image of the congruent rectangle further out.
    \draw[thick, fill=blue!20]
      (10:1.5) -- (10:2.0) arc (10:45:2.0) -- (45:1.5) arc (45:10:1.5) -- cycle;
  \end{scope}

  \node[align=center, font=\small] at (4.4, -1.7)
    {two congruent rectangles, two images of \emph{different} area\\
     $\Rightarrow$ we must compensate with the (Jacobian) determinant};
\end{tikzpicture}
\end{center}

The Jacobian matrix is
\[
\Jac\Phi(x,y) = \begin{pmatrix} \frac{x}{\sqrt{x^2+y^2}} & \frac{y}{\sqrt{x^2+y^2}} \\[1ex] \frac{-y}{x^2+y^2} & \frac{x}{x^2+y^2} \end{pmatrix}.
\]
Its determinant is
\[
\det \Jac\Phi(x,y) = \frac{x^2}{(x^2+y^2)^{3/2}} + \frac{y^2}{(x^2+y^2)^{3/2}} = \frac{1}{\sqrt{x^2+y^2}} = \frac{1}{r}.
\]
Applying the theorem (ignoring the null set $\{(0,0)\}$ where $\Phi$ is undefined), we get
\[
I^2 = \int_{\Phi(\mathbb{R}^2 \setminus \{(0,0)\})} \frac{e^{-r^2}}{\lvert1/r\rvert} \,dr\,d\varphi = \int_{-\pi}^\pi d\varphi \int_0^\infty r e^{-r^2}\,dr = 2\pi \cdot \left[ -\frac{1}{2} e^{-r^2} \right]_0^\infty = 2\pi \cdot \frac{1}{2} = \pi.
\]
Since $I > 0$, we conclude that $I = \sqrt{\pi}$.
\end{example}

% Supplement: Sascha Brack/Class Notes/Week_08_Notes_Friday.pdf, pp. 2--3
\begin{example}[Change of variables on a diamond]
\label{ex:change_of_variables_diamond}
Compute $I := \int_E e^{x+y}(x-y) \mathop{d\text{vol}}(x,y)$, where $E$ is the region bounded by the lines $y=x$, $y=-x+2$, $y=x-2$, and $y=-x$.

Writing out the inequalities, $E = \{ y \leq x \} \cap \{ y \leq -x+2 \} \cap \{ y \geq -x \} \cap \{ y \geq x-2 \}$.
Rearranging these gives $0 \leq x-y \leq 2$ and $0 \leq x+y \leq 2$.
This suggests the substitution $u := x+y$ and $v := x-y$. The new domain is the square $V = [0,2] \times [0,2]$.
The inverse map is $\Phi^{-1}(u,v) = (x,y) = \big(\frac{u+v}{2}, \frac{u-v}{2}\big)$.
Its Jacobian matrix is $\Jac\Phi^{-1}(u,v) = \begin{pmatrix} 1/2 & 1/2 \\ 1/2 & -1/2 \end{pmatrix}$, so $\big|\det\Jac\Phi^{-1}\big| = \big|-\frac{1}{4} - \frac{1}{4}\big| = \frac{1}{2}$.
By the change of variables formula:
\[I = \int_0^2 \int_0^2 e^u v \left|\det\Jac\Phi^{-1}\right| \,du\,dv = \frac{1}{2} \int_0^2 \int_0^2 e^u v \,du\,dv = \frac{1}{2} \left[ e^u \right]_0^2 \left[ \frac{v^2}{2} \right]_0^2 = \frac{1}{2} (e^2-1) \cdot 2 = e^2-1.\]
\end{example}

\begin{example}[Improper integral via polar coordinates]
\label{ex:improper_integral_polar}
Compute $\int_{\mathbb{R}^2} \frac{1}{(1+x^2+y^2)^2} \mathop{d\text{vol}}(x,y)$.
This integral is improper because the domain is unbounded. Since the integrand is non-negative and continuous, its Riemann improper integral is well-defined, and we can compute it using polar coordinates $x = r\cos\theta, y = r\sin\theta$.
The Jacobian determinant is $r$. The domain $\mathbb{R}^2$ corresponds to $r \in [0, \infty)$ and $\theta \in [0, 2\pi)$.
\[
\int_{\mathbb{R}^2} \frac{1}{(1+x^2+y^2)^2} \,dx\,dy
= \int_0^{2\pi} \int_0^\infty \frac{r}{(1+r^2)^2} \,dr\,d\theta.
\]
Substituting $u = r^2$, $du = 2r\,dr$, the inner integral is $\frac{1}{2} \int_0^\infty \frac{1}{(1+u)^2} \,du = \frac{1}{2} \left[ \frac{-1}{1+u} \right]_0^\infty = \frac{1}{2}$.
The outer integral then gives $\frac{1}{2} \int_0^{2\pi} d\theta = \pi$.
\end{example}

% Supplement: Sascha Brack/Class Notes/Week_08_Notes_Monday.pdf, pp. 8--9
\begin{example}[Volume of a sphere]
\label{ex:volume_sphere}
Let $A := \{(x,y,z) \in \mathbb{R}^3 : x^2+y^2+z^2 \leq 1\}$ be the unit ball. Given that $A$ is Jordan measurable, its volume is $\mu(A) = \int_A 1 \mathop{d\text{vol}}(x,y,z)$.
We can compute this using spherical coordinates:
\[\Phi(r,\theta,\varphi) = \begin{pmatrix} r\sin\theta\cos\varphi \\ r\sin\theta\sin\varphi \\ r\cos\theta \end{pmatrix}.\]
Up to a Jordan null set (which does not change the volume), the domain $A$ corresponds to the box $V = (0,1) \times (0,\pi) \times (-\pi,\pi)$.
The Jacobian determinant is $\det\Jac\Phi(r,\theta,\varphi) = r^2\sin\theta$, which is positive on this domain.
Applying the change of variables formula and Fubini's theorem (which allows slicing up the integral):
\begin{align*}
\mu(A) &= \int_{\Phi^{-1}(A)} \big|\det\Jac\Phi(r,\theta,\varphi)\big| \,dr\,d\theta\,d\varphi \\
&= \int_{-\pi}^\pi \int_0^\pi \int_0^1 r^2\sin\theta \,dr\,d\theta\,d\varphi \\
&= \left(\int_{-\pi}^\pi d\varphi\right) \left(\int_0^\pi \sin\theta \,d\theta\right) \left(\int_0^1 r^2 \,dr\right) \\
&= (2\pi) \cdot (2) \cdot \left(\frac{1}{3}\right) = \frac{4\pi}{3}.
\end{align*}
\end{example}

% Quelle: Linus Lüchinger/Slides/Slides-05-11.pdf, slide 3
% Generator: Claude Opus 5 (High)
\begin{remark}[Polar coordinates as slicing into shells]
\label{rem:polar_as_shells}
The factor $r^2$ in the computation above, and the $r$ in the two-dimensional case, are usually
produced by grinding out a Jacobian determinant. There is a reading that gets them without any
computation, and that generalises to $\mathbb{R}^n$ at no extra cost.

Slice $\mathbb{R}^n$ into spheres centred at the origin. The sphere of radius $r$ is an
$(n-1)$-dimensional submanifold, and it is the unit sphere scaled by $r$; since scaling by $r$
multiplies $(n-1)$-dimensional volume by $r^{n-1}$, its volume is
\[\vol_{n-1}\big(r\,\mathbb{S}^{n-1}\big) = C_n\,r^{n-1},
  \qquad C_n := \vol_{n-1}\big(\mathbb{S}^{n-1}\big).\]
A shell between radii $r$ and $r+dr$ is that sphere thickened by $dr$, so it contributes
$C_n r^{n-1}\,dr$ of $n$-dimensional volume. Adding the shells,
\[\int_{B_R(0)} f(\lvert x\rvert)\,dx = C_n\int_0^R f(r)\,r^{n-1}\,dr
  \qquad\text{for radial } f.\]

The powers now need no memorising: $n=2$ gives $r^{1}$ and $n=3$ gives $r^{2}$, matching the two
Jacobians above, and $C_2 = 2\pi$, $C_3 = 4\pi$ are the circumference of the unit circle and the
area of the unit sphere. Taking $f \equiv 1$, $n=3$, $R=1$ reproduces
$4\pi\int_0^1 r^2\,dr = 4\pi/3$ in one line.

This is a genuine shortcut only for \emph{radial} integrands, where the angular integral is a
constant that can be pulled out; when $f$ depends on the angles as well, the full change of
variables is unavoidable. But it explains where the radial power comes from, which the
determinant computation does not.
\end{remark}

% Supplement: Analysis_II_Script_v1.pdf, pp. 113--114
\begin{remark}[Reference: the standard coordinate Jacobians]
\label{rem:coordinate_jacobians_table}
The polar, cylindrical and spherical substitutions recur throughout this chapter and the next
several; collected here in one place, with $r \geq 0$ throughout and each parametrizing its
target domain up to a Jordan null set:
\begin{center}
\begin{tabular}{@{}lll@{}}
\toprule
\textbf{Coordinates} & \textbf{Transformation} & $\lvert\det\Jac\Phi\rvert$ \\
\midrule
Polar ($\mathbb{R}^2$) &
  $\Phi(r,\varphi) = (r\cos\varphi,\ r\sin\varphi)$ & $r$ \\
Cylindrical ($\mathbb{R}^3$) &
  $\Phi(r,\varphi,z) = (r\cos\varphi,\ r\sin\varphi,\ z)$ & $r$ \\
Spherical ($\mathbb{R}^3$) &
  $\Phi(r,\theta,\varphi) = (r\sin\theta\cos\varphi,\ r\sin\theta\sin\varphi,\ r\cos\theta)$
  & $r^2\sin\theta$ \\
\bottomrule
\end{tabular}
\end{center}
The angular ranges are $\varphi \in (-\pi,\pi)$ throughout, and $\theta \in (0,\pi)$
(\newterm{colatitude}, measured from the positive $z$-axis) for the spherical case. Cylindrical
coordinates are just polar coordinates in the first two variables, carried along unchanged in
$z$, which is exactly why the Jacobian is the same $r$ as in the planar case: the extra $dz$
contributes a factor of $1$. Spherical coordinates are used already in
\cref{ex:volume_sphere} above; cylindrical coordinates are the natural choice whenever a region
has rotational symmetry about an axis without being a ball, e.g.\ a solid cylinder or cone.
\end{remark}

% Source: old_exams/examFS24.pdf (21 August 2024), Wahr/Falsch-Fragen, Aufgabe 8, p. 5
% Extractor: Claude Opus 5 (High)
% Worked out here rather than set as an exercise: it is a checklist against the table immediately
% above, and its three-dimensional counterpart ex:examHS24_TF_spherical is already an exercise
% further down this file. Note the exam's argument order (theta, r), which is the reverse of the
% table's; kept as the exam has it, since the sign of the determinant depends on it.
\begin{example}[Polar coordinates, checked item by item]
\label{ex:polar_coordinates_checklist}
% Generator: Claude Opus 5 (High) --- the four statements are the examiner's; the working is
% written for these notes.
Consider $\Psi : \mathbb{R}\times[0,\infty) \to \mathbb{R}^2$,
$\Psi(\theta, r) := (r\cos\theta,\ r\sin\theta)$, and note the argument order: the angle first.
Four statements, of which two are true.
\begin{enumerate}[label=\textbf{(\alph*)}]
  \item \textbf{$\Psi$ restricted to $(0,2\pi]\times(0,\infty)$ is injective: true.} Suppose
    $\Psi(\theta_1,r_1) = \Psi(\theta_2,r_2)$. Taking Euclidean norms gives $r_1 = r_2 =: r > 0$,
    and dividing by $r$ leaves $(\cos\theta_1,\sin\theta_1) = (\cos\theta_2,\sin\theta_2)$, so
    $\theta_1 - \theta_2 \in 2\pi\mathbb{Z}$. A half-open interval of length $2\pi$ contains at
    most one representative of each such class, so $\theta_1 = \theta_2$.
  \item \textbf{$\Psi$ restricted to $(0,2\pi]\times(0,\infty)$ is surjective: false.} Its image is
    $\mathbb{R}^2\setminus\{(0,0)\}$: every non-zero point is reached, by \textbf{(a)} exactly
    once, and the origin is reached only from $r = 0$, which the restriction excludes.
  \item \textbf{$\det\Jac\Psi(\theta,r) = r^2\sin\theta\cos\theta$: false.} The columns of the
    Jacobi matrix are the partial derivatives in $\theta$ and in $r$, in that order:
    \[\Jac\Psi(\theta,r) = \begin{pmatrix} -r\sin\theta & \cos\theta \\
      r\cos\theta & \sin\theta\end{pmatrix}, \qquad
      \det\Jac\Psi(\theta,r) = -r\sin^2\theta - r\cos^2\theta = -r .\]
    The magnitude $\lvert\det\Jac\Psi\rvert = r$ is the entry in
    \cref{rem:coordinate_jacobians_table}; the minus sign appears only because the exam lists the
    angle before the radius, which swaps two columns.
  \item \textbf{$\vol_2(B_R) = 2\pi\int_0^R r\,dr$ for $B_R = \{x^2+y^2 < R^2\}$: true.} By
    \cref{thm:change_of_variables} with the weight $\lvert\det\Jac\Psi\rvert = r$, and since the
    integrand $1$ does not depend on $\theta$, the angular integral contributes a factor $2\pi$.
    The right-hand side is $2\pi\cdot R^2/2 = \pi R^2$, as it should be. This is also the $n=2$
    case of \cref{rem:polar_as_shells}.
\end{enumerate}
The proposed formula in \textbf{(c)} is worth a second look, because it is not a random wrong
answer: $r^2\sin\theta$ is the \emph{spherical} Jacobian, and $\cos\theta$ is what a misremembered
product rule adds to it. There is a cheap guard against importing one row of the table into
another. Every system in it has the form $\Phi(r,\omega) = r\,u(\omega)$ for a map $u$ into the
unit sphere, so $\Phi$ is homogeneous of degree $1$ in $r$. The radial column of $\Jac\Phi$ is then
$u(\omega)$, free of $r$, while each of the $n-1$ angular columns is $r\,\partial u/\partial\omega_i$
and so carries exactly one factor $r$. Multilinearity of the determinant pulls all $n-1$ of them
out, leaving $\lvert\det\Jac\Phi\rvert = r^{n-1}$ times a function of the angles alone. In
$\mathbb{R}^2$ that power is $r^1$, and the proposal carries $r^2$.
\exinfo{This example is taken from the exam of 21 August 2024 (Prof.\ Joaquim Serra), where it appears as the eighth block
of true/false questions.}
\end{example}

% Source: Jérôme Paschoud/Notizen/Woche 9.1 Variablenwechsel.pdf, p. 2
% Generator: Gemini 3.1 Pro (High)
\begin{example}[Volume of an ellipse]
\label{ex:volume_ellipse}
Compute the area (2-dimensional volume) of the ellipse $E := \{(x,y) \in \mathbb{R}^2 \mid (x/a)^2 + (y/b)^2 < 1\}$ for $a,b > 0$.
We use the transformation $\Phi(r,\varphi) = (ar\cos\varphi, br\sin\varphi)$. This maps the rectangle $V = (0,1) \times (0,2\pi)$ bijectively onto $E$ minus a null set (the positive $x$-axis).
The Jacobian matrix is
\[\Jac\Phi(r,\varphi) = \begin{pmatrix} a\cos\varphi & -ar\sin\varphi \\ b\sin\varphi & br\cos\varphi \end{pmatrix},\]
and its determinant is $\det\Jac\Phi(r,\varphi) = abr(\cos^2\varphi + \sin^2\varphi) = abr$.
Thus, by the change of variables formula,
\[\int_E 1 \,dx\,dy = \int_0^{2\pi} \int_0^1 abr \,dr\,d\varphi = \left(\int_0^{2\pi} ab \,d\varphi\right) \left(\int_0^1 r \,dr\right) = 2\pi ab \cdot \frac{1}{2} = \pi ab.\]
\end{example}

% Source: Jérôme Paschoud/Notizen/Woche 9.1 Variablenwechsel.pdf, p. 2
% Generator: Gemini 3.1 Pro (High)
\begin{example}[Area of a cardioid]
\label{ex:area_cardioid}
Compute the area of the region $\Omega$ given in polar coordinates by $\Omega := \{(r,\varphi) \mid \varphi \in [0, 2\pi),\ 0 \leq r \leq 1 + \cos\varphi\}$.
Using polar coordinates $\Phi(r,\varphi) = (r\cos\varphi, r\sin\varphi)$, the area is
\[\int_\Omega 1 \,dx\,dy = \int_0^{2\pi} \int_0^{1+\cos\varphi} r \,dr\,d\varphi = \int_0^{2\pi} \left[ \frac{r^2}{2} \right]_{r=0}^{r=1+\cos\varphi} \,d\varphi = \frac{1}{2} \int_0^{2\pi} (1+\cos\varphi)^2 \,d\varphi.\]
Expanding the integrand: $(1+\cos\varphi)^2 = 1 + 2\cos\varphi + \cos^2\varphi$. We know $\int_0^{2\pi} \cos\varphi \,d\varphi = 0$ and $\int_0^{2\pi} \cos^2\varphi \,d\varphi = \pi$. Thus,
\[\int_\Omega 1 \,dx\,dy = \frac{1}{2} (2\pi + 0 + \pi) = \frac{3\pi}{2}.\]
\end{example}

% Source: Jérôme Paschoud/Notizen/Woche 9.1 Variablenwechsel.pdf, p. 3
% Generator: Gemini 3.1 Pro (High)
\begin{example}[Combining substitution and Fubini]
\label{ex:substitution_and_fubini_rational}
Compute $\int_\Omega (x^3 y + x y^3) \,dx\,dy$ where $\Omega := \{(x,y) \in \mathbb{R}^2 \mid x,y > 0,\ 1/x \le y \le 3/x,\ 1 \le x^2 - y^2 \le 4\}$.
The integrand factors as $xy(x^2+y^2)$. The boundary inequalities are equivalent to $1 \le xy \le 3$ and $1 \le x^2-y^2 \le 4$, which strongly suggests defining $u := xy$ and $v := x^2 - y^2$.
The domain in the $(u,v)$-plane is the rectangle $R = [1,3] \times [1,4]$.
Let $\Psi(x,y) = (xy, x^2-y^2) = (u,v)$. The Jacobian of this inverse transformation is
\[\Jac\Psi(x,y) = \begin{pmatrix} y & x \\ 2x & -2y \end{pmatrix}, \qquad \det\Jac\Psi(x,y) = -2y^2 - 2x^2 = -2(x^2+y^2).\]
For the forward transformation $\Phi = \Psi^{-1}$, we have $\big|\det\Jac\Phi(u,v)\big| = \frac{1}{\lvert\det\Jac\Psi(x,y)\rvert} = \frac{1}{2(x^2+y^2)}$.
Applying the change of variables formula:
\[\int_\Omega xy(x^2+y^2) \,dx\,dy = \int_R u(x^2+y^2) \frac{1}{2(x^2+y^2)} \,du\,dv = \frac{1}{2} \int_1^4 \int_1^3 u \,du\,dv = \frac{1}{2} \cdot 3 \cdot \left[ \frac{u^2}{2} \right]_1^3 = \frac{3}{2} \cdot 4 = 6.\]
\end{example}


% Source: Corsin Nick/Class Notes/Week 8.pdf, p. 7


% --- exercises relocated here from the dissolved sheet 8 ---

% Originally: content/week-08.tex, \exercisesheet{8}, problem 8.5
% Source: exercises/Ex8_Analysis2_eng.pdf

% Source: exercises/Ex8_Analysis2_eng.pdf, pp. 1--2
\begin{exercise}[Change of variables and Jacobians \textnormal{(important)}]
\label{ex:change_of_variables_computations}
For each of the following domains and change of variables find the jacobian and the appropriate
transformed domain. There is no need to actually compute the integrals!
\begin{enumerate}[label=\textbf{(\alph*)}]
  \item $A := \{(x_1,x_2) \in \mathbb{R}^2 : x_1 > 0,\ x_2 < x_1,\ 1 < x_1^2+x_2^2 < 4\}$ and
    $x_1 = r\cos\theta$, $x_2 = r\sin\theta$. Using the change of variables formula, complete
    the dots in the following formula
    \[\int_A x_1^2\sin(x_2)\,dx_1dx_2 = \int_{\dots}\dots\,dr\,d\theta\]
  \item $B := \{(x,y) \in \mathbb{R}^2 \mid 1 < xy < 2,\ x^2 < y < 2x^2\}$ and $u := xy$,
    $v := x^2$. Using the change of variables formula, complete the dots in the following
    formula
    \[\int_B y^2e^{-xy}\,dx\,dy = \int_{\dots}\dots\,du\,dv\]
  \item $C := \{(x,y,z) \in \mathbb{R}^3 \mid 1 < z-2y < 2,\ 0 < z < 1,\ -2 < 3x+y+z < 0\}$ and
    $u := z$, $v := z-2y$, $w := 3x+y+z$. Using the change of variables formula, complete the
    dots in the following formula
    \[\int_C xyz\,dx\,dy = \int_{\dots}\dots\,du\,dv\,dw.\]
\end{enumerate}

% Supplement: Sascha Brack/Ex Sheet Hints/Ex8_Analysis2_hints.pdf, p. 1
\exhint[Sascha Brack's hint on 8.5.2]{First pin down the signs. The constraints force
$x \neq 0$ and $y \neq 0$; then $y > 0$ forces $x > 0$, so $u > 0$ and $v > 0$ throughout $B$.
That is what lets you divide by $x$ and $y$ and rearrange the inequalities without their
directions flipping.}
\exinfo{This exercise is Problem 8.5 of Problem Sheet 8, Analysis II, Spring Semester 2026.}
\exsol{sol:change_of_variables_computations}
\end{exercise}

\begin{ainote}
The integral in \textbf{8.5.3} is written $\int_C xyz\,dx\,dy$ on the official sheet, but $C$ is
a subset of $\mathbb{R}^3$ and the right-hand side carries three differentials, so this must
read $dx\,dy\,dz$. Quoted verbatim above, as the sheet is a primary source.
\end{ainote}

% An AI-Note stood here until 2026-08-09 announcing, above the exercise, that part (c)'s domain is
% the box (0,1)x(1,2)x(-2,0) and that part (b)'s second constraint becomes v^{3/2} < u < 2v^{3/2}.
% That is the answer to two of the three parts. Moved into 99-solutions.tex per the spoiler rule
% in style.md; the surviving half duplicated the AI-Note already sitting beside the solution.

\begin{aiexercise}[A substitution read off the integrand]
% Generator: Claude Opus 5 (High)
\label{ex:ai_triangle_substitution}
Let $T$ be the closed triangle with vertices $(0,0)$, $(1,0)$ and $(0,1)$, and consider
\[I := \int_T \exp\!\left(\frac{y-x}{y+x}\right)dx\,dy.\]
\begin{enumerate}[label=\textbf{(\alph*)}]
  \item Substitute $u := y-x$ and $v := y+x$. Show that the interior of $T$ corresponds to
    $\{(u,v) : 0 < v < 1,\ -v < u < v\}$, and find the factor relating $dx\,dy$ to $du\,dv$.
  \item Evaluate $I$. The inner integral is the one to do first, and it is elementary.
  \item The integrand is undefined at the origin, which lies in $T$. Why does this not affect
    the answer?
\end{enumerate}
\exsol{sol:ai_triangle_substitution}
\end{aiexercise}

% Source: old_exams/examHS24.pdf (13 February 2025), Kurzproblem 1, part 1, p. 8
% Extractor: Gemini 3.6 Flash (High)
% Correction: Claude Opus 5 (High) --- \operatorname{vol} -> \vol and J\Phi -> \Jac\Phi per the
% declared-macro rule.
\begin{exercise}[Volume of an ellipsoid via linear change of variables]
\label{ex:volume_ellipsoid_change_variables}
Let $a, b, c > 0$ and consider the solid ellipsoid
\[U := \left\{(x,y,z) \in \mathbb{R}^3 \ \middle|\ \left(\frac{x}{a}\right)^2 + \left(\frac{y}{b}\right)^2 + \left(\frac{z}{c}\right)^2 \leq 1\right\}.\]
\begin{enumerate}[label=\textbf{(\alph*)}]
  \item Define a linear change of variables $\Phi : \mathbb{R}^3 \to \mathbb{R}^3$ mapping the closed unit ball $\overline{B_1(0)} \subseteq \mathbb{R}^3$ bijectively onto $U$.
  \item Compute the Jacobian determinant $\det \Jac\Phi(u,v,w)$ of this transformation.
  \item Compute the three-dimensional Jordan volume $\vol_3(U)$ using the change of variables formula and the known volume of the unit ball.
\end{enumerate}
\exinfo{This exercise is taken from the exam of 13 February 2025 (Prof.\ Joaquim Serra), Short Problem 1, part 1, which
also asks for a labelled 3D sketch of $U$.}
\exsol{sol:volume_ellipsoid_change_variables}
\end{exercise}

\begin{ainote}
The original exam writes the ellipsoid with an equality, $(x/a)^2+(y/b)^2+(z/c)^2 = 1$, and then
asks for its volume --- but that set is the ellipsoid \emph{surface}, whose Jordan volume is $0$.
The intended object is the solid region it bounds, so the inequality is used here.
\end{ainote}

% Source: old_exams/fs2023/Prüfung.pdf (9 August 2021), Teil A, Frage 5, p. 6
% Extractor: Claude Opus 5 (High)
\begin{exercise}[A cubic shear, and the area it sweeps out]
\label{ex:cubic_shear_area}
Let $\varphi : \mathbb{R}^2 \to \mathbb{R}^2$ be given by $\varphi(x) := (x_1^3 - x_2,\; x_2 + x_1)$,
where $x = (x_1,x_2)$.
\begin{enumerate}[label=\textbf{(\alph*)}]
  \item Compute the derivative $D_x\varphi$ for every $x \in \mathbb{R}^2$.
  \item Show that $\varphi : \mathbb{R}^2 \to \varphi(\mathbb{R}^2)$ is a diffeomorphism.
  \item What is the area of the image $\varphi(Q)$ of the square $Q := [0,1] \times [0,1]$?
\end{enumerate}
\exinfo{This exercise is taken from the Analysis I \& II exam of 9 August 2021 (Prof.\ Giovanni Felder), Part A, Problem 5.}
\exsol{sol:cubic_shear_area}
\end{exercise}

% Source: old_exams/examHS24.pdf (13 February 2025), Wahr/Falsch-Fragen, Aufgabe 9, p. 5
% Extractor: Claude Opus 5 (High)
\begin{exercise}[A Gaussian second moment in polar coordinates]
\label{ex:gaussian_second_moment_polar}
For $\alpha > 0$, compute the integral
\[I_\alpha := \int_{\mathbb{R}^2} x^2 e^{-\alpha(x^2+y^2)} \, dx\,dy .\]
\exhint[Official hint]{Notice that, by symmetry, $I_\alpha = \int_{\mathbb{R}^2} y^2
e^{-\alpha(x^2+y^2)}\,dx\,dy$ as well. Compute $2I_\alpha$ using polar coordinates.}
\exinfo{This exercise is taken from the exam of 13 February 2025 (Prof.\ Joaquim Serra), where it appears as the ninth
block of true/false questions, there phrased as a series of checkpoints along this computation.}
\exsol{sol:gaussian_second_moment_polar}
\end{exercise}

% Source: old_exams/examFS24.pdf (21 August 2024), Wahr/Falsch-Fragen, Aufgabe 9, p. 5
% Extractor: Claude Opus 5 (High)
% Filed next to its February 2025 sibling above deliberately: both are one integral over R^2 in
% polar coordinates, and the pair is worth doing together because the weight differs (Gaussian
% against exponential-of-the-radius) while the method does not.
\begin{exercise}[An exponential of the radius, and how it decays in $\alpha$]
\label{ex:radial_exponential_integral}
For $\alpha > 0$, compute
\[I_\alpha := \int_{\mathbb{R}^2} e^{-\alpha\sqrt{x^2+y^2}}\,dx\,dy,\]
and then decide whether each of the following is true or false: $I_1 = \pi$; $I_2 = \pi/2$;
$\lim_{\alpha\to+\infty}\alpha I_\alpha = 0$.
\exhint[Official hint]{You may use polar coordinates and integration by parts.}
\exinfo{This exercise is taken from the exam of 21 August 2024 (Prof.\ Joaquim Serra), where it appears as the ninth block
of true/false questions. The instruction to compute $I_\alpha$ is the exam's own preamble to those
three statements.}
\exsol{sol:radial_exponential_integral}
\end{exercise}

% Source: old_exams/examHS24.pdf (13 February 2025), Wahr/Falsch-Fragen, Aufgabe 6, p. 4
% Extractor: Claude Opus 5 (High)
% Filed in this chapter rather than in 15-inverse-function-theorem, where a Gemini 3.1 Pro pass had
% put a different (and invented) block under this citation: only part (a) below is the inverse
% function theorem, and parts (b) and (c) are change of variables.
\begin{exercise}[True or False \textnormal{(the image of a small set under a $C^1$ map)}]
\label{ex:examHS24_TF_image_of_small_sets}
Consider $\Phi \in C^1(\mathbb{R}^2, \mathbb{R}^2)$ with
\[\Phi(0,0) = \begin{pmatrix} 0 \\ 0 \end{pmatrix} \quad\text{and}\quad
  \Jac\Phi(0,0) = \begin{pmatrix} 2 & 1 \\ -1 & 1 \end{pmatrix},\]
and let $X := \Phi^{-1}\big((0,0)\transp\big)$. Decide whether each of the following statements is
true or false.
\begin{enumerate}[label=\textbf{(\alph*)}]
  \item For $r > 0$ small enough, the set $\Phi(B_r(0,0))$ is open.
  \item $\vol_2\big(\Phi(B_r(0,0))\big) = 2\pi r^2 + o(r^2)$ as $r \to 0^+$.
  \item $\vol_2\big(\Phi([0,r] \times [0, ar^{1/2}])\big) = 3ar^{3/2} + o(r^{3/2})$ as
    $r \to 0^+$, where $a > 0$ is a constant.
\end{enumerate}
\exinfo{This exercise is taken from the exam of 13 February 2025 (Prof.\ Joaquim Serra), where it appears as the sixth
block of true/false questions.}
\exsol{sol:image_of_small_sets_tf}
\end{exercise}

% Source: old_exams/examHS24.pdf (13 February 2025), Wahr/Falsch-Fragen, Aufgabe 7, p. 4
% Extractor: Claude Opus 5 (High)
% Replaces a block that a Gemini 3.1 Pro pass filed under this same citation. That block asked for
% the volume of an ellipsoid spanned by three vectors of determinant 2, which is not on p. 4 and is
% nowhere on this exam; it survives, honestly relabelled, as ex:ai_spanned_ellipsoid_volume below.
\begin{exercise}[True or False \textnormal{(volumes under an affine map)}]
\label{ex:examHS24_TF_affine_volume_scaling}
Consider the map
\[\Psi : \mathbb{R}^3 \to \mathbb{R}^3, \qquad
  \Psi(x,y,z) := \big(-x + e^\pi,\; y + \alpha z,\; \alpha y + z\big),\]
where $\alpha \in \mathbb{R}$ is a parameter. Decide whether each of the following statements is
true or false.
\begin{enumerate}[label=\textbf{(\alph*)}]
  \item $\vol_3(\Psi(E)) = 0$ for all Jordan measurable $E \subseteq \mathbb{R}^3$ if and only if
    $\alpha = 1$.
  \item For $\alpha = -2$, $\vol_3(\Psi^k(E)) = 3^k \vol_3(E)$ for all Jordan measurable
    $E \subseteq \mathbb{R}^3$ and all $k \ge 0$, where $\Psi^0 := \id$ and
    $\Psi^{k+1} := \Psi \circ \Psi^k$.
  \item For $\alpha = 0$, $\vol_3(\Psi^{2025}(E)) = \vol_3(E)$ for all Jordan measurable
    $E \subseteq \mathbb{R}^3$.
  \item For all $\alpha \in (0,1)$, $\vol_3\big(\Psi^{-1}([0,1]^3)\big) = -1/(1-\alpha^2)$.
\end{enumerate}
\exinfo{This exercise is taken from the exam of 13 February 2025 (Prof.\ Joaquim Serra), where it appears as the seventh
block of true/false questions.}
\exsol{sol:affine_volume_scaling_tf}
\end{exercise}

% Source: old_exams/examFS24.pdf (21 August 2024), Wahr/Falsch-Fragen, Aufgabe 7, pp. 4--5
% Extractor: Claude Opus 5 (High)
% Worked out rather than set as an exercise, because the exercise directly above is the same
% question with a parameter in it and is strictly harder. What this block adds is the composition
% and inverse rules read off one determinant, which the February 2025 version states only for
% iterates of a single map.
\begin{example}[One determinant, three volume statements]
\label{ex:affine_volume_one_determinant}
% Generator: Claude Opus 5 (High) --- the three statements are the examiner's; the working is
% written for these notes.
Consider $\Psi : \mathbb{R}^3 \to \mathbb{R}^3$, $\Psi(x,y,z) := (z-1,\ x+2y,\ x-y)$. Its linear
part has matrix
\[A = \begin{pmatrix} 0 & 0 & 1 \\ 1 & 2 & 0 \\ 1 & -1 & 0\end{pmatrix},
  \qquad \det A = 1\cdot\det\begin{pmatrix} 1 & 2 \\ 1 & -1\end{pmatrix} = -3,\]
expanding along the first row. The translation by $(-1,0,0)$ does not affect volumes, so
$\vol_3(\Psi(E)) = 3\vol_3(E)$ for every Jordan measurable $E$, by
\cref{thm:change_of_variables}. Everything else follows from the factor $3$.
\begin{enumerate}[label=\textbf{(\alph*)}]
  \item \textbf{$\vol_3(\Psi^{-1}(E)) = \vol_3(E)$ for all Jordan measurable $E$: false.} The
    inverse scales by $\lvert\det A^{-1}\rvert = 1/3$, so $\vol_3(\Psi^{-1}(E)) = \vol_3(E)/3$.
  \item \textbf{$\vol_3(\Psi(E)) = 2\vol_3(E)$: false.} The factor is $3$, not $2$.
  \item \textbf{$\vol_3(\Psi(\Psi(E))) = 9\vol_3(E)$: true.} Applying \textbf{(b)}'s correct factor
    twice gives $3 \cdot 3 = 9$, which is $\lvert\det A^2\rvert = \lvert\det A\rvert^2$.
\end{enumerate}
Only \textbf{(b)} requires the determinant to be computed. Items \textbf{(a)} and \textbf{(c)} then
follow from multiplicativity of the determinant alone, without recomputing anything, which is the
point of the block.
\exhint[Official hint]{To answer these three, it can be useful to compute the Jacobian determinant
of $\Psi$ first.}
\exinfo{This example is taken from the exam of 21 August 2024 (Prof.\ Joaquim Serra), where it appears as the seventh
block of true/false questions.}
\end{example}

% Source: old_exams/examHS24.pdf (13 February 2025), Wahr/Falsch-Fragen, Aufgabe 8, pp. 4--5
% Extractor: Gemini 3.1 Pro (High)
\begin{exercise}[True or False \textnormal{(spherical coordinates)}]
\label{ex:examHS24_TF_spherical}
Consider the polar (spherical) coordinates $\Psi \colon (0, 2\pi) \times (0, \pi) \times (0, \infty)
\to \mathbb{R}^3$ defined by
\[ \Psi(\phi, \theta, r) = (r \cos \phi \sin \theta, r \sin \phi \sin \theta, r \cos \theta). \]
Decide whether each of the following statements is true or false.
\begin{enumerate}[label=\textbf{(\alph*)}]
  \item $\Psi$ is injective.
  \item $\Psi$ is surjective.
  \item The Jacobian determinant $\det \Jac\Psi(\theta, r)$ is given by $r^2 \sin \theta$.
  \item $\vol_3(B_R \setminus B_r) = \pi(R^3 - r^3)$.
  \item $\vol_3(B_R) = 2\pi \int_0^\pi \int_0^R r^2 \sin \theta \, dr \, d\theta$.
  \item $\lim_{R \to 1^+} \frac{\vol_3(B_R \setminus B_1)}{R - 1} = 4\pi$.
\end{enumerate}
\exinfo{This exercise is taken from the exam of 13 February 2025 (Prof.\ Joaquim Serra), where it appears as the eighth
block of true/false questions.}
\exsol{sol:spherical_coordinates_integral_tf}
\end{exercise}

% Source: old_exams/examHS24.pdf (13 February 2025), Kästchenaufgabe 2, p. 7
% Extractor: Gemini 3.1 Pro (High)
\begin{exercise}[Box question: Linear map scaling volume]
\label{ex:examHS24_box_linear_vol}
Give an explicit example of a linear map $\Phi \colon \mathbb{R}^2 \to \mathbb{R}^2$ such that
\[ \vol_2(\Phi(E)) = 2 \vol_2(E) \quad \text{for all } E \subset \mathbb{R}^2 \text{ Jordan measurable.} \]
\exinfo{This exercise is taken from the exam of 13 February 2025 (Prof.\ Joaquim Serra), Box Problem 2, where only the
final answer was graded.}
\exsol{sol:linear_map_volume_scaling_box}
\end{exercise}

% Source: old_exams/examFS24.pdf (21 August 2024), Kästchenaufgabe 2, p. 7
% Extractor: Claude Opus 5 (High)
% The August 2024 counterpart of the box problem directly above: same one-line format, factor 1
% instead of factor 2, and the added clause "different from the identity", which is what stops the
% obvious answer and makes the problem worth having next to its sibling.
\begin{exercise}[Box question: a linear map that preserves every volume]
\label{ex:box_linear_volume_preserving}
Give an explicit example of a linear map $\Phi : \mathbb{R}^2 \to \mathbb{R}^2$, \textbf{different
from the identity}, such that
\[\vol_2(\Phi(E)) = \vol_2(E) \quad \text{for all Jordan measurable } E \subseteq \mathbb{R}^2 .\]
\exinfo{This exercise is taken from the exam of 21 August 2024 (Prof.\ Joaquim Serra), Box Problem 2, where only the final
answer was graded.}
\exsol{sol:box_linear_volume_preserving}
\end{exercise}

% Correction: Claude Opus 5 (High) --- provenance repair. A Gemini 3.1 Pro pass filed the statement
% below as an exam problem, citing old_exams/examHS24.pdf, Aufgabe 7, p. 4. That page carries
% ex:examHS24_TF_affine_volume_scaling instead, and nothing resembling this problem is anywhere on
% that exam. The mathematics is sound, so the block is kept as what it actually is: an exercise
% authored for these notes. Its % Source:, % Extractor: and \exinfo lines have been removed.
\begin{aiexercise}[The volume of a spanned ellipsoid]
\label{ex:ai_spanned_ellipsoid_volume}
% Generator: Gemini 3.1 Pro (High)
Let $v_1, v_2, v_3 \in \mathbb{R}^3$ be vectors with $\det(v_1, v_2, v_3) = 2$, and consider the
region
\[ A := \big\{ t_1 v_1 + t_2 v_2 + t_3 v_3 \ \big|\ (t_1, t_2, t_3) \in \mathbb{R}^3
  \text{ with } t_1^2 + t_2^2 + t_3^2 \le 1 \big\}. \]
Decide whether each of the following statements is true or false.
\begin{enumerate}[label=\textbf{(\alph*)}]
  \item $\vol_3(A) = \frac{8\pi}{3}$.
  \item $A$ is Jordan measurable.
  \item The volume of $A$ depends on the choice of the vectors $v_1, v_2, v_3$.
\end{enumerate}
\exsol{sol:ai_spanned_ellipsoid_volume}
\end{aiexercise}

% Source: old_exams/HS18.pdf (Analysis I & II, winter 2018/2019, Prof. Manfred Einsiedler),
% Aufgabe 13, p. 4
% Extractor: Claude Opus 5 (High)
% This is the proof of thm:jordan_measure_linear_transformation, which chapter 19 states without
% one. Filed here rather than beside the theorem because the exam frames it as the linear case of
% the substitution rule, which is this chapter's subject.
\begin{exercise}[The substitution rule in the linear case]
\label{ex:linear_substitution_elementary_matrices}
Let $A$ be an invertible $n\times n$ matrix. \Cref{thm:jordan_measure_linear_transformation}
asserts that for every Jordan measurable $B \subseteq \mathbb{R}^n$ the image $A(B)$ is again
Jordan measurable, with
\begin{equation}
\label{eq:linear_substitution_rule}
\vol\big(A(B)\big) = \lvert\det A\rvert \, \vol(B).
\end{equation}
This exercise proves it.
\begin{enumerate}[label=\textbf{(\alph*)}]
  \item Prove \eqref{eq:linear_substitution_rule} directly in the case that $A$ is a diagonal
    matrix.
  \item Explain in a few sentences why it now suffices to prove \eqref{eq:linear_substitution_rule}
    for permutation matrices and for shear matrices.
  \item Prove \eqref{eq:linear_substitution_rule} when $A$ is a shear matrix. You may assume for
    simplicity that $B = Q$ is a closed axis-parallel box.
\end{enumerate}
\exinfo{This exercise is taken from the Analysis I \& II examination of winter 2018/2019
(Prof.\ Manfred Einsiedler), Problem 13.}
\exsol{sol:linear_substitution_elementary_matrices}
\end{exercise}

% Source: old_exams/HS17.pdf (Analysis I & II, winter 2017/2018, Prof. Manfred Einsiedler),
% Aufgabe 13, pp. 3--4
% Extractor: Claude Opus 5 (High)
% Kept as a worked example, not an exercise: this is the one genuinely hard lemma behind
% thm:change_of_variables, which this document states without proof, and it is well beyond what
% any exercise here asks. The exam splits it into four graded parts; those are preserved as the
% four steps below.
% Notation normalised to the house convention: the exam writes ||.||_2 for the Euclidean norm,
% which is single bars here (style.md, "Norms"). ||.||_infty keeps its double bars.
\begin{example}[The covering lemma behind the change of variables formula]
\label{ex:cube_covering_lemma_substitution}
\Cref{thm:change_of_variables} is stated in this chapter without proof. The step that makes the
proof hard is not the Jacobian factor but a prior, purely geometric one: a diffeomorphism that is
close to the identity cannot let a cube shrink away. Made precise, that is the following.

\textbf{Claim.} Let $\ell > 0$, let $Q_0 := [-\ell,\ell]^n \subseteq \mathbb{R}^n$, let
$X \supseteq Q_0$ be open and let $\Phi : X \to \Phi(X) \subseteq \mathbb{R}^n$ be a
diffeomorphism with $\Phi(0) = 0$ and $D_0\Phi = I_n$. Suppose $\sigma > 0$ is such that
\[\big\lvert (D_x\Phi - I_n)v \big\rvert \le \sigma \lvert v\rvert
\qquad \text{for all } x \in Q_0,\ v \in \mathbb{R}^n,\]
and assume $s := \sigma\sqrt{n} < 1$. Then
\[(1-s)\,Q_0 \subseteq \Phi(Q_0).\]

Fix $y \in (1-s)Q_0$ and define $F_y : Q_0 \to \mathbb{R}^n$ by
$F_y(x) := y - \big(\Phi(x) - x\big)$.

\begin{enumerate}[label=\textbf{(\alph*)}]
  \item \textbf{Why a fixed point is what we want.} By construction $F_y(x) = x$ holds if and only
    if $y - \Phi(x) + x = x$, that is, if and only if $\Phi(x) = y$. So a fixed point of $F_y$ in
    $Q_0$ is exactly a preimage of $y$ under $\Phi$ lying in $Q_0$, which is what
    $y \in \Phi(Q_0)$ means.

  \item \textbf{The displacement $\Phi - \id$ is a contraction.} Write $\Psi(x) := \Phi(x) - x$, so
    that $D_x\Psi = D_x\Phi - I_n$ and the hypothesis reads
    $\lvert D_x\Psi(v)\rvert \le \sigma\lvert v\rvert$ on $Q_0$, that is,
    $\lVert D_x\Psi\rVert_{\mathrm{op}} \le \sigma$ there. Since $Q_0$ is convex, the segment from
    $x_1$ to $x_2$ stays inside it, and the mean value inequality
    (\cref{thm:mean_value_inequality}) gives
    \[\big\lvert \Phi(x_2)-x_2 - \big(\Phi(x_1)-x_1\big)\big\rvert
      = \lvert \Psi(x_2)-\Psi(x_1)\rvert \le \sigma \lvert x_2-x_1\rvert\]
    for all $x_1,x_2 \in Q_0$.

  \item \textbf{$F_y$ maps $Q_0$ into itself.} Applying \textbf{(b)} with $x_1 := 0$, and using
    $\Psi(0) = \Phi(0)-0 = 0$, we get $\lvert \Psi(x)\rvert \le \sigma\lvert x\rvert$ for
    $x \in Q_0$. For $x \in Q_0$ we have $\lVert x\rVert_\infty \le \ell$ and hence
    $\lvert x\rvert \le \sqrt{n}\,\ell$, so
    \[\lVert \Psi(x)\rVert_\infty \le \lvert \Psi(x)\rvert \le \sigma\sqrt{n}\,\ell = s\ell.\]
    Since $y \in (1-s)Q_0$ means $\lVert y\rVert_\infty \le (1-s)\ell$, the triangle inequality
    gives
    \[\lVert F_y(x)\rVert_\infty \le \lVert y\rVert_\infty + \lVert \Psi(x)\rVert_\infty
      \le (1-s)\ell + s\ell = \ell,\]
    that is, $F_y(x) \in Q_0$. So $F_y : Q_0 \to Q_0$.

  \item \textbf{Conclusion.} $Q_0$ is a closed subset of $\mathbb{R}^n$, hence complete, and by
    \textbf{(b)} the map $F_y$ satisfies
    $\lvert F_y(x_2)-F_y(x_1)\rvert = \lvert \Psi(x_2)-\Psi(x_1)\rvert \le
    \sigma\lvert x_2-x_1\rvert$ with $\sigma \le s < 1$. It is therefore a contraction of the
    complete metric space $Q_0$ into itself, and the Banach fixed point theorem
    (\cref{thm:banach_fixed_point}) supplies a fixed point $x \in Q_0$. By \textbf{(a)},
    $\Phi(x) = y$, so $y \in \Phi(Q_0)$. As $y \in (1-s)Q_0$ was arbitrary,
    $(1-s)Q_0 \subseteq \Phi(Q_0)$.
\end{enumerate}

\begin{remark}[What the lemma is doing in the proof of the change of variables formula]
The formula is proved by cutting the domain into tiny cubes and comparing $\vol(\Phi(Q))$ with
$\lvert\det D\Phi\rvert\vol(Q)$ on each. The upper bound is the easy half: $\Phi$ is Lipschitz on a
compact piece, so $\Phi(Q)$ cannot be much bigger than $Q$. The lower bound is where a naive
argument stalls, because nothing about a Lipschitz map prevents an image from collapsing. This
lemma is that lower bound. First normalise so that $D_0\Phi = I_n$, which is always possible by
composing with a linear map, the case already settled in
\cref{ex:linear_substitution_elementary_matrices}. The lemma then says the image of a cube still
contains a cube shrunk only by the factor $1-s$, and $s \to 0$ as the cubes shrink, because
$D\Phi$ is continuous. Letting the two bounds close is the proof.

Notice which two theorems do the work: the mean value inequality on a convex set in \textbf{(b)},
and the Banach fixed point theorem in \textbf{(d)}. The same pair proves the inverse function
theorem (\cref{thm:inverse_function_theorem}), and that is not a coincidence. This claim is a
quantitative version of the statement that $\Phi$ is locally surjective near a point where its
differential is invertible.
\end{remark}
\exinfo{This example is taken from the Analysis I \& II examination of winter 2017/2018
(Prof.\ Manfred Einsiedler), Problem 13, where the four steps above are set as four graded parts.
The framing and the closing remark are added here.}
\end{example}


% Originally: content/week-08.tex, \subsection{Fubini's theorem} --- promoted to a section

\section{Fubini's theorem}
\label{sec:fubini}

\begin{theorem}[Fubini]
\label{thm:fubini}
Let $X \subseteq \mathbb{R}^m$ and $Y \subseteq \mathbb{R}^n$ be intervals (\qt{boxes}), and let
$f : X\times Y \to \mathbb{R}$, $(x,y)\mapsto f(x,y)$, be continuous. Then
\[\int_{X\times Y} f(x,y)\,dx\,dy
  = \int_X\left(\int_Y f(x,y)\,dy\right)dx
  = \int_Y\left(\int_X f(x,y)\,dx\right)dy,\]
where we look at the expression in brackets as a function, e.g.\
$x \mapsto \int_Y f(x,y)\,dy$.
\end{theorem}

% Source: Diego Torres Tejeda/Notes - 27.04 - Diego Torres Tejeda.pdf, p. 1
% Generator: Gemini 3.6 Flash (High)
\begin{proposition}[Cavalieri's principle]
\label{prop:cavalieri_principle}
Let $A, B \subseteq \mathbb{R}^{n+1}$ be Jordan measurable sets. For each height $t \in \mathbb{R}$, denote by $A_t := \{x \in \mathbb{R}^n \mid (x,t) \in A\}$ and $B_t := \{x \in \mathbb{R}^n \mid (x,t) \in B\}$ the $n$-dimensional cross-sections at $t$. If the cross-sectional volumes agree for all $t \in \mathbb{R}$, i.e.\
\[
\mu_n(A_t) = \mu_n(B_t) \quad \forall t \in \mathbb{R},
\]
then the total $(n+1)$-dimensional volumes are equal: $\mu_{n+1}(A) = \mu_{n+1}(B)$.
\end{proposition}

\begin{proof}
By Fubini's Theorem (\cref{thm:fubini}), integrating the indicator function over slices gives:
\[
\mu_{n+1}(A) = \int_{\mathbb{R}^{n+1}} \chi_A(x,t) \,dx\,dt = \int_{\mathbb{R}} \left( \int_{\mathbb{R}^n} \chi_{A_t}(x) \,dx \right) dt = \int_{\mathbb{R}} \mu_n(A_t) \,dt.
\]
Replacing $\mu_n(A_t)$ with $\mu_n(B_t)$ yields $\mu_{n+1}(B)$.
\end{proof}

\begin{theorem}[Volume of a cone over a measurable base]
\label{thm:cone_volume_formula}
Let $\Omega \subseteq \mathbb{R}^n$ be a bounded, Jordan measurable domain. The \newterm{cone over $\Omega$} \germanfor{Kegel \"uber $\Omega$} with vertex at $(0,1) \in \mathbb{R}^n \times \mathbb{R}$ and base $\Omega \times \{0\}$ is defined by
\[
C\Omega := \big\{(x, t) \in \mathbb{R}^n \times [0,1] \mid x \in (1-t)\Omega\big\}.
\]
Then $C\Omega$ is Jordan measurable in $\mathbb{R}^{n+1}$, and its volume is given by
\[
\mu_{n+1}(C\Omega) = \frac{\mu_n(\Omega)}{n+1}.
\]
\end{theorem}

\begin{proof}
For any $t \in [0,1]$, the cross-section of $C\Omega$ at height $t$ is $(C\Omega)_t = (1-t)\Omega \subseteq \mathbb{R}^n$. By the scaling property of $n$-dimensional Jordan measure, $\mu_n\big((1-t)\Omega\big) = (1-t)^n \mu_n(\Omega)$. Applying Cavalieri's Principle (\cref{prop:cavalieri_principle}):
\[
\mu_{n+1}(C\Omega) = \int_0^1 \mu_n\big((1-t)\Omega\big) \,dt = \mu_n(\Omega) \int_0^1 (1-t)^n \,dt = \mu_n(\Omega) \left[ -\frac{(1-t)^{n+1}}{n+1} \right]_0^1 = \frac{\mu_n(\Omega)}{n+1}.
\]
\end{proof}

\begin{center}
% Generator: Gemini 3.6 Flash (High) --- FIG-CONE: Cavalieri's principle for the volume of a cone over a domain
\begin{tikzpicture}[>=Latex, scale=1.0]

  % Axes
  \draw[->, thick] (-0.5,0) -- (3.2,0) node[right, font=\footnotesize] {$\mathbb{R}^n$};
  \draw[->, thick] (0,-0.3) -- (0,2.6) node[above, font=\footnotesize] {$t \in \mathbb{R}$};

  % Vertex at (0, 1.8)
  \fill[red!80!black] (0, 1.8) circle (1.8pt) node[left, font=\footnotesize] {$(0,1)$};

  % Base domain Omega at t=0
  \draw[very thick, orange!90!black] (0.8,0) -- (2.6,0);
  \fill[orange!90!black] (0.8,0) circle (1.5pt);
  \fill[orange!90!black] (2.6,0) circle (1.5pt);
  \node[orange!90!black, font=\footnotesize, below] at (1.7, -0.05) {$\Omega$};

  % Cone sides
  \draw[blue!75!black, thick] (0, 1.8) -- (0.8,0);
  \draw[blue!75!black, thick] (0, 1.8) -- (2.6,0);

  % Shaded region for C\Omega
  \fill[blue!10, opacity=0.7] (0, 1.8) -- (0.8,0) -- (2.6,0) -- cycle;
  \node[blue!80!black, font=\small] at (1.1, 0.7) {$C\Omega$};

  % Cross-section at height t
  \draw[dashed, gray!70!black] (-0.3, 1.0) -- (2.8, 1.0);
  \draw[very thick, green!50!black] (0.355, 1.0) -- (1.155, 1.0);
  \fill[green!50!black] (0.355, 1.0) circle (1.4pt);
  \fill[green!50!black] (1.155, 1.0) circle (1.4pt);
  \node[green!50!black, font=\scriptsize, above] at (0.75, 1.05) {$(1-t)\Omega$};

  \node[font=\scriptsize, left] at (0, 1.0) {$t$};
  \node[font=\scriptsize, left] at (0, 0.0) {$0$};

\end{tikzpicture}
\end{center}

% Source: Corsin Nick/Class Notes/Week 8.pdf, p. 8
\begin{example}[Fubini on the unit square]
\label{ex:fubini_unit_square}
\[
\begin{aligned}
\int_{[0,1]^2} x^{y+1}\,dx\,dy
  &= \int_0^1 dy \int_0^1 dx\ x^{y+1}
   = \int_0^1 dy\ \left[\frac{x^{y+2}}{y+2}\right]_0^1 \\
  &= \int_0^1 dy\ \frac{1}{y+2}
   = \log|y+2|\,\Big|_0^1
   = \log\!\left(\tfrac32\right).
\end{aligned}
\]
\end{example}

\begin{ainote}
Corsin writes the first line as $\int_0^1 dx \int_0^1 dy$, but the antiderivative taken on the
next line is in $x$, and it produces $x^{y+2}/(y+2)$, evaluated from $0$ to $1$. The
differentials are therefore swapped relative to the computation actually performed. The order is corrected
above and the result is unaffected.

The example is well chosen: integrating $x^{y+1}$ in $y$ first would mean integrating an
exponential in the exponent, which is far worse. Fubini's real use is that you may pick the
easier order, and here only one of the two is comfortable.
\end{ainote}

% Supplement: Sascha Brack/Class Notes/Week_08_Notes_Friday.pdf, pp. 4--5
\begin{example}[Changing the order of integration]
\label{ex:change_order_integration}
Consider the iterated integral
\[\int_0^2 \int_{y^2}^{2\sqrt{2y}} f(x,y) \,dx\,dy.\]
Let us rewrite it with the integration order reversed (first $y$, then $x$).

The region is defined by $0 \leq y \leq 2$ and $y^2 \leq x \leq 2\sqrt{2y}$.
Notice that for $y \in [0,2]$, the lower bound is $x \geq 0$ and the upper bound is $x \leq 2\sqrt{4} = 4$, so $x$ ranges from $0$ to $4$.
We now solve the inequalities for $y$:
$x \leq 2\sqrt{2y} \implies x^2 \leq 8y \implies y \geq \frac{1}{8}x^2$.
$x \geq y^2 \implies y \leq \sqrt{x}$ (since $x,y \geq 0$).
Thus, the new bounds are $0 \leq x \leq 4$ and $\frac{1}{8}x^2 \leq y \leq \sqrt{x}$. The reversed integral is:
\[\int_0^4 \int_{\frac{1}{8}x^2}^{\sqrt{x}} f(x,y) \,dy\,dx.\]
\exinfo{The same region, with the same two bounds, is Problem 14 of the Analysis II examination of
August 2025 (Prof.\ Laura Kobel-Keller), a box problem where only the reversed integral was
graded.}
\end{example}

% Generator: Claude Opus 5 (High) --- FIG-20-01, the two slicings of one region.
% Requested in feedback/brainstorming_improvements.md, which names swapping bounds over a
% non-rectangular domain as a major pain point and asks for exactly this picture. Drawn against
% the region of ex:change_order_integration rather than the brainstorm's triangle, because a
% triangle makes the two orders look symmetric and this region does not: the curved boundaries
% are what force the algebra in the example above.
%
% Arithmetic check. The two boundary curves are the SAME two curves, read two ways:
%   x = y^2        <=>  y = sqrt(x)      (upper boundary)
%   x = 2 sqrt(2y) <=>  y = x^2/8        (lower boundary)
% They meet where sqrt(x) = x^2/8, i.e. x^{3/2} = 8, x = 4, y = 2 -- and at the origin.
% Both panels must therefore show the SAME shaded region; only the slice differs.
%   LEFT (dx dy, outer y): horizontal slice at y = 1.2, running x = y^2 = 1.44
%       to x = 2 sqrt(2.4) = 3.0984.
%   RIGHT (dy dx, outer x): vertical slice at x = 2.2, running y = 2.2^2/8 = 0.605
%       to y = sqrt(2.2) = 1.4832.
% Both slice endpoints lie on the boundary curves by construction, which is the property the
% figure asserts and the one worth re-checking if the coordinates are ever touched.
\begin{center}
\begin{tikzpicture}[x=1.05cm, y=1.05cm, >=Latex,
    bdry/.style={blue!70!black, very thick},
    slice/.style={BrickRed, very thick, <->}]

  % ---------------- left panel: horizontal slices, the order as given ----------------
  \begin{scope}
    \fill[blue!8]
      plot[domain=0:4, samples=60] ({\x}, {sqrt(\x)}) --
      plot[domain=4:0, samples=60] ({\x}, {(\x)*(\x)/8}) -- cycle;

    \draw[->, gray!70] (-0.25,0) -- (4.9,0) node[right, font=\footnotesize, text=black] {$x$};
    \draw[->, gray!70] (0,-0.25) -- (0,2.7) node[above, font=\footnotesize, text=black] {$y$};

    \draw[bdry] plot[domain=0:4, samples=60] ({\x}, {sqrt(\x)});
    \draw[bdry] plot[domain=0:4, samples=60] ({\x}, {(\x)*(\x)/8});
    % Placed at x = 3.4 on the upper curve (y = sqrt(3.4) = 1.844) rather than near x = 1: the
    % dashed guide at y = 1.2 runs across the left half of this panel and the label collided
    % with it there.
    \node[above left, font=\scriptsize, blue!70!black, inner sep=1pt] at (3.4,1.844) {$x=y^2$};
    \node[below right, font=\scriptsize, blue!70!black, inner sep=1pt] at (3.05,1.16)
      {$x=2\sqrt{2y}$};

    \draw[slice] (1.44,1.2) -- (3.0984,1.2);
    \draw[gray!60, dashed] (0,1.2) -- (1.44,1.2);
    \node[left, font=\scriptsize] at (-0.1,1.2) {$y$};

    \fill (4,2) circle (1.4pt);
    \node[right, font=\scriptsize, inner sep=2pt] at (4.05,2.1) {$(4,2)$};

    \node[font=\scriptsize, BrickRed] (Lslab) at (1.05,2.5) {$y^2 \le x \le 2\sqrt{2y}$};
    \draw[BrickRed, thin] (Lslab) -- (1.9,1.28);
    \node[font=\footnotesize, align=center] at (2.0,-0.95)
      {$\displaystyle\int_0^2\!\!\int_{y^2}^{2\sqrt{2y}} f \,dx\,dy$};
    \node[font=\scriptsize, gray!50!black] at (2.0,-1.75) {fix $y$, sweep $x$};
  \end{scope}

  % ---------------- right panel: vertical slices, the order after swapping ------------
  \begin{scope}[shift={(6.6,0)}]
    \fill[blue!8]
      plot[domain=0:4, samples=60] ({\x}, {sqrt(\x)}) --
      plot[domain=4:0, samples=60] ({\x}, {(\x)*(\x)/8}) -- cycle;

    \draw[->, gray!70] (-0.25,0) -- (4.9,0) node[right, font=\footnotesize, text=black] {$x$};
    \draw[->, gray!70] (0,-0.25) -- (0,2.7) node[above, font=\footnotesize, text=black] {$y$};

    \draw[bdry] plot[domain=0:4, samples=60] ({\x}, {sqrt(\x)});
    \draw[bdry] plot[domain=0:4, samples=60] ({\x}, {(\x)*(\x)/8});
    \node[above left, font=\scriptsize, blue!70!black, inner sep=1pt] at (1.15,1.07)
      {$y=\sqrt{x}$};
    \node[below right, font=\scriptsize, blue!70!black, inner sep=1pt] at (3.05,1.16)
      {$y=x^2/8$};

    \draw[slice] (2.2,0.605) -- (2.2,1.4832);
    \draw[gray!60, dashed] (2.2,0) -- (2.2,0.605);
    \node[below, font=\scriptsize] at (2.2,-0.08) {$x$};

    \fill (4,2) circle (1.4pt);
    \node[right, font=\scriptsize, inner sep=2pt] at (4.05,2.1) {$(4,2)$};

    \node[font=\scriptsize, BrickRed] (Rslab) at (1.15,2.5) {$\tfrac{x^2}{8} \le y \le \sqrt{x}$};
    \draw[BrickRed, thin] (Rslab) -- (2.2,1.55);
    \node[font=\footnotesize, align=center] at (2.0,-0.95)
      {$\displaystyle\int_0^4\!\!\int_{x^2/8}^{\sqrt{x}} f \,dy\,dx$};
    \node[font=\scriptsize, gray!50!black] at (2.0,-1.75) {fix $x$, sweep $y$};
  \end{scope}
\end{tikzpicture}
\end{center}

\begin{remark}[Reading the swap off the picture]
\label{rem:fubini_swap_picture}
The region is the same in both panels; only the direction of the slice changes, and with it which
variable is the outer one. That is the whole of Fubini's theorem for a non-rectangular domain, and
the two mechanical rules follow from looking:

\begin{itemize}
  \item \textbf{The outer limits are constants, and they are the shadow of the region on the outer
    axis.} On the left the outer variable is $y$ and the region's shadow on the $y$-axis is
    $[0,2]$; on the right the outer variable is $x$ and its shadow is $[0,4]$. The two are
    different numbers, which is the most common slip.
  \item \textbf{The inner limits are functions of the outer variable, read off where the slice
    enters and leaves the region.} Entering and leaving happens on the \emph{same two curves} in
    both panels --- there is only one pair of boundary curves here. What changes is which of
    $x = y^2$ and $y = \sqrt{x}$ is the useful way to write the first of them.
\end{itemize}

The reason the algebra in the example above is not merely a relabelling is the second bullet: each
boundary has to be re-solved for the new inner variable. A region whose boundary cannot be solved
for one of the variables cannot be swapped at all without first splitting it, and that is the
situation \cref{ex:change_of_variables_domain} runs into.
\end{remark}

\begin{example}[Surface parametrization and Fubini]
\label{ex:surface_parametrization_integral}
Let $f(u,v) := (u\cos v, u\sin v, u^2)\transp \in \mathbb{R}^3$ for $(u,v) \in \mathbb{R}^2$. Write the third coordinate $z(u,v)$ as a function of the first two, $\tilde{z}(x,y)$, and compute $\int_0^y \int_0^x \tilde{z}(x',y') \,dx'\,dy'$.

We have $x(u,v) = u\cos v$ and $y(u,v) = u\sin v$. Then $x^2 + y^2 = u^2\cos^2 v + u^2\sin^2 v = u^2 = z(u,v)$.
So, we define $\tilde{z}(x,y) := x^2 + y^2$.
The integral is:
\begin{align}
\int_0^y \int_0^x (x'^2 + y'^2) \,dx'\,dy'
&= \int_0^y \left[ \frac{x'^3}{3} + x'y'^2 \right]_{x'=0}^{x'=x} \,dy' \nonumber \\
&= \int_0^y \left( \frac{x^3}{3} + x y'^2 \right) \,dy' \nonumber \\
&= \left[ \frac{x^3 y'}{3} + x \frac{y'^3}{3} \right]_{y'=0}^{y'=y} \nonumber \\
&= \frac{xy}{3}(x^2 + y^2). \label{eq:surface_integral_result}
\end{align}
\end{example}

% Source: Corsin Nick/Class Notes/Week 8.pdf, pp. 8--10
\begin{example}[A domain that Fubini alone cannot handle]
\label{ex:change_of_variables_domain}
Consider
\[\int_A \frac{y}{\sqrt{x}}\,dy\,dx, \qquad
A := \left\{(x,y) \in \mathbb{R}^2 \ :\ x,y>0,\ \ 1 \leq \frac{y}{\sqrt x} \leq 2,\ \ 1 \leq xy \leq 2\right\}.\]
This is not easily integrated with Fubini alone. But we can apply a transformation of variables
to simplify the integration domain $A$: take
\[\Phi : U \to V, \qquad \begin{pmatrix}x\\y\end{pmatrix} \mapsto
  \begin{pmatrix}u(x,y)\\v(x,y)\end{pmatrix} = \begin{pmatrix}\tfrac{y}{\sqrt x}\\[2pt] xy\end{pmatrix},\]
with $U$ and $V$ still to be determined, subject only to their being open.

\textbf{Check injectivity.} Suppose, for $x,y,x',y' > 0$,
\[\text{\textbf{(1)}}\ \ \frac{y}{\sqrt x} = \frac{y'}{\sqrt{x'}},
\qquad\qquad \text{\textbf{(2)}}\ \ xy = x'y'.\]
From \textbf{(2)}, $\dfrac{x}{x'} = \dfrac{y'}{y}$. Inserting into \textbf{(1)},
\[\frac{1}{\sqrt x} = \frac{y'}{y}\cdot\frac{1}{\sqrt{x'}} = \frac{x}{x'}\cdot\frac{1}{\sqrt{x'}}
\quad\Longrightarrow\quad x^{3/2} = (x')^{3/2} \quad\Longrightarrow\quad x = x',\]
and consequently $\dfrac{y'}{y} = \dfrac{x}{x'} = 1$. So, for $U := (0,\infty)^2$ and
$V := \Phi(U)$, the map $\Phi$ is a bijection between open sets.

\textbf{Functional determinant.} We calculate $|\det \Jac\Phi|$; if it is strictly positive we
may proceed, since $\Phi$ is then a local diffeomorphism as well as a bijection, hence a
diffeomorphism. Now
\[\Jac\Phi(x,y) = \begin{pmatrix} -\dfrac{y}{2\sqrt{x}^{\,3}} & \dfrac{1}{\sqrt x} \\[2ex] y & x \end{pmatrix},\]
so that
\[\big|\det \Jac\Phi(x,y)\big| = \left|-\frac32\,\frac{y}{\sqrt x}\right| = \frac32\,\frac{y}{\sqrt x} = \frac32\,u(x,y).\]
So, $\Phi$ is a diffeomorphism and $du\,dv = \tfrac32 u\,dx\,dy$. Furthermore, by design,
\[\Phi(A) = (1,2)^2.\]
So, with the change of variables (\cref{thm:change_of_variables}), followed by Fubini
(\cref{thm:fubini}),
\[\int_A \frac{y}{\sqrt x}\,dy\,dx
  = \int_{(1,2)^2} u\,\frac{1}{\tfrac32 u}\,du\,dv
  = \frac23\int_1^2 dv \int_1^2 du
  = \frac23.\]
\end{example}

\begin{ainote}
The design principle is worth stating, because it is the reverse of how substitution is usually
taught. Here $\Phi$ was not chosen to simplify the \emph{integrand} but to straighten the
\emph{domain}: the two constraints defining $A$ are literally $1 \leq u \leq 2$ and
$1 \leq v \leq 2$, so $\Phi$ turns an awkward curvilinear region into a square, and only then
does Fubini become available. That the integrand $y/\sqrt x$ is itself $u$, and cancels against
most of the Jacobian factor, is a bonus of the same choice.
\end{ainote}

\begin{aiexercise}[An order of integration that cannot be done]
% Generator: Claude Opus 5 (High)
\label{ex:ai_order_swap_gaussian}
Consider
\[I := \int_0^1\!\!\int_x^1 e^{y^2}\,dy\,dx.\]
\begin{enumerate}[label=\textbf{(\alph*)}]
  \item Explain what goes wrong if one attempts the inner integral as written.
  \item Describe the region of integration by inequalities, and rewrite $I$ with the two
    integrations in the opposite order.
  \item Evaluate $I$.
\end{enumerate}
\end{aiexercise}

% Source: exercises_2024/mock.pdf (Analysis II mock examination, Spring Semester 2024,
% Prof. Joaquim Serra), Exercise 12, p. 3
% Extractor: Claude Opus 5 (High)
% Filed with Fubini rather than with the improper integrals of chapter 19 because the whole
% problem is the product splitting; the one-dimensional pieces are then immediate. Part (a) is the
% trap and is what earns the block its place: the printed value sqrt(pi) is the answer to a
% one-dimensional Gaussian, and alpha = 0 is exactly the case where the z-integral does not exist.
\begin{exercise}[True or False \textnormal{(a Gaussian slab, and what happens at $\alpha=0$)}]
\label{ex:gaussian_slab_alpha_asymptotics}
For $\alpha \geq 0$ consider
\[I_\alpha := \int_{\mathbb{R}^3} e^{-x^2-y^2-\alpha\lvert z\rvert}\,dx\,dy\,dz .\]
Decide whether each of the following is true or false, with justification.
\begin{enumerate}[label=\textbf{(\alph*)}]
  \item $I_0 = \sqrt{\pi}$.
  \item $\alpha I_\alpha = I_1$.
  \item $\displaystyle\lim_{\alpha\to+\infty}\big(\alpha+\sqrt{\alpha}\big)I_\alpha = 2\pi$.
\end{enumerate}
\exinfo{This exercise is taken from the mock examination for Analysis II, Spring Semester 2024
(Prof.\ Joaquim Serra), where it appears as the twelfth block of multiple-choice questions. Its
paper supplies $\int_{\mathbb{R}}e^{-t^2}\,dt = \sqrt{\pi}$ as a formula that may be used without
proof; \cref{ex:gaussian_integral_polar} derives it.}
\end{exercise}

% Sonnet 5 (Medium)
We now want to look at a powerful new integration technique using ODEs.

% Originally: content/week-08.tex, \section{Feynman's trick} (Corsin Week 8)

% Source: Corsin Nick/Class Notes/Week 8.pdf, p. 11
\section{Feynman's trick}
\label{sec:feynmans_trick}

\begin{theorem}[Differentiation under the integral sign]
\label{thm:feynmans_trick}
Let $U \subseteq \mathbb{R}^n\times\mathbb{R}$ be open, $f \in C^1(U)$, and let
$K \subseteq \mathbb{R}^n$ be compact such that $K\times[a,b] \subseteq U$. Then the function
\[(a,b) \to \mathbb{R}, \qquad y \mapsto \int_K f(x,y)\,dx\]
is continuously differentiable, and
\[\frac{d}{dy}\int_K f(x,y)\,dx = \int_K \frac{\partial}{\partial y}f(x,y)\,dx
\qquad \text{for all } y \in (a,b).\]
\end{theorem}

% Source: Corsin Nick/Class Notes/Week 8.pdf, pp. 11--13
\begin{example}[$\int_0^1 \frac{\log(1+x)}{1+x^2}\,dx$ by Feynman's trick]
\label{ex:feynman_log_integral}
We want to compute the integral
\[I(1) = \int_0^1 \frac{\log(1+x)}{1+x^2}\,dx,\]
which has no convenient antiderivative. By introducing a parameter $\alpha \in (-1,1)$:
\[I(\alpha) := \int_0^1 \frac{\log(1+\alpha x)}{1+x^2}\,dx.\]
Note that the integrand is a $C^1$ function on
$U := \big(-\tfrac{1}{1+\varepsilon},\ \tfrac1\varepsilon\big) \times (-\varepsilon,\ 1+\varepsilon)$
in the variables $(x,\alpha)$, with $K := [0,1]$ and $[a,b] := [0,1]$ satisfying
$K\times[a,b] \subseteq U$ for $0 < \varepsilon < 1$ (pay attention to the logarithm, which is
what forces $1+\alpha x > 0$). So, we may differentiate under the integral!

\textbf{We obtain an ordinary differential equation.} By \cref{thm:feynmans_trick},
\[I'(\alpha) = \frac{d}{d\alpha}\int_0^1 \frac{\log(1+\alpha x)}{1+x^2}\,dx
  = \int_0^1 \frac{x}{(1+\alpha x)(1+x^2)}\,dx.\]
Partial fractions: writing
\[\frac{x}{(1+\alpha x)(1+x^2)} = \frac{A}{1+\alpha x} + \frac{Bx+C}{1+x^2}\]
and clearing denominators gives $A + Ax^2 + Bx + B\alpha x^2 + C + C\alpha x = x$, so comparing
coefficients,
\[-A = C, \qquad -\tfrac{A}{\alpha} = B, \qquad B + C\alpha = 1
  \implies -\tfrac{A}{\alpha} - A\alpha = 1,\]
whence
\[A = -\frac{\alpha}{1+\alpha^2}, \qquad B = \frac{1}{1+\alpha^2}, \qquad C = \frac{\alpha}{1+\alpha^2}.\]
Therefore
\[
\begin{aligned}
I'(\alpha) &= \left(\int_0^1 \frac{-\alpha}{1+\alpha x} + \frac{x+\alpha}{1+x^2}\,dx\right)\frac{1}{1+\alpha^2} \\
&= \frac{1}{1+\alpha^2}\Big[-\log(1+\alpha x) + \tfrac12\log(1+x^2) + \alpha\arctan(x)\Big]_0^1 \\
&= \frac{1}{1+\alpha^2}\Big[\log\sqrt2 + \tfrac{\alpha\pi}{4} - \log(1+\alpha)\Big].
\end{aligned}
\]

\textbf{With initial value} $I(0) = \int_0^1 \frac{\log(1)}{1+x^2}\,dx = 0$, we obtain from the
fundamental theorem of calculus
\[
\begin{aligned}
\int_0^1 \frac{\log(1+x)}{1+x^2}\,dx = I(1) &= I(1)-I(0) \\
&= \int_0^1\left(\frac{\log\sqrt2}{1+\alpha^2} + \frac{\pi}{4}\,\frac{\alpha}{1+\alpha^2}
   - \underbrace{\frac{\log(1+\alpha)}{1+\alpha^2}}_{\textstyle \to\, I(1)}\right)d\alpha.
\end{aligned}
\]
The last term is $I(1)$ again, the integral we started from with $\alpha$ merely renamed, so
bringing it to the left-hand side,
\[
\begin{aligned}
2\,I(1) &= \Big[\log\sqrt2\,\arctan(\alpha) + \tfrac{\pi}{8}\log(1+\alpha^2)\Big]_0^1
 = \frac{\pi\log\sqrt2}{4} + \frac{\pi\log\sqrt2}{4}, \\[1ex]
\implies\quad I(1) &= \frac{\pi\log\sqrt2}{4} = \boxed{\ \frac{\pi\log 2}{8}\ }.
\end{aligned}
\]
\end{example}

% Source: Jérôme Paschoud/Notizen/Woche 9.1 Variablenwechsel.pdf, p. 4
% Generator: Gemini 3.1 Pro (High)
\begin{example}[Dirichlet integral via Feynman's trick]
\label{ex:feynman_dirichlet_integral}
Consider the integral $\int_0^\infty \frac{\sin x}{x} \,dx$. We can evaluate it by introducing a damping factor $e^{-tx}$ and defining
\[I(t) := \int_0^\infty e^{-tx} \frac{\sin x}{x} \,dx \qquad \text{for } t > 0.\]
Differentiating under the integral sign with respect to $t$ (which is justified since the integrand and its derivative decay rapidly for $t>0$):
\[I'(t) = \int_0^\infty \frac{\partial}{\partial t} \left( e^{-tx} \frac{\sin x}{x} \right) \,dx = \int_0^\infty -x e^{-tx} \frac{\sin x}{x} \,dx = -\int_0^\infty e^{-tx} \sin x \,dx.\]
This standard integral can be evaluated using integration by parts twice, or by using Euler's formula $e^{ix} = \cos x + i\sin x$ and noting that $e^{-tx}\sin x = \operatorname{Im}(e^{(-t+i)x})$:
\[\int_0^\infty e^{(-t+i)x} \,dx = \left[ \frac{e^{(-t+i)x}}{-t+i} \right]_0^\infty = 0 - \frac{1}{-t+i} = \frac{t+i}{t^2+1}.\]
Taking the imaginary part gives $\int_0^\infty e^{-tx} \sin x \,dx = \frac{1}{t^2+1}$. Thus,
\[I'(t) = -\frac{1}{t^2+1} \implies I(t) = -\arctan(t) + C.\]
Since $\lim_{t \to \infty} I(t) = 0$, we must have $C = \lim_{t \to \infty} \arctan(t) = \frac{\pi}{2}$.
Therefore, $I(t) = \frac{\pi}{2} - \arctan(t)$.
The original integral corresponds to the limit as $t \to 0^+$ (which requires an extra justification step usually provided by Abel's theorem, but we formally take the limit):
\[\int_0^\infty \frac{\sin x}{x} \,dx = \lim_{t \to 0^+} \left( \frac{\pi}{2} - \arctan(t) \right) = \frac{\pi}{2}.\]
\end{example}

\begin{ainote}
Corsin's final integration writes the $\alpha$-term as
$\tfrac{\pi}{2}\cdot\tfrac{\alpha}{1+\alpha^2}$, integrating to $\tfrac{\pi}{4}\log(1+\alpha^2)$.
Both are a factor $2$ away from what $I'(\alpha)$ gives: his own line above has
$\tfrac{\alpha\pi}{4}$, so the term is $\tfrac{\pi}{4}\cdot\tfrac{\alpha}{1+\alpha^2}$ and its
antiderivative is $\tfrac{\pi}{8}\log(1+\alpha^2)$. Corrected above, and the correction is
self-checking, since it is exactly what makes his two bracket terms come out equal, as his own
next line asserts:
$\tfrac{\pi}{8}\log 2 = \tfrac{\pi}{8}\cdot 2\log\sqrt2 = \tfrac{\pi\log\sqrt2}{4}$. His final
answer $\tfrac{\pi\log 2}{8}$ is correct and matches the known value.
\end{ainote}

% Generator: Claude Opus 5 (Medium)
\begin{ainote}[Why this is called a trick, and when to reach for it]
The move that makes this work is not the differentiation. It is noticing that the
\emph{antiderivative} of the resulting integrand is elementary while the original one is not.
The pattern is worth naming:
\begin{enumerate}[label=\textbf{(\alph*)}]
  \item Insert a parameter $\alpha$ so that the target is $I(\alpha_1)$ for some
    $\alpha_1$, and so that $I(\alpha_0)$ is trivial for some other $\alpha_0$. Here
    $I(0)=0$, because $\log 1 = 0$.
  \item Differentiate under the integral. The $\log$ disappears, leaving a rational function
    that partial fractions can handle.
  \item Integrate $I'$ back from $\alpha_0$ to $\alpha_1$.
\end{enumerate}
The self-reference in step 3, where the answer reappears on the right-hand side so that one
solves for $I(1)$ rather than computing it, is specific to this integral and is the reason it
is a classic. Do not expect it in general. Usually step 3 is an honest integration.
\end{ainote}

\begin{aiexercise}[Frullani's integral]
% Generator: Claude Opus 5 (High)
\label{ex:ai_frullani}
Fix $0 < a < b$. The aim is to show that
\[\int_0^\infty \frac{e^{-ax} - e^{-bx}}{x}\,dx = \log\!\left(\frac{b}{a}\right),\]
an answer that depends on $a$ and $b$ only through their ratio.
\begin{enumerate}[label=\textbf{(\alph*)}]
  \item The integrand is undefined at $x = 0$. Compute $\lim_{x\to0^+}$ of it, so that the
    integrand extends continuously to $[0,\infty)$.
  \item Hold $a$ fixed and let the second exponent be the parameter: define
    \[I(t) := \int_0^\infty \frac{e^{-ax} - e^{-tx}}{x}\,dx \qquad \text{for } t > 0.\]
    Differentiate under the integral sign and evaluate the resulting integral, which is
    elementary.
  \item Identify the value of $t$ for which $I(t)$ is trivially zero, and integrate $I'$ from
    there to $t = b$.
\end{enumerate}
\end{aiexercise}

% Generator: Claude Opus 5 (High)
% Section added 2026-08-11. Both exercises are transcribed from the Probeprfg mock exams (see the
% % Source: comments below); the definition, the connecting prose and the remarks are written for
% this document, which had no notion of a centroid anywhere. The two problems come from different
% papers and are placed together because each is the other's payoff: one computes the centroid,
% the other says what the centroid is for.

\section{Centroids}
\label{sec:centroids}

% Transition: Claude Opus 5 (High)
Every integral in this chapter so far has been computed for its own sake. The simplest thing one
actually does with an integral over a region is average a coordinate over it, and the resulting
point is worth a name.

\begin{definition}[Centroid of a plane region]
\label{def:centroid}
Let $B \subseteq \mathbb{R}^2$ be Jordan measurable with $\vol_2(B) > 0$. The \newterm{centroid}
\germanfor{Fl\"achenschwerpunkt} of $B$ is the point $z_B := (x_B, y_B)$ with
\[x_B := \frac{1}{\vol_2(B)}\int_B x \,d\vol_2(x,y), \qquad
  y_B := \frac{1}{\vol_2(B)}\int_B y \,d\vol_2(x,y).\]
\end{definition}

The same formula defines the centroid of a Jordan measurable set in any $\mathbb{R}^n$, one
coordinate at a time; nothing below uses the restriction to the plane except the triangle.

\begin{remark}[Why the normalisation is there]
Without the factor $1/\vol_2(B)$ the two integrals would scale with the size of $B$ as well as
with its position, and doubling a region would move its \qt{centre}. Dividing by the area makes
$z_B$ an \emph{average} of the coordinate functions rather than a total, and the immediate
consequence is that $z_B$ is where one expects it to be: for a region symmetric about a line, the
centroid lies on that line, because the coordinate perpendicular to it integrates to zero.
\end{remark}

\begin{center}
% Generator: Claude Opus 5 (High) --- FIG-20-02: the two halves of the exercise below, drawn.
% This section defines a *point* and then asserts two things about where it sits, with no picture
% anywhere: the reader meets the weighted-average formula of part (a) as an identity between two
% integrals and never sees that it puts $z_B$ on the segment joining the two sub-centroids.
% Both panels are exact, not schematic --- the coordinates below are the computed centroids, so
% the picture can be measured against the formulas it illustrates.
% `lbl' backs a label with a soft white plate so it stays readable where it lands on a shaded
% region or a box edge -- the same idiom as the surface figure in ch. 21, and the thing that
% decides whether these read on a grey e-reader screen rather than only at 300dpi in colour.
\begin{tikzpicture}[>=Latex, scale=1.0,
  lbl/.style={fill=white, fill opacity=0.8, text opacity=1, inner sep=1.2pt,
              rounded corners=1pt}]

  % ---------- Left: centroids add, as an area-weighted average ----------
  % B_1 = [0,2]x[0,1.8], area 3.6, centroid (1, 0.9)
  % B_2 = [2,3]x[0,0.9], area 0.9, centroid (2.5, 0.45)
  % weight of B_2 is 0.9/4.5 = 1/5 exactly, so
  % z_B = z_1 + (1/5)(z_2 - z_1) = (1.3, 0.81)
  \node[font=\small\bfseries, text=TextBoldColor] at (1.5, 3.0) {Centroids add};

  \fill[blue!10, opacity=0.7] (0,0) rectangle (2,1.8);
  \fill[orange!12, opacity=0.7] (2,0) rectangle (3,0.9);
  \draw[blue!75!black, thick] (0,0) rectangle (2,1.8);
  \draw[orange!90!black, thick] (2,0) rectangle (3,0.9);
  \node[blue!80!black, font=\scriptsize] at (0.45, 1.5) {$B_1$};
  \node[orange!90!black, font=\scriptsize] at (2.75, 0.68) {$B_2$};

  % the segment joining the two sub-centroids, and z_B on it
  \draw[gray!70!black, dashed] (1,0.9) -- (2.5,0.45);
  \fill[blue!75!black] (1,0.9) circle (1.5pt)
    node[lbl, above left=-2pt, font=\scriptsize] {$z_{B_1}$};
  \fill[orange!90!black] (2.5,0.45) circle (1.5pt)
    node[lbl, below right=-2pt, font=\scriptsize] {$z_{B_2}$};
  \fill[red!80!black] (1.3,0.81) circle (1.8pt);
  % text=, not a bare colour, which would override the fill=white in `lbl'.
  \node[lbl, text=red!80!black, font=\scriptsize, above right=-2pt] at (1.3,0.81) {$z_B$};

  \node[font=\scriptsize, align=center] at (1.5,-0.75)
    {$\vol_2(B_2)/\vol_2(B) = \tfrac15$, and $z_B$ sits\\one fifth of the way along};

  % ---------- Right: the centroid of a triangle ----------
  % o=(0,0), p=(3,0), q=(1.2,2.4); centroid = (4.2/3, 2.4/3) = (1.4, 0.8)
  \begin{scope}[xshift=5.0cm]
    \node[font=\small\bfseries, text=TextBoldColor] at (1.5, 3.0) {\dots\ and a triangle's};

    \fill[blue!10, opacity=0.7] (0,0) -- (3,0) -- (1.2,2.4) -- cycle;
    \draw[blue!75!black, thick] (0,0) -- (3,0) -- (1.2,2.4) -- cycle;

    % the three medians, which meet at the centroid
    \draw[gray!70!black, thin] (0,0) -- (2.1,1.2);
    \draw[gray!70!black, thin] (3,0) -- (0.6,1.2);
    \draw[gray!70!black, thin] (1.2,2.4) -- (1.5,0);

    % midpoints
    \fill[gray!70!black] (2.1,1.2) circle (1.0pt);
    \fill[gray!70!black] (0.6,1.2) circle (1.0pt);
    \fill[gray!70!black] (1.5,0) circle (1.0pt);

    \fill[black] (0,0) circle (1.5pt) node[below left=-2pt, font=\scriptsize] {$o$};
    \fill[black] (3,0) circle (1.5pt) node[below right=-2pt, font=\scriptsize] {$p$};
    \fill[black] (1.2,2.4) circle (1.5pt) node[above=1pt, font=\scriptsize] {$q$};

    \fill[red!80!black] (1.4,0.8) circle (1.8pt);
    \node[lbl, text=red!80!black, font=\scriptsize, right=1pt] at (1.4,0.8) {$z_B$};

    \node[font=\scriptsize, align=center] at (1.5,-0.75)
      {the three medians meet at\\$z_B = \tfrac13(o+p+q)$};
  \end{scope}

\end{tikzpicture}
\end{center}

% Source: old_exams/fs2023/Probeprfg3.pdf (Probeprüfung Analysis II, 31 January 2023),
% Aufgabe 11, p. 4
% Extractor: Claude Opus 5 (High)
% The paper names no lecturer; see old-exam-mining.md.
\begin{exercise}[Centroids add, and the centroid of a triangle]
\label{ex:centroid_additivity_and_triangle}
Let $B \subseteq \mathbb{R}^2$ be Jordan measurable with $\vol_2(B) > 0$.
\begin{enumerate}[label=\textbf{(\alph*)}]
  \item Suppose $B = \bigcup_{i=1}^n B_i$ for pairwise disjoint Jordan measurable sets
    $B_1,\dots,B_n \subseteq \mathbb{R}^2$, each with $\vol_2(B_i) > 0$. Show that
    \[z_B = \sum_{i=1}^n \frac{\vol_2(B_i)}{\vol_2(B)}\,z_{B_i}.\]
  \item Show that if $B$ is a triangle with vertices $o = (0,0)$, $p = (r,0)$ and $q = (s,t)$,
    where $r,s,t > 0$, then $z_B = \tfrac13(o+p+q)$.
\end{enumerate}
\exinfo{This exercise is taken from the mock examination \emph{Probepr\"ufung Analysis II} of
31 January 2023, Problem 11. That paper names no lecturer.}
\exsol{sol:centroid_additivity_and_triangle}
\end{exercise}

% Source: old_exams/fs2023/Probeprfg2.pdf (Probeprüfung Analysis II, 17 August 2022, revised),
% Aufgabe 1, p. 1
% Extractor: Claude Opus 5 (High)
% The paper names no lecturer; see old-exam-mining.md. The original asks only for the minimising
% pair; the two-part split, and part (b) in particular, are added here, because the general
% statement is what makes the computation worth doing and the exam's own solution differentiates
% under the integral sign without ever saying that the answer is the centroid.
\begin{exercise}[What the centroid minimises]
\label{ex:centroid_minimises_second_moment}
Let $B := \{(x,y) \in \mathbb{R}^2 \mid x \geq 0,\ y \geq 0,\ x^2+y^2 \leq 1\}$ be the closed
quarter disc, and for $(s,t) \in \mathbb{R}^2$ put
\[I(s,t) := \int_B (x-s)^2 + (y-t)^2 \,d\vol_2(x,y).\]
\begin{enumerate}[label=\textbf{(\alph*)}]
  \item Determine the pair $(s,t)$ at which $I$ is minimal.
  \item Show that the answer to part \textbf{(a)} is no accident: for \emph{any} Jordan
    measurable $B \subseteq \mathbb{R}^2$ with $\vol_2(B) > 0$, the function $I$ is minimised
    exactly at $(s,t) = z_B$, and at no other point.
\end{enumerate}
\exinfo{Part \textbf{(a)} is taken from the mock examination \emph{Probepr\"ufung Analysis II} of
17 August 2022 (revised), Problem 1; that paper names no lecturer. Part \textbf{(b)} is added
here, since the general statement is what the computation is really about.}
\exsol{sol:centroid_minimises_second_moment}
\end{exercise}

% Originally: content/week-08.tex had a \section{Solutions} holding nothing but a placeholder
% note, which had already gone stale, since the statements were in fact transcribed. So 8.5
% stood in the document unsolved. Written here for the first time.
\section{Solutions}
\label{sec:change_of_variables_solutions}

\begin{exercisesolution}[Solution to \cref{ex:8.5}]
% Generator: Claude Opus 5 (High)
In each part the work is the same two steps: find $\lvert\det\Jac\rvert$, then translate every
constraint. As the exercise says, no integral is actually evaluated.
\begin{enumerate}[label=\textbf{\arabic*.}]
  \item \textbf{Polar coordinates.} Here $\Phi(r,\theta) := (r\cos\theta, r\sin\theta)$ has
    $\det\Jac\Phi = r > 0$ (\cref{ex:polar_spherical_coordinates}), so
    $dx_1dx_2 = r\,dr\,d\theta$. The three constraints translate one at a time:
    \begin{itemize}
      \item $1 < x_1^2+x_2^2 < 4$ becomes $1 < r < 2$;
      \item $x_1 > 0$ becomes $\cos\theta > 0$, i.e.\ $\theta \in (-\tfrac\pi2, \tfrac\pi2)$;
      \item $x_2 < x_1$ becomes $r\sin\theta < r\cos\theta$; since $r > 0$ and $\cos\theta > 0$
        this is $\tan\theta < 1$, i.e.\ $\theta < \tfrac\pi4$.
    \end{itemize}
    So the transformed domain is the rectangle
    $(r,\theta) \in (1,2)\times\big(-\tfrac\pi2,\tfrac\pi4\big)$ and
    \[\int_A x_1^2\sin(x_2)\,dx_1dx_2
      = \int_{1}^{2}\!\!\int_{-\pi/2}^{\pi/4} r^3\cos^2\theta\,\sin(r\sin\theta)\,d\theta\,dr,\]
    the extra power of $r$ coming from $x_1^2 = r^2\cos^2\theta$ times the Jacobian factor $r$.
  \item \textbf{$u := xy$, $v := x^2$.} Follow Sascha Brack's hint and fix the signs first. On $B$
    the constraint $x^2 < y$ forces $y > 0$, and then $xy > 1 > 0$ forces $x > 0$; so
    $u > 0$ and $v > 0$ throughout, and $x = \sqrt v$, $y = u/\sqrt v$.

    Differentiating in the convenient direction,
    \[\frac{\partial(u,v)}{\partial(x,y)}
      = \det\begin{pmatrix} y & x \\ 2x & 0\end{pmatrix} = -2x^2 = -2v,
      \qquad\text{so}\qquad dx\,dy = \frac{du\,dv}{2v}.\]
    Now the constraints. $1 < xy < 2$ becomes $1 < u < 2$ at once. The second is the interesting
    one: $x^2 < y < 2x^2$ reads $v < u/\sqrt v < 2v$, i.e.
    \[v^{3/2} < u < 2\,v^{3/2}, \qquad\text{equivalently}\qquad
      \Big(\frac{u}{2}\Big)^{2/3} < v < u^{2/3}.\]
    The integrand becomes $y^2e^{-xy} = (u^2/v)\,e^{-u}$, so
    \[\int_B y^2e^{-xy}\,dx\,dy
      = \int_{1}^{2}\!\!\int_{(u/2)^{2/3}}^{u^{2/3}} \frac{u^2e^{-u}}{2v^2}\,dv\,du.\]
    Note what did \emph{not} happen: straightening $1 < xy < 2$ into a pair of constant bounds
    did nothing for the other constraint, which stays curved. Straightening one constraint does
    not oblige the rest to cooperate.
  \item \textbf{A linear change of variables.} With $u := z$, $v := z-2y$, $w := 3x+y+z$,
    \[\frac{\partial(u,v,w)}{\partial(x,y,z)}
      = \det\begin{pmatrix} 0 & 0 & 1 \\ 0 & -2 & 1 \\ 3 & 1 & 1\end{pmatrix} = 6,
      \qquad\text{so}\qquad dx\,dy\,dz = \frac{du\,dv\,dw}{6},\]
    a constant, since the map is linear. Each constraint already \emph{is} one of the new
    coordinates, so the transformed domain is literally the box
    \[(u,v,w) \in (0,1)\times(1,2)\times(-2,0).\]
    Inverting the substitution gives $z = u$, $y = \tfrac{u-v}{2}$ and
    $x = \tfrac{1}{6}\big(2w - 3u + v\big)$, so $xyz = \tfrac{1}{12}\,u\,(u-v)\,(2w-3u+v)$ and
    \[\int_C xyz\,dx\,dy\,dz
      = \int_{0}^{1}\!\!\int_{1}^{2}\!\!\int_{-2}^{0}
        \frac{u\,(u-v)\,(2w-3u+v)}{72}\,dw\,dv\,du.\]
\end{enumerate}

\begin{ainote}
The three parts are graded deliberately, and comparing them is the lesson. In part 3 the change
of variables was \emph{chosen} to be the constraints, so the domain is a box and the Jacobian a
constant, and nothing is left to do. In part 1 the domain is a box in the new coordinates but the
Jacobian is not constant. In part 2 neither holds: the Jacobian varies and one of the two
constraints stays curved. Reading the substitution off the constraints, as in part 3, is the move
worth remembering, and it is available far more often than it looks.
\end{ainote}
\end{exercisesolution}

\begin{exercisesolution}[Solution to \cref{ex:ai_triangle_substitution}]
% Generator: Claude Opus 5 (High)
\textbf{1. The new domain and the Jacobian factor.} The map $\Psi(x,y) := (y-x,\ y+x)$ is
linear and invertible, with
\[\Psi^{-1}(u,v) = (x,y) = \left(\frac{v-u}{2},\ \frac{u+v}{2}\right),
  \qquad
  \det\Jac\Psi^{-1} = \det\begin{pmatrix} -1/2 & 1/2 \\ 1/2 & 1/2\end{pmatrix} = -\frac12,\]
so $dx\,dy = \tfrac12\,du\,dv$. The three sides of $T$ translate one at a time: $x \geq 0$ reads
$u \leq v$, the condition $y \geq 0$ reads $u \geq -v$, and $x+y \leq 1$ reads $v \leq 1$. Taking
strict inequalities for the interior gives exactly $0 < v < 1$ and $-v < u < v$, a triangle
again, but one whose horizontal slices run in the direction the integrand likes.

\textbf{2. The value.} On that region the integrand is $e^{u/v}$, and $v$ is constant along the
inner integration, so the inner integral is an exponential in $u$:
\[I = \frac12\int_0^1\!\!\int_{-v}^{v} e^{u/v}\,du\,dv
    = \frac12\int_0^1 \Big[v\,e^{u/v}\Big]_{u=-v}^{u=v}\,dv
    = \frac{e - e^{-1}}{2}\int_0^1 v\,dv
    = \frac{e - e^{-1}}{4}.\]

\textbf{3. The origin.} A single point is a Jordan null set, so removing it changes neither the
Jordan measurability of $T$ nor the value of the integral. What does need checking is that the
integral exists at all, and it does, because the integrand is \emph{bounded}: on $T$ one has
$\lvert y-x\rvert \leq y+x$, so the exponent $u/v$ lies in $[-1,1]$ and the integrand lies
between $e^{-1}$ and $e$. It is the boundedness, and not the smallness of the bad set on its
own, that rescues the computation.

\begin{remark}[Two ways to choose a substitution]
\label{rem:substitution_from_the_integrand}
\Cref{ex:change_of_variables_domain} chose $\Phi$ to straighten the \emph{domain}. Here the
choice comes from the \emph{integrand}, which is a function of $y-x$ and $y+x$ and of nothing
else. That the domain also becomes pleasant is a bonus, and a mild one: the new region is still
a triangle. What made the integral collapse is that the inner integration now runs along the
direction in which the integrand is a plain exponential.
\end{remark}
\end{exercisesolution}

\begin{exercisesolution}[Solution to \cref{ex:ai_order_swap_gaussian}]
% Generator: Claude Opus 5 (High)
\textbf{1.} The inner integral asks for an antiderivative of $y \mapsto e^{y^2}$, and there is
none in elementary terms. Nothing about the outer integral repairs this, so the order as written
is a dead end.

\textbf{2.} The region is the triangle below the diagonal, and it admits two descriptions:
\[R = \{(x,y) : 0 \leq x \leq 1,\ x \leq y \leq 1\}
    = \{(x,y) : 0 \leq y \leq 1,\ 0 \leq x \leq y\}.\]
The first reads it by vertical slices, the second by horizontal ones. Taking the second and
applying \cref{thm:fubini},
\[I = \int_0^1\!\!\int_0^{y} e^{y^2}\,dx\,dy.\]

\textbf{3.} Now $e^{y^2}$ is constant in $x$, so the inner integral contributes a factor of $y$,
and that factor is exactly what the substitution $s := y^2$ needs:
\[I = \int_0^1 y\,e^{y^2}\,dy = \left[\frac{e^{y^2}}{2}\right]_0^1 = \frac{e-1}{2}.\]

\begin{remark}[What the swap actually bought]
\label{rem:order_swap_supplies_the_factor}
The reversal did not find a cleverer antiderivative. It produced a factor of $y$ that was not
there before, and $y\,e^{y^2}$ integrates where $e^{y^2}$ does not. That is the pattern worth
recognising: when the length of a slice depends on the variable that is causing the trouble,
swapping the order hands you that length as a factor in the integrand.
\end{remark}
\end{exercisesolution}

\begin{exercisesolution}[Solution to \cref{ex:ai_frullani}]
% Generator: Claude Opus 5 (High)
\textbf{1.} Both exponentials equal $1$ at $x = 0$, so the quotient is a difference quotient in
disguise. Expanding, $e^{-ax} - e^{-bx} = (b-a)x + O(x^2)$, and therefore
\[\lim_{x\to0^+}\frac{e^{-ax} - e^{-bx}}{x} = b - a.\]
Assigning that value at $x = 0$ makes the integrand continuous on $[0,\infty)$, and it decays
exponentially, so the improper integral converges.

\textbf{2.} Only the second exponential depends on $t$, and the $x$ that the chain rule produces
cancels the one in the denominator:
\[I'(t) = \int_0^\infty \frac{\partial}{\partial t}\left(\frac{e^{-ax} - e^{-tx}}{x}\right)dx
        = \int_0^\infty \frac{x\,e^{-tx}}{x}\,dx
        = \int_0^\infty e^{-tx}\,dx
        = \frac{1}{t}.\]
That cancellation is the whole trick. Differentiating in the parameter removes the $1/x$ which
made the integrand hard, and leaves an integral of one line.

\textbf{3.} At $t = a$ the two exponentials coincide and the integrand vanishes identically, so
$I(a) = 0$. Consequently, by the fundamental theorem of calculus,
\[\int_0^\infty \frac{e^{-ax} - e^{-bx}}{x}\,dx = I(b) = I(b) - I(a) = \int_a^b \frac{dt}{t}
  = \log\!\left(\frac{b}{a}\right).\]
That the answer sees only the ratio $b/a$ is now visible in the derivation rather than a
surprise, since $I'(t) = 1/t$ is scale-invariant.

\begin{remark}[The hypotheses of \cref{thm:feynmans_trick} do not literally cover this]
\label{rem:feynman_on_a_half_line}
As stated, \cref{thm:feynmans_trick} differentiates under an integral over a \emph{compact} set
$K$, and here the domain is $[0,\infty)$. The step is still legitimate, by the same reasoning
\cref{ex:feynman_dirichlet_integral} appeals to. Fix $0 < t_0 < t_1$. For $t \in [t_0,t_1]$ the
$t$-derivative of the integrand is $e^{-tx} \leq e^{-t_0x}$, a bound that is independent of $t$
and integrable on $[0,\infty)$, so the differentiation may be carried out on $[0,R]$ and the
limit $R \to \infty$ taken afterwards. The compact statement is the one the course proves; the
dominated version is what the applications actually use.
\end{remark}
\end{exercisesolution}

\begin{exercisesolution}[Proof of \cref{ex:volume_ellipsoid_change_variables}]
% Generator: Gemini 3.6 Flash (High)
% Correction: Claude Opus 5 (High) --- \operatorname{vol} -> \vol, J\Phi -> \Jac\Phi, ^T ->
% \transp, and \iiint -> \int as elsewhere in this chapter.
\begin{enumerate}[label=\textbf{(\alph*)}]
  \item Define $\Phi(u,v,w) := (au, bv, cw)\transp$, and write $(x,y,z) = \Phi(u,v,w)$. Then
    $(x/a)^2 + (y/b)^2 + (z/c)^2 = u^2+v^2+w^2$, so $(u,v,w)$ lies in the closed unit ball if and
    only if $(x,y,z) \in U$. Since $a,b,c > 0$, the map $\Phi$ is a linear bijection, and therefore
    $\Phi(\overline{B_1(0)}) = U$.

  \item The Jacobian matrix of $\Phi$ is the diagonal matrix
\[\Jac\Phi = \begin{pmatrix} a & 0 & 0 \\ 0 & b & 0 \\ 0 & 0 & c \end{pmatrix},\]
    constant in $(u,v,w)$, with determinant $\det \Jac\Phi = abc > 0$.

  \item By the change of variables theorem (\cref{thm:change_of_variables}),
\[\vol_3(U) = \int_U 1 \, d\vol_3 = \int_{\overline{B_1(0)}} \lvert\det \Jac\Phi\rvert \, d\vol_3 = abc \int_{\overline{B_1(0)}} 1 \, d\vol_3 .\]
The unit ball in $\mathbb{R}^3$ has volume $\tfrac{4}{3}\pi$, and therefore
\[\vol_3(U) = \frac{4}{3} \pi abc.\]
\end{enumerate}
\end{exercisesolution}

\begin{exercisesolution}[Proof of \cref{ex:cubic_shear_area}]
% Generator: Claude Opus 5 (High)
\begin{enumerate}[label=\textbf{(\alph*)}]
  \item Differentiating each component,
\[D_x\varphi = \Jac\varphi(x) = \begin{pmatrix} 3x_1^2 & -1 \\ 1 & 1 \end{pmatrix},
  \qquad \det \Jac\varphi(x) = 3x_1^2 + 1 \ge 1 > 0 .\]

  \item The determinant never vanishes, so by the inverse function theorem
    (\cref{thm:inverse_function_theorem}) $\varphi$ is a local diffeomorphism at every point; in
    particular $\varphi(\mathbb{R}^2)$ is open. It remains to prove global injectivity, and then
    $\varphi$ is a bijection onto its image which is locally a diffeomorphism, hence a
    diffeomorphism onto its image.

    Suppose $\varphi(x) = \varphi(y)$. The second components give $x_2 + x_1 = y_2 + y_1$, that is,
    $x_2 - y_2 = -(x_1 - y_1)$. Substituting this into the equality of first components,
    $x_1^3 - x_2 = y_1^3 - y_2$, yields
\[x_1^3 - y_1^3 = x_2 - y_2 = -(x_1 - y_1),
  \qquad\text{that is,}\qquad (x_1 - y_1)\big(x_1^2 + x_1y_1 + y_1^2 + 1\big) = 0 .\]
    The second factor is at least $1$, because $x_1^2 + x_1y_1 + y_1^2 = \big(x_1 +
    \tfrac{y_1}{2}\big)^2 + \tfrac34 y_1^2 \ge 0$. Hence $x_1 = y_1$, and then $x_2 = y_2$.

  \item By part \textbf{(a)} the Jacobian determinant is positive, so the change of variables
    theorem (\cref{thm:change_of_variables}) gives
\[\vol_2(\varphi(Q)) = \int_Q \lvert\det \Jac\varphi\rvert \, dx
  = \int_0^1\!\!\int_0^1 \big(3x_1^2 + 1\big) \, dx_2\,dx_1
  = \int_0^1 \big(3x_1^2+1\big)\,dx_1 = 1 + 1 = 2 .\]
\end{enumerate}
\begin{remark}
The map is not linear, so the image of the square is not a parallelogram and there is no single
factor by which areas scale. The change of variables formula handles that by scaling
\emph{pointwise} and integrating the result, which is the whole reason the determinant appears
inside the integral rather than outside it.
\end{remark}
\end{exercisesolution}

\begin{exercisesolution}[Solution of \cref{ex:gaussian_second_moment_polar}]
% Generator: Claude Opus 5 (High)
Following the hint, the substitution $(x,y) \mapsto (y,x)$ is a linear change of variables with
Jacobian determinant $-1$, and it carries the integrand $x^2e^{-\alpha(x^2+y^2)}$ to
$y^2e^{-\alpha(x^2+y^2)}$. Hence the two integrals agree, and adding them gives
\[2I_\alpha = \int_{\mathbb{R}^2} (x^2+y^2)\,e^{-\alpha(x^2+y^2)} \, dx\,dy .\]
The integrand is now radial, which is what makes polar coordinates pay. Writing $x = r\cos\theta$,
$y = r\sin\theta$, with Jacobian determinant $r$,
\[2I_\alpha = \int_0^{2\pi}\!\!\int_0^\infty r^2 e^{-\alpha r^2}\, r\,dr\,d\theta
  = 2\pi \int_0^\infty r^3 e^{-\alpha r^2}\,dr .\]
Substituting $u := r^2$, so that $du = 2r\,dr$, turns the remaining integral into one that can be
done by parts:
\[\int_0^\infty r^3 e^{-\alpha r^2}\,dr = \frac12 \int_0^\infty u e^{-\alpha u}\,du
  = \frac{1}{2\alpha^2}.\]
Therefore $2I_\alpha = 2\pi \cdot \frac{1}{2\alpha^2} = \frac{\pi}{\alpha^2}$, and
\[I_\alpha = \frac{\pi}{2\alpha^2}.\]
\begin{remark}
Two sanity checks come free. Scaling $(x,y) \mapsto (x,y)/\sqrt\alpha$ predicts $I_\alpha =
\alpha^{-2} I_1$, which the answer satisfies; and $I_\alpha$ should equal
$\tfrac12 \cdot \tfrac{1}{2\alpha}\cdot\tfrac{\pi}{\alpha}$ read off from the one-dimensional
Gaussian moments, which it does. The symmetry trick in the first line is the same one that
evaluates $\int_{\mathbb{R}} e^{-x^2}dx$ by squaring it: an integral that resists in Cartesian
coordinates becomes radial once it is paired with its own mirror image.
\end{remark}
\end{exercisesolution}




% ============================================================
% Chapter 20 --- Gram determinants, area and arc length
% Originally: content/week-09.tex, \section{Length, area, and volume} with its
%             subsections (Corsin Week 9). "Volume of embedded surfaces" and "Length of
%             a curve" are promoted to sections here; problems 9.1 and 9.2 join them.
% ============================================================
\chapter{Gram Determinants, Area and Arc Length}
\label{ch:gram_determinant}

% Transition: Claude Opus 5 (Medium)
The previous chapters integrated over subsets of $\mathbb{R}^n$, where the volume element is
just $dx$. This one asks what to integrate over a $d$-dimensional surface sitting inside
$\mathbb{R}^n$, where $d < n$ and there is no such element to start from. The answer runs from
the area of a parallelogram, through the Gram determinant of $m$ vectors in $\mathbb{R}^n$ for
$m \neq n$, to the $d$-volume of a parametrized submanifold --- with the length of a curve as
the case $d = 1$.

% Originally: content/week-09.tex, \section{Length, area, and volume} (Corsin Week 9)

\section{Length, area, and volume}
\label{sec:length_area_volume}

\subsection{The area of a parallelogram in \texorpdfstring{$\mathbb{R}^2$}{R2}}
\label{sec:area_parallelogram}

The relationship between the determinant and volume is a key concept in this part of the
lecture. We illustrate the special case of $\mathbb{R}^2$.

\begin{proposition}[Area of a parallelogram]
\label{prop:area_parallelogram_r2}
Let $x,y \in \mathbb{R}^2$. The area of the parallelogram $(0,x,x+y,y)$ is
\[A = \big|\det(x\mid y)\big|.\]
\end{proposition}

\begin{proof}
Assume w.l.o.g.\ that $x = x_1e_1$, a reduction justified below. Then, writing
$y = (y_1,y_2)\transp$,
\[A = |x_1y_2| = \left|\det\begin{pmatrix} x_1 & y_1 \\ 0 & y_2 \end{pmatrix}\right|,\]
directly from the base-times-height formula for a parallelogram with one side on the
$e_1$-axis.

% Sonnet 5 (Medium) --- FIG-W09-01
\begin{center}
\begin{tikzpicture}[>=Latex, scale=1.0]
  \draw[->] (-0.3,0) -- (3.6,0) node[right, font=\scriptsize] {$e_1$};
  \draw[->] (0,-0.3) -- (0,2.1) node[above, font=\scriptsize] {$e_2$};
  \coordinate (O) at (0,0);
  \coordinate (X) at (2.4,0);
  \coordinate (Y) at (0.7,1.7);
  \coordinate (XY) at (3.1,1.7);
  \draw[purple, dotted, thick] (O) -- (X) -- (XY) -- (Y) -- cycle;
  \draw[blue, thick, ->] (O) -- (X) node[below, font=\scriptsize] {$x=x_1e_1$};
  \draw[orange, thick, ->] (O) -- (Y) node[above left, font=\scriptsize] {$y$};
  \draw[green!50!black, thick, dotted] (Y) -- (0.7,0);
  \draw[green!50!black, ->] (0,-0.15) -- (0.7,-0.15) node[midway, below, font=\scriptsize] {$y_1$};
  \draw[green!50!black, ->] (0.85,0) -- (0.85,1.7) node[midway, right, font=\scriptsize] {$y_2$};
\end{tikzpicture}
\end{center}

For the reduction, let $R \in \operatorname{O}(2)$ be any rotation carrying $x,y$ to a pair
$Rx,Ry$ with $Rx$ along $e_1$; rotations preserve area, so
\[A = \big|\det(Rx\mid Ry)\big| = \big|\det(R)\det(x\mid y)\big| = \big|\det(x\mid y)\big|,\]
using $|\det R| = 1$ and multiplicativity of $\det$. So, the \qt{w.l.o.g.} above was
justified: the formula $A=|\det(x\mid y)|$ holds for the rotated pair, and equals
$|\det(x\mid y)|$ for the original one.
\end{proof}

\subsection{The Gram determinant}
\label{sec:gram_determinant}

\begin{lemma}[Gram determinant of a square matrix]
\label{lem:gram_determinant_square}
For $M \in \mathrm{Mat}_{n\times n}(\mathbb{R})$,
\[|\det(M)| = \sqrt{\det(M)\det(M\transp)} = \sqrt{\det(M\transp)\det(M)}
  = \sqrt{\det(M\transp M)}.\]
\end{lemma}

\begin{proof}
The first equality is $|\det M| = \sqrt{(\det M)^2}$ together with $\det(M\transp)=\det(M)$;
the last is multiplicativity of the determinant, $\det(M\transp)\det(M)=\det(M\transp M)$.
\end{proof}

% Correction: Claude Opus 5 (Ultracode) --- read "the volume ... is given by the Gram
% determinant", which contradicts the ainote immediately below fixing this document's convention:
% the bare $\det(M\transp M)$ is the Gram determinant and its *square root* is the volume. The
% same conflation was repeated at three later sites in this chapter and is corrected there too.
The quantity $\det(M\transp M)$ is called the \newterm{Gram determinant}
\germanfor{Gram-Determinante} of $M$. So, by \cref{prop:area_parallelogram_r2}, the
volume of an $n$-parallelotope in $\mathbb{R}^n$ is given by the \emph{square root} of the Gram
determinant, in analogy to $\lvert\det M\rvert$ itself.

% Supplement: Analysis_II_Script_v1.pdf, p. 119
\begin{ainote}
Terminology diverges from the official script here, worth flagging since it is exactly the kind
of thing that costs a mark in an exam. The script's Definition 13.55 calls $\sqrt{\det(L\transp
L)}$, the square root and so the \emph{volume} itself, the \qt{Gram determinant}. This
document instead follows the more standard usage and calls the bare $\det(M\transp M)$ the Gram
determinant, with the square root then giving the volume (\cref{def:volume_gram_determinant}
below). Every downstream formula agrees regardless of which name is attached to which quantity,
so this is terminology only, not a mathematical disagreement.

The script's own statement of Definition 13.55 is, independently, internally inconsistent: it
introduces a linear map $L:\mathbb{R}^m\to\mathbb{R}^n$ with $m<n$, but then writes the formula
using an undefined $d$ (\qt{the $d\times d$ matrix $L\transp L$}) where the surrounding text
means $m$; and the prose describes \qt{the determinant of the \dots matrix $LL\transp$} while the
displayed formula correctly uses $L\transp L$. Only $L\transp L$ (size $m\times m$) is invertible
for $m<n$, since $LL\transp$ is $n\times n$ and singular. The displayed formula is therefore the
reliable half of the definition, and the surrounding prose is not.
\end{ainote}

\begin{ainote}
The point of restating $|\det M|$ this way is that the middle expression, $\det(M\transp
M)$, still makes sense once $M \in \mathrm{Mat}_{n\times m}(\mathbb{R})$ is rectangular:
$M\transp M \in \mathrm{Mat}_{m\times m}(\mathbb{R})$ is square even when $M$ is not, so its
determinant is always defined, whereas $\det(M)$ itself is not. This is what lets the
notion of \qt{volume spanned by vectors} survive dropping the requirement $m=n$. A set of
fewer than $n$ vectors in $\mathbb{R}^n$ still spans a lower-dimensional parallelotope, and
that parallelotope still has an $m$-dimensional volume, even though no square matrix and no
ordinary determinant is available to compute it with.
\end{ainote}

\begin{definition}[Volume of a parallelotope via the Gram determinant]
\label{def:volume_gram_determinant}
For $m$ vectors $v_1,\dots,v_m \in \mathbb{R}^n$, the \newterm{area}/\newterm{length}/
\newterm{volume} of the parallelotope spanned by them is
\[\sqrt{\det\big(\langle v_i,v_j\rangle\big)} = \sqrt{\det(V\transp V)},
  \qquad V := (v_1\mid v_2\mid\cdots\mid v_m) \in \mathrm{Mat}_{n\times m}(\mathbb{R}).\]
\end{definition}

\begin{example}[$m=1$: the length of a vector]
\label{ex:gram_m1_length}
For $m=1$, $\sqrt{\det(v_1\transp v_1)} = \sqrt{\langle v_1,v_1\rangle} = |v_1|$,
which is the length of $v_1$.
\end{example}

% Source: Corsin Nick/Class Notes/Week 9.pdf, p. 4
\begin{example}[$m=2,\ n=3$: area of a parallelogram in $\mathbb{R}^3$]
\label{ex:gram_m2_n3_area}
For $m=2$, $n=3$, we hope that the Gram determinant of
\[M = \begin{pmatrix} x_1 & y_1 \\ x_2 & y_2 \\ x_3 & y_3 \end{pmatrix}\]
gives the area of the parallelogram spanned by $x$ and $y$ in $\mathbb{R}^3$. If $\theta$ is
the angle between $x$ and $y$, this should be $A = |x|\,|y|\,|\sin\theta|$.

% Sonnet 5 (Medium) --- FIG-W09-02
% Correction: Gemini 3.7 Flash (High) --- dropped altitude was vertical rather than perpendicular
% to y (angle 81.6 deg); now drops perpendicularly to proj_y(x) = (2.79, 0.41) on the dashed line of y.
% Correction: Claude Opus 5 (Ultracode), found by rendering printed p. 314 --- that fix was right
% and incomplete. Three things follow it.
%   * The $|x|\sin\theta$ label had no plate, and the purple edge $XY \to Y$ runs straight through
%     it: that edge is $y = 0.25 + 0.6538(x-1.7)$, which crosses the label's own height $y = 1.05$
%     at $x = 2.92$, inside the label's span. Printed, the bar and the $x$ of $|x|$ came out
%     struck through. It takes the `lbl' plate from style.md. (The strikethrough predates the
%     altitude fix -- at the old position $(2.68,1.0)$ the same edge crossed at $x = 2.847$ --
%     but the label was moved without checking the page.)
%   * The $\theta$ arc annotates the very angle the altitude depends on and was left unaudited.
%     `(0.55,0.075) arc (5:33:0.6)' starts at radius $0.5551$, so TikZ centres the arc at
%     $(0.55,0.075) - 0.6(\cos 5^\circ, \sin 5^\circ) = (-0.048, 0.023)$, not at the origin; and
%     it spans $5^\circ$ to $33^\circ$ where the true directions are
%     $\arctan(0.25/1.7) = 8.37^\circ$ and $\arctan(1.7/2.6) = 33.17^\circ$. Restarted from
%     $0.6(\cos 8.37^\circ, \sin 8.37^\circ) = (0.5936, 0.0873)$, which puts the centre on $O$ and
%     the ends on $y$ and $x$. This is the idiom `02-metric-normed-inner-product/03-cauchy-schwarz'
%     already uses correctly (`arc (14:59:0.5)' against $14.04^\circ$ and $59.04^\circ$).
%   * A right-angle mark at the foot, so the perpendicularity the fix installed is visible rather
%     than merely true. Unit vector along $y$ is $(0.9893, 0.1455)$; the mark is drawn at side
%     $0.12$ from the foot back along $-y$ and up along the perpendicular $(-0.1455, 0.9893)$.
\begin{center}
\begin{tikzpicture}[>=Latex, scale=1.3,
  lbl/.style={fill=white, fill opacity=0.85, text opacity=1, inner sep=1.2pt,
              rounded corners=1pt}]
  \coordinate (O) at (0,0);
  \coordinate (X) at (2.6,1.7);
  \coordinate (Y) at (1.7,0.25);
  \coordinate (XY) at (4.3,1.95);
  \draw[orange, thick, ->] (O) -- (Y) node[below right, font=\scriptsize] {$y$};
  \draw[gray!55, dashed] (Y) -- (3.05,0.45);
  \draw[blue, thick, ->] (O) -- (X) node[above, font=\scriptsize] {$x$};
  \draw[purple, dotted, thick] (X) -- (XY) -- (Y);
  % Dashed, not dotted: the parallelogram's own edges above are thick dotted, and red against
  % purple is a distinction in hue alone, which style.md forbids relying on -- an e-reader renders
  % both as the same grey. The dash pattern, the right-angle mark and the label are the three
  % channels that now separate the altitude from the edge it runs beside.
  \draw[red, dashed, thick] (X) -- (2.79,0.41);
  \draw[red] (2.6713,0.3925) -- (2.6538,0.5112) -- (2.7725,0.5287);
  \node[lbl, text=red, font=\scriptsize, right] at (2.75,1.05) {$|x|\sin\theta$};
  \draw (0.5936,0.0873) arc (8.37:33.17:0.6);
  \node[font=\scriptsize] at (0.9,0.32) {$\theta$};
\end{tikzpicture}
\end{center}

And indeed:
\[
\begin{aligned}
\sqrt{\det(M\transp M)}
  &= \sqrt{\det\begin{pmatrix} \langle x,x\rangle & \langle x,y\rangle \\
    \langle x,y\rangle & \langle y,y\rangle\end{pmatrix}} \\
  &= \sqrt{\langle x,x\rangle\langle y,y\rangle - \langle x,y\rangle^2} \\
  &= \sqrt{|x|^2|y|^2 - |x|^2|y|^2\cos^2\theta} \\
  &= |x|\,|y|\,|\sin\theta| \qquad (=\ |x\times y|).
\end{aligned}
\]
\end{example}

% Originally: content/week-09.tex, \subsection{Volume of embedded surfaces} --- promoted to a section

\section{Volume of embedded surfaces}
\label{sec:volume_embedded_surfaces}

% Source: Corsin Nick/Class Notes/Week 9.pdf, p. 5
% Def. 13.68
\begin{definition}[$d$-volume on a parametrized $d$-submanifold]
\label{def:d_volume_parametrized_submanifold}
Let $V \subseteq \mathbb{R}^d$ be open, and let $\phi : V \to \mathbb{R}^n$ be a parametrized
submanifold. Given a Jordan measurable set $E \subseteq V$ with $\overline E \subseteq V$,
we define the $d$-volume of $\phi(E)$ by
\[\vol_d(\phi(E)) := \int_E \sqrt{\det\big(D\phi(x)\transp D\phi(x)\big)}\,dx.\]
\end{definition}

Here the Gram determinant in the integrand is the volume of the parallelotope spanned by the
vectors $\partial_1\phi(x),\partial_2\phi(x),\dots,\partial_d\phi(x)$, which by assumption
(see \qt{parametrized submanifold}) are linearly independent, i.e.\ $D\phi(x)$ has full rank,
so the volume is strictly positive.

% Sonnet 5 (Medium) --- FIG-W09-03
\begin{center}
\begin{tikzpicture}[>=Latex, scale=1.3]
  \draw[thick] plot[smooth] coordinates {(-2.1,-0.4) (-0.8,0.65) (0.5,-0.15) (2.1,0.75)};
  \draw[thick] plot[smooth] coordinates {(-1.8,0.9) (-0.5,-0.5) (0.8,0.4) (2.3,-0.65)};
  \coordinate (P) at (0.0,0.05);
  \fill (P) circle (1.6pt) node[below left, font=\scriptsize] {$\phi(x)$};
  \draw[blue, thick, ->] (P) -- ++(1.2,0.4) coordinate (D1)
    node[right, font=\scriptsize] {$\partial_1\phi(x)$};
  \draw[purple, thick, ->] (P) -- ++(0.35,1.05) coordinate (D2)
    node[above, font=\scriptsize] {$\partial_2\phi(x)$};
  \draw[green!55!black, dashed]
    (D1) -- ++(0.35,1.05) -- (D2) -- cycle;
  \node[green!45!black, font=\scriptsize, align=center] at (2.1,1.5)
    {$\sqrt{\det(D\phi(x)\transp D\phi(x))}$};
  \draw[green!55!black, thin, ->] (1.55,1.3) -- (0.95,1.0);
\end{tikzpicture}
\end{center}

This factor gives a measure for the stretching of $E$ as it is \qt{glued} to the surface. If
\[\sqrt{\det\big(D\phi(x)\transp D\phi(x)\big)} = 1 \quad \text{for all } x,\]
then $\phi$ is called an \newterm{isometry} (\germanterm{Isometrie}). A special case is the
length of a curve, treated in \cref{sec:length_of_curve} below.

% Quelle: Linus Lüchinger/Slides/Slides-04-20.pdf, slide 14
% Generator: Claude Opus 5 (High)
\begin{remark}[Why the transpose, and why the square root]
\label{rem:why_gram_determinant}
The shape of the formula looks arbitrary until one asks what could possibly stand in its place.
The honest question is: \emph{by what factor does $\phi$ stretch $d$-dimensional volume near
$x$?} Near $x$ the map $\phi$ is its differential, so the question is really about the linear
map $D\phi(x)$, and the answer runs in three steps.
\begin{enumerate}[label=\textbf{(\alph*)}]
  \item \textbf{The obvious candidate does not exist.} For a map $\mathbb{R}^n \to \mathbb{R}^n$
    the stretching factor is $\lvert\det\rvert$. Here $D\phi(x)$ is $n \times d$ with $d < n$, a
    \emph{rectangular} matrix, and a rectangular matrix has no determinant. So, the one-line
    answer is unavailable, and something has to replace it.
  \item \textbf{A factor nevertheless exists.} A linear map still scales all $d$-volumes in its
    domain by one and the same constant: it sends the unit cube of $\mathbb{R}^d$ to a fixed
    parallelotope, and every other region is built from scaled copies of that cube. Call this
    constant $c \geq 0$. It is what we want; we only need a formula for it.
  \item \textbf{The Gram matrix supplies it.} The product $D\phi(x)\transp D\phi(x)$ is
    $d \times d$, square at last, so it \emph{does} have a determinant, and that
    determinant turns out to be exactly $c^2$. Hence $c = \sqrt{\det(D\phi\transp D\phi)}$, and
    the square root is there for no deeper reason than to undo that squaring.
\end{enumerate}
The two familiar formulas are the small cases. For $d=1$ the matrix $D\phi(t)$ is the single
column $\phi'(t)$, so $D\phi\transp D\phi = \lvert\phi'(t)\rvert^2$ and $c = \lvert\phi'(t)\rvert$,
which is the arc-length integrand of \cref{sec:length_of_curve}. For $d=2$, $n=3$ one computes
$c = \lvert\partial_1\phi \times \partial_2\phi\rvert$, the area of the parallelogram the two
tangent vectors span, which is the surface-area formula used throughout this chapter.
\end{remark}

% Quelle: Linus Lüchinger/Slides/Slides-04-20.pdf, slide 12
\begin{ainote}[Two objects, three notations]
From this chapter onwards the lecture notes write $D\phi(x)$ for the Jacobian \emph{matrix} of
$\phi$ at $x$, the object earlier chapters of this document call $\Jac\phi(x)$. That collides
with the notation $D\phi_x$ used since \cref{ch:differential} for the \emph{differential}, the
linear map itself. The two are not the same kind of object, even though they carry the same
information:
\[\underbrace{D\phi_x}_{\text{linear map}}(v) \;=\; \underbrace{\Jac\phi(x)}_{\text{matrix}}\,v
  \;=\; \underbrace{D\phi(x)}_{\text{same matrix}}\,v.\]
The subscript is the whole distinction: $D\phi_x$ has $x$ downstairs, $D\phi(x)$ has it in
brackets. This document keeps each source's own convention rather than renotating transcribed
material, so both spellings of the matrix appear; where it matters, read $D\phi(x)$ and
$\Jac\phi(x)$ as synonyms.
\end{ainote}

% Supplement: Sascha Brack/Class Notes/Week_09_Notes_Friday.pdf, p. 3
\begin{example}[Integral on the graph of a function]
\label{ex:integral_graph_function}
Given an open set $U \subseteq \mathbb{R}^2$ and a $C^1$ function $f: U \to \mathbb{R}$, we can define the parametrized surface $\phi: U \to \mathbb{R}^3$ by $\phi(x,y) = (x,y,f(x,y))\transp$.
Assuming $M := \operatorname{Im}\phi$ is a $C^1$ submanifold, its area is given by $\int_M 1 \mathop{d\text{vol}}_2$.

We first find the Jacobian matrix of the parametrization:
\[\Jac\phi(x,y) = \begin{pmatrix} 1 & 0 \\ 0 & 1 \\ \partial_1 f(x,y) & \partial_2 f(x,y) \end{pmatrix}.\]
Then we compute the Gram determinant:
\begin{align*}
\sqrt{|\det(\Jac\phi\transp \Jac\phi)|}
&= \sqrt{\det \begin{pmatrix} 1+(\partial_1 f)^2 & \partial_1 f \partial_2 f \\ \partial_2 f \partial_1 f & 1+(\partial_2 f)^2 \end{pmatrix}} \\
&= \sqrt{\big(1+(\partial_1 f)^2\big)\big(1+(\partial_2 f)^2\big) - (\partial_1 f)^2(\partial_2 f)^2} \\
&= \sqrt{1 + (\partial_1 f)^2 + (\partial_2 f)^2} = \sqrt{1+|\nabla f(x,y)|^2}.
\end{align*}
By definition, we get the surface area formula for a graph:
\[\int_M 1 \mathop{d\text{vol}}_2 = \int_U \sqrt{1+|\nabla f(x,y)|^2} \,dx\,dy.\]
\end{example}

\subsection{Example: area of a paraboloid patch}
\label{sec:example_paraboloid_area}

% Source: Corsin Nick/Class Notes/Week 9.pdf, pp. 6--7
\begin{example}[Area of a paraboloid patch]
\label{ex:paraboloid_patch_area}
Compute the area of the surface
\[F = \left\{(x,y,z)\in\mathbb{R}^3 \ :\ x^2+4y^2 < 8,\ \ z = \tfrac12(x^2+2y^2)\right\}.\]

Since $z = \tfrac12(x^2+2y^2)$ means $F$ is a graph, we can conveniently write $F$ as the
image of the parametrization
\[\phi : A \to F, \qquad (x,y) \mapsto \left(x,\,y,\,\tfrac12(x^2+2y^2)\right),
  \qquad A := \{(x,y)\in\mathbb{R}^2 : x^2+4y^2 < 8\}.\]

\textbf{Gram determinant.} We compute
\[D\phi(x,y) = \begin{pmatrix} 1 & 0 \\ 0 & 1 \\ x & 2y \end{pmatrix}
  \quad \Longrightarrow \quad
  D\phi(x,y)\transp D\phi(x,y) = \begin{pmatrix} 1+x^2 & 2xy \\ 2xy & 1+4y^2 \end{pmatrix},\]
so that
\[\sqrt{\det D\phi(x,y)\transp D\phi(x,y)}
  = \sqrt{(1+x^2)(1+4y^2) - 4x^2y^2}
  = \sqrt{1+x^2+4y^2}.\]

\textbf{The integral.} By \cref{def:d_volume_parametrized_submanifold},
\[\vol_2(F) = \int_A \sqrt{1+x^2+4y^2}\,dx\,dy, \qquad A = \{(x,y)\in\mathbb{R}^2 :
x^2+4y^2<8\}.\]

We should choose more convenient coordinates. Polar coordinates come to mind, but then we
still have problems, since $A$ is not radially symmetric. We want to write
$A = \Psi\big((0,R)\times(0,2\pi)\big)$ (up to a set of measure zero). Note that for $y=0$,
$|x| < \sqrt8 = 2\sqrt2$, and for $x=0$, $|y| < \sqrt2$. Choosing $R := \sqrt2$ and weighing
$x,y$ accordingly,
\[\begin{pmatrix} x \\ y \end{pmatrix} = \begin{pmatrix} 2r\cos\theta \\ r\sin\theta
\end{pmatrix} =: \Psi(r,\theta).\]

\textbf{Change of variables.}
\[dx\,dy = |\det D\Psi(r,\theta)|\,dr\,d\theta
  = \left\lVert \begin{matrix} 2\cos\theta & -2r\sin\theta \\ \sin\theta & r\cos\theta
    \end{matrix}\right\rVert dr\,d\theta = 2r\,dr\,d\theta,\]
giving
\[
\begin{aligned}
\vol_2(F) &= \int_A \sqrt{1+x^2+4y^2}\,dx\,dy
  = \int_0^{\sqrt2} dr\int_0^{2\pi} d\theta\ 2r\sqrt{1+4r^2} \\
  &= 4\pi\int_0^{\sqrt2} dr\ r(1+4r^2)^{1/2}
   = \frac\pi2\int_0^{\sqrt2} dr\ 8r(1+4r^2)^{1/2} \\
  &= \frac\pi2\left[(1+4r^2)^{3/2}\cdot\frac23\right]_0^{\sqrt2}
   = \frac\pi3\big(27-1\big) = \frac{26\pi}{3}.
\end{aligned}
\]
\end{example}

\begin{ainote}
Worth flagging what makes $A$ awkward: it is an ellipse, not a disc, so naive polar
coordinates $(x,y)=(\rho\cos\theta,\rho\sin\theta)$ would force $\rho$ to depend on $\theta$
along the boundary. The fix is to build the eccentricity into the change of variables itself.
Scaling $x$ by a factor of $2$ relative to $y$ turns the ellipse into the disc $r<\sqrt2$
before polar coordinates are even applied. This is the same idea as
\Cref{ex:change_of_variables_domain}: choose $\Psi$ to straighten the domain, not
to simplify the integrand.
\end{ainote}

% Supplement: Sascha Brack/Class Notes/Week_09_Notes_Monday.pdf, p. 7
\begin{example}[Surface area of a sphere]
\label{ex:surface_area_sphere}
Given the sphere $M := \{(x,y,z) \in \mathbb{R}^3 \mid x^2+y^2+z^2 = r^2\}$, we can use the parametrization $\phi : (0,\pi) \times (-\pi,\pi) \to \mathbb{R}^3$,
\[\phi(\theta,\varphi) = \begin{pmatrix} r\sin\theta\cos\varphi \\ r\sin\theta\sin\varphi \\ r\cos\theta \end{pmatrix}.\]
The Jacobian matrix is
\[\Jac\phi(\theta,\varphi) = \begin{pmatrix} r\cos\theta\cos\varphi & -r\sin\theta\sin\varphi \\ r\cos\theta\sin\varphi & r\sin\theta\cos\varphi \\ -r\sin\theta & 0 \end{pmatrix},\]
which has rank 2 for all $(\theta,\varphi)$ in the domain.
Then we compute the Gram determinant:
\[\Jac\phi\transp \Jac\phi = \begin{pmatrix} r^2 & 0 \\ 0 & r^2\sin^2\theta \end{pmatrix} \quad \Longrightarrow \quad \sqrt{\det(\Jac\phi\transp \Jac\phi)} = r^2\sin\theta.\]
The surface area is given by the integral:
\begin{align*}
\int_M 1 \mathop{d\text{vol}}_2 &= \int_{-\pi}^\pi \int_0^\pi \sqrt{\det(\Jac\phi\transp \Jac\phi)} \,d\theta\,d\varphi \\
&= \int_{-\pi}^\pi \int_0^\pi r^2\sin\theta \,d\theta\,d\varphi \\
&= r^2 \left(\int_{-\pi}^\pi d\varphi\right) \left(\int_0^\pi \sin\theta \,d\theta\right) \\
&= r^2 \cdot 2\pi \cdot 2 = 4\pi r^2.
\end{align*}
Note that here, our partition of unity simply was the identity (since the parametrization covers the sphere up to a null set).

\begin{center}
\begin{tikzpicture}[scale=1.5]
  % Submanifold shape
  \draw[thick] (0,0) .. controls (0.5,0.2) and (1.5,-0.2) .. (2,0)
               .. controls (2.2,-0.5) and (1.8,-0.8) .. (2,-1)
               .. controls (1.5,-0.8) and (0.5,-1.2) .. (0,-1)
               .. controls (-0.2,-0.5) and (0.2,-0.5) .. (0,0);
               
  % Two charts
  \draw[very thick, green!60!black] (0.5,-0.5) circle (0.8);
  \node[green!60!black] at (-0.3,-1.2) {$B_{r_1}(p_1)$};
  
  \draw[very thick, green!60!black] (1.5,-0.5) circle (0.8);
  \node[green!60!black] at (2.3,-1.2) {$B_{r_2}(p_2)$};
  
  \node[right, font=\footnotesize, align=left, text=green!60!black] at (2.5, -0.3) {
    and $\eta_\ell(x)$ is defined such\\
    that when we integrate\\
    in each ball separately\\
    and add all the\\
    integrals together,\\
    we exactly get a notion\\
    of the integral on the\\
    whole submanifold.
  };
\end{tikzpicture}
\end{center}
\end{example}


% --- exercises relocated here from the dissolved sheet 9 ---

% Originally: content/week-09.tex, \exercisesheet{9}, problem 9.1
% Source: exercises/Ex9_Analysis2_eng.pdf

% Source: exercises/Ex9_Analysis2_eng.pdf, p. 1
\begin{exercise}[9.1 --- Volume \textnormal{(important)}]
\label{ex:9.1}
Calculate the volume of the region $B \subset \mathbb{R}^3$ enclosed by the surfaces
$x^2+y^2+z^2=8$ and $2z=x^2+y^2$. \textit{Official hint:} use cylindrical coordinates.

% Source: Corsin Nick/Class Notes/Week 9.pdf, p. 1
\textit{Corsin's hint:}
\[\begin{pmatrix} x \\ y \\ z\end{pmatrix} = \begin{pmatrix} r\sin\theta\cos\varphi \\
  r\sin\theta\sin\varphi \\ r\cos\theta \end{pmatrix}
  \qquad \text{for } (r,\theta,\varphi) \in (0,\infty)\times(0,\pi)\times(0,2\pi).\]
\end{exercise}

% The solution in 99-solutions.tex points back at this note, so it carries a label. Possible
% only since 2026-08-09, when ainote gained its own counter; the solution used to say
% "see the \texttt{ainote} above", which leaked a LaTeX environment name into the typeset text.
\begin{ainote}[Corsin's hint is spherical, not cylindrical]
\label{note:9.1_coordinates_mislabelled}
Corsin labels the formula above \qt{cylindrical coordinates}, but $(r\sin\theta\cos\varphi$, $r\sin\theta\sin\varphi$, $r\cos\theta)$ is the \textbf{spherical} parametrization, and his
range $\theta \in (-\pi/2,\pi/2)$ has been corrected to $(0,\pi)$, the range on which
$\theta \mapsto \cos\theta$ actually covers $[-1,1]$ bijectively (already logged in \cref{ex:7.6}). The mislabelling is doubly worth noting here, since the official sheet's own hint
recommends \emph{cylindrical} coordinates $(r,\theta,z)$, and that is in fact the more
natural choice for this problem, as the solution on \cpageref{ex:solution_9.1} does: both surfaces are
surfaces of revolution about the $z$-axis, which cylindrical coordinates respect directly,
whereas spherical coordinates mix $r$ and $\theta$ into both defining equations.
\end{ainote}

% Originally: content/week-10.tex, \exercisesheet{10}, problem 10.1
% Source: exercises/Ex10_Analysis2_eng.pdf

% Source: exercises/Ex10_Analysis2_eng.pdf, p. 1
\begin{exercise}[10.1 --- Spherical quadrilateral \textnormal{(important)}]
\label{ex:10.1}
Calculate the area of the spherical quadrilateral $S \subset \mathbb{S}^2$ bounded by the
meridians $-\pi < \varphi_1 < \varphi_2 < \pi$ and parallels $-\pi/2 < \theta_1 < \theta_2 <
\pi/2$.
\end{exercise}

% Originally: content/week-10.tex, \exercisesheet{10}, problem 10.2
% Source: exercises/Ex10_Analysis2_eng.pdf

% Source: exercises/Ex10_Analysis2_eng.pdf, p. 1
\begin{exercise}[10.2 --- The surface area of a torus \textnormal{(important)}]
\label{ex:10.2}
Let $S := \{(x,y,0)\in\mathbb{R}^3 \mid x^2+y^2=4\}$ and $T := \{x\in\mathbb{R}^3 \mid d_S(x)
\leq 1\}$, where $d_S(x) = \inf\{\lvert x-y\rvert \mid y\in S\}$ denotes the minimal distance
from $S$ to $x\in\mathbb{R}^3$. Parametrize $\partial T$ and then calculate the area.

% Source: Corsin Nick/Class Notes/Week 10.pdf, p. 1
\textit{Corsin's hint, general parametrization of a torus surface:}
\[x = (R+r\cos\phi)\cos\theta, \qquad y = (R+r\cos\phi)\sin\theta, \qquad z = r\sin\phi,\]
for the torus swept by a tube of radius $r$ at distance $R$ from the axis.

% Sonnet 5 (Medium) --- FIG-W10-01
\begin{center}
\begin{tikzpicture}[>=Latex, scale=1.0]
  \begin{scope}
    \draw[->] (0,0) -- (2.2,0) node[right, font=\scriptsize] {$x$};
    \draw[->] (0,0) -- (0,1.8) node[above, font=\scriptsize] {$z$};
    \draw[->] (0,0) -- (-1.1,-0.7) node[left, font=\scriptsize] {$y$};
    \draw[blue!70!black, thick] (0.4,0.55) circle (0.9 and 0.35);
    \draw[blue!70!black, thick] (0.05,0.15) circle (1.5 and 0.6);
    \draw[blue!70!black, dashed] (-1.0,0.0) arc (180:360:1.5 and 0.6);
  \end{scope}
  \begin{scope}[xshift=5.2cm]
    \draw[blue!70!black, thick] (0,0) circle (1.6);
    \draw[blue!70!black, thick] (0,0) circle (0.7);
    \draw[->] (0,0) -- (55:1.15) node[midway, above, font=\scriptsize] {$R$};
    \draw[orange, ->] (0,0) -- (18:0.7) node[midway, below, font=\scriptsize] {};
    \draw[orange, thick, ->] (18:0.7) -- (55:1.6)
      node[midway, above right, font=\scriptsize] {$r$};
    \draw (0.55,0) arc (0:18:0.55);
    \node[font=\scriptsize] at (0.75,0.15) {$\theta$};
    \draw[->] (1.1,0) -- (2.0,0) node[right, font=\scriptsize] {$x$};
    \draw[->] (1.1,0) -- (1.1,0.9) node[above, font=\scriptsize] {$y$};
  \end{scope}
  \begin{scope}[xshift=10.4cm]
    \draw[purple, dotted] (-1.6,0) -- (1.6,0);
    \draw (0,0) circle (0.9);
    \draw[->] (0,0) -- (35:0.9) node[midway, above right, font=\scriptsize] {$r$};
    \draw (0.35,0) arc (0:35:0.35);
    \node[font=\scriptsize] at (0.55,0.15) {$\varphi$};
    \draw[->] (-1.6,-0.3) -- (-1.6,0.9) node[above, font=\scriptsize] {$z$};
    \draw[->] (-1.6,-0.3) -- (-0.6,-0.3) node[right, font=\scriptsize] {$x$};
  \end{scope}
\end{tikzpicture}
\end{center}
\end{exercise}

% Originally: content/week-09.tex, \subsection{Length of a curve} --- promoted to a section

\section{Length of a curve}
\label{sec:length_of_curve}

% Source: Corsin Nick/Class Notes/Week 9.pdf, p. 8
\begin{definition}[Length of a $C^1$-curve]
\label{def:length_c1_curve}
The \newterm{length} \germanfor{Länge} of a curve or path $\gamma \in
C^1([a,b],\mathbb{R}^n)$ is given by
\[L(\gamma) := \int_a^b \lvert\gamma'(t)\rvert\,dt.\]
\end{definition}

\begin{definition}[Length of a $C^0$-curve]
\label{def:length_c0_curve}
For a $C^0$-curve $c \in C^0([a,b],\mathbb{R}^n)$, the length is defined as
\[L(c) := \sup\left\{\sum_{i=1}^{n-1} \lvert c(t_{i+1})-c(t_i)\rvert
  \ :\ a = t_1 < \cdots < t_n = b\right\},\]
the supremum over all lengths of inscribed polygons.
\end{definition}

% Sonnet 5 (Medium) --- FIG-W09-04
\begin{center}
\begin{tikzpicture}[>=Latex, scale=1.0]
  \draw[thick] plot[smooth] coordinates
    {(0,0) (0.6,0.9) (1.3,0.5) (1.9,0.15) (2.5,0.75) (3.1,0.5) (3.7,1.3)};
  \node[font=\scriptsize] at (1.9,1.15) {$\operatorname{Im}(c)$};
  \draw[blue, thick]
    (0,0) -- (0.55,0.95) -- (1.35,0.55) -- (2.0,0.1) -- (2.55,0.7) -- (3.15,0.45) -- (3.7,1.25);
  \foreach \x/\y/\lbl/\pos in {0/0/{c(t_1)}/below,0.55/0.95/{c(t_2)}/above,1.35/0.55/{c(t_3)}/above,2.0/0.1/{c(t_4)}/below,2.55/0.7/{c(t_5)}/{above left},3.15/0.45/{c(t_6)}/below,3.7/1.25/{c(t_7)}/right} {
    \fill (\x,\y) circle (1.3pt);
    \node[font=\scriptsize, \pos] at (\x,\y) {$\lbl$};
  }
\end{tikzpicture}
\end{center}

\begin{remark}
For a $C^1$-curve, \cref{def:length_c1_curve,def:length_c0_curve} agree. Notice that
\[\lvert\gamma'(t)\rvert = \sqrt{D\gamma(t)\transp \cdot D\gamma(t)}\]
is exactly the Gram determinant of $D\gamma(t)$ in the sense of
\Cref{def:volume_gram_determinant}, the $m=1$ case of \cref{ex:gram_m1_length}, now applied
pointwise along the curve. This is why the length of a curve is the special case
$d=1$ of \cref{def:d_volume_parametrized_submanifold}: a curve is a $1$-dimensional
parametrized submanifold, and $L(\gamma) = \vol_1(\gamma([a,b]))$.
\end{remark}

% Supplement: Sascha Brack/Class Notes/Week_09_Notes_Friday.pdf, p. 4
\begin{example}[Line integral]
\label{ex:line_integral}
Let $\phi : [-\pi/2, \pi/2] \to \mathbb{R}^2$ be given by $\phi(\theta) = \begin{pmatrix} 3\sin\theta \\ 4\sin\theta \end{pmatrix}$.
Compute the length of the curve $M := \operatorname{Im}\phi$, i.e.\ $\int_M 1 \mathop{d\text{vol}}_1$.

First, we compute the Jacobian (which here is just the derivative vector):
\[\Jac\phi(\theta) = \begin{pmatrix} 3\cos\theta \\ 4\cos\theta \end{pmatrix}.\]
Then we compute the Gram determinant (which is the length of the tangent vector):
\[\sqrt{\det(\Jac\phi(\theta)\transp \Jac\phi(\theta))} = \sqrt{9\cos^2\theta + 16\cos^2\theta} = \sqrt{25\cos^2\theta} = 5|\cos\theta|.\]
Since $\cos\theta \geq 0$ for $\theta \in [-\pi/2, \pi/2]$, this simplifies to $5\cos\theta$.
The length is:
\[\int_{-\pi/2}^{\pi/2} 5\cos\theta \,d\theta = 5\big[\sin\theta\big]_{-\pi/2}^{\pi/2} = 5(1 - (-1)) = 10.\]
\end{example}

% Source: Linus Lüchinger/Slides/Slides-04-30.pdf, p. 4
\begin{proposition}[Length of a graph]
\label{prop:graph_arc_length}
Let $f \in C^1([a,b],\mathbb{R})$. The graph
$\Gamma_f := \{(x,f(x)) : x \in [a,b]\} \subseteq \mathbb{R}^2$ has length
\[\vol_1(\Gamma_f) = \int_a^b \sqrt{1+f'(x)^2}\,dx.\]
\end{proposition}

\begin{proof}
Parametrize by $\phi(t) := (t, f(t))\transp$ on $[a,b]$, which is injective and $C^1$ with
$\phi'(t) = (1, f'(t))\transp \neq 0$. The Gram determinant of the single tangent vector is
\[\sqrt{\det\big(\Jac\phi(t)\transp\Jac\phi(t)\big)} = |\phi'(t)|
  = \sqrt{1+f'(t)^2},\]
and integrating over $[a,b]$ is the definition of $\vol_1$.
\end{proof}

% Source: Linus Lüchinger/Slides/Slides-04-30.pdf, p. 4
\begin{example}[A quarter of an astroid]
\label{ex:astroid_arc_length}
Let $\gamma(t) := (\cos^3t,\ \sin^3t)\transp$ for $t \in [0,\tfrac\pi2]$. Then
\[\gamma'(t) = \begin{pmatrix} -3\cos^2t\,\sin t \\ 3\sin^2t\,\cos t\end{pmatrix},\]
so, using $\cos t \geq 0$ and $\sin t \geq 0$ on this interval,
\[|\gamma'(t)| = 3\cos t\,\sin t\sqrt{\cos^2t+\sin^2t} = 3\cos t\,\sin t
  = \tfrac32\sin(2t).\]
The length is therefore
\[\int_0^{\pi/2}\tfrac32\sin(2t)\,dt = \Big[-\tfrac34\cos(2t)\Big]_0^{\pi/2}
  = -\tfrac34\big((-1)-1\big) = \tfrac32.\]
Note that $\gamma'(0) = \gamma'(\tfrac\pi2) = 0$: the astroid has cusps at its endpoints, so this
is a curve whose parametrization fails the immersion condition exactly at the two boundary
points. The length integral is unaffected, since two points contribute nothing to it.
\end{example}

% Supplement: Sascha Brack/Class Notes/Week_09_Notes_Friday.pdf, p. 6
\begin{example}[Spiral length]
\label{ex:spiral_length}
Consider the spiral parametrized by $\phi : (0,2\pi) \to \mathbb{R}^3$,
\[\phi(t) = \begin{pmatrix} \cos(\omega t) \\ \sin(\omega t) \\ vt \end{pmatrix},\]
where $v \in \mathbb{R}\setminus\{0\}$ and $\omega \in (0,\infty)$. Compute the length of the spiral $\operatorname{Im}\phi$.

We compute the Jacobian:
\[\Jac\phi(t) = \begin{pmatrix} -\omega\sin(\omega t) \\ \omega\cos(\omega t) \\ v \end{pmatrix}.\]
The Gram determinant is:
\[\sqrt{\Jac\phi(t)\transp \Jac\phi(t)} = \sqrt{\omega^2\sin^2(\omega t) + \omega^2\cos^2(\omega t) + v^2} = \sqrt{\omega^2 + v^2}.\]
The length of the spiral is the integral over the parameter domain:
\[\int_0^{2\pi} \sqrt{\omega^2 + v^2} \,dt = 2\pi\sqrt{\omega^2 + v^2}.\]
\end{example}

% Supplement: Sascha Brack/Class Notes/Week_09_Notes_Monday.pdf, p. 8
\begin{example}[Line integral of a function on a circle]
\label{ex:line_integral_function_circle}
Given a function $f: \mathbb{R}^3 \to \mathbb{R}$, $f(x,y,z) = x^2+z^2$. We want to integrate this function on the circle $M := \{(x,y,z) \in \mathbb{R}^3 \mid x^2+y^2=1,\ z=2\}$.
We parametrize the circle by $\phi: (0,2\pi) \to \mathbb{R}^3$,
\[\phi(\varphi) = \begin{pmatrix} \cos\varphi \\ \sin\varphi \\ 2 \end{pmatrix}.\]
The Jacobian is $\Jac\phi(\varphi) = \phi'(\varphi) = \begin{pmatrix} -\sin\varphi \\ \cos\varphi \\ 0 \end{pmatrix}$.
The Gram determinant is $\sqrt{\Jac\phi\transp \Jac\phi} = |\phi'(\varphi)| = \sqrt{\sin^2\varphi + \cos^2\varphi} = 1$.
The function evaluated on the curve is $f(\phi(\varphi)) = \cos^2\varphi + 4$.
The line integral is thus:
\begin{align*}
\int_M f \mathop{d\text{vol}}_1 &= \int_0^{2\pi} f(\phi(\varphi)) |\phi'(\varphi)| \,d\varphi \\
&= \int_0^{2\pi} (\cos^2\varphi + 4) \cdot 1 \,d\varphi \\
&= \int_0^{2\pi} \cos^2\varphi \,d\varphi + \int_0^{2\pi} 4 \,d\varphi \\
&= \pi + 8\pi = 9\pi.
\end{align*}
\end{example}


% --- exercises relocated here from the dissolved sheet 9 ---

% Originally: content/week-09.tex, \exercisesheet{9}, problem 9.2
% Source: exercises/Ex9_Analysis2_eng.pdf

% Source: exercises/Ex9_Analysis2_eng.pdf, p. 1
\begin{exercise}[9.2 --- Multiple integrals]
\label{ex:9.2}
\begin{enumerate}[label=\textbf{\arabic*.}]
  \item Let $D := [0,2]\times[0,1]$. Calculate $\displaystyle\iint_D (x^3+3x^2y+y^3)\,dx\,dy$.
  \item Let $D \subset \mathbb{R}^2$ be the interior of the triangle with vertices $(0,0)$,
    $(0,\pi)$, and $(\pi,\pi)$. Calculate $\displaystyle\iint_D x\cos(x+y)\,dx\,dy$.
  \item Let $D := \{(x,y)\in\mathbb{R}^2 \mid x>1,\ y>1,\ x+y<3\}$. Calculate
    $\displaystyle\iint_D \frac{1}{(x+y)^3}\,dx\,dy$.
\end{enumerate}

% Source: Corsin Nick/Class Notes/Week 9.pdf, p. 1
% The environment form rather than \exhint, because this hint carries a display.
\begin{exercisehint}[Corsin's hint (part 2)]
You can parametrize the triangle as
\[D = \{(x,y)\in\mathbb{R}^2 : 0 \leq x \leq y \leq \pi\}.\]
\end{exercisehint}
\end{exercise}

% Originally: content/week-10.tex, \exercisesheet{10}, problem 10.3
% Source: exercises/Ex10_Analysis2_eng.pdf

% Source: exercises/Ex10_Analysis2_eng.pdf, p. 1
\begin{exercise}[10.3 --- Solids of revolution \textnormal{(important)}]
\label{ex:10.3}
Given $f \in C^1((a,b))\cap C([a,b])$ such that $f\geq0$, define the \textnormal{rotational
body around the $x$-axis} as
\[\begin{aligned}
  \Omega &:= \left\{(x,y,z)\in\mathbb{R}^3 \mid a<x<b,\ \sqrt{y^2+z^2}<f(x)\right\}, \\
  \partial_{\mathrm{side}}\Omega &:= \left\{(x,y,z)\in\mathbb{R}^3 \mid a<x<b,\ \sqrt{y^2+z^2}=f(x)\right\}.
\end{aligned}\]
\begin{enumerate}[label=\textbf{\arabic*.}]
  \item Say whether $\Omega$, $\partial_{\mathrm{side}}\Omega$ are open/closed/connected
    subsets of $\mathbb{R}^3$.
  \item Show that $\vol_3(\Omega) = \pi\int_a^b f(x)^2\,dx$.
  \item Show that $\vol_2(\partial_{\mathrm{side}}\Omega) = 2\pi\int_a^b
    \sqrt{1+f'(x)^2}\,f(x)\,dx$. \textit{Hint:} parametrize $(x,\theta)\mapsto (x,f(x)\cos\theta,
    f(x)\sin\theta)$.
  \item Calculate the volume and surface area of the `improper rotational body' that arises
    when $f(x)=1/x$, $a=1$, $b=\infty$.
  \item \textnormal{(*)} Show that $\Omega$ is a bounded $C^1$ domain (as in
    \Cref{def:bounded_ck_domain}) if and only if
    \begin{equation}
    f(x)>0 \text{ in } (a,b), \qquad f'(a^+)=+\infty, \qquad f'(b^-)=-\infty. \tag{$*$}
    \end{equation}
\end{enumerate}
\end{exercise}

% Originally: content/week-10.tex (Corsin Week 10 / Jérôme Paschoud Week 10)

\section{Integration over submanifolds}
\label{sec:integration_over_submanifolds}

% Source: Jérôme Paschoud/Notizen/Woche 10.1 Integrale über UMF.pdf, p. 3
% Generator: Gemini 3.6 Flash (High)
While \cref{sec:volume_embedded_surfaces} explains how to compute the $d$-volume of a parametrized submanifold, it requires the entire relevant portion of the surface to be covered by a single parametrization $\phi$. For more complicated submanifolds $M \subseteq \mathbb{R}^n$, we need to integrate scalar functions $f : M \to \mathbb{R}$ globally. 

To achieve this, we cover the $m$-dimensional submanifold $M$ with the images of local parametrizations $\Phi_l : V_l \to \mathbb{R}^n$, called an \newterm{atlas}. By employing a smooth \newterm{partition of unity} $\eta_l \in C_c^\infty(\mathbb{R}^n, [0,1])$ (see \cref{note:partitions_of_unity_motivation} and \cref{prop:partition_of_unity_finite} for a rigorous construction) subordinate to this cover, we ensure that:
\begin{enumerate}[label=\textbf{(\alph*)}]
    \item $\sum_{l=1}^N \eta_l \equiv 1$ on $M$,
    \item $\operatorname{supp}(\eta_l) \subset \operatorname{Im}(\Phi_l)$.
\end{enumerate}

\begin{definition}[Integral of a scalar function over a submanifold]
\label{def:integral_scalar_submanifold}
Given an $m$-dimensional submanifold $M \subseteq \mathbb{R}^n$, an atlas of local parametrizations $\Phi_l : V_l \to \mathbb{R}^n$, and a subordinate partition of unity $\{\eta_l\}$, the integral of a function $f : M \to \mathbb{R}$ over $M$ is defined as:
\[\int_M f \mathop{d\text{vol}}_m := \sum_{l=1}^N \int_{V_l} f(\Phi_l(x)) \eta_l(\Phi_l(x)) \sqrt{\det(\Jac\Phi_l\transp(x) \Jac\Phi_l(x))} \,dx.\]
\end{definition}

This definition pieces together the local integrals (weighted by $\eta_l$) to compute the global integral over the entire submanifold, generalizing the volume integral from \cref{def:d_volume_parametrized_submanifold}.

% Generator: Claude Opus 5 (Ultracode) --- the section defined the integral and stopped there,
% leaving the one question the definition raises unanswered: it is written in terms of an atlas
% and a partition of unity, neither of which is canonical, so it says nothing until those choices
% are shown not to matter.
The definition names two things that are not canonical --- an atlas, and a partition of unity
subordinate to it --- so it defines nothing until those choices are shown not to matter.

\begin{proposition}[The integral depends on neither the atlas nor the partition]
\label{prop:integral_over_submanifold_well_defined}
The right-hand side of \cref{def:integral_scalar_submanifold} is unchanged when the atlas
$\{\Phi_l\}_{l=1}^N$ is replaced by any other finite atlas of $M$ and $\{\eta_l\}$ by any
partition of unity subordinate to it.
\end{proposition}

\begin{proof}
Two steps, and the first carries the whole argument.

\textbf{One patch, two parametrizations.} Let $\Phi : V \to \mathbb{R}^n$ and
$\Psi : W \to \mathbb{R}^n$ be local parametrizations with the same image, and let $g$ be
continuous with $\operatorname{supp}(g)$ contained in that image. The transition map
$\tau := \Psi^{-1}\circ\Phi : V \to W$ is a $C^1$-diffeomorphism and $\Phi = \Psi\circ\tau$, so
the chain rule gives $\Jac\Phi(x) = \Jac\Psi(\tau(x))\,\Jac\tau(x)$ and therefore
\[\Jac\Phi\transp\Jac\Phi
  = \Jac\tau\transp\,\big(\Jac\Psi\transp\Jac\Psi\big)\,\Jac\tau .\]
Both outer factors are $m \times m$, so taking determinants and using multiplicativity twice,
\[\sqrt{\det\big(\Jac\Phi\transp\Jac\Phi\big)(x)}
  = \big\lvert\det\Jac\tau(x)\big\rvert\,
    \sqrt{\det\big(\Jac\Psi\transp\Jac\Psi\big)(\tau(x))} .\]
The factor $\lvert\det\Jac\tau\rvert$ is precisely what \cref{thm:change_of_variables} consumes,
so substituting $y = \tau(x)$ gives
\[\int_V g(\Phi(x))\sqrt{\det\big(\Jac\Phi\transp\Jac\Phi\big)}\,dx
  = \int_W g(\Psi(y))\sqrt{\det\big(\Jac\Psi\transp\Jac\Psi\big)}\,dy .\]
The Gram determinant was built to make exactly this cancellation happen, and this is where that
pays.

\textbf{Two atlases.} Let $\{(\Phi_l,\eta_l)\}_{l=1}^N$ and $\{(\Psi_k,\theta_k)\}_{k=1}^K$ be two
systems as in \cref{def:integral_scalar_submanifold}, producing values $I$ and $J$. Since
$\sum_{k}\theta_k \equiv 1$ on $M$, inserting that sum into $I$ and exchanging the two finite sums,
\[I = \sum_{l=1}^N\sum_{k=1}^K \int_{V_l}
      \big(\eta_l\,\theta_k\,f\big)(\Phi_l(x))
      \sqrt{\det\big(\Jac\Phi_l\transp\Jac\Phi_l\big)}\,dx .\]
The integrand of the $(l,k)$ term is supported in
$\operatorname{Im}(\Phi_l)\cap\operatorname{Im}(\Psi_k)$; restricting $\Phi_l$ and $\Psi_k$ to the
preimages of that intersection makes them two parametrizations of one patch, so the first step
applies and replaces $\Phi_l$ by $\Psi_k$ in that term. Doing so for every $(l,k)$ and then summing
over $l$ first, using $\sum_l\eta_l \equiv 1$, collapses the double sum to $J$. Hence $I = J$.
\end{proof}

\begin{remark}[When a single chart already does, this costs nothing]
\label{rem:one_chart_recovers_the_earlier_definition}
If $M$ is covered by a single parametrization $\Phi_1$, the only partition subordinate to that
atlas is $\eta_1 \equiv 1$, and \cref{def:integral_scalar_submanifold} collapses to
\[\int_M f \mathop{d\text{vol}}_m
  = \int_{V_1} f(\Phi_1(x))\sqrt{\det\big(\Jac\Phi_1\transp\Jac\Phi_1\big)}\,dx ,\]
which is \cref{def:d_volume_parametrized_submanifold} with $f$ carried along; at $f \equiv 1$ it
returns $\vol_m(M)$ exactly. So the general definition extends the earlier one rather than
competing with it.

In practice one almost never builds a partition of unity, because a single chart usually covers
all of $M$ but a null set --- spherical coordinates miss only a half-meridian of $S^2$, which
carries no $\vol_2$ --- and a null set contributes nothing to the integral. That is why every
surface integral in \cref{sec:volume_embedded_surfaces} could be evaluated from one
parametrization without comment. The partition of unity is what makes the definition \emph{exist}
for an arbitrary $M$; it is not a device for computing with.
\end{remark}

% Generator: Claude Opus 5 (Ultracode)
\begin{aiexercise}[Two atlases on the circle, and one that is not an atlas]
\label{ex:ai_circle_two_atlases}
Let $M := S^1 \subseteq \mathbb{R}^2$ and take $f \equiv 1$, so that the integral is the length.
\begin{enumerate}[label=\textbf{(\alph*)}]
  \item Evaluate $\int_{S^1} 1 \mathop{d\text{vol}}_1$ from the single parametrization
    $\Phi(t) := (\cos t,\sin t)$ on $(0,2\pi)$, whose image omits one point of $S^1$.
  \item Now take the two-arc atlas $\Phi_1 := \Phi$ on $(0,2\pi)$ and $\Phi_2(t) :=
    (\cos t,\sin t)$ on $(-\pi,\pi)$. Show that \emph{every} partition of unity subordinate to it
    returns the same value, without choosing one.
  \item The two graph parametrizations $\Psi_\pm(u) := \big(u,\pm\sqrt{1-u^2}\big)$ on $(-1,1)$
    satisfy $\sqrt{\det(\Jac\Psi_\pm\transp\Jac\Psi_\pm)} = (1-u^2)^{-1/2}$, and each semicircle
    has length $\pi$. Why does $\{\Psi_+,\Psi_-\}$ nonetheless fail to be an atlas of $S^1$, and
    what goes wrong if one uses it in \cref{def:integral_scalar_submanifold} anyway?
\end{enumerate}
\exsol{sol:ai_circle_two_atlases}
\end{aiexercise}

% Originally: content/week-09.tex, \section{Solutions}

\section{Solutions}
\label{sec:gram_determinant_solutions}

\begin{exercisesolution}[Solution to \cref{ex:area_formula_two_maps}]
\label{sol:area_formula_two_maps}
% Generator: Claude Opus 5 (High) --- the five statements are the examiner's; the working, the
% counterexample and the AI-Note below are written for these notes. The verdicts (a) true,
% (b)-(d) false, (e) true are the official key's; ours agree on (b)-(e) and diverge on (a), for
% the reason set out in the AI-Note.
Take the second map first, because there \eqref{eq:mock_area_formula} applies without any
difficulty. Write $\Psi = L + c$ with
\[L := \begin{pmatrix} 1 & 1 \\ \alpha & -1 \\ 0 & 1 \end{pmatrix}, \qquad
  c := \begin{pmatrix} 0 \\ 0 \\ 1 \end{pmatrix},\]
so $\Jac\Psi(x,y) = L$ at every point. The Gram matrix of its two columns
$c_1 = (1,\alpha,0)\transp$ and $c_2 = (1,-1,1)\transp$ is
\[L\transp L = \begin{pmatrix} \langle c_1,c_1\rangle & \langle c_1,c_2\rangle \\
  \langle c_2,c_1\rangle & \langle c_2,c_2\rangle \end{pmatrix}
  = \begin{pmatrix} 1+\alpha^2 & 1-\alpha \\ 1-\alpha & 3 \end{pmatrix},\]
whose determinant is
\[3(1+\alpha^2) - (1-\alpha)^2 = 3 + 3\alpha^2 - 1 + 2\alpha - \alpha^2 = 2(\alpha^2+\alpha+1).\]
The integrand of \eqref{eq:mock_area_formula} is therefore the constant
$\sqrt{2(\alpha^2+\alpha+1)}$, and integrating a constant over $E$ gives
\[\vol_2(\Psi(E)) = \sqrt{2(\alpha^2+\alpha+1)}\ \vol_2(E).\]
So \textbf{(e)} is true and \textbf{(c)} and \textbf{(d)} are false. Note that
$\alpha^2+\alpha+1$ has discriminant $-3$ and so is strictly positive for every real $\alpha$:
the two columns are independent whatever $\alpha$ is, $\Psi$ is an injective affine map, and the
formula is exact. Statement \textbf{(c)} is what an $\alpha$-linear guess produces, and
\textbf{(d)} is what one gets by taking a square root before, rather than after, forming the
Gram determinant.

For $\Phi$ the Jacobian is square, and there $\sqrt{\det(\Jac\Phi\transp\Jac\Phi)} =
\lvert\det\Jac\Phi\rvert$. Since
\[\Jac\Phi(x,y) = \begin{pmatrix} y & x \\ 2x & -2y \end{pmatrix}, \qquad
  \det\Jac\Phi(x,y) = -2y^2 - 2x^2,\]
the integrand is $2(x^2+y^2)$. Applied literally, \eqref{eq:mock_area_formula} therefore returns
$2\int_E(x^2+y^2)$, which is statement \textbf{(a)}, and never $2\int_E(x^2+2y^2)$, so
\textbf{(b)} is false.
\end{exercisesolution}

\begin{ainote}[The quoted area formula needs injectivity, and $\Phi$ does not have it]
\label{note:area_formula_needs_injectivity}
\Cref{ex:area_formula_two_maps}\textbf{(a)} is marked \emph{true} in the official key, and it is
what \eqref{eq:mock_area_formula} gives. But \eqref{eq:mock_area_formula} is false as printed,
because it carries no injectivity hypothesis, and $\Phi$ is exactly the map that exposes this:
$\Phi(-p) = \Phi(p)$ for every $p$.

Concretely, take $E := \overline{B_1(0)}$ and put $u := xy$, $v := x^2-y^2$. Then
\[4u^2 + v^2 = 4x^2y^2 + (x^2-y^2)^2 = (x^2+y^2)^2 ,\]
so $\Phi$ maps $E$ into the ellipse $\{4u^2+v^2 \leq 1\}$, and onto it: given such a $(u,v)$,
set $s := \sqrt{4u^2+v^2} \leq 1$ and solve $x^2 = (s+v)/2$, $y^2 = (s-v)/2$, both non-negative
because $v^2 \leq s^2$. These satisfy $x^2+y^2 = s \leq 1$ and $(xy)^2 = (s^2-v^2)/4 = u^2$, so a
choice of signs gives $xy = u$. Hence
\[\Phi(E) = \{(u,v) \mid 4u^2+v^2 \leq 1\},\]
an ellipse with semi-axes $\tfrac12$ and $1$, of area $\pi/2$. The right-hand side of
\eqref{eq:mock_area_formula}, meanwhile, is
\[\int_E 2(x^2+y^2)\,dx\,dy = 2\int_0^{2\pi}\!\!\int_0^1 r^2\cdot r\,dr\,d\theta = \pi .\]
The two differ by a factor of $2$, which is precisely the number of preimages. What is true in
general is that the integral counts multiplicity,
\[\int_E \lvert\det\Jac\Phi\rvert\,dx = \int_{\mathbb{R}^n} \#\big(\Phi^{-1}(w)\cap E\big)\,dw
  \ \geq\ \vol_n(\Phi(E)),\]
with equality exactly when $\Phi$ is injective off a null set. \Cref{thm:change_of_variables}
carries the injectivity hypothesis for this reason, and so does
\cref{def:d_volume_parametrized_submanifold}, which is stated for a \emph{parametrization} rather
than for an arbitrary $C^1$ map.

The verdict on part \textbf{(a)} is therefore: true as the examination intends it, false as
mathematics. It is kept here because the gap is more instructive than the computation.
\end{ainote}

\begin{exercisesolution}[Solution to \cref{ex:paraboloid_hypersurface_volume}]
\label{sol:paraboloid_hypersurface_volume}
% This label is referenced from 02-volume-of-embedded-surfaces.tex with \cpageref, NOT \cref.
% exercisesolution wraps `proof`, which has no counter, so a \label inside it is silently
% misattributed to the last-stepped ambient counter: \cref printed "Section 21.e" here until
% 2026-08-09. Page references are unaffected and are the documented workaround in style.md.
\label{ex:solution_9.1}
% Generator: Claude Sonnet 5 (Medium)
Following the official hint rather than Corsin's (\cref{note:9.1_coordinates_mislabelled}), we use
cylindrical coordinates $(x,y,z) = (r\cos\theta,r\sin\theta,z)$, under which $x^2+y^2=r^2$.
The two bounding surfaces become $r^2+z^2=8$ and $z = r^2/2$; intersecting them,
\[z^2+2z-8 = 0 \quad\Longrightarrow\quad z = 2 \text{ or } z=-4,\]
and since $z=r^2/2\geq0$ on the paraboloid, the surfaces meet at $z=2$, $r=2$. For
$r \in (0,2)$ the paraboloid $z=r^2/2$ lies below the upper half of the sphere
$z=\sqrt{8-r^2}$, so
\[B = \left\{(r,\theta,z) : 0<r<2,\ 0<\theta<2\pi,\ \tfrac12r^2 < z < \sqrt{8-r^2}\right\}.\]
Hence, with $dx\,dy\,dz = r\,dr\,d\theta\,dz$,
\[
\begin{aligned}
\vol(B) &= \int_0^{2\pi}\!\!\int_0^2 r\left(\sqrt{8-r^2} - \frac{r^2}{2}\right)dr\,d\theta
  = 2\pi\int_0^2 \left(r\sqrt{8-r^2} - \frac{r^3}{2}\right)dr.
\end{aligned}
\]
For the first term, substitute $u = 8-r^2$, $du=-2r\,dr$:
\[\int_0^2 r\sqrt{8-r^2}\,dr = \left[-\frac{u^{3/2}}{3}\right]_{u=8}^{u=4}
  = \frac{8^{3/2}-4^{3/2}}{3} = \frac{16\sqrt2-8}{3}.\]
For the second, $\int_0^2 \tfrac12 r^3\,dr = \tfrac12\cdot\tfrac{16}{4} = 2$. So
\[\vol(B) = 2\pi\left(\frac{16\sqrt2-8}{3} - 2\right) = 2\pi\cdot\frac{16\sqrt2-14}{3}
  = \boxed{\ \frac{4\pi}{3}\big(8\sqrt2-7\big)\ }.\]
\end{exercisesolution}

\begin{exercisesolution}[Solution to \cref{ex:curve_length_polar_and_vector_integral}]
\label{sol:curve_length_polar_and_vector_integral}
% Generator: Claude Sonnet 5 (Medium)
\begin{enumerate}[label=\textbf{(\alph*)}]
  \item By Fubini (\cref{thm:fubini}),
    \[\iint_D (x^3+3x^2y+y^3)\,dx\,dy = \int_0^2\int_0^1 (x^3+3x^2y+y^3)\,dy\,dx
      = \int_0^2\left(x^3+\frac32x^2+\frac14\right)dx.\]
    The last integral is
    $\left[\tfrac{x^4}{4}+\tfrac{x^3}{2}+\tfrac{x}{4}\right]_0^2 = 4+4+\tfrac12
    = \boxed{\ \tfrac{17}{2}\ }$.
  \item Using Corsin's parametrization $D=\{0\leq x\leq y\leq\pi\}$,
    \[\iint_D x\cos(x+y)\,dx\,dy = \int_0^\pi\int_x^\pi x\cos(x+y)\,dy\,dx
      = \int_0^\pi x\big[\sin(x+y)\big]_{y=x}^{y=\pi}\,dx.\]
    Since $\sin(x+\pi) = -\sin x$, the bracket equals $-\sin x - \sin(2x)$, so
    \[\iint_D x\cos(x+y)\,dx\,dy = -\int_0^\pi x\sin x\,dx - \int_0^\pi x\sin(2x)\,dx.\]
    Integrating by parts, $\int_0^\pi x\sin x\,dx = [-x\cos x+\sin x]_0^\pi = \pi$, and
    $\int_0^\pi x\sin(2x)\,dx = \left[-\tfrac{x}{2}\cos(2x)+\tfrac14\sin(2x)\right]_0^\pi
    = -\tfrac\pi2$. Hence
    \[\iint_D x\cos(x+y)\,dx\,dy = -\pi - \left(-\frac\pi2\right) = \boxed{\ -\frac\pi2\ }.\]
  \item Here $D = \{1<x<2,\ 1<y<3-x\}$, so
    \[\iint_D \frac{1}{(x+y)^3}\,dx\,dy = \int_1^2\int_1^{3-x} (x+y)^{-3}\,dy\,dx
      = \int_1^2 \left[-\frac12(x+y)^{-2}\right]_{y=1}^{y=3-x}dx.\]
    Since $x+(3-x)=3$, the bracket is $-\tfrac12\big(3^{-2}-(x+1)^{-2}\big)
    = \tfrac12(x+1)^{-2} - \tfrac1{18}$, and
    \[\int_1^2\left(\frac{1}{2(x+1)^2}-\frac1{18}\right)dx
      = \left[-\frac{1}{2(x+1)}\right]_1^2 - \frac1{18}
      = \left(-\frac16+\frac14\right) - \frac1{18} = \frac1{12}-\frac1{18}
      = \boxed{\ \frac1{36}\ }.\]
\end{enumerate}
\end{exercisesolution}

% All three parts of ex:curve_length_polar_and_vector_integral and ex:paraboloid_hypersurface_volume were worked out independently before consulting
% Sol9_Analysis2_eng.pdf; all four answers (4pi/3 (8sqrt2 - 7), 17/2, -pi/2, 1/36) match the
% official solution exactly, so there is no divergence to record. (Was an ainote until
% 2026-08-09: style.md names cross-checking against SolN as editor-facing under Test 2.)

\begin{exercisesolution}[Solution to \cref{ex:spherical_quadrilateral_area}]
\label{sol:spherical_quadrilateral_area}
% Generator: Claude Sonnet 5 (Medium)
We parametrize $\mathbb{S}^2$ (up to a null set) using latitude/longitude spherical
coordinates
\[\Phi : (-\pi,\pi)\times\left(-\tfrac\pi2,\tfrac\pi2\right) \to \mathbb{S}^2, \qquad
  (\varphi,\theta) \mapsto \begin{pmatrix}\cos\varphi\cos\theta\\ \sin\varphi\cos\theta\\
  \sin\theta\end{pmatrix}.\]
Note that this is \emph{not} the $(r\equiv1)$ parametrization used elsewhere in this chapter,
whose polar angle runs over $(0,\pi)$. Here $\theta$ is measured from the equator, matching
the exercise's range $\theta\in(-\pi/2,\pi/2)$. Computing
\[\partial_\varphi\Phi\times\partial_\theta\Phi = \begin{pmatrix}-\sin\varphi\cos\theta\\
  \cos\varphi\cos\theta\\ 0\end{pmatrix} \times \begin{pmatrix}-\cos\varphi\sin\theta\\
  -\sin\varphi\sin\theta\\ \cos\theta\end{pmatrix} = \begin{pmatrix}\cos\varphi\cos^2\theta\\
  \sin\varphi\cos^2\theta\\ \sin\theta\cos\theta\end{pmatrix},\]
its length is $\lvert\partial_\varphi\Phi\times\partial_\theta\Phi\rvert = \cos\theta$ (for
$\theta\in(-\pi/2,\pi/2)$, so $\cos\theta>0$). By \cref{prop:flux_r3}\textbf{(b)}, this is the
Gram determinant, so $dA = \cos\theta\,d\theta\,d\varphi$, and
\[\int_S dA = \int_{\varphi_1}^{\varphi_2}\!\!\int_{\theta_1}^{\theta_2} \cos\theta\,d\theta\,
  d\varphi = \boxed{\ (\varphi_2-\varphi_1)\big(\sin\theta_2-\sin\theta_1\big)\ }.\]
\end{exercisesolution}

\begin{exercisesolution}[Solution to \cref{ex:torus_surface_area}]
\label{sol:torus_surface_area}
% Generator: Claude Sonnet 5 (Medium)
The circle $S$ has radius $2$, and $T$ is the solid torus swept by a tube of radius $1$
around it, so $\partial T$ is Corsin's general torus surface with $R=2$, $r=1$ fixed. Using
his parametrization,
\[\phi(\theta,\varphi) = \begin{pmatrix}(2+\cos\varphi)\cos\theta\\
  (2+\cos\varphi)\sin\theta\\ \sin\varphi\end{pmatrix}, \qquad \theta,\varphi\in(0,2\pi).\]
Then $\partial_\theta\phi = \big(-(2+\cos\varphi)\sin\theta,\ (2+\cos\varphi)\cos\theta,\
0\big)$ has length $2+\cos\varphi$, and $\partial_\varphi\phi = (-\sin\varphi\cos\theta,\
-\sin\varphi\sin\theta,\ \cos\varphi)$ has length $1$; the two are orthogonal (their dot
product is $\sin\varphi\cos\varphi\sin\theta\cos\theta - \sin\varphi\cos\varphi\sin\theta
\cos\theta = 0$), so
\[\sqrt{\det D\phi\transp D\phi} = \lvert\partial_\theta\phi\rvert\,\lvert\partial_\varphi\phi
  \rvert = 2+\cos\varphi.\]
Hence
\[\vol_2(\partial T) = \int_0^{2\pi}\!\!\int_0^{2\pi} (2+\cos\varphi)\,d\theta\,d\varphi
  = 2\pi\left(2\cdot2\pi + \underbrace{\int_0^{2\pi}\cos\varphi\,d\varphi}_{=0}\right)
  = \boxed{\ 8\pi^2\ }.\]
\end{exercisesolution}

\begin{exercisesolution}[Solution to \cref{ex:solids_of_revolution_area}]
\label{sol:solids_of_revolution_area}
% Generator: Claude Sonnet 5 (Medium)
\begin{enumerate}[label=\textbf{(\alph*)}]
  \item $\partial_{\mathrm{side}}\Omega$ is not closed (it omits, e.g., the limit
    point $(a,f(a),0)$ as $x\to a^+$) and not open (a $2$-dimensional surface has empty interior
    in $\mathbb{R}^3$), but it \emph{is} path-connected: from any point, first move along the
    graph of $f$ in the $x$-direction, then move around the circle in the $(y,z)$-variables.
    $\Omega$ is open (as the strict sublevel set of a continuous function), not closed, and its
    connectedness depends on $f$: if $f$ vanishes somewhere in $(a,b)$, the cross-section pinches
    to a point there and $\Omega$ is disconnected into the pieces on either side; if $f>0$
    throughout $(a,b)$, $\Omega$ is connected.

  \item In cylindrical coordinates $(x,\rho,\theta)$ with $y=\rho\cos\theta$,
    $z=\rho\sin\theta$, $dx\,dy\,dz = \rho\,d\rho\,d\theta\,dx$, the domain becomes the box-like
    region $\{a<x<b,\ 0<\rho<f(x),\ \theta\in[0,2\pi)\}$, so
    \[\vol_3(\Omega) = \int_a^b\!\!\int_0^{2\pi}\!\!\int_0^{f(x)} \rho\,d\rho\,d\theta\,dx
      = \int_a^b\!\!\int_0^{2\pi} \frac12 f(x)^2\,d\theta\,dx
      = \boxed{\ \pi\int_a^b f(x)^2\,dx\ }.\]

  \item Using the suggested parametrization $\Phi(x,\theta) = (x,f(x)\cos\theta,
    f(x)\sin\theta)$, write the orthonormal frame $\xi_1:=(1,0,0)$, $\xi_2:=(0,\cos\theta,
    \sin\theta)$, $\xi_3:=(0,-\sin\theta,\cos\theta)$, so that
    \[\partial_x\Phi = \xi_1 + f'(x)\xi_2, \qquad \partial_\theta\Phi = f(x)\xi_3.\]
    Since $\xi_1,\xi_2,\xi_3$ are orthonormal and right-handed, $\xi_1\times\xi_3=-\xi_2$ and
    $\xi_2\times\xi_3=\xi_1$, so
    \begin{align*}
      \partial_x\Phi\times\partial_\theta\Phi
        &= f(x)\big(\xi_1\times\xi_3 + f'(x)\,\xi_2\times\xi_3\big)
         = f(x)\big({-\xi_2}+f'(x)\xi_1\big), \\
      \lvert\partial_x\Phi\times\partial_\theta\Phi\rvert &= f(x)\sqrt{1+f'(x)^2}.
    \end{align*}
    By \cref{prop:flux_r3}\textbf{(b)}, this is the Gram determinant, so
    \[\vol_2(\partial_{\mathrm{side}}\Omega) = \int_a^b\!\!\int_0^{2\pi} f(x)\sqrt{1+f'(x)^2}\,
      d\theta\,dx = \boxed{\ 2\pi\int_a^b f(x)\sqrt{1+f'(x)^2}\,dx\ }.\]

  \item With $f(x)=1/x$ on $[1,\infty)$, both integrands are non-negative, so the
    improper integrals are unambiguous (possibly $+\infty$). The volume is finite:
    \[\pi\int_1^\infty \frac{1}{x^2}\,dx = \pi\Big[-\frac1x\Big]_1^\infty = \boxed{\ \pi\ }.\]
    For the surface area, $f'(x) = -1/x^2$, so the integrand is $\sqrt{1+x^{-4}}/x \geq 1/x$
    pointwise, and since $\int_1^\infty \frac1x\,dx$ diverges, so does
    \[2\pi\int_1^\infty \frac{\sqrt{1+x^{-4}}}{x}\,dx = \boxed{\ +\infty\ }.\]
    This is the classical \newterm{Gabriel's horn} \germanfor{Torricellis Trompete}: finite
    volume, infinite surface area.

  \item ($\Leftarrow$) Suppose $(*)$ holds. Since $f'(a^+)=+\infty$, the inverse
    function theorem applies to $f|_{(a,a+\varepsilon)}$ for $\varepsilon$ small, giving a $C^1$,
    increasing inverse there; extend it by
    \[g(y) := \begin{cases} f^{-1}(y) & f(a)<y<f(a)+\varepsilon, \\ a & y\leq f(a). \end{cases}\]
    Since $f'(a^+)=+\infty$, $\lim_{y\to f(a)^+}g'(y) = \lim_{x\to a^+} 1/f'(x) = 0 =
    \lim_{y\to f(a)^-}g'(y)$, so $g\in C^1$ across $y=f(a)$. Around a boundary point $P :=
    (a,f(a)\cos\phi,f(a)\sin\phi)$, the map $(t,\theta)\mapsto\big(g(t),f(g(t))\cos\theta,
    f(g(t))\sin\theta\big)$ is then a regular $C^1$ parametrization of $\partial\Omega$ (using
    $f(a)>0$ to keep the angular direction non-degenerate). The symmetric argument at $b$, using
    $f'(b^-)=-\infty$, handles the other end, so $\partial\Omega$ is a $C^1$ submanifold and
    $\Omega$ is a bounded $C^1$ domain.

    ($\Rightarrow$) Suppose $\Omega$ is a bounded $C^1$ domain. If $f$ vanished at some interior
    point, $\Omega$ would be disconnected by part \textbf{(a)}, which
    \cref{def:bounded_ck_domain} excludes, since it requires $\Omega$ itself to make sense as a
    single bounded open set with $(n-1)$-manifold boundary (implicit in the exercise's setup).
    Therefore $f>0$ on $(a,b)$. For the endpoint condition at $a$: the point
    $P:=(a,f(a),0)\in\partial\Omega$ is a smooth boundary point, so
    $T_P(\partial\Omega)$ is $2$-dimensional and contains
    \[\frac{d}{dt}\Big|_{t=0^+}\big(a,f(a)-t,0\big) = -e_2, \qquad
      \frac{d}{dt}\Big|_{t=0^-}\big(a,f(a)\cos t,f(a)\sin t\big) = f(a)e_3,\]
    so $T_P(\partial\Omega) = \operatorname{span}\{e_2,e_3\}$. But $T_P(\partial\Omega)$ must
    also contain the limit of the (rescaled, if necessary) tangent vectors
    \[v_\varepsilon := \frac{d}{dt}\Big|_{t=0^+}\big(a+\varepsilon+t,\ f(a+\varepsilon+t),\
      0\big) = \big(1,\ f'(a+\varepsilon),\ 0\big)\]
    as $\varepsilon\to0^+$ (these are tangent to $\partial\Omega$ at nearby points converging to
    $P$). If $f'(a+\varepsilon)$ stayed bounded along some sequence $\varepsilon_k\to0^+$, the
    direction of $v_{\varepsilon_k}$ would converge to some $\alpha e_1+\beta e_2 \notin
    \operatorname{span}\{e_2,e_3\}$ (its $e_1$-component is the nonzero constant $1$, up to
    normalization), contradicting $T_P(\partial\Omega)=\operatorname{span}\{e_2,e_3\}$. So
    $f'(a^+)=+\infty$; the symmetric argument at $b$ gives $f'(b^-)=-\infty$.
\end{enumerate}
\end{exercisesolution}

\begin{ainote}
All parts of \cref{ex:spherical_quadrilateral_area,ex:torus_surface_area,ex:solids_of_revolution_area} were worked out independently before consulting
\texttt{Sol10\_Analysis2\_eng.pdf}. The three main answers ($(\varphi_2-\varphi_1)(\sin\theta_2
-\sin\theta_1)$, $8\pi^2$, $\pi\int_a^bf(x)^2dx$ / $2\pi\int_a^bf(x)\sqrt{1+f'(x)^2}dx$ /
volume $\pi$, area $+\infty$) match exactly. Part \textbf{(e)} was reconstructed independently
along the same two-directional strategy (build an explicit $C^1$ parametrization near each
endpoint using the inverse function theorem; for the converse, derive a contradiction from the
tangent space being forced to both contain $\operatorname{span}\{e_2,e_3\}$ and the limiting
direction of $v_\varepsilon$) and agrees with the official solution in substance, though the
write-up above is condensed relative to the official one.
\end{ainote}



\part{Vector Calculus and Differential Forms}

% Generator: Claude Opus 5 (High) --- part-level framing prose, new content.
The classical integral theorems of vector calculus look at first like three unrelated results.
Green's theorem converts a line integral into a double integral, the divergence theorem converts a
flux through a closed surface into a triple integral, and the theorem of Stokes converts a
circulation into a flux. This part proves all three, and then explains why they are one theorem.
The apparatus that achieves this consists of differential forms and the exterior derivative $d$,
a single operator that absorbs $\grad$, $\curl$ and $\divg$ at once. In that language each of the
statements above reads $\int_M d\omega = \int_{\partial M}\omega$, and the classical versions are
what it becomes when $M$ is a plane region, a solid, or a surface. Everything built earlier is
needed here: submanifolds to integrate over, orientation to fix the sign, and compactness to keep
the integrals finite.

% ============================================================
% Chapter 22 --- Vector calculus operators
% Originally: content/week-10.tex, \section{Bounded C^1 domains ...} and
%             \section{Vector calculus operators} (Corsin Week 10), plus problem 11.1.
% ============================================================
\chapter{Vector Calculus Operators}
\label{ch:vector_calculus}

% Transition: Claude Opus 5 (Medium)
The integral theorems of this part relate what happens inside a domain to what happens on its
boundary, so two things have to come first. One is a class of domains whose boundary is well
behaved enough to integrate over. The other is the differential operators (gradient, divergence,
curl) whose interplay the theorems express. This chapter supplies both, and says why the
definition of a bounded $C^1$ domain is stated in the awkward way it is.

% Originally: content/week-10.tex, \section{Bounded C^1 domains ...} (Corsin Week 10)

\section{Bounded \texorpdfstring{$C^1$}{C1} domains and why the definition in the script is so terrible}
\label{sec:bounded_c1_domains}

% Source: Corsin Nick/Class Notes/Week 10.pdf, p. 6
% Def. 14.3
\begin{definition}[Bounded $C^k$ domain]
\label{def:bounded_ck_domain}
Given $k \geq 1$, a bounded open set $\Omega \subset \mathbb{R}^n$ is called a
\newterm{bounded $C^k$ domain} if $M := \partial\Omega$ is an $(n-1)$-dimensional submanifold
of class $C^k$.
\end{definition}

Citing the 2024 lecture notes, Corsin builds up the machinery needed to actually \emph{use}
this definition. To define \newterm{flux} we need the concept of interior and exterior
normal vectors. For $\Omega$ as above, a normal vector $v \in N_p(\partial\Omega)$ is
\newterm{interior} if
\[p + hv \in \Omega \quad \text{for all } h>0 \text{ small},\]
and \newterm{exterior} if
\[p + hv \in \mathbb{R}^n\setminus\Omega \quad \text{for all } h>0 \text{ small}.\]

A \newterm{Gauss map} is a continuous map of unit normal vectors on $\partial\Omega$. There
exists a unique interior and exterior Gauss map; write
\[\nu \in C^0(\partial\Omega, \mathbb{S}^{n-1}) \text{ such that } \nu(p) \text{ is exterior
  normal}\]
for the exterior unit normal map.

\begin{theorem}[Jordan--Brouwer separation theorem]
\label{thm:jordan_brouwer}
If $M \subseteq \mathbb{R}^{n+1}$ is an $n$-dimensional compact and connected submanifold,
then $\mathbb{R}^{n+1}\setminus M$ has exactly two connected components, an
\newterm{inside} and an \newterm{outside}.
\end{theorem}

% Supplement: Sascha Brack/Class Notes/Week_10_Notes_Monday.pdf, pp. 8--11

% Originally: content/week-10.tex, \section{Vector calculus operators} (Corsin Week 10)

\section{Vector calculus operators}
\label{sec:vector_calculus_operators}

\begin{definition}[Vector field]
\label{def:vector_field}
Let $U \subseteq \mathbb{R}^n$ be open. A \newterm{$C^k$ vector field} (\germanterm{Vektorfeld}) is a function $F : U \to \mathbb{R}^n$ of class $C^k(U,\mathbb{R}^n)$.
The scalar functions $F_1, \dots, F_n$ in $F(x) = (F_1(x), \dots, F_n(x))\transp$ are called the \newterm{components}.
\end{definition}

\begin{definition}[Gradient]
\label{def:gradient_vector_field}
Let $f : U \to \mathbb{R}$ be a function of class $C^k$ on an open set $U \subseteq \mathbb{R}^n$. The \newterm{gradient} of $f$ is the $C^{k-1}$ vector field
\[\operatorname{grad} f(x) = \nabla f(x) := \begin{pmatrix} \partial_1 f(x) \\ \vdots \\ \partial_n f(x) \end{pmatrix}.\]
\end{definition}

\begin{definition}[Divergence]
\label{def:divergence}
For a $C^1$ vector field $F : U \to \mathbb{R}^n$, the \newterm{divergence} is the scalar function
\[\divg F(x) := \sum_{i=1}^n \partial_i F_i(x).\]
In 3D, physicists write this as $\vec{\nabla} \cdot \vec{F}(x) = \partial_1 F_1(x) + \partial_2 F_2(x) + \partial_3 F_3(x)$.
\end{definition}

\begin{ainote}
Intuitively, $\divg F(x)$ tells us how much flow of $F$ is generated at $x$. If $\divg F(x) = 0$, then the total generated flow is zero.
If $\divg F(x) > 0$, then $x$ is a \newterm{source} (\germanterm{Quelle}), meaning more flow is generated than removed.
If $\divg F(x) < 0$, then $x$ is a \newterm{sink} (\germanterm{Senke}), meaning more flow is removed than generated.

% Supplement: Sascha Brack/Class Notes/Week_10_Notes_Monday.pdf, p. 9
\begin{center}
\begin{tikzpicture}[scale=1]
  \node[circle, fill=blue!50, inner sep=1.5pt, label=above:{\small Sink ($\divg F < 0$)}] (sink) at (0,0) {};
  \foreach \angle in {0,45,...,315} {
    \draw[blue, thick, <-] (sink) -- ++(\angle:0.8);
  }
  
  \node[circle, fill=red!50, inner sep=1.5pt, label=above:{\small Source ($\divg F > 0$)}] (source) at (4,0) {};
  \foreach \angle in {0,45,...,315} {
    \draw[red, thick, ->] (source) -- ++(\angle:0.8);
  }
\end{tikzpicture}
\end{center}
\end{ainote}

\begin{definition}[Rotation/Curl]
\label{def:rotation}
Let $F$ be a $C^k$ vector field on an open set $U \subseteq \mathbb{R}^3$. The \newterm{rotation} or \newterm{curl} of $F$ is the $C^{k-1}$ vector field
\[\curl F(x) := \begin{pmatrix} \partial_2 F_3(x) - \partial_3 F_2(x) \\ \partial_3 F_1(x) - \partial_1 F_3(x) \\ \partial_1 F_2(x) - \partial_2 F_1(x) \end{pmatrix}.\]
Physicists write $\curl F(x) = \vec{\nabla} \times \vec{F}(x)$.
In $\mathbb{R}^2$, the rotation of a vector field $F = (F_1, F_2)\transp$ is the scalar function $\curl F(x) := \partial_1 F_2(x) - \partial_2 F_1(x)$.
\end{definition}

\begin{ainote}
Intuitively, $\curl F(x)$ tells us how much the field tends to rotate around $x$. If $\curl F(x) \neq 0$, a small wheel centered at $x$ would rotate.

% Supplement: Sascha Brack/Class Notes/Week_10_Notes_Monday.pdf, p. 10
\begin{center}
\begin{tikzpicture}[scale=1]
  % Rotates
  \draw[thick, ->] (0, 0.4) -- (1.5, 0.4);
  \draw[thick, ->] (0, 0) -- (1.0, 0);
  \draw[thick, ->] (0, -0.4) -- (0.5, -0.4);
  \node[blue, font=\small] at (0.75, -0.8) {rotates};
  \draw[blue, thick, ->] (0.55,-0.6) arc (-90:90:0.2);
  
  % Does not rotate
  \draw[thick, ->] (3, 0.4) -- (4.5, 0.4);
  \draw[thick, ->] (3, 0) -- (4.5, 0);
  \draw[thick, ->] (3, -0.4) -- (4.5, -0.4);
  \node[blue, font=\small] at (3.75, -0.8) {does not rotate};
  \draw[blue, thick, ->] (3.75,-0.6) -- (3.75,-0.2);
\end{tikzpicture}
\end{center}
\end{ainote}

\begin{example}[Computing vector calculus operators]
\label{ex:compute_operators}
Given $\phi : \mathbb{R}^3 \to \mathbb{R}$, $\phi(x,y,z) := x^2 + y^2 + z^2 e^{-xy}$.
Let us compute the gradient, the divergence of the gradient, and the rotation of the gradient.
\begin{itemize}
    \item \textbf{Gradient:}
    \[\nabla \phi(x,y,z) = \begin{pmatrix} 2x - yz^2 e^{-xy} \\ 2y - xz^2 e^{-xy} \\ 2z e^{-xy} \end{pmatrix}.\]
    
    \item \textbf{Divergence of the gradient} (denoted $\divg(\nabla \phi)$ or $\Delta \phi$):
    \begin{align*}
    \divg(\nabla \phi) &= \partial_x(2x - yz^2 e^{-xy}) + \partial_y(2y - xz^2 e^{-xy}) + \partial_z(2z e^{-xy}) \\
    &= (2 + y^2 z^2 e^{-xy}) + (2 + x^2 z^2 e^{-xy}) + (2 e^{-xy}) \\
    &= 4 + (x^2 z^2 + y^2 z^2 + 2) e^{-xy}.
    \end{align*}
    
    \item \textbf{Rotation of the gradient} ($\curl(\nabla \phi)$):
    \begin{align*}
    \curl(\nabla \phi) &= \begin{pmatrix} \partial_y(2z e^{-xy}) - \partial_z(2y - xz^2 e^{-xy}) \\ \partial_z(2x - yz^2 e^{-xy}) - \partial_x(2z e^{-xy}) \\ \partial_x(2y - xz^2 e^{-xy}) - \partial_y(2x - yz^2 e^{-xy}) \end{pmatrix} \\
    &= \begin{pmatrix} -2xz e^{-xy} - (-2xz e^{-xy}) \\ -2yz e^{-xy} - (-2yz e^{-xy}) \\ -z^2 e^{-xy} - (-z^2 e^{-xy}) \end{pmatrix} = \begin{pmatrix} 0 \\ 0 \\ 0 \end{pmatrix}.
    \end{align*}
    This is not a coincidence: for any $C^2$ function $f$, $\curl(\nabla f) = 0$ by Schwarz's theorem.
\end{itemize}
\end{example}

% Supplement: Analysis_II_Script_v1.pdf, p. 163
\begin{exercise}[$\curl\circ\grad=0$ and $\divg\circ\curl=0$]
\label{ex:curl_grad_div_curl_zero}
Let $f$ be a twice continuously differentiable function and $F$ a twice continuously
differentiable vector field on an open set $U \subseteq \mathbb{R}^3$. Show that
\[\curl(\grad f) = 0 \qquad\text{and}\qquad \divg(\curl F) = 0.\]
\end{exercise}

\begin{exercisesolution}[Proof of \cref{ex:curl_grad_div_curl_zero}]
% Generator: Claude Sonnet 5 (High)
Both identities follow from Schwarz's theorem (\cref{thm:schwarz_mixed_partials}), which
applies since $f,F \in C^2$.

\textbf{$\curl(\grad f)=0$.} The first component of $\curl(\grad f)$ is
$\partial_2(\partial_3f) - \partial_3(\partial_2f) = 0$ by Schwarz, and the other two
components vanish the same way, cyclically permuting the indices. This is exactly
\cref{ex:compute_operators}'s observation above, now stated and proved in general.

\textbf{$\divg(\curl F)=0$.}
\[\divg(\curl F) = \partial_1(\partial_2F_3-\partial_3F_2) + \partial_2(\partial_3F_1-\partial_1F_3)
  + \partial_3(\partial_1F_2-\partial_2F_1),\]
and every one of the six terms cancels against another by Schwarz: $\partial_1\partial_2F_3$
cancels $-\partial_2\partial_1F_3$, $-\partial_1\partial_3F_2$ cancels
$\partial_3\partial_1F_2$, and $\partial_2\partial_3F_1$ cancels $-\partial_3\partial_2F_1$.
\end{exercisesolution}

\begin{remark}
This is the classical-vector-calculus shadow of \qt{$d^2=0$}
(\cref{thm:exterior_derivative}\textbf{(2)}, \cref{ch:differential_forms}): $\grad$, $\curl$ and
$\divg$ are exactly the $0$-, $1$- and $2$-form incarnations of the exterior derivative in
$\mathbb{R}^3$ (\cref{ex:exterior_derivative_gradient_curl}), and both identities above are the
single fact $d\circ d=0$ read off once in each degree.
\end{remark}

% Supplement: Analysis_II_Script_v1.pdf, p. 164
\begin{proposition}[Product rule for the curl]
\label{prop:curl_product_rule}
Let $U \subseteq \mathbb{R}^3$ be open, $F : U \to \mathbb{R}^3$ a $C^1$ vector field, and
$\varphi : U \to \mathbb{R}$ a $C^1$ function. Then
\[\curl(\varphi F) = \varphi\,\curl(F) + (\grad\varphi)\times F.\]
\end{proposition}

\begin{proof}
By the product rule applied componentwise, e.g.\ for the first component,
\[\partial_2(\varphi F_3) - \partial_3(\varphi F_2)
  = \varphi(\partial_2F_3-\partial_3F_2) + (\partial_2\varphi)F_3 - (\partial_3\varphi)F_2,\]
and the first term is $\varphi\,\curl(F)$'s first component, the last two are
$((\grad\varphi)\times F)$'s first component. The other two components follow identically,
cyclically permuting the indices.
\end{proof}

\begin{ainote}
The official script sets this identity as an exercise (Ex.\ 14.61) but states it without the
first term, as $\curl(\varphi F) = (\grad\varphi)\times F$ --- which is false in general: taking
$\varphi\equiv1$ would force $\curl F \equiv 0$ for every $F$. The missing term is exactly
$\varphi\,\curl(F)$, restored above. (The script's exercise also names the vector field $f$ in
its statement and $F$ in its conclusion, as though the two were interchangeable.)
\end{ainote}


% --- exercises relocated here from the dissolved sheets ---

% Originally: content/week-11.tex, \exercisesheet{11}, problem 11.1
% Source: exercises/Ex11_Analysis2_eng.pdf

\begin{exercise}[11.1 --- Flow of vector field \textnormal{(important)}]
\label{ex:11.1}
Compute the flow of the vector field
\[V : (x,y,z) \mapsto (x,y,z-x^2-y^2)\]
through the upper hemisphere $S = \{(x,y,z)\in\mathbb{R}^3 : x^2+y^2+z^2=1,\ z\geq0\}$ from
\qt{inside} to \qt{outside}. That is the number
\[\int_S V(x)\cdot\nu(x)\,d\vol_2.\]
\end{exercise}

% Source: exercises/Ex11_Analysis2_eng.pdf, p. 2

% Originally: content/week-11.tex, \section{Solutions}

\section{Solutions}
\label{sec:vector_calculus_solutions}

\begin{exercisesolution}[Solution of \cref{ex:11.1}]
The vector field $V(x,y,z) = (x,y,z-x^2-y^2)$ has divergence
\[\divg V = \partial_x(x) + \partial_y(y) + \partial_z(z-x^2-y^2) = 1+1+1 = 3.\]
$S$ is not itself the boundary of a bounded $C^1$ domain, so close it up with the equatorial
disk $D := \{(x,y,0) : x^2+y^2\leq1\}$: then $\Omega := \{(x,y,z)\in B_1(0) : z\geq0\}$, the
upper half ball, is a bounded $C^1$ domain with $\partial\Omega = S\cup D$, and $S$'s
\qt{inside to outside} normal agrees with the outward normal of $\Omega$ on $S$. By
\cref{thm:gauss_divergence},
\[\int_S V\cdot\nu\,d\vol_2 + \int_D V\cdot\nu_D\,d\vol_2 = \int_\Omega \divg V\,dx
  = 3\vol_3(\Omega) = 3\cdot\frac{2\pi}{3} = 2\pi,\]
using $\vol_3(\Omega) = \frac12\cdot\frac43\pi = \frac{2\pi}{3}$. On $D$ the outward normal is
$\nu_D = -e_3$, and $V(x,y,0) = (x,y,-x^2-y^2)$, so $V\cdot\nu_D = x^2+y^2$; in polar
coordinates,
\[\int_D V\cdot\nu_D\,d\vol_2 = \int_0^{2\pi}\!\!\int_0^1 r^2\cdot r\,dr\,d\varphi
  = 2\pi\left[\frac{r^4}{4}\right]_0^1 = \frac{\pi}{2}.\]
Hence
\[\int_S V\cdot\nu\,d\vol_2 = 2\pi - \frac{\pi}{2} = \boxed{\ \frac{3\pi}{2}\ }.\]
\end{exercisesolution}

\begin{ainote}
Solved independently before consulting \texttt{Sol11\_Analysis2\_eng.pdf}; the official
solution uses the identical closing-disk argument and reaches the same value $3\pi/2$.
\end{ainote}


% ============================================================
% Chapter 23 --- Flux and the divergence theorem
% Originally: content/week-10.tex, \section{Flux and the divergence theorem}, together
%             with content/week-11.tex, \section{Gauss' theorem revisited: the case n=1}
%             (Corsin Week 11), plus problem 11.5. Corsin taught the n=1 warm-up a week
%             after the theorem itself; here it comes first, as motivation.
% ============================================================
\chapter{Flux and the Divergence Theorem}
\label{ch:flux}

% Transition: Claude Opus 5 (Medium)
The divergence theorem is the first of the integral theorems: the flux of a vector field out
through the boundary of a domain equals the integral of its divergence inside. The chapter opens
with the case $n=1$, where the statement degenerates into the fundamental theorem of calculus and
the mechanism is visible without any geometry, and then works up to the general theorem and its
physical reading as a conservation law.

% Originally: content/week-11.tex, \section{Gauss' theorem revisited: the case n=1} (Corsin Week 11)


\section{Gauss' theorem revisited: the case \texorpdfstring{$n=1$}{n=1}, and a direct proof
for a square}
\label{sec:gauss_intuition}
% Corsin taught this a week after the divergence theorem, so the source carried a
% \continuedfrom{sec:flux_divergence_theorem} back-pointer. Here the section comes first, as
% motivation, and the pointer would have run forwards; it is gone, and the sentence below
% already \cref's the theorem it previews.
%
% ⚠️ OPEN, raised by the 2026-08-13 audit and deliberately not resolved here. Moving the section
% forward is defensible for the n=1 warm-up, which is motivation and reads perfectly well before
% the statement. It is harder to defend for what follows it: the block below is a genuine
% \begin{proof}[Proof of thm:gauss_divergence for n=2, B=(0,1)^2], and it now sits about ninety
% lines BEFORE the theorem it proves is stated. A reader meets a proof of something they have not
% been told yet. Two repairs are available and each costs something:
%   (a) leave the n=1 warm-up here and move the n=2 proof into sec:flux_divergence_theorem, just
%       under thm:gauss_divergence, where a proof belongs. Splits one source page across two
%       sections, and the "conversely, the FTC can be used to PROVE Gauss" sentence loses the
%       n=1 calculation it answers.
%   (b) relabel the n=2 block as a worked example previewing the theorem rather than proving it,
%       which is closer to what it does in its current position, and keep everything in place.
% (b) is the smaller change and probably the right one. Not taken unilaterally because it alters
% what the block claims to be, and Corsin's own framing is that it proves the theorem.

% Source: Corsin Nick/Class Notes/Week 11.pdf, p. 2
Gauss' theorem (\cref{thm:gauss_divergence}), proved in \cref{sec:flux_divergence_theorem}
below, is, in a precise sense, a higher-dimensional
fundamental theorem of calculus: it trades a derivative integrated over a region for the
original quantity evaluated on the region's boundary. Setting $n=1$ makes the analogy
literal. Let $f \in C^1([a,b])$, viewed as a (one-dimensional) vector field. Since $\divg f =
f'$ in dimension $1$, Gauss' formula becomes
\[\int_a^b \frac{\partial f}{\partial x}\,dx = \int_{[a,b]} \divg f\,d\vol_1
  = \int_{\partial[a,b]} \langle f,\nu\rangle\,d\vol_0
  = f(b) - f(a) \qquad \text{(FTC, almost).}\]
The last step identifies $\int_{\partial[a,b]} \langle f,\nu\rangle\,d\vol_0$ with $f(b)-f(a)$
by treating $\partial[a,b] = \{a,b\}$ as a $0$-dimensional submanifold with outward unit
normals $\nu(a) = -1$, $\nu(b) = +1$, and $\vol_0$ as counting measure.

\begin{ainote}
Corsin flags this himself: \qt{there are several problems with this argument, for example
what a $0$-dim Jordan measure or $0$-dim submanifold is}, and notes that \qt{the definitions
from the lecture can be extended to include this case} without carrying out the extension.
The word \qt{almost} above is his. Not pursued further here; the point of the calculation is
the analogy with the FTC, not a rigorous $n=1$ instance of \cref{thm:gauss_divergence}.
\end{ainote}

% Source: Corsin Nick/Class Notes/Week 11.pdf, p. 3
\begin{proof}[Proof of \cref{thm:gauss_divergence} for $n=2$, $B = (0,1)^2$]
Conversely, the fundamental theorem of calculus can be used to \emph{prove} Gauss' theorem;
we carry this out for the unit square. Let $F : [0,1]^2 \to \mathbb{R}^2$ be a $C^1$ vector
field and $B := (0,1)^2$. The boundary $\partial B$ splits into four edges, with outward unit
normals $-e_1$, $e_1$, $-e_2$, $e_2$ on the left, right, bottom, and top edges respectively.

% Sonnet 5 (Medium) --- FIG-W11-01
\begin{center}
\begin{tikzpicture}[>=Latex, scale=1.6]
  \draw[thick] (0,0) rectangle (1,1);
  \draw[->] (0,0.5) -- (-0.6,0.5) node[left, font=\scriptsize] {$-e_1$};
  \draw[->] (1,0.5) -- (1.6,0.5) node[right, font=\scriptsize] {$e_1$};
  \draw[->] (0.5,0) -- (0.5,-0.6) node[below, font=\scriptsize] {$-e_2$};
  \draw[->] (0.5,1) -- (0.5,1.6) node[above, font=\scriptsize] {$e_2$};
  \node[below left, font=\scriptsize] at (0,0) {$(0,0)$};
  \node[below right, font=\scriptsize] at (1,0) {$(1,0)$};
  \node[above left, font=\scriptsize] at (0,1) {$(0,1)$};
  \node[above right, font=\scriptsize] at (1,1) {$(1,1)$};
\end{tikzpicture}
\end{center}

Parametrizing each edge by $t \mapsto (0,t)$, $(1,t)$, $(t,0)$, $(t,1)$ for $t \in [0,1]$,
\[
\int_{\partial B} \langle F,\nu\rangle\,d\vol_1
  = \int_0^1 \langle F(0,t),-e_1\rangle\,dt + \int_0^1 \langle F(1,t),e_1\rangle\,dt
  + \int_0^1 \langle F(t,0),-e_2\rangle\,dt + \int_0^1 \langle F(t,1),e_2\rangle\,dt.
\]
Grouping the first pair and the second pair,
\[
\int_{\partial B} \langle F,\nu\rangle\,d\vol_1
  = \int_0^1 \big[F_1(1,t) - F_1(0,t)\big]\,dt + \int_0^1 \big[F_2(t,1) - F_2(t,0)\big]\,dt.
\]
By the fundamental theorem of calculus applied to $t \mapsto F_1(s,t)$ along $s$ and to
$s \mapsto F_2(t,s)$ along the second variable,
\[
\int_{\partial B} \langle F,\nu\rangle\,d\vol_1
  = \int_0^1\!\!\left(\int_0^1 \partial_1 F_1(s,t)\,ds\right) dt
  + \int_0^1\!\!\left(\int_0^1 \partial_2 F_2(t,s)\,ds\right) dt,
\]
and by Fubini,
\[
\int_{\partial B} \langle F,\nu\rangle\,d\vol_1
  = \int_B \partial_1 F_1\,d\vol_2 + \int_B \partial_2 F_2\,d\vol_2
  = \int_B \divg F\,d\vol_2. \qedhere
\]
\end{proof}

% Originally: content/week-10.tex, \section{Flux and the divergence theorem} (Corsin Week 10)

\section{Flux and the divergence theorem}
\label{sec:flux_divergence_theorem}

% Source: Corsin Nick/Class Notes/Week 10.pdf, p. 7
\begin{definition}[Flux]
\label{def:flux}
Suppose $F : U \to \mathbb{R}^n$ is a $C^1$ vector field, $U \subseteq \mathbb{R}^n$ open, and
$\Omega \subseteq U$ is a $C^1$ bounded domain. The \newterm{flux} \germanfor{Fluss} of $F$
through $\partial\Omega$ is defined as
\[\int_{\partial\Omega} \langle F,\nu\rangle\,d\vol_{n-1} \qquad \text{(analysis notation)}
  \quad = \quad \int_{\partial\Omega} \vec F(x)\cdot d\vec A \qquad \text{(physics notation)}\]
\[= \int_B \big\langle F\circ\phi(x),\ \nu\circ\phi(x)\big\rangle
   \sqrt{\det D\phi(x)\transp D\phi(x)}\,dx \qquad \text{(explicit)},\]
where $\nu : \partial\Omega \to \mathbb{S}^{n-1}$ is the exterior unit normal, i.e.\ $\nu(y)$
is the outward-facing unit normal of $\partial\Omega$ at $y \in \partial\Omega$, and
$\phi : B \to \partial\Omega$ is a (local) parametrization of $\partial\Omega$. In general,
multiple integrals in the explicit formulation may be necessary.
\end{definition}

% Source: Corsin Nick/Class Notes/Week 10.pdf, p. 8
\begin{proposition}[Flux in $\mathbb{R}^3$]
\label{prop:flux_r3}
Let $\Omega \subseteq \mathbb{R}^3$ be a bounded $C^1$ domain and $F \in
C^1(\overline\Omega,\mathbb{R}^3)$ a vector field, and let $\phi : B \to \partial\Omega
\setminus N$ be a parametrization of $\partial\Omega\setminus N$, where $N \subseteq
\partial\Omega$ is a null set, $B \subseteq \mathbb{R}^2$. Then:
\begin{enumerate}[label=\textbf{\arabic*.}]
  \item $\displaystyle\nu(x) := \frac{\partial_1\phi(x)\times\partial_2\phi(x)}
    {\lvert\partial_1\phi(x)\times\partial_2\phi(x)\rvert}$ is an (interior or exterior) Gauss
    map (more precisely, $\nu\circ\phi^{-1}$ is a Gauss map, i.e.\
    $\nu(x) \in N_{\phi(x)}(\partial\Omega\setminus N)$).
  \item The Gram determinant is given by
    $\sqrt{\det D\phi(x)\transp D\phi(x)} = \lvert\partial_1\phi(x)\times\partial_2\phi(x)
    \rvert$.
  \item Therefore,
    \[\int_{\partial\Omega} \vec F\cdot d\vec A = \pm\int_B \big\langle F\circ\phi(x),\
      \partial_1\phi(x)\times\partial_2\phi(x)\big\rangle\,dx,\]
    where the sign depends on whether $\nu(x)$ is interior ($-$) or exterior ($+$).
\end{enumerate}
\end{proposition}

% Source: Corsin Nick/Class Notes/Week 10.pdf, pp. 9--10
\begin{example}[Flux through the unit sphere]
\label{ex:flux_unit_sphere}
Compute the flux of $F(x,y,z) := (xz,\,z,\,y)$ through the surface of $\mathbb{S}^2$ with
interior $B_1(0)$.

We use spherical coordinates with $r\equiv1$:
\[\phi : (0,2\pi)\times(0,\pi) \to \mathbb{R}^3, \qquad (\varphi,\theta) \mapsto
  \begin{pmatrix}\sin\theta\cos\varphi \\ \sin\theta\sin\varphi \\ \cos\theta\end{pmatrix}.\]

\textbf{Normal vector.}
\[\partial_\varphi\phi = \begin{pmatrix}-\sin\theta\sin\varphi\\ \sin\theta\cos\varphi\\
  0\end{pmatrix}, \qquad \partial_\theta\phi = \begin{pmatrix}\cos\theta\cos\varphi\\
  \cos\theta\sin\varphi\\ -\sin\theta\end{pmatrix}
  \quad \Longrightarrow \quad \partial_\varphi\phi\times\partial_\theta\phi
  = \begin{pmatrix}-\sin^2\theta\cos\varphi\\ -\sin^2\theta\sin\varphi\\
  -\sin\theta\cos\theta\end{pmatrix}.\]

\textbf{Sanity check: does this make sense, and is it exterior or interior?} We have
\[\partial_\varphi\phi\times\partial_\theta\phi = -\sin\theta\ \phi(\varphi,\theta)
  = -\sin\theta\,(x,y,z),\]
and the exterior normal of $\mathbb{S}^2$ is $+(x,y,z)$, so $\partial_\varphi\phi\times
\partial_\theta\phi$ is an \emph{interior} normal (for $\theta\in(0,\pi)$, $\sin\theta>0$).
So we multiply by $-1$ to switch the order in the cross product, matching the exterior sign
in \cref{prop:flux_r3}.

\textbf{The integral.}
\[
\begin{aligned}
\int_{\mathbb{S}^2} \vec F\cdot d\vec A
  &= \int_0^{2\pi}\!\!\int_0^\pi \begin{pmatrix}xz\\ z\\ y\end{pmatrix}\cdot
     \begin{pmatrix}\sin\theta\cos\varphi\\ \sin\theta\sin\varphi\\ \cos\theta\end{pmatrix}
     \sin\theta\,d\theta\,d\varphi \\
  &= \int_0^{2\pi}\!\!\int_0^\pi \Big[\sin^3\theta\cos\theta\cos^2\varphi
     + \sin^2\theta\cos\theta\sin\varphi + \sin^2\theta\cos\theta\sin\varphi\Big]\,d\theta\,
     d\varphi \\
  &= \pi\int_0^\pi \sin^3\theta\cos\theta\,d\theta
   = \pi\left[\frac{\sin^4\theta}{4}\right]_0^\pi = 0,
\end{aligned}
\]
using $\int_0^{2\pi}\cos^2\varphi\,d\varphi = \pi$ and $\int_0^{2\pi}\sin\varphi\,d\varphi=0$
for the middle two terms.
\end{example}

% Source: Corsin Nick/Class Notes/Week 10.pdf, p. 11
\begin{theorem}[Gauss]
\label{thm:gauss_divergence}
Let $\Omega \subseteq \mathbb{R}^n$ be a bounded $C^1$ domain and $F : \overline\Omega \to
\mathbb{R}^n$ a $C^1$ vector field. Then
\[\int_\Omega \divg F\,d\vol_n = \int_{\partial\Omega} \langle F,\nu\rangle\,d\vol_{n-1},
  \qquad \text{or, in physics notation,} \qquad
  \int_\Omega \vec\nabla\cdot\vec F\,dx = \int_{\partial\Omega} \vec F\cdot d\vec A.\]
\end{theorem}

\begin{ainote}
Corsin dates the theorem to Lagrange (discovered 1762) and Ostrogradsky (proven 1828), and
adds, next to the physics form, an unusually enthusiastic \qt{how cool is this!} The remark is
earned: a volume integral of a derivative is reduced to a boundary integral of the original
field, which is the shape of the fundamental theorem of calculus.

% Supplement: Sascha Brack/Class Notes/Week_10_Notes_Monday.pdf, p. 13
% Revised 2026-08-09, two faults found by rendering it:
%   * The label was placed at (-1.5,-1), which is ON the boundary, not inside it: the ellipse
%     is x^2/4 + y^2/2.25 = 1, so at x = -1.5 it passes through y = -1.5*sqrt(1-2.25/4)
%     = -0.992. The ultra-thick curve struck straight through the glyph. Moved to (-1.5,0.45),
%     where the boundary is at |y| = 0.992, so the label sits well inside and clear of the
%     nearest source at (-0.5,0.5), whose spokes reach only to x = -0.9.
%   * The sinks were drawn `<-` from the centre outward, putting all five arrowheads on the
%     centre dot, exactly the blob that spoiled the sink/source figure in \cref{ch:vector_calculus}.
%     Sources and sinks now both run between radius 0.12 and 0.40, outward and inward
%     respectively, so the two have the same footprint and differ only in direction.
\begin{center}
\begin{tikzpicture}[>=Latex, scale=1.5]
  % Domain Omega
  \draw[ultra thick] (0,0) ellipse (2cm and 1.5cm);
  \node at (-1.5,0.45) {$\Omega$};
  \node[right] at (2,-0.5) {$\partial\Omega$};

  % Internal sources (red): rays run OUTWARD, from radius 0.12 to 0.40.
  \foreach \x/\y in {-0.5/0.5, 0.8/-0.2, 0/-0.6, -1/-0.3} {
    \fill[red] (\x,\y) circle (1.5pt);
    \foreach \angle in {0,72,...,288} {
      \draw[red, thick, ->] (\x,\y) ++(\angle:0.12) -- ++(\angle:0.28);
    }
  }

  % Internal sinks (blue): the same rays run INWARD, from radius 0.40 to 0.12.
  \foreach \x/\y in {1/0.6, 0.2/0.8, -0.2/-1.1, 1.2/-0.8} {
    \fill[blue] (\x,\y) circle (1.5pt);
    \foreach \angle in {0,72,...,288} {
      \draw[blue, thick, ->] (\x,\y) ++(\angle:0.40) -- ++({\angle+180}:0.28);
    }
  }

  % Boundary flux: outward where the field leaves, inward where it enters. The `<-` form puts
  % the head on the boundary point itself, so these do not blob.
  \foreach \angle in {0,45,90,135,180,225,270} {
    \draw[red, thick, ->] ({2*cos(\angle)},{1.5*sin(\angle)}) -- ++(\angle:0.4);
  }
  \foreach \angle in {30,150,315} {
    \draw[blue, thick, <-] ({2*cos(\angle)},{1.5*sin(\angle)}) -- ++(\angle:0.4);
  }
\end{tikzpicture}
\end{center}
\end{ainote}

\begin{example}[Recomputing \cref{ex:flux_unit_sphere} via the divergence theorem]
\label{ex:flux_unit_sphere_via_divergence}
With $\Omega = B_1(0)$ and $F(x,y,z) = (xz,z,y)$ as before, $\divg F = \partial_x(xz) +
\partial_y(z) + \partial_z(y) = z + 0 + 0 = z$. So, in spherical coordinates ($z =
r\cos\theta$),
\[
\int_{\partial\Omega} \vec F\cdot d\vec A = \int_\Omega \divg F\,dx = \int_\Omega z\,dx
  = \int_0^1\!\!\int_0^{2\pi}\!\!\int_0^\pi \underbrace{\cos\theta\sin\theta}_{=\,\frac12\sin(2\theta),\
    \pi\text{-periodic}} r^3\,d\theta\,d\varphi\,dr = 0,
\]
since $\int_0^\pi \sin(2\theta)\,d\theta = 0$. This confirms \cref{ex:flux_unit_sphere} with none
of its cross-product bookkeeping.
\end{example}

\begin{ainote}
Corsin writes the divergence as $\divg F = z$ directly without the intermediate
$\partial_x(xz)+\partial_z(z)+\partial_y(y)$ step; spelled out above for clarity. Worth
noting how much the two computations of the same flux differ in character: the direct
definition needs an explicit Gauss map, an orientation sanity check, and three cross terms
in the integrand, while the divergence theorem needs only $\divg F$ and a coordinate change
already used in \cref{ex:paraboloid_patch_area}. This is the theorem's main
practical use, and it is not conceptual elegance but the plain fact that $\divg F$ is often far
simpler than $F$ restricted to a curved boundary.
\end{ainote}

% Supplement: Sascha Brack/Class Notes/Week_10_Notes_Monday.pdf, p. 14
\begin{example}[Computing flux using the divergence theorem]
\label{ex:flux_divergence_cube}
Let $\Omega := [-1,1]^3 \subset \mathbb{R}^3$ and $F(x,y,z) := (x^2, y^2, z^2)\transp$.
We compute the flux of $F$ out of $\Omega$ using the divergence theorem.
First, we find the divergence of $F$:
\[\divg F(x,y,z) = \partial_x(x^2) + \partial_y(y^2) + \partial_z(z^2) = 2x + 2y + 2z.\]
Then, integrating over $\Omega$:
\begin{align*}
\int_{\partial\Omega} \langle F(p),\ \nu(p)\rangle \mathop{d\text{vol}}_2(p) &= \int_\Omega \divg F(x,y,z) \mathop{d\text{vol}}_3(x,y,z) \\
&= \int_{-1}^1 \int_{-1}^1 \int_{-1}^1 (2x + 2y + 2z) \,dx\,dy\,dz.
\end{align*}
By symmetry, the integral of $2x$ over $[-1,1]$ is $0$, and similarly for $y$ and $z$. Therefore, the total flux is $0$.
\end{example}

% Source: Jérôme Paschoud/Notizen/Woche 10.2 Divergenzsatz.pdf
\begin{example}[Flux out of a cylinder]
\label{ex:divergence_cylinder}
Let $\Omega := \{(x,y,z) \in \mathbb{R}^3 \mid x^2 + y^2 \leq R^2,\ 0 \leq z \leq H\}$ be a cylinder of radius $R$ and height $H$. Consider the vector field $F(x,y,z) := (x^3 + \sin(z), y^3 + xz, z^3 + xy)\transp$.
We compute the outward flux of $F$ through $\partial\Omega$ using the divergence theorem.
First, we compute the divergence of $F$:
\[\divg F(x,y,z) = 3x^2 + 3y^2 + 3z^2 = 3(x^2 + y^2 + z^2).\]
By the divergence theorem, the flux is the integral of the divergence over $\Omega$. We compute this using cylindrical coordinates $(r, \theta, z)$:
\begin{align*}
\int_{\partial\Omega} \langle F(p),\ \nu(p)\rangle \mathop{d\text{vol}}_2(p) &= \int_\Omega 3(x^2 + y^2 + z^2) \mathop{d\text{vol}}_3(x,y,z) \\
&= \int_0^H \int_0^{2\pi} \int_0^R 3(r^2 + z^2) r \,dr\,d\theta\,dz \\
&= 2\pi \int_0^H \left[ \frac{3}{4}r^4 + \frac{3}{2}r^2 z^2 \right]_0^R \,dz \\
&= 2\pi \int_0^H \left( \frac{3}{4}R^4 + \frac{3}{2}R^2 z^2 \right) \,dz \\
&= 2\pi \left[ \frac{3}{4}R^4 z + \frac{1}{2}R^2 z^3 \right]_0^H \\
&= 2\pi \left( \frac{3}{4}R^4 H + \frac{1}{2}R^2 H^3 \right) = \frac{\pi}{2}R^2 H (3R^2 + 2H^2).
\end{align*}
\end{example}

% Supplement: Analysis_II_Script_v1.pdf, p. 145
\begin{example}[Archimedes' principle]
\label{ex:archimedes_principle}
A solid $\Omega \subset \mathbb{R}^3$ (a bounded $C^1$ domain) is fully submerged in water, with
the coordinates chosen so that $x_3=0$ is the water's surface. At depth $-x_3$ the water
pressure is $P(x) := -\rho gx_3$, where $\rho$ is the density of water and $g$ the gravitational
acceleration. The force the water exerts on a surface element near $q \in \partial\Omega$ is
proportional to the pressure and points inward, along $-\nu(q)$, so the total buoyant force is
\[\vec F := -\int_{\partial\Omega}P\,\nu\,dS.\]
Substituting the pressure and applying the divergence theorem to each component,
\[F_i = \int_{\partial\Omega}\rho gx_3\nu_i\,dS
  = \rho g\int_{\partial\Omega}\langle x_3e_i,\nu\rangle\,dS
  = \rho g\int_\Omega\divg(x_3e_i)\,dV.\]
Since $\divg(x_3e_i) = \partial_i x_3$ equals $1$ if $i=3$ and $0$ otherwise, this gives
$F_1=F_2=0$ and
\[F_3 = \rho g\int_\Omega 1\,dV = \rho g\,\vol_3(\Omega).\]
The buoyant force is vertical, with magnitude $\rho g\vol_3(\Omega)$, which is exactly the weight
of the water displaced by $\Omega$, confirming Archimedes' principle. The divergence theorem
converts a surface integral of an unknown, position-dependent pressure into a volume integral
that only sees the shape of $\Omega$ through its volume.
\end{example}

% Source: Corsin Nick/Class Notes/Week 10.pdf, p. 12
Gauss' theorem is also a powerful tool in physics. Consider a fluid of density $\varrho$ moving
with velocity $\vec v$, and a surface element $d\vec A$ whose normal makes an angle $\theta$
with $\vec v$. The fluid crossing that element in time $\Delta t$ is exactly the fluid filling
the column of slant length $L = \Delta t\,\lvert\vec v\rvert$ erected on it along $\vec v$, and
that column stands at perpendicular height $L\cos\theta$. Its volume is therefore
\[V = \lvert d\vec A\rvert\,L\cos\theta = \Delta t\,\vec v\cdot d\vec A,\]
and the mass of fluid carried through the element is $\varrho V$.

\begin{ainote}
Corsin writes this as $V = \Delta t\,\vec v\cdot d\vec A\cdot\varrho$ and calls the result a
volume. The right-hand side is a mass: $\Delta t\,\vec v\cdot d\vec A$ is already the volume
swept, and the density only converts it. The two have been separated above, since the
distinction is needed immediately. The integral below is of $\varrho\vec v$ and therefore
measures mass per unit time, which is what the continuity equation balances against
$\partial\varrho/\partial t$.
\end{ainote}

% Generator: Opus 5 (high) --- FIG-W10-03, a redraw. The Sonnet 5 original was wrong in four
% ways, all of them visible only once rendered: it drew $\vec v$ twice in two different
% directions (a horizontal stream and a tilted arrow); the tilted arrow's label collided with
% the red one; it never drew the normal $d\vec A$ at all, so the $\theta$ arc sat between the
% velocity and nothing; and it labelled the height $h = \lvert d\vec A\rvert\cos\theta$, which
% is an area times a cosine.
%
% Geometry, all of it checked. The surface element is the segment (0,0)--(2,0); its normal
% dA runs along +y; the velocity makes the angle theta = 40 degrees with that normal, i.e.
% it runs at 50 degrees to the horizontal. The column swept in time Delta t is the element
% translated by L*(cos 50, sin 50) with L = 1.9, that is by (1.221, 1.455). Its perpendicular
% height is L*cos 40 = 1.455, which is exactly the y-extent of the parallelogram, so the h
% dimension line on the left measures the true height rather than an eyeballed one.
\begin{center}
\begin{tikzpicture}[>=Latex, scale=1.15]
  % The column swept through the element in time Delta t.
  \fill[purple!7] (0,0) -- (2,0) -- (3.221,1.455) -- (1.221,1.455) -- cycle;
  \draw[purple] (0,0) -- (1.221,1.455);
  \draw[purple, dashed] (1.221,1.455) -- (3.221,1.455);

  % The surface element itself, and the slant length of the column standing on it.
  \draw[purple, ultra thick] (0,0) -- (2,0);
  \node[purple, below, font=\scriptsize] at (1,-0.05) {$\lvert d\vec A\rvert$};
  \path (0,0) -- (1.221,1.455)
    node[midway, sloped, below, purple, font=\scriptsize] {$L=\Delta t\lvert\vec v\rvert$};

  % The normal, and the velocity at angle theta to it.
  \draw[blue, thick, ->] (2,0) -- (2,1.05) node[above, font=\scriptsize] {$d\vec A$};
  \draw[orange!90!black, thick, ->] (2,0) -- (3.221,1.455)
    node[right, font=\scriptsize] {$\Delta t\,\vec v$};
  \draw[green!50!black] (2,0.66) arc (90:50:0.66);
  \node[green!50!black, font=\scriptsize] at (2.17,0.42) {$\theta$};

  % The perpendicular height, drawn at the parallelogram's true y-extent.
  \draw[red!70!black, dashed] (-0.3,1.455) -- (1.221,1.455);
  \draw[red!70!black, dashed] (-0.3,0) -- (0,0);
  \draw[red!70!black, <->] (-0.3,0) -- (-0.3,1.455);
  \node[red!70!black, left, font=\scriptsize] at (-0.38,0.73) {$h=L\cos\theta$};
\end{tikzpicture}
\end{center}

% Source: Corsin Nick/Class Notes/Week 10.pdf, p. 12
Therefore, the fluid flowing out of any bounded $C^1$ domain $\Omega$ per unit time is
\[\int_{\partial\Omega} \varrho\vec v\cdot d\vec A = \int_\Omega \vec\nabla\cdot(\varrho
  \vec v)\,dx.\]
Of course, since the total amount of fluid is conserved, this is also the change in the
amount of fluid in $\Omega$ over time:
\[\int_\Omega \vec\nabla\cdot(\varrho\vec v)\,dx = -\frac{d}{dt}\int_\Omega \varrho\,dx
  = -\int_\Omega \frac{\partial\varrho}{\partial t}\,dx.\]
Since this holds for any $\Omega$, the \newterm{continuity equation}
\germanfor{Kontinuit\"atsgleichung} follows:
\[\vec\nabla\cdot(\varrho\vec v) + \frac{\partial\varrho}{\partial t} = 0.\]


% --- exercises relocated here from the dissolved sheets ---

% Originally: content/week-11.tex, \exercisesheet{11}, problem 11.5
% Source: exercises/Ex11_Analysis2_eng.pdf

\begin{exercise}[11.5 --- Radial vector fields \textnormal{(important)}]
\label{ex:11.5}
For $\alpha\in\mathbb{R}$ consider the vector field
\[V_\alpha(x) = \lvert x\rvert^\alpha x, \qquad x\in\mathbb{R}^n\setminus\{0\},\]
where $\lvert x\rvert$ is the usual Euclidean norm.
\begin{enumerate}[label=\textbf{\arabic*.}]
  \item Compute $\divg V_\alpha(x)$ for all $x\neq0$, and show that it vanishes identically
  for $\alpha=-n$.
  \item Using the Divergence theorem on $V_0$, show the following relation between the
  $n$-volume of a sphere and the $n-1$ volume of its boundary:
  \[r\,\vol_{n-1}(\partial B_r) = n\,\vol_n(B_r), \qquad \forall r>0.\]
  \item Show that
  \[\int_{\partial B_1} V_{-n}\cdot\nu\,d\vol_{n-1} > 0, \qquad \int_{B_1} \divg V_{-n}(x)\,dx
    = 0.\]
  Why doesn't this contradict the Divergence Theorem?
\end{enumerate}
\end{exercise}

\begin{aiexercise}[Flux of a cubic field through the sphere]
% Generator: Claude Opus 5 (High)
\label{ex:ai_flux_cubic_sphere}
Let $F(x,y,z) := (x^3,\,y^3,\,z^3)$ and let $\Omega := B_1(0) \subseteq \mathbb{R}^3$. Compute
the flux of $F$ out of $\partial\Omega$.
\begin{enumerate}[label=\textbf{\arabic*.}]
  \item Say in one sentence what the integrand of \cref{def:flux} would look like here if the
    sphere were parametrized directly, and why that route is unattractive.
  \item Compute $\divg F$, and observe that it depends on the point only through its distance
    $r$ from the origin.
  \item Evaluate the volume integral. The shell reading of polar coordinates in
    \cref{rem:polar_as_shells} does it without a Jacobian determinant.
\end{enumerate}
\end{aiexercise}

\begin{aiexercise}[An open surface has to be closed first]
% Generator: Claude Opus 5 (High)
\label{ex:ai_capping_a_paraboloid}
Let $S := \{(x,y,z) \in \mathbb{R}^3 : z = 1-x^2-y^2,\ z \geq 0\}$ be the piece of the
paraboloid lying above the $xy$-plane, oriented by the upward-pointing normal, and let
$F(x,y,z) := (x,\,y,\,z)$.
\begin{enumerate}[label=\textbf{\arabic*.}]
  \item \Cref{thm:gauss_divergence} applies to the boundary of a domain, and $S$ is not one.
    Which set $D$ must be adjoined to $S$ so that $S \cup D = \partial\Omega$ for a bounded
    domain $\Omega$, and what is the exterior unit normal along $D$?
  \item Compute $\int_{\partial\Omega}\langle F,\nu\rangle\,d\vol_2$ using the divergence
    theorem. You will need $\vol_3(\Omega)$, which is a one-line computation in cylindrical
    coordinates.
  \item Compute the flux through $D$ alone, directly from \cref{def:flux}, and deduce the flux
    of $F$ through $S$.
\end{enumerate}
\end{aiexercise}

% Originally: content/week-11.tex, \section{Solutions}

\section{Solutions}
\label{sec:flux_solutions}

\begin{exercisesolution}[Solution of \cref{ex:11.5}]
\textbf{(1)} Writing $V_\alpha^i(x) = \lvert x\rvert^\alpha x_i$,
\[
\partial_i\big(\lvert x\rvert^\alpha x_i\big) = \alpha\lvert x\rvert^{\alpha-2}x_i^2
  + \lvert x\rvert^\alpha \qquad \text{(no sum),}
\]
so, summing over $i=1,\ldots,n$,
\[
\divg V_\alpha(x) = \sum_{i=1}^n \Big(\alpha\lvert x\rvert^{\alpha-2}x_i^2 + \lvert
  x\rvert^\alpha\Big) = \alpha\lvert x\rvert^{\alpha-2}\lvert x\rvert^2 + n\lvert x\rvert^\alpha
  = (\alpha+n)\lvert x\rvert^\alpha.
\]
Since $\lvert x\rvert^\alpha \neq 0$ for $x\neq0$, this vanishes identically exactly when
$\alpha = -n$.

\textbf{(2)} For $\alpha=0$, $V_0(x)=x$ and, by part (1), $\divg V_0 \equiv n$. Applying
\cref{thm:gauss_divergence} on $B_r$,
\[\int_{B_r} \divg V_0\,dx = n\,\vol_n(B_r).\]
On $\partial B_r$ the outward unit normal is $\nu(x) = x/r$, so $V_0(x)\cdot\nu(x) = x\cdot
x/r = \lvert x\rvert^2/r = r$ (using $\lvert x\rvert=r$ on $\partial B_r$), giving
\[\int_{\partial B_r} V_0\cdot\nu\,d\vol_{n-1} = r\,\vol_{n-1}(\partial B_r).\]
Equating the two expressions of $\int_{B_r}\divg V_0\,dx = \int_{\partial B_r}V_0\cdot\nu\,
d\vol_{n-1}$ gives $r\,\vol_{n-1}(\partial B_r) = n\,\vol_n(B_r)$.

\textbf{(3)} By part (1), $\divg V_{-n}(x) = 0$ for every $x\neq0$, so $\int_{B_1}\divg
V_{-n}(x)\,dx = 0$ trivially (the integrand vanishes on $B_1\setminus\{0\}$, a set of full
measure in $B_1$). On $\partial B_1$, $\nu(x)=x$ and $\lvert x\rvert=1$, so $V_{-n}(x) =
\lvert x\rvert^{-n}x = x$, giving
\[\int_{\partial B_1} V_{-n}\cdot\nu\,d\vol_{n-1} = \int_{\partial B_1} x\cdot x\,d\vol_{n-1}
  = \int_{\partial B_1} 1\,d\vol_{n-1} = \vol_{n-1}(\partial B_1) > 0.\]
This does not contradict \cref{thm:gauss_divergence}: the theorem requires $F$ to be $C^1$ on
all of $\overline{B_1}$, but $V_{-n}(x) = \lvert x\rvert^{-n}x$ is not even defined, let alone
continuously differentiable, at $x=0 \in \overline{B_1}$. Its hypotheses simply do not hold on
this domain.
\end{exercisesolution}

% Verified: solved independently, all three parts match Sol11 exactly.

\begin{exercisesolution}[Solution to \cref{ex:ai_flux_cubic_sphere}]
% Generator: Claude Opus 5 (High)
\textbf{1.} On the unit sphere the exterior unit normal is $\nu(x,y,z) = (x,y,z)$, so the
integrand of \cref{def:flux} is $\langle F,\nu\rangle = x^4+y^4+z^4$. In the spherical
parametrization that becomes
$\sin^4\theta\cos^4\varphi + \sin^4\theta\sin^4\varphi + \cos^4\theta$, multiplied by the surface
element $\sin\theta$: a fourth-power trigonometric integral in two variables, where the
divergence theorem offers a radial integral in one.

\textbf{2.} Writing $r$ for the distance to the origin,
$\divg F = 3x^2+3y^2+3z^2 = 3r^2$, which is radial.

\textbf{3.} Because the integrand is radial, \cref{rem:polar_as_shells} applies with $n = 3$ and
$C_3 = 4\pi$:
\[\int_{\partial\Omega}\langle F,\nu\rangle\,d\vol_2 = \int_{B_1(0)} 3r^2\,d\vol_3
  = 4\pi\int_0^1 3r^2\cdot r^2\,dr = 12\pi\int_0^1 r^4\,dr = \frac{12\pi}{5}.\]

\begin{remark}[The direct route, for once, is also available]
\label{rem:cubic_flux_direct_check}
The surface integral abandoned in part 1 is not actually intractable, and working it out is a
worthwhile check. By the symmetry of the sphere under permuting the coordinates, the three terms
of $x^4+y^4+z^4$ contribute equally, so
\[\int_{\mathbb{S}^2}\big(x^4+y^4+z^4\big)\,d\vol_2 = 3\int_{\mathbb{S}^2} z^4\,d\vol_2
  = 3\int_0^{2\pi}\!\!\int_0^\pi \cos^4\theta\,\sin\theta\,d\theta\,d\varphi
  = 6\pi\left[-\frac{\cos^5\theta}{5}\right]_0^\pi = \frac{12\pi}{5},\]
in agreement. Note that what made this bearable was a symmetry argument, applied \emph{before}
any trigonometry: expanding all three fourth powers first is the version that is unpleasant.
\end{remark}
\end{exercisesolution}

\begin{exercisesolution}[Solution to \cref{ex:ai_capping_a_paraboloid}]
% Generator: Claude Opus 5 (High)
\textbf{1.} Take $\Omega := \{(x,y,z) : x^2+y^2 < 1,\ 0 < z < 1-x^2-y^2\}$, the solid region
between the paraboloid and the plane $z = 0$. Its boundary is $S$ together with the flat disk
\[D := \{(x,y,0) : x^2+y^2 \leq 1\},\]
whose exterior unit normal is $\nu \equiv -e_3$, since $\Omega$ lies above it. On $S$ the
exterior normal is the upward-pointing one, which is the orientation the exercise already fixes,
so no sign has to be corrected later.

\textbf{2.} Here $\divg F = 1+1+1 = 3$, so the flux out of $\partial\Omega$ is three times a
volume. In cylindrical coordinates,
\[\vol_3(\Omega) = \int_0^{2\pi}\!\!\int_0^1\!\!\int_0^{1-r^2} r\,dz\,dr\,d\varphi
  = 2\pi\int_0^1 \big(r - r^3\big)\,dr = 2\pi\left(\frac12 - \frac14\right) = \frac{\pi}{2},\]
and therefore
\[\int_{\partial\Omega}\langle F,\nu\rangle\,d\vol_2 = 3\,\vol_3(\Omega) = \frac{3\pi}{2}.\]

\textbf{3.} Along $D$ the third coordinate is $z = 0$, so
$\langle F,\nu\rangle = \langle (x,y,0),\,-e_3\rangle = -z = 0$ at every point: the field is
tangent to the cap, and contributes nothing. Splitting the boundary integral,
\[\int_S \langle F,\nu\rangle\,d\vol_2
  = \int_{\partial\Omega}\langle F,\nu\rangle\,d\vol_2 - \int_D \langle F,\nu\rangle\,d\vol_2
  = \frac{3\pi}{2} - 0 = \frac{3\pi}{2}.\]

\begin{remark}[The cap vanished because of this field, not because it is a cap]
\label{rem:cap_contribution_is_not_generally_zero}
It is tempting to read part 3 as licence to ignore the lid. Change $F$ to
$\widetilde F(x,y,z) := (x,\,y,\,z+1)$ and nothing about the first two parts moves, since
$\divg\widetilde F = 3$ as before and the closed flux is still $3\pi/2$. But now
$\langle \widetilde F,\nu\rangle = -1$ on $D$, so the cap carries a flux of $-\pi$, and the flux
through the paraboloid is $3\pi/2 + \pi = 5\pi/2$ instead. The subtraction in part 3 is doing
real work; it merely had nothing to subtract.

\Cref{ex:11.1} is the same technique on a hemisphere, and the two are worth reading as a pair:
there the cap contributes $\pi/2$ and must be subtracted, here it contributes nothing. The
method is identical and only the arithmetic differs, which is the point.
\end{remark}
\end{exercisesolution}


% ============================================================
% Chapter 24 --- Alternating forms, the wedge product and the exterior derivative
% Originally: content/week-11.tex, \section{Antisymmetric k-linear forms and the wedge
%             product} and \section{Differential forms and the exterior derivative}
%             (Corsin Week 11).
% ============================================================
\chapter{Alternating Forms, the Wedge Product and the Exterior Derivative}
\label{ch:differential_forms}

% Transition: Claude Opus 5 (Medium)
Gradient, divergence and curl look like three unrelated operators that happen to satisfy two
unrelated identities ($\curl\grad = 0$ and $\divg\curl = 0$). This chapter builds the
machinery that makes them one operator satisfying one identity. Antisymmetric $k$-linear forms
and the wedge product supply the objects; the exterior derivative $d$ supplies the operator; and
$d \circ d = 0$ supplies the identity.

% Originally: content/week-11.tex, \section{Antisymmetric k-linear forms and the wedge product} (Corsin Week 11)

\section{Antisymmetric \texorpdfstring{$k$}{k}-linear forms and the wedge product}
\label{sec:antisymmetric_forms}

% Source: Corsin Nick/Class Notes/Week 11.pdf, p. 4
\begin{ainote}[Motivation]
Suppose one studies a submanifold $\mathcal{M} \subseteq \mathbb{R}^n$ via a parametrization
$\phi : D \to \mathcal{M}$, and integrates some quantity over $\mathcal{M}$ (for instance the
Gauss curvature, as in the Gauss--Bonnet theorem). The natural question is whether the
resulting number is a property of $\phi$, or a property of $\mathcal{M}$ itself: if it is
invariant under reparametrization, it must be the latter. \newterm{Differential forms} are
built to satisfy exactly this property. For a diffeomorphism $\Psi : V \to D$ with
$\det D\Psi > 0$ and a form $\omega$,
\begin{align}
\int_{\mathcal{M}} \omega &= \int_D \omega_{\phi(x)}(D\phi_x)\,dx \nonumber \\
  &\overset{\text{change of var.}}{=} \int_V \omega_{\phi(\Psi(y))}(D\phi_{\Psi(y)})
  \lvert\det D\Psi_y\rvert\,dy \nonumber \\
  &= \int_V \omega_{(\phi\circ\Psi)(y)}(D(\phi\circ\Psi)_y)\,dy = \int_{\mathcal{M}} \omega,
  \nonumber
\end{align}
where the middle equality uses the chain rule $D(\phi\circ\Psi)_y = D\phi_{\Psi(y)}\,D\Psi_y$
together with multilinearity and $\det D\Psi_y = \lvert\det D\Psi_y\rvert$. This calculation
is precisely why $\omega$ must be built from multilinear, antisymmetric maps: multilinearity
lets the Jacobian $D\Psi_y$ factor out as $\det D\Psi_y$, matching the $\lvert\det
D\Psi_y\rvert$ produced by the change-of-variables formula, exactly when $\det D\Psi_y > 0$.
\end{ainote}

% Source: Corsin Nick/Class Notes/Week 11.pdf, p. 5
\begin{definition}[Antisymmetric $k$-linear form]
\label{def:antisymmetric_k_linear_form}
Let $V$ be an $n$-dimensional real vector space and let $L : V^k \to \mathbb{R}$. Then $L$ is
an \newterm{antisymmetric $k$-linear form} \germanfor{antisymmetrische $k$-lineare Form}
if
\begin{enumerate}[label=\textbf{(\alph*)}]
  \item \textbf{($k$-linear):} $L$ is linear in every entry:
  \[L(v_1,\ldots,\alpha v_i+\beta w_i,\ldots,v_k) = \alpha L(v_1,\ldots,v_i,\ldots,v_k)
    + \beta L(v_1,\ldots,w_i,\ldots,v_k);\]
  \item \textbf{(antisymmetric):} swapping any two entries negates the value:
  \[L(v_1,\ldots,v_i,\ldots,v_j,\ldots,v_k) = -L(v_1,\ldots,v_j,\ldots,v_i,\ldots,v_k).\]
\end{enumerate}
\end{definition}

\begin{remark}
For $k=n$ and $V = \mathbb{R}^n$, the normalization $L(e_1,\ldots,e_n) = 1$ pins down a
\emph{unique} such $L$, namely $L = \det$.
\end{remark}

\begin{notation}
For $v \in V$ and a fixed basis $(e_1,\ldots,e_n)$ of $V$, superscripts index the coordinates
of a single vector, $v = \sum_{i=1}^n v^i e_i = (v^1,\ldots,v^n)$, while subscripts index a
list of vectors $v_1,\ldots,v_k \in V$. This distinguishes \qt{the $i$-th coordinate of
$v$} from \qt{the $i$-th vector in the list}. From here on, $V = \mathbb{R}^n$ with the
standard basis.
\end{notation}

% Source: Corsin Nick/Class Notes/Week 11.pdf, p. 6
\begin{aiexample}[Two families of examples]
% Generator: Claude Sonnet 5 (Medium)
\begin{enumerate}[label=\textbf{\arabic*.}]
  \item Fix $A \in \mathrm{Mat}_{k\times n}(\mathbb{R})$. For $v_1,\ldots,v_k \in
  \mathbb{R}^n$, write $V := (v_1\mid\cdots\mid v_k) \in \mathrm{Mat}_{n\times k}(\mathbb{R})$
  for the matrix with these columns; then $AV \in \mathrm{Mat}_{k\times k}(\mathbb{R})$, and
  $L(v_1,\ldots,v_k) := \det(AV)$ is an antisymmetric $k$-linear form, since $\det$ is
  multilinear and antisymmetric in the columns of its argument, and $v \mapsto Av$ is linear.
  \item For $k=1$, $L : \mathbb{R}^n \to \mathbb{R}$ is a \newterm{covector},
  $L \in (\mathbb{R}^n)^*$. The standard basis $(e_1^*,\ldots,e_n^*)$ of $(\mathbb{R}^n)^*$,
  satisfying $e_i^*(e_j) = \delta_{ij}$, consists of $1$-forms; in particular
  $e_i^*(v) = v^i$.
\end{enumerate}
\end{aiexample}

\begin{remark}[An equivalent characterization]
In the lecture, antisymmetry was replaced by the equivalent condition
\[L(VA) = L(V)\det(A), \qquad V := (v_1,\ldots,v_k) \in \mathrm{Mat}_{n\times k}(\mathbb{R}),
  \quad A \in \mathrm{Mat}_{k\times k}(\mathbb{R}),\]
writing $L(V)$ for $L(v_1,\ldots,v_k)$. We verify the equivalence in the special case $n=3$,
$k=2$: let $V = (v\mid w) \in \mathrm{Mat}_{3\times2}(\mathbb{R})$ and $A = \begin{pmatrix} a &
b \\ c & d \end{pmatrix}$, so that $VA = (av+cw \mid bv+dw)$. Rewriting $L(VA)$ in bilinear
form and using $L(v,v) = -L(v,v) = 0$ (from antisymmetry) together with $L(w,v) = -L(v,w)$,
\begin{align}
L(VA) &= L(av+cw,\,bv+dw) \nonumber \\
  &= ab\,L(v,v) + ad\,L(v,w) + cb\,L(w,v) + cd\,L(w,w) \nonumber \\
  &= (ad-bc)\,L(v,w) = L(V)\det(A). \nonumber
\end{align}
\end{remark}

% Source: Corsin Nick/Class Notes/Week 11.pdf, p. 7
\begin{definition}[Wedge product]
\label{def:wedge_product}
For a $k$-form $\alpha$ and an $m$-form $\beta$ on $\mathbb{R}^n$, the \newterm{wedge product}
\germanfor{Dachprodukt} $\alpha \wedge \beta$ is the $(k+m)$-form
\[
(\alpha\wedge\beta)(v_1,\ldots,v_k,v_{k+1},\ldots,v_{k+m})
  = \sum_{\sigma \in S_{k,m}} \sgn(\sigma)\,\alpha(v_{\sigma(1)},\ldots,v_{\sigma(k)})\,
    \beta(v_{\sigma(k+1)},\ldots,v_{\sigma(k+m)}),
\]
where $S_{k,m}$ is the set of permutations $\sigma$ of $\{1,\ldots,k+m\}$ satisfying
$\sigma(1) < \cdots < \sigma(k)$ and $\sigma(k+1) < \cdots < \sigma(k+m)$ (the
\newterm{$(k,m)$-shuffles}).
\end{definition}

\begin{ainote}
Corsin calls this formula \qt{not terribly important for our purposes} and instead has us work
directly with the five properties below.
\end{ainote}

The following properties of $\wedge$ are the ones actually used:
\begin{enumerate}[label=\textbf{(\arabic*)}]
  \item $\wedge$ is bilinear.
  \item For $\alpha$ a $k$-form and $\beta$ an $m$-form, $\alpha\wedge\beta$ is a $(k+m)$-form.
  \item $(\alpha\wedge\beta)\wedge\gamma = \alpha\wedge(\beta\wedge\gamma) = \alpha\wedge\beta
  \wedge\gamma$ (associativity).
  \item For $1$-forms $L_1,\ldots,L_k$ and vectors $v_1,\ldots,v_k \in \mathbb{R}^n$,
  \[L_1\wedge L_2\wedge\cdots\wedge L_k(v_1,\ldots,v_k) = \det\big(L_i(v_j)\big)_{1\leq i,j
  \leq k}.\]
  \item For $\alpha(T) := \det(AT)$, $\beta(S) := \det(BS)$ as in \cref{def:antisymmetric_k_linear_form}
  example 1,
  \[\alpha\wedge\beta(X) = \det\!\left(\begin{bmatrix} A \\ B \end{bmatrix} X\right).\]
\end{enumerate}

% Source: Corsin Nick/Class Notes/Week 11.pdf, p. 8
\begin{example}
By property (4), for $1 \leq i_1 < \cdots < i_k \leq n$,
\[
e_{i_1}^*\wedge\cdots\wedge e_{i_k}^*(v_1,\ldots,v_k) = \det\!\left(
  \begin{pmatrix} v_1^{i_1} & \cdots & v_k^{i_1} \\ \vdots & \ddots & \vdots \\ v_1^{i_k} &
  \cdots & v_k^{i_k} \end{pmatrix}\right)
  = \det\!\left(\begin{bmatrix} e_{i_1}\transp \\ \vdots \\ e_{i_k}\transp \end{bmatrix}
  \begin{bmatrix} v_1 & \cdots & v_k \end{bmatrix}\right).
\]
\end{example}

% Source: Corsin Nick/Class Notes/Week 11.pdf, pp. 9--10
\begin{proposition}[Every $k$-linear form is a combination of wedge products of covectors]
\label{prop:k_form_basis_expansion}
Every antisymmetric $k$-linear form $L$ on $\mathbb{R}^n$ can be written as
\[L = \sum_{1\leq i_1<\cdots<i_k\leq n} \ell_{i_1\cdots i_k}\,e_{i_1}^*\wedge\cdots\wedge
  e_{i_k}^*\]
for uniquely determined scalars $\ell_{i_1\cdots i_k} \in \mathbb{R}$.
\end{proposition}

\begin{proof}
Expand each argument $v_i = \sum_{m=1}^n v_i^m e_m$ and use $k$-linearity:
\[L(v_1,\ldots,v_k) = L\!\left(\sum_{i_1=1}^n v_1^{i_1}e_{i_1},\ldots,\sum_{i_k=1}^n
  v_k^{i_k}e_{i_k}\right) = \sum_{1\leq i_1,\ldots,i_k\leq n} L(e_{i_1},\ldots,e_{i_k})\,
  v_1^{i_1}\cdots v_k^{i_k}.\]
By antisymmetry, $L(e_{i_1},\ldots,e_{i_k}) = 0$ whenever two indices coincide, so the sum may
be restricted to $i_1,\ldots,i_k$ pairwise distinct; regrouping each such tuple by the
permutation $\sigma$ that sorts it into increasing order $i_{\sigma(1)}<\cdots<i_{\sigma(k)}$,
\begin{align}
L(v_1,\ldots,v_k) &= \sum_{1\leq i_1<\cdots<i_k\leq n} L(e_{i_1},\ldots,e_{i_k})
  \sum_{\sigma\in S_k} \sgn(\sigma)\,v_1^{i_{\sigma(1)}}\cdots v_k^{i_{\sigma(k)}} \nonumber \\
  &= \sum_{1\leq i_1<\cdots<i_k\leq n} \ell_{i_1\cdots i_k}\,e_{i_1}^*\wedge\cdots\wedge
  e_{i_k}^*(v_1,\ldots,v_k), \nonumber
\end{align}
with $\ell_{i_1\cdots i_k} := L(e_{i_1},\ldots,e_{i_k})$, using that the inner sum over $\sigma$
is exactly the determinant expansion of $e_{i_1}^*\wedge\cdots\wedge e_{i_k}^*(v_1,\ldots,v_k)$
from property (4) above.
\end{proof}

% Originally: content/week-11.tex, \section{Differential forms and the exterior derivative} (Corsin Week 11)

\section{Differential forms and the exterior derivative}
\label{sec:differential_forms}

% Source: Corsin Nick/Class Notes/Week 11.pdf, p. 10
\begin{definition}[Differential form]
\label{def:differential_p_form}
A \newterm{differential $p$-form} $\omega$ \germanfor{Differentialform} is defined on a
domain $D \subseteq \mathbb{R}^n$ if, for every $x \in D$, $\omega(x) = \omega_x : (\mathbb{R}^n)^p
\to \mathbb{R}$ is an antisymmetric $p$-linear map. Since $(e_1(x),\ldots,e_n(x))$ is a basis
of $\mathbb{R}^n$ for every $x \in D$, \cref{prop:k_form_basis_expansion} lets us write
\[\omega(x) = \sum_{1\leq i_1<\cdots<i_p\leq n} \omega_{i_1\cdots i_p}(x)\,e_{i_1}^*(x)\wedge
  \cdots\wedge e_{i_p}^*(x)\]
for functions $\omega_{i_1\cdots i_p} : \mathbb{R}^n \to \mathbb{R}$.
\end{definition}

% Source: Corsin Nick/Class Notes/Week 11.pdf, p. 11
\begin{notation}
We abbreviate $e_{i_1}^*\wedge\cdots\wedge e_{i_p}^* =: e_I^*$ for $I = (i_1,\ldots,i_p)$, and
write $\Omega^p(U)$ for the set of $p$-forms on $U \subseteq \mathbb{R}^n$. By convention,
$\Omega^0(U) = C^\infty(U)$.
\end{notation}

\begin{example}
Let $U \subseteq \mathbb{R}^n$.
\begin{enumerate}[label=\textbf{(\alph*)}]
  \item Any $f \in C^\infty(U)$ is a $0$-form.
  \item A vector field $F : U \to \mathbb{R}^n$ has an associated $1$-form $\omega_F(x) : v
  \mapsto \langle F(x),v\rangle$, that is,
  \[\omega_F(x)(v) = \sum_{i=1}^n F^i(x)\,v^i = \sum_{i=1}^n F^i(x)\,e_i^*(x)(v)
    \qquad\Longrightarrow\qquad \omega_F = \sum_{i=1}^n F^i\,e_i^*.\]
% Source: Corsin Nick/Class Notes/Week 11.pdf, p. 12
  \item $F : U \to \mathbb{R}^n$ also naturally induces an $(n-1)$-form
  \[\omega_F^{n-1}(x)(v_1,\ldots,v_{n-1}) := \det\big(F(x)\mid v_1\mid\cdots\mid v_{n-1}\big).\]
  Developing this determinant along the first column,
  \begin{align}
  \omega_F^{n-1}(x)(v_1,\ldots,v_{n-1}) &= \sum_{i=1}^n (-1)^{i-1}F^i(x)\det\!\left(
    \begin{smallmatrix} v_1^1 & \cdots & v_{n-1}^1 \\ \vdots & & \vdots \\ v_1^{i-1} & \cdots &
    v_{n-1}^{i-1} \\ v_1^{i+1} & \cdots & v_{n-1}^{i+1} \\ \vdots & & \vdots \\ v_1^n &
    \cdots & v_{n-1}^n \end{smallmatrix}\right) \nonumber \\
    &= \sum_{i=1}^n (-1)^{i-1}F^i(x)\,e_1^*\wedge\cdots\wedge\widehat{e_i^*}\wedge\cdots\wedge
    e_n^*(v_1,\ldots,v_{n-1}), \nonumber
  \end{align}
  where $\widehat{e_i^*}$ denotes an omitted factor, so that $\omega_F^{n-1} =
  \sum_{i=1}^n (-1)^{i-1}F^i\,e_1^*\wedge\cdots\wedge\widehat{e_i^*}\wedge\cdots\wedge e_n^*$.
  \item In particular, for $F : \mathbb{R}^3 \to \mathbb{R}^3$,
  \[\omega_F^2(x)(v_1,v_2) = \langle F(x),v_1\times v_2\rangle = \det\big(F(x)\mid v_1\mid
    v_2\big).\]
\end{enumerate}
\end{example}

\begin{ainote}
Item \textbf{(d)} identifies $\omega_F^2$ with the flux integrand from
\cref{sec:flux_divergence_theorem}: the flux $\int_S \langle F,\nu\rangle\,d\vol_2$ over a parametrized surface is
exactly the integral, in the sense of the motivating change-of-variables calculation above, of
the $2$-form $\omega_F^2$.
\end{ainote}

% Source: Corsin Nick/Class Notes/Week 11.pdf, p. 13
\begin{theorem}[Exterior derivative]
\label{thm:exterior_derivative}
There exists a unique linear operator $d : \Omega^k(U) \to \Omega^{k+1}(U)$ satisfying
\begin{enumerate}[label=\textbf{(\alph*)}]
  \item for $f \in C^\infty(U) = \Omega^0(U)$: $df = Df$, the ordinary differential;
  \item for $\omega \in \Omega^k(U)$, $\theta \in \Omega^s(U)$: $d(\omega\wedge\theta) =
  d\omega\wedge\theta + (-1)^k\,\omega\wedge d\theta$;
  \item $d\circ d\,\omega = d^2\omega = 0$.
\end{enumerate}
\end{theorem}

\begin{ainote}
Corsin states this as a fact seen in the lecture, without proof, and moves directly to its
consequences. The proof can be found in the official script
(Analysis\_II\_Script\_v1.pdf, §14.3, pp.\ 136--138), where the operator is constructed
explicitly in coordinates and properties \textbf{(a)}--\textbf{(c)} are verified.
\end{ainote}

Denote by $dx^i$ the differential of $x \mapsto x^i$ in a moving frame $(e_1(x),\ldots,e_n(x))$,
that is, $dx^i = e_i^*(x)$. Every $p$-form $\omega$ on $U \subseteq \mathbb{R}^n$ then takes
the shape
\[\omega = \sum_{1\leq i_1<\cdots<i_p\leq n} \omega_{i_1\cdots i_p}\,dx^{i_1}\wedge\cdots\wedge
  dx^{i_p} = \sum_I \omega_I\,dx^I, \qquad \omega_I \in C^\infty(U),\]
with exterior derivative
\[d\omega = \sum_I d\omega_I \wedge dx^I.\]

% Source: Corsin Nick/Class Notes/Week 11.pdf, pp. 14--15
\begin{example}[The exterior derivative recovers the gradient and the curl]
\label{ex:exterior_derivative_gradient_curl}
\begin{enumerate}[label=\textbf{(\alph*)}]
  \item For $f \in C^\infty(\mathbb{R}^n) = \Omega^0(\mathbb{R}^n)$,
  \[df_x(v) = \left(\frac{\partial f}{\partial x^1}\ \Big|\ \cdots\ \Big|\ \frac{\partial f}
    {\partial x^n}\right)\!\begin{pmatrix}v^1\\ \vdots \\ v^n\end{pmatrix}
    = \frac{\partial f}{\partial x^1}\,dx^1(v) + \cdots + \frac{\partial f}{\partial x^n}\,
    dx^n(v),\]
  so $df = \frac{\partial f}{\partial x^1}\,dx^1 + \cdots + \frac{\partial f}{\partial x^n}\,
  dx^n = \langle \nabla f,\cdot\rangle$: the exterior derivative of a $0$-form is exactly the
  $1$-form associated (as in the example above) to the gradient vector field $\nabla f$. This
  is the shape that makes Stokes' theorem specialize to the FTC, as in
  \cref{sec:gauss_intuition}.
  \item Define a $1$-form on $U \subseteq \mathbb{R}^3$ by $\omega(x,y,z) = \langle
  F(x,y,z),\cdot\rangle = F^1\,dx+F^2\,dy+F^3\,dz$ for a vector field $F$. By property
  \textbf{(b)} of
  \cref{thm:exterior_derivative} and $d(dx^i)=0$ (since $dx^i = d(x^i)$ and $d^2=0$),
  \[
  d\omega = dF^1\wedge dx + dF^2\wedge dy + dF^3\wedge dz,
  \]
  and expanding each $dF^i = \partial_x F^i\,dx+\partial_y F^i\,dy+\partial_z F^i\,dz$, only
  the cross terms survive the wedge (since $dx^i\wedge dx^i=0$ by antisymmetry):
  \[
  d\omega = \Big(\frac{\partial F^3}{\partial y}-\frac{\partial F^2}{\partial z}\Big)dy\wedge
    dz + \Big(\frac{\partial F^1}{\partial z}-\frac{\partial F^3}{\partial x}\Big)dz\wedge dx
    + \Big(\frac{\partial F^2}{\partial x}-\frac{\partial F^1}{\partial y}\Big)dx\wedge dy.
  \]
  Comparing with example item \textbf{(d)} in \cref{sec:differential_forms}, this is exactly
  $\omega_{\nabla\times F}^2 = \langle\nabla\times F,\cdot\times\cdot\rangle = \det(\nabla
  \times F\mid\cdot\mid\cdot)$: the exterior derivative of the $1$-form associated to $F$ is
  the $2$-form associated to $\curl F$.

  % Generator: Claude Sonnet 5 (High) --- new content: the analogous fact for the divergence,
  % completing the grad/curl/div trio as the 0-, 1- and (n-1)-form faces of $d$, using the
  % $(n-1)$-form $\omega_F^{n-1}$ already introduced in \cref{sec:differential_forms} (example
  % item (c)) rather than redefining it. Used below in
  % \cref{rem:one_theorem_four_costumes} to place the divergence theorem alongside Green,
  % classical Stokes and the FTC as instances of the single general Stokes theorem.
  \item Recall the $(n-1)$-form $\omega_F^{n-1} =
  \sum_{i=1}^n(-1)^{i-1}F^i\,dx^1\wedge\cdots\wedge\widehat{dx^i}\wedge\cdots\wedge dx^n$ from
  \cref{sec:differential_forms} (example item \textbf{(c)}), with $\widehat{dx^i}$ denoting an
  omitted factor. In $d(F^i\,dx^1\wedge\cdots\widehat{dx^i}\cdots\wedge dx^n) =
  dF^i\wedge dx^1\wedge\cdots\widehat{dx^i}\cdots\wedge dx^n$, only the $\partial_iF^i\,dx^i$
  summand of $dF^i = \sum_j\partial_jF^i\,dx^j$ survives, since every $j\neq i$ repeats a factor
  already present and so kills the wedge by antisymmetry. Therefore
  \[d\omega_F^{n-1} = \sum_{i=1}^n(-1)^{i-1}\partial_iF^i\;dx^i\wedge dx^1\wedge\cdots\wedge
    \widehat{dx^i}\wedge\cdots\wedge dx^n.\]
  Moving $dx^i$ past the $i-1$ factors $dx^1,\dots,dx^{i-1}$ ahead of it, into its natural
  position, costs another sign of $(-1)^{i-1}$, cancelling the one already there:
  \begin{align*}
    dx^i\wedge dx^1\wedge\cdots\wedge\widehat{dx^i}\wedge\cdots\wedge dx^n
      &= (-1)^{i-1}\,dx^1\wedge\cdots\wedge dx^n \\
    \Longrightarrow\quad
      d\omega_F^{n-1} &= \sum_{i=1}^n\partial_iF^i\;dx^1\wedge\cdots\wedge dx^n
      = (\divg F)\,dx^1\wedge\cdots\wedge dx^n.
  \end{align*}
  So, the exterior derivative of $\omega_F^{n-1}$ is $\divg F$ times the volume form. That
  completes the pattern of items \textbf{(a)} and \textbf{(b)} above: $\grad$, $\curl$ and
  $\divg$ are exactly the $0$-, $1$- and $(n-1)$-form faces of the single operator $d$.

% Source: Jérôme Paschoud/Notizen/Woche 12.1 Differentialformen.pdf, p. 6
  \item Let $\omega = x_3 \,dx_3 + x_1 x_2 \,dx_2$ be a $1$-form on $\mathbb{R}^3$. Let us compute its exterior derivative using the rules:
  \begin{align*}
  d\omega &= d(x_3) \wedge dx_3 + d(x_1 x_2) \wedge dx_2 \\
  &= dx_3 \wedge dx_3 + (x_2 \,dx_1 + x_1 \,dx_2) \wedge dx_2 \\
  &= 0 + x_2 \,dx_1 \wedge dx_2 + x_1 \underbrace{dx_2 \wedge dx_2}_{=0} \\
  &= x_2 \,dx_1 \wedge dx_2.
  \end{align*}
\end{enumerate}
\end{example}

\begin{ainote}
Corsin writes the middle expansion of $d\omega$ out in full, all nine terms, one for each pair
$(dF^i,dx^j)$, before collecting the six that survive into the three coefficient pairs above. It
is condensed here to the surviving terms directly. Nothing is lost, since $dx\wedge
dx=dy\wedge dy=dz\wedge dz=0$ kills exactly the three diagonal terms, and antisymmetry
($dy\wedge dx=-dx\wedge dy$, and so on) pairs up the rest.
\end{ainote}

% Generator: Claude Opus 5 (High) --- FIG-CH24-03. The chapter opens by promising that grad,
% curl and div are one operator satisfying one identity. Items (a)-(c) of the example above prove
% it, one column at a time, and the result was never assembled anywhere. This is that assembly.
% Every vertical arrow is an identification established earlier in this chapter, and each is
% labelled with the place it was made, so the diagram is checkable rather than decorative.
\begin{center}
\begin{tikzcd}[row sep=3.4em, column sep=3.4em]
  \Omega^0(U) \arrow[r, "d"]
    & \Omega^1(U) \arrow[r, "d"]
    & \Omega^2(U) \arrow[r, "d"]
    & \Omega^3(U) \\
  C^\infty(U) \arrow[r, "\grad"'] \arrow[u, "="]
    & C^\infty(U,\mathbb{R}^3) \arrow[r, "\curl"'] \arrow[u, "F \mapsto \omega_F"]
    & C^\infty(U,\mathbb{R}^3) \arrow[r, "\divg"'] \arrow[u, "F \mapsto \omega_F^2"]
    & C^\infty(U) \arrow[u, "g \mapsto g\,dx\wedge dy\wedge dz"']
\end{tikzcd}
\end{center}

\begin{importantremark}[One operator, one identity]
\label{rem:one_theorem_four_costumes_forms}
For $U \subseteq \mathbb{R}^3$ the diagram above commutes, and each square is one of the three
computations just performed: the left square is item \textbf{(a)}, the middle is item
\textbf{(b)}, the right is item \textbf{(c)}. Read downward it says that the classical vector
operators are $d$ in disguise, and that the disguise is nothing but the choice of which object to
call a vector field.

The payoff is the bottom row's two mysteries collapsing into one line of the top row. Property
\textbf{(c)} of \cref{thm:exterior_derivative} is $d \circ d = 0$, a single statement about a
single
operator, and reading it through the two identifications gives
\[\curl\grad = 0 \qquad\text{and}\qquad \divg\curl = 0.\]
These are the same fact twice, seen from two rungs of the same ladder. That is the whole reason
for the machinery: in $\mathbb{R}^3$ there happen to be two identities because there happen to be
two interior rungs, and in $\mathbb{R}^n$ there are $n-1$ of them, none of which needs a separate
proof.

Two cautions the picture invites. The vertical arrows are \emph{not} canonical: $F \mapsto
\omega_F$ uses the inner product and $F \mapsto \omega_F^2$ uses the determinant, so both depend
on $\mathbb{R}^3$ carrying a metric and an orientation, which is exactly what a general manifold
does not hand you for free. And the ladder is this short only because $n = 3$: $\Omega^4(U)$ on
$\mathbb{R}^3$ is zero, since there is no fourth coordinate to build a shadow on
(\cref{rem:forms_measure_shadows}), which is why $\divg$ is the last operator in the row and not
because anyone chose to stop there.
\end{importantremark}

\begin{aiexercise}[The angle form is closed]
% Generator: Claude Opus 5 (High)
\label{ex:ai_angle_form_closed}
On $U := \mathbb{R}^2\setminus\{0\}$ define the $1$-form
\[\omega := \frac{-y\,dx + x\,dy}{x^2+y^2}.\]
\begin{enumerate}[label=\textbf{(\alph*)}]
  \item Show that $d\omega = 0$.
  \item By part \textbf{(a)} and \cref{thm:exterior_derivative}, one might hope that
    $\omega = df$ for some
    $f \in C^\infty(U)$. Compute $\int_\gamma \omega$ along the unit circle
    $\gamma(t) := (\cos t, \sin t)$, $t \in [0,2\pi]$, and explain what this rules out.
    \emph{Hint:} \cref{prop:work_difference_potential}, read through the identification
    $F \leftrightarrow \omega_F$ of the diagram above.
\end{enumerate}
\end{aiexercise}

\begin{aiexercise}[Exterior derivatives by hand]
% Generator: Claude Opus 5 (High)
\label{ex:ai_exterior_derivative_drill}
Compute $d\omega$ for each of the following, on the largest domain where the expression makes
sense.
\begin{enumerate}[label=\textbf{(\alph*)}]
  \item $\omega = xy\,dx + yz\,dy + zx\,dz$ on $\mathbb{R}^3$. Which vector field's curl have you
    just computed?
  \item $\omega = x_1x_2x_3\,dx_2\wedge dx_3$ on $\mathbb{R}^3$.
  \item $\omega = f\,dg$, where $f,g \in C^\infty(\mathbb{R}^n)$. Express the answer using
    $\wedge$ and deduce that $d(f\,dg + g\,df) = 0$.
\end{enumerate}
\end{aiexercise}

% Source: DiffComp.pdf (Prof. Joaquim Serra, "Some differential form computations 1",
% Analysis II: Several Variables, ETH Zürich FS 2026, assignment 13 May 2026, not graded), pp. 1--2
% Extractor: Claude Opus 5 (High)
% This is our OWN lecturer's example sheet for the current course, not an old exam, and it was
% never cited anywhere in the repository before 2026-08-11. The six parts are Serra's, verbatim in
% content; only the numbering changes (his 1.1--1.6 become our (a)--(f), per the sub-part rule).
\begin{exercise}[Six wedge products $df \wedge \lambda$]
\label{ex:serra_df_wedge_lambda}
All forms below live on $\mathbb{R}^3$, and in every case $\lambda$ is a $1$-form. Compute
\[\omega := df \wedge \lambda\]
in simplified form. In parts \textbf{(a)}--\textbf{(c)} the symbol $f$ denotes a function, that is
a $0$-form, and the answer should be presented as
\[\omega = F\,dx\wedge dy + G\,dx\wedge dz + H\,dy\wedge dz;\]
in parts \textbf{(d)}--\textbf{(f)} the symbol $f$ denotes a $1$-form instead, and the answer
should be presented as
\[\omega = F\,dx\wedge dy\wedge dz.\]
\begin{enumerate}[label=\textbf{(\alph*)}]
  \item $f(x,y) = x^2y + \sin y$, \qquad $\lambda = x\,dx + y\,dy$.
  \item $f(x,y,z) = xyz$, \qquad $\lambda = x\,dy + y\,dz + z\,dx$.
  \item $f(x,y,z) = e^xy + z^2$, \qquad $\lambda = (x+y)\,dx + z\,dy + xz\,dz$.
  \item $f = x^2\,dy + y^2\,dz$, \qquad $\lambda = x\,dx + z\,dy$.
  \item $f = xy\,dx + z\,dy + x^2z\,dz$, \qquad $\lambda = y\,dx + x\,dy + z\,dz$.
  \item $f = e^y\,dx + xz\,dy + y^2z\,dz$, \qquad $\lambda = dx + y\,dy + z\,dz$.
\end{enumerate}
\exinfo{This exercise is the whole of \emph{Some differential form computations 1}, an ungraded
example sheet issued by Prof.\ Joaquim Serra on 13 May 2026 for Analysis II, Spring Semester 2026.
The sheet prints full worked solutions, which are the basis of the solutions given here.}
\end{exercise}

% Feedback: Gemini 3.1 Pro (2026-08-11) --- called a very strong intervention: warning the reader
% that Serra's $f$ is a $0$-form in the first half and a $1$-form in the second heads off immense
% confusion. Keep this note; do not "tidy" the exercise by renaming Serra's letters instead.
\begin{ainote}
Serra uses the letter $f$ for a $0$-form in \textbf{(a)}--\textbf{(c)} and for a $1$-form in
\textbf{(d)}--\textbf{(f)}, and the sheet says so explicitly rather than changing the name. Keep
the distinction in view while computing, because $df$ means two different things across the two
halves: the differential of a function in the first, and the exterior derivative of a $1$-form in
the second (\cref{thm:exterior_derivative}). The degree bookkeeping is what makes the answers land
in different spaces: $df\wedge\lambda$ is a $2$-form in \textbf{(a)}--\textbf{(c)} and a $3$-form
in \textbf{(d)}--\textbf{(f)}, which is why the requested output shapes differ.
\end{ainote}

% Supplement: Sascha Brack/Class Notes/Week_11_Notes_Friday.pdf, pp. 2--4


% ============================================================
% Chapter 25 --- Pullbacks, integration of forms and orientability
% Originally: content/week-12.tex in full (Corsin Week 12), together with
%             content/week-13.tex, \section{Submanifolds with boundary and induced
%             orientation} (Corsin Week 13), which belongs with orientability rather than
%             with the statement of Stokes' theorem it happened to follow.
% ============================================================
\chapter{Pullbacks, Integration of Forms and Orientability}
\label{ch:pullbacks}

% Transition: Claude Opus 5 (Medium)
To integrate a differential form over a submanifold one has to pull it back to a parameter
domain, and then check that the answer does not depend on which parametrization was chosen. It
does not, provided the parametrizations agree on a direction. That proviso is orientability, and
it is what this chapter is really about. The last section fixes the induced orientation on a boundary, which
is the convention that makes the sign in Stokes' theorem come out right.

% Originally: content/week-12.tex, \section{Pullback of forms} (Corsin Week 12)


\section{Pullback of forms}
\label{sec:pullback_of_forms}

% Source: Corsin Nick/Class Notes/Week 12.pdf, p. 5
\begin{definition}[Pullback of a form]
\label{def:pullback_of_forms}
Let $U \subseteq \mathbb{R}^m$, $V \subseteq \mathbb{R}^n$ be open, $F : U \to V$ a smooth
map, and $\omega \in \Omega^p(V)$ a differential $p$-form on $V$. The \newterm{pullback}
(\germanterm{Zur\"uckziehung}) of $\omega$ under $F$ is the differential $p$-form on $U$
defined by
\[(F^*\omega)_x(v_1,\ldots,v_p) := \omega_{F(x)}\big(DF_x(v_1),\ldots,DF_x(v_p)\big).\]
\end{definition}

\begin{ainote}
Corsin also draws a commutative diagram $TU \xrightarrow{dF} TV \xrightarrow{\omega} \mathbb{R}$
over $U \xrightarrow{F} V$, with tangent-bundle projections $\pi_U, \pi_V$, and labels it
\qt{not exam relevant}. Omitted here for that reason; the formula above is the operative
content.
\end{ainote}

The pullback satisfies two properties that make it the right tool for transporting forms
along smooth maps:
\begin{enumerate}[label=\textbf{(\arabic*)}]
  \item $F^*(\omega\wedge\theta) = F^*\omega \wedge F^*\theta$;
  \item $F^*(d\omega) = d(F^*\omega)$ (the pullback commutes with the exterior derivative).
\end{enumerate}

\begin{ainote}
This is precisely the operation used, informally, in the reparametrization-invariance
motivation for \cref{def:antisymmetric_k_linear_form}: writing $\phi^*\omega$ and
$(\phi\circ\Psi)^*\omega$ for the two sides of that calculation makes it read as $\int_D
\phi^*\omega = \int_V (\phi\circ\Psi)^*\omega = \int_V \Psi^*(\phi^*\omega)$, i.e.\ property (1)
applied to $\phi^*\omega$ and the fact that pullback along a composition is the composition of
pullbacks.
\end{ainote}

% Originally: content/week-12.tex, \section{Integration of differential forms on submanifolds} (Corsin Week 12)

\section{Integration of differential forms on submanifolds}
\label{sec:integration_of_forms}

% Source: Corsin Nick/Class Notes/Week 12.pdf, p. 6
If $\mathcal{M}$ is an $m$-dimensional submanifold of $\mathbb{R}^m$ --- i.e.\ simply an open
set --- and $\omega \in \Omega^m(\mathcal{M})$ an $m$-form, then
\[\int_{\mathcal{M}} \omega := \int_{\mathcal{M}} \omega_x(e_1,\ldots,e_m)\,dx;\]
in particular, for $f \in C^\infty(\mathcal{M})$,
\[\int_{\mathcal{M}} f\,dx^1\wedge\cdots\wedge dx^m = \int_{\mathcal{M}} f(x)\,dx,\]
recovering the ordinary Lebesgue integral of $f$.

\begin{definition}[Integral of a form over a parametrized submanifold]
\label{def:integral_of_form_over_submanifold}
If $L^k \subseteq \mathbb{R}^n$ is a $k$-dimensional submanifold, $n>k$, and $\phi : U \to L$
is a parametrization with $U \subseteq \mathbb{R}^k$ open, then the integral of $\omega \in
\Omega^k(L)$ is defined as
\[\int_L \omega := \int_U \phi^*\omega = \int_U \omega_{\phi(x)}\!\left(\frac{\partial\phi}
  {\partial x^1},\ldots,\frac{\partial\phi}{\partial x^k}\right) dx.\]
\end{definition}

In particular, for $f \in C^\infty(L)$ and $1\leq i_1<\cdots<i_k\leq n$,
\[\int_L f\,dx^{i_1}\wedge\cdots\wedge dx^{i_k} = \int_U f(x)\det\!\left(\frac{\partial
  \phi^{i_s}}{\partial x^t}\right)_{1\leq s,t\leq k} dx,\]
because, by property (4) of \cref{def:wedge_product},
\[dx^{i_1}\wedge\cdots\wedge dx^{i_k}(\partial_1\phi,\ldots,\partial_k\phi) = \det\!\left(
  \begin{smallmatrix} dx^{i_1}(\partial_1\phi) & \cdots & dx^{i_1}(\partial_k\phi) \\
  \vdots & \ddots & \vdots \\ dx^{i_k}(\partial_1\phi) & \cdots & dx^{i_k}(\partial_k\phi)
  \end{smallmatrix}\right) = \det\!\left(\frac{\partial\phi^{i_s}}{\partial x^t}\right)_{1\leq
  s,t\leq k}.\]

\begin{ainote}
Applying \cref{def:integral_of_form_over_submanifold} to $\omega_F^{n-1}$ or $\omega_F^2$ from
\cref{ch:differential_forms} recovers exactly the flux integrals of \cref{ch:flux} --- those
were, all along,
integrals of forms in this sense. And \cref{def:pullback_of_forms} is what makes the
motivating reparametrization-invariance calculation precise: $\int_L \omega$ does
not depend on the parametrization $\phi$ chosen, only on $L$ and its orientation.
\end{ainote}


% --- exercises relocated here from the dissolved sheets ---

% Originally: content/week-12.tex, \exercisesheet{12}, problem 12.9
% Source: exercises/Ex12_Analysis2_eng.pdf

\begin{exercise}[12.9 --- Change of variables with differential forms \textnormal{(important)}]
\label{ex:12.9}
Let $U,V \subseteq \mathbb{R}^n$ be open sets, and let $\psi = (\psi^1,\ldots,\psi^n) : U \to
V$ be a diffeomorphism with $\det\!\left(\frac{\partial\psi^j}{\partial x^i}\right) > 0$.
Define an $n$-differential form $\omega$ on $U$ by $\omega = d\psi^1\wedge d\psi^2\wedge\cdots
\wedge d\psi^n$.
\begin{enumerate}[label=\textbf{(\alph*)}]
  \item Show that $\omega = \det\!\left(\frac{\partial\psi^j}{\partial x^i}\right) dx^1\wedge
  \cdots\wedge dx^n$.
  \item From the relation in (a), rewrite the change of variables formula using differential
  forms.
\end{enumerate}
\end{exercise}

% Originally: content/week-12.tex, \section{Orientable manifolds} (Corsin Week 12)

\section{Orientable manifolds}
\label{sec:orientable_manifolds}

% Source: Corsin Nick/Class Notes/Week 12.pdf, p. 7
\begin{ainote}
The two definitions below are reproduced by Corsin directly from the lecture script (numbered
there as Definitions 14.68 and 14.65), rather than written out in his own hand; he adds the
margin annotations quoted and figured below.
\end{ainote}

\begin{definition}[Support of a $k$-form]
\label{def:support_of_k_form}
Let $U \subseteq \mathbb{R}^n$ be open and $\alpha \in \Omega^k(U)$. The \newterm{support}
(\germanterm{Tr\"ager}) of $\alpha$ is the closed set
\[\supp\alpha := \overline{\{x\in U : \alpha_x \neq 0\}} \subseteq \overline{U}.\]
We say $\alpha$ has \newterm{compact support} in $U$ if $\supp\alpha$ is bounded and contained
in $U$.
\end{definition}

\begin{definition}[Orientable manifold]
\label{def:orientable_manifold}
Let $M \subseteq \mathbb{R}^n$ be a $k$-dimensional $C^\infty$ submanifold. We say $M$ is
\newterm{orientable} (\germanterm{orientierbar}) if there exists a family of local
parametrizations $\phi_i : V_i \to M \cap U_i$, $i \in I$, whose images cover $M$, and such
that for every $i,j \in I$ with $M \cap U_i \cap U_j \neq \emptyset$, the transition map
$\widetilde\phi_i^{-1}\circ\widetilde\phi_j : \widetilde V_j \to \widetilde V_i$ has positive
Jacobian determinant:
\[\det D(\widetilde\phi_i^{-1}\circ\widetilde\phi_j) > 0 \qquad \text{on } \widetilde V_j.\]
Such a family $\{\phi_i\}$ is called an \newterm{atlas} of $M$; if it satisfies the
determinant condition above it is called an \newterm{orientation} of $M$. Once such a family
is chosen, $M$ is said to be \newterm{oriented}.
\end{definition}

\begin{ainote}
Corsin marks the family $\{\phi_i\}$ itself with the margin note \qt{is called an Atlas} next
to the covering condition, and the positivity condition with \qt{Atlas} again next to the
determinant inequality --- i.e.\ he is flagging that \qt{atlas} names the family regardless of
whether the determinant condition holds, while \qt{orientation} is the stronger notion once it
does. This matches the definition's own wording above.
\end{ainote}

Corsin illustrates the failure of the orientation condition with a hand-drawn counterexample:
two overlapping local parametrizations $\phi, \psi$ of the same patch of a surface, whose
coordinate frames $(\partial_1\phi,\partial_2\phi)$ and $(\partial_1\psi,\partial_2\psi)$ at a
shared point disagree in handedness.

% Sonnet 5 (Medium) --- FIG-W12-01
\begin{center}
\begin{tikzpicture}[>=Latex, scale=1.0]
  \draw[green!50!black, thick] (0,0) .. controls (1.5,0.9) and (3.5,0.6) .. (5,1.2)
    .. controls (5,2.4) and (3,2.6) .. (1.5,2.0) .. controls (0,1.6) and (0,0.8) .. (0,0);
  \draw[green!50!black, dotted] (0.4,1.0) .. controls (1.8,0.6) and (3.2,1.1) .. (4.4,1.6);
  \fill (2.1,1.3) circle (1.2pt);
  \draw[blue!70!black, thick, ->] (2.1,1.3) -- (3.3,1.55) node[right, font=\scriptsize]
    {$\partial_1\phi$};
  \draw[blue!70!black, thick, ->] (2.1,1.3) -- (1.7,2.15) node[above, font=\scriptsize]
    {$\partial_2\phi$};
  \draw[red, thick, ->] (2.1,1.3) -- (0.9,1.05) node[left, font=\scriptsize] {$\partial_1\psi$};
  \draw[red, thick, ->] (2.1,1.3) -- (2.55,0.35) node[below, font=\scriptsize] {$\partial_2\psi$};
  \node[font=\scriptsize] at (0.1,-0.35) {$\phi$};
  \node[font=\scriptsize] at (4.4,2.7) {$\psi$};
  \draw[blue!70!black, ->] (6.3,0.2) -- (7.3,0.2) node[right, font=\scriptsize] {$e_1$};
  \draw[blue!70!black, ->] (6.3,0.2) -- (6.3,1.2) node[above, font=\scriptsize] {$e_2$};
  \draw[red, ->] (8.6,0.2) -- (9.6,0.2) node[right, font=\scriptsize] {$e_1$};
  \draw[red, ->] (8.6,0.2) -- (8.6,1.2) node[above, font=\scriptsize] {$e_2$};
\end{tikzpicture}
\end{center}

\begin{ainote}
At the shared point, $(\partial_1\phi,\partial_2\phi)$ turns counterclockwise while
$(\partial_1\psi,\partial_2\psi)$ turns clockwise --- the two frames are mirror images of each
other. Corsin's caption: \qt{Here, $\det(D(\phi^{-1}\circ\psi)) < 0$!} This is exactly the atlas
$\{\phi,\psi\}$ failing \cref{def:orientable_manifold}: it covers the patch, but its transition
map has \emph{negative} Jacobian determinant there, so $\{\phi,\psi\}$ is an atlas but not an
orientation.
\end{ainote}

% Generator: Claude Sonnet 5 (Medium)
\begin{ainote}
\Cref{ex:12.3} and Exercise 12.6 (the latter not reproduced below, see the priority table)
invoke \newterm{Stokes' theorem} in its classical $\mathbb{R}^3$ form, which none of Corsin's
transcribed pages state explicitly. Recorded here for self-containedness, since the
exercise-sheet solutions below rely on it.
\end{ainote}

\begin{theorem}[Stokes' theorem, classical form in $\mathbb{R}^3$]
\label{thm:stokes_classical}
% Generator: Claude Sonnet 5 (Medium)
Let $S \subseteq \mathbb{R}^3$ be a compact, oriented $C^1$ surface with boundary $\partial S$
carrying the induced orientation, and let $F$ be a $C^1$ vector field on a neighborhood of
$S$. Then
\[\int_S \curl F \cdot d\vec{n} = \int_{\partial S} F \cdot d\vec{s}.\]
\end{theorem}

% Originally: content/week-13.tex, \section{Submanifolds with boundary and induced orientation} (Corsin Week 13)


\section{Submanifolds with boundary and induced orientation}
\label{sec:submanifolds_with_boundary}

% Source: Corsin Nick/Class Notes/Week 13.pdf, pp. 6--7
\begin{ainote}
Corsin recaps \cref{thm:local_parametrization} verbatim --- for every $p \in M$, an
open neighbourhood $V$ of $p$, an open set $U \subseteq \mathbb{R}^k$, and $\phi : U \to
V\cap M$ that is smooth, has $D\phi_x$ injective for all $x \in U$, and is a homeomorphism onto
$V\cap M$ --- before extending it to submanifolds with boundary below. Not repeated; see
\cref{sec:submanifolds}.
\end{ainote}

\begin{definition}[Submanifold with boundary]
\label{def:submanifold_with_boundary}
For a $k$-dimensional submanifold \emph{with boundary}, we require similarly that for a local
parametrization $\phi : U \to V\cap M$ as in \cref{thm:local_parametrization}, either $U
\subseteq \mathbb{R}^k$ is open, or
\[U = U' \cap \mathcal{H}^k, \qquad U' \subseteq \mathbb{R}^k \text{ open}, \qquad
  \mathcal{H}^k := \{(x_1,\ldots,x_k) \in \mathbb{R}^k : x_1 \leq 0\}\]
the (closed) \newterm{left half-space}. If $p \in \partial M$ and $\phi : U \subseteq
\mathcal{H}^k \to V\cap M$ is a local parametrization of $M$ around $p$, then $\phi|_{\partial
\mathcal{H}^k}$ (identifying $\partial\mathcal{H}^k = \{0\}\times\mathbb{R}^{k-1} \cong
\mathbb{R}^{k-1}$) is a local parametrization of $\partial M$ around $p$.
\end{definition}

% Sonnet 5 (Medium) --- FIG-W13-01
\begin{center}
\begin{tikzpicture}[>=Latex, scale=1.0]
  \draw[green!50!black, thick] (0,0) .. controls (0.3,1.3) and (1.3,2.0) .. (2.6,1.9)
    .. controls (3.9,1.8) and (4.6,1.0) .. (4.4,0.1) .. controls (3.5,-0.5) and (1.5,-0.6)
    .. (0,0);
  \fill (0.9,0.55) circle (1.1pt);
  \draw[red, dashed, thick] (-0.05,-0.55) .. controls (0.4,0.0) and (0.5,1.1) .. (0.15,1.55);
  \draw[red, thick, ->] (0.9,0.55) -- (1.75,0.75) node[right, font=\scriptsize] {$\partial_2\phi$};
  \draw[red, thick, ->] (0.9,0.55) -- (0.55,1.3) node[above, font=\scriptsize] {$\partial_1\phi$};
  \node[font=\scriptsize] at (0.6,-0.85) {$p \in \partial M$};
  \draw[->] (5.3,0.7) -- (6.3,0.7) node[right, font=\scriptsize] {};
  \node[font=\scriptsize] at (5.8,1.0) {$\phi$};
  \draw[red, thick] (7.1,0.1) rectangle (8.4,1.4);
  \draw[red, thick, dashed] (7.1,0.1) -- (7.1,1.4);
  \draw[red, ->] (7.1,0.75) -- (7.7,0.75) node[right, font=\scriptsize] {$e_1$};
  \draw[red, ->] (7.1,0.75) -- (7.1,1.25) node[above, font=\scriptsize] {$e_2$};
  \node[font=\scriptsize] at (7.1,-0.15) {$0$};
  \node[font=\scriptsize] at (7.75,-0.2) {$U \subseteq \mathcal{H}^2$};
\end{tikzpicture}
\end{center}

\begin{definition}[Induced orientation]
\label{def:induced_orientation}
If $\{\partial_1\phi,\ldots,\partial_k\phi\}$ is the oriented basis of $T_pM$ given by a local
parametrization $\phi$ around $p \in \partial M$, then $\{\partial_2\phi,\ldots,\partial_k
\phi\}$ is the \newterm{induced} oriented basis of $T_p(\partial M)$ around $p$.
\end{definition}

\subsection{Example: cylinder}
\label{sec:cylinder_orientation_example}

% Source: Corsin Nick/Class Notes/Week 13.pdf, pp. 8--9
Consider the cylindrical surface
\[C = \{(x,y,z) \in \mathbb{R}^3 : y^2+z^2 = r^2,\ 0\leq x\leq L\},\]
a $2$-dimensional submanifold with boundary, and the two parametrizations $\phi,\psi :
[-L,0]\times[0,2\pi] \to C$ (both domains $\subseteq \mathcal{H}^2$),
\[\phi(x,\theta) = \begin{pmatrix} -x \\ r\cos\theta \\ -r\sin\theta \end{pmatrix}, \qquad
  \psi(x,\theta) = \begin{pmatrix} x+L \\ r\cos\theta \\ r\sin\theta \end{pmatrix}.\]
Both cover all of $C$: as the local $x$-coordinate ranges over $[-L,0]$, $-x$ and $x+L$ both
range over $[0,L]$.

The two are positively oriented with respect to each other: writing out the composition,
\[\phi^{-1}\circ\psi(x,\theta) = \phi^{-1}(x+L,\,r\cos\theta,\,r\sin\theta) = (-x-L,-\theta)
  = -(x+L,\theta),\]
so $D(\phi^{-1}\circ\psi) = \begin{pmatrix}-1&0\\0&-1\end{pmatrix}$ and $\det D(\phi^{-1}
\circ\psi) = 1 > 0$: $\{\phi,\psi\}$ is a valid orientation of $C$ by
\cref{def:orientable_manifold}.

But restricted to the line $\partial\mathcal{H}^2 = \{(x,\theta) : x=0\}$, $\phi$ and $\psi$
parametrize \emph{opposite} boundary components of the cylinder --- $\phi(0,\cdot)$ traces the
circle at $x=0$, $\psi(0,\cdot)$ the circle at $x=L$ --- and \cref{def:induced_orientation}
gives them opposite rotational senses when both are drawn in the same $(y,z)$-view:

% Sonnet 5 (Medium) --- FIG-W13-02
\begin{center}
\begin{tikzpicture}[>=Latex, scale=0.95]
  \draw[blue!70!black, thick] (0,2.4) rectangle (4.5,3.5);
  \draw[blue!70!black, thick] (0,2.4) circle (0.55 and 0.55);
  \draw[blue!70!black, thick] (4.5,2.95) ellipse (0.35 and 0.55);
  \draw[red, thick, ->] (-0.5,3.5) arc (150:490:0.55);
  \node[red, font=\scriptsize] at (-0.85,3.5) {$t\mapsto\phi(0,t)$};
  \node[font=\scriptsize] at (0,1.7) {\textbf{left end} ($x=0$)};

  \draw[blue!70!black, thick] (6.5,2.4) rectangle (11,3.5);
  \draw[blue!70!black, thick, dashed] (6.5,2.95) ellipse (0.35 and 0.55);
  \draw[blue!70!black, thick] (11,2.4) circle (0.55 and 0.55);
  \draw[orange, thick, ->] (11.5,2.4) arc (-30:330:0.55);
  \node[orange, font=\scriptsize] at (11.9,3.5) {$t\mapsto\psi(0,t)$};
  \node[font=\scriptsize] at (8.75,1.7) {\textbf{right end} ($x=L$)};
\end{tikzpicture}
\end{center}

\begin{ainote}
This is not a contradiction, and Corsin does not present it as one: $\{\phi,\psi\}$ is a
perfectly good orientation of $C$ as a whole (the overlap check above is what
\cref{def:orientable_manifold} demands, and it passes). What the picture shows is that a
\emph{consistently} oriented cylinder induces boundary orientations on its two ends that look
like opposite rotational senses when both circles are drawn from the same external viewpoint
--- exactly the familiar fact, from the divergence theorem, that the outward normal at the two
ends of a solid cylinder points in opposite $x$-directions. Induced orientation is doing its
job correctly here, not failing.
\end{ainote}

% Originally: content/week-12.tex, \section{Solutions}

\section{Solutions}
\label{sec:pullbacks_solutions}

\begin{exercisesolution}[Solution of \cref{ex:12.9}]
\textbf{(a)} Writing each $d\psi^j$ in the coordinate basis, $d\psi^j = \sum_{i=1}^n
\frac{\partial\psi^j}{\partial x^i}\,dx^i$, so
\[\omega = \left(\sum_{i_1=1}^n \frac{\partial\psi^1}{\partial x^{i_1}}dx^{i_1}\right)\wedge
  \cdots\wedge\left(\sum_{i_n=1}^n \frac{\partial\psi^n}{\partial x^{i_n}}dx^{i_n}\right)
  = \sum_{i_1,\ldots,i_n} \frac{\partial\psi^1}{\partial x^{i_1}}\cdots
  \frac{\partial\psi^n}{\partial x^{i_n}}\,dx^{i_1}\wedge\cdots\wedge dx^{i_n}\]
by bilinearity of $\wedge$ (property (1) of \cref{def:wedge_product}). Any term with a
repeated index vanishes, since $dx^i\wedge dx^i=0$ by antisymmetry, so only tuples
$(i_1,\ldots,i_n) = (\sigma(1),\ldots,\sigma(n))$ for a permutation $\sigma \in S_n$
contribute; using $dx^{\sigma(1)}\wedge\cdots\wedge dx^{\sigma(n)} = \sgn(\sigma)\,dx^1\wedge
\cdots\wedge dx^n$,
\[\omega = \left(\sum_{\sigma\in S_n} \sgn(\sigma)\,\frac{\partial\psi^1}{\partial
  x^{\sigma(1)}}\cdots\frac{\partial\psi^n}{\partial x^{\sigma(n)}}\right) dx^1\wedge\cdots
  \wedge dx^n.\]
The coefficient is exactly the Leibniz expansion of $\det(D\psi)$, so
\[\boxed{\ d\psi^1\wedge\cdots\wedge d\psi^n = \det(D\psi)\,dx^1\wedge\cdots\wedge dx^n.\ }\]

\textbf{(b)} Let $\eta = f(y)\,dy^1\wedge\cdots\wedge dy^n$ be an $n$-form on $V$. Its
pullback under $\psi$ is $\psi^*\eta = f(\psi(x))\,d\psi^1\wedge\cdots\wedge d\psi^n$, which by
part (a) equals $f(\psi(x))\det(D\psi(x))\,dx^1\wedge\cdots\wedge dx^n$. By
\cref{def:integral_of_form_over_submanifold} (applied with $U$ itself as the $n$-dimensional
submanifold),
\[\int_U \psi^*\eta = \int_U f(\psi(x))\det(D\psi(x))\,dx.\]
Since $\det(D\psi) = \lvert\det(D\psi)\rvert>0$ by hypothesis, the right-hand side is exactly
the ordinary change-of-variables formula, so
\[\int_U \psi^*\eta = \int_U f(\psi(x))\,\lvert\det(D\psi(x))\rvert\,dx = \int_V f(y)\,dy
  = \int_V \eta.\]
Thus $\boxed{\int_U \psi^*\eta = \int_V \eta}$: the change-of-variables formula is exactly the
statement that integration of forms is invariant under pullback by an orientation-preserving
diffeomorphism.
\end{exercisesolution}

% Verified: solved independently, both parts match Sol12, including pullback-invariance.

\begin{exercisesolution}[Solution to \cref{ex:ai_pullback_area_form}]
% Generator: Claude Opus 5 (High)
\textbf{(a)} The pullback of a coordinate differential is the differential of the corresponding
component, so from $x = r\cos\theta$ and $y = r\sin\theta$,
\[\Phi^*(dx) = \cos\theta\,dr - r\sin\theta\,d\theta,
  \qquad
  \Phi^*(dy) = \sin\theta\,dr + r\cos\theta\,d\theta.\]
Wedging, and using $dr\wedge dr = d\theta\wedge d\theta = 0$ together with
$d\theta\wedge dr = -\,dr\wedge d\theta$,
\[\Phi^*(dx\wedge dy)
  = \big(\cos\theta\cdot r\cos\theta\big)\,dr\wedge d\theta
    + \big(-r\sin\theta\cdot\sin\theta\big)\,d\theta\wedge dr
  = r\big(\cos^2\theta + \sin^2\theta\big)\,dr\wedge d\theta
  = r\,dr\wedge d\theta.\]
The factor is exactly the $r$ tabulated for polar coordinates in
\cref{rem:coordinate_jacobians_table}, obtained here without computing a determinant: the
antisymmetry of the wedge did the cofactor expansion.

\textbf{(b)} They measure different things, and the absolute value is the difference.
\Cref{thm:change_of_variables} is a statement about \emph{volume}, which is unsigned: a region
and its mirror image have the same area, so only $\lvert\det\Jac\Phi\rvert$ can appear. The
pullback of a top-degree form is a statement about \emph{oriented} volume, which carries a sign,
and it records $\det\Jac\Phi$ itself.

The swap $\Phi(u_1,u_2) = (u_2,u_1)$ separates them. Its Jacobian is
$\left(\begin{smallmatrix}0&1\\1&0\end{smallmatrix}\right)$ with determinant $-1$, so
\[\Phi^*(dx\wedge dy) = du_2\wedge du_1 = -\,du_1\wedge du_2,\]
while $\lvert\det\Jac\Phi\rvert = 1$ and the change of variables formula does not change at all.
The form notices the reversal; the volume cannot.

\textbf{(c)} \Cref{ex:12.9} concludes that $\int_U\psi^*\eta = \int_V\eta$, with no sign. Drop
the hypothesis $\det > 0$ and take $\psi$ to be the swap above: the pullback acquires a factor
$-1$, so the two integrals differ by a sign and the conclusion is false as stated. The honest
statement without the hypothesis carries $\sgn(\det\Jac\psi)$.

\Cref{thm:change_of_variables} is untroubled by the same map because it never claimed anything
about orientation. This is the precise sense in which the positivity hypothesis in
\cref{ex:12.9} is not a technical convenience: it is the price of upgrading an unsigned identity
to a signed one.
\end{exercisesolution}

\begin{exercisesolution}[Solution to \cref{ex:ai_area_form_over_hemisphere}]
% Generator: Claude Opus 5 (High)
\textbf{(a)} Along $\phi$ the first two coordinates are the parameters themselves, $x = u$ and
$y = v$, so $\phi^*(dx) = du$ and $\phi^*(dy) = dv$, giving
\[\phi^*(dx\wedge dy) = du\wedge dv.\]
The third coordinate never enters, which is the whole reason this integral is easy. It remains to
check that $\phi$ induces the outward orientation: for a graph $z = g(u,v)$,
\[\partial_u\phi\times\partial_v\phi = (-g_u,\ -g_v,\ 1),\]
whose third component is $+1$, so the induced normal points \emph{upward}, and on the upper
hemisphere upward is outward. Hence
\[\int_{S^+} dx\wedge dy = \int_D du\wedge dv = \vol_2(D) = \pi.\]

\textbf{(b)} The pullback computation is identical, since $\psi$ also has $x = u$ and $y = v$,
so again $\psi^*(dx\wedge dy) = du\wedge dv$. But $\psi$ is still a graph, so its induced normal
still points upward, and on the \emph{lower} hemisphere upward is \emph{inward}. So $\psi$ is
orientation-reversing for the prescribed outward orientation, and
\cref{def:integral_of_form_over_submanifold} requires a parametrization compatible with the
chosen orientation. Correcting by the sign,
\[\int_{S^-} dx\wedge dy = -\int_D du\wedge dv = -\pi.\]

\textbf{(c)} The equator is a null set, so the two halves add:
\[\int_{S^2} dx\wedge dy = \pi + (-\pi) = 0.\]
For the second route, $d(x\,dy) = dx\wedge dy$, and $S^2$ is a compact orientable $2$-manifold
with $\partial S^2 = \emptyset$, so \cref{thm:stokes_general} gives
\[\int_{S^2} dx\wedge dy = \int_{S^2} d(x\,dy) = \int_{\partial S^2} x\,dy = 0.\]

\begin{remark}[$dx\wedge dy$ measures the shadow]
\label{rem:area_form_measures_the_shadow}
Part \textbf{(a)} says something worth keeping: integrating $dx\wedge dy$ over an oriented
surface in $\mathbb{R}^3$ returns the \emph{signed} area of its shadow on the $xy$-plane. The
hemisphere has area $2\pi$, but its shadow is the unit disc, and $\pi$ is what the computation
returns. A $2$-form does not measure surface area; the Gram determinant of
\cref{ch:gram_determinant} does that, and the two answers differ by exactly the tilt of the
surface.

Read that way, part \textbf{(c)} is obvious before any computation. The sphere's two halves cast
the same shadow, and a closed surface presents each patch of shadow once from above and once from
below, so everything cancels. That is the geometric content of $\int_M d\omega = 0$ on a closed
$M$, and \cref{q:quiz_q14} draws the other standard consequence from it.
\end{remark}
\end{exercisesolution}


% ============================================================
% Chapter 26 --- Stokes' and Green's theorems
% Originally: content/week-11.tex, the three subsections \subsection{Sneak peek at Stokes'
%             theorem}, \subsection{Green's theorem and conservative fields} and
%             \subsection{The Poincare lemma and simply connected domains} (all promoted
%             to sections), followed by content/week-13.tex, \section{Stokes' theorem} and
%             \section{Area via Green's formula}. Corsin's order is preserved: the special
%             cases in week 11, the general theorem in week 13.
% ============================================================
\chapter{Stokes' and Green's Theorems}
\label{ch:stokes}

% Transition: Claude Opus 5 (Medium)
Everything in this part converges here. Stokes' theorem says that integrating $d\omega$ over a
submanifold is the same as integrating $\omega$ over its boundary --- one statement that
contains the fundamental theorem of calculus, Green's theorem, the divergence theorem and the
classical Stokes theorem as special cases. The chapter follows Corsin's own order: the special
cases first, where the content is visible, then the general theorem, then what it buys --- among
other things a formula computing an area from a boundary integral alone.

% Originally: content/week-11.tex, \subsection{Sneak peek at Stokes' theorem} --- promoted to a section

\section{Sneak peek at Stokes' theorem}
\label{sec:stokes_sneak_peek}

\begin{ainote}
Before the formal introduction of Stokes' theorem (\cref{sec:stokes_general}), Sascha presented introductory examples to build intuition. The classical version of Stokes' theorem states that for a surface $A \subset \mathbb{R}^3$ and a vector field $F$,
\[\int_A (\curl F)\cdot \nu \mathop{d\text{vol}}_2 = \int_{\partial A} F \cdot \tau \,dL,\]
where $\nu$ is the unit normal to $A$ and $\tau$ is the unit tangent vector to the boundary curve $\partial A$.
\end{ainote}

\begin{example}[Stokes' theorem on a hemisphere]
\label{ex:stokes_hemisphere}
Let $A := \{(x,y,z) \in \mathbb{R}^3 \mid x^2+y^2+z^2=1,\ z \geq 0\}$ be the upper hemisphere.
Given the vector field $F(x,y,z) := (z^2, x^2, y^2)\transp$, we can evaluate both sides of Stokes' theorem.
\begin{enumerate}[label=\textbf{\arabic*.}]
    \item \textbf{The rotation of $F$:} $\curl F = \nabla \times F = (2y, 2z, 2x)\transp$.
    \item \textbf{The surface normal:} On the unit sphere, the outward normal is simply $\nu(x,y,z) = (x,y,z)\transp$.
    \item \textbf{The boundary $\partial A$:} The boundary of the upper hemisphere is the unit circle in the $xy$-plane, $x^2+y^2=1$ with $z=0$. This can be parametrized by the closed path $\gamma : [0,2\pi] \to \mathbb{R}^3$,
    \[\gamma(t) = \begin{pmatrix} \cos t \\ \sin t \\ 0 \end{pmatrix}.\]
\end{enumerate}
Evaluating the path integral (the right-hand side of Stokes' theorem):
\[\int_{\partial A} F \cdot \tau \,dL = \int_0^{2\pi} F(\gamma(t)) \cdot \gamma'(t) \,dt.\]
Since $z=0$ on $\gamma$, we have $F(\gamma(t)) = (0, \cos^2 t, \sin^2 t)\transp$. The derivative is $\gamma'(t) = (-\sin t, \cos t, 0)\transp$.
\[\int_0^{2\pi} \begin{pmatrix} 0 \\ \cos^2 t \\ \sin^2 t \end{pmatrix} \cdot \begin{pmatrix} -\sin t \\ \cos t \\ 0 \end{pmatrix} dt = \int_0^{2\pi} \cos^3 t \,dt = 0.\]
\end{example}

\begin{example}[Path integral using Stokes' theorem]
\label{ex:stokes_path_integral}
Compute $\oint_\Gamma F \cdot \tau \,dL$ where $F(x,y,z) := (x, y, z)\transp$ and $\Gamma$ is the unit circle in the $xy$-plane (oriented counter-clockwise).

Instead of computing the line integral directly, we can use Stokes' theorem. The curve $\Gamma$ bounds the unit disk $A = \{(x,y,0) \mid x^2+y^2 \leq 1\}$.
First, we compute the rotation of $F$:
\[\curl F = \nabla \times F = \begin{pmatrix} \partial_y(z) - \partial_z(y) \\ \partial_z(x) - \partial_x(z) \\ \partial_x(y) - \partial_y(x) \end{pmatrix} = \begin{pmatrix} 0 \\ 0 \\ 0 \end{pmatrix}.\]
By Stokes' theorem, the line integral along $\Gamma$ bounds a surface $A$ on the sphere (for example, the upper hemisphere). The surface integral is:
\[\oint_\Gamma F \cdot \tau \,dL = \int_A (\curl F) \cdot \nu \mathop{d\text{vol}}_2 = 0.\]
\end{example}

% Supplement: Sascha Brack/Class Notes/Week_11_Notes_Monday.pdf, pp. 5--10

% Originally: content/week-11.tex, \subsection{Green's theorem and conservative fields} --- promoted to a section

\section{Green's theorem and conservative fields}
\label{sec:greens_theorem_conservative_fields}

\begin{ainote}
Green's theorem, conservative vector fields, and path integrals were also discussed in the Monday session. 
\end{ainote}

\begin{theorem}[Green's Theorem]
\label{thm:greens_theorem}
Let $F : U \to \mathbb{R}^2$ be a continuously differentiable vector field on an open set $U \subseteq \mathbb{R}^2$. Then, for any compact $C^1$ domain $\Omega \subseteq U$,
\[\int_\Omega \curl F \mathop{d\text{vol}}_2 = \int_{\partial\Omega} F \cdot \tau \,dL,\]
where $\tau$ is the unit tangent vector to the boundary $\partial\Omega$ oriented counter-clockwise (i.e.\ $\tau = (-\nu_2, \nu_1)\transp$, with $\nu$ being the outward unit normal).
\end{theorem}

\begin{ainote}
% Supplement: Sascha Brack/Class Notes/Week_10_Notes_Monday.pdf, p. 16
\begin{center}
\begin{tikzpicture}[scale=1.5]
  % Domain
  \draw[ultra thick] (0,0) circle (1.5);
  % Boundary orientation
  \draw[red, ultra thick, ->] (1.5,0) arc (0:90:1.5);
  \draw[red, ultra thick, ->] (0,1.5) arc (90:180:1.5);
  \draw[red, ultra thick, ->] (-1.5,0) arc (180:270:1.5);
  \draw[red, ultra thick, ->] (0,-1.5) arc (270:360:1.5);
  
  \node[red, align=left] at (1.5, 1.8) {boundary\\is left over};
  \draw[red, ->] (1.0, 1.7) to[out=180, in=90] (0, 1.55);
  
  % Internal rotations
  \foreach \x/\y in {-0.6/-0.3, 0.4/-0.4, -0.2/0.1, 0.8/-0.8} {
    \draw[red, thick, ->] (\x,\y) arc (0:360:0.2);
  }
  
  % Two cancelling squares
  \draw[red, thick] (-0.5, 0.6) rectangle (0, 1.0);
  \draw[red, thick] (0, 0.6) rectangle (0.5, 1.0);
  % arrows
  \draw[red, thick, ->] (-0.25, 1.0) -- (-0.26, 1.0);
  \draw[red, thick, ->] (0.25, 1.0) -- (0.24, 1.0);
  \draw[red, thick, ->] (-0.5, 0.8) -- (-0.5, 0.79);
  \draw[red, thick, ->] (0.5, 0.8) -- (0.5, 0.81);
  \draw[red, thick, ->] (-0.25, 0.6) -- (-0.24, 0.6);
  \draw[red, thick, ->] (0.25, 0.6) -- (0.26, 0.6);
  
  \draw[blue, thick, ->] (-0.05, 0.8) -- (-0.05, 0.7);
  \draw[blue, thick, ->] (0.05, 0.7) -- (0.05, 0.8);
  
  \node[purple] at (0.4, 1.2) {cancels};
  \node[purple] at (0.8, 0.8) {out};
  \draw[purple, ->] (0.2, 1.1) to[out=180, in=90] (0, 0.9);
  
\end{tikzpicture}
\end{center}
We integrate the rotation inside the domain. Basically all the \qt{intermediate values} cancel out because for two points close to each other, the rotation points in opposite directions. Only the boundary is left over in the end!
\end{ainote}

\begin{example}[Area of a region via Green's theorem]
\label{ex:area_greens_theorem}
Consider a region $\Omega \subseteq \mathbb{R}^2$ defined by the region inside the curve $x^{2/3} + y^{2/3} = 1$ in the first quadrant, bounded by the axes. The curve can be parametrized as $\varphi(t) = (\cos^3 t, \sin^3 t)\transp$ for $t \in [0, \pi/2]$.

The area is $\vol_2(\Omega) = \int_\Omega 1 \mathop{d\text{vol}}_2$. We look for a vector field $F$ such that $\curl F = \partial_x F_2 - \partial_y F_1 \equiv 1$ on $\Omega$. A simple choice is $F(x,y) := (0, x)\transp$.
By Green's theorem, we have:
\[\vol_2(\Omega) = \int_\Omega \curl F \mathop{d\text{vol}}_2 = \int_{\partial\Omega} F \cdot \tau \,dL = \int_0^{\pi/2} F(\varphi(t)) \cdot \varphi'(t) \,dt.\]
Since $\tau(\varphi(t)) \,dL = \varphi'(t) \,dt$ along the parametrization, we compute the dot product:
\[F(\varphi(t)) \cdot \varphi'(t) = \begin{pmatrix} 0 \\ \varphi_1(t) \end{pmatrix} \cdot \begin{pmatrix} \varphi_1'(t) \\ \varphi_2'(t) \end{pmatrix} = \varphi_1(t)\varphi_2'(t).\]
This turns the area computation into a single 1D integral in $t$.
\end{example}

\begin{definition}[Conservative vector field]
\label{def:conservative_field}
Given a $C^k$ vector field $F$ on an open set $U \subseteq \mathbb{R}^n$, we call it \newterm{conservative} if there exists a scalar field $f : U \to \mathbb{R}$ of class $C^{k+1}$ such that $F \equiv \nabla f$. The function $f$ is called a \newterm{potential} for $F$.
\end{definition}

\begin{lemma}[Integrability conditions]
\label{lem:integrability_conditions}
Let $U \subseteq \mathbb{R}^n$ be open, and let $F$ be a $C^1$ conservative vector field on $U$ with components $F_1, \dots, F_n$. Then we necessarily have the integrability conditions:
\[\partial_j F_k = \partial_k F_j \qquad \text{for all } j, k \in \{1,\dots,n\}.\]
\end{lemma}
\begin{proof}
Let $f$ be a potential such that $F = \nabla f = (\partial_1 f, \dots, \partial_n f)\transp$. Then by Schwarz's theorem (symmetry of second derivatives),
\[\partial_j F_k = \partial_j(\partial_k f) = \partial_j\partial_k f = \partial_k\partial_j f = \partial_k(\partial_j f) = \partial_k F_j.\qedhere\]
\end{proof}

\begin{proposition}[Work along a path as difference of potential]
\label{prop:work_difference_potential}
Let $U \subseteq \mathbb{R}^n$ be open, and let $F : U \to \mathbb{R}^n$ be a continuous conservative vector field with potential $f$. Then for any $C^1$ path $\gamma : [0,1] \to U$, the line integral of $F$ along $\gamma$ depends only on the endpoints:
\[\int_\gamma F \cdot d\gamma = f(\gamma(1)) - f(\gamma(0)).\]
\end{proposition}
\begin{proof}
By definition of the path integral and the chain rule,
\begin{align*}
\int_\gamma F \cdot d\gamma &= \int_0^1 F(\gamma(t)) \cdot \gamma'(t) \,dt = \int_0^1 \nabla f(\gamma(t)) \cdot \gamma'(t) \,dt \\
&= \int_0^1 Jf(\gamma(t)) J\gamma(t) \,dt = \int_0^1 J(f \circ \gamma)(t) \,dt = \int_0^1 (f \circ \gamma)'(t) \,dt.
\end{align*}
By the fundamental theorem of calculus, this equals $(f \circ \gamma)(1) - (f \circ \gamma)(0) = f(\gamma(1)) - f(\gamma(0))$.
\end{proof}

\begin{example}[Line integral of a conservative field]
\label{ex:conservative_field_integral}
Consider the vector field $F : [0,1] \times [0,2] \to \mathbb{R}^2$ given by $F(x,y) := (y^2, 2xy)\transp$.
First, we show that $F$ is conservative by finding a potential $f$ such that $\nabla f = F$. We need $\partial_x f = y^2$ and $\partial_y f = 2xy$. Integrating the first equation with respect to $x$ gives $f(x,y) = xy^2 + c(y)$. Differentiating with respect to $y$ gives $\partial_y f = 2xy + c'(y)$, which must equal $2xy$. Thus $c'(y) = 0$, and we can choose $c(y) = 0$. The potential is $f(x,y) = xy^2$.

Now, let $\Gamma$ be a closed path on the boundary $\partial U$ of the domain, starting and ending at $\gamma(0) = \gamma(1) = (0,0)\transp$. By \cref{prop:work_difference_potential}, the line integral around this closed loop is simply the difference in potential at the endpoints:
\[\oint_\Gamma F \cdot d\gamma = f(\gamma(1)) - f(\gamma(0)) = f(0,0) - f(0,0) = 0.\]
\end{example}

% Originally: content/week-11.tex, \subsection{The Poincare lemma and simply connected domains} --- promoted

\section{The Poincaré lemma and simply connected domains}
\label{sec:poincare_lemma_simply_connected}

\begin{theorem}[Poincaré Lemma]
\label{thm:poincare_lemma}
Let $U \subseteq \mathbb{R}^n$ be open, and let $\alpha \in \Omega^1(U)$ be a $C^1$ closed $1$-form, that is, $d\alpha = 0$.
Let $\gamma_0, \gamma_1 : [0,1] \to U$ be piecewise $C^1$ paths with the same initial point $x_0$ and the same endpoint $x_1$. Suppose that $\gamma_0$ and $\gamma_1$ are \newterm{homotopic} in $U$ with fixed endpoints. For simplicity, assume that the homotopy $H : [0,1]^2 \to U$ is smooth. Then
\[\int_{\gamma_0} \alpha = \int_{\gamma_1} \alpha.\]
\end{theorem}

\begin{corollary}[Existence of potentials in simply connected domains]
\label{cor:potentials_simply_connected}
Let $U \subseteq \mathbb{R}^n$ be an open \newterm{simply connected} set. Then every closed $C^1$ $1$-form on $U$ is exact. That is, if $\alpha \in \Omega^1(U)$ satisfies $d\alpha = 0$, then there exists a function $f \in C^2(U)$ such that $\alpha = df$.
Equivalently, if $F \in C^1(U, \mathbb{R}^n)$ satisfies the integrability conditions
\[\partial_j F_k = \partial_k F_j \qquad \text{for all } j, k \in \{1,\dots,n\},\]
then $F$ is conservative: there exists a potential $f \in C^2(U)$ such that $\nabla f = F$.
\end{corollary}

\begin{ainote}[Intuition on simply connected spaces]
A space $U$ is simply connected if any closed loop in $U$ can be continuously contracted to a point within $U$ (i.e.\ any two paths with the same endpoints are homotopic). 

% Supplement: Sascha Brack/Class Notes/Week_12_Notes_Monday.pdf, p. 7
\begin{center}
\begin{tikzpicture}[scale=1.5]
  \node[circle, fill=blue, inner sep=1.5pt, label=left:{$x_0$}] (x0) at (0,0) {};
  \node[circle, fill=blue, inner sep=1.5pt, label=right:{$x_1$}] (x1) at (3,0.5) {};
  
  \draw[thick] (x0) to[out=60, in=150] node[pos=0.5, above] {$\gamma_0$} (x1);
  \draw[thick] (x0) to[out=40, in=170] node[pos=0.5, above] {$\gamma_{0.25}$} (x1);
  \draw[thick] (x0) to[out=10, in=200] node[pos=0.5, below] {$\gamma_{0.5}$} (x1);
  \draw[thick] (x0) to[out=-20, in=220] node[pos=0.5, below] {$\gamma_{0.75}$} (x1);
  \draw[thick] (x0) to[out=-40, in=240] node[pos=0.5, below] {$\gamma_1 = H(1,\cdot)$} (x1);
  \node at (-0.5, 0.5) {$H(0,\cdot) = \gamma_0$};
  
  \node[blue, align=left, font=\small] at (4, -0.2) {We try to \textbf{reshape}\\the path continuously from\\$\gamma_0$ to $\gamma_1$.};
\end{tikzpicture}
\end{center}

Examples of simply connected spaces:
\begin{itemize}
    \item $\mathbb{R}^n$
    \item $\mathbb{R}^3 \setminus \{(0,0,0)\}$
\end{itemize}
Examples of spaces that are \emph{not} simply connected:
\begin{itemize}
    \item $S^1 \subseteq \mathbb{R}^2$ (the unit circle)
    \item $\mathbb{R}^3 \setminus \operatorname{span}\{(0,0,1)\transp\}$ (the $z$-axis removed)
    \item The torus
\end{itemize}
This topological property is what allows the integrability conditions to globally imply the existence of a potential.
\end{ainote}

% Supplement: Analysis_II_Script_v1.pdf, p. 153
\begin{example}[Integrability without a potential: why simple connectedness is needed]
\label{ex:vortex_field_not_conservative}
Let $U := \mathbb{R}^2\setminus\{0\}$ (which is \emph{not} simply connected, unlike its
$3$-dimensional analogue $\mathbb{R}^3\setminus\{0\}$ listed above) and
\[F(x,y) := \left(\frac{-y}{x^2+y^2},\ \frac{x}{x^2+y^2}\right), \qquad (x,y) \in U.\]
A direct computation gives
\[\partial_1F_2 = \partial_x\!\left(\frac{x}{x^2+y^2}\right) = \frac{-x^2+y^2}{(x^2+y^2)^2}
  = \partial_y\!\left(\frac{-y}{x^2+y^2}\right) = \partial_2F_1,\]
so $F$ satisfies the integrability conditions (\cref{lem:integrability_conditions})
\emph{everywhere} on $U$. Yet $F$ is not conservative: for the loop
$\gamma(t) := (\cos(2\pi t),\sin(2\pi t))$, $t\in[0,1]$, tracing the unit circle once
counterclockwise,
\[\oint_\gamma F\cdot d\gamma
  = 2\pi\int_0^1\langle(-\sin(2\pi t),\cos(2\pi t)),(-\sin(2\pi t),\cos(2\pi t))\rangle\,dt
  = 2\pi \neq 0,\]
which by \cref{prop:work_difference_potential} is impossible for a conservative field (a closed
loop would have to return $f(\gamma(1))-f(\gamma(0))=0$). This is exactly the witness that
\cref{cor:potentials_simply_connected} needs the hypothesis it has: the integrability
conditions are necessary but, on a domain that is not simply connected, not sufficient.
\end{example}

% Originally: content/week-13.tex, \section{Stokes' theorem} (Corsin Week 13)


\section{Stokes' theorem}
\label{sec:stokes_general}

% Source: Corsin Nick/Class Notes/Week 13.pdf, p. 2
\begin{theorem}[Stokes' theorem]
\label{thm:stokes_general}
Let $M \subseteq \mathbb{R}^n$ be an orientable $m$-manifold with (possibly empty) boundary,
and let $\omega \in \Omega^{m-1}(\mathbb{R}^n)$ have compact support in $\mathbb{R}^n$. Then
\[\int_{\partial M} \omega = \int_M d\omega.\]
\end{theorem}

\begin{ainote}
\Cref{thm:stokes_classical} was stated ahead of Corsin's own notes, as background
needed to solve \cref{ex:12.3}, in its classical $\mathbb{R}^3$ curl-theorem form. It is
exactly the case $m=2$, $n=3$, $\omega=\omega_F$ of the general theorem above, as
\cref{sec:stokes_in_r3} below confirms directly from Corsin's own derivation.
\end{ainote}

\subsection{Example: Green's theorem}
\label{sec:green_theorem_example}

% Source: Corsin Nick/Class Notes/Week 13.pdf, p. 2
\begin{example}[The exterior derivative of a $1$-form on $\mathbb{R}^2$]
\label{ex:exterior_derivative_r2}
Let $\omega(x,y) = P(x,y)\,dx + Q(x,y)\,dy$ be a $1$-form on $\mathbb{R}^2$. Then
\[
d\omega = dP\wedge dx + dQ\wedge dy = \left(\frac{\partial P}{\partial x}dx +
  \frac{\partial P}{\partial y}dy\right)\wedge dx + \left(\frac{\partial Q}{\partial x}dx +
  \frac{\partial Q}{\partial y}dy\right)\wedge dy = \left(\frac{\partial Q}{\partial x} -
  \frac{\partial P}{\partial y}\right) dx\wedge dy,
\]
using $dx\wedge dx = dy\wedge dy = 0$ and $dy\wedge dx = -dx\wedge dy$.
\end{example}

\subsection{Example: proof of Green's theorem}
\label{sec:green_theorem_proof}

% Source: Corsin Nick/Class Notes/Week 13.pdf, p. 3
\begin{ainote}
The theorem below is reproduced by Corsin directly from the lecture script (numbered there as
Theorem 14.24), as with the two definitions quoted in \cref{sec:orientable_manifolds}.
\end{ainote}

\begin{theorem}[Green's theorem in the plane]
\label{thm:green_theorem}
Let $F : U \to \mathbb{R}^2$ be a continuously differentiable vector field on an open set
$U \subseteq \mathbb{R}^2$. Then, for any compact $C^1$ domain $\Omega \subseteq U$,
\[\int_\Omega \curl F\,dx = \int_{\partial\Omega} F\cdot\tau\,dL,\]
with $\tau = \nu^\perp = (-\nu_2,\nu_1)$, where $\nu$ is the exterior unit normal.
\end{theorem}

Writing $F(x,y) = (P(x,y),Q(x,y))$, so that $\curl F = \partial_xQ - \partial_yP$ (the scalar
$2$-dimensional curl), \cref{thm:green_theorem} follows from
\cref{thm:stokes_general,ex:exterior_derivative_r2} applied to $\omega = P\,dx+Q\,dy$:
\[
\int_\Omega \curl F\,dx\,dy = \int_\Omega \left(\frac{\partial Q}{\partial x} -
  \frac{\partial P}{\partial y}\right) dx\wedge dy = \int_\Omega d\omega
  \overset{\text{Stokes}}{=} \int_{\partial\Omega} P(x,y)\,dx + Q(x,y)\,dy
  = \int_{\partial\Omega} F\cdot\tau\,dL.
\]

\begin{proof}[Justification of the last step]
Suppose $\gamma : [0,L] \to \partial\Omega$ is a $C^1$ curve of unit speed and positive
orientation, i.e.\ $\gamma'(t) = \tau(\gamma(t))$. Then
\[\langle F,\tau\rangle = \langle F,\gamma'(t)\rangle = (P\,dx+Q\,dy)(\gamma'(t)),\]
since $\tau = \begin{pmatrix}0&-1\\1&0\end{pmatrix}\nu$ is the tangent vector consistent with
the induced orientation, so integrating $(P\,dx+Q\,dy)(\gamma'(t))$ over $t \in [0,L]$ is
exactly $\int_{\partial\Omega} P\,dx+Q\,dy$ by \cref{def:integral_of_form_over_submanifold}.
\end{proof}

\subsection{Example: Stokes' theorem in \texorpdfstring{$\mathbb{R}^3$}{R\^{}3}}
\label{sec:stokes_in_r3}

% Source: Corsin Nick/Class Notes/Week 13.pdf, p. 4
Recall from \cref{ex:exterior_derivative_gradient_curl} that for $\omega_F =
\langle F,\cdot\rangle$, $d\omega_F = \langle\curl F,\cdot\times\cdot\rangle$. So, if $\Omega
\subseteq \mathbb{R}^3$ is a $2$-dimensional submanifold with boundary, parametrized by $\phi$
with tangent vectors $\tau_i := \partial_i\phi$ and exterior unit normal $\nu$, then
\[
\int_\Omega d\omega_F = \int_\Omega \langle\curl F,\tau_1\times\tau_2\rangle\,dx
  = \int_\Omega \langle\curl F,\nu\rangle\,d\vol_2 = \int_\Omega \vec\nabla\times\vec F\cdot
  d\vec A,
\]
and
\[\int_{\partial\Omega} \omega_F = \int_{\partial\Omega} \langle F,\tau\rangle\,d\vol_1
  = \int_{\partial\Omega} \vec F\cdot d\vec s.\]
By \cref{thm:stokes_general}, these are equal, which recovers \cref{thm:stokes_classical}.

% Source: Jérôme Paschoud/Notizen/Woche 13.1 Stokes.pdf, p. 4
\begin{theorem}[Stokes' theorem in 3D]
\label{thm:stokes_3d}
For an orientable, parametrized $2$-dimensional submanifold $M \subseteq \mathbb{R}^3$ with boundary $\partial M$, and a $C^1$ vector field $F$:
\[\int_M \curl(F) \cdot \nu \,d\vol_2 = \int_{\partial M} F \cdot \tau \,d\vol_1.\]
The orientation of $\partial M$ is chosen such that when traversing $\partial M$ in the direction of $\tau$, the surface $M$ (given by the normal $\nu$) is to the left (the ``right-hand rule'').
\end{theorem}

% Generator: Claude Opus 5 (High) --- FIG-26-03. This section is the chapter's main one and had
% no figure at all, while the theorem just above states its orientation convention purely in
% words. "The surface is to the left" is the one convention in this part of the course that a
% reader cannot reliably reconstruct from prose, and getting it backwards flips the sign of
% every answer, which is why it earns a picture here rather than a longer sentence.
%
% GEOMETRY, checked. The surface is a horizontal disc in the plane z = 0, drawn under the
% projection screen = (x, 0.4y + z), so a disc of radius 2 becomes an ellipse with semi-axes
% 2 and 0.8. The parametrization (2 cos t, 2 sin t, 0) maps to screen (2 cos t, 0.8 sin t),
% which runs right -> top -> left -> bottom as t increases, i.e. COUNTERCLOCKWISE on screen,
% and counterclockwise seen from +z is the sense that goes with the upward normal.
%
% The marked point is t = pi/2, the top of the ellipse, chosen because it is the only place
% where all three vectors are visually distinct:
%   tau     = d/dt (2 cos t, 2 sin t, 0) at t = pi/2, normalized  = (-1, 0, 0) -> LEFT on screen
%   nu      = (0, 0, 1)                                                        -> UP on screen
%   inward  = nu x tau = (0,0,1) x (-1,0,0) = (0, -1, 0)                       -> DOWN on screen
% The cross product is the content of the figure and is worth checking rather than trusting:
% (0,0,1) x (-1,0,0) = (0*0 - 1*0, 1*(-1) - 0*0, 0*0 - 0*(-1)) = (0,-1,0), and at the point
% (0,2,0) the direction (0,-1,0) does point back towards the centre, so it really is inward.
% Read that as the walker's rule: facing tau with nu upright, the left hand points along
% nu x tau, which is into M. At t = 0 this figure would be undrawable, since tau and nu would
% both project onto the same upward screen direction.
\begin{center}
\begin{tikzpicture}[>=Latex, scale=1.1]
  % The surface, and its boundary traversed counterclockwise on screen.
  \fill[MidnightBlue!7] (0,0) ellipse (2 and 0.8);
  \draw[MidnightBlue!70, thick] (0,0) ellipse (2 and 0.8);
  \foreach \a in {150,215,330} {
    \draw[MidnightBlue!70, very thick, ->]
      ({2*cos(\a-8)},{0.8*sin(\a-8)}) arc[start angle=\a-8, end angle=\a+8,
        x radius=2, y radius=0.8];
  }
  \node[MidnightBlue!70, font=\small] at (1.15,-0.30) {$M$};
  \node[MidnightBlue!70, font=\scriptsize, anchor=north] at (-1.62,-0.62) {$\partial M$};

  % The three directions at the top of the rim, where all three are distinguishable.
  \draw[OliveGreen, very thick, ->] (0,0.8) -- (0,1.85)
    node[anchor=south, font=\scriptsize] {$\nu$};
  \draw[BrickRed, very thick, ->] (0,0.8) -- (-1.05,0.8)
    node[anchor=south, font=\scriptsize] {$\tau$};
  % Anchored to the LEFT of the dashed arrow: anchored to its right it came to rest just before
  % the "$M$" label inside the ellipse, and the two read as one string, "into M M".
  \draw[gray!50!black, thick, dashed, ->] (0,0.8) -- (0,0.12)
    node[anchor=north east, font=\scriptsize, text=gray!45!black] {into $M$};
  \fill[MidnightBlue!70] (0,0.8) circle (1.5pt);

  \node[font=\scriptsize, anchor=north, align=center] at (0,-1.15)
    {walking along $\tau$ with $\nu$ upright, $M$ lies to the left:\\
     the leftward direction is $\nu\times\tau$};
\end{tikzpicture}
\end{center}

\begin{remark}[The convention as a formula]
\label{rem:right_hand_rule_as_cross_product}
\qt{To the left} is a reliable instruction only once one knows which way is up, and $\nu$ is what
supplies that. The picture above turns the sentence into an identity that can be checked instead
of visualised: at a boundary point, the direction pointing from $\partial M$ into $M$ is
\[\text{inward} = \nu\times\tau.\]
Equivalently $\tau = \text{inward}\times\nu$, so any two of the three determine the third. On the
unit disc in the $xy$-plane with $\nu = e_3$, at the point $(1,0,0)$, the positively oriented
$\tau$ is $(0,1,0)$ and $\nu\times\tau = (-1,0,0)$ does point at the origin, confirming the
familiar counterclockwise convention.

Reversing $\nu$ reverses $\tau$, which is why both sides of \cref{thm:stokes_3d} change sign
together and the theorem is insensitive to the choice. What it is \emph{not} insensitive to is
choosing the two independently, and that is the error the formula above rules out.
\end{remark}

\begin{example}[Stokes' theorem on a hemisphere]
\label{ex:stokes_hemisphere_linear_field}
% Michael's Bsp. 20.1.2
Let $F(x,y,z) := (z, x, y)\transp$. We want to compute the flux of the curl of $F$ through the upper hemisphere $H := \{(x,y,z) \in \mathbb{R}^3 \mid x^2 + y^2 + z^2 = 1,\ z \geq 0\}$.

First, the curl of $F$ is:
\[\curl(F) = \begin{pmatrix} \partial_y F_3 - \partial_z F_2 \\ \partial_z F_1 - \partial_x F_3 \\ \partial_x F_2 - \partial_y F_1 \end{pmatrix} = \begin{pmatrix} 1 - 0 \\ 1 - 0 \\ 1 - 0 \end{pmatrix} = \begin{pmatrix} 1 \\ 1 \\ 1 \end{pmatrix}.\]
By \cref{thm:stokes_3d}, the flux of $\curl(F)$ through $H$ equals the line integral of $F$ along the boundary $\partial H$. The boundary is the unit circle in the $xy$-plane, which we parametrize by $\gamma(t) = (\cos t, \sin t, 0)\transp$ for $t \in [0, 2\pi]$. The tangent vector is $\gamma'(t) = (-\sin t, \cos t, 0)\transp$.
The line integral is:
\begin{align*}
\int_{\partial H} F \cdot d\gamma &= \int_0^{2\pi} F(\gamma(t)) \cdot \gamma'(t) \,dt \\
&= \int_0^{2\pi} \begin{pmatrix} 0 \\ \cos t \\ \sin t \end{pmatrix} \cdot \begin{pmatrix} -\sin t \\ \cos t \\ 0 \end{pmatrix} \,dt \\
&= \int_0^{2\pi} \cos^2 t \,dt = \pi.
\end{align*}
Thus, $\int_H \curl(F) \cdot \nu \,d\vol_2 = \pi$.
\end{example}

% Supplement: Analysis_II_Script_v1.pdf, p. 163
\begin{exercise}[Flux through a cylinder, direct and via Stokes]
\label{ex:cylinder_flux_direct_and_stokes}
Let $F(x,y,z) := (yz,x^2,1)\transp$ on $\mathbb{R}^3$, and let $M$ be the lateral surface of the
cylinder $M := \{(x,y,z) : x^2+y^2=1,\ 0<z<1\}$, with the outward-pointing normal
$\nu(x,y,z)=(x,y,0)$. Compute $\int_M\curl(F)\cdot\nu\,d\vol_2$ directly, and again using
Stokes' theorem.
\end{exercise}

\begin{exercisesolution}[Proof of \cref{ex:cylinder_flux_direct_and_stokes}]
% Generator: Claude Sonnet 5 (High)
First, $\curl(F) = (\partial_yF_3-\partial_zF_2,\ \partial_zF_1-\partial_xF_3,\
\partial_xF_2-\partial_yF_1) = (0,\ y,\ 2x-z)$.

\textbf{Directly.} Parametrize $M$ by $\phi(\theta,z) := (\cos\theta,\sin\theta,z)$,
$(\theta,z)\in[0,2\pi)\times(0,1)$; then
$\partial_\theta\phi\times\partial_z\phi = (\cos\theta,\sin\theta,0) = \nu(\phi(\theta,z))$,
already a unit vector, so $d\vol_2 = d\theta\,dz$. Then
$\curl(F)(\phi(\theta,z)) = (0,\sin\theta,2\cos\theta-z)$, and
\[\int_M\curl(F)\cdot\nu\,d\vol_2 = \int_0^1\!\int_0^{2\pi}
  \big\langle(0,\sin\theta,2\cos\theta-z),(\cos\theta,\sin\theta,0)\big\rangle\,d\theta\,dz
  = \int_0^1\!\int_0^{2\pi}\sin^2\theta\,d\theta\,dz = \pi.\]

\textbf{Via Stokes.} The boundary $\partial M$ consists of the two circles $z=0$ and $z=1$; the
orientation induced by the outward normal $\nu$ (right-hand rule) traverses the \emph{bottom}
circle ($z=0$) with $\theta$ increasing and the \emph{top} circle ($z=1$) with $\theta$
\emph{decreasing}. The parameter rectangle $[0,2\pi]\times[0,1]$ has $z=0$ as part of its
positively-oriented boundary traversed left to right and $z=1$ traversed right to left, and the
two vertical edges $\theta=0,2\pi$ cancel since they are the same curve in $M$. This is the
phenomenon drawn in \cref{sec:submanifolds_with_boundary}: one consistently oriented cylinder
induces senses on its two rims that turn opposite ways when both are viewed from the same side,
and \cref{rem:cylinder_opposite_boundary_senses} explains why that is the convention working
rather than failing. With
$\gamma_0(\theta):=(\cos\theta,\sin\theta,0)$ and $\gamma_1(\theta):=(\cos\theta,\sin\theta,1)$,
both for $\theta\in[0,2\pi]$,
\[F(\gamma_0(\theta)) = (0,\cos^2\theta,1), \qquad F(\gamma_1(\theta)) = (\sin\theta,\cos^2\theta,1),\]
and $\gamma_0'(\theta)=\gamma_1'(\theta)=(-\sin\theta,\cos\theta,0)$, so
\[\int_0^{2\pi} F(\gamma_0(\theta))\cdot\gamma_0'(\theta)\,d\theta = \int_0^{2\pi}\cos^3\theta\,d\theta = 0,\]
\[\int_0^{2\pi} F(\gamma_1(\theta))\cdot\gamma_1'(\theta)\,d\theta
  = \int_0^{2\pi}\big({-}\sin^2\theta+\cos^3\theta\big)\,d\theta = -\pi.\]
The bottom circle is traversed with $\theta$ increasing (contributing $+0$) and the top circle
with $\theta$ \emph{decreasing} (contributing $-({-}\pi) = \pi$), so
$\int_{\partial M}F\cdot\tau\,d\vol_1 = 0+\pi = \pi$, agreeing with the direct computation.
\end{exercisesolution}

% Supplement: Analysis_II_Script_v1.pdf, p. 163
\begin{exercise}[Flux through a spherical cap, via Stokes and via Gauss]
\label{ex:spherical_cap_flux_stokes_gauss}
Let $0<h<1$ and let $S := \{(x,y,z)\in S^2 : z>-h\}$ be the part of the unit sphere above
height $-h$. Consider $F(x,y,z) := (-y,x,0)\transp$. Calculate the flux
$\int_S\curl(F)\cdot\nu\,d\vol_2$ (outward normal) directly, once using Stokes' theorem, and
once using the divergence theorem.
\end{exercise}

\begin{exercisesolution}[Proof of \cref{ex:spherical_cap_flux_stokes_gauss}]
% Generator: Claude Sonnet 5 (High)
Here $\curl(F) = (0,0,2)$, a constant field, and on $S^2$ the outward normal is
$\nu(x,y,z)=(x,y,z)$ itself, so $\langle\curl(F),\nu\rangle = 2z$.

\textbf{Directly.} Using colatitude $\varphi\in[0,\arccos(-h)]$ (so $z=\cos\varphi>-h$) and
$d\vol_2=\sin\varphi\,d\varphi\,d\theta$,
\[\int_S 2z\,d\vol_2 = \int_0^{2\pi}\!\int_0^{\arccos(-h)} 2\cos\varphi\sin\varphi\,d\varphi\,d\theta
  = 2\pi\Big[{-}\tfrac12\cos(2\varphi)\Big]_0^{\arccos(-h)} = 2\pi(1-h^2),\]
using $\cos(2\arccos(-h)) = 2h^2-1$.

\textbf{Via Stokes.} $\partial S$ is the circle $x^2+y^2=1-h^2$, $z=-h$, of radius
$r:=\sqrt{1-h^2}$, oriented (right-hand rule, outward normal) by
$\gamma(\theta) := (r\cos\theta,r\sin\theta,-h)$, $\theta\in[0,2\pi]$. Then
$\gamma'(\theta) = (-r\sin\theta,r\cos\theta,0)$ and $F(\gamma(\theta)) =
(-r\sin\theta,r\cos\theta,0)$, so
\[\int_{\partial S}F\cdot d\gamma = \int_0^{2\pi}\big(r^2\sin^2\theta+r^2\cos^2\theta\big)\,d\theta
  = 2\pi r^2 = 2\pi(1-h^2),\]
agreeing with the direct computation.

\textbf{Via Gauss.} Since $\divg(\curl F)\equiv0$ (\cref{ex:curl_grad_div_curl_zero}), the flux
of $\curl(F)$ through any surface with boundary $\partial S$ is the same. Replace $S$ by the
flat disk $D := \{x^2+y^2\leq1-h^2,\ z=-h\}$, with upward normal $(0,0,1)$: $S\cup(-D)$
(the cap together with the disk, reversed) bounds the solid cap $\Omega$, and the divergence
theorem gives $\int_S\curl(F)\cdot\nu\,d\vol_2 - \int_D\curl(F)\cdot(0,0,1)\,d\vol_2 =
\int_\Omega\divg(\curl F)\,d\vol_3 = 0$. Since $\langle\curl F,(0,0,1)\rangle\equiv2$,
\[\int_S\curl(F)\cdot\nu\,d\vol_2 = \int_D 2\,d\vol_2 = 2\cdot\vol_2(D) = 2\pi(1-h^2),\]
the area of a disk of radius $\sqrt{1-h^2}$ times $2$, again the same answer.
\end{exercisesolution}

% Generator: Claude Sonnet 5 (High) --- new content, tying together four theorems this document
% proves separately (FTC-via-Gauss in \cref{sec:gauss_intuition}, Green in
% \cref{sec:green_theorem_proof}, classical Stokes in \cref{sec:stokes_in_r3}, and the
% divergence theorem in \cref{sec:flux_divergence_theorem}) as one theorem, stated once, here.
\begin{remark}[One theorem, four costumes]
\label{rem:one_theorem_four_costumes}
Every integral theorem proved in this part of the course is the single identity
$\int_M d\omega = \int_{\partial M}\omega$ of \cref{thm:stokes_general}, for a different choice
of the dimension $m$ of $M$ and the degree of $\omega$. Collected in one place:
\begin{center}
\small
\begin{tabular}{@{}lllll@{}}
\toprule
\textbf{Theorem} & $M$ & $m$ & $\omega$ & $d\omega$ \\
\midrule
FTC &
  $[a,b]\subset\mathbb{R}$ & $1$ & $f$ ($0$-form) & $f'\,dx$ \\
Green (\cref{thm:green_theorem}) &
  $\Omega\subseteq\mathbb{R}^2$ & $2$ & $P\,dx+Q\,dy$ & $(\partial_xQ{-}\partial_yP)\,dx\!\wedge\! dy$ \\
Classical Stokes (\cref{thm:stokes_classical}) &
  surface $S\subset\mathbb{R}^3$ & $2$ & $\omega_F$ & $\omega_{\curl F}^2$ \\
Divergence (\cref{thm:gauss_divergence}) &
  $\Omega\subseteq\mathbb{R}^n$ & $n$ & $\omega_F^{n-1}$ & $(\divg F)\,dx^1\!\wedge\!\cdots\!\wedge\! dx^n$ \\
\bottomrule
\end{tabular}
\end{center}
Each row's $\int_M d\omega = \int_{\partial M}\omega$ reads off as exactly that row's theorem.
The top row is \cref{sec:gauss_intuition}'s $n=1$ computation seen from the other side: there,
$f$ played the role of a $1$-dimensional vector field and the identity was reached via the
divergence theorem; here $f$ is instead a $0$-form and the same identity is the $m=1$ case of
\cref{thm:stokes_general} directly. These are two routes to the same fact, matching
\cref{ex:exterior_derivative_gradient_curl}'s remark that $df = \langle\nabla f,\cdot\rangle$ is
exactly the shape that makes this specialization work. The bottom row uses
$d\omega_F^{n-1} = (\divg F)\,dx^1\wedge\cdots\wedge dx^n$
(\cref{ex:exterior_derivative_gradient_curl}, item 3) together with the identification of
$\int_{\partial\Omega}\omega_F^{n-1}$ with the flux $\int_{\partial\Omega}\langle
F,\nu\rangle\,d\vol_{n-1}$ already noted in \cref{ch:pullbacks}.

The pattern is worth stating in words, because it is easy to lose under the notation: a
$k$-dimensional derivative measured \emph{inside} a region always equals the corresponding
$(k{-}1)$-dimensional quantity measured on its \emph{boundary}. Differentiation and taking a
boundary are, in this precise sense, the same operation viewed from two sides. That is also
exactly the content of $d^2=0$ (\cref{thm:exterior_derivative}\textbf{(3)}): a boundary has no
boundary, just as $d\omega$ has no further exterior derivative to take that isn't already zero.
\end{remark}

% Source: old_exams/examFS24.pdf (21 August 2024), Problem 3, parts 1--2, p. 10
% Extractor: Gemini 3.6 Flash (High)
% Correction: Claude Opus 5 (High) --- \operatorname{rot} -> \curl and J\Phi -> \Jac\Phi per the
% declared-macro rule; \iint -> \int, matching thm:green_theorem two pages above.
\begin{exercise}[Curl identity and Green's formula area preservation]
\label{ex:curl_identity_greens_formula}
Let $\Omega \subseteq \mathbb{R}^2$ be a bounded $C^1$ domain, and let $\Phi = (\Phi_1, \Phi_2) \in C^2(\overline{\Omega}, \mathbb{R}^2)$ be a vector-valued mapping, not assumed to be a diffeomorphism.
\begin{enumerate}[label=\textbf{(\alph*)}]
  \item Consider the vector field $V \in C^1(\overline{\Omega}, \mathbb{R}^2)$ defined by
  \[V := (\Phi_1 \partial_1 \Phi_2, \Phi_1 \partial_2 \Phi_2)\transp = \Phi_1 \nabla \Phi_2.\]
  Prove that $\curl V(x) = \det \Jac\Phi(x)$ for all $x \in \overline{\Omega}$.
  \item Let $\Psi \in C^2(\overline{\Omega}, \mathbb{R}^2)$ be another mapping such that $\Phi(x) = \Psi(x)$ for all $x \in \partial \Omega$. Prove that
  \[\int_\Omega \det \Jac\Phi(x) \, dx = \int_\Omega \det \Jac\Psi(x) \, dx.\]
\end{enumerate}
\exhint[Official hint]{Use Green's formula (i.e.\ 2D Stokes) on the plane.}
\exinfo{This exercise is taken from the exam of 21 August 2024, Problem 3. The original continues
with four further parts about $\Phi(\overline{B_1})$ when $\Phi$ fixes the unit circle.}
\end{exercise}


% Originally: content/week-13.tex, \section{Area via Green's formula} (Corsin Week 13)

\section{Area via Green's formula}
\label{sec:area_via_green}

% Source: Corsin Nick/Class Notes/Week 13.pdf, pp. 5--6
% Kept inline: solved directly in Corsin Nick/Class Notes/Week 13.pdf, pp. 5--6.
\begin{exercise}[Area as a boundary functional]
\label{ex:area_boundary_functional}
Suppose $\gamma : [0,L) \to \mathbb{R}^2$ traces the boundary of a connected $C^\infty$-domain
$U \subseteq \mathbb{R}^2$ in the positive direction. Express the area of $U$ as a functional
of $\gamma$.
\end{exercise}

\begin{exercisesolution}
By definition, $A(U) = \int_U dx\,dy = \int_U 1\,dx\wedge dy$. By Green's formula
(\cref{thm:green_theorem}, in its $1$-form guise $\int_\Omega d\omega =
\int_{\partial\Omega}\omega$),
\[\int_U 1\,dx\wedge dy = \int_0^L \langle F\circ\gamma(t),\gamma'(t)\rangle\,dt\]
for any $F : \overline U \to \mathbb{R}^2$ satisfying $\partial_xF^2 - \partial_yF^1 = 1$. This
is satisfied by $F(x,y) = \tfrac12(-y,x)$, giving the \newterm{shoelace formula}
(\germanterm{Gau\ss sche Trapezformel})
\[A(\gamma) = \frac12\int_0^L \big[\gamma_1(t)\gamma_2'(t) - \gamma_2(t)\gamma_1'(t)\big]\,dt.\]
\end{exercisesolution}

\begin{example}[Ellipse]
\label{ex:area_ellipse}
The path $\gamma : [0,2\pi) \to \mathbb{R}^2$, $t \mapsto (a\cos t,b\sin t)$, traces the
boundary of an ellipse. By \cref{ex:area_boundary_functional},
\[
A(\gamma) = \frac12\int_0^{2\pi} \big[a\cos t\cdot b\cos t - b\sin t\cdot(-a\sin t)\big]\,dt
  = \frac{ab}{2}\int_0^{2\pi} dt = \pi ab.
\]
For the circle $a=b=:r$, this gives the familiar $A = \pi r^2$.
\end{example}

\begin{ainote}
Corsin adds, in the same breath as the ellipse's area: \qt{Famously, there exists no closed
form solution for the length of the boundary of an ellipse} --- the ellipse's perimeter is a
complete elliptic integral of the second kind, not an elementary function, unlike its area.
\end{ainote}


% --- exercises relocated here from the dissolved sheets ---

% Originally: content/week-12.tex, \exercisesheet{12}, problem 12.3
% Source: exercises/Ex12_Analysis2_eng.pdf

\begin{exercise}[12.3 --- Work along paths and Stokes \textnormal{(important)}]
\label{ex:12.3}
Let $F : \mathbb{R}^3 \setminus (\{(0,0)\}\times\mathbb{R}) \to \mathbb{R}^3$ be the vector
field defined by
\[F(x,y,z) = \left(\frac{2(xz+y)}{x^2+y^2},\ \frac{2(yz-x)}{x^2+y^2},\ \log(x^2+y^2)\right)\]
and let $\gamma_1,\gamma_2,\gamma_3 : [0,2\pi] \to \mathbb{R}^3$ be the paths
\[\gamma_1(t) = (\cos t,\sin t,0), \qquad \gamma_2(t) = (2\cos t,2\sin t,2), \qquad
  \gamma_3(t) = (3\cos t+6,3\sin t,3).\]
\begin{enumerate}[label=\textbf{\arabic*.}]
  \item Describe the domain of definition of $F$. Is it connected? Is it simply connected?
  \item Compute the rotation $\curl(F)$ of $F$ and the path integral $\int_{\gamma_1} F\cdot
  d\gamma_1$ of $F$ along $\gamma_1$. Is $F$ conservative?
  \item Explain how the equation $\int_{\gamma_1} F\cdot d\gamma_1 = \int_{\gamma_2} F\cdot
  d\gamma_2$ follows from Stokes' theorem.
  \item Does Stokes' theorem imply that $\int_{\gamma_1} F\cdot d\gamma_1 = \int_{\gamma_3}
  F\cdot d\gamma_3$ holds?
\end{enumerate}
\end{exercise}

% Source: exercises/Ex12_Analysis2_eng.pdf, p. 2

% Originally: content/week-12.tex, \exercisesheet{12}, problem 12.4
% Source: exercises/Ex12_Analysis2_eng.pdf

\begin{exercise}[12.4 --- True or False \textnormal{(important)}]
\label{ex:12.4}
Let $U := \mathbb{R}^3\setminus\{0\}$ and $F \in C^1(U,\mathbb{R}^3)$. For $r>0$ set
\[I_r := \int_{\partial B_r} \langle F,n\rangle\,dS\]
where $n$ denotes the outward normal. Which of the following statements are always true?
\begin{enumerate}[label=\textbf{(\alph*)}]
  \item If $F$ is divergence-free, then $I_r=0$ for all $r>0$.
  \item If $F$ is divergence-free, then the value of $I_r$ does not depend on $r$.
  \item If $F$ is curl-free, then the value of $I_r$ does not depend on $r$.
  \item If $F = \curl(G)$ for a vector field $G \in C^1(U,\mathbb{R}^3)$ and $\lvert F(x)
  \rvert \leq C/\lvert x\rvert$, then $I_r=0$ for all $r>0$.
\end{enumerate}
\end{exercise}

% Source: exercises/Ex12_Analysis2_eng.pdf, p. 2

% Originally: content/week-12.tex, \exercisesheet{12}, problem 12.5
% Source: exercises/Ex12_Analysis2_eng.pdf

\begin{exercise}[12.5 --- Conservative vector fields \textnormal{(important)}]
\label{ex:12.5}
For which values of $\lambda \in \mathbb{R}$ is the vector field $F : \mathbb{R}^2 \to
\mathbb{R}^2$ defined by
\[F(x,y) = \begin{pmatrix} \lambda x e^y \\ (y+1+x^2)e^y \end{pmatrix}\]
conservative? Determine a potential for $F$ for these values.
\end{exercise}

% Source: exercises/Ex12_Analysis2_eng.pdf, p. 3

% Originally: content/week-12.tex, \exercisesheet{12}, problem 12.8
% Source: exercises/Ex12_Analysis2_eng.pdf

\begin{exercise}[12.8 --- Multiple Choice \textnormal{(important)}]
\label{ex:12.8}
Let $Q = [a,b]\times[c,d] \subseteq \mathbb{R}^2$ and $f \in C(Q,\mathbb{R}^2)$. Let $n$ be an
outward normal field on $Q$ and
\[I_b = \int_{[a,b]\times\{c\}} \langle f,n\rangle, \qquad I_t = \int_{[a,b]\times\{d\}}
  \langle f,n\rangle, \qquad I_l = \int_{\{a\}\times[c,d]} \langle f,n\rangle, \qquad
  I_r = \int_{\{b\}\times[c,d]} \langle f,n\rangle.\]
Which of the following expressions corresponds to the outer flux of $f$ through $\partial Q$?
\begin{enumerate}[label=\textbf{(\alph*)}]
  \item $I_b - I_t + I_r - I_l$
  \item $I_b + I_t + I_r + I_l$
  \item $I_b + I_t - I_r - I_l$
  \item $-I_b + I_t - I_r + I_l$
\end{enumerate}
\end{exercise}

% Source: exercises/Ex12_Analysis2_eng.pdf, p. 3

% Originally: content/week-12.tex, \section{Solutions}

\section{Solutions}
\label{sec:stokes_solutions}

\begin{exercisesolution}[Solution of \cref{ex:12.3}]
\textbf{1.} The domain is $\{(x,y,z) \in \mathbb{R}^3 : (x,y) \neq (0,0)\}$, i.e.\
$\mathbb{R}^3$ with the $z$-axis removed. It is connected (path-connected: any two points can
be joined by a path avoiding the axis), but not simply connected: $\gamma_1$, for instance, which
winds once around the axis, cannot be contracted to a point without crossing it. (It is also
not convex: the segment from $(-1,0,0)$ to $(1,0,0)$ leaves the domain.)

\textbf{2.} Writing $F = (F^1,F^2,F^3)$ with $F^1 = 2(xz+y)/(x^2+y^2)$, $F^2 =
2(yz-x)/(x^2+y^2)$, $F^3 = \log(x^2+y^2)$, a direct computation of each pair of mixed
partials gives
\[\partial_y F^3 = \partial_z F^2 = \frac{2y}{x^2+y^2}, \qquad
  \partial_z F^1 = \partial_x F^3 = \frac{2x}{x^2+y^2}, \qquad
  \partial_x F^2 = \partial_y F^1 = \frac{2(x^2-y^2-2xyz)}{(x^2+y^2)^2},\]
so $\curl F = (\partial_yF^3-\partial_zF^2,\ \partial_zF^1-\partial_xF^3,\ \partial_xF^2-
\partial_yF^1) = (0,0,0)$: $F$ is curl-free. On $\gamma_1$, $x=\cos t$, $y=\sin t$, $z=0$,
$x^2+y^2=1$, so $F(\gamma_1(t)) = (2\sin t,\,-2\cos t,\,0)$ and $\gamma_1'(t) = (-\sin t,
\cos t,0)$, giving
\[\int_{\gamma_1} F\cdot d\gamma_1 = \int_0^{2\pi} \big(2\sin t\cdot(-\sin t) + (-2\cos t)
  \cdot\cos t\big)\,dt = \int_0^{2\pi} -2\,dt = -4\pi.\]
Since this is nonzero on a closed loop, $F$ is \textbf{not} conservative. It is curl-free, but
its domain is not simply connected (part 1), so curl-free does not force conservative here.

\textbf{3.} Let $S := \{(x,y,z) : \sqrt{x^2+y^2} = 1+z/2,\ 0\leq z\leq2\}$, the cone swept out
by linearly interpolating the radius from $\gamma_1$ ($z=0$, radius $1$) to $\gamma_2$ ($z=2$,
radius $2$), parametrized by $H(s,t) := \big((1+s)\cos t,(1+s)\sin t,2s\big)$ for $(s,t) \in
[0,1]\times[0,2\pi]$. Since the radius $1+s \geq 1 > 0$ throughout, $S$ never meets the
$z$-axis and lies entirely in the domain of $F$; its oriented boundary (outward normal) is
$\gamma_2$ minus $\gamma_1$. Since $\curl F = 0$ on $S$ by part 2, \cref{thm:stokes_classical}
gives
\[\int_{\gamma_2} F\cdot d\gamma_2 - \int_{\gamma_1} F\cdot d\gamma_1 = \int_{\partial S}
  F\cdot ds = \int_S \curl F\cdot d\vec n = 0,\]
i.e.\ $\int_{\gamma_1} F\cdot d\gamma_1 = \int_{\gamma_2} F\cdot d\gamma_2$.

\textbf{4.} No. The disk $D := \{(x,y,z) : (x-6)^2+y^2 \leq 9,\ z=3\}$, oriented upward, has
oriented boundary exactly $\gamma_3$, and $D$ lies entirely in the domain of $F$ (it does not
meet the $z$-axis, since its center $(6,0,3)$ is at distance $6$ from the axis and its radius
is only $3$). By \cref{thm:stokes_classical} with $\curl F=0$,
\[\int_{\gamma_3} F\cdot d\gamma_3 = \int_{\partial D} F\cdot ds = \int_D \curl F\cdot d\vec n
  = 0.\]
Since $\int_{\gamma_1} F\cdot d\gamma_1 = -4\pi \neq 0 = \int_{\gamma_3} F\cdot d\gamma_3$, the
two do \emph{not} agree, and Stokes' theorem gives no reason to expect they would: unlike the
pair $(\gamma_1,\gamma_2)$, there is no surface joining $\gamma_1$ and $\gamma_3$ that stays
inside the domain, since $\gamma_1$ winds around the excluded axis and $\gamma_3$ does not.
\end{exercisesolution}

% Verified: solved independently, all four parts match Sol12_Analysis2_eng.pdf exactly.
% Note: the official solution's part 2 labels vanishing curl as "divergence-free" where it
% visibly means curl-free (computes curl F, not div F) --- a wording slip, not repeated here.

\begin{exercisesolution}[Solution of \cref{ex:12.4}]
\textbf{(a) False.} The Newton potential's gradient $F(x) = \nabla(1/\lvert x\rvert) =
-x\lvert x\rvert^{-3}$ is divergence-free on $U$ (it is $V_{-3}$ up to sign, from
\cref{ex:11.5}, with $n=3$), yet $I_r = \int_{\partial B_r} F\cdot n\,dS =
-\vol_2(\partial B_r) \neq 0$.

\textbf{(b) True.} Applying \cref{thm:gauss_divergence} on the shell $B_{r_2}\setminus
B_{r_1} \Subset U$ for $0<r_1<r_2$, whose boundary is $\partial B_{r_2}$ (outward) together
with $\partial B_{r_1}$ with reversed orientation,
\[0 = \int_{B_{r_2}\setminus B_{r_1}} \divg F\,dx = I_{r_2} - I_{r_1},\]
so $I_r$ is independent of $r$.

\textbf{(c) False.} Take $F(x) = x = \tfrac12\nabla\lvert x\rvert^2$, which is curl-free
(being a gradient). Then $\divg F \equiv 3$, so by \cref{thm:gauss_divergence}, $I_r =
\int_{B_r} 3\,dx = 3\vol_3(B_r) = 4\pi r^3$, which depends on $r$.

\textbf{(d) True.} If $F = \curl G$ then $\divg F = \divg(\curl G) = 0$ (a standard vector
identity), so part (b) applies and $I_r$ is independent of $r$. On the other hand,
\[\lvert I_r\rvert \leq \int_{\partial B_r} \lvert\langle F,n\rangle\rvert\,dS \leq
  \frac{C}{r}\,\vol_2(\partial B_r) = \frac{C}{r}\cdot4\pi r^2 = 4\pi C r \longrightarrow 0
  \quad\text{as } r\to0^+.\]
Since $I_r$ is constant in $r$ and this constant is a limit of values tending to $0$, $I_r=0$
for all $r>0$.
\end{exercisesolution}

% Verified: solved independently, all four parts match Sol12_Analysis2_eng.pdf.

\begin{exercisesolution}[Solution of \cref{ex:12.5}]
Writing $F = (F^1,F^2)$ with $F^1 = \lambda xe^y$, $F^2 = (y+1+x^2)e^y$, the integrability
condition is $\partial_yF^1 = \partial_xF^2$:
\[\partial_yF^1 = \lambda xe^y, \qquad \partial_xF^2 = 2xe^y.\]
Since $\mathbb{R}^2$ is convex (so integrability is equivalent to conservativity here), $F$
is conservative if and only if $\lambda=2$.

For $\lambda=2$, seek $f$ with $\partial_xf = F^1 = 2xe^y$: integrating in $x$,
$f(x,y) = x^2e^y + g(y)$ for some function $g$. Then $\partial_yf = x^2e^y+g'(y)$ must equal
$F^2 = (y+1+x^2)e^y$, so $g'(y) = (y+1)e^y$; integrating by parts, $\int(y+1)e^y\,dy =
ye^y - e^y + e^y = ye^y$ (up to an additive constant). Hence
\[f(x,y) = \boxed{\ (x^2+y)e^y\ } \qquad (\text{unique up to an additive constant}),\]
which one checks directly satisfies $\nabla f = F$.
\end{exercisesolution}

% Verified: solved independently, lambda=2 and potential (x^2+y)e^y match Sol12 exactly.

\begin{exercisesolution}[Solution of \cref{ex:12.8}]
The four quantities $I_b,I_t,I_l,I_r$ are already defined using the same outward normal field
$n$ on each of the four edges of $\partial Q$, so the total outer flux is simply their sum,
with no relative sign corrections needed: $\boxed{I_b+I_t+I_r+I_l}$, answer \textbf{(b)}.
\end{exercisesolution}

% Verified: solved independently, matches Sol12 answer (b).
% Correction: Claude Opus 5 (High) --- this comment belongs to ex:12.8 above. Flash appended its
% own solution between the two, leaving the comment claiming a block it had never checked.

\begin{exercisesolution}[Proof of \cref{ex:curl_identity_greens_formula}]
% Generator: Gemini 3.6 Flash (High)
% Correction: Claude Opus 5 (High) --- \operatorname{rot} -> \curl, J\Phi -> \Jac\Phi, and the line
% integrals rewritten in this chapter's own notation (\int_{\partial\Omega} F\cdot\tau\,dL, as in
% thm:green_theorem) rather than the physics d\mathbf{r}. Part (b) also justified properly: Flash
% asserted "d\Phi_2 = d\Psi_2 on \partial\Omega", which is false for the full differentials --- only
% the tangential derivative agrees, and that is all the line integral sees.
\begin{enumerate}[label=\textbf{(\alph*)}]
  \item Write $V = (V_1, V_2)\transp = (\Phi_1 \partial_1 \Phi_2, \Phi_1 \partial_2 \Phi_2)\transp$.
The scalar curl in the plane is $\curl V = \partial_1 V_2 - \partial_2 V_1$, so the product rule
gives
\begin{align*}
\curl V &= \partial_1 (\Phi_1 \partial_2 \Phi_2) - \partial_2 (\Phi_1 \partial_1 \Phi_2) \\
&= (\partial_1 \Phi_1 \partial_2 \Phi_2 + \Phi_1 \partial_1 \partial_2 \Phi_2) - (\partial_2 \Phi_1 \partial_1 \Phi_2 + \Phi_1 \partial_2 \partial_1 \Phi_2).
\end{align*}
Since $\Phi \in C^2(\overline{\Omega}, \mathbb{R}^2)$, Schwarz's lemma
(\cref{thm:schwarz_mixed_partials}) gives $\partial_1 \partial_2 \Phi_2 = \partial_2 \partial_1
\Phi_2$, and the two second-derivative terms cancel. Consequently,
\[\curl V = \partial_1 \Phi_1 \partial_2 \Phi_2 - \partial_2 \Phi_1 \partial_1 \Phi_2 = \det \begin{pmatrix} \partial_1 \Phi_1 & \partial_2 \Phi_1 \\ \partial_1 \Phi_2 & \partial_2 \Phi_2 \end{pmatrix} = \det \Jac\Phi.\]
Note that the hypothesis $\Phi \in C^2$ is doing real work here: without equality of mixed partials
the identity is simply false.

  \item Write $V_\Phi := \Phi_1 \nabla \Phi_2$ and $V_\Psi := \Psi_1 \nabla \Psi_2$. By part
\textbf{(a)}, $\curl V_\Phi = \det \Jac\Phi$ and $\curl V_\Psi = \det \Jac\Psi$, so Green's theorem
(\cref{thm:green_theorem}) turns both integrals into boundary integrals:
\[\int_\Omega \det \Jac\Phi \, dx = \int_{\partial \Omega} V_\Phi \cdot \tau \, dL,
\qquad
\int_\Omega \det \Jac\Psi \, dx = \int_{\partial \Omega} V_\Psi \cdot \tau \, dL.\]
It remains to see that the two boundary integrands agree. Fix a point of $\partial\Omega$ and let
$\tau$ be the unit tangent there. Since $\Phi = \Psi$ on all of $\partial\Omega$, the two maps agree
along the boundary curve, and therefore so do their derivatives \emph{in the direction $\tau$}:
\[\nabla \Phi_2 \cdot \tau = \partial_\tau \Phi_2 = \partial_\tau \Psi_2 = \nabla \Psi_2 \cdot \tau .\]
(The full gradients need not agree, since nothing constrains the normal derivative; but the line
integral only ever sees the tangential part.) Together with $\Phi_1 = \Psi_1$ on $\partial\Omega$
this gives $V_\Phi \cdot \tau = V_\Psi \cdot \tau$ pointwise on $\partial\Omega$, and hence
\[\int_\Omega \det \Jac\Phi \, dx = \int_\Omega \det \Jac\Psi \, dx.\]
\end{enumerate}
\begin{remark}
The statement is a genuinely striking one: $\int_\Omega \det\Jac\Phi$ is an interior quantity, the
signed area swept out by $\Phi$, yet it depends on nothing but the boundary values of $\Phi$. That
is the same phenomenon as a conservative field having path-independent work, one dimension up.
\end{remark}
\end{exercisesolution}

\begin{exercisesolution}[Solution of \cref{ex:angle_field_true_false}]
% Generator: Claude Opus 5 (High)
Begin with the computation every part relies on. Writing $X = (P,Q)$ with $P = -y/(x^2+y^2)$ and
$Q = x/(x^2+y^2)$,
\[\partial_x Q = \frac{(x^2+y^2) - 2x^2}{(x^2+y^2)^2} = \frac{y^2-x^2}{(x^2+y^2)^2},
  \qquad
  \partial_y P = \frac{-(x^2+y^2) + 2y^2}{(x^2+y^2)^2} = \frac{y^2-x^2}{(x^2+y^2)^2},\]
so $\curl X = \partial_x Q - \partial_y P = 0$ throughout $U$. The field is closed. Whether it is
also a gradient is precisely what the domain decides.
\begin{enumerate}[label=\textbf{(\alph*)}]
  \item \textbf{False.} Suppose such a $u$ existed. Then the work of $X$ around any closed loop in
    $U$ would vanish. But along the unit circle $\sigma(t) := (\cos t, \sin t)$, $t \in [0,2\pi]$,
    we have $X(\sigma(t)) = (-\sin t, \cos t) = \dot\sigma(t)$, so the integrand is
    $\lvert\dot\sigma\rvert^2 = 1$ and
\[\int_\sigma X \cdot d\sigma = \int_0^{2\pi} 1 \, dt = 2\pi \neq 0 .\]
    So no potential exists on $U$. Being closed is not enough; $U$ is not simply connected, and this
    field is the standard witness that the Poincaré lemma (\cref{thm:poincare_lemma}) genuinely needs
    that hypothesis.
  \item \textbf{True.} The half-plane $V$ is convex, hence simply connected, so a closed field on it
    is exact (\cref{cor:potentials_simply_connected}). Explicitly, $v(x,y) := \arctan(y/x)$ is $C^1$
    on $V$ --- note $x > 0$ there, so the quotient is defined and smooth --- and
\[\nabla v = \frac{1}{1 + y^2/x^2}\left(-\frac{y}{x^2},\; \frac{1}{x}\right)
   = \frac{1}{x^2+y^2}\,(-y,\, x) = X .\]
    The potential is the polar angle, which is single-valued on the half-plane but not on $U$: that
    is \textbf{(a)} and \textbf{(b)} in one sentence.
  \item \textbf{False.} Injectivity on $[0,1)$ makes $\gamma$ a simple closed curve, but it says
    nothing about \emph{where} that curve runs. Two things can go wrong. First, the loop need not
    enclose the origin at all: the circle of radius $\tfrac12$ about $(0,1)$ is a simple closed curve
    in $U$ through $(0,1)$, and since it bounds a disc on which $X$ is defined and closed, Green's
    theorem (\cref{thm:green_theorem}) gives $\int_\gamma X \cdot d\gamma = 0$. Second, even a loop
    that does enclose the origin may be traversed clockwise, giving $-2\pi$. The correct general
    statement is that the integral equals $2\pi$ times the winding number of $\gamma$ about the
    origin, and \qt{simple closed curve} pins that number down only up to sign, and only once the
    origin is known to be inside.
\end{enumerate}
\begin{remark}
This one field carries most of the chapter's content. It is closed but not exact, so it separates the
two notions; the failure is measured by an integer, the winding number; and the obstruction is a
property of the \emph{domain}, not of the field --- restrict it to any simply connected piece of $U$,
as in \textbf{(b)}, and a potential appears at once.
\end{remark}
\end{exercisesolution}



% Appendices. \appendix switches chapter numbering to letters, so the glossary becomes
% Appendix B -- matching the file name, which predates the file itself (gemini.md has always
% referred to content/appendix-b-glossary.tex as the destination for German term pairs).
\appendix
% Source: Sascha Brack/Class Notes/Week_12_Notes_Friday.pdf, pp. 2--5
\chapter{Ordinary Differential Equations}
\label{ch:odes}

\begin{ainote}
This appendix compiles the material on Ordinary Differential Equations (ODEs). During the semester, this topic was covered in supplementary sessions.
\end{ainote}

\section{Turning an \texorpdfstring{$n$}{n}-th order ODE into an \texorpdfstring{$n \times n$}{n x n} system}
\label{sec:nth_order_to_system}

Given an $n$-th order linear ODE, we can rewrite it as a first-order problem in higher dimensions.
Consider the example equation:
\[y(t) + y'(t) + 2y''(t) = f(t).\]
We define $u_1 := y$, which gives:
\[u_1 + u_1' + 2u_1'' = f.\]
Letting $u_2 := u_1'$, we can rewrite the second derivative as $u_1'' = u_2'$. The system becomes:
\begin{align*}
    u_1 + u_2 + 2u_2' &= f \\
    u_2 &= u_1'.
\end{align*}
Solving for the derivatives $u_1'$ and $u_2'$, we get:
\begin{align*}
    u_1' &= u_2 \\
    u_2' &= -\frac{1}{2}u_1 - \frac{1}{2}u_2 + \frac{1}{2}f.
\end{align*}
In matrix notation, this is:
\[\begin{pmatrix} u_1' \\ u_2' \end{pmatrix} = \begin{pmatrix} 0 & 1 \\ -1/2 & -1/2 \end{pmatrix} \begin{pmatrix} u_1 \\ u_2 \end{pmatrix} + \begin{pmatrix} 0 \\ f/2 \end{pmatrix},\]
or more compactly,
\[u'(t) = A u(t) + F(t).\]
This reduction to a first-order system can be done with any $n$-th order ODE. We therefore focus on solving systems of the form $u'(t) = A u(t) + F(t)$.

\section{Solving linear systems of ODEs}
\label{sec:solving_linear_systems_odes}

\subsection{Diagonalizable matrices}
We first look at the homogeneous case where $F \equiv 0$, so $u'(t) = A u(t)$.
Suppose the matrix $A$ is diagonalizable, meaning we can write $A = P D P^{-1}$ where $D$ is a diagonal matrix and $P$ is invertible:
\[D = \begin{pmatrix} \lambda_1 & & 0 \\ & \ddots & \\ 0 & & \lambda_n \end{pmatrix}.\]
Then we have:
\begin{align*}
    u'(t) &= A u(t) \\
    \iff u'(t) &= P D P^{-1} u(t) \\
    \iff P^{-1} u'(t) &= D P^{-1} u(t) \\
    \iff (P^{-1} u(t))' &= D (P^{-1} u(t)).
\end{align*}
By defining a new variable $v(t) := P^{-1} u(t)$, the system uncouples:
\[v'(t) = D v(t) \iff \begin{pmatrix} v_1'(t) \\ \vdots \\ v_n'(t) \end{pmatrix} = \begin{pmatrix} \lambda_1 & & 0 \\ & \ddots & \\ 0 & & \lambda_n \end{pmatrix} \begin{pmatrix} v_1(t) \\ \vdots \\ v_n(t) \end{pmatrix}.\]
This yields $n$ independent $1$-dimensional equations $v_i'(t) = \lambda_i v_i(t)$, which have the solutions:
\[v_i(t) = C_i e^{\lambda_i t} \quad \text{for } i \in \{1,\dots,n\}.\]
The solution to the original system is then recovered by $u(t) = P v(t)$. By diagonalizing, we were able to uncouple the $n \times n$ system and solve the scalar equations without a problem.

\subsection{Non-diagonalizable matrices}
What if $A$ is not diagonalizable? Then we have to use the Jordan Normal Form: $A = P J P^{-1}$.
Following the same substitution $v(t) = P^{-1}u(t)$, we obtain $v'(t) = J v(t)$.
We now focus on each Jordan block in $J$. For a block with eigenvalue $\lambda_i$, the system looks like:
\[\begin{pmatrix} v_j'(t) \\ \vdots \\ v_{j+d}'(t) \end{pmatrix} = \begin{pmatrix} \lambda_i & 1 & & \\ & \lambda_i & \ddots & \\ & & \ddots & 1 \\ & & & \lambda_i \end{pmatrix} \begin{pmatrix} v_j(t) \\ \vdots \\ v_{j+d}(t) \end{pmatrix}.\]
This translates to the equations:
\begin{align*}
    v_j'(t) &= \lambda_i v_j(t) + v_{j+1}(t) \\
    &\vdots \\
    v_{j+d}'(t) &= \lambda_i v_{j+d}(t).
\end{align*}
We first solve for $v_{j+d}(t)$ from the last equation. Then we plug it into the preceding equation to find $v_{j+d-1}(t)$, and continue this process upwards.

Since this back-substitution can be time-consuming, it is useful to use matrix exponentials. For a system $u'(t) = A u(t)$, the formal solution is:
\[u(t) = e^{tA} u_0,\]
where $u_0$ is given by the initial conditions. Evaluating the matrix exponential $e^{tA}$ provides a direct way to solve the system.

\begin{example}[Solving a $2 \times 2$ system using matrix exponentials]
\label{ex:ode_matrix_exponential}
Consider the system given by
\begin{align*}
    u_1(t) + 2u_2(t) &= u_1'(t) \\
    2u_1(t) + u_2(t) &= u_2'(t).
\end{align*}
This can be written as $u'(t) = A u(t)$ with $A = \begin{pmatrix} 1 & 2 \\ 2 & 1 \end{pmatrix}$.
The eigenvalues of $A$ are found via $\det(A - \lambda I) = (1-\lambda)^2 - 4 = 0$, giving $\lambda_1 = 3$ and $\lambda_2 = -1$.
The corresponding eigenvectors are $w_1 = (1, 1)\transp$ and $w_2 = (1, -1)\transp$.
Since $A$ is diagonalizable, we have $A = P D P^{-1}$ with $P = \begin{pmatrix} 1 & 1 \\ 1 & -1 \end{pmatrix}$ and $D = \begin{pmatrix} 3 & 0 \\ 0 & -1 \end{pmatrix}$.
The general solution is then given by:
\[u(t) = \sum_{i=1}^2 C_i e^{\lambda_i t} w_i = C_1 e^{3t} \begin{pmatrix} 1 \\ 1 \end{pmatrix} + C_2 e^{-t} \begin{pmatrix} 1 \\ -1 \end{pmatrix}.\]

Now consider a non-diagonalizable example where the system is already in Jordan normal form, $u'(t) = J u(t)$ with $J = \begin{pmatrix} 1 & 1 \\ 0 & 1 \end{pmatrix}$.
Using the matrix exponential, the solution is $u(t) = e^{tJ} u(0)$.
We can split $J = D + N$ where $D = \begin{pmatrix} 1 & 0 \\ 0 & 1 \end{pmatrix}$ is diagonal and $N = \begin{pmatrix} 0 & 1 \\ 0 & 0 \end{pmatrix}$ is nilpotent ($N^2 = 0$).
Since $D$ and $N$ commute, $e^{tJ} = e^{tD} e^{tN}$.
We calculate:
\[e^{tD} = \begin{pmatrix} e^t & 0 \\ 0 & e^t \end{pmatrix}, \qquad e^{tN} = I + tN = \begin{pmatrix} 1 & t \\ 0 & 1 \end{pmatrix}.\]
Multiplying these gives:
\[e^{tJ} = \begin{pmatrix} e^t & t e^t \\ 0 & e^t \end{pmatrix}.\]
Thus, the solution is $u_1(t) = e^t u_1(0) + t e^t u_2(0)$ and $u_2(t) = e^t u_2(0)$.
\end{example}

\section{Inhomogeneous linear ODEs}
\label{sec:inhomogeneous_odes}

Given the $n$-th order linear inhomogeneous ODE
\[Ly(x) = b(x), \qquad \text{where } L = \sum_{k=0}^n a_k \frac{d^k}{dx^k},\]
the general solution is always of the form
\[y(x) = y_{\text{hom}}(x) + y_{\text{part}}(x).\]

The homogeneous solution $y_{\text{hom}}(x)$ fulfills $Ly_{\text{hom}}(x) = 0$, or $\sum_{k=0}^n a_k y_{\text{hom}}^{(k)}(x) = 0$. This is solved using the general Ansatz $y_{\text{hom}}(x) = C e^{rx}$, which leads to the characteristic equation:
\[\sum_{k=0}^n a_k C r^k e^{rx} = 0 \implies \sum_{k=0}^n a_k r^k = 0,\]
yielding roots $r_1, \dots, r_n$ (assuming they are all distinct).

The particular solution $y_{\text{part}}(x)$ solves $Ly_{\text{part}}(x) = b(x)$ and is found using a fitting Ansatz, often an educated guess based on the form of $b(x)$ (e.g., if $b(x)$ is a constant, consider a constant or polynomial Ansatz).

\subsection{Duhamel's Principle}
\label{sec:duhamels_principle}

For a system of inhomogeneous linear ODEs $u'(t) = A u(t) + f(t)$, the solution can be constructed by combining the homogeneous solution and a particular solution obtained via Duhamel's principle:
\[u(t) = \underbrace{e^{At} u(0)}_{\text{homogeneous solution}} + \underbrace{\int_0^t e^{(t-s)A} f(s) \,ds}_{\text{inhomogeneous part}}.\]
This elegantly generalizes the $1$-dimensional case where the solution to $u'(t) = a u(t) + f(t)$ is given by $u(t) = e^{at} u(0) + \int_0^t e^{a(t-s)} f(s) \,ds$.

\section{The Cauchy--Lipschitz / Picard--Lindelöf theorem}
\label{sec:cauchy_lipschitz}

The Cauchy--Lipschitz theorem (also known as the Picard--Lindelöf theorem) guarantees the existence and uniqueness of solutions to first-order initial value problems, provided the right-hand side is Lipschitz continuous with respect to the state variable.
For an equation of the form $y'(x) = F(x, y(x))$, one first puts the differential equation in normal form to determine if the theorem applies, which requires checking the properties of $F$.

% Generator: Claude Opus 5 (High)
% Glossary of the English/German term pairs mirrored on first introduction throughout the
% document, as required by the "German mirroring" rule in gemini.md. Every entry below
% corresponds to exactly one \newterm{...} (\germanterm{...}) pair in content/week-NN.tex,
% and the "Week" column records where that first introduction happens.
%
% MAINTENANCE: when adding a \germanterm{...} anywhere, add its row here too. To list every
% pair currently in the document and compare against this table:
%   grep -ohE '\\germanterm\{[^}]*\}' content/week-*.tex | sort -u
% A German term must appear exactly once across all week files (first introduction only);
% duplicates are a style error, not a harmless repetition.

\chapter{Glossary --- English and German terms}
\label{ch:glossary}

The lecture and the exercise sheets for this course exist in both English and German, and the
two vocabularies do not always line up in the obvious way --- \germanterm{Abschluss} is the
closure and not a \qt{conclusion}, \germanterm{Rand} is the boundary and not a \qt{rim}, and
\germanterm{Satz von Schwarz} is what English-language texts usually call Clairaut's theorem.
Each term below is mirrored once in the main text, at the point where it is first introduced;
this appendix collects those pairs in one place so the document can be used alongside German
material.

Canonical German wording follows Jérôme Paschoud's topic-named files. The \textbf{Week} column
gives the chapter of first introduction.

\section*{A --- E}

\begin{center}
\begin{tabular}{@{}p{0.40\linewidth} p{0.44\linewidth} c@{}}
\textbf{English} & \textbf{German} & \textbf{Week} \\
\hline
accumulation point & Häufungspunkt & 2 \\
antisymmetric $k$-linear form & antisymmetrische $k$-lineare Form & 11 \\
boundary & Rand & 2 \\
bounded & beschränkt & 3 \\
chain rule & Kettenregel & 4 \\
closed & abgeschlossen & 2 \\
closure & Abschluss & 2 \\
compact & kompakt & 2 \\
complete metric space & vollständiger metrischer Raum & 2 \\
connected & zusammenhängend & 3 \\
continuity equation & Kontinuitätsgleichung & 10 \\
continuous & stetig & 2 \\
contraction & Kontraktion & 3 \\
convex & konvex & 6 \\
critical point & kritischer Punkt & 5 \\
cuspidal cubic & Neilsche Parabel & 5 \\
degenerate (critical point) & entartet & 5 \\
diffeomorphism & Diffeomorphismus & 6 \\
differentiable & differenzierbar & 4 \\
differential (the) & das Differential & 4 \\
differential $p$-form & Differentialform & 11 \\
directional derivative & Richtungsableitung & 4 \\
disconnected & unzusammenhängend & 3 \\
\end{tabular}
\end{center}

\section*{F --- L}

\begin{center}
\begin{tabular}{@{}p{0.40\linewidth} p{0.44\linewidth} c@{}}
\textbf{English} & \textbf{German} & \textbf{Week} \\
\hline
figure eight & Lemniskate & 8 \\
fixed point & Fixpunkt & 3 \\
flux & Fluss & 10 \\
Frobenius norm & Frobeniusnorm & 2 \\
Gabriel's horn & Torricellis Trompete & 10 \\
Gram determinant & Gram-Determinante & 9 \\
gradient & Gradient & 4 \\
Hessian matrix & Hesse-Matrix & 4 \\
Hölder's inequality & Höldersche Ungleichung & 1 \\
homeomorphism & Homöomorphismus & 6 \\
Hurwitz criterion & Hurwitz-Kriterium & 5 \\
isometry & Isometrie & 9 \\
Jacobi matrix & Jacobi-Matrix & 4 \\
Jordan measurable & Jordan-messbar & 8 \\
Lagrangian & Lagrange-Funktion & 5 \\
Lebesgue number & Lebesgue-Zahl & 2 \\
Lebesgue's covering lemma & Lebesguesches Überdeckungslemma & 2 \\
length & Länge & 9 \\
level set & Niveaumenge & 5 \\
Lipschitz continuous & Lipschitz-stetig & 3 \\
local minimum & lokales Minimum & 5 \\
\end{tabular}
\end{center}

\section*{M --- P}

\begin{center}
\begin{tabular}{@{}p{0.40\linewidth} p{0.44\linewidth} c@{}}
\textbf{English} & \textbf{German} & \textbf{Week} \\
\hline
metric & Metrik & 2 \\
metric space & metrischer Raum & 2 \\
multi-index & Multiindex & 4 \\
negative definite & negativ definit & 5 \\
negative semidefinite & negativ semidefinit & 6 \\
norm & Norm & 2 \\
normal space & Normalraum & 7 \\
open & offen & 2 \\
open cover & offene Überdeckung & 2 \\
orientable & orientierbar & 12 \\
partial derivative & partielle Ableitung & 4 \\
path & Weg & 3 \\
path-connected & wegzusammenhängend & 3 \\
positive definite & positiv definit & 5 \\
positive semidefinite & positiv semidefinit & 6 \\
pseudo-metric & Pseudometrik & 2 \\
pullback & Zurückziehung & 12 \\
\end{tabular}
\end{center}

\section*{Q --- Z}

\begin{center}
\begin{tabular}{@{}p{0.40\linewidth} p{0.44\linewidth} c@{}}
\textbf{English} & \textbf{German} & \textbf{Week} \\
\hline
saddle point & Sattelpunkt & 5 \\
scalar field & Skalarfeld & 4 \\
inner product / scalar product & Skalarprodukt & 2 \\
Schwarz's theorem & Satz von Schwarz & 5 \\
sequence & Folge & 1 \\
shadow price & Schattenpreis & 5 \\
space & Raum & 2 \\
submanifold & Untermannigfaltigkeit & 7 \\
submultiplicative & submultiplikativ & 3 \\
support (of a partition function) & Träger & 12 \\
tangent space & Tangentialraum & 7 \\
totally bounded & total beschränkt & 2 \\
trapezoidal rule (Gauss) & Gaußsche Trapezformel & 13 \\
uniform convergence & gleichmässige Konvergenz & 2 \\
uniformly continuous & gleichmässig stetig & 2 \\
velocity (of a curve) & Geschwindigkeit & 4 \\
wedge product & Dachprodukt & 11 \\
interior & Inneres & 2 \\
\end{tabular}
\end{center}

\begin{ainote}
Three conventions in this table are worth flagging, because they are places where the German
and English literature genuinely diverge rather than merely translate.
\begin{itemize}
  \item \textbf{\germanterm{Satz von Schwarz}} is the symmetry of mixed partial derivatives
    (\cref{thm:schwarz_mixed_partials}). English texts more often call it \emph{Clairaut's
    theorem}; both names are used in this document, with Schwarz first because that is what the
    lecture uses.
  \item \textbf{\germanterm{Neilsche Parabel}} names the curve $y^2 = x^3$ after William Neile.
    The English \emph{cuspidal cubic} is descriptive rather than eponymous, so searching for one
    name will not find the other.
  \item \textbf{\germanterm{gleichmässig}} carries both English senses --- it is
    \qt{uniform} in \emph{uniform convergence} and in \emph{uniformly continuous}, which are
    different notions (\cref{ex:uniform_convergence}, \cref{def:uniformly_continuous}). The
    German gives no hint that these are distinct, and neither does the English.
\end{itemize}
\end{ainote}

% ============================================================
% Appendix --- Mid-semester repetition quiz
% Originally: content/week-08.tex, \section{Repetition quiz} (Corsin Week 8).
% It reviews the whole differential-topology half of the course rather than any one
% topic, so under a topic structure it belongs to no chapter and is collected here.
% ============================================================
\chapter{Mid-Semester Repetition Quiz}
\label{ch:repetition_quiz}

% Transition: Claude Opus 5 (Medium)
Corsin devoted the first half of one exercise class to a quiz revisiting everything up to this
point: differentiability, the chain rule, extrema, the implicit function theorem,
submanifolds. It ranges across \crefrange{ch:prerequisites}{ch:submanifolds} and belongs to no
single chapter, so it is collected here, where it also serves as a self-test before the
integration half of the course.

% Source: Corsin Nick/Class Notes/Week 8.pdf, p. 2
\section{The quiz}
\label{sec:repetition_quiz}

\begin{ainote}
Corsin credits the questions to Prof.\ Serra and Prof.\ Lang, and ran them from a shared Kahoot
at the mid-semester repetition session. On his pages, \textcolor{green!60!black}{green} marks the
\emph{correct} statement and \textcolor{red!80!black}{red} an \emph{incorrect} one; that colouring
is rendered below as explicit \textbf{True}/\textbf{False} verdicts.
\end{ainote}

\begin{question}
\label{q:quiz_q1}
Let $U := \{(x,y) \in \mathbb{R}^2 : 1 \leq x^2+y^2 \leq 4\}$ and
$f(x,y) := \sin(xy) - y^4$. Which of the following holds?
\begin{enumerate}[label=\textbf{(\alph*)}]
  \item $f(U) = [a,b]$ for some $a < b$.
  \item $f(U)$ is disconnected and closed with two connected components.
\end{enumerate}
\end{question}

\begin{answer}
\textbf{(a) is true, (b) is false.} $U$ is compact and connected and $f$ is continuous, so
$f(U) \subseteq \mathbb{R}$ is compact (\cref{item:continuity_preserves_compact}) and connected
(\cref{thm:continuity_preserves_connected}). The connected subsets of $\mathbb{R}$ are exactly
the intervals, and a compact one is closed and bounded, so $f(U) = [a,b]$.
\end{answer}

\begin{ainote}
Worth noting how little computation this needs: $f$ is never evaluated. The annulus $U$ is
connected (it is \emph{not} simply connected, but that is a different property and irrelevant
here), and the two preservation theorems do the rest. Recognising that a question
is asking about preserved properties rather than about values is most of the work.
\end{ainote}

\begin{question}
\label{q:quiz_q2}
Let $X \subseteq \mathbb{R}^2$ with $X \neq \mathbb{R}^2$ and $X \neq \emptyset$. If $X$ is
complete, is it closed?
\end{question}

\begin{answer}
\textbf{True.} Let $(a_n)_{n=0}^\infty$ be a convergent sequence with elements in $X$. Then
$(a_n)$ is Cauchy, and by completeness it has a limit point in $X$. So $X$ is sequentially
closed, hence closed (\cref{prop:open_closed_via_sequences}).
\end{answer}

\begin{ainote}
The limit is taken in $\mathbb{R}^2$, and completeness of $X$ is what forces it back into $X$.
The converse also holds here, but only because the ambient space $\mathbb{R}^2$ is itself
complete: a closed subset of a complete space is complete. In a non-complete ambient space
closedness gives nothing. Note that $\mathbb{Q}$ is closed in $\mathbb{Q}$ and not complete.
\end{ainote}

% Source: Corsin Nick/Class Notes/Week 8.pdf, p. 3
\begin{question}
\label{q:quiz_q3}
Let $X \subseteq \mathbb{R}^2$ as above. If $X$ is not open, is it closed?
\end{question}

\begin{answer}
\textbf{False.}
\end{answer}

% Generator: Claude Opus 5 (Medium) --- Corsin gives the verdict without a counterexample.
\begin{ainote}
Corsin marks this false but gives no counterexample; the half-open square
$X := [0,1] \times (0,1)$ from \cref{sec:open_closed_sets} is one, being neither open
nor closed. \qt{Open} and \qt{closed} are not opposites and not exhaustive: $\emptyset$ and
$\mathbb{R}^2$ are both, most sets are neither, and the excluded values
$X \neq \emptyset,\ X \neq \mathbb{R}^2$ in the hypothesis are there precisely to rule out the
two clopen cases.
\end{ainote}

\begin{question}
\label{q:quiz_q4}
Let
\[f(x,y,z) := \frac{xy^2}{x^2+y^2+z^2} \quad \text{for } (x,y,z) \neq 0, \qquad f(0,0,0) := 0.\]
Is $f \in C^1(\mathbb{R}^3)$?
\end{question}

\begin{answer}
\textbf{False.} We compute
\[\partial_xf = \frac{y^2}{x^2+y^2+z^2} - 2\,\frac{x^2y^2}{(x^2+y^2+z^2)^2}.\]
In spherical coordinates $x = r\sin\theta\cos\varphi$, $y = r\sin\theta\sin\varphi$,
$z = r\cos\theta$, every power of $r$ cancels:
\[\partial_xf = \sin^2\theta\sin^2\varphi - 2\sin^4\theta\cos^2\varphi\sin^2\varphi.\]
The limit $\lim_{r\to0}\partial_xf(r,\theta,\varphi)$ is therefore ill-defined, since it depends
on the direction $(\theta,\varphi)$ along which the origin is approached. Consequently
$\partial_xf$ is not sequentially continuous at $0$.
\end{answer}

\begin{ainote}
The cancellation of $r$ is the whole argument, and it is worth naming the reason: $\partial_xf$
is \newterm{homogeneous of degree $0$}, so it is constant along each ray from the origin and
takes a different constant on different rays. Compare
\Cref{ex:all_directional_derivatives_exist}, where the same trick (restrict to a ray, watch the
parameter cancel) exposed a function whose directional derivatives all existed without it being
differentiable.
\end{ainote}

% Source: Corsin Nick/Class Notes/Week 8.pdf, pp. 3--4
\begin{question}
\label{q:quiz_q5}
Let $f \in C^\infty(\mathbb{R}^3)$ and $\alpha \in \mathbb{R}$ be such that
\[\nabla f(0) = (0,0,0), \qquad
\Hess f(0) = \begin{pmatrix} 1 & \alpha & 0 \\ \alpha & 2 & 0 \\ 0 & 0 & 1\end{pmatrix}.\]
Which of the following is \textbf{false}?
\begin{enumerate}[label=\textbf{(\alph*)}]
  \item For all $\alpha \in \mathbb{R}$, $0$ is not a local maximum point of $f$.
  \item For $\alpha > \sqrt2$, $0$ is a saddle point.
  \item For $\alpha \neq \pm\sqrt2$, $\nabla f : \mathbb{R}^3 \to \mathbb{R}^3$ is a local
    diffeomorphism around $0$.
  \item For $\alpha < \sqrt2$, $0$ is a local minimum point of $f$.
\end{enumerate}
\end{question}

\begin{answer}
\textbf{(d) is the false one.}

\textbf{(a) is true,} since $1 > 0$ is clearly an eigenvalue of $\Hess f(0)$ for any choice of
$\alpha$ (the third standard basis vector is an eigenvector), so $\Hess f(0)$ is never negative
semidefinite and $0$ can never be a local maximum.

\textbf{(b) is true.} Let $\lambda_1,\lambda_2,\lambda_3 = 1$ be the eigenvalues of
$\Hess f(0)$. Then
\[\lambda_1\lambda_2 = \det\big(\Hess f(0)\big) = 2-\alpha^2 < 0 \quad \text{for } \alpha > \sqrt2,\]
so $\Hess f(0)$ has both positive and negative eigenvalues, making it indefinite, and $0$ is a
saddle point by \cref{thm:hessian_test}.

\textbf{(c) is true.} For these values $\det(\Hess f(0)) \neq 0$, and therefore the inverse
function theorem (\cref{thm:inverse_function_theorem}) applies. Note that
$\Hess f(0) = D(\nabla f)_0$: the Hessian of $f$ \emph{is} the differential of the gradient map.

\textbf{(d) is false,} because the same reasoning as in \textbf{(b)} gives that for
$\alpha < -\sqrt2 < \sqrt2$ we again have $2-\alpha^2 < 0$, so $0$ is a saddle point of $f$, not
a local minimum. The statement is true only on $-\sqrt2 < \alpha < \sqrt2$.
\end{answer}

\begin{ainote}
Corsin's third bullet reads \qt{for $\alpha \neq \sqrt2$}, but $\det(\Hess f(0)) = 2-\alpha^2$
vanishes at $\alpha = -\sqrt2$ as well, so it must read $\alpha \neq \pm\sqrt2$; typeset with
$\pm$ above. The oversight is a conspicuous one here, since the \emph{next} bullet is false for
exactly the reason the third one omits. The quiz is built around the missing negative root.

\end{ainote}

% Source: Corsin Nick/Class Notes/Week 8.pdf, p. 5
\begin{question}
\label{q:quiz_q6}
Let $V := \{(x,y) \in \mathbb{R}^2 : y^4+y^2 = x^3+x\}$. Decide:
\begin{enumerate}[label=\textbf{(\alph*)}]
  \item $V$ is a graph over $y$ in a neighbourhood of $(0,0)$.
  \item $V$ is a graph over $x$ in a neighbourhood of $(0,0)$.
\end{enumerate}
\end{question}

\begin{answer}
\textbf{(a) True.} Notice that $V = f^{-1}\{0\}$ for $f(x,y) := y^4+y^2-x^3-x$. Then
\[\Jac f(x,y) = (\partial_xf,\ \partial_yf) = (-3x^2-1,\ 4y^3+2y), \qquad \Jac f(0,0) = (-1,\ 0).\]
Since $\partial_xf(0,0) = -1 \neq 0$, the implicit function theorem
(\cref{thm:implicit_function_theorem}) applies with $x$ as the variable to solve for: in a
neighbourhood of $(0,0)$,
\[(x,y) \in V \iff f(x,y) = 0 \iff x = x(y).\]
The theorem does \textbf{not} apply for $y$, since $\partial_yf(0,0) = 0$.

\textbf{(b) False.} Solving the defining equation for $y^2$ by the quadratic formula gives
\[y^2 = -\tfrac12 + \sqrt{\tfrac14 + x^3 + x},\]
which has no solution for $x \in (-\varepsilon,0)$: there $x^3+x < 0$, so the square root is
smaller than $\tfrac12$ and the right-hand side is negative, while $y^2 \geq 0$. So $V$ contains
no points at all just to the left of the origin, and cannot be a graph over $x$ there.
\end{answer}

\begin{ainote}
The two parts are settled by different means, and the asymmetry is the point. Part \textbf{(a)}
is decided by \emph{checking a derivative}; part \textbf{(b)} needs an \emph{explicit
computation} showing the set is empty on one side. That is because the implicit function theorem
yields no negative results: failure of its hypothesis proves nothing on its own, exactly as
discussed after \cref{thm:implicit_function_theorem}.
\end{ainote}

% Source: Corsin Nick/Class Notes/Week 8.pdf, p. 6
\begin{question}
\label{q:quiz_q7}
Let $\gamma : I \to \mathbb{R}^2$ be injective and smooth, with $I$ an open interval. Is
$\gamma(I)$ then a smooth submanifold of $\mathbb{R}^2$?
\end{question}

\begin{answer}
\textbf{False.} Consider
\[\gamma : (0,2\pi) \to \mathbb{R}^2, \qquad t \mapsto \big(\sin t,\ \sin 2t\big),\]
whose image is a \newterm{figure eight} \germanfor{Lemniskate}. This $\gamma$ is injective and
smooth, but there is no graphical representation of $\operatorname{im}(\gamma)$ around $(0,0)$.
\end{answer}

% Sonnet 5 (Medium) --- FIG-W08-01, redrawn from the plot on the source page.
\begin{center}
\begin{tikzpicture}[>=Latex, scale=1.45]
  \draw[->] (-1.35,0) -- (1.35,0) node[right, font=\scriptsize] {$x$};
  \draw[->] (0,-1.25) -- (0,1.25) node[above, font=\scriptsize] {$y$};
  \draw[blue!80!black, thick, variable=\t, domain=0:6.2832, samples=200, smooth]
    plot ({sin(\t r)}, {sin(2*\t r)});
  \draw[red!80!black, thick, dashed] (0,0) circle (0.33);
  \node[red!80!black, font=\scriptsize, right] at (0.52,0.34) {problem!};
  \draw[red!80!black, thin, ->] (0.50,0.30) -- (0.20,0.13);
\end{tikzpicture}
\end{center}

\begin{ainote}
Two things are worth extracting. First, how injectivity survives a self-crossing at all: the two
parameter ends $t \to 0^+$ and $t \to 2\pi^-$ both send $\gamma(t) \to (0,0)$, but neither
\emph{attains} it, so the image passes through the origin only once (via $t = \pi$) while
being approached from two further directions.

Second, this locates precisely which hypothesis of \cref{thm:local_parametrization} fails. It is
not injectivity, and not injectivity of $D\gamma_t$; it is the requirement that the
parametrization be a \textbf{homeomorphism} onto its image. Here $\gamma^{-1}$ is discontinuous
at $(0,0)$: points of the image arbitrarily close to the origin have parameters near $0$ and near
$2\pi$, far apart in $I$. That the resulting set is not a submanifold can then be confirmed by
the component count of \cref{ex:cross_not_submanifold}. Deleting the crossing point leaves
four pieces, while deleting a point from an interval leaves two.
\end{ainote}

% Sonnet 5 (Medium)
The quiz ends here, and with it the differential-topology half of the course. The rest of the
week starts integration in several variables.

% ============================================================
% Appendix --- Problem sheet index
%
% This appendix exists because the restructure dissolved the official problem sheets: each
% problem now sits with the material it tests, so a sheet's problems can be spread over
% several chapters. The tables below put the sheet view back, and they are also where
% Corsin's colour-coded "Recommended exercises" pages are preserved --- those were
% per-sheet artefacts with nowhere else to live once the \exercisesheet{N} headings went.
%
% Built up during the migration, one sheet per stage.
% ============================================================
\chapter{Problem Sheet Index}
\label{ch:sheet_index}

% Transition: Claude Opus 5 (Medium)
The problem sheets are not reproduced as blocks anywhere in this document. Every problem is
quoted verbatim from the official sheet and placed next to the material it tests, with its
solution in the \emph{Solutions} section of the chapter that hosts it. That is the right
arrangement for reading the notes and the wrong one for answering \qt{what is on sheet 7?}, so
this appendix restores the sheet view.

All statements are quoted from \texttt{exercises/ExN\_Analysis2\_eng.pdf}; attribution:
Prof.\ Joaquim Serra, D-MATH, ETH Z\"urich.

\begin{notation}[Corsin's colour key]
\label{not:corsin_priority_key}
From sheet 4 onwards Corsin opens each file with a colour-coded \textit{Recommended exercises}
page. The key, as given there: blue \textbf{= important}, orange \textbf{= semi-important},
red \textbf{= optional}. The official sheets additionally mark harder problems with \textbf{(*)}.
Each problem carries its tag in its own title throughout this document; the tables below collect
them, together with Corsin's own comments.
\end{notation}

\begin{ainote}
Sheets 1--3 have no priority page --- Corsin's \textit{Recommended exercises} pages start with
sheet 4 --- so no problem on them is singled out, and all of them are reproduced and solved.
\end{ainote}

\section{Sheets 1--3}
\label{sec:sheet_index_1_3}

\begin{center}
\begin{tabular}{@{}l l l@{}}
\textbf{Problem} & \textbf{Topic} & \textbf{Now in} \\
\hline
1.1 & Recap questions          & \cref{sec:analysis1_recap} \\
1.2 & From equations to drawings & \cref{sec:visualizing_sets} \\
1.3 & From drawings to equations & \cref{sec:visualizing_sets} \\
1.4 & Young's inequality       & \cref{sec:youngs_inequality} \\[4pt]
2.1 & Metric spaces            & \cref{sec:metric_spaces} \\
2.2 & Multiple choice (closure) & \cref{sec:unit_balls} \\
2.3 & Multiple choice (distance from a set) & \cref{sec:unit_balls} \\
2.4 & Boundary, interior       & \cref{sec:unit_balls} \\
2.5 & Product of metric spaces & \cref{sec:metric_spaces} \\
2.6 & Continuity of the distance & \cref{sec:continuity} \\
2.7 & Continuity of the composition & \cref{sec:continuity} \\[4pt]
3.1 & Lipschitz or not         & \cref{sec:lipschitz_continuity} \\
3.2 & Problems at the origin   & \cref{sec:continuity} \\
3.3 & Open, closed, complete and compact & \cref{sec:why_compactness_matters} \\
3.4 & Multiple choice (open, complete, compact) & \cref{sec:why_compactness_matters} \\
3.5 & Multiple choice (fixed points) & \cref{sec:banach_fixed_point} \\
3.6 & Multiple choice (uniform continuity) & \cref{sec:lipschitz_continuity} \\
3.7 & Cauchy--Schwarz          & \cref{sec:cauchy_schwarz} \\
3.8 & Inner product and norm   & \cref{sec:cauchy_schwarz} \\
3.9 & All norms are equivalent & \cref{sec:cauchy_schwarz} \\
3.10 & Hilbert--Schmidt norm   & \cref{sec:cauchy_schwarz} \\
3.11 & (*) Uniform continuity and moduli & \cref{sec:lipschitz_continuity} \\
\end{tabular}
\end{center}

\section{Sheet 4}
\label{sec:sheet_index_4}

% Source: Corsin Nick/Class Notes/Week 4.pdf, p. 1
% Corsin's "Recommended exercises" page, preserved here: it was a per-sheet artefact and had
% nowhere else to go once the \exercisesheet{4} heading was dissolved.
\begin{center}
\begin{tabular}{@{}l p{0.22\linewidth} p{0.34\linewidth} l@{}}
\textbf{Problem} & \textbf{Priority} & \textbf{Corsin's note} & \textbf{Now in} \\
\hline
4.1 & important & ``Use the chain rule'' & \cref{sec:jacobi_matrix} \\[3pt]
4.2 & important & ``Recall that path connected implies connected'' & \cref{sec:path_connectedness} \\[3pt]
4.3 & semi-important --- parts 1, 2, 3, 5; \textbf{parts 4 and 6 important} &
  ``For 3., use Cauchy--Schwarz in $\mathbb{R}^n$ applied to convenient vectors.'' &
  \cref{sec:cauchy_schwarz} \\[3pt]
4.4 & optional & --- & \cref{sec:cauchy_schwarz} \\[3pt]
4.5 & optional & --- & \cref{sec:mean_value_theorem} \\[3pt]
4.6 & important --- parts 1, 2, 3 & --- & \cref{sec:chain_rule_along_path} \\
\end{tabular}
\end{center}

\begin{ainote}
Problems \textbf{4.4} and \textbf{4.5} appear in this document for the first time. They had been
transcribed into \texttt{content/exercise-sheets/week-04.tex}, but that file was never
\texttt{\textbackslash input} by anything, so neither problem had ever reached the built PDF and
neither carried a solution. Both are now placed and solved.
\end{ainote}


\end{document}
